%% Generated by Sphinx.
\def\sphinxdocclass{jupyterBook}
\documentclass[letterpaper,10pt,english]{jupyterBook}
\ifdefined\pdfpxdimen
   \let\sphinxpxdimen\pdfpxdimen\else\newdimen\sphinxpxdimen
\fi \sphinxpxdimen=.75bp\relax
\ifdefined\pdfimageresolution
    \pdfimageresolution= \numexpr \dimexpr1in\relax/\sphinxpxdimen\relax
\fi
%% let collapsible pdf bookmarks panel have high depth per default
\PassOptionsToPackage{bookmarksdepth=5}{hyperref}
%% turn off hyperref patch of \index as sphinx.xdy xindy module takes care of
%% suitable \hyperpage mark-up, working around hyperref-xindy incompatibility
\PassOptionsToPackage{hyperindex=false}{hyperref}
%% memoir class requires extra handling
\makeatletter\@ifclassloaded{memoir}
{\ifdefined\memhyperindexfalse\memhyperindexfalse\fi}{}\makeatother

\PassOptionsToPackage{booktabs}{sphinx}
\PassOptionsToPackage{colorrows}{sphinx}

\PassOptionsToPackage{warn}{textcomp}

\catcode`^^^^00a0\active\protected\def^^^^00a0{\leavevmode\nobreak\ }
\usepackage{cmap}
\usepackage{fontspec}
\defaultfontfeatures[\rmfamily,\sffamily,\ttfamily]{}
\usepackage{amsmath,amssymb,amstext}
\usepackage{polyglossia}
\setmainlanguage{english}



\setmainfont{FreeSerif}[
  Extension      = .otf,
  UprightFont    = *,
  ItalicFont     = *Italic,
  BoldFont       = *Bold,
  BoldItalicFont = *BoldItalic
]
\setsansfont{FreeSans}[
  Extension      = .otf,
  UprightFont    = *,
  ItalicFont     = *Oblique,
  BoldFont       = *Bold,
  BoldItalicFont = *BoldOblique,
]
\setmonofont{FreeMono}[
  Extension      = .otf,
  UprightFont    = *,
  ItalicFont     = *Oblique,
  BoldFont       = *Bold,
  BoldItalicFont = *BoldOblique,
]



\usepackage[Bjarne]{fncychap}
\usepackage[,numfigreset=1,mathnumfig]{sphinx}

\fvset{fontsize=\small}
\usepackage{geometry}


% Include hyperref last.
\usepackage{hyperref}
% Fix anchor placement for figures with captions.
\usepackage{hypcap}% it must be loaded after hyperref.
% Set up styles of URL: it should be placed after hyperref.
\urlstyle{same}


\usepackage{sphinxmessages}
\setcounter{tocdepth}{2}


        % Start of preamble defined in sphinx-jupyterbook-latex %
         \usepackage[Latin,Greek]{ucharclasses}
        \usepackage{unicode-math}
        % fixing title of the toc
        \addto\captionsenglish{\renewcommand{\contentsname}{Contents}}
        \hypersetup{
            pdfencoding=auto,
            psdextra
        }
        % End of preamble defined in sphinx-jupyterbook-latex %
        

\title{Statistical Foundations of Machine Learning - Practicals handbook}
\date{Sep 16, 2025}
\release{}
\author{Tribel Pascal, Simar Cédric, Bontempi Gianluca}
\newcommand{\sphinxlogo}{\vbox{}}
\renewcommand{\releasename}{}
\makeindex
\begin{document}

\pagestyle{empty}
\sphinxmaketitle
\pagestyle{plain}
\sphinxtableofcontents
\pagestyle{normal}
\phantomsection\label{\detokenize{frontpage::doc}}


\sphinxAtStartPar
Welcome. In the following pages, you will find the practicals reference for the course of \sphinxhref{https://www.ulb.be/en/programme/2025-info-f422}{INFO\sphinxhyphen{}F422 \sphinxhyphen{} Statistical Foundations of Machine Learning} given at the \sphinxhref{http://www.ulb.be}{Université Libre de Bruxelles}.

\begin{DUlineblock}{0em}
\item[] \sphinxstylestrong{\Large Teaching staff}
\end{DUlineblock}

\sphinxAtStartPar
The 2024\sphinxhyphen{}2025 teaching staff is made of:
\begin{itemize}
\item {} 
\sphinxAtStartPar
\sphinxhref{mailto:gianluca.bontempi@ulb.be}{✉️ Gianluca Bontempi}: Teacher

\item {} 
\sphinxAtStartPar
\sphinxhref{mailto:cedric.simar@ulb.be}{✉️ Cédric Simar}: Teaching Assistant

\item {} 
\sphinxAtStartPar
\sphinxhref{mailto:pascal.tribel@ulb.be}{✉️ Pascal Tribel}: Teaching Assistant

\item {} 
\sphinxAtStartPar
\sphinxhref{mailto:gian.marco.paldino@ulb.be}{✉️ Gian Marco Paldino}: PhD Student

\end{itemize}

\begin{DUlineblock}{0em}
\item[] \sphinxstylestrong{\large Reference}
\end{DUlineblock}
\begin{itemize}
\item {} 
\sphinxAtStartPar
\sphinxhref{https://www.researchgate.net/publication/242692234\_Statistical\_foundations\_of\_machine\_learning\_the\_handbook}{📜 Statistical Foundations of Machine Learning: the handbook} by \sphinxstylestrong{Gianluca Bontempi}, Machine Learning Group, Computer Science Department, ULB

\end{itemize}
\begin{itemize}
\item {} 
\sphinxAtStartPar
{\hyperref[\detokenize{introduction::doc}]{\sphinxcrossref{Introduction}}}

\item {} 
\sphinxAtStartPar
{\hyperref[\detokenize{01::doc}]{\sphinxcrossref{Introduction to Probabilistic Methods and Monte Carlo Simulations}}}

\item {} 
\sphinxAtStartPar
{\hyperref[\detokenize{02::doc}]{\sphinxcrossref{Linear Models}}}

\item {} 
\sphinxAtStartPar
{\hyperref[\detokenize{03::doc}]{\sphinxcrossref{Data Preprocessing, Local Models and Tree\sphinxhyphen{}Based Models}}}

\item {} 
\sphinxAtStartPar
{\hyperref[\detokenize{04::doc}]{\sphinxcrossref{Neural Networks for Regression}}}

\item {} 
\sphinxAtStartPar
{\hyperref[\detokenize{05::doc}]{\sphinxcrossref{Ensembles of models and feature selection}}}

\item {} 
\sphinxAtStartPar
{\hyperref[\detokenize{06::doc}]{\sphinxcrossref{Classification}}}

\item {} 
\sphinxAtStartPar
{\hyperref[\detokenize{conclusion::doc}]{\sphinxcrossref{Conclusion}}}

\item {} 
\sphinxAtStartPar
{\hyperref[\detokenize{bibliography::doc}]{\sphinxcrossref{Bibliography}}}

\end{itemize}

\begin{DUlineblock}{0em}
\item[] \sphinxstylestrong{\large Mistakes, Remarks, or Compliments}
\end{DUlineblock}

\sphinxAtStartPar
For anything you would like to share to the team, feel free to write an e\sphinxhyphen{}mail to \sphinxhref{mailto:pascal.tribel@ulb.be}{Pascal Tribel}.

\sphinxstepscope


\chapter{Introduction}
\label{\detokenize{introduction:introduction}}\label{\detokenize{introduction::doc}}

\section{Welcome}
\label{\detokenize{introduction:welcome}}
\sphinxAtStartPar
Welcome to the \sphinxstyleemphasis{Statistical Foundations of Machine Learning} course, part of the Master’s program in Computer Science at ULB. This course provides the essential mathematical and algorithmic foundations behind the modern techniques used in machine learning, with a strong focus on statistics, probabilistic modeling, and regression/classification methods. It is the perfect starting point for anyone aiming to understand, build, and evaluate predictive models rigorously.


\bigskip\hrule\bigskip



\section{Why This Course?}
\label{\detokenize{introduction:why-this-course}}
\sphinxAtStartPar
Machine learning is not magic. It is grounded in solid mathematical ideas, especially probability theory, statistics, linear algebra, and optimization. Understanding \sphinxstylestrong{how and why} algorithms work is crucial for designing good systems that are not only performant but also interpretable and reliable.

\sphinxAtStartPar
This course will:
\begin{itemize}
\item {} 
\sphinxAtStartPar
Build your \sphinxstylestrong{intuition} for probabilistic reasoning and statistical inference.

\item {} 
\sphinxAtStartPar
Introduce \sphinxstylestrong{core regression and classification models}, from linear regression to neural networks.

\item {} 
\sphinxAtStartPar
Teach you how to \sphinxstylestrong{implement algorithms}, visualize results, and understand performance metrics.

\item {} 
\sphinxAtStartPar
Provide \sphinxstylestrong{practical tools and code exercises} using Python and popular libraries such as NumPy, Matplotlib, scikit\sphinxhyphen{}learn, and PyTorch.

\item {} 
\sphinxAtStartPar
Develop your ability to \sphinxstylestrong{choose, analyze, and compare models}, considering trade\sphinxhyphen{}offs like bias/variance, overfitting, and interpretability.

\end{itemize}

\sphinxAtStartPar
Whether your goal is to become a data scientist, ML engineer, or researcher, these foundations will serve as the most important base for more advanced courses and projects.


\bigskip\hrule\bigskip



\section{Course Overview}
\label{\detokenize{introduction:course-overview}}
\sphinxAtStartPar
Here’s what you can expect:


\subsection{Introduction to Probabilistic Methods and Monte Carlo Simulations}
\label{\detokenize{introduction:introduction-to-probabilistic-methods-and-monte-carlo-simulations}}
\sphinxAtStartPar
We start with the building blocks of randomness and probability:
\begin{itemize}
\item {} 
\sphinxAtStartPar
Simulating data with \sphinxcode{\sphinxupquote{NumPy}}

\item {} 
\sphinxAtStartPar
Visualizing distributions with \sphinxcode{\sphinxupquote{Matplotlib}}

\item {} 
\sphinxAtStartPar
Understanding Monte Carlo simulations

\item {} 
\sphinxAtStartPar
Working with \sphinxstylestrong{multivariate Gaussians}

\item {} 
\sphinxAtStartPar
Learning about \sphinxstylestrong{minimization} in statistical modeling

\item {} 
\sphinxAtStartPar
Exploring \sphinxstylestrong{Gaussian Mixture Models}

\end{itemize}


\bigskip\hrule\bigskip



\subsection{Linear Models}
\label{\detokenize{introduction:linear-models}}
\sphinxAtStartPar
From classical regression to more flexible transformations:
\begin{itemize}
\item {} 
\sphinxAtStartPar
Ordinary Least Squares (OLS) regression

\item {} 
\sphinxAtStartPar
Analyzing residuals and model assumptions

\item {} 
\sphinxAtStartPar
Polynomial and \sphinxstylestrong{Radial Basis Function (RBF)} transformations

\item {} 
\sphinxAtStartPar
Ridge regression and the role of regularization

\end{itemize}


\bigskip\hrule\bigskip



\subsection{Preprocessing, Local Models and Tree\sphinxhyphen{}Based Models}
\label{\detokenize{introduction:preprocessing-local-models-and-tree-based-models}}
\sphinxAtStartPar
Preparing your data and working with interpretable machine learning models:
\begin{itemize}
\item {} 
\sphinxAtStartPar
Data preprocessing: normalization, handling missing data

\item {} 
\sphinxAtStartPar
\sphinxstylestrong{k\sphinxhyphen{}Nearest Neighbors (kNN)} for regression

\item {} 
\sphinxAtStartPar
\sphinxstylestrong{Decision trees} and \sphinxstylestrong{random forests}

\item {} 
\sphinxAtStartPar
Using \sphinxstylestrong{scikit\sphinxhyphen{}learn pipelines}

\item {} 
\sphinxAtStartPar
Understanding the \sphinxstylestrong{bias\sphinxhyphen{}variance tradeoff} in practice

\end{itemize}


\bigskip\hrule\bigskip



\subsection{Neural Networks for Regression}
\label{\detokenize{introduction:neural-networks-for-regression}}
\sphinxAtStartPar
Deep learning in its most intuitive form:
\begin{itemize}
\item {} 
\sphinxAtStartPar
Normalization and standardization for stable learning

\item {} 
\sphinxAtStartPar
MLPs (Multilayer Perceptrons) for regression tasks

\item {} 
\sphinxAtStartPar
Cost functions and model training

\item {} 
\sphinxAtStartPar
\sphinxstylestrong{Matrix\sphinxhyphen{}based backpropagation}

\item {} 
\sphinxAtStartPar
Introduction to \sphinxstylestrong{CNNs}

\item {} 
\sphinxAtStartPar
Dealing with \sphinxstylestrong{overfitting}, \sphinxstylestrong{cross\sphinxhyphen{}validation}

\item {} 
\sphinxAtStartPar
Model persistence: saving and loading in \sphinxstylestrong{PyTorch}

\end{itemize}


\bigskip\hrule\bigskip



\subsection{Ensemble Models and Feature Selection}
\label{\detokenize{introduction:ensemble-models-and-feature-selection}}
\sphinxAtStartPar
Combining models and selecting the most useful inputs:
\begin{itemize}
\item {} 
\sphinxAtStartPar
Error decomposition and performance metrics

\item {} 
\sphinxAtStartPar
\sphinxstylestrong{Filter, wrapper, and embedded} feature selection

\item {} 
\sphinxAtStartPar
Building \sphinxstylestrong{ensembles} of models for improved accuracy

\end{itemize}


\bigskip\hrule\bigskip



\subsection{Classification}
\label{\detokenize{introduction:classification}}
\sphinxAtStartPar
From regression to classification, we shift the focus:
\begin{itemize}
\item {} 
\sphinxAtStartPar
Evaluation metrics for classification (accuracy, precision, recall, F1)

\item {} 
\sphinxAtStartPar
\sphinxstylestrong{Bayes classifier} and probabilistic foundations

\item {} 
\sphinxAtStartPar
\sphinxstylestrong{Logistic regression} and ROC curves

\item {} 
\sphinxAtStartPar
Working with \sphinxstylestrong{unequal costs}

\item {} 
\sphinxAtStartPar
Alternative classification models

\end{itemize}


\bigskip\hrule\bigskip



\section{Practical Exercises}
\label{\detokenize{introduction:practical-exercises}}
\sphinxAtStartPar
This handbook serves as the practicals reference for the course. The theoretical handbook remains the main reference. We want to focus here in the implementations, the algorithms and the existing libraries for implementing the topics covered in the course. The student should refer as much as possible to the theoretical handbook before delving into the coding. Here, we try, as much as possible, to include:
\begin{itemize}
\item {} 
\sphinxAtStartPar
\sphinxstylestrong{Hands\sphinxhyphen{}on coding exercises}

\item {} 
\sphinxAtStartPar
\sphinxstylestrong{Mini\sphinxhyphen{}projects and simulations}

\item {} 
\sphinxAtStartPar
\sphinxstylestrong{Graphical visualizations} of models and results

\item {} 
\sphinxAtStartPar
Opportunities to \sphinxstylestrong{experiment} with real\sphinxhyphen{}world datasets

\end{itemize}

\sphinxAtStartPar
Doing the exercises is \sphinxstylestrong{not optional} if you want to truly master the material. They are designed to:
\begin{itemize}
\item {} 
\sphinxAtStartPar
Reinforce theoretical concepts

\item {} 
\sphinxAtStartPar
Develop your implementation skills

\item {} 
\sphinxAtStartPar
Sharpen your critical thinking about model performance and assumptions

\end{itemize}


\bigskip\hrule\bigskip



\section{Final Notes}
\label{\detokenize{introduction:final-notes}}
\sphinxAtStartPar
This is a \sphinxstylestrong{theory\sphinxhyphen{}meets\sphinxhyphen{}practice} course. You’ll build up solid statistical thinking and programming fluency, while also learning to critically assess and apply machine learning models. Every part of the course prepares you for future work in AI, Data Science, or scientific computing.

\sphinxAtStartPar
\sphinxstylestrong{Take it seriously. Be curious. Challenge your understanding. And above all: have fun exploring the possibilities of statistical thinking in machine learning.}

\sphinxAtStartPar
Let’s begin!

\sphinxstepscope


\chapter{Introduction to Probabilistic Methods and Monte Carlo Simulations}
\label{\detokenize{01:introduction-to-probabilistic-methods-and-monte-carlo-simulations}}\label{\detokenize{01::doc}}
\sphinxAtStartPar
Probabilistic methods and Monte Carlo simulations represent fundamental techniques in computational mathematics, offering powerful approaches to solve complex problems that would be analytically intractable. These methods have found widespread applications across diverse domains including statistical physics, financial mathematics, quantum mechanics, machine learning, and computational biology. The course of INFO\sphinxhyphen{}F305 \sphinxhyphen{} \sphinxstyleemphasis{Modélisation et Simulation} proposes a good introduction to Monte Carlo simulations.

\sphinxAtStartPar
This chapter introduces the fundamental concepts, tools, and techniques of probabilistic methods with a particular focus on Monte Carlo simulations. We will explore how to generate random variables from different distributions, estimate integrals, solve optimization problems, and analyze statistical properties of random processes. Throughout the chapter, we will emphasize both the theoretical foundations and practical implementations using Python and its scientific computing ecosystem.

\sphinxAtStartPar
Note, the proper understanding of this set of methods is extremely important for having a good approach of many modern Machine Learning problems. Real\sphinxhyphen{}life data collection can indeed, in many cases, be modelized by a Monte Carlo simulation, and the statistical/probabilistic study of the simulated problem often offers valuable insights about the real\sphinxhyphen{}life dynamics.

\sphinxAtStartPar
By the end of this chapter, you will have developed a solid understanding of how to apply Monte Carlo techniques to solve a variety of mathematical and statistical problems, setting the foundation for more advanced applications in subsequent chapters.


\section{Necessary Libraries}
\label{\detokenize{01:necessary-libraries}}
\sphinxAtStartPar
To run this practical work, the following libraries must be installed:
\begin{itemize}
\item {} 
\sphinxAtStartPar
\sphinxcode{\sphinxupquote{NumPy}}: Vector and matrix manipulation, random number generation,

\item {} 
\sphinxAtStartPar
\sphinxcode{\sphinxupquote{SciPy}}: For advanced numerical statistical methods,

\item {} 
\sphinxAtStartPar
\sphinxcode{\sphinxupquote{Pandas}}: For advanced data manipulation

\item {} 
\sphinxAtStartPar
\sphinxcode{\sphinxupquote{Matplotlib}}, \sphinxcode{\sphinxupquote{seaborn}}: For data visualization,

\item {} 
\sphinxAtStartPar
\sphinxcode{\sphinxupquote{tqdm}}: For beautiful loading bars.

\end{itemize}

\begin{sphinxShadowBox}
\sphinxstylesidebartitle{}

\sphinxAtStartPar
If the libraries are already installed, you can skip this step.
\end{sphinxShadowBox}

\sphinxAtStartPar
If you are using Jupyter Notebook, you can install the dependencies directly by running the cell below:

\begin{sphinxuseclass}{cell}\begin{sphinxVerbatimInput}

\begin{sphinxuseclass}{cell_input}
\begin{sphinxVerbatim}[commandchars=\\\{\}]
\PYG{c+ch}{\PYGZsh{}!pip install numpy matplotlib tqdm pandas seaborn  numpyarray\PYGZus{}to\PYGZus{}latex}
\end{sphinxVerbatim}

\end{sphinxuseclass}\end{sphinxVerbatimInput}

\end{sphinxuseclass}
\begin{sphinxShadowBox}
\sphinxstylesidebartitle{}

\sphinxAtStartPar
This \sphinxcode{\sphinxupquote{pprint}} function will allow to render the output of some codes in LaTeX, such as the vectors and the matrices. Remark that calling this function only works in \sphinxcode{\sphinxupquote{Jupyter}}. If you code in plain \sphinxcode{\sphinxupquote{Python}}, you should rather use the default \sphinxcode{\sphinxupquote{str}} and \sphinxcode{\sphinxupquote{print}} functions.
\end{sphinxShadowBox}

\sphinxAtStartPar
We make some imports and define a \sphinxcode{\sphinxupquote{pprint}} function.

\begin{sphinxuseclass}{cell}\begin{sphinxVerbatimInput}

\begin{sphinxuseclass}{cell_input}
\begin{sphinxVerbatim}[commandchars=\\\{\}]
\PYG{k+kn}{import}\PYG{+w}{ }\PYG{n+nn}{numpy}\PYG{+w}{ }\PYG{k}{as}\PYG{+w}{ }\PYG{n+nn}{np}
\PYG{k+kn}{from}\PYG{+w}{ }\PYG{n+nn}{IPython}\PYG{n+nn}{.}\PYG{n+nn}{display}\PYG{+w}{ }\PYG{k+kn}{import} \PYG{n}{display}\PYG{p}{,} \PYG{n}{Math}
\PYG{k+kn}{from}\PYG{+w}{ }\PYG{n+nn}{numpyarray\PYGZus{}to\PYGZus{}latex}\PYG{n+nn}{.}\PYG{n+nn}{jupyter}\PYG{+w}{ }\PYG{k+kn}{import} \PYG{n}{to\PYGZus{}ltx}
\PYG{k}{def}\PYG{+w}{ }\PYG{n+nf}{pprint}\PYG{p}{(}\PYG{o}{*}\PYG{n}{args}\PYG{p}{)}\PYG{p}{:}
    \PYG{n}{res} \PYG{o}{=} \PYG{l+s+s2}{\PYGZdq{}}\PYG{l+s+s2}{\PYGZdq{}}
    \PYG{k}{for} \PYG{n}{i} \PYG{o+ow}{in} \PYG{n}{args}\PYG{p}{:}
        \PYG{k}{if} \PYG{n+nb}{type}\PYG{p}{(}\PYG{n}{i}\PYG{p}{)} \PYG{o}{==} \PYG{n}{np}\PYG{o}{.}\PYG{n}{ndarray}\PYG{p}{:}
            \PYG{n}{res} \PYG{o}{+}\PYG{o}{=} \PYG{n}{to\PYGZus{}ltx}\PYG{p}{(}\PYG{n}{i}\PYG{p}{,} \PYG{n}{brackets}\PYG{o}{=}\PYG{l+s+s1}{\PYGZsq{}}\PYG{l+s+s1}{[]}\PYG{l+s+s1}{\PYGZsq{}}\PYG{p}{)}
        \PYG{k}{elif} \PYG{n+nb}{type}\PYG{p}{(}\PYG{n}{i}\PYG{p}{)} \PYG{o}{==} \PYG{n+nb}{str}\PYG{p}{:}
            \PYG{n}{res} \PYG{o}{+}\PYG{o}{=} \PYG{n}{i}
    \PYG{n}{display}\PYG{p}{(}\PYG{n}{Math}\PYG{p}{(}\PYG{n}{res}\PYG{p}{)}\PYG{p}{)}
\end{sphinxVerbatim}

\end{sphinxuseclass}\end{sphinxVerbatimInput}

\end{sphinxuseclass}

\section{Introduction to \sphinxstyleliteralintitle{\sphinxupquote{NumPy}}}
\label{\detokenize{01:introduction-to-numpy}}
\sphinxAtStartPar
\sphinxcode{\sphinxupquote{NumPy}} is an essential library for scientific computations in Python. It offers:
\begin{itemize}
\item {} 
\sphinxAtStartPar
\sphinxcode{\sphinxupquote{ndarray}} objects to represent multi\sphinxhyphen{}dimensional arrays.

\item {} 
\sphinxAtStartPar
A collection of mathematical functions to perform fast operations on these arrays.

\end{itemize}


\subsection{Manipulation of \sphinxstyleliteralintitle{\sphinxupquote{NumPy}} Arrays}
\label{\detokenize{01:manipulation-of-numpy-arrays}}
\sphinxAtStartPar
Here are some examples for creating and manipulating \sphinxcode{\sphinxupquote{NumPy}} arrays:

\begin{sphinxuseclass}{cell}\begin{sphinxVerbatimInput}

\begin{sphinxuseclass}{cell_input}
\begin{sphinxVerbatim}[commandchars=\\\{\}]
\PYG{n}{v0} \PYG{o}{=} \PYG{n}{np}\PYG{o}{.}\PYG{n}{array}\PYG{p}{(}\PYG{p}{[}\PYG{l+m+mi}{1}\PYG{p}{,} \PYG{l+m+mi}{2}\PYG{p}{,} \PYG{l+m+mi}{3}\PYG{p}{]}\PYG{p}{)}
\PYG{n}{v1} \PYG{o}{=} \PYG{n}{np}\PYG{o}{.}\PYG{n}{arange}\PYG{p}{(}\PYG{l+m+mi}{0}\PYG{p}{,} \PYG{l+m+mi}{3}\PYG{p}{,} \PYG{l+m+mf}{0.4}\PYG{p}{)}
\PYG{n}{v2} \PYG{o}{=} \PYG{n}{np}\PYG{o}{.}\PYG{n}{linspace}\PYG{p}{(}\PYG{l+m+mi}{0}\PYG{p}{,} \PYG{l+m+mi}{5}\PYG{p}{,} \PYG{l+m+mi}{10}\PYG{p}{)}
\PYG{n}{v3} \PYG{o}{=} \PYG{n}{np}\PYG{o}{.}\PYG{n}{linspace}\PYG{p}{(}\PYG{l+m+mi}{0}\PYG{p}{,} \PYG{l+m+mi}{2}\PYG{p}{,} \PYG{l+m+mi}{10}\PYG{p}{)}
\PYG{n}{M} \PYG{o}{=} \PYG{n}{np}\PYG{o}{.}\PYG{n}{array}\PYG{p}{(}\PYG{p}{[}\PYG{p}{[}\PYG{l+m+mi}{1}\PYG{p}{,} \PYG{l+m+mi}{2}\PYG{p}{,} \PYG{l+m+mi}{3}\PYG{p}{]}\PYG{p}{,} \PYG{p}{[}\PYG{l+m+mi}{4}\PYG{p}{,} \PYG{l+m+mi}{5}\PYG{p}{,} \PYG{l+m+mi}{6}\PYG{p}{]}\PYG{p}{,} \PYG{p}{[}\PYG{l+m+mi}{7}\PYG{p}{,} \PYG{l+m+mi}{8}\PYG{p}{,} \PYG{l+m+mi}{9}\PYG{p}{]}\PYG{p}{]}\PYG{p}{)}
\PYG{n}{pprint}\PYG{p}{(}\PYG{l+s+s2}{\PYGZdq{}}\PYG{l+s+s2}{v\PYGZus{}0=}\PYG{l+s+s2}{\PYGZdq{}}\PYG{p}{,} \PYG{n}{v0}\PYG{p}{)}
\PYG{n}{pprint}\PYG{p}{(}\PYG{l+s+s2}{\PYGZdq{}}\PYG{l+s+s2}{v\PYGZus{}1=}\PYG{l+s+s2}{\PYGZdq{}}\PYG{p}{,} \PYG{n}{v1}\PYG{p}{)}
\PYG{n}{pprint}\PYG{p}{(}\PYG{l+s+s2}{\PYGZdq{}}\PYG{l+s+s2}{v\PYGZus{}2=}\PYG{l+s+s2}{\PYGZdq{}}\PYG{p}{,} \PYG{n}{v2}\PYG{p}{)}
\PYG{n}{pprint}\PYG{p}{(}\PYG{l+s+s2}{\PYGZdq{}}\PYG{l+s+s2}{v\PYGZus{}3=}\PYG{l+s+s2}{\PYGZdq{}}\PYG{p}{,} \PYG{n}{v3}\PYG{p}{)}
\PYG{n}{pprint}\PYG{p}{(}\PYG{l+s+s2}{\PYGZdq{}}\PYG{l+s+s2}{M=}\PYG{l+s+s2}{\PYGZdq{}}\PYG{p}{,} \PYG{n}{M}\PYG{p}{)}
\end{sphinxVerbatim}

\end{sphinxuseclass}\end{sphinxVerbatimInput}
\begin{sphinxVerbatimOutput}

\begin{sphinxuseclass}{cell_output}\begin{equation*}
\begin{split}\displaystyle v_0=\left[
\begin{array}{}
  1.0000 &  2.0000 &  3.0000
\end{array}
\right]\end{split}
\end{equation*}\begin{equation*}
\begin{split}\displaystyle v_1=\left[
\begin{array}{}
  0.0000 &  0.4000 &  0.8000 &  1.2000 &  1.6000 &  2.0000 &  2.4000 &  2.8000
\end{array}
\right]\end{split}
\end{equation*}\begin{equation*}
\begin{split}\displaystyle v_2=\left[
\begin{array}{}
  0.0000 &  0.5556 &  1.1111 &  1.6667 &  2.2222 &  2.7778 &  3.3333 &  3.8889 &  4.4444 &  5.0000
\end{array}
\right]\end{split}
\end{equation*}\begin{equation*}
\begin{split}\displaystyle v_3=\left[
\begin{array}{}
  0.0000 &  0.2222 &  0.4444 &  0.6667 &  0.8889 &  1.1111 &  1.3333 &  1.5556 &  1.7778 &  2.0000
\end{array}
\right]\end{split}
\end{equation*}\begin{equation*}
\begin{split}\displaystyle M=\left[
\begin{array}{}
  1.0000 &  2.0000 &  3.0000\\
  4.0000 &  5.0000 &  6.0000\\
  7.0000 &  8.0000 &  9.0000
\end{array}
\right]\end{split}
\end{equation*}
\end{sphinxuseclass}\end{sphinxVerbatimOutput}

\end{sphinxuseclass}

\subsection{Basic Operations}
\label{\detokenize{01:basic-operations}}
\begin{sphinxShadowBox}
\sphinxstylesidebartitle{}

\sphinxAtStartPar
Remark, those functions (and basically, most of the \sphinxcode{\sphinxupquote{NumPy}} functions) are \sphinxstyleemphasis{vectorial}. This means that they are applied in parallel on all the components of the vectors. Later on, when we will present the \sphinxcode{\sphinxupquote{PyTorch}} library, we will show how this parallelisation can benefit from using specialized hardware, such as the GPU’s.
\end{sphinxShadowBox}

\begin{sphinxuseclass}{cell}\begin{sphinxVerbatimInput}

\begin{sphinxuseclass}{cell_input}
\begin{sphinxVerbatim}[commandchars=\\\{\}]
\PYG{n}{s} \PYG{o}{=} \PYG{n}{np}\PYG{o}{.}\PYG{n}{sum}\PYG{p}{(}\PYG{n}{v0}\PYG{p}{)}
\PYG{n}{m} \PYG{o}{=} \PYG{n}{np}\PYG{o}{.}\PYG{n}{mean}\PYG{p}{(}\PYG{n}{M}\PYG{p}{)}
\PYG{n}{p} \PYG{o}{=} \PYG{n}{np}\PYG{o}{.}\PYG{n}{dot}\PYG{p}{(}\PYG{n}{v2}\PYG{p}{,} \PYG{n}{v3}\PYG{p}{)}
\PYG{n+nb}{print}\PYG{p}{(}\PYG{l+s+s2}{\PYGZdq{}}\PYG{l+s+s2}{Sum of v0 components:}\PYG{l+s+s2}{\PYGZdq{}}\PYG{p}{,} \PYG{n}{s}\PYG{p}{)}
\PYG{n+nb}{print}\PYG{p}{(}\PYG{l+s+s2}{\PYGZdq{}}\PYG{l+s+s2}{Average value in M:}\PYG{l+s+s2}{\PYGZdq{}}\PYG{p}{,} \PYG{n}{m}\PYG{p}{)}
\PYG{n+nb}{print}\PYG{p}{(}\PYG{l+s+s2}{\PYGZdq{}}\PYG{l+s+s2}{Scalar product of v2 and v3:}\PYG{l+s+s2}{\PYGZdq{}}\PYG{p}{,} \PYG{n}{p}\PYG{p}{)}
\end{sphinxVerbatim}

\end{sphinxuseclass}\end{sphinxVerbatimInput}
\begin{sphinxVerbatimOutput}

\begin{sphinxuseclass}{cell_output}
\begin{sphinxVerbatim}[commandchars=\\\{\}]
Sum of v0 components: 6
Average value in M: 5.0
Scalar product of v2 and v3: 35.18518518518518
\end{sphinxVerbatim}

\end{sphinxuseclass}\end{sphinxVerbatimOutput}

\end{sphinxuseclass}
\sphinxAtStartPar
We encourage you to consult the official \sphinxhref{https://numpy.org/doc/stable/}{\sphinxcode{\sphinxupquote{NumPy}} documentation} to explore its many features.
A larger introduction to \sphinxcode{\sphinxupquote{NumPy}} and \sphinxcode{\sphinxupquote{MatPlotLib}} is available in the \(5^{\text{th}}\) practical of the course INFO\sphinxhyphen{}F305 \sphinxhyphen{} \sphinxstyleemphasis{Modélisation et Simulation}.


\section{Probabilistic Foundations}
\label{\detokenize{01:probabilistic-foundations}}
\sphinxAtStartPar
The course of MATH\sphinxhyphen{}F315 \sphinxhyphen{} \sphinxstyleemphasis{Probabilités et Statistiques} has presented all the necessary materials used in this chapter. This is let as reference for the formal demonstrations of the tools and theorems used here. However, we propose a small review of the most important notions that will help you understand the exercises at the end of the chapter. Remember: the exercises are not meant to make you familiar with one specific practical implementation of the techniques and tools, they are intended to make you apply some well understood and completely digested statistical and probabilistic concepts. We strongly recommend the reader to delve in depth in the theoretical concepts related to the exercises.


\subsection{Expectation and Variance}
\label{\detokenize{01:expectation-and-variance}}
\sphinxAtStartPar
The expectation of a random variable \(\mathbf X\), denoted \(E[\mathbf X]\), is defined as:
\begin{equation*}
\begin{split}
E[\mathbf X] = \begin{cases}
    \int_{-\infty}^{\infty} x f_{\mathbf X}(x) dx & \text{for continuous $\mathbf X$} \\
    \sum_{x} x p_{\mathbf X}(x) & \text{for discrete $\mathbf X$}
\end{cases}
\end{split}
\end{equation*}
\sphinxAtStartPar
Statistically, the average is an unbiased estimator of the expected value. In \textbackslash{}codeword\{NumPy\}, the average value in an array can be computed with the \sphinxcode{\sphinxupquote{numpy.mean}} function.
The variance of \(\mathbf X\), denoted \(\text{Var}(\mathbf X)\), measures the dispersion around the mean:
\begin{equation*}
\begin{split}
\text{Var}(\mathbf X) = E[(\mathbf X - E[\mathbf X])^2] = E[\mathbf X^2] - (E[\mathbf X])^2
\end{split}
\end{equation*}
\sphinxAtStartPar
The variance is the square of the standard deviation.
So, in \sphinxcode{\sphinxupquote{NumPy}}, the variance of an array can be computed by squaring the result of the \sphinxcode{\sphinxupquote{numpy.std}} function.
For two random variables \(\mathbf X\) and \textbackslash{}mathbf \(Y\), their covariance is defined as:
\begin{equation*}
\begin{split}
\text{Cov}(\mathbf X, \mathbf Y) = E[(\mathbf X - E[\mathbf X])(\mathbf Y - E[\mathbf Y])] = E[\mathbf X\mathbf Y] - E[\mathbf X]E[\mathbf Y]
\end{split}
\end{equation*}
\sphinxAtStartPar
In \sphinxcode{\sphinxupquote{NumPy}}, the covariance between two vectors can be computed with the \sphinxcode{\sphinxupquote{numpy.cov}} function.
The correlation coefficient, which measures the linear dependence between \(\mathbf X\) and \(\mathbf Y\), is:
\begin{equation*}
\begin{split}
\rho_{\mathbf X,\mathbf Y} = \frac{\text{Cov}(\mathbf X, \mathbf Y)}{\sqrt{\text{Var}(\mathbf X) \cdot \text{Var}(\mathbf Y)}}
\end{split}
\end{equation*}
\sphinxAtStartPar
In \sphinxcode{\sphinxupquote{NumPy}}, the correlation coefficient between two vectors can be computed with the `numpy.corrcoef function.


\subsection{Common Probability Distributions}
\label{\detokenize{01:common-probability-distributions}}

\subsubsection{Uniform Distribution}
\label{\detokenize{01:uniform-distribution}}
\sphinxAtStartPar
A random variable \(\mathbf X\) follows a uniform distribution on the interval \([a, b]\), denoted \(\mathbf X \sim \mathcal{U}(a, b)\), if its PDF is:
\begin{equation*}
\begin{split}
    f_{\mathbf X}(x) = \begin{cases}
    \frac{1}{b - a} & \text{if } x \in [a, b] \\
    0 & \text{otherwise}
    \end{cases}
\end{split}
\end{equation*}
\sphinxAtStartPar
The mean and variance are \(E[\mathbf X] = \frac{a + b}{2}\) and \(\text{Var}\mathbf (X) = \frac{(b - a)^2}{12}\), respectively.
In \sphinxcode{\sphinxupquote{NumPy}}, the uniform distribution can be sampled with the \sphinxcode{\sphinxupquote{numpy.random.uniform}} function.


\subsubsection{Normal (Gaussian) Distribution}
\label{\detokenize{01:normal-gaussian-distribution}}
\sphinxAtStartPar
A random variable \(\mathbf X\) follows a normal distribution with parameters \(\mu\) and \(\sigma^2\), denoted \(X \sim \mathcal{N}(\mu, \sigma^2)\), if its PDF is:
\begin{equation*}
\begin{split}
    f_{\mathbf X}(x) = \frac{1}{\sigma\sqrt{2\pi}} e^{-\frac{(x-\mu)^2}{2\sigma^2}}
\end{split}
\end{equation*}
\sphinxAtStartPar
where \(\mu\) is the mean and \(\sigma^2\) is the variance.
In \sphinxcode{\sphinxupquote{NumPy}}, the uniform distribution can be sampled with the \sphinxcode{\sphinxupquote{numpy.random.normal}} function.


\section{Introduction to the \sphinxstyleliteralintitle{\sphinxupquote{numpy.random}} Module}
\label{\detokenize{01:introduction-to-the-numpy-random-module}}
\sphinxAtStartPar
The \sphinxcode{\sphinxupquote{NumPy}} library includes a \sphinxcode{\sphinxupquote{numpy.random}} submodule dedicated to random number generation. This module offers:
\begin{itemize}
\item {} 
\sphinxAtStartPar
Random number generators for various distributions (uniform, normal, binomial, etc.).

\item {} 
\sphinxAtStartPar
Tools for drawing samples and manipulating random sequences.

\end{itemize}


\subsection{Reproducibility of Results}
\label{\detokenize{01:reproducibility-of-results}}
\sphinxAtStartPar
To ensure the reproducibility of simulations, it is crucial to set a random seed.

\begin{sphinxuseclass}{cell}\begin{sphinxVerbatimInput}

\begin{sphinxuseclass}{cell_input}
\begin{sphinxVerbatim}[commandchars=\\\{\}]
\PYG{n}{np}\PYG{o}{.}\PYG{n}{random}\PYG{o}{.}\PYG{n}{seed}\PYG{p}{(}\PYG{l+m+mi}{42}\PYG{p}{)}
\PYG{n}{x} \PYG{o}{=} \PYG{n}{np}\PYG{o}{.}\PYG{n}{random}\PYG{o}{.}\PYG{n}{rand}\PYG{p}{(}\PYG{l+m+mi}{5}\PYG{p}{)}
\PYG{n}{pprint}\PYG{p}{(}\PYG{l+s+s2}{\PYGZdq{}}\PYG{l+s+s2}{x=}\PYG{l+s+s2}{\PYGZdq{}}\PYG{p}{,} \PYG{n}{x}\PYG{p}{)}
\end{sphinxVerbatim}

\end{sphinxuseclass}\end{sphinxVerbatimInput}
\begin{sphinxVerbatimOutput}

\begin{sphinxuseclass}{cell_output}\begin{equation*}
\begin{split}\displaystyle x=\left[
\begin{array}{}
  0.3745 &  0.9507 &  0.7320 &  0.5987 &  0.1560
\end{array}
\right]\end{split}
\end{equation*}
\end{sphinxuseclass}\end{sphinxVerbatimOutput}

\end{sphinxuseclass}

\subsection{Available Distributions}
\label{\detokenize{01:available-distributions}}
\sphinxAtStartPar
Here are some common distributions:
\begin{itemize}
\item {} 
\sphinxAtStartPar
\sphinxcode{\sphinxupquote{numpy.random.uniform(low, high, size)}} : Uniform distribution.

\item {} 
\sphinxAtStartPar
\sphinxcode{\sphinxupquote{numpy.random.normal(loc, scale, size)}} : Normal (Gaussian) distribution.

\item {} 
\sphinxAtStartPar
\sphinxcode{\sphinxupquote{numpy.random.binomial(n, p, size)}} : Binomial distribution.

\end{itemize}


\subsection{Sampling}
\label{\detokenize{01:sampling}}
\sphinxAtStartPar
Sampling is the process of randomly selecting elements from a population. For example:

\begin{sphinxuseclass}{cell}\begin{sphinxVerbatimInput}

\begin{sphinxuseclass}{cell_input}
\begin{sphinxVerbatim}[commandchars=\\\{\}]
\PYG{n}{population} \PYG{o}{=} \PYG{n}{np}\PYG{o}{.}\PYG{n}{array}\PYG{p}{(}\PYG{p}{[}\PYG{l+m+mi}{1}\PYG{p}{,} \PYG{l+m+mi}{2}\PYG{p}{,} \PYG{l+m+mi}{3}\PYG{p}{,} \PYG{l+m+mi}{4}\PYG{p}{,} \PYG{l+m+mi}{5}\PYG{p}{]}\PYG{p}{)}
\PYG{n}{s} \PYG{o}{=} \PYG{n}{np}\PYG{o}{.}\PYG{n}{random}\PYG{o}{.}\PYG{n}{choice}\PYG{p}{(}\PYG{n}{population}\PYG{p}{,} \PYG{n}{size}\PYG{o}{=}\PYG{p}{(}\PYG{l+m+mi}{3}\PYG{p}{,} \PYG{l+m+mi}{3}\PYG{p}{)}\PYG{p}{,} \PYG{n}{replace}\PYG{o}{=}\PYG{k+kc}{True}\PYG{p}{)}
\PYG{n}{pprint}\PYG{p}{(}\PYG{l+s+s2}{\PYGZdq{}}\PYG{l+s+s2}{s=}\PYG{l+s+s2}{\PYGZdq{}}\PYG{p}{,} \PYG{n}{s}\PYG{p}{)}
\end{sphinxVerbatim}

\end{sphinxuseclass}\end{sphinxVerbatimInput}
\begin{sphinxVerbatimOutput}

\begin{sphinxuseclass}{cell_output}\begin{equation*}
\begin{split}\displaystyle s=\left[
\begin{array}{}
  3.0000 &  3.0000 &  3.0000\\
  5.0000 &  4.0000 &  3.0000\\
  5.0000 &  2.0000 &  4.0000
\end{array}
\right]\end{split}
\end{equation*}
\end{sphinxuseclass}\end{sphinxVerbatimOutput}

\end{sphinxuseclass}

\subsection{Statistical Analysis}
\label{\detokenize{01:statistical-analysis}}
\sphinxAtStartPar
\sphinxcode{\sphinxupquote{Scipy}} offers tools to analyse statistical distributions. As an example, if wan to get density, repartition function and quantiles of a normal distribution \(\mathcal{N}(0, 1)\):

\begin{sphinxuseclass}{cell}\begin{sphinxVerbatimInput}

\begin{sphinxuseclass}{cell_input}
\begin{sphinxVerbatim}[commandchars=\\\{\}]
\PYG{k+kn}{import}\PYG{+w}{ }\PYG{n+nn}{scipy}\PYG{+w}{ }\PYG{k}{as}\PYG{+w}{ }\PYG{n+nn}{sp}
\end{sphinxVerbatim}

\end{sphinxuseclass}\end{sphinxVerbatimInput}

\end{sphinxuseclass}
\sphinxAtStartPar
The probability density function of a normal distribution is given by
\begin{equation*}
\begin{split}f(x) = \frac{(\frac{e^{\frac{-(x-\mu)^2}{2}}}{\sqrt{2\pi}})}{\sigma}\end{split}
\end{equation*}
\sphinxAtStartPar
for \(\mu\) and \(\sigma\) the mean and the standard deviation of the distribution.

\begin{sphinxuseclass}{cell}\begin{sphinxVerbatimInput}

\begin{sphinxuseclass}{cell_input}
\begin{sphinxVerbatim}[commandchars=\\\{\}]
\PYG{n}{mu}\PYG{p}{,} \PYG{n}{sigma} \PYG{o}{=} \PYG{l+m+mi}{0}\PYG{p}{,} \PYG{l+m+mi}{1}
\PYG{n+nb}{print}\PYG{p}{(}\PYG{n}{sp}\PYG{o}{.}\PYG{n}{stats}\PYG{o}{.}\PYG{n}{norm}\PYG{o}{.}\PYG{n}{pdf}\PYG{p}{(}\PYG{l+m+mi}{1}\PYG{p}{,} \PYG{n}{loc}\PYG{o}{=}\PYG{n}{mu}\PYG{p}{,} \PYG{n}{scale}\PYG{o}{=}\PYG{n}{sigma}\PYG{p}{)}\PYG{p}{)}
\end{sphinxVerbatim}

\end{sphinxuseclass}\end{sphinxVerbatimInput}
\begin{sphinxVerbatimOutput}

\begin{sphinxuseclass}{cell_output}
\begin{sphinxVerbatim}[commandchars=\\\{\}]
0.24197072451914337
\end{sphinxVerbatim}

\end{sphinxuseclass}\end{sphinxVerbatimOutput}

\end{sphinxuseclass}
\sphinxAtStartPar
The quantiles can be obtained using the \sphinxcode{\sphinxupquote{ppf}} function:

\begin{sphinxuseclass}{cell}\begin{sphinxVerbatimInput}

\begin{sphinxuseclass}{cell_input}
\begin{sphinxVerbatim}[commandchars=\\\{\}]
\PYG{n+nb}{print}\PYG{p}{(}\PYG{n}{sp}\PYG{o}{.}\PYG{n}{stats}\PYG{o}{.}\PYG{n}{norm}\PYG{o}{.}\PYG{n}{ppf}\PYG{p}{(}\PYG{l+m+mf}{0.75}\PYG{p}{,} \PYG{n}{loc}\PYG{o}{=}\PYG{n}{mu}\PYG{p}{,} \PYG{n}{scale}\PYG{o}{=}\PYG{n}{sigma}\PYG{p}{)}\PYG{p}{)}
\end{sphinxVerbatim}

\end{sphinxuseclass}\end{sphinxVerbatimInput}
\begin{sphinxVerbatimOutput}

\begin{sphinxuseclass}{cell_output}
\begin{sphinxVerbatim}[commandchars=\\\{\}]
0.6744897501960817
\end{sphinxVerbatim}

\end{sphinxuseclass}\end{sphinxVerbatimOutput}

\end{sphinxuseclass}
\sphinxAtStartPar
Finally, the cumulative distribution function can be obtained with \sphinxcode{\sphinxupquote{cdf}}:

\begin{sphinxuseclass}{cell}\begin{sphinxVerbatimInput}

\begin{sphinxuseclass}{cell_input}
\begin{sphinxVerbatim}[commandchars=\\\{\}]
\PYG{n+nb}{print}\PYG{p}{(}\PYG{n}{sp}\PYG{o}{.}\PYG{n}{stats}\PYG{o}{.}\PYG{n}{norm}\PYG{o}{.}\PYG{n}{cdf}\PYG{p}{(}\PYG{l+m+mf}{0.5}\PYG{p}{,} \PYG{n}{loc}\PYG{o}{=}\PYG{n}{mu}\PYG{p}{,} \PYG{n}{scale}\PYG{o}{=}\PYG{n}{sigma}\PYG{p}{)}\PYG{p}{)}
\end{sphinxVerbatim}

\end{sphinxuseclass}\end{sphinxVerbatimInput}
\begin{sphinxVerbatimOutput}

\begin{sphinxuseclass}{cell_output}
\begin{sphinxVerbatim}[commandchars=\\\{\}]
0.6914624612740131
\end{sphinxVerbatim}

\end{sphinxuseclass}\end{sphinxVerbatimOutput}

\end{sphinxuseclass}

\subsection{Exercise: Histogram}
\label{\detokenize{01:exercise-histogram}}
\sphinxAtStartPar
Using the previous \sphinxcode{\sphinxupquote{pdf}} function, plot the histogram of a sampled normal \(\mathcal{N}(0, 1)\) and check that it behaves as the analytical probability density function. Show on the same graph the three quarter quantiles (0.25, 0.5, 0.75).


\subsubsection{Solution}
\label{\detokenize{01:solution}}
\sphinxAtStartPar
In the following code, we sample a normal \(\mathbf{N}(0, 1)\), and, using \sphinxcode{\sphinxupquote{Matplotlib}}, we display the histogram of the sampled values, along the probability density function provided by \sphinxcode{\sphinxupquote{Scipy}}. We also display the three first quarters using the \sphinxcode{\sphinxupquote{scipy.stats.norm.ppf}} function.

\begin{sphinxuseclass}{cell}\begin{sphinxVerbatimInput}

\begin{sphinxuseclass}{cell_input}
\begin{sphinxVerbatim}[commandchars=\\\{\}]
\PYG{k+kn}{import}\PYG{+w}{ }\PYG{n+nn}{matplotlib}\PYG{n+nn}{.}\PYG{n+nn}{pyplot}\PYG{+w}{ }\PYG{k}{as}\PYG{+w}{ }\PYG{n+nn}{plt}
\PYG{n}{plt}\PYG{o}{.}\PYG{n}{style}\PYG{o}{.}\PYG{n}{use}\PYG{p}{(}\PYG{p}{[}\PYG{l+s+s1}{\PYGZsq{}}\PYG{l+s+s1}{seaborn\PYGZhy{}v0\PYGZus{}8\PYGZhy{}muted}\PYG{l+s+s1}{\PYGZsq{}}\PYG{p}{,} \PYG{l+s+s1}{\PYGZsq{}}\PYG{l+s+s1}{practicals.mplstyle}\PYG{l+s+s1}{\PYGZsq{}}\PYG{p}{]}\PYG{p}{)}
\PYG{n}{pprint}\PYG{p}{(}\PYG{l+s+s2}{\PYGZdq{}}\PYG{l+s+s2}{s=}\PYG{l+s+s2}{\PYGZdq{}}\PYG{p}{,} \PYG{n}{sp}\PYG{o}{.}\PYG{n}{stats}\PYG{o}{.}\PYG{n}{norm}\PYG{o}{.}\PYG{n}{ppf}\PYG{p}{(}\PYG{n}{np}\PYG{o}{.}\PYG{n}{array}\PYG{p}{(}\PYG{p}{[}\PYG{l+m+mf}{0.25}\PYG{p}{,} \PYG{l+m+mf}{0.5}\PYG{p}{,} \PYG{l+m+mf}{0.75}\PYG{p}{]}\PYG{p}{)}\PYG{p}{,} \PYG{n}{loc}\PYG{o}{=}\PYG{n}{mu}\PYG{p}{,} \PYG{n}{scale}\PYG{o}{=}\PYG{n}{sigma}\PYG{p}{)}\PYG{p}{)}
\end{sphinxVerbatim}

\end{sphinxuseclass}\end{sphinxVerbatimInput}
\begin{sphinxVerbatimOutput}

\begin{sphinxuseclass}{cell_output}\begin{equation*}
\begin{split}\displaystyle s=\left[
\begin{array}{}
 -0.6745 &  0.0000 &  0.6745
\end{array}
\right]\end{split}
\end{equation*}
\end{sphinxuseclass}\end{sphinxVerbatimOutput}

\end{sphinxuseclass}
\begin{sphinxShadowBox}
\sphinxstylesidebartitle{}

\sphinxAtStartPar
Try changing the number of bins. What happens if it is too low ?
\end{sphinxShadowBox}

\begin{sphinxuseclass}{cell}\begin{sphinxVerbatimInput}

\begin{sphinxuseclass}{cell_input}
\begin{sphinxVerbatim}[commandchars=\\\{\}]
\PYG{n}{n\PYGZus{}bins} \PYG{o}{=} \PYG{l+m+mi}{100}
\PYG{n}{x} \PYG{o}{=} \PYG{n}{np}\PYG{o}{.}\PYG{n}{linspace}\PYG{p}{(}\PYG{o}{\PYGZhy{}}\PYG{l+m+mi}{5}\PYG{p}{,} \PYG{l+m+mi}{5}\PYG{p}{,} \PYG{l+m+mi}{200}\PYG{p}{)}
\PYG{n}{plt}\PYG{o}{.}\PYG{n}{figure}\PYG{p}{(}\PYG{n}{figsize}\PYG{o}{=}\PYG{p}{(}\PYG{l+m+mi}{10}\PYG{p}{,} \PYG{l+m+mi}{5}\PYG{p}{)}\PYG{p}{)}
\PYG{n}{plt}\PYG{o}{.}\PYG{n}{plot}\PYG{p}{(}\PYG{n}{x}\PYG{p}{,} \PYG{n}{sp}\PYG{o}{.}\PYG{n}{stats}\PYG{o}{.}\PYG{n}{norm}\PYG{o}{.}\PYG{n}{pdf}\PYG{p}{(}\PYG{n}{x}\PYG{p}{,} \PYG{n}{loc}\PYG{o}{=}\PYG{n}{mu}\PYG{p}{,} \PYG{n}{scale}\PYG{o}{=}\PYG{n}{sigma}\PYG{p}{)}\PYG{p}{)}
\PYG{n}{plt}\PYG{o}{.}\PYG{n}{hist}\PYG{p}{(}\PYG{n}{np}\PYG{o}{.}\PYG{n}{random}\PYG{o}{.}\PYG{n}{normal}\PYG{p}{(}\PYG{n}{size}\PYG{o}{=}\PYG{p}{(}\PYG{l+m+mi}{25000}\PYG{p}{,}\PYG{p}{)}\PYG{p}{)}\PYG{p}{,} \PYG{n}{bins}\PYG{o}{=}\PYG{n}{n\PYGZus{}bins}\PYG{p}{,} \PYG{n}{density}\PYG{o}{=}\PYG{k+kc}{True}\PYG{p}{)}
\PYG{k}{for} \PYG{n}{q} \PYG{o+ow}{in} \PYG{n}{sp}\PYG{o}{.}\PYG{n}{stats}\PYG{o}{.}\PYG{n}{norm}\PYG{o}{.}\PYG{n}{ppf}\PYG{p}{(}\PYG{n}{np}\PYG{o}{.}\PYG{n}{array}\PYG{p}{(}\PYG{p}{[}\PYG{l+m+mf}{0.25}\PYG{p}{,} \PYG{l+m+mf}{0.5}\PYG{p}{,} \PYG{l+m+mf}{0.75}\PYG{p}{]}\PYG{p}{)}\PYG{p}{,} \PYG{n}{loc}\PYG{o}{=}\PYG{n}{mu}\PYG{p}{,} \PYG{n}{scale}\PYG{o}{=}\PYG{n}{sigma}\PYG{p}{)}\PYG{p}{:}
    \PYG{n}{plt}\PYG{o}{.}\PYG{n}{vlines}\PYG{p}{(}\PYG{n}{q}\PYG{p}{,} \PYG{l+m+mi}{0}\PYG{p}{,} \PYG{l+m+mf}{0.6}\PYG{p}{)}
\PYG{n}{plt}\PYG{o}{.}\PYG{n}{grid}\PYG{p}{(}\PYG{l+s+s2}{\PYGZdq{}}\PYG{l+s+s2}{on}\PYG{l+s+s2}{\PYGZdq{}}\PYG{p}{)}
\PYG{n}{plt}\PYG{o}{.}\PYG{n}{title}\PYG{p}{(}\PYG{l+s+s2}{\PYGZdq{}}\PYG{l+s+s2}{Probability density function of a N(0, 1)}\PYG{l+s+s2}{\PYGZdq{}}\PYG{p}{)}
\PYG{n}{plt}\PYG{o}{.}\PYG{n}{show}\PYG{p}{(}\PYG{p}{)}
\end{sphinxVerbatim}

\end{sphinxuseclass}\end{sphinxVerbatimInput}
\begin{sphinxVerbatimOutput}

\begin{sphinxuseclass}{cell_output}
\noindent\sphinxincludegraphics{{fcd24af6cda6466895f40f7649f784837a77df253623eba163b682f3155ed432}.png}

\end{sphinxuseclass}\end{sphinxVerbatimOutput}

\end{sphinxuseclass}

\section{Monte Carlo Simulation}
\label{\detokenize{01:monte-carlo-simulation}}
\sphinxAtStartPar
The core principle of Monte Carlo methods lies in using random sampling to approximate numerical results. Rather than attempting to solve problems through analytical formulas, Monte Carlo techniques leverage the law of large numbers to estimate solutions through repeated random sampling. This approach is particularly valuable when dealing with high\sphinxhyphen{}dimensional spaces, complex geometries, or systems with multiple degrees of freedom. The simplest visual example, to understand the notion of Monte Carlo, is the following. You want to measure the surface of a lake. Its shape is complex, difficult to measure and to compute the area it covers. However, you can easily surround it by a square for which the side can accurately be measured. You then decide to randomly throw stones with a catapult on the square. The proportion of number of stones that fall into the lake compared to the total number of thrown stones (if the number of stones is infinite) is the same proportion as the proportion of the surface of the lake compared to the square surface. The more stone you throw, the more accurate the estimation is. This idea is close to the one used for Monte Carlo integration.


\subsection{Monte Carlo Integration}
\label{\detokenize{01:monte-carlo-integration}}
\sphinxAtStartPar
Consider the problem of evaluating the integral:
\begin{equation*}
\begin{split}
    I = \int_{\Omega} f(\mathbf{x}) d\mathbf{x}
\end{split}
\end{equation*}
\sphinxAtStartPar
The Monte Carlo estimate of this integral is:
\begin{equation*}
\begin{split}
    \hat{I}_n = \frac{|\Omega|}{n} \sum_{i=1}^{n} f(\mathbf{X}_i)
\end{split}
\end{equation*}
\sphinxAtStartPar
where \(\mathbf{X}_1, \mathbf{X}_2, \dots, \mathbf{X}_n\) are i.i.d. random variables uniformly distributed over \(\Omega\), and \(|\Omega|\) is the measure (volume) of \(\Omega\).
For a definite integral over \([a, b]\), the Monte Carlo estimate is:
\begin{equation*}
\begin{split}
    \int_{a}^{b} f(x) dx \approx (b - a) \cdot \frac{1}{n} \sum_{i=1}^{n} f(X_i)
\end{split}
\end{equation*}
\sphinxAtStartPar
where \(X_i \sim \mathcal{U}(a, b)\).


\subsection{Exercise: Integral Estimation}
\label{\detokenize{01:exercise-integral-estimation}}
\sphinxAtStartPar
Implement a Monte Carlo simulation to estimate the integral of the function \(f(x) = x^2\) over the interval \([0, 1]\).


\subsubsection{Solution}
\label{\detokenize{01:id1}}
\sphinxAtStartPar
The indefinite integral is given by
\begin{equation*}
\begin{split}\int x^2 \,dx = \frac{x^3}{3}\end{split}
\end{equation*}
\sphinxAtStartPar
and therefore the definite interval is
\begin{equation*}
\begin{split}\int_{0}^{1} x^2 \,dx = \frac{1^3}{3} - \frac{0^3}{3} = \frac13\end{split}
\end{equation*}
\begin{sphinxuseclass}{cell}\begin{sphinxVerbatimInput}

\begin{sphinxuseclass}{cell_input}
\begin{sphinxVerbatim}[commandchars=\\\{\}]
\PYG{n}{n\PYGZus{}points} \PYG{o}{=} \PYG{l+m+mi}{10000}
\PYG{n}{approximations} \PYG{o}{=} \PYG{p}{[}\PYG{p}{]}
\PYG{n}{points} \PYG{o}{=} \PYG{n}{np}\PYG{o}{.}\PYG{n}{random}\PYG{o}{.}\PYG{n}{rand}\PYG{p}{(}\PYG{n}{n\PYGZus{}points}\PYG{p}{,} \PYG{l+m+mi}{2}\PYG{p}{)}
\PYG{n}{under\PYGZus{}curve} \PYG{o}{=} \PYG{n}{np}\PYG{o}{.}\PYG{n}{array}\PYG{p}{(}\PYG{n}{points}\PYG{p}{[}\PYG{p}{:}\PYG{p}{,} \PYG{l+m+mi}{1}\PYG{p}{]}\PYG{p}{)} \PYG{o}{\PYGZlt{}}\PYG{o}{=} \PYG{n}{np}\PYG{o}{.}\PYG{n}{array}\PYG{p}{(}\PYG{n}{points}\PYG{p}{[}\PYG{p}{:}\PYG{p}{,} \PYG{l+m+mi}{0}\PYG{p}{]}\PYG{p}{)}\PYG{o}{*}\PYG{o}{*}\PYG{l+m+mi}{2}
\PYG{n}{estimation} \PYG{o}{=} \PYG{n}{np}\PYG{o}{.}\PYG{n}{sum}\PYG{p}{(}\PYG{n}{under\PYGZus{}curve}\PYG{p}{)} \PYG{o}{/} \PYG{n}{n\PYGZus{}points}
\PYG{n}{plt}\PYG{o}{.}\PYG{n}{figure}\PYG{p}{(}\PYG{n}{figsize}\PYG{o}{=}\PYG{p}{(}\PYG{l+m+mi}{10}\PYG{p}{,} \PYG{l+m+mi}{10}\PYG{p}{)}\PYG{p}{)}
\PYG{n}{plt}\PYG{o}{.}\PYG{n}{scatter}\PYG{p}{(}\PYG{n}{points}\PYG{p}{[}\PYG{n}{under\PYGZus{}curve}\PYG{p}{,} \PYG{l+m+mi}{0}\PYG{p}{]}\PYG{p}{,} \PYG{n}{points}\PYG{p}{[}\PYG{n}{under\PYGZus{}curve}\PYG{p}{,} \PYG{l+m+mi}{1}\PYG{p}{]}\PYG{p}{,} \PYG{n}{marker}\PYG{o}{=}\PYG{l+s+s1}{\PYGZsq{}}\PYG{l+s+s1}{.}\PYG{l+s+s1}{\PYGZsq{}}\PYG{p}{,} \PYG{n}{label}\PYG{o}{=}\PYG{l+s+s1}{\PYGZsq{}}\PYG{l+s+s1}{Under the curve}\PYG{l+s+s1}{\PYGZsq{}}\PYG{p}{)}
\PYG{n}{plt}\PYG{o}{.}\PYG{n}{scatter}\PYG{p}{(}\PYG{n}{points}\PYG{p}{[}\PYG{o}{\PYGZti{}}\PYG{n}{under\PYGZus{}curve}\PYG{p}{,} \PYG{l+m+mi}{0}\PYG{p}{]}\PYG{p}{,} \PYG{n}{points}\PYG{p}{[}\PYG{o}{\PYGZti{}}\PYG{n}{under\PYGZus{}curve}\PYG{p}{,} \PYG{l+m+mi}{1}\PYG{p}{]}\PYG{p}{,} \PYG{n}{marker}\PYG{o}{=}\PYG{l+s+s1}{\PYGZsq{}}\PYG{l+s+s1}{.}\PYG{l+s+s1}{\PYGZsq{}}\PYG{p}{,} \PYG{n}{label}\PYG{o}{=}\PYG{l+s+s1}{\PYGZsq{}}\PYG{l+s+s1}{Above the curve}\PYG{l+s+s1}{\PYGZsq{}}\PYG{p}{)}
\PYG{n}{plt}\PYG{o}{.}\PYG{n}{title}\PYG{p}{(}\PYG{l+s+sa}{f}\PYG{l+s+s1}{\PYGZsq{}}\PYG{l+s+s1}{Estimated integral: }\PYG{l+s+si}{\PYGZob{}}\PYG{n}{estimation}\PYG{l+s+si}{:}\PYG{l+s+s1}{.5f}\PYG{l+s+si}{\PYGZcb{}}\PYG{l+s+s1}{\PYGZsq{}}\PYG{p}{)}
\PYG{n}{plt}\PYG{o}{.}\PYG{n}{legend}\PYG{p}{(}\PYG{n}{loc}\PYG{o}{=}\PYG{l+s+s2}{\PYGZdq{}}\PYG{l+s+s2}{upper left}\PYG{l+s+s2}{\PYGZdq{}}\PYG{p}{)}
\PYG{n}{plt}\PYG{o}{.}\PYG{n}{show}\PYG{p}{(}\PYG{p}{)}
\end{sphinxVerbatim}

\end{sphinxuseclass}\end{sphinxVerbatimInput}
\begin{sphinxVerbatimOutput}

\begin{sphinxuseclass}{cell_output}
\noindent\sphinxincludegraphics{{a22258fa0707635cd2a18df0837809354138c55694c8840e0a92fd3b3715554d}.png}

\end{sphinxuseclass}\end{sphinxVerbatimOutput}

\end{sphinxuseclass}

\subsection{Exercise: Probability Estimation}
\label{\detokenize{01:exercise-probability-estimation}}
\sphinxAtStartPar
Use a normal distribution to model a random variable, with parameters \(\mu = 0\) and \(\sigma^2 = \frac12\) and calculate the probability that it takes a value in the interval \([-0.25, 0.25]\).

\begin{sphinxShadowBox}
\sphinxstylesidebartitle{}

\sphinxAtStartPar
\sphinxstylestrong{Don’t skip this part!} Rather, take a sheet of paper, and understand properly every step. You should be able to redo this computation yourself without the solution. Try also changing the values: don’t worry, then, you will be able to validate your theoretical answer by the simulation.
\end{sphinxShadowBox}


\subsubsection{Solution}
\label{\detokenize{01:id2}}
\sphinxAtStartPar
Let us first show how to analytically compute the requested probablity.
The probability that \(X\) lies within the interval is given by:
\begin{equation*}
\begin{split}P(-0.25 \leq \mathbf X \leq 0.25) = \int_{-0.25}^{0.25} f_{\mathbf X}(x) \, dx\end{split}
\end{equation*}
\sphinxAtStartPar
where the Probability Density Function \(f_{\mathbf X}(x)\) is:
\begin{equation*}
\begin{split}f_{\mathbf X}(x) = \frac{1}{\sqrt{2 \pi \sigma^2}} e^{-\frac{(x - \mu)^2}{2 \sigma^2}}\end{split}
\end{equation*}
\sphinxAtStartPar
In our case, \(\mu = 0\) and \(\sigma^2 = 1/2\), so:
\begin{equation*}
\begin{split}f_{\mathbf X}(x) = \frac{1}{\sqrt{\pi}} e^{-2x^2}\end{split}
\end{equation*}
\sphinxAtStartPar
To simplify the calculation, we convert \(X\) to the standard normal variable \(\mathbf Z\):
\begin{equation*}
\begin{split}\mathbf Z = \frac{\mathbf X - \mu}{\sigma}, \quad \text{where } \sigma = \sqrt{1/2} = \frac{\sqrt{2}}{2}\end{split}
\end{equation*}
\sphinxAtStartPar
For the bounds of the interval:
\begin{equation*}
\begin{split}z_{-0.25} = \frac{-0.25 - 0}{\sigma} = -0.25 \cdot \frac{2}{\sqrt{2}} = -\frac{\sqrt{2}}{4}, \quad
z_{0.25} = \frac{0.25 - 0}{\sigma} = \frac{\sqrt{2}}{4}\end{split}
\end{equation*}
\sphinxAtStartPar
The probability is then expressed using the repartition function \(\Phi(z)\) of the standard normal distribution:
\begin{equation*}
\begin{split}P(-0.25 \leq \mathbf X \leq 0.25) = P\left(-\frac{\sqrt{2}}{4} \leq \mathbf Z \leq \frac{\sqrt{2}}{4}\right) = \Phi\left(\frac{\sqrt{2}}{4}\right) - \Phi\left(-\frac{\sqrt{2}}{4}\right)\end{split}
\end{equation*}
\sphinxAtStartPar
Using the symmetry of the standard normal distribution:
\begin{equation*}
\begin{split}\Phi(-z) = 1 - \Phi(z)\end{split}
\end{equation*}
\sphinxAtStartPar
This simplifies to:
\begin{equation*}
\begin{split}P(-0.25 \leq \mathbf X \leq 0.25) = 2 \cdot \Phi\left(\frac{\sqrt{2}}{4}\right) - 1\end{split}
\end{equation*}
\sphinxAtStartPar
To find \(\Phi\left(\frac{\sqrt{2}}{4}\right)\), we can use values in tables. It yields:
\begin{equation*}
\begin{split}\frac{\sqrt{2}}{4} \approx 0.3536\end{split}
\end{equation*}
\sphinxAtStartPar
From standard normal tables or software (we can use \sphinxcode{\sphinxupquote{SciPy}}):

\begin{sphinxuseclass}{cell}\begin{sphinxVerbatimInput}

\begin{sphinxuseclass}{cell_input}
\begin{sphinxVerbatim}[commandchars=\\\{\}]
\PYG{k+kn}{import}\PYG{+w}{ }\PYG{n+nn}{scipy}
\PYG{n}{scipy}\PYG{o}{.}\PYG{n}{stats}\PYG{o}{.}\PYG{n}{norm}\PYG{o}{.}\PYG{n}{cdf}\PYG{p}{(}\PYG{p}{(}\PYG{l+m+mi}{2}\PYG{o}{*}\PYG{o}{*}\PYG{l+m+mf}{0.5}\PYG{p}{)}\PYG{o}{/}\PYG{l+m+mi}{4}\PYG{p}{)}
\end{sphinxVerbatim}

\end{sphinxuseclass}\end{sphinxVerbatimInput}
\begin{sphinxVerbatimOutput}

\begin{sphinxuseclass}{cell_output}
\begin{sphinxVerbatim}[commandchars=\\\{\}]
0.6381631950841185
\end{sphinxVerbatim}

\end{sphinxuseclass}\end{sphinxVerbatimOutput}

\end{sphinxuseclass}\begin{equation*}
\begin{split}\Phi(0.3536) \approx 0.6388\end{split}
\end{equation*}
\sphinxAtStartPar
Thus:
\begin{equation*}
\begin{split}P(-0.25 \leq X \leq 0.25)\end{split}
\end{equation*}
\sphinxAtStartPar
is given by

\begin{sphinxuseclass}{cell}\begin{sphinxVerbatimInput}

\begin{sphinxuseclass}{cell_input}
\begin{sphinxVerbatim}[commandchars=\\\{\}]
\PYG{l+m+mi}{2}\PYG{o}{*}\PYG{n}{scipy}\PYG{o}{.}\PYG{n}{stats}\PYG{o}{.}\PYG{n}{norm}\PYG{o}{.}\PYG{n}{cdf}\PYG{p}{(}\PYG{p}{(}\PYG{l+m+mi}{2}\PYG{o}{*}\PYG{o}{*}\PYG{l+m+mf}{0.5}\PYG{p}{)}\PYG{o}{/}\PYG{l+m+mi}{4}\PYG{p}{)}\PYG{o}{\PYGZhy{}}\PYG{l+m+mi}{1}
\end{sphinxVerbatim}

\end{sphinxuseclass}\end{sphinxVerbatimInput}
\begin{sphinxVerbatimOutput}

\begin{sphinxuseclass}{cell_output}
\begin{sphinxVerbatim}[commandchars=\\\{\}]
0.2763263901682369
\end{sphinxVerbatim}

\end{sphinxuseclass}\end{sphinxVerbatimOutput}

\end{sphinxuseclass}
\sphinxAtStartPar
The Monte Carlo simulation yields a pretty close value:

\begin{sphinxuseclass}{cell}\begin{sphinxVerbatimInput}

\begin{sphinxuseclass}{cell_input}
\begin{sphinxVerbatim}[commandchars=\\\{\}]
\PYG{n}{n\PYGZus{}points} \PYG{o}{=} \PYG{l+m+mi}{100000}
\PYG{n}{approximations} \PYG{o}{=} \PYG{p}{[}\PYG{p}{]}
\PYG{n}{points} \PYG{o}{=} \PYG{n}{np}\PYG{o}{.}\PYG{n}{random}\PYG{o}{.}\PYG{n}{normal}\PYG{p}{(}\PYG{l+m+mi}{0}\PYG{p}{,} \PYG{l+m+mf}{0.5}\PYG{o}{*}\PYG{o}{*}\PYG{l+m+mf}{0.5}\PYG{p}{,} \PYG{n}{size}\PYG{o}{=}\PYG{p}{(}\PYG{n}{n\PYGZus{}points}\PYG{p}{,}\PYG{p}{)}\PYG{p}{)}
\PYG{n}{in\PYGZus{}interval} \PYG{o}{=} \PYG{p}{(}\PYG{n}{np}\PYG{o}{.}\PYG{n}{abs}\PYG{p}{(}\PYG{n}{np}\PYG{o}{.}\PYG{n}{array}\PYG{p}{(}\PYG{n}{points}\PYG{p}{)}\PYG{p}{)} \PYG{o}{\PYGZlt{}}\PYG{o}{=} \PYG{l+m+mf}{0.25}\PYG{p}{)}
\PYG{n}{estimation} \PYG{o}{=} \PYG{n}{np}\PYG{o}{.}\PYG{n}{sum}\PYG{p}{(}\PYG{n}{in\PYGZus{}interval}\PYG{p}{)} \PYG{o}{/} \PYG{n}{n\PYGZus{}points}
\PYG{n+nb}{print}\PYG{p}{(}\PYG{n}{estimation}\PYG{p}{)}
\end{sphinxVerbatim}

\end{sphinxuseclass}\end{sphinxVerbatimInput}
\begin{sphinxVerbatimOutput}

\begin{sphinxuseclass}{cell_output}
\begin{sphinxVerbatim}[commandchars=\\\{\}]
0.27771
\end{sphinxVerbatim}

\end{sphinxuseclass}\end{sphinxVerbatimOutput}

\end{sphinxuseclass}

\subsection{Exercise: Monte Carlo Approximation of Parameters}
\label{\detokenize{01:exercise-monte-carlo-approximation-of-parameters}}
\sphinxAtStartPar
Use Monte Carlo simulation to approximate numerically the mean of a Normal r.v. \({\mathbf Z} \sim N(\mu,\sigma^2)\) and of the variable \({\mathbf Y}= {\mathbf Z}^2\). Verify that MC returns a good approximation of the analytical result. As a hint, note that \(\mbox{Var}[{\mathbf Z}]=E[{\mathbf Z}^2]-E[{\mathbf Z}]^2\)


\subsubsection{Solution}
\label{\detokenize{01:id3}}
\sphinxAtStartPar
First, analytically, we can note that, if \(\mu\) is the expected value of \(\mathbf{Z}\), and  \(\sigma^2\) is its variance, then the expected value of \(\mathbf{Y}=\mathbf{Z}^2\) is given by
\begin{equation*}
\begin{split}
    E[\mathbf Z^2] = Var[Z] + (E[\mathbf Z])^2 = \sigma^2 + \mu^2
\end{split}
\end{equation*}
\sphinxAtStartPar
Secondly, the variance of \(Y=Z^2\) is given by
\begin{equation*}
\begin{split}
    Var[\mathbf Z^2] = E[\mathbf Z^4]-(E[\mathbf Z^2])^2
\end{split}
\end{equation*}
\sphinxAtStartPar
The second term is already known (\(E[\mathbf Z^2] = \sigma^2+\mu^2\)). The first term is the fourth moment (for the development, see the notion \sphinxstyleemphasis{moment generating function} in the course of \sphinxstyleemphasis{Probabilités et Statistiques}) of a normal variable, i.e.:
\begin{equation*}
\begin{split}
    E[\mathbf Z^4] = 3\sigma^4 + 6\sigma^2\mu^2 + \mu^4
\end{split}
\end{equation*}
\sphinxAtStartPar
If we study a concrete example, let us fix \(\mu=\frac12\) and \(\sigma^2=1.7\):

\begin{sphinxuseclass}{cell}\begin{sphinxVerbatimInput}

\begin{sphinxuseclass}{cell_input}
\begin{sphinxVerbatim}[commandchars=\\\{\}]
\PYG{n}{R} \PYG{o}{=} \PYG{l+m+mi}{100000} 
\PYG{n}{Ez} \PYG{o}{=} \PYG{l+m+mf}{0.5}
\PYG{n}{Vz} \PYG{o}{=} \PYG{l+m+mf}{1.7}
\PYG{n}{z} \PYG{o}{=} \PYG{n}{np}\PYG{o}{.}\PYG{n}{random}\PYG{o}{.}\PYG{n}{normal}\PYG{p}{(}\PYG{n}{Ez}\PYG{p}{,} \PYG{n}{np}\PYG{o}{.}\PYG{n}{sqrt}\PYG{p}{(}\PYG{n}{Vz}\PYG{p}{)}\PYG{p}{,} \PYG{n}{R}\PYG{p}{)}
\PYG{n}{y} \PYG{o}{=} \PYG{n}{z}\PYG{o}{*}\PYG{o}{*}\PYG{l+m+mi}{2}
\PYG{n}{pprint}\PYG{p}{(}\PYG{l+s+sa}{f}\PYG{l+s+s2}{\PYGZdq{}}\PYG{l+s+s2}{E[}\PYG{l+s+se}{\PYGZbs{}\PYGZbs{}}\PYG{l+s+s2}{mathbf Z]=}\PYG{l+s+si}{\PYGZob{}}\PYG{n}{Ez}\PYG{l+s+si}{\PYGZcb{}}\PYG{l+s+s2}{;}\PYG{l+s+s2}{\PYGZbs{}}\PYG{l+s+s2}{ MC}\PYG{l+s+s2}{\PYGZbs{}}\PYG{l+s+s2}{ E[}\PYG{l+s+se}{\PYGZbs{}\PYGZbs{}}\PYG{l+s+s2}{mathbf Z]=}\PYG{l+s+si}{\PYGZob{}}\PYG{n}{np}\PYG{o}{.}\PYG{n}{mean}\PYG{p}{(}\PYG{n}{z}\PYG{p}{)}\PYG{l+s+si}{:}\PYG{l+s+s2}{.4f}\PYG{l+s+si}{\PYGZcb{}}\PYG{l+s+s2}{\PYGZdq{}}\PYG{p}{)}
\PYG{n}{pprint}\PYG{p}{(}\PYG{l+s+sa}{f}\PYG{l+s+s2}{\PYGZdq{}}\PYG{l+s+s2}{Var[}\PYG{l+s+se}{\PYGZbs{}\PYGZbs{}}\PYG{l+s+s2}{mathbf Z]=}\PYG{l+s+si}{\PYGZob{}}\PYG{n}{Vz}\PYG{l+s+si}{\PYGZcb{}}\PYG{l+s+s2}{;}\PYG{l+s+s2}{\PYGZbs{}}\PYG{l+s+s2}{ MC}\PYG{l+s+s2}{\PYGZbs{}}\PYG{l+s+s2}{ Var[}\PYG{l+s+se}{\PYGZbs{}\PYGZbs{}}\PYG{l+s+s2}{mathbf Z]=}\PYG{l+s+si}{\PYGZob{}}\PYG{n}{np}\PYG{o}{.}\PYG{n}{var}\PYG{p}{(}\PYG{n}{z}\PYG{p}{)}\PYG{l+s+si}{:}\PYG{l+s+s2}{.4f}\PYG{l+s+si}{\PYGZcb{}}\PYG{l+s+s2}{\PYGZdq{}}\PYG{p}{)}
\PYG{n}{pprint}\PYG{p}{(}\PYG{l+s+sa}{f}\PYG{l+s+s2}{\PYGZdq{}}\PYG{l+s+s2}{E[}\PYG{l+s+se}{\PYGZbs{}\PYGZbs{}}\PYG{l+s+s2}{mathbf Z\PYGZca{}2]=}\PYG{l+s+si}{\PYGZob{}}\PYG{n}{Vz}\PYG{o}{+}\PYG{n}{Ez}\PYG{o}{*}\PYG{o}{*}\PYG{l+m+mi}{2}\PYG{l+s+si}{\PYGZcb{}}\PYG{l+s+s2}{;}\PYG{l+s+s2}{\PYGZbs{}}\PYG{l+s+s2}{ MC}\PYG{l+s+s2}{\PYGZbs{}}\PYG{l+s+s2}{ E[}\PYG{l+s+se}{\PYGZbs{}\PYGZbs{}}\PYG{l+s+s2}{mathbf Y=}\PYG{l+s+se}{\PYGZbs{}\PYGZbs{}}\PYG{l+s+s2}{mathbf Z\PYGZca{}2]=}\PYG{l+s+si}{\PYGZob{}}\PYG{n}{np}\PYG{o}{.}\PYG{n}{mean}\PYG{p}{(}\PYG{n}{y}\PYG{p}{)}\PYG{l+s+si}{:}\PYG{l+s+s2}{.4f}\PYG{l+s+si}{\PYGZcb{}}\PYG{l+s+s2}{\PYGZdq{}}\PYG{p}{)}
\PYG{n}{pprint}\PYG{p}{(}\PYG{l+s+sa}{f}\PYG{l+s+s2}{\PYGZdq{}}\PYG{l+s+s2}{Var[}\PYG{l+s+se}{\PYGZbs{}\PYGZbs{}}\PYG{l+s+s2}{mathbf Z\PYGZca{}2]=}\PYG{l+s+si}{\PYGZob{}}\PYG{p}{(}\PYG{l+m+mi}{3}\PYG{o}{*}\PYG{n}{Vz}\PYG{o}{*}\PYG{o}{*}\PYG{l+m+mi}{2}\PYG{o}{+}\PYG{l+m+mi}{6}\PYG{o}{*}\PYG{n}{Vz}\PYG{o}{*}\PYG{n}{Ez}\PYG{o}{*}\PYG{o}{*}\PYG{l+m+mi}{2}\PYG{o}{+}\PYG{n}{Ez}\PYG{o}{*}\PYG{o}{*}\PYG{l+m+mi}{4}\PYG{p}{)}\PYG{+w}{ }\PYG{o}{\PYGZhy{}}\PYG{+w}{ }\PYG{p}{(}\PYG{n}{Vz}\PYG{o}{+}\PYG{n}{Ez}\PYG{o}{*}\PYG{o}{*}\PYG{l+m+mi}{2}\PYG{p}{)}\PYG{o}{*}\PYG{o}{*}\PYG{l+m+mi}{2}\PYG{l+s+si}{:}\PYG{l+s+s2}{.4f}\PYG{l+s+si}{\PYGZcb{}}\PYG{l+s+s2}{;}\PYG{l+s+s2}{\PYGZbs{}}\PYG{l+s+s2}{ MC}\PYG{l+s+s2}{\PYGZbs{}}\PYG{l+s+s2}{ Var[}\PYG{l+s+se}{\PYGZbs{}\PYGZbs{}}\PYG{l+s+s2}{mathbf Y=}\PYG{l+s+se}{\PYGZbs{}\PYGZbs{}}\PYG{l+s+s2}{mathbf Z\PYGZca{}2]=}\PYG{l+s+si}{\PYGZob{}}\PYG{n}{np}\PYG{o}{.}\PYG{n}{std}\PYG{p}{(}\PYG{n}{y}\PYG{p}{)}\PYG{o}{*}\PYG{o}{*}\PYG{l+m+mi}{2}\PYG{l+s+si}{:}\PYG{l+s+s2}{.4f}\PYG{l+s+si}{\PYGZcb{}}\PYG{l+s+s2}{\PYGZdq{}}\PYG{p}{)}
\end{sphinxVerbatim}

\end{sphinxuseclass}\end{sphinxVerbatimInput}
\begin{sphinxVerbatimOutput}

\begin{sphinxuseclass}{cell_output}\begin{equation*}
\begin{split}\displaystyle E[\mathbf Z]=0.5;\ MC\ E[\mathbf Z]=0.5015\end{split}
\end{equation*}\begin{equation*}
\begin{split}\displaystyle Var[\mathbf Z]=1.7;\ MC\ Var[\mathbf Z]=1.6931\end{split}
\end{equation*}\begin{equation*}
\begin{split}\displaystyle E[\mathbf Z^2]=1.95;\ MC\ E[\mathbf Y=\mathbf Z^2]=1.9446\end{split}
\end{equation*}\begin{equation*}
\begin{split}\displaystyle Var[\mathbf Z^2]=7.4800;\ MC\ Var[\mathbf Y=\mathbf Z^2]=7.4843\end{split}
\end{equation*}
\end{sphinxuseclass}\end{sphinxVerbatimOutput}

\end{sphinxuseclass}

\subsection{Exercise: Absolute Value}
\label{\detokenize{01:exercise-absolute-value}}
\sphinxAtStartPar
Use MC to approximate numerically \(E[ {\mathbf K}]\) and \(\mbox{Var}[ {\mathbf K}]\) where \({\mathbf K} = |{\mathbf Z}|\) and check that the approximation is good w.r.t. the analytical results in \sphinxhref{https://www.quora.com/If-Y-X-where-X-has-normal-distribution-N-0-1-what-is-the-density-function-expectation-and-variance-of-Y}{this article}.


\subsubsection{Solution}
\label{\detokenize{01:id4}}
\sphinxAtStartPar
The process is truly similar to the previous one.

\begin{sphinxuseclass}{cell}\begin{sphinxVerbatimInput}

\begin{sphinxuseclass}{cell_input}
\begin{sphinxVerbatim}[commandchars=\\\{\}]
\PYG{n}{Ez} \PYG{o}{=} \PYG{l+m+mi}{0}
\PYG{n}{Vz} \PYG{o}{=} \PYG{l+m+mi}{1}
\PYG{n}{z} \PYG{o}{=} \PYG{n}{np}\PYG{o}{.}\PYG{n}{random}\PYG{o}{.}\PYG{n}{normal}\PYG{p}{(}\PYG{n}{Ez}\PYG{p}{,} \PYG{n}{np}\PYG{o}{.}\PYG{n}{sqrt}\PYG{p}{(}\PYG{n}{Vz}\PYG{p}{)}\PYG{p}{,} \PYG{n}{R}\PYG{p}{)}
\PYG{n}{k} \PYG{o}{=} \PYG{n}{np}\PYG{o}{.}\PYG{n}{abs}\PYG{p}{(}\PYG{n}{z}\PYG{p}{)}

\PYG{n}{Ek} \PYG{o}{=} \PYG{n}{np}\PYG{o}{.}\PYG{n}{sqrt}\PYG{p}{(}\PYG{l+m+mi}{2}\PYG{o}{/}\PYG{n}{np}\PYG{o}{.}\PYG{n}{pi}\PYG{p}{)}
\PYG{n}{Vk} \PYG{o}{=} \PYG{l+m+mi}{1} \PYG{o}{\PYGZhy{}} \PYG{l+m+mi}{2}\PYG{o}{/}\PYG{n}{np}\PYG{o}{.}\PYG{n}{pi}

\PYG{n}{pprint}\PYG{p}{(}\PYG{l+s+sa}{f}\PYG{l+s+s2}{\PYGZdq{}}\PYG{l+s+s2}{E[}\PYG{l+s+se}{\PYGZbs{}\PYGZbs{}}\PYG{l+s+s2}{mathbf K] = }\PYG{l+s+si}{\PYGZob{}}\PYG{n}{Ek}\PYG{l+s+si}{:}\PYG{l+s+s2}{.4f}\PYG{l+s+si}{\PYGZcb{}}\PYG{l+s+s2}{;}\PYG{l+s+s2}{\PYGZbs{}}\PYG{l+s+s2}{ MC}\PYG{l+s+s2}{\PYGZbs{}}\PYG{l+s+s2}{ E[}\PYG{l+s+se}{\PYGZbs{}\PYGZbs{}}\PYG{l+s+s2}{mathbf K=|}\PYG{l+s+se}{\PYGZbs{}\PYGZbs{}}\PYG{l+s+s2}{mathbf Z|] = }\PYG{l+s+si}{\PYGZob{}}\PYG{n}{np}\PYG{o}{.}\PYG{n}{mean}\PYG{p}{(}\PYG{n}{k}\PYG{p}{)}\PYG{l+s+si}{:}\PYG{l+s+s2}{.4f}\PYG{l+s+si}{\PYGZcb{}}\PYG{l+s+s2}{\PYGZdq{}}\PYG{p}{)}
\PYG{n}{pprint}\PYG{p}{(}\PYG{l+s+sa}{f}\PYG{l+s+s2}{\PYGZdq{}}\PYG{l+s+s2}{Var[}\PYG{l+s+se}{\PYGZbs{}\PYGZbs{}}\PYG{l+s+s2}{mathbf K] = }\PYG{l+s+si}{\PYGZob{}}\PYG{n}{Vk}\PYG{l+s+si}{:}\PYG{l+s+s2}{.4f}\PYG{l+s+si}{\PYGZcb{}}\PYG{l+s+s2}{;}\PYG{l+s+s2}{\PYGZbs{}}\PYG{l+s+s2}{ MC}\PYG{l+s+s2}{\PYGZbs{}}\PYG{l+s+s2}{ Var[}\PYG{l+s+se}{\PYGZbs{}\PYGZbs{}}\PYG{l+s+s2}{mathbf K=|}\PYG{l+s+se}{\PYGZbs{}\PYGZbs{}}\PYG{l+s+s2}{mathbf Z|] = }\PYG{l+s+si}{\PYGZob{}}\PYG{n}{np}\PYG{o}{.}\PYG{n}{var}\PYG{p}{(}\PYG{n}{k}\PYG{p}{)}\PYG{l+s+si}{:}\PYG{l+s+s2}{.4f}\PYG{l+s+si}{\PYGZcb{}}\PYG{l+s+s2}{\PYGZdq{}}\PYG{p}{)}
\end{sphinxVerbatim}

\end{sphinxuseclass}\end{sphinxVerbatimInput}
\begin{sphinxVerbatimOutput}

\begin{sphinxuseclass}{cell_output}\begin{equation*}
\begin{split}\displaystyle E[\mathbf K] = 0.7979;\ MC\ E[\mathbf K=|\mathbf Z|] = 0.7998\end{split}
\end{equation*}\begin{equation*}
\begin{split}\displaystyle Var[\mathbf K] = 0.3634;\ MC\ Var[\mathbf K=|\mathbf Z|] = 0.3633\end{split}
\end{equation*}
\end{sphinxuseclass}\end{sphinxVerbatimOutput}

\end{sphinxuseclass}

\subsection{Exercise: Expected Value for any Function}
\label{\detokenize{01:exercise-expected-value-for-any-function}}
\sphinxAtStartPar
Use MC to approximate numerically the value of \(E[f({\mathbf X})]\) and \(f(E[{\mathbf X}])\) for a given deterministic function \(f\), where \({\mathbf X}\) is a normal random variable.


\subsubsection{Solution}
\label{\detokenize{01:id5}}
\sphinxAtStartPar
The following script show the MC process to estimate the expected value of any Python function.

\begin{sphinxShadowBox}
\sphinxstylesidebartitle{}

\sphinxAtStartPar
Change the function \sphinxcode{\sphinxupquote{f}}. Is the process valid in all the cases ? Try to find an example of function for which this process could fail. Hint: some functions might not be defined depending on their input.
\end{sphinxShadowBox}

\begin{sphinxuseclass}{cell}\begin{sphinxVerbatimInput}

\begin{sphinxuseclass}{cell_input}
\begin{sphinxVerbatim}[commandchars=\\\{\}]
\PYG{n}{np}\PYG{o}{.}\PYG{n}{random}\PYG{o}{.}\PYG{n}{seed}\PYG{p}{(}\PYG{l+m+mi}{0}\PYG{p}{)}
\PYG{n}{R} \PYG{o}{=} \PYG{l+m+mi}{10000}

\PYG{k}{def}\PYG{+w}{ }\PYG{n+nf}{f}\PYG{p}{(}\PYG{n}{x}\PYG{p}{)}\PYG{p}{:}
    \PYG{k}{return} \PYG{o}{\PYGZhy{}}\PYG{n}{x}\PYG{o}{*}\PYG{o}{*}\PYG{l+m+mi}{2}

\PYG{n}{FX} \PYG{o}{=} \PYG{p}{[}\PYG{p}{]}
\PYG{n}{X} \PYG{o}{=} \PYG{p}{[}\PYG{p}{]}

\PYG{n}{a} \PYG{o}{=} \PYG{l+m+mi}{0}
\PYG{n}{b} \PYG{o}{=} \PYG{l+m+mi}{1}

\PYG{k}{for} \PYG{n}{r} \PYG{o+ow}{in} \PYG{n+nb}{range}\PYG{p}{(}\PYG{n}{R}\PYG{p}{)}\PYG{p}{:}
    \PYG{n}{x} \PYG{o}{=} \PYG{n}{np}\PYG{o}{.}\PYG{n}{random}\PYG{o}{.}\PYG{n}{normal}\PYG{p}{(}\PYG{n}{a}\PYG{p}{,} \PYG{n}{b}\PYG{p}{)}
    \PYG{n}{FX}\PYG{o}{.}\PYG{n}{append}\PYG{p}{(}\PYG{n}{f}\PYG{p}{(}\PYG{n}{x}\PYG{p}{)}\PYG{p}{)}
    \PYG{n}{X}\PYG{o}{.}\PYG{n}{append}\PYG{p}{(}\PYG{n}{x}\PYG{p}{)}

\PYG{n}{muX} \PYG{o}{=} \PYG{n}{np}\PYG{o}{.}\PYG{n}{mean}\PYG{p}{(}\PYG{n}{X}\PYG{p}{)}

\PYG{n}{pprint}\PYG{p}{(}\PYG{l+s+sa}{f}\PYG{l+s+s2}{\PYGZdq{}}\PYG{l+s+s2}{E[f(x)]= }\PYG{l+s+si}{\PYGZob{}}\PYG{n}{np}\PYG{o}{.}\PYG{n}{mean}\PYG{p}{(}\PYG{n}{FX}\PYG{p}{)}\PYG{l+s+si}{:}\PYG{l+s+s2}{.4f}\PYG{l+s+si}{\PYGZcb{}}\PYG{l+s+s2}{\PYGZdq{}}\PYG{p}{)}
\PYG{n}{pprint}\PYG{p}{(}\PYG{l+s+sa}{f}\PYG{l+s+s2}{\PYGZdq{}}\PYG{l+s+s2}{f(E[x])= }\PYG{l+s+si}{\PYGZob{}}\PYG{n}{f}\PYG{p}{(}\PYG{n}{muX}\PYG{p}{)}\PYG{l+s+si}{:}\PYG{l+s+s2}{.4f}\PYG{l+s+si}{\PYGZcb{}}\PYG{l+s+s2}{\PYGZdq{}}\PYG{p}{)}
\end{sphinxVerbatim}

\end{sphinxuseclass}\end{sphinxVerbatimInput}
\begin{sphinxVerbatimOutput}

\begin{sphinxuseclass}{cell_output}\begin{equation*}
\begin{split}\displaystyle E[f(x)]= -0.9756\end{split}
\end{equation*}\begin{equation*}
\begin{split}\displaystyle f(E[x])= -0.0003\end{split}
\end{equation*}
\end{sphinxuseclass}\end{sphinxVerbatimOutput}

\end{sphinxuseclass}

\subsection{Exercise: Variance of Uniform Random Variable}
\label{\detokenize{01:exercise-variance-of-uniform-random-variable}}
\sphinxAtStartPar
Use MC to approximate numerically the value of the variance \(\mbox{Var}[{\mathbf Z}]\) where \({\mathbf Z} \sim U(a,b)\) is uniformly distributed. Check that the result is a good approximation of the analytical value.


\subsubsection{Solution}
\label{\detokenize{01:id6}}
\sphinxAtStartPar
The following script estimates the variance of an uniformly distributed variable by Monte Carlo. The distribution is uniform between \(-1\) and \(1\).

\begin{sphinxuseclass}{cell}\begin{sphinxVerbatimInput}

\begin{sphinxuseclass}{cell_input}
\begin{sphinxVerbatim}[commandchars=\\\{\}]
\PYG{n}{R} \PYG{o}{=} \PYG{l+m+mi}{50000}
\PYG{n}{a} \PYG{o}{=} \PYG{o}{\PYGZhy{}}\PYG{l+m+mf}{1.0}
\PYG{n}{b} \PYG{o}{=} \PYG{l+m+mf}{1.0}
\PYG{n}{muz} \PYG{o}{=} \PYG{p}{(}\PYG{n}{b} \PYG{o}{+} \PYG{n}{a}\PYG{p}{)} \PYG{o}{/} \PYG{l+m+mi}{2}

\PYG{n}{Z} \PYG{o}{=} \PYG{n}{np}\PYG{o}{.}\PYG{n}{zeros}\PYG{p}{(}\PYG{n}{R}\PYG{p}{)}
\PYG{k}{for} \PYG{n}{r} \PYG{o+ow}{in} \PYG{n+nb}{range}\PYG{p}{(}\PYG{n}{R}\PYG{p}{)}\PYG{p}{:}
    \PYG{n}{z} \PYG{o}{=} \PYG{n}{np}\PYG{o}{.}\PYG{n}{random}\PYG{o}{.}\PYG{n}{uniform}\PYG{p}{(}\PYG{n}{a}\PYG{p}{,} \PYG{n}{b}\PYG{p}{)}
    \PYG{n}{Z}\PYG{p}{[}\PYG{n}{r}\PYG{p}{]} \PYG{o}{=} \PYG{p}{(}\PYG{n}{z} \PYG{o}{\PYGZhy{}} \PYG{n}{muz}\PYG{p}{)}\PYG{o}{*}\PYG{o}{*}\PYG{l+m+mi}{2}

\PYG{n}{pprint}\PYG{p}{(}\PYG{l+s+sa}{f}\PYG{l+s+s2}{\PYGZdq{}}\PYG{l+s+s2}{Var[}\PYG{l+s+se}{\PYGZbs{}\PYGZbs{}}\PYG{l+s+s2}{mathbf Z]= }\PYG{l+s+si}{\PYGZob{}}\PYG{p}{(}\PYG{n}{b}\PYG{o}{\PYGZhy{}}\PYG{n}{a}\PYG{p}{)}\PYG{o}{*}\PYG{o}{*}\PYG{l+m+mi}{2}\PYG{o}{/}\PYG{l+m+mi}{12}\PYG{l+s+si}{:}\PYG{l+s+s2}{.4f}\PYG{l+s+si}{\PYGZcb{}}\PYG{l+s+s2}{;}\PYG{l+s+s2}{\PYGZbs{}}\PYG{l+s+s2}{ MC}\PYG{l+s+s2}{\PYGZbs{}}\PYG{l+s+s2}{ Var[}\PYG{l+s+se}{\PYGZbs{}\PYGZbs{}}\PYG{l+s+s2}{mathbf Z]= }\PYG{l+s+si}{\PYGZob{}}\PYG{n}{np}\PYG{o}{.}\PYG{n}{mean}\PYG{p}{(}\PYG{n}{Z}\PYG{p}{)}\PYG{l+s+si}{:}\PYG{l+s+s2}{.4f}\PYG{l+s+si}{\PYGZcb{}}\PYG{l+s+s2}{\PYGZdq{}}\PYG{p}{)}
\end{sphinxVerbatim}

\end{sphinxuseclass}\end{sphinxVerbatimInput}
\begin{sphinxVerbatimOutput}

\begin{sphinxuseclass}{cell_output}\begin{equation*}
\begin{split}\displaystyle Var[\mathbf Z]= 0.3333;\ MC\ Var[\mathbf Z]= 0.3340\end{split}
\end{equation*}
\end{sphinxuseclass}\end{sphinxVerbatimOutput}

\end{sphinxuseclass}

\subsection{Exercise: Covariance of a Multiple of a Random Variable}
\label{\detokenize{01:exercise-covariance-of-a-multiple-of-a-random-variable}}
\sphinxAtStartPar
Use MC to approximate the value of the covariance of \({\mathbf X}\) and \({\mathbf Y}=K{\mathbf X}\) and verify that the result is a good approximation of the analytical derivation. Check for different distributions of \({\mathbf X}\). As a hint, note that
\begin{equation*}
\begin{split}
\begin{align}
\mbox{Cov}({\mathbf X},{\mathbf Y})&=E[{\mathbf X}{\mathbf Y}]-K E[{\mathbf X}]E[{\mathbf X}]\\
&=E[K*{\mathbf X}^2]-K(E[{\mathbf X}]^2)\\
&=K(\mbox{Var}({\mathbf X})+E[{\mathbf X}]^2)-K (E[{\mathbf X}])^2\\
&=K\mbox{Var}({\mathbf X})
\end{align}\end{split}
\end{equation*}

\subsubsection{Solution}
\label{\detokenize{01:id7}}
\sphinxAtStartPar
The following code shows the covariance estimation of a sampled random variable with a multiple of this random variable. The estimated covariance is close to the analytical value.

\begin{sphinxuseclass}{cell}\begin{sphinxVerbatimInput}

\begin{sphinxuseclass}{cell_input}
\begin{sphinxVerbatim}[commandchars=\\\{\}]
\PYG{n}{R} \PYG{o}{=} \PYG{l+m+mi}{50000}
\PYG{n}{distr} \PYG{o}{=} \PYG{l+s+s2}{\PYGZdq{}}\PYG{l+s+s2}{uniform}\PYG{l+s+s2}{\PYGZdq{}}
\PYG{n}{XY} \PYG{o}{=} \PYG{p}{[}\PYG{p}{]}
\PYG{n}{X} \PYG{o}{=} \PYG{p}{[}\PYG{p}{]}
\PYG{n}{K} \PYG{o}{=} \PYG{l+m+mi}{2}
\PYG{k}{for} \PYG{n}{r} \PYG{o+ow}{in} \PYG{n+nb}{range}\PYG{p}{(}\PYG{n}{R}\PYG{p}{)}\PYG{p}{:}
    \PYG{k}{if} \PYG{n}{distr} \PYG{o}{==} \PYG{l+s+s2}{\PYGZdq{}}\PYG{l+s+s2}{uniform}\PYG{l+s+s2}{\PYGZdq{}}\PYG{p}{:}
        \PYG{n}{a}\PYG{p}{,} \PYG{n}{b} \PYG{o}{=} \PYG{l+m+mi}{1}\PYG{p}{,} \PYG{l+m+mi}{20}
        \PYG{n}{x} \PYG{o}{=} \PYG{n}{np}\PYG{o}{.}\PYG{n}{random}\PYG{o}{.}\PYG{n}{uniform}\PYG{p}{(}\PYG{n}{a}\PYG{p}{,} \PYG{n}{b}\PYG{p}{)}
        \PYG{n}{VX} \PYG{o}{=} \PYG{l+m+mi}{1}\PYG{o}{/}\PYG{l+m+mi}{12} \PYG{o}{*} \PYG{p}{(}\PYG{n}{b}\PYG{o}{\PYGZhy{}}\PYG{n}{a}\PYG{p}{)}\PYG{o}{*}\PYG{o}{*}\PYG{l+m+mi}{2}
    \PYG{k}{else}\PYG{p}{:}
        \PYG{n}{mu}\PYG{p}{,} \PYG{n}{sigma} \PYG{o}{=} \PYG{l+m+mi}{1}\PYG{p}{,} \PYG{l+m+mi}{1}
        \PYG{n}{x} \PYG{o}{=} \PYG{n}{np}\PYG{o}{.}\PYG{n}{random}\PYG{o}{.}\PYG{n}{normal}\PYG{p}{(}\PYG{n}{mu}\PYG{p}{,} \PYG{n}{sigma}\PYG{o}{*}\PYG{o}{*}\PYG{l+m+mi}{2}\PYG{p}{)}
    \PYG{n}{y} \PYG{o}{=} \PYG{n}{K} \PYG{o}{*} \PYG{n}{x}
    \PYG{n}{XY}\PYG{o}{.}\PYG{n}{append}\PYG{p}{(}\PYG{n}{x} \PYG{o}{*} \PYG{n}{y}\PYG{p}{)}
    \PYG{n}{X}\PYG{o}{.}\PYG{n}{append}\PYG{p}{(}\PYG{n}{x}\PYG{p}{)}
\PYG{n+nb}{print}\PYG{p}{(}\PYG{l+s+sa}{f}\PYG{l+s+s2}{\PYGZdq{}}\PYG{l+s+s2}{Analytical covariance: }\PYG{l+s+si}{\PYGZob{}}\PYG{n}{K}\PYG{o}{*}\PYG{n}{VX}\PYG{l+s+si}{\PYGZcb{}}\PYG{l+s+s2}{\PYGZdq{}}\PYG{p}{)}
\PYG{n+nb}{print}\PYG{p}{(}\PYG{l+s+sa}{f}\PYG{l+s+s2}{\PYGZdq{}}\PYG{l+s+s2}{Monte Carlo covariance: }\PYG{l+s+si}{\PYGZob{}}\PYG{n}{np}\PYG{o}{.}\PYG{n}{mean}\PYG{p}{(}\PYG{n}{XY}\PYG{p}{)}\PYG{+w}{ }\PYG{o}{\PYGZhy{}}\PYG{+w}{ }\PYG{n}{K}\PYG{o}{*}\PYG{n}{np}\PYG{o}{.}\PYG{n}{mean}\PYG{p}{(}\PYG{n}{X}\PYG{p}{)}\PYG{o}{*}\PYG{o}{*}\PYG{l+m+mi}{2}\PYG{l+s+si}{\PYGZcb{}}\PYG{l+s+s2}{\PYGZdq{}}\PYG{p}{)}
\end{sphinxVerbatim}

\end{sphinxuseclass}\end{sphinxVerbatimInput}
\begin{sphinxVerbatimOutput}

\begin{sphinxuseclass}{cell_output}
\begin{sphinxVerbatim}[commandchars=\\\{\}]
Analytical covariance: 60.166666666666664
Monte Carlo covariance: 60.455461621858774
\end{sphinxVerbatim}

\end{sphinxuseclass}\end{sphinxVerbatimOutput}

\end{sphinxuseclass}

\section{Multivariate Gaussian Distributions}
\label{\detokenize{01:multivariate-gaussian-distributions}}
\sphinxAtStartPar
Multivariate Gaussian distributions are multidimensional generalizations of the normal distribution (that you already know well).
They can be described by mean vector \(\mu\) and a covariance matrix \(\Sigma\), which shows the covariance between each of the distribution’s dimensions. As an example, if this matrix is diagonal, the variables corresponding to each dimension are independant. In the same way, this matrix has to be diagonal (can you tell why ?).


\subsection{Exercise: 2D Monte Carlo Simulation}
\label{\detokenize{01:exercise-2d-monte-carlo-simulation}}
\sphinxAtStartPar
Use MC simulation approximate the mean and covairance of 2\sphinxhyphen{}dimensional Gaussian random data with the following parameters:
\begin{itemize}
\item {} 
\sphinxAtStartPar
Mean vector: \([2, 3]\)

\item {} 
\sphinxAtStartPar
Covariance matrix: \(\begin{bmatrix}2&0.8\\0.8&1\end{bmatrix}\)

\end{itemize}

\sphinxAtStartPar
Visualize the used data as a scatter plot. Compare the sample mean and the given mean vector.
Hint: use \sphinxcode{\sphinxupquote{scipy.stats.multivariate\_normal}} for generation and numpy to calculate the sample mean, and use \sphinxcode{\sphinxupquote{numpy.cov}} for covariance estimation.
Ensure you understand the difference between population and sample covariance.


\subsubsection{Solution}
\label{\detokenize{01:id8}}
\sphinxAtStartPar
The \sphinxcode{\sphinxupquote{rvs}} from \sphinxcode{\sphinxupquote{scipy.stats.multivariate\textbackslash{}\_normal}} (for \sphinxstyleemphasis{random variable sample} allows to sample a bivariate Gaussian adequately defined. The following code shows how to sample such kind of distribution, how to display it in \sphinxcode{\sphinxupquote{Matplotlib}}, and how to estimate the covariance matrix.

\begin{sphinxuseclass}{cell}\begin{sphinxVerbatimInput}

\begin{sphinxuseclass}{cell_input}
\begin{sphinxVerbatim}[commandchars=\\\{\}]
\PYG{k+kn}{from}\PYG{+w}{ }\PYG{n+nn}{scipy}\PYG{n+nn}{.}\PYG{n+nn}{stats}\PYG{+w}{ }\PYG{k+kn}{import} \PYG{n}{multivariate\PYGZus{}normal} \PYG{k}{as} \PYG{n}{mvn}

\PYG{n}{mean} \PYG{o}{=} \PYG{n}{np}\PYG{o}{.}\PYG{n}{array}\PYG{p}{(}\PYG{p}{[}\PYG{l+m+mi}{2}\PYG{p}{,} \PYG{l+m+mi}{3}\PYG{p}{]}\PYG{p}{)}
\PYG{n}{cov} \PYG{o}{=} \PYG{n}{np}\PYG{o}{.}\PYG{n}{array}\PYG{p}{(}\PYG{p}{[}\PYG{p}{[}\PYG{l+m+mi}{1}\PYG{p}{,} \PYG{l+m+mf}{0.8}\PYG{p}{]}\PYG{p}{,} \PYG{p}{[}\PYG{l+m+mf}{0.8}\PYG{p}{,} \PYG{l+m+mi}{1}\PYG{p}{]}\PYG{p}{]}\PYG{p}{)}
\PYG{n}{n\PYGZus{}samples} \PYG{o}{=} \PYG{l+m+mi}{1500}

\PYG{n}{data} \PYG{o}{=} \PYG{n}{mvn}\PYG{p}{(}\PYG{n}{mean}\PYG{o}{=}\PYG{n}{mean}\PYG{p}{,} \PYG{n}{cov}\PYG{o}{=}\PYG{n}{cov}\PYG{p}{)}\PYG{o}{.}\PYG{n}{rvs}\PYG{p}{(}\PYG{n}{size}\PYG{o}{=}\PYG{n}{n\PYGZus{}samples}\PYG{p}{)}

\PYG{n}{sample\PYGZus{}mean} \PYG{o}{=} \PYG{n}{np}\PYG{o}{.}\PYG{n}{mean}\PYG{p}{(}\PYG{n}{data}\PYG{p}{,} \PYG{n}{axis}\PYG{o}{=}\PYG{l+m+mi}{0}\PYG{p}{)}
\PYG{n}{plt}\PYG{o}{.}\PYG{n}{figure}\PYG{p}{(}\PYG{n}{figsize}\PYG{o}{=}\PYG{p}{(}\PYG{l+m+mi}{10}\PYG{p}{,} \PYG{l+m+mi}{6}\PYG{p}{)}\PYG{p}{)}
\PYG{n}{plt}\PYG{o}{.}\PYG{n}{scatter}\PYG{p}{(}\PYG{n}{data}\PYG{p}{[}\PYG{p}{:}\PYG{p}{,} \PYG{l+m+mi}{0}\PYG{p}{]}\PYG{p}{,} \PYG{n}{data}\PYG{p}{[}\PYG{p}{:}\PYG{p}{,} \PYG{l+m+mi}{1}\PYG{p}{]}\PYG{p}{,} \PYG{n}{alpha}\PYG{o}{=}\PYG{l+m+mf}{0.5}\PYG{p}{)}
\PYG{n}{plt}\PYG{o}{.}\PYG{n}{title}\PYG{p}{(}\PYG{l+s+s2}{\PYGZdq{}}\PYG{l+s+s2}{Scatter Plot of Generated Data}\PYG{l+s+s2}{\PYGZdq{}}\PYG{p}{)}
\PYG{n}{plt}\PYG{o}{.}\PYG{n}{xlabel}\PYG{p}{(}\PYG{l+s+s2}{\PYGZdq{}}\PYG{l+s+s2}{Dimension 1}\PYG{l+s+s2}{\PYGZdq{}}\PYG{p}{)}
\PYG{n}{plt}\PYG{o}{.}\PYG{n}{ylabel}\PYG{p}{(}\PYG{l+s+s2}{\PYGZdq{}}\PYG{l+s+s2}{Dimension 2}\PYG{l+s+s2}{\PYGZdq{}}\PYG{p}{)}
\PYG{n}{plt}\PYG{o}{.}\PYG{n}{grid}\PYG{p}{(}\PYG{l+s+s2}{\PYGZdq{}}\PYG{l+s+s2}{on}\PYG{l+s+s2}{\PYGZdq{}}\PYG{p}{)}
\PYG{n}{plt}\PYG{o}{.}\PYG{n}{show}\PYG{p}{(}\PYG{p}{)}

\PYG{n}{sample\PYGZus{}cov} \PYG{o}{=} \PYG{n}{np}\PYG{o}{.}\PYG{n}{cov}\PYG{p}{(}\PYG{n}{data}\PYG{p}{,} \PYG{n}{rowvar}\PYG{o}{=}\PYG{k+kc}{False}\PYG{p}{)}

\PYG{n}{pprint}\PYG{p}{(}\PYG{l+s+s2}{\PYGZdq{}}\PYG{l+s+se}{\PYGZbs{}\PYGZbs{}}\PYG{l+s+s2}{mu=}\PYG{l+s+s2}{\PYGZdq{}}\PYG{p}{,} \PYG{n}{mean}\PYG{p}{)}
\PYG{n}{pprint}\PYG{p}{(}\PYG{l+s+s2}{\PYGZdq{}}\PYG{l+s+se}{\PYGZbs{}\PYGZbs{}}\PYG{l+s+s2}{hat}\PYG{l+s+s2}{\PYGZob{}}\PYG{l+s+se}{\PYGZbs{}\PYGZbs{}}\PYG{l+s+s2}{mu\PYGZcb{}=}\PYG{l+s+s2}{\PYGZdq{}}\PYG{p}{,} \PYG{n}{sample\PYGZus{}mean}\PYG{p}{)}
\PYG{n}{pprint}\PYG{p}{(}\PYG{l+s+s2}{\PYGZdq{}}\PYG{l+s+s2}{Cov:}\PYG{l+s+s2}{\PYGZdq{}}\PYG{p}{,} \PYG{n}{np}\PYG{o}{.}\PYG{n}{array}\PYG{p}{(}\PYG{n}{cov}\PYG{p}{)}\PYG{p}{)}
\PYG{n}{pprint}\PYG{p}{(}\PYG{l+s+s2}{\PYGZdq{}}\PYG{l+s+s2}{MC}\PYG{l+s+s2}{\PYGZbs{}}\PYG{l+s+s2}{ Cov:}\PYG{l+s+s2}{\PYGZdq{}}\PYG{p}{,} \PYG{n}{sample\PYGZus{}cov}\PYG{p}{)}
\end{sphinxVerbatim}

\end{sphinxuseclass}\end{sphinxVerbatimInput}
\begin{sphinxVerbatimOutput}

\begin{sphinxuseclass}{cell_output}
\noindent\sphinxincludegraphics{{22a0568a5d3418b46d1efe7645e2c83b4c3cd707b9daaa40ea3c7c1e45a238c4}.png}
\begin{equation*}
\begin{split}\displaystyle \mu=\left[
\begin{array}{}
  2.0000 &  3.0000
\end{array}
\right]\end{split}
\end{equation*}\begin{equation*}
\begin{split}\displaystyle \hat{\mu}=\left[
\begin{array}{}
  1.9620 &  2.9493
\end{array}
\right]\end{split}
\end{equation*}\begin{equation*}
\begin{split}\displaystyle Cov:\left[
\begin{array}{}
  1.0000 &  0.8000\\
  0.8000 &  1.0000
\end{array}
\right]\end{split}
\end{equation*}\begin{equation*}
\begin{split}\displaystyle MC\ Cov:\left[
\begin{array}{}
  0.9969 &  0.7940\\
  0.7940 &  0.9862
\end{array}
\right]\end{split}
\end{equation*}
\end{sphinxuseclass}\end{sphinxVerbatimOutput}

\end{sphinxuseclass}

\subsection{Exercise:}
\label{\detokenize{01:exercise}}
\sphinxAtStartPar
We define a bivariate gaussian distribution \((\mathbf Z_1, \mathbf Z_2)\) with the following parameters:
\begin{itemize}
\item {} 
\sphinxAtStartPar
Mean vector: \([0.25, 1.0]\)

\item {} 
\sphinxAtStartPar
Covariance matrix: \(\begin{bmatrix}1.0&0.5\\0.5&1.0\end{bmatrix}\)

\end{itemize}

\sphinxAtStartPar
Plot the regressions of \(\mathbf Z_1|\mathbf Z_2\) and \(\mathbf Z_2|\mathbf Z_1\). What is their intersection ? Hint: those are given by expressing \(\mathbf Z_1\) in terms of \(\mathbf Z_2\) (or the opposite): \(\mathbf Z_1 = a_1 + b_1\mathbf Z_2\).


\subsubsection{Solution}
\label{\detokenize{01:id9}}
\sphinxAtStartPar
For the following derivation, we redirect the student to the theoretical handbook, pages 64\sphinxhyphen{}65. The following code generates a set of samples from a bivariate normal distribution with mean vector \(\mu = \begin{bmatrix} 0.25 \\ 1 \end{bmatrix}\) and covariance matrix \(\Sigma = \begin{bmatrix} 1 & 0.5 \\ 0.5 & 1 \end{bmatrix}\), then computes and visualizes the regression lines corresponding to the conditional expectations \(\mathbb{E}[\mathbf Z_1 \mid \mathbf Z_2]\) and \(\mathbb{E}[\mathbf Z_2 \mid \mathbf Z_1]\). The regression coefficients are computed as \(b_1 = \frac{\Sigma_{12}}{\Sigma_{22}}\) and \(b_2 = \frac{\Sigma_{12}}{\Sigma_{11}}\), with intercepts \(a_1 = \mu_1 - b_1 \mu_2\) and \(a_2 = \mu_2 - b_2 \mu_1\). The resulting lines are plotted on the \((\mathbf Z_2, \mathbf Z_1)\) plane, showing the best linear predictors of each variable given the other.

\begin{sphinxuseclass}{cell}\begin{sphinxVerbatimInput}

\begin{sphinxuseclass}{cell_input}
\begin{sphinxVerbatim}[commandchars=\\\{\}]
\PYG{n}{mean} \PYG{o}{=} \PYG{n}{np}\PYG{o}{.}\PYG{n}{array}\PYG{p}{(}\PYG{p}{[}\PYG{l+m+mf}{0.25}\PYG{p}{,} \PYG{l+m+mf}{1.}\PYG{p}{]}\PYG{p}{)}
\PYG{n}{cov} \PYG{o}{=} \PYG{n}{np}\PYG{o}{.}\PYG{n}{array}\PYG{p}{(}\PYG{p}{[}\PYG{p}{[}\PYG{l+m+mf}{1.}\PYG{p}{,} \PYG{l+m+mf}{0.5}\PYG{p}{]}\PYG{p}{,} \PYG{p}{[}\PYG{l+m+mf}{0.5}\PYG{p}{,} \PYG{l+m+mf}{1.}\PYG{p}{]}\PYG{p}{]}\PYG{p}{)}

\PYG{n}{data} \PYG{o}{=} \PYG{n}{mvn}\PYG{p}{(}\PYG{n}{mean}\PYG{o}{=}\PYG{n}{mean}\PYG{p}{,} \PYG{n}{cov}\PYG{o}{=}\PYG{n}{cov}\PYG{p}{)}\PYG{o}{.}\PYG{n}{rvs}\PYG{p}{(}\PYG{n}{n\PYGZus{}points}\PYG{p}{)}

\PYG{n}{b\PYGZus{}1} \PYG{o}{=} \PYG{n}{cov}\PYG{p}{[}\PYG{l+m+mi}{0}\PYG{p}{,} \PYG{l+m+mi}{1}\PYG{p}{]} \PYG{o}{/} \PYG{n}{cov}\PYG{p}{[}\PYG{l+m+mi}{1}\PYG{p}{,} \PYG{l+m+mi}{1}\PYG{p}{]}
\PYG{n}{b\PYGZus{}2} \PYG{o}{=} \PYG{n}{cov}\PYG{p}{[}\PYG{l+m+mi}{0}\PYG{p}{,} \PYG{l+m+mi}{1}\PYG{p}{]} \PYG{o}{/} \PYG{n}{cov}\PYG{p}{[}\PYG{l+m+mi}{0}\PYG{p}{,} \PYG{l+m+mi}{0}\PYG{p}{]}

\PYG{n}{a1} \PYG{o}{=} \PYG{n}{mean}\PYG{p}{[}\PYG{l+m+mi}{0}\PYG{p}{]} \PYG{o}{\PYGZhy{}} \PYG{n}{b\PYGZus{}1} \PYG{o}{*} \PYG{n}{mean}\PYG{p}{[}\PYG{l+m+mi}{1}\PYG{p}{]}
\PYG{n}{a2} \PYG{o}{=} \PYG{n}{mean}\PYG{p}{[}\PYG{l+m+mi}{1}\PYG{p}{]} \PYG{o}{\PYGZhy{}} \PYG{n}{b\PYGZus{}2} \PYG{o}{*} \PYG{n}{mean}\PYG{p}{[}\PYG{l+m+mi}{0}\PYG{p}{]}

\PYG{n}{pprint}\PYG{p}{(}\PYG{l+s+s2}{\PYGZdq{}}\PYG{l+s+se}{\PYGZbs{}\PYGZbs{}}\PYG{l+s+s2}{text}\PYG{l+s+s2}{\PYGZob{}}\PYG{l+s+s2}{Regression of \PYGZcb{}Z\PYGZus{}1}\PYG{l+s+se}{\PYGZbs{}\PYGZbs{}}\PYG{l+s+s2}{text}\PYG{l+s+s2}{\PYGZob{}}\PYG{l+s+s2}{ on \PYGZcb{}Z\PYGZus{}2: a\PYGZus{}1 =}\PYG{l+s+s2}{\PYGZdq{}}\PYG{p}{,} \PYG{n+nb}{str}\PYG{p}{(}\PYG{n}{a1}\PYG{p}{)}\PYG{p}{,} \PYG{l+s+s2}{\PYGZdq{}}\PYG{l+s+s2}{,}\PYG{l+s+s2}{\PYGZbs{}}\PYG{l+s+s2}{ b\PYGZus{}1 =}\PYG{l+s+s2}{\PYGZdq{}}\PYG{p}{,} \PYG{n+nb}{str}\PYG{p}{(}\PYG{n}{b\PYGZus{}1}\PYG{p}{)}\PYG{p}{)}
\PYG{n}{pprint}\PYG{p}{(}\PYG{l+s+s2}{\PYGZdq{}}\PYG{l+s+se}{\PYGZbs{}\PYGZbs{}}\PYG{l+s+s2}{text}\PYG{l+s+s2}{\PYGZob{}}\PYG{l+s+s2}{Regression of \PYGZcb{}Z\PYGZus{}2}\PYG{l+s+se}{\PYGZbs{}\PYGZbs{}}\PYG{l+s+s2}{text}\PYG{l+s+s2}{\PYGZob{}}\PYG{l+s+s2}{ on \PYGZcb{}Z\PYGZus{}1: a\PYGZus{}2 =}\PYG{l+s+s2}{\PYGZdq{}}\PYG{p}{,} \PYG{n+nb}{str}\PYG{p}{(}\PYG{n}{a2}\PYG{p}{)}\PYG{p}{,} \PYG{l+s+s2}{\PYGZdq{}}\PYG{l+s+s2}{,}\PYG{l+s+s2}{\PYGZbs{}}\PYG{l+s+s2}{ b\PYGZus{}2 =}\PYG{l+s+s2}{\PYGZdq{}}\PYG{p}{,} \PYG{n+nb}{str}\PYG{p}{(}\PYG{n}{b\PYGZus{}2}\PYG{p}{)}\PYG{p}{)}

\PYG{n}{z2} \PYG{o}{=} \PYG{n}{np}\PYG{o}{.}\PYG{n}{linspace}\PYG{p}{(}\PYG{o}{\PYGZhy{}}\PYG{l+m+mi}{4}\PYG{p}{,} \PYG{l+m+mi}{4}\PYG{p}{,} \PYG{l+m+mi}{100}\PYG{p}{)}

\PYG{n}{z1\PYGZus{}pred} \PYG{o}{=} \PYG{n}{a1} \PYG{o}{+} \PYG{n}{b\PYGZus{}1} \PYG{o}{*} \PYG{n}{z2}
\PYG{n}{z2\PYGZus{}pred} \PYG{o}{=} \PYG{p}{(}\PYG{n}{z2} \PYG{o}{\PYGZhy{}} \PYG{n}{a2}\PYG{p}{)} \PYG{o}{/} \PYG{n}{b\PYGZus{}2}

\PYG{n}{fig}\PYG{p}{,} \PYG{n}{ax} \PYG{o}{=} \PYG{n}{plt}\PYG{o}{.}\PYG{n}{subplots}\PYG{p}{(}\PYG{n}{figsize}\PYG{o}{=}\PYG{p}{(}\PYG{l+m+mi}{10}\PYG{p}{,} \PYG{l+m+mi}{10}\PYG{p}{)}\PYG{p}{)}
\PYG{n}{ax}\PYG{o}{.}\PYG{n}{set\PYGZus{}xlim}\PYG{p}{(}\PYG{o}{\PYGZhy{}}\PYG{l+m+mi}{4}\PYG{p}{,} \PYG{l+m+mi}{4}\PYG{p}{)}
\PYG{n}{ax}\PYG{o}{.}\PYG{n}{set\PYGZus{}ylim}\PYG{p}{(}\PYG{o}{\PYGZhy{}}\PYG{l+m+mi}{4}\PYG{p}{,} \PYG{l+m+mi}{4}\PYG{p}{)}
\PYG{n}{ax}\PYG{o}{.}\PYG{n}{axhline}\PYG{p}{(}\PYG{l+m+mi}{0}\PYG{p}{,} \PYG{n}{color}\PYG{o}{=}\PYG{l+s+s1}{\PYGZsq{}}\PYG{l+s+s1}{black}\PYG{l+s+s1}{\PYGZsq{}}\PYG{p}{,} \PYG{n}{lw}\PYG{o}{=}\PYG{l+m+mi}{1}\PYG{p}{)}
\PYG{n}{ax}\PYG{o}{.}\PYG{n}{axvline}\PYG{p}{(}\PYG{l+m+mi}{0}\PYG{p}{,} \PYG{n}{color}\PYG{o}{=}\PYG{l+s+s1}{\PYGZsq{}}\PYG{l+s+s1}{black}\PYG{l+s+s1}{\PYGZsq{}}\PYG{p}{,} \PYG{n}{lw}\PYG{o}{=}\PYG{l+m+mi}{1}\PYG{p}{)}

\PYG{n}{plt}\PYG{o}{.}\PYG{n}{plot}\PYG{p}{(}\PYG{n}{z2}\PYG{p}{,} \PYG{n}{z1\PYGZus{}pred}\PYG{p}{,} \PYG{n}{label}\PYG{o}{=}\PYG{l+s+s2}{\PYGZdq{}}\PYG{l+s+s2}{\PYGZdl{}E[Z\PYGZus{}1 | Z\PYGZus{}2]\PYGZdl{}}\PYG{l+s+s2}{\PYGZdq{}}\PYG{p}{)}
\PYG{n}{plt}\PYG{o}{.}\PYG{n}{plot}\PYG{p}{(}\PYG{n}{z2}\PYG{p}{,} \PYG{n}{z2\PYGZus{}pred}\PYG{p}{,} \PYG{n}{label}\PYG{o}{=}\PYG{l+s+s2}{\PYGZdq{}}\PYG{l+s+s2}{\PYGZdl{}E[Z\PYGZus{}2 | Z\PYGZus{}1]\PYGZdl{}}\PYG{l+s+s2}{\PYGZdq{}}\PYG{p}{)}
\PYG{n}{plt}\PYG{o}{.}\PYG{n}{xlabel}\PYG{p}{(}\PYG{l+s+s2}{\PYGZdq{}}\PYG{l+s+s2}{\PYGZdl{}Z\PYGZus{}2\PYGZdl{}}\PYG{l+s+s2}{\PYGZdq{}}\PYG{p}{)}
\PYG{n}{plt}\PYG{o}{.}\PYG{n}{ylabel}\PYG{p}{(}\PYG{l+s+s2}{\PYGZdq{}}\PYG{l+s+s2}{\PYGZdl{}Z\PYGZus{}1\PYGZdl{}}\PYG{l+s+s2}{\PYGZdq{}}\PYG{p}{)}
\PYG{n}{plt}\PYG{o}{.}\PYG{n}{legend}\PYG{p}{(}\PYG{p}{)}
\PYG{n}{plt}\PYG{o}{.}\PYG{n}{grid}\PYG{p}{(}\PYG{l+s+s2}{\PYGZdq{}}\PYG{l+s+s2}{on}\PYG{l+s+s2}{\PYGZdq{}}\PYG{p}{)}
\PYG{n}{plt}\PYG{o}{.}\PYG{n}{title}\PYG{p}{(}\PYG{l+s+s2}{\PYGZdq{}}\PYG{l+s+s2}{Regression Lines for Bivariate Normal Distribution}\PYG{l+s+s2}{\PYGZdq{}}\PYG{p}{)}
\PYG{n}{plt}\PYG{o}{.}\PYG{n}{show}\PYG{p}{(}\PYG{p}{)}
\end{sphinxVerbatim}

\end{sphinxuseclass}\end{sphinxVerbatimInput}
\begin{sphinxVerbatimOutput}

\begin{sphinxuseclass}{cell_output}\begin{equation*}
\begin{split}\displaystyle \text{Regression of }Z_1\text{ on }Z_2: a_1 =-0.25,\ b_1 =0.5\end{split}
\end{equation*}\begin{equation*}
\begin{split}\displaystyle \text{Regression of }Z_2\text{ on }Z_1: a_2 =0.875,\ b_2 =0.5\end{split}
\end{equation*}
\noindent\sphinxincludegraphics{{edef9cfedf97e72c5f46c36a9b411b5a3314ff0ba06cbeffeb00d8e34242f14d}.png}

\end{sphinxuseclass}\end{sphinxVerbatimOutput}

\end{sphinxuseclass}

\subsection{Exercise: Trivariate Gaussian Data}
\label{\detokenize{01:exercise-trivariate-gaussian-data}}
\sphinxAtStartPar
Finally, generate 1000 samples of 3\sphinxhyphen{}dimensional Gaussian random data with the following parameters:
\begin{itemize}
\item {} 
\sphinxAtStartPar
Mean vector: \([1, 4, 7]\)

\item {} 
\sphinxAtStartPar
Covariance matrix:\(\begin{bmatrix}3&1.5&0.8\\1.5&2&0.5\\0.8&0.5&1\end{bmatrix}\)

\end{itemize}

\sphinxAtStartPar
Estimate the covariance matrix using the generated data. Visualize the pairwise scatter plots of the dimensions (use a scatter plot matrix).
Hints: use \sphinxcode{\sphinxupquote{pandas.plotting.scatter\_matrix}} or \sphinxcode{\sphinxupquote{seaborn.pairplot}} for visualization.
Pay attention to the interpretation of correlations between dimensions.


\subsubsection{Solution}
\label{\detokenize{01:id10}}
\sphinxAtStartPar
The following code generates a 3\sphinxhyphen{}dimensional dataset of 1000 samples from a multivariate normal distribution with a specified mean vector \([1, 4, 7]\) and covariance matrix \(\begin{bmatrix} 3 & 1.5 & 0.8 \\ 1.5 & 2 & 0.5 \\ 0.8 & 0.5 & 1 \end{bmatrix}\). It stores the generated samples in a Pandas DataFrame with columns labeled “Dimension 1”, “Dimension 2”, and “Dimension 3”. It then computes the empirical covariance matrix of the generated data using \sphinxcode{\sphinxupquote{np.cov}}, prints both the given and estimated covariance matrices, and visualizes the pairwise relationships between the three dimensions using \sphinxcode{\sphinxupquote{Seaborn}}’s \sphinxcode{\sphinxupquote{pairplot}}.

\begin{sphinxuseclass}{cell}\begin{sphinxVerbatimInput}

\begin{sphinxuseclass}{cell_input}
\begin{sphinxVerbatim}[commandchars=\\\{\}]
\PYG{k+kn}{import}\PYG{+w}{ }\PYG{n+nn}{pandas}\PYG{+w}{ }\PYG{k}{as}\PYG{+w}{ }\PYG{n+nn}{pd}
\PYG{k+kn}{import}\PYG{+w}{ }\PYG{n+nn}{seaborn}\PYG{+w}{ }\PYG{k}{as}\PYG{+w}{ }\PYG{n+nn}{sns}

\PYG{n}{mean\PYGZus{}3d} \PYG{o}{=} \PYG{p}{[}\PYG{l+m+mi}{1}\PYG{p}{,} \PYG{l+m+mi}{4}\PYG{p}{,} \PYG{l+m+mi}{7}\PYG{p}{]}
\PYG{n}{cov\PYGZus{}3d} \PYG{o}{=} \PYG{p}{[}\PYG{p}{[}\PYG{l+m+mi}{3}\PYG{p}{,} \PYG{l+m+mf}{1.5}\PYG{p}{,} \PYG{l+m+mf}{0.8}\PYG{p}{]}\PYG{p}{,} 
          \PYG{p}{[}\PYG{l+m+mf}{1.5}\PYG{p}{,} \PYG{l+m+mi}{2}\PYG{p}{,} \PYG{l+m+mf}{0.5}\PYG{p}{]}\PYG{p}{,} 
          \PYG{p}{[}\PYG{l+m+mf}{0.8}\PYG{p}{,} \PYG{l+m+mf}{0.5}\PYG{p}{,} \PYG{l+m+mi}{1}\PYG{p}{]}\PYG{p}{]}

\PYG{n}{n\PYGZus{}samples\PYGZus{}3d} \PYG{o}{=} \PYG{l+m+mi}{1000}

\PYG{n}{data\PYGZus{}3d} \PYG{o}{=} \PYG{n}{mvn}\PYG{p}{(}\PYG{n}{mean}\PYG{o}{=}\PYG{n}{mean\PYGZus{}3d}\PYG{p}{,} \PYG{n}{cov}\PYG{o}{=}\PYG{n}{cov\PYGZus{}3d}\PYG{p}{)}\PYG{o}{.}\PYG{n}{rvs}\PYG{p}{(}\PYG{n}{size}\PYG{o}{=}\PYG{n}{n\PYGZus{}samples\PYGZus{}3d}\PYG{p}{)}

\PYG{n}{df} \PYG{o}{=} \PYG{n}{pd}\PYG{o}{.}\PYG{n}{DataFrame}\PYG{p}{(}\PYG{n}{data\PYGZus{}3d}\PYG{p}{,} \PYG{n}{columns}\PYG{o}{=}\PYG{p}{[}\PYG{l+s+s2}{\PYGZdq{}}\PYG{l+s+s2}{Dimension 1}\PYG{l+s+s2}{\PYGZdq{}}\PYG{p}{,} \PYG{l+s+s2}{\PYGZdq{}}\PYG{l+s+s2}{Dimension 2}\PYG{l+s+s2}{\PYGZdq{}}\PYG{p}{,} \PYG{l+s+s2}{\PYGZdq{}}\PYG{l+s+s2}{Dimension 3}\PYG{l+s+s2}{\PYGZdq{}}\PYG{p}{]}\PYG{p}{)}

\PYG{n}{sample\PYGZus{}cov\PYGZus{}3d} \PYG{o}{=} \PYG{n}{np}\PYG{o}{.}\PYG{n}{cov}\PYG{p}{(}\PYG{n}{data\PYGZus{}3d}\PYG{p}{,} \PYG{n}{rowvar}\PYG{o}{=}\PYG{k+kc}{False}\PYG{p}{)}

\PYG{n}{pprint}\PYG{p}{(}\PYG{l+s+s2}{\PYGZdq{}}\PYG{l+s+s2}{Cov:}\PYG{l+s+s2}{\PYGZdq{}}\PYG{p}{,} \PYG{n}{np}\PYG{o}{.}\PYG{n}{array}\PYG{p}{(}\PYG{n}{cov\PYGZus{}3d}\PYG{p}{)}\PYG{p}{)}
\PYG{n}{pprint}\PYG{p}{(}\PYG{l+s+s2}{\PYGZdq{}}\PYG{l+s+s2}{MC}\PYG{l+s+s2}{\PYGZbs{}}\PYG{l+s+s2}{ Cov:}\PYG{l+s+s2}{\PYGZdq{}}\PYG{p}{,} \PYG{n}{sample\PYGZus{}cov\PYGZus{}3d}\PYG{p}{)}

\PYG{n}{sns}\PYG{o}{.}\PYG{n}{pairplot}\PYG{p}{(}\PYG{n}{df}\PYG{p}{)}
\PYG{n}{plt}\PYG{o}{.}\PYG{n}{show}\PYG{p}{(}\PYG{p}{)}
\end{sphinxVerbatim}

\end{sphinxuseclass}\end{sphinxVerbatimInput}
\begin{sphinxVerbatimOutput}

\begin{sphinxuseclass}{cell_output}\begin{equation*}
\begin{split}\displaystyle Cov:\left[
\begin{array}{}
  3.0000 &  1.5000 &  0.8000\\
  1.5000 &  2.0000 &  0.5000\\
  0.8000 &  0.5000 &  1.0000
\end{array}
\right]\end{split}
\end{equation*}\begin{equation*}
\begin{split}\displaystyle MC\ Cov:\left[
\begin{array}{}
  3.1791 &  1.6075 &  0.9758\\
  1.6075 &  2.0246 &  0.6169\\
  0.9758 &  0.6169 &  1.0465
\end{array}
\right]\end{split}
\end{equation*}
\noindent\sphinxincludegraphics{{024fd0dc3544d213a9cddf643aba1af08c5051d4e0a4cd6db708ec6059525a94}.png}

\end{sphinxuseclass}\end{sphinxVerbatimOutput}

\end{sphinxuseclass}

\subsection{Exercise: Partial Correlation}
\label{\detokenize{01:exercise-partial-correlation}}
\sphinxAtStartPar
Given two normal \(\mathcal{N}(0, 1)\) r.v. \(Z_1, Z_2\) and a third r.v. \(Y = Z_1 + Z_2\), compute the correlation between \(Z_1\) and \(Z_2\), then compute the partial correlation between \(Z_1\) and \(Z_2\) given \(Y\). Hint: use the function \sphinxcode{\sphinxupquote{corrcoef}} of \sphinxcode{\sphinxupquote{Numpy}} to compute the correlation coefficients. What happens to the partial correlation between \(Z_1\) and \(Z_2\) given \(Y\) if now  \(Y = Z_1 + Z_2 + Z_3\), for \(Z_3\) another \(\mathcal{N}(0, 1)\) r.v. ?


\subsubsection{Solution}
\label{\detokenize{01:id11}}
\sphinxAtStartPar
This code defines a function \sphinxcode{\sphinxupquote{partial\_corr(X, Y, Z)}} that computes the partial correlation between two variables \(\mathbf X\) and \(\mathbf Y\) given a third variable \(\mathbf Z\), using the formula \(\rho_{\mathbf X\mathbf Y \cdot \mathbf Z} = \frac{\rho_{\mathbf X\mathbf Y} - \rho_{\mathbf X\mathbf Z}\rho_{\mathbf Y\mathbf Z}}{\sqrt{(1 - \rho_{\mathbf X\mathbf Z}^2)(1 - \rho_{\mathbf Y\mathbf Z}^2)}}\). It then generates one million samples of three independent standard normal variables \(\mathbf Z_1, \mathbf Z_2, \mathbf Z_3\), defines \(\mathbf Y = \mathbf Z_1 + \mathbf Z_2\), and prints the empirical correlation between \(\mathbf Z_1\) and \(\mathbf Z_2\), as well as the partial correlations between \(\mathbf Z_1\) and \(\mathbf Z_2\) given \(\mathbf Y\) and given \(\mathbf Y + \mathbf Z_3\), illustrating that the dependence between \(\mathbf Z_1\) and \(\mathbf Z_2\) vanishes when conditioning on their sum.

\begin{sphinxuseclass}{cell}\begin{sphinxVerbatimInput}

\begin{sphinxuseclass}{cell_input}
\begin{sphinxVerbatim}[commandchars=\\\{\}]
\PYG{k}{def}\PYG{+w}{ }\PYG{n+nf}{partial\PYGZus{}corr}\PYG{p}{(}\PYG{n}{X}\PYG{p}{,}\PYG{n}{Y}\PYG{p}{,}\PYG{n}{Z}\PYG{p}{)}\PYG{p}{:}
    \PYG{n}{rhoXY} \PYG{o}{=} \PYG{n}{np}\PYG{o}{.}\PYG{n}{corrcoef}\PYG{p}{(}\PYG{n}{X}\PYG{o}{.}\PYG{n}{T}\PYG{p}{,}\PYG{n}{Y}\PYG{o}{.}\PYG{n}{T}\PYG{p}{)}\PYG{p}{[}\PYG{l+m+mi}{0}\PYG{p}{,}\PYG{l+m+mi}{1}\PYG{p}{]}
    \PYG{n}{rhoXZ} \PYG{o}{=} \PYG{n}{np}\PYG{o}{.}\PYG{n}{corrcoef}\PYG{p}{(}\PYG{n}{X}\PYG{o}{.}\PYG{n}{T}\PYG{p}{,}\PYG{n}{Z}\PYG{o}{.}\PYG{n}{T}\PYG{p}{)}\PYG{p}{[}\PYG{l+m+mi}{0}\PYG{p}{,}\PYG{l+m+mi}{1}\PYG{p}{]}
    \PYG{n}{rhoYZ} \PYG{o}{=} \PYG{n}{np}\PYG{o}{.}\PYG{n}{corrcoef}\PYG{p}{(}\PYG{n}{Y}\PYG{o}{.}\PYG{n}{T}\PYG{p}{,}\PYG{n}{Z}\PYG{o}{.}\PYG{n}{T}\PYG{p}{)}\PYG{p}{[}\PYG{l+m+mi}{0}\PYG{p}{,}\PYG{l+m+mi}{1}\PYG{p}{]}    
    \PYG{k}{return} \PYG{p}{(}\PYG{n}{rhoXY}\PYG{o}{\PYGZhy{}}\PYG{n}{rhoXZ}\PYG{o}{*}\PYG{n}{rhoYZ}\PYG{p}{)}\PYG{o}{/}\PYG{p}{(}\PYG{n}{np}\PYG{o}{.}\PYG{n}{sqrt}\PYG{p}{(}\PYG{p}{(}\PYG{l+m+mi}{1}\PYG{o}{\PYGZhy{}}\PYG{n}{rhoXZ}\PYG{o}{*}\PYG{o}{*}\PYG{l+m+mi}{2}\PYG{p}{)}\PYG{o}{*}\PYG{p}{(}\PYG{l+m+mi}{1}\PYG{o}{\PYGZhy{}}\PYG{n}{rhoYZ}\PYG{o}{*}\PYG{o}{*}\PYG{l+m+mi}{2}\PYG{p}{)}\PYG{p}{)}\PYG{p}{)}

\PYG{n}{N} \PYG{o}{=} \PYG{n+nb}{int}\PYG{p}{(}\PYG{l+m+mf}{1e6}\PYG{p}{)}

\PYG{n}{Z1} \PYG{o}{=} \PYG{n}{np}\PYG{o}{.}\PYG{n}{random}\PYG{o}{.}\PYG{n}{normal}\PYG{p}{(}\PYG{n}{size}\PYG{o}{=}\PYG{p}{(}\PYG{n}{N}\PYG{p}{,} \PYG{l+m+mi}{1}\PYG{p}{)}\PYG{p}{)}
\PYG{n}{Z2} \PYG{o}{=} \PYG{n}{np}\PYG{o}{.}\PYG{n}{random}\PYG{o}{.}\PYG{n}{normal}\PYG{p}{(}\PYG{n}{size}\PYG{o}{=}\PYG{p}{(}\PYG{n}{N}\PYG{p}{,} \PYG{l+m+mi}{1}\PYG{p}{)}\PYG{p}{)}
\PYG{n}{Z3} \PYG{o}{=} \PYG{n}{np}\PYG{o}{.}\PYG{n}{random}\PYG{o}{.}\PYG{n}{normal}\PYG{p}{(}\PYG{n}{size}\PYG{o}{=}\PYG{p}{(}\PYG{n}{N}\PYG{p}{,} \PYG{l+m+mi}{1}\PYG{p}{)}\PYG{p}{)}

\PYG{n}{Y} \PYG{o}{=} \PYG{n}{Z1} \PYG{o}{+} \PYG{n}{Z2}

\PYG{n}{pprint}\PYG{p}{(}\PYG{l+s+s2}{\PYGZdq{}}\PYG{l+s+se}{\PYGZbs{}\PYGZbs{}}\PYG{l+s+s2}{text}\PYG{l+s+s2}{\PYGZob{}}\PYG{l+s+s2}{Correlation between \PYGZcb{}}\PYG{l+s+se}{\PYGZbs{}\PYGZbs{}}\PYG{l+s+s2}{mathbf Z\PYGZus{}1}\PYG{l+s+se}{\PYGZbs{}\PYGZbs{}}\PYG{l+s+s2}{text}\PYG{l+s+s2}{\PYGZob{}}\PYG{l+s+s2}{ and \PYGZcb{}}\PYG{l+s+se}{\PYGZbs{}\PYGZbs{}}\PYG{l+s+s2}{mathbf Z\PYGZus{}2: }\PYG{l+s+s2}{\PYGZdq{}}\PYG{p}{,} \PYG{n+nb}{str}\PYG{p}{(}\PYG{n}{np}\PYG{o}{.}\PYG{n}{corrcoef}\PYG{p}{(}\PYG{n}{Z1}\PYG{o}{.}\PYG{n}{T}\PYG{p}{,}\PYG{n}{Z2}\PYG{o}{.}\PYG{n}{T}\PYG{p}{)}\PYG{p}{[}\PYG{l+m+mi}{0}\PYG{p}{,} \PYG{l+m+mi}{1}\PYG{p}{]}\PYG{p}{)}\PYG{p}{[}\PYG{p}{:}\PYG{l+m+mi}{8}\PYG{p}{]}\PYG{p}{)}
\PYG{n}{pprint}\PYG{p}{(}\PYG{l+s+s2}{\PYGZdq{}}\PYG{l+s+se}{\PYGZbs{}\PYGZbs{}}\PYG{l+s+s2}{text}\PYG{l+s+s2}{\PYGZob{}}\PYG{l+s+s2}{Partial Correlation between \PYGZcb{}}\PYG{l+s+se}{\PYGZbs{}\PYGZbs{}}\PYG{l+s+s2}{mathbf Z\PYGZus{}1}\PYG{l+s+se}{\PYGZbs{}\PYGZbs{}}\PYG{l+s+s2}{text}\PYG{l+s+s2}{\PYGZob{}}\PYG{l+s+s2}{ and \PYGZcb{}}\PYG{l+s+se}{\PYGZbs{}\PYGZbs{}}\PYG{l+s+s2}{mathbf Z\PYGZus{}2}\PYG{l+s+se}{\PYGZbs{}\PYGZbs{}}\PYG{l+s+s2}{text}\PYG{l+s+s2}{\PYGZob{}}\PYG{l+s+s2}{ given \PYGZcb{}}\PYG{l+s+se}{\PYGZbs{}\PYGZbs{}}\PYG{l+s+s2}{mathbf Y=}\PYG{l+s+se}{\PYGZbs{}\PYGZbs{}}\PYG{l+s+s2}{mathbf Z\PYGZus{}1+}\PYG{l+s+se}{\PYGZbs{}\PYGZbs{}}\PYG{l+s+s2}{mathbf Z\PYGZus{}2:}\PYG{l+s+s2}{\PYGZdq{}}\PYG{p}{,} \PYG{n+nb}{str}\PYG{p}{(}\PYG{n}{partial\PYGZus{}corr}\PYG{p}{(}\PYG{n}{Z1}\PYG{p}{,} \PYG{n}{Z2}\PYG{p}{,} \PYG{n}{Y}\PYG{p}{)}\PYG{p}{)}\PYG{p}{[}\PYG{p}{:}\PYG{l+m+mi}{8}\PYG{p}{]}\PYG{p}{)}
\PYG{n}{pprint}\PYG{p}{(}\PYG{l+s+s2}{\PYGZdq{}}\PYG{l+s+se}{\PYGZbs{}\PYGZbs{}}\PYG{l+s+s2}{text}\PYG{l+s+s2}{\PYGZob{}}\PYG{l+s+s2}{Partial Correlation between \PYGZcb{}}\PYG{l+s+se}{\PYGZbs{}\PYGZbs{}}\PYG{l+s+s2}{mathbf Z\PYGZus{}1}\PYG{l+s+se}{\PYGZbs{}\PYGZbs{}}\PYG{l+s+s2}{text}\PYG{l+s+s2}{\PYGZob{}}\PYG{l+s+s2}{ and \PYGZcb{}}\PYG{l+s+se}{\PYGZbs{}\PYGZbs{}}\PYG{l+s+s2}{mathbf Z\PYGZus{}2}\PYG{l+s+se}{\PYGZbs{}\PYGZbs{}}\PYG{l+s+s2}{text}\PYG{l+s+s2}{\PYGZob{}}\PYG{l+s+s2}{ given \PYGZcb{}}\PYG{l+s+se}{\PYGZbs{}\PYGZbs{}}\PYG{l+s+s2}{mathbf Y=}\PYG{l+s+se}{\PYGZbs{}\PYGZbs{}}\PYG{l+s+s2}{mathbf Z\PYGZus{}1+}\PYG{l+s+se}{\PYGZbs{}\PYGZbs{}}\PYG{l+s+s2}{mathbf Z\PYGZus{}2+}\PYG{l+s+se}{\PYGZbs{}\PYGZbs{}}\PYG{l+s+s2}{mathbf Z\PYGZus{}3:}\PYG{l+s+s2}{\PYGZdq{}}\PYG{p}{,} \PYG{n+nb}{str}\PYG{p}{(}\PYG{n}{partial\PYGZus{}corr}\PYG{p}{(}\PYG{n}{Z1}\PYG{p}{,} \PYG{n}{Z2}\PYG{p}{,} \PYG{n}{Z1} \PYG{o}{+} \PYG{n}{Z2} \PYG{o}{+} \PYG{n}{Z3}\PYG{p}{)}\PYG{p}{)}\PYG{p}{[}\PYG{p}{:}\PYG{l+m+mi}{8}\PYG{p}{]}\PYG{p}{)}
\end{sphinxVerbatim}

\end{sphinxuseclass}\end{sphinxVerbatimInput}
\begin{sphinxVerbatimOutput}

\begin{sphinxuseclass}{cell_output}\begin{equation*}
\begin{split}\displaystyle \text{Correlation between }\mathbf Z_1\text{ and }\mathbf Z_2: -0.00073\end{split}
\end{equation*}\begin{equation*}
\begin{split}\displaystyle \text{Partial Correlation between }\mathbf Z_1\text{ and }\mathbf Z_2\text{ given }\mathbf Y=\mathbf Z_1+\mathbf Z_2:-0.99999\end{split}
\end{equation*}\begin{equation*}
\begin{split}\displaystyle \text{Partial Correlation between }\mathbf Z_1\text{ and }\mathbf Z_2\text{ given }\mathbf Y=\mathbf Z_1+\mathbf Z_2+\mathbf Z_3:-0.50070\end{split}
\end{equation*}
\end{sphinxuseclass}\end{sphinxVerbatimOutput}

\end{sphinxuseclass}

\section{Minimization}
\label{\detokenize{01:minimization}}
\sphinxAtStartPar
\sphinxcode{\sphinxupquote{Scipy}} offers a convenient tool to perform the minimization of a derivable univariate function: in the module \sphinxcode{\sphinxupquote{optimize}}, we have the \sphinxcode{\sphinxupquote{minimize}} function.


\subsection{Exercise: Minimization}
\label{\detokenize{01:exercise-minimization}}
\sphinxAtStartPar
Define a function
\begin{equation*}
\begin{split}f(x) = x^2 + (\sin{x} + 25)^2\end{split}
\end{equation*}
\sphinxAtStartPar
and use the \sphinxcode{\sphinxupquote{sp.optimize.minimize}} function to find different local minima. Plot the function an show the different minimas. On which parameter of the \sphinxcode{\sphinxupquote{minimize}} function can you play?


\subsubsection{Solution}
\label{\detokenize{01:id12}}
\sphinxAtStartPar
The following code defines a function \(f(x) = x^2 + (\sin(x) + 25)^2\), then plots it over the interval \([-15, 15]\). For several initial guesses \(x_0 \in \{-12, -4, 1, 7, 14\}\), the script uses the \sphinxcode{\sphinxupquote{SciPy}} optimization routine \sphinxcode{\sphinxupquote{minimize}} to find a local minimum of \(f(x)\) starting from \(x_0\). Each minimization path is illustrated as a dashed line connecting the initial point to the corresponding minimum, with different colors representing different initial guesses. The plot includes a legend, grid, and title to clearly display the function and minimization results.

\begin{sphinxuseclass}{cell}\begin{sphinxVerbatimInput}

\begin{sphinxuseclass}{cell_input}
\begin{sphinxVerbatim}[commandchars=\\\{\}]
\PYG{k}{def}\PYG{+w}{ }\PYG{n+nf}{f}\PYG{p}{(}\PYG{n}{x}\PYG{p}{)}\PYG{p}{:}
    \PYG{k}{return} \PYG{p}{(}\PYG{n}{x}\PYG{o}{*}\PYG{o}{*}\PYG{l+m+mi}{2} \PYG{o}{+} \PYG{p}{(}\PYG{n}{np}\PYG{o}{.}\PYG{n}{sin}\PYG{p}{(}\PYG{n}{x}\PYG{p}{)}\PYG{o}{+}\PYG{l+m+mi}{25}\PYG{p}{)}\PYG{o}{*}\PYG{o}{*}\PYG{l+m+mi}{2}\PYG{p}{)}

\PYG{n}{x\PYGZus{}vec} \PYG{o}{=} \PYG{n}{np}\PYG{o}{.}\PYG{n}{linspace}\PYG{p}{(}\PYG{o}{\PYGZhy{}}\PYG{l+m+mi}{15}\PYG{p}{,} \PYG{l+m+mi}{15}\PYG{p}{,} \PYG{l+m+mi}{200}\PYG{p}{)}
\PYG{n}{plt}\PYG{o}{.}\PYG{n}{figure}\PYG{p}{(}\PYG{n}{figsize}\PYG{o}{=}\PYG{p}{(}\PYG{l+m+mi}{10}\PYG{p}{,} \PYG{l+m+mi}{6}\PYG{p}{)}\PYG{p}{)}
\PYG{n}{plt}\PYG{o}{.}\PYG{n}{plot}\PYG{p}{(}\PYG{n}{x\PYGZus{}vec}\PYG{p}{,} \PYG{n}{f}\PYG{p}{(}\PYG{n}{x\PYGZus{}vec}\PYG{p}{)}\PYG{p}{,} \PYG{n}{linestyle}\PYG{o}{=}\PYG{l+s+s1}{\PYGZsq{}}\PYG{l+s+s1}{dashed}\PYG{l+s+s1}{\PYGZsq{}}\PYG{p}{,} \PYG{n}{label}\PYG{o}{=}\PYG{l+s+s2}{\PYGZdq{}}\PYG{l+s+s2}{f(x)}\PYG{l+s+s2}{\PYGZdq{}}\PYG{p}{)}
\PYG{k}{for} \PYG{n}{c} \PYG{o+ow}{in} \PYG{p}{(}\PYG{n}{c\PYGZus{}i}\PYG{o}{:=}\PYG{p}{\PYGZob{}}\PYG{l+s+s2}{\PYGZdq{}}\PYG{l+s+s2}{red}\PYG{l+s+s2}{\PYGZdq{}}\PYG{p}{:} \PYG{o}{\PYGZhy{}}\PYG{l+m+mi}{12}\PYG{p}{,} \PYG{l+s+s2}{\PYGZdq{}}\PYG{l+s+s2}{green}\PYG{l+s+s2}{\PYGZdq{}}\PYG{p}{:} \PYG{o}{\PYGZhy{}}\PYG{l+m+mi}{4}\PYG{p}{,} \PYG{l+s+s2}{\PYGZdq{}}\PYG{l+s+s2}{orange}\PYG{l+s+s2}{\PYGZdq{}}\PYG{p}{:} \PYG{l+m+mi}{1}\PYG{p}{,} \PYG{l+s+s2}{\PYGZdq{}}\PYG{l+s+s2}{magenta}\PYG{l+s+s2}{\PYGZdq{}}\PYG{p}{:} \PYG{l+m+mi}{7}\PYG{p}{,} \PYG{l+s+s2}{\PYGZdq{}}\PYG{l+s+s2}{lightgreen}\PYG{l+s+s2}{\PYGZdq{}}\PYG{p}{:} \PYG{l+m+mi}{14}\PYG{p}{\PYGZcb{}}\PYG{p}{)}\PYG{p}{:}   
    \PYG{n}{minimum} \PYG{o}{=} \PYG{n}{sp}\PYG{o}{.}\PYG{n}{optimize}\PYG{o}{.}\PYG{n}{minimize}\PYG{p}{(}\PYG{n}{f}\PYG{p}{,} \PYG{n}{x0}\PYG{o}{=}\PYG{n}{c\PYGZus{}i}\PYG{p}{[}\PYG{n}{c}\PYG{p}{]}\PYG{p}{)}
    \PYG{n}{line\PYGZus{}x}\PYG{p}{,} \PYG{n}{line\PYGZus{}y} \PYG{o}{=} \PYG{p}{[}\PYG{n}{c\PYGZus{}i}\PYG{p}{[}\PYG{n}{c}\PYG{p}{]}\PYG{p}{,} \PYG{n}{minimum}\PYG{o}{.}\PYG{n}{x}\PYG{p}{[}\PYG{l+m+mi}{0}\PYG{p}{]}\PYG{p}{]}\PYG{p}{,} \PYG{p}{[}\PYG{n}{f}\PYG{p}{(}\PYG{n}{c\PYGZus{}i}\PYG{p}{[}\PYG{n}{c}\PYG{p}{]}\PYG{p}{)}\PYG{p}{,} \PYG{n}{minimum}\PYG{o}{.}\PYG{n}{fun}\PYG{p}{]}
    \PYG{n}{plt}\PYG{o}{.}\PYG{n}{plot}\PYG{p}{(}\PYG{n}{line\PYGZus{}x}\PYG{p}{,} \PYG{n}{line\PYGZus{}y}\PYG{p}{,} \PYG{n}{color}\PYG{o}{=}\PYG{n}{c}\PYG{p}{,} \PYG{n}{label}\PYG{o}{=}\PYG{l+s+sa}{f}\PYG{l+s+s2}{\PYGZdq{}}\PYG{l+s+s2}{Initial condition: x=}\PYG{l+s+si}{\PYGZob{}}\PYG{n}{c\PYGZus{}i}\PYG{p}{[}\PYG{n}{c}\PYG{p}{]}\PYG{l+s+si}{\PYGZcb{}}\PYG{l+s+s2}{\PYGZdq{}}\PYG{p}{)}
\PYG{n}{plt}\PYG{o}{.}\PYG{n}{legend}\PYG{p}{(}\PYG{p}{)}
\PYG{n}{plt}\PYG{o}{.}\PYG{n}{grid}\PYG{p}{(}\PYG{l+s+s2}{\PYGZdq{}}\PYG{l+s+s2}{on}\PYG{l+s+s2}{\PYGZdq{}}\PYG{p}{)}
\PYG{n}{plt}\PYG{o}{.}\PYG{n}{title}\PYG{p}{(}\PYG{l+s+s2}{\PYGZdq{}}\PYG{l+s+s2}{Minimization examples from different initial conditions}\PYG{l+s+s2}{\PYGZdq{}}\PYG{p}{)}
\PYG{n}{plt}\PYG{o}{.}\PYG{n}{show}\PYG{p}{(}\PYG{p}{)}
\end{sphinxVerbatim}

\end{sphinxuseclass}\end{sphinxVerbatimInput}
\begin{sphinxVerbatimOutput}

\begin{sphinxuseclass}{cell_output}
\noindent\sphinxincludegraphics{{5ba328d5be31b0853ab3c4ffdc99d8cf70488f1f0d61fe5169a33b5da18b3671}.png}

\end{sphinxuseclass}\end{sphinxVerbatimOutput}

\end{sphinxuseclass}

\section{Mixture of Gaussians}
\label{\detokenize{01:mixture-of-gaussians}}
\sphinxAtStartPar
Mixture of Gaussians are particular cases of linear combinations of Gaussians r.v. Formally, they are defined as
\begin{equation*}
\begin{split}p = \sum_{i=1}^m w_i p_i\end{split}
\end{equation*}
\sphinxAtStartPar
where \(p_i\) are normal r.v. and \(\sum_i w_1=1\).
As an example, we show here how to define (and plot) the mixture of three Gaussians:

\begin{sphinxuseclass}{cell}\begin{sphinxVerbatimInput}

\begin{sphinxuseclass}{cell_input}
\begin{sphinxVerbatim}[commandchars=\\\{\}]
\PYG{n}{means} \PYG{o}{=} \PYG{p}{[}\PYG{o}{\PYGZhy{}}\PYG{l+m+mi}{2}\PYG{p}{,} \PYG{l+m+mi}{0}\PYG{p}{,} \PYG{l+m+mi}{3}\PYG{p}{]}  
\PYG{n}{stds} \PYG{o}{=} \PYG{p}{[}\PYG{l+m+mf}{0.8}\PYG{p}{,} \PYG{l+m+mf}{1.2}\PYG{p}{,} \PYG{l+m+mf}{0.5}\PYG{p}{]}  
\PYG{n}{weights} \PYG{o}{=} \PYG{p}{[}\PYG{l+m+mi}{1}\PYG{o}{/}\PYG{l+m+mi}{3}\PYG{p}{,} \PYG{l+m+mi}{1}\PYG{o}{/}\PYG{l+m+mi}{3}\PYG{p}{,} \PYG{l+m+mi}{1}\PYG{o}{/}\PYG{l+m+mi}{3}\PYG{p}{]}  

\PYG{n}{x} \PYG{o}{=} \PYG{n}{np}\PYG{o}{.}\PYG{n}{linspace}\PYG{p}{(}\PYG{o}{\PYGZhy{}}\PYG{l+m+mi}{5}\PYG{p}{,} \PYG{l+m+mi}{6}\PYG{p}{,} \PYG{l+m+mi}{1000}\PYG{p}{)}

\PYG{n}{y} \PYG{o}{=} \PYG{n}{np}\PYG{o}{.}\PYG{n}{zeros\PYGZus{}like}\PYG{p}{(}\PYG{n}{x}\PYG{p}{)}
\PYG{k}{for} \PYG{n}{mean}\PYG{p}{,} \PYG{n}{std}\PYG{p}{,} \PYG{n}{weight} \PYG{o+ow}{in} \PYG{n+nb}{zip}\PYG{p}{(}\PYG{n}{means}\PYG{p}{,} \PYG{n}{stds}\PYG{p}{,} \PYG{n}{weights}\PYG{p}{)}\PYG{p}{:}
    \PYG{n}{y} \PYG{o}{+}\PYG{o}{=} \PYG{n}{weight} \PYG{o}{*} \PYG{n}{sp}\PYG{o}{.}\PYG{n}{stats}\PYG{o}{.}\PYG{n}{norm}\PYG{o}{.}\PYG{n}{pdf}\PYG{p}{(}\PYG{n}{x}\PYG{p}{,} \PYG{n}{mean}\PYG{p}{,} \PYG{n}{std}\PYG{p}{)}

\PYG{n}{plt}\PYG{o}{.}\PYG{n}{figure}\PYG{p}{(}\PYG{n}{figsize}\PYG{o}{=}\PYG{p}{(}\PYG{l+m+mi}{10}\PYG{p}{,} \PYG{l+m+mi}{6}\PYG{p}{)}\PYG{p}{)}
\PYG{k}{for} \PYG{n}{mean}\PYG{p}{,} \PYG{n}{std}\PYG{p}{,} \PYG{n}{weight} \PYG{o+ow}{in} \PYG{n+nb}{zip}\PYG{p}{(}\PYG{n}{means}\PYG{p}{,} \PYG{n}{stds}\PYG{p}{,} \PYG{n}{weights}\PYG{p}{)}\PYG{p}{:}
    \PYG{n}{plt}\PYG{o}{.}\PYG{n}{plot}\PYG{p}{(}\PYG{n}{x}\PYG{p}{,} \PYG{n}{weight} \PYG{o}{*} \PYG{n}{sp}\PYG{o}{.}\PYG{n}{stats}\PYG{o}{.}\PYG{n}{norm}\PYG{o}{.}\PYG{n}{pdf}\PYG{p}{(}\PYG{n}{x}\PYG{p}{,} \PYG{n}{mean}\PYG{p}{,} \PYG{n}{std}\PYG{p}{)}\PYG{p}{,} \PYG{l+s+s1}{\PYGZsq{}}\PYG{l+s+s1}{\PYGZhy{}\PYGZhy{}}\PYG{l+s+s1}{\PYGZsq{}}\PYG{p}{,} \PYG{n}{label}\PYG{o}{=}\PYG{l+s+sa}{f}\PYG{l+s+s1}{\PYGZsq{}}\PYG{l+s+s1}{N(}\PYG{l+s+si}{\PYGZob{}}\PYG{n}{mean}\PYG{l+s+si}{\PYGZcb{}}\PYG{l+s+s1}{, }\PYG{l+s+si}{\PYGZob{}}\PYG{n}{std}\PYG{l+s+si}{\PYGZcb{}}\PYG{l+s+s1}{\PYGZca{}2)}\PYG{l+s+s1}{\PYGZsq{}}\PYG{p}{)}

\PYG{n}{plt}\PYG{o}{.}\PYG{n}{plot}\PYG{p}{(}\PYG{n}{x}\PYG{p}{,} \PYG{n}{y}\PYG{p}{,} \PYG{n}{label}\PYG{o}{=}\PYG{l+s+s1}{\PYGZsq{}}\PYG{l+s+s1}{Mixture Density}\PYG{l+s+s1}{\PYGZsq{}}\PYG{p}{,} \PYG{n}{linewidth}\PYG{o}{=}\PYG{l+m+mi}{2}\PYG{p}{)}
\PYG{n}{plt}\PYG{o}{.}\PYG{n}{legend}\PYG{p}{(}\PYG{p}{)}
\PYG{n}{plt}\PYG{o}{.}\PYG{n}{xlabel}\PYG{p}{(}\PYG{l+s+s1}{\PYGZsq{}}\PYG{l+s+s1}{x}\PYG{l+s+s1}{\PYGZsq{}}\PYG{p}{)}
\PYG{n}{plt}\PYG{o}{.}\PYG{n}{ylabel}\PYG{p}{(}\PYG{l+s+s1}{\PYGZsq{}}\PYG{l+s+s1}{Density}\PYG{l+s+s1}{\PYGZsq{}}\PYG{p}{)}
\PYG{n}{plt}\PYG{o}{.}\PYG{n}{grid}\PYG{p}{(}\PYG{l+s+s2}{\PYGZdq{}}\PYG{l+s+s2}{on}\PYG{l+s+s2}{\PYGZdq{}}\PYG{p}{)}
\PYG{n}{plt}\PYG{o}{.}\PYG{n}{title}\PYG{p}{(}\PYG{l+s+s1}{\PYGZsq{}}\PYG{l+s+s1}{Mixture of Three Gaussians}\PYG{l+s+s1}{\PYGZsq{}}\PYG{p}{)}
\PYG{n}{plt}\PYG{o}{.}\PYG{n}{show}\PYG{p}{(}\PYG{p}{)}
\end{sphinxVerbatim}

\end{sphinxuseclass}\end{sphinxVerbatimInput}
\begin{sphinxVerbatimOutput}

\begin{sphinxuseclass}{cell_output}
\noindent\sphinxincludegraphics{{a9c05fc866f7b26d0995eda8c57d9f396c87f8ce0caaad9fbfffc7bd3609aa66}.png}

\end{sphinxuseclass}\end{sphinxVerbatimOutput}

\end{sphinxuseclass}

\section{Conclusion}
\label{\detokenize{01:conclusion}}
\sphinxAtStartPar
Monte Carlo methods represent a powerful paradigm in computational mathematics, offering solutions to problems that would otherwise be analytically intractable. Throughout this chapter, we’ve explored the fundamental principles of probabilistic methods and Monte Carlo simulations, covering essential concepts from probability theory and their practical applications.

\sphinxAtStartPar
We’ve seen how random sampling can be leveraged to approximate integrals, estimate statistical properties, and solve complex problems in high\sphinxhyphen{}dimensional spaces. The exercises provided hands\sphinxhyphen{}on experience with implementing these techniques using \sphinxcode{\sphinxupquote{Python}}’s scientific computing ecosystem, particularly \sphinxcode{\sphinxupquote{NumPy}} and \sphinxcode{\sphinxupquote{SciPy}}.

\sphinxAtStartPar
The versatility of Monte Carlo methods makes them invaluable across numerous fields, from physics and engineering to finance and machine learning. By understanding both the theoretical foundations and practical implementations, you now possess a valuable computational tool that can be applied to a wide range of problems.

\sphinxAtStartPar
As computational resources continue to advance, Monte Carlo techniques will remain at the forefront of scientific computing, enabling researchers and practitioners to tackle increasingly complex challenges. The principles covered in this chapter provide a solid foundation for exploring more advanced Monte Carlo techniques such as Markov Chain Monte Carlo (MCMC), importance sampling, particle filters… But, in our case, Monte Carlo processes can be seen as good abstraction of real\sphinxhyphen{}life processes, and Machine Learning techniques then serve as tools to study the outcomes of Monte Carlo experiments.

\sphinxAtStartPar
Remember that Monte Carlo methods are not just computational techniques but represent a probabilistic way of thinking about problems. This mindset of using randomness to solve deterministic problems opens up new approaches to complex challenges across all quantitative disciplines.

\sphinxstepscope


\chapter{Linear Models}
\label{\detokenize{02:linear-models}}\label{\detokenize{02::doc}}

\section{Introduction}
\label{\detokenize{02:introduction}}
\sphinxAtStartPar
Linear models form one of the foundational pillars of statistical learning. Their simplicity, interpretability, and strong theoretical underpinnings make them an essential first step when analyzing data and building predictive models. Despite their apparent simplicity, linear models offer a surprising level of flexibility, especially when combined with feature transformations such as polynomials or radial basis functions.

\sphinxAtStartPar
In this chapter, we explore the principles and practical implementation of linear models. We begin with the generation of synthetic multivariate data to simulate a controlled environment for experimentation. We then dive into the manual computation of the Ordinary Least Squares (OLS) solution using matrix algebra, before comparing our results to those produced by well\sphinxhyphen{}established Python libraries such as \sphinxcode{\sphinxupquote{statsmodels}}.

\sphinxAtStartPar
The chapter also highlights the importance of residual analysis for model diagnostics, offering tools to detect violations of assumptions like linearity, homoscedasticity, and normality. We further extend the linear model framework through input transformations, allowing us to capture non\sphinxhyphen{}linear patterns while preserving the model’s linear structure in parameters.

\sphinxAtStartPar
Finally, we introduce \sphinxstylestrong{Ridge Regression}, a regularized variant of linear regression, and examine its role in mitigating overfitting—particularly in scenarios involving multicollinearity or high\sphinxhyphen{}dimensional input spaces.

\sphinxAtStartPar
We now assume basic familiarity with matrix algebra and Python packages such as \sphinxcode{\sphinxupquote{NumPy}}, \sphinxcode{\sphinxupquote{Pandas}}, and \sphinxcode{\sphinxupquote{Statsmodels}}.

\begin{sphinxShadowBox}
\sphinxstylesidebartitle{}

\sphinxAtStartPar
Don’t forget: if some of those libraries are not installed on your computer, you can always install them using \sphinxhref{http://pypi.org}{Pip}.
\end{sphinxShadowBox}

\sphinxAtStartPar
We define the necessary imports as well as the \sphinxcode{\sphinxupquote{pprint}} function.

\begin{sphinxuseclass}{cell}\begin{sphinxVerbatimInput}

\begin{sphinxuseclass}{cell_input}
\begin{sphinxVerbatim}[commandchars=\\\{\}]
\PYG{k+kn}{import}\PYG{+w}{ }\PYG{n+nn}{numpy}\PYG{+w}{ }\PYG{k}{as}\PYG{+w}{ }\PYG{n+nn}{np}
\PYG{k+kn}{import}\PYG{+w}{ }\PYG{n+nn}{pandas}\PYG{+w}{ }\PYG{k}{as}\PYG{+w}{ }\PYG{n+nn}{pd}
\PYG{k+kn}{import}\PYG{+w}{ }\PYG{n+nn}{seaborn}\PYG{+w}{ }\PYG{k}{as}\PYG{+w}{ }\PYG{n+nn}{sns}
\PYG{k+kn}{import}\PYG{+w}{ }\PYG{n+nn}{statsmodels}\PYG{n+nn}{.}\PYG{n+nn}{api}\PYG{+w}{ }\PYG{k}{as}\PYG{+w}{ }\PYG{n+nn}{sm}
\PYG{k+kn}{from}\PYG{+w}{ }\PYG{n+nn}{scipy}\PYG{n+nn}{.}\PYG{n+nn}{stats}\PYG{+w}{ }\PYG{k+kn}{import} \PYG{n}{probplot}
\PYG{k+kn}{from}\PYG{+w}{ }\PYG{n+nn}{IPython}\PYG{n+nn}{.}\PYG{n+nn}{display}\PYG{+w}{ }\PYG{k+kn}{import} \PYG{n}{display}\PYG{p}{,} \PYG{n}{Math}
\PYG{k+kn}{from}\PYG{+w}{ }\PYG{n+nn}{numpyarray\PYGZus{}to\PYGZus{}latex}\PYG{n+nn}{.}\PYG{n+nn}{jupyter}\PYG{+w}{ }\PYG{k+kn}{import} \PYG{n}{to\PYGZus{}ltx}
\PYG{k+kn}{import}\PYG{+w}{ }\PYG{n+nn}{matplotlib}\PYG{n+nn}{.}\PYG{n+nn}{pyplot}\PYG{+w}{ }\PYG{k}{as}\PYG{+w}{ }\PYG{n+nn}{plt}
\PYG{n}{plt}\PYG{o}{.}\PYG{n}{style}\PYG{o}{.}\PYG{n}{use}\PYG{p}{(}\PYG{p}{[}\PYG{l+s+s1}{\PYGZsq{}}\PYG{l+s+s1}{seaborn\PYGZhy{}v0\PYGZus{}8\PYGZhy{}muted}\PYG{l+s+s1}{\PYGZsq{}}\PYG{p}{,} \PYG{l+s+s1}{\PYGZsq{}}\PYG{l+s+s1}{practicals.mplstyle}\PYG{l+s+s1}{\PYGZsq{}}\PYG{p}{]}\PYG{p}{)}
\PYG{k}{def}\PYG{+w}{ }\PYG{n+nf}{pprint}\PYG{p}{(}\PYG{o}{*}\PYG{n}{args}\PYG{p}{)}\PYG{p}{:}
    \PYG{n}{res} \PYG{o}{=} \PYG{l+s+s2}{\PYGZdq{}}\PYG{l+s+s2}{\PYGZdq{}}
    \PYG{k}{for} \PYG{n}{i} \PYG{o+ow}{in} \PYG{n}{args}\PYG{p}{:}
        \PYG{k}{if} \PYG{n+nb}{type}\PYG{p}{(}\PYG{n}{i}\PYG{p}{)} \PYG{o}{==} \PYG{n}{np}\PYG{o}{.}\PYG{n}{ndarray}\PYG{p}{:}
            \PYG{n}{res} \PYG{o}{+}\PYG{o}{=} \PYG{n}{to\PYGZus{}ltx}\PYG{p}{(}\PYG{n}{i}\PYG{p}{,} \PYG{n}{brackets}\PYG{o}{=}\PYG{l+s+s1}{\PYGZsq{}}\PYG{l+s+s1}{[]}\PYG{l+s+s1}{\PYGZsq{}}\PYG{p}{)}
        \PYG{k}{elif} \PYG{n+nb}{type}\PYG{p}{(}\PYG{n}{i}\PYG{p}{)} \PYG{o}{==} \PYG{n+nb}{str}\PYG{p}{:}
            \PYG{n}{res} \PYG{o}{+}\PYG{o}{=} \PYG{n}{i}
    \PYG{n}{display}\PYG{p}{(}\PYG{n}{Math}\PYG{p}{(}\PYG{n}{res}\PYG{p}{)}\PYG{p}{)}
\end{sphinxVerbatim}

\end{sphinxuseclass}\end{sphinxVerbatimInput}

\end{sphinxuseclass}

\section{Generating Synthetic Multivariate Data}
\label{\detokenize{02:generating-synthetic-multivariate-data}}
\sphinxAtStartPar
We’ll create a dataset where \sphinxcode{\sphinxupquote{Y}} depends on three features: \sphinxcode{\sphinxupquote{X1, X2, X3}}.
The true data\sphinxhyphen{}generating process (DGP) is:
\begin{equation*}
\begin{split}
  Y = \beta_0 + \beta_1 X_1 + \beta_2 X_2 + \beta_3 X_3 + \epsilon,
\end{split}
\end{equation*}
\sphinxAtStartPar
where \(\epsilon \sim \mathcal{N}(0, \sigma^2)\).


\subsection{Exercise}
\label{\detokenize{02:exercise}}\begin{enumerate}
\sphinxsetlistlabels{\arabic}{enumi}{enumii}{}{.}%
\item {} 
\sphinxAtStartPar
Create \sphinxcode{\sphinxupquote{N}} observations of \sphinxcode{\sphinxupquote{(X1, X2, X3)}}.

\item {} 
\sphinxAtStartPar
Define true coefficients \sphinxcode{\sphinxupquote{beta\_true}}.

\item {} 
\sphinxAtStartPar
Generate noise and produce \sphinxcode{\sphinxupquote{Y}}.

\item {} 
\sphinxAtStartPar
Store everything in a DataFrame for convenience.

\end{enumerate}


\subsubsection{Solution}
\label{\detokenize{02:solution}}
\sphinxAtStartPar
The following code generates a synthetic dataset for a linear regression problem with three explanatory variables (\(\mathbf X_1\), \(\mathbf X_2\), \(\mathbf X_3\)) and one target variable (\(\mathbf Y\)). It sets a random seed for reproducibility and defines true regression coefficients (\sphinxcode{\sphinxupquote{beta\_true}}) and a standard deviation for the noise (\sphinxcode{\sphinxupquote{sigma\_true}}). The explanatory variables are drawn from different distributions: uniform for \(\mathbf X_1\) and \(\mathbf X_2\), and normal for \(\mathbf X_3\). The true response \sphinxcode{\sphinxupquote{Y\_true}} is computed as a linear combination of the variables with the defined coefficients, and Gaussian noise is added to simulate observed values of \(\mathbf Y\). The resulting dataset, stored in a \sphinxcode{\sphinxupquote{pandas}} DataFrame, includes the predictors and the noisy target variable, and the first few rows are displayed using \sphinxcode{\sphinxupquote{df.head()}}.

\begin{sphinxuseclass}{cell}\begin{sphinxVerbatimInput}

\begin{sphinxuseclass}{cell_input}
\begin{sphinxVerbatim}[commandchars=\\\{\}]
\PYG{n}{np}\PYG{o}{.}\PYG{n}{random}\PYG{o}{.}\PYG{n}{seed}\PYG{p}{(}\PYG{l+m+mi}{123}\PYG{p}{)}

\PYG{n}{N} \PYG{o}{=} \PYG{l+m+mi}{200}
\PYG{n}{sigma\PYGZus{}true} \PYG{o}{=} \PYG{l+m+mf}{2.0}

\PYG{n}{beta\PYGZus{}true} \PYG{o}{=} \PYG{n}{np}\PYG{o}{.}\PYG{n}{array}\PYG{p}{(}\PYG{p}{[}\PYG{l+m+mf}{5.0}\PYG{p}{,} \PYG{l+m+mf}{1.5}\PYG{p}{,} \PYG{o}{\PYGZhy{}}\PYG{l+m+mf}{2.0}\PYG{p}{,} \PYG{l+m+mf}{0.5}\PYG{p}{]}\PYG{p}{)}

\PYG{n}{X1} \PYG{o}{=} \PYG{n}{np}\PYG{o}{.}\PYG{n}{random}\PYG{o}{.}\PYG{n}{uniform}\PYG{p}{(}\PYG{o}{\PYGZhy{}}\PYG{l+m+mi}{5}\PYG{p}{,} \PYG{l+m+mi}{5}\PYG{p}{,} \PYG{n}{size}\PYG{o}{=}\PYG{n}{N}\PYG{p}{)}
\PYG{n}{X2} \PYG{o}{=} \PYG{n}{np}\PYG{o}{.}\PYG{n}{random}\PYG{o}{.}\PYG{n}{uniform}\PYG{p}{(}\PYG{o}{\PYGZhy{}}\PYG{l+m+mi}{10}\PYG{p}{,} \PYG{l+m+mi}{10}\PYG{p}{,} \PYG{n}{size}\PYG{o}{=}\PYG{n}{N}\PYG{p}{)}
\PYG{n}{X3} \PYG{o}{=} \PYG{n}{np}\PYG{o}{.}\PYG{n}{random}\PYG{o}{.}\PYG{n}{normal}\PYG{p}{(}\PYG{l+m+mi}{0}\PYG{p}{,} \PYG{l+m+mi}{2}\PYG{p}{,} \PYG{n}{size}\PYG{o}{=}\PYG{n}{N}\PYG{p}{)}

\PYG{n}{df} \PYG{o}{=} \PYG{n}{pd}\PYG{o}{.}\PYG{n}{DataFrame}\PYG{p}{(}\PYG{p}{\PYGZob{}}\PYG{l+s+s1}{\PYGZsq{}}\PYG{l+s+s1}{X1}\PYG{l+s+s1}{\PYGZsq{}}\PYG{p}{:}\PYG{n}{X1}\PYG{p}{,} \PYG{l+s+s1}{\PYGZsq{}}\PYG{l+s+s1}{X2}\PYG{l+s+s1}{\PYGZsq{}}\PYG{p}{:}\PYG{n}{X2}\PYG{p}{,} \PYG{l+s+s1}{\PYGZsq{}}\PYG{l+s+s1}{X3}\PYG{l+s+s1}{\PYGZsq{}}\PYG{p}{:}\PYG{n}{X3}\PYG{p}{\PYGZcb{}}\PYG{p}{)}

\PYG{n}{Y\PYGZus{}true} \PYG{o}{=} \PYG{p}{(}\PYG{n}{beta\PYGZus{}true}\PYG{p}{[}\PYG{l+m+mi}{0}\PYG{p}{]}
          \PYG{o}{+} \PYG{n}{beta\PYGZus{}true}\PYG{p}{[}\PYG{l+m+mi}{1}\PYG{p}{]}\PYG{o}{*}\PYG{n}{df}\PYG{p}{[}\PYG{l+s+s1}{\PYGZsq{}}\PYG{l+s+s1}{X1}\PYG{l+s+s1}{\PYGZsq{}}\PYG{p}{]}
          \PYG{o}{+} \PYG{n}{beta\PYGZus{}true}\PYG{p}{[}\PYG{l+m+mi}{2}\PYG{p}{]}\PYG{o}{*}\PYG{n}{df}\PYG{p}{[}\PYG{l+s+s1}{\PYGZsq{}}\PYG{l+s+s1}{X2}\PYG{l+s+s1}{\PYGZsq{}}\PYG{p}{]}
          \PYG{o}{+} \PYG{n}{beta\PYGZus{}true}\PYG{p}{[}\PYG{l+m+mi}{3}\PYG{p}{]}\PYG{o}{*}\PYG{n}{df}\PYG{p}{[}\PYG{l+s+s1}{\PYGZsq{}}\PYG{l+s+s1}{X3}\PYG{l+s+s1}{\PYGZsq{}}\PYG{p}{]}\PYG{p}{)}

\PYG{n}{noise} \PYG{o}{=} \PYG{n}{np}\PYG{o}{.}\PYG{n}{random}\PYG{o}{.}\PYG{n}{normal}\PYG{p}{(}\PYG{l+m+mi}{0}\PYG{p}{,} \PYG{n}{sigma\PYGZus{}true}\PYG{p}{,} \PYG{n}{size}\PYG{o}{=}\PYG{n}{N}\PYG{p}{)}
\PYG{n}{df}\PYG{p}{[}\PYG{l+s+s1}{\PYGZsq{}}\PYG{l+s+s1}{Y}\PYG{l+s+s1}{\PYGZsq{}}\PYG{p}{]} \PYG{o}{=} \PYG{n}{Y\PYGZus{}true} \PYG{o}{+} \PYG{n}{noise}

\PYG{n}{df}\PYG{o}{.}\PYG{n}{head}\PYG{p}{(}\PYG{p}{)}
\end{sphinxVerbatim}

\end{sphinxuseclass}\end{sphinxVerbatimInput}
\begin{sphinxVerbatimOutput}

\begin{sphinxuseclass}{cell_output}
\begin{sphinxVerbatim}[commandchars=\\\{\}]
         X1        X2        X3          Y
0  1.964692  0.852719 \PYGZhy{}3.510804   5.687964
1 \PYGZhy{}2.138607 \PYGZhy{}8.664511 \PYGZhy{}0.697215  16.509223
2 \PYGZhy{}2.731485  3.067297 \PYGZhy{}0.385230  \PYGZhy{}4.050829
3  0.513148  9.921727  0.898271 \PYGZhy{}12.453060
4  2.194690  5.387947 \PYGZhy{}0.290727  \PYGZhy{}1.697682
\end{sphinxVerbatim}

\end{sphinxuseclass}\end{sphinxVerbatimOutput}

\end{sphinxuseclass}

\subsection{Quick Data Inspection}
\label{\detokenize{02:quick-data-inspection}}
\sphinxAtStartPar
The following code performs exploratory data analysis on a DataFrame \sphinxcode{\sphinxupquote{df}}. It first prints summary statistics of the DataFrame using \sphinxcode{\sphinxupquote{df.describe()}}, which includes metrics like mean, standard deviation, min, and max for each numerical column. Then, it creates a pairplot using Seaborn with regression lines (\sphinxcode{\sphinxupquote{kind='reg'}}) to visualize pairwise relationships between all numerical features in the DataFrame. The regression lines are colored red, as specified in the \sphinxcode{\sphinxupquote{plot\_kws}}. A title is added slightly above the plot using \sphinxcode{\sphinxupquote{plt.suptitle()}}, and finally, the complete visualization is displayed with \sphinxcode{\sphinxupquote{plt.show()}}.

\begin{sphinxShadowBox}
\sphinxstylesidebartitle{}

\sphinxAtStartPar
Doing a first data inspection is extremely important. In real\sphinxhyphen{}life cases, you could discover patterns in the data, discard particular features, or even, sometimes, solve the problem you are tackling just by extracting visual information by the data. This step also helps understanding the problem you are dealing with. The more visualization of the data you make, the more expert you become on the dataset. Never underestimate the importance of the base knowledge of the problem: \sphinxstyleemphasis{data is more than just numbers}.
\end{sphinxShadowBox}

\begin{sphinxuseclass}{cell}\begin{sphinxVerbatimInput}

\begin{sphinxuseclass}{cell_input}
\begin{sphinxVerbatim}[commandchars=\\\{\}]
\PYG{n+nb}{print}\PYG{p}{(}\PYG{n}{df}\PYG{o}{.}\PYG{n}{describe}\PYG{p}{(}\PYG{p}{)}\PYG{p}{)}
\PYG{n}{sns}\PYG{o}{.}\PYG{n}{pairplot}\PYG{p}{(}\PYG{n}{df}\PYG{p}{,} \PYG{n}{kind}\PYG{o}{=}\PYG{l+s+s1}{\PYGZsq{}}\PYG{l+s+s1}{reg}\PYG{l+s+s1}{\PYGZsq{}}\PYG{p}{,} \PYG{n}{plot\PYGZus{}kws}\PYG{o}{=}\PYG{p}{\PYGZob{}}\PYG{l+s+s1}{\PYGZsq{}}\PYG{l+s+s1}{line\PYGZus{}kws}\PYG{l+s+s1}{\PYGZsq{}}\PYG{p}{:}\PYG{p}{\PYGZob{}}\PYG{l+s+s1}{\PYGZsq{}}\PYG{l+s+s1}{color}\PYG{l+s+s1}{\PYGZsq{}}\PYG{p}{:}\PYG{l+s+s1}{\PYGZsq{}}\PYG{l+s+s1}{green}\PYG{l+s+s1}{\PYGZsq{}}\PYG{p}{\PYGZcb{}}\PYG{p}{\PYGZcb{}}\PYG{p}{)}
\PYG{n}{plt}\PYG{o}{.}\PYG{n}{suptitle}\PYG{p}{(}\PYG{l+s+s2}{\PYGZdq{}}\PYG{l+s+s2}{Pairplot of the Generated Data}\PYG{l+s+s2}{\PYGZdq{}}\PYG{p}{,} \PYG{n}{y}\PYG{o}{=}\PYG{l+m+mf}{1.02}\PYG{p}{)}
\PYG{n}{plt}\PYG{o}{.}\PYG{n}{show}\PYG{p}{(}\PYG{p}{)}
\end{sphinxVerbatim}

\end{sphinxuseclass}\end{sphinxVerbatimInput}
\begin{sphinxVerbatimOutput}

\begin{sphinxuseclass}{cell_output}
\begin{sphinxVerbatim}[commandchars=\\\{\}]
               X1          X2          X3           Y
count  200.000000  200.000000  200.000000  200.000000
mean     0.032708   \PYGZhy{}0.193044   \PYGZhy{}0.112239    5.572700
std      2.663947    5.907073    1.963533   12.473036
min     \PYGZhy{}4.973119   \PYGZhy{}9.929356   \PYGZhy{}5.262876  \PYGZhy{}20.342999
25\PYGZpc{}     \PYGZhy{}1.953937   \PYGZhy{}5.254828   \PYGZhy{}1.286318   \PYGZhy{}4.258475
50\PYGZpc{}      0.207716   \PYGZhy{}0.714124   \PYGZhy{}0.101728    5.417761
75\PYGZpc{}      2.181854    5.168943    1.206797   14.861563
max      4.953585    9.921727    4.501352   30.849943
\end{sphinxVerbatim}

\noindent\sphinxincludegraphics{{456981c4f9bfe9b9d8b1c0e7cdae067f5cdd8df79b0ad78dd5a71c595220be19}.png}

\end{sphinxuseclass}\end{sphinxVerbatimOutput}

\end{sphinxuseclass}

\section{Ordinary Least Squares (OLS) – Manual Computation Using Matrix Algebra}
\label{\detokenize{02:ordinary-least-squares-ols-manual-computation-using-matrix-algebra}}
\sphinxAtStartPar
We begin by specifying the standard linear regression model in matrix form:
\begin{equation*}
\begin{split}
Y = X \boldsymbol{\beta} + \boldsymbol{\epsilon},
\end{split}
\end{equation*}
\sphinxAtStartPar
where:
\begin{itemize}
\item {} 
\sphinxAtStartPar
\(\mathbf{Y}\) is an \(N \times 1\) vector of observed responses,

\item {} 
\sphinxAtStartPar
\(X\) is an \(N \times p\) \sphinxstyleemphasis{design matrix}, which includes a column of ones to account for the intercept,

\item {} 
\sphinxAtStartPar
\(\boldsymbol{\beta}\) is a \(p \times 1\) vector of regression coefficients (including the intercept),

\item {} 
\sphinxAtStartPar
\(\boldsymbol{\epsilon}\) is an \(N \times 1\) vector of errors or residuals.

\end{itemize}

\sphinxAtStartPar
The goal of Ordinary Least Squares is to find the vector \(\hat{\boldsymbol{\beta}}\) that minimizes the sum of squared residuals. There is a well\sphinxhyphen{}known closed\sphinxhyphen{}form solution for this optimization problem, given by:
\begin{equation*}
\begin{split}
\hat{\boldsymbol{\beta}} = (X^\top X)^{-1} X^\top Y.
\end{split}
\end{equation*}
\sphinxAtStartPar
This formula provides the best linear unbiased estimator (BLUE) of the coefficients under standard assumptions (linearity, independence, homoscedasticity, and normally distributed errors).

\sphinxAtStartPar
In the following steps, we will:
\begin{enumerate}
\sphinxsetlistlabels{\arabic}{enumi}{enumii}{}{.}%
\item {} 
\sphinxAtStartPar
Manually construct the design matrix \(X\), ensuring it includes a column of ones for the intercept.

\item {} 
\sphinxAtStartPar
Compute the estimated coefficients \(\hat{\boldsymbol{\beta}}\) using NumPy’s matrix operations, specifically \sphinxcode{\sphinxupquote{np.linalg.inv(...)}}.

\item {} 
\sphinxAtStartPar
Compare the computed estimates with the known or true values of the coefficients (\sphinxcode{\sphinxupquote{beta\_true}}) used to generate the data.

\item {} 
\sphinxAtStartPar
Validate our manual computation by comparing it to the output from the \sphinxcode{\sphinxupquote{statsmodels}} OLS implementation.

\end{enumerate}

\begin{sphinxuseclass}{cell}\begin{sphinxVerbatimInput}

\begin{sphinxuseclass}{cell_input}
\begin{sphinxVerbatim}[commandchars=\\\{\}]
\PYG{n}{X\PYGZus{}mat} \PYG{o}{=} \PYG{n}{np}\PYG{o}{.}\PYG{n}{column\PYGZus{}stack}\PYG{p}{(}\PYG{p}{[}
    \PYG{n}{np}\PYG{o}{.}\PYG{n}{ones}\PYG{p}{(}\PYG{n}{N}\PYG{p}{)}\PYG{p}{,}
    \PYG{n}{df}\PYG{p}{[}\PYG{l+s+s1}{\PYGZsq{}}\PYG{l+s+s1}{X1}\PYG{l+s+s1}{\PYGZsq{}}\PYG{p}{]}\PYG{o}{.}\PYG{n}{values}\PYG{p}{,}
    \PYG{n}{df}\PYG{p}{[}\PYG{l+s+s1}{\PYGZsq{}}\PYG{l+s+s1}{X2}\PYG{l+s+s1}{\PYGZsq{}}\PYG{p}{]}\PYG{o}{.}\PYG{n}{values}\PYG{p}{,}
    \PYG{n}{df}\PYG{p}{[}\PYG{l+s+s1}{\PYGZsq{}}\PYG{l+s+s1}{X3}\PYG{l+s+s1}{\PYGZsq{}}\PYG{p}{]}\PYG{o}{.}\PYG{n}{values}
\PYG{p}{]}\PYG{p}{)} 

\PYG{n}{y\PYGZus{}vec} \PYG{o}{=} \PYG{n}{df}\PYG{p}{[}\PYG{l+s+s1}{\PYGZsq{}}\PYG{l+s+s1}{Y}\PYG{l+s+s1}{\PYGZsq{}}\PYG{p}{]}\PYG{o}{.}\PYG{n}{values}\PYG{o}{.}\PYG{n}{reshape}\PYG{p}{(}\PYG{o}{\PYGZhy{}}\PYG{l+m+mi}{1}\PYG{p}{,}\PYG{l+m+mi}{1}\PYG{p}{)}

\PYG{n+nb}{print}\PYG{p}{(}\PYG{l+s+s2}{\PYGZdq{}}\PYG{l+s+s2}{X\PYGZus{}mat shape:}\PYG{l+s+s2}{\PYGZdq{}}\PYG{p}{,} \PYG{n}{X\PYGZus{}mat}\PYG{o}{.}\PYG{n}{shape}\PYG{p}{)}
\PYG{n+nb}{print}\PYG{p}{(}\PYG{l+s+s2}{\PYGZdq{}}\PYG{l+s+s2}{y\PYGZus{}vec shape:}\PYG{l+s+s2}{\PYGZdq{}}\PYG{p}{,} \PYG{n}{y\PYGZus{}vec}\PYG{o}{.}\PYG{n}{shape}\PYG{p}{)}
\end{sphinxVerbatim}

\end{sphinxuseclass}\end{sphinxVerbatimInput}
\begin{sphinxVerbatimOutput}

\begin{sphinxuseclass}{cell_output}
\begin{sphinxVerbatim}[commandchars=\\\{\}]
X\PYGZus{}mat shape: (200, 4)
y\PYGZus{}vec shape: (200, 1)
\end{sphinxVerbatim}

\end{sphinxuseclass}\end{sphinxVerbatimOutput}

\end{sphinxuseclass}
\begin{sphinxuseclass}{cell}\begin{sphinxVerbatimInput}

\begin{sphinxuseclass}{cell_input}
\begin{sphinxVerbatim}[commandchars=\\\{\}]
\PYG{n}{XTX} \PYG{o}{=} \PYG{n}{X\PYGZus{}mat}\PYG{o}{.}\PYG{n}{T} \PYG{o}{@} \PYG{n}{X\PYGZus{}mat}
\PYG{n}{XTX\PYGZus{}inv} \PYG{o}{=} \PYG{n}{np}\PYG{o}{.}\PYG{n}{linalg}\PYG{o}{.}\PYG{n}{inv}\PYG{p}{(}\PYG{n}{XTX}\PYG{p}{)}
\PYG{n}{XTy} \PYG{o}{=} \PYG{n}{X\PYGZus{}mat}\PYG{o}{.}\PYG{n}{T} \PYG{o}{@} \PYG{n}{y\PYGZus{}vec}
\PYG{n}{beta\PYGZus{}hat\PYGZus{}manual} \PYG{o}{=} \PYG{n}{XTX\PYGZus{}inv} \PYG{o}{@} \PYG{n}{XTy}

\PYG{n}{beta\PYGZus{}hat\PYGZus{}manual} \PYG{o}{=} \PYG{n}{beta\PYGZus{}hat\PYGZus{}manual}\PYG{o}{.}\PYG{n}{flatten}\PYG{p}{(}\PYG{p}{)}

\PYG{n}{pprint}\PYG{p}{(}\PYG{l+s+s2}{\PYGZdq{}}\PYG{l+s+se}{\PYGZbs{}\PYGZbs{}}\PYG{l+s+s2}{text}\PYG{l+s+s2}{\PYGZob{}}\PYG{l+s+s2}{Manual OLS estimates\PYGZcb{}:}\PYG{l+s+s2}{\PYGZdq{}}\PYG{p}{,} \PYG{n}{beta\PYGZus{}hat\PYGZus{}manual}\PYG{p}{)}
\PYG{n}{pprint}\PYG{p}{(}\PYG{l+s+s2}{\PYGZdq{}}\PYG{l+s+se}{\PYGZbs{}\PYGZbs{}}\PYG{l+s+s2}{text}\PYG{l+s+s2}{\PYGZob{}}\PYG{l+s+s2}{True coefficients\PYGZcb{}:}\PYG{l+s+s2}{\PYGZdq{}}\PYG{p}{,} \PYG{n}{beta\PYGZus{}true}\PYG{p}{)}
\end{sphinxVerbatim}

\end{sphinxuseclass}\end{sphinxVerbatimInput}
\begin{sphinxVerbatimOutput}

\begin{sphinxuseclass}{cell_output}\begin{equation*}
\begin{split}\displaystyle \text{Manual OLS estimates}:\left[
\begin{array}{}
  5.2043 &  1.4772 & -1.9629 &  0.5243
\end{array}
\right]\end{split}
\end{equation*}\begin{equation*}
\begin{split}\displaystyle \text{True coefficients}:\left[
\begin{array}{}
  5.0000 &  1.5000 & -2.0000 &  0.5000
\end{array}
\right]\end{split}
\end{equation*}
\end{sphinxuseclass}\end{sphinxVerbatimOutput}

\end{sphinxuseclass}
\sphinxAtStartPar
We see the manually computed OLS estimates are close to the true ones.


\subsection{Comparison with statsmodels}
\label{\detokenize{02:comparison-with-statsmodels}}
\sphinxAtStartPar
We’ll let \sphinxcode{\sphinxupquote{statsmodels}} compute OLS, then compare the results in detail.

\begin{sphinxShadowBox}
\sphinxstylesidebartitle{}

\sphinxAtStartPar
Note, concretely, you are expected to be able to compute all the algorithms \sphinxstyleemphasis{by hand}, meaning, \sphinxstyleemphasis{in plain \sphinxcode{\sphinxupquote{Python}} code}. However, in real\sphinxhyphen{}life, you can of course use the existing libraries, and knowing the most common ones is part of the basic knowledge expected from a modern data scientist.
\end{sphinxShadowBox}

\begin{sphinxuseclass}{cell}\begin{sphinxVerbatimInput}

\begin{sphinxuseclass}{cell_input}
\begin{sphinxVerbatim}[commandchars=\\\{\}]
\PYG{n}{X\PYGZus{}design\PYGZus{}sm} \PYG{o}{=} \PYG{n}{sm}\PYG{o}{.}\PYG{n}{add\PYGZus{}constant}\PYG{p}{(}\PYG{n}{df}\PYG{p}{[}\PYG{p}{[}\PYG{l+s+s1}{\PYGZsq{}}\PYG{l+s+s1}{X1}\PYG{l+s+s1}{\PYGZsq{}}\PYG{p}{,}\PYG{l+s+s1}{\PYGZsq{}}\PYG{l+s+s1}{X2}\PYG{l+s+s1}{\PYGZsq{}}\PYG{p}{,}\PYG{l+s+s1}{\PYGZsq{}}\PYG{l+s+s1}{X3}\PYG{l+s+s1}{\PYGZsq{}}\PYG{p}{]}\PYG{p}{]}\PYG{p}{)}
\PYG{n}{model\PYGZus{}ols} \PYG{o}{=} \PYG{n}{sm}\PYG{o}{.}\PYG{n}{OLS}\PYG{p}{(}\PYG{n}{df}\PYG{p}{[}\PYG{l+s+s1}{\PYGZsq{}}\PYG{l+s+s1}{Y}\PYG{l+s+s1}{\PYGZsq{}}\PYG{p}{]}\PYG{p}{,} \PYG{n}{X\PYGZus{}design\PYGZus{}sm}\PYG{p}{)}\PYG{o}{.}\PYG{n}{fit}\PYG{p}{(}\PYG{p}{)}
\PYG{n+nb}{print}\PYG{p}{(}\PYG{n}{model\PYGZus{}ols}\PYG{o}{.}\PYG{n}{summary}\PYG{p}{(}\PYG{p}{)}\PYG{p}{)}

\PYG{n}{beta\PYGZus{}hat\PYGZus{}sm} \PYG{o}{=} \PYG{n}{model\PYGZus{}ols}\PYG{o}{.}\PYG{n}{params}\PYG{o}{.}\PYG{n}{values}
\PYG{n+nb}{print}\PYG{p}{(}\PYG{l+s+s2}{\PYGZdq{}}\PYG{l+s+s2}{statsmodels OLS:}\PYG{l+s+s2}{\PYGZdq{}}\PYG{p}{,} \PYG{n}{beta\PYGZus{}hat\PYGZus{}sm}\PYG{p}{)}
\PYG{n+nb}{print}\PYG{p}{(}\PYG{l+s+s2}{\PYGZdq{}}\PYG{l+s+s2}{Manual OLS:}\PYG{l+s+s2}{\PYGZdq{}}\PYG{p}{,} \PYG{n}{beta\PYGZus{}hat\PYGZus{}manual}\PYG{p}{)}
\end{sphinxVerbatim}

\end{sphinxuseclass}\end{sphinxVerbatimInput}
\begin{sphinxVerbatimOutput}

\begin{sphinxuseclass}{cell_output}
\begin{sphinxVerbatim}[commandchars=\\\{\}]
                            OLS Regression Results                            
==============================================================================
Dep. Variable:                      Y   R\PYGZhy{}squared:                       0.976
Model:                            OLS   Adj. R\PYGZhy{}squared:                  0.976
Method:                 Least Squares   F\PYGZhy{}statistic:                     2657.
Date:                Tue, 16 Sep 2025   Prob (F\PYGZhy{}statistic):          2.01e\PYGZhy{}158
Time:                        13:31:42   Log\PYGZhy{}Likelihood:                \PYGZhy{}415.03
No. Observations:                 200   AIC:                             838.1
Df Residuals:                     196   BIC:                             851.2
Df Model:                           3                                         
Covariance Type:            nonrobust                                         
==============================================================================
                 coef    std err          t      P\PYGZgt{}|t|      [0.025      0.975]
\PYGZhy{}\PYGZhy{}\PYGZhy{}\PYGZhy{}\PYGZhy{}\PYGZhy{}\PYGZhy{}\PYGZhy{}\PYGZhy{}\PYGZhy{}\PYGZhy{}\PYGZhy{}\PYGZhy{}\PYGZhy{}\PYGZhy{}\PYGZhy{}\PYGZhy{}\PYGZhy{}\PYGZhy{}\PYGZhy{}\PYGZhy{}\PYGZhy{}\PYGZhy{}\PYGZhy{}\PYGZhy{}\PYGZhy{}\PYGZhy{}\PYGZhy{}\PYGZhy{}\PYGZhy{}\PYGZhy{}\PYGZhy{}\PYGZhy{}\PYGZhy{}\PYGZhy{}\PYGZhy{}\PYGZhy{}\PYGZhy{}\PYGZhy{}\PYGZhy{}\PYGZhy{}\PYGZhy{}\PYGZhy{}\PYGZhy{}\PYGZhy{}\PYGZhy{}\PYGZhy{}\PYGZhy{}\PYGZhy{}\PYGZhy{}\PYGZhy{}\PYGZhy{}\PYGZhy{}\PYGZhy{}\PYGZhy{}\PYGZhy{}\PYGZhy{}\PYGZhy{}\PYGZhy{}\PYGZhy{}\PYGZhy{}\PYGZhy{}\PYGZhy{}\PYGZhy{}\PYGZhy{}\PYGZhy{}\PYGZhy{}\PYGZhy{}\PYGZhy{}\PYGZhy{}\PYGZhy{}\PYGZhy{}\PYGZhy{}\PYGZhy{}\PYGZhy{}\PYGZhy{}\PYGZhy{}\PYGZhy{}
const          5.2043      0.138     37.720      0.000       4.932       5.476
X1             1.4772      0.052     28.497      0.000       1.375       1.579
X2            \PYGZhy{}1.9629      0.023    \PYGZhy{}83.872      0.000      \PYGZhy{}2.009      \PYGZhy{}1.917
X3             0.5243      0.070      7.446      0.000       0.385       0.663
==============================================================================
Omnibus:                        1.025   Durbin\PYGZhy{}Watson:                   2.042
Prob(Omnibus):                  0.599   Jarque\PYGZhy{}Bera (JB):                1.117
Skew:                          \PYGZhy{}0.117   Prob(JB):                        0.572
Kurtosis:                       2.719   Cond. No.                         5.91
==============================================================================

Notes:
[1] Standard Errors assume that the covariance matrix of the errors is correctly specified.
statsmodels OLS: [ 5.20431058  1.47721924 \PYGZhy{}1.96286883  0.52430039]
Manual OLS: [ 5.20431058  1.47721924 \PYGZhy{}1.96286883  0.52430039]
\end{sphinxVerbatim}

\end{sphinxuseclass}\end{sphinxVerbatimOutput}

\end{sphinxuseclass}

\subsection{Observations}
\label{\detokenize{02:observations}}\begin{itemize}
\item {} 
\sphinxAtStartPar
The \sphinxstyleemphasis{plain \sphinxcode{\sphinxupquote{Python}}} results match \sphinxcode{\sphinxupquote{statsmodels}} results (modulo tiny numerical differences).

\item {} 
\sphinxAtStartPar
Both are near the true coefficients \sphinxcode{\sphinxupquote{(5.0, 1.5, \sphinxhyphen{}2.0, 0.5)}}.

\end{itemize}


\section{Residual Analysis and Tests}
\label{\detokenize{02:residual-analysis-and-tests}}
\sphinxAtStartPar
Once we have the fitted model, let’s analyze \sphinxstylestrong{residuals}:
\begin{equation*}
\begin{split}
  e_i = y_i - \hat{y}_i.
\end{split}
\end{equation*}
\begin{sphinxuseclass}{cell}\begin{sphinxVerbatimInput}

\begin{sphinxuseclass}{cell_input}
\begin{sphinxVerbatim}[commandchars=\\\{\}]
\PYG{n}{residuals} \PYG{o}{=} \PYG{n}{model\PYGZus{}ols}\PYG{o}{.}\PYG{n}{resid}
\PYG{n}{fitted\PYGZus{}vals} \PYG{o}{=} \PYG{n}{model\PYGZus{}ols}\PYG{o}{.}\PYG{n}{fittedvalues}
\end{sphinxVerbatim}

\end{sphinxuseclass}\end{sphinxVerbatimInput}

\end{sphinxuseclass}
\sphinxAtStartPar
When evaluating the quality of a regression model, it is crucial to check whether the underlying assumptions hold. Two common residual diagnostic plots are:
\begin{itemize}
\item {} 
\sphinxAtStartPar
The residuals vs. fitted plot

\item {} 
\sphinxAtStartPar
The QQ\sphinxhyphen{}plot.

\end{itemize}

\sphinxAtStartPar
We present both in details in the following sections.


\subsection{Residuals vs. Fitted Plot}
\label{\detokenize{02:residuals-vs-fitted-plot}}\begin{itemize}
\item {} 
\sphinxAtStartPar
\sphinxstylestrong{Purpose}: This plot helps detect \sphinxstylestrong{non\sphinxhyphen{}linearity} and \sphinxstylestrong{heteroscedasticity} (unequal variance of residuals).

\item {} 
\sphinxAtStartPar
\sphinxstylestrong{How it works}:
\begin{itemize}
\item {} 
\sphinxAtStartPar
The \sphinxstylestrong{x\sphinxhyphen{}axis} represents the fitted (predicted) values from the regression model.

\item {} 
\sphinxAtStartPar
The \sphinxstylestrong{y\sphinxhyphen{}axis} represents the residuals (errors).

\item {} 
\sphinxAtStartPar
Ideally, the residuals should be randomly scattered around zero, forming a \sphinxstylestrong{homogeneous “cloud”}.

\end{itemize}

\item {} 
\sphinxAtStartPar
\sphinxstylestrong{Interpretation}:
\begin{itemize}
\item {} 
\sphinxAtStartPar
If a clear pattern (e.g., a curve) appears, it suggests \sphinxstylestrong{non\sphinxhyphen{}linearity}, meaning the model might be missing key features or transformations.

\item {} 
\sphinxAtStartPar
If the spread of residuals \sphinxstylestrong{increases or decreases systematically}, it indicates \sphinxstylestrong{heteroscedasticity}, meaning the variance of errors is not constant, which violates regression assumptions.

\end{itemize}

\end{itemize}

\begin{sphinxuseclass}{cell}\begin{sphinxVerbatimInput}

\begin{sphinxuseclass}{cell_input}
\begin{sphinxVerbatim}[commandchars=\\\{\}]
\PYG{n}{plt}\PYG{o}{.}\PYG{n}{figure}\PYG{p}{(}\PYG{n}{figsize}\PYG{o}{=}\PYG{p}{(}\PYG{l+m+mi}{10}\PYG{p}{,} \PYG{l+m+mi}{6}\PYG{p}{)}\PYG{p}{)}
\PYG{n}{plt}\PYG{o}{.}\PYG{n}{scatter}\PYG{p}{(}\PYG{n}{fitted\PYGZus{}vals}\PYG{p}{,} \PYG{n}{residuals}\PYG{p}{,} \PYG{n}{alpha}\PYG{o}{=}\PYG{l+m+mf}{0.5}\PYG{p}{)}
\PYG{n}{plt}\PYG{o}{.}\PYG{n}{axhline}\PYG{p}{(}\PYG{n}{y}\PYG{o}{=}\PYG{l+m+mi}{0}\PYG{p}{,} \PYG{n}{color}\PYG{o}{=}\PYG{l+s+s1}{\PYGZsq{}}\PYG{l+s+s1}{green}\PYG{l+s+s1}{\PYGZsq{}}\PYG{p}{,} \PYG{n}{linestyle}\PYG{o}{=}\PYG{l+s+s1}{\PYGZsq{}}\PYG{l+s+s1}{\PYGZhy{}\PYGZhy{}}\PYG{l+s+s1}{\PYGZsq{}}\PYG{p}{)}
\PYG{n}{plt}\PYG{o}{.}\PYG{n}{xlabel}\PYG{p}{(}\PYG{l+s+s2}{\PYGZdq{}}\PYG{l+s+s2}{Fitted Values}\PYG{l+s+s2}{\PYGZdq{}}\PYG{p}{)}
\PYG{n}{plt}\PYG{o}{.}\PYG{n}{ylabel}\PYG{p}{(}\PYG{l+s+s2}{\PYGZdq{}}\PYG{l+s+s2}{Residuals}\PYG{l+s+s2}{\PYGZdq{}}\PYG{p}{)}
\PYG{n}{plt}\PYG{o}{.}\PYG{n}{title}\PYG{p}{(}\PYG{l+s+s2}{\PYGZdq{}}\PYG{l+s+s2}{Residuals vs. Fitted}\PYG{l+s+s2}{\PYGZdq{}}\PYG{p}{)}
\PYG{n}{plt}\PYG{o}{.}\PYG{n}{grid}\PYG{p}{(}\PYG{l+s+s2}{\PYGZdq{}}\PYG{l+s+s2}{on}\PYG{l+s+s2}{\PYGZdq{}}\PYG{p}{)}
\PYG{n}{plt}\PYG{o}{.}\PYG{n}{show}\PYG{p}{(}\PYG{p}{)}
\end{sphinxVerbatim}

\end{sphinxuseclass}\end{sphinxVerbatimInput}
\begin{sphinxVerbatimOutput}

\begin{sphinxuseclass}{cell_output}
\noindent\sphinxincludegraphics{{1921734cb73ae8ebd1a30b01ed0947980b222fdd39ece0f80010ebae317d706c}.png}

\end{sphinxuseclass}\end{sphinxVerbatimOutput}

\end{sphinxuseclass}

\subsection{Normal Q\sphinxhyphen{}Q Plot (Quantile\sphinxhyphen{}Quantile Plot)}
\label{\detokenize{02:normal-q-q-plot-quantile-quantile-plot}}\begin{itemize}
\item {} 
\sphinxAtStartPar
\sphinxstylestrong{Purpose}: Checks whether the residuals follow a \sphinxstylestrong{normal distribution}, a key assumption in many regression analyses.

\item {} 
\sphinxAtStartPar
\sphinxstylestrong{How it works}:
\begin{itemize}
\item {} 
\sphinxAtStartPar
The \sphinxstylestrong{x\sphinxhyphen{}axis} shows the theoretical quantiles (expected values under normality).

\item {} 
\sphinxAtStartPar
The \sphinxstylestrong{y\sphinxhyphen{}axis} shows the actual quantiles from the residuals.

\item {} 
\sphinxAtStartPar
If the residuals are normally distributed, they should closely follow a \sphinxstylestrong{straight diagonal line}.

\end{itemize}

\item {} 
\sphinxAtStartPar
\sphinxstylestrong{Interpretation}:
\begin{itemize}
\item {} 
\sphinxAtStartPar
\sphinxstylestrong{Straight\sphinxhyphen{}line alignment}: Residuals are normally distributed.

\item {} 
\sphinxAtStartPar
\sphinxstylestrong{Deviations (e.g., S\sphinxhyphen{}shape, heavy tails, skewed distribution)}: Suggests non\sphinxhyphen{}normal residuals, which might indicate missing variables, skewed data, or influential outliers.

\end{itemize}

\end{itemize}

\begin{sphinxuseclass}{cell}\begin{sphinxVerbatimInput}

\begin{sphinxuseclass}{cell_input}
\begin{sphinxVerbatim}[commandchars=\\\{\}]
\PYG{n}{plt}\PYG{o}{.}\PYG{n}{figure}\PYG{p}{(}\PYG{n}{figsize}\PYG{o}{=}\PYG{p}{(}\PYG{l+m+mi}{10}\PYG{p}{,} \PYG{l+m+mi}{6}\PYG{p}{)}\PYG{p}{)}
\PYG{n}{probplot}\PYG{p}{(}\PYG{n}{residuals}\PYG{p}{,} \PYG{n}{dist}\PYG{o}{=}\PYG{l+s+s2}{\PYGZdq{}}\PYG{l+s+s2}{norm}\PYG{l+s+s2}{\PYGZdq{}}\PYG{p}{,} \PYG{n}{plot}\PYG{o}{=}\PYG{n}{plt}\PYG{p}{)}
\PYG{n}{plt}\PYG{o}{.}\PYG{n}{title}\PYG{p}{(}\PYG{l+s+s2}{\PYGZdq{}}\PYG{l+s+s2}{Normal Q\PYGZhy{}Q Plot of Residuals}\PYG{l+s+s2}{\PYGZdq{}}\PYG{p}{)}
\PYG{n}{plt}\PYG{o}{.}\PYG{n}{grid}\PYG{p}{(}\PYG{l+s+s2}{\PYGZdq{}}\PYG{l+s+s2}{on}\PYG{l+s+s2}{\PYGZdq{}}\PYG{p}{)}
\PYG{n}{plt}\PYG{o}{.}\PYG{n}{show}\PYG{p}{(}\PYG{p}{)}
\end{sphinxVerbatim}

\end{sphinxuseclass}\end{sphinxVerbatimInput}
\begin{sphinxVerbatimOutput}

\begin{sphinxuseclass}{cell_output}
\noindent\sphinxincludegraphics{{b12fe6781548d7942f02ad291a9bc90f78ff77e5523831dd2353856a7bf1db05}.png}

\end{sphinxuseclass}\end{sphinxVerbatimOutput}

\end{sphinxuseclass}

\subsubsection{Comments}
\label{\detokenize{02:comments}}\begin{itemize}
\item {} 
\sphinxAtStartPar
Ideally, residuals vs. fitted values have \sphinxstylestrong{no clear pattern} and appear \sphinxstyleemphasis{cloud\sphinxhyphen{}like}.

\item {} 
\sphinxAtStartPar
The Q\sphinxhyphen{}Q plot should be fairly linear if they are normally distributed.

\end{itemize}

\sphinxAtStartPar
We can also check the standard deviation or run further tests (e.g., Shapiro\sphinxhyphen{}Wilk, Breusch\sphinxhyphen{}Pagan for heteroscedasticity) if desired.


\section{Input Transformations \sphinxhyphen{} Polynomial Features}
\label{\detokenize{02:input-transformations-polynomial-features}}
\sphinxAtStartPar
When we talk about \sphinxstylestrong{linear regression}, we mean that the model is \sphinxstylestrong{linear in parameters}, not necessarily in the input variables \(X\). That is, the target variable \(y\) is modeled as a \sphinxstylestrong{linear combination of parameters}:
\begin{equation*}
\begin{split}
y = \beta_0 + \beta_1 f_1(X) + \beta_2 f_2(X) + \dots + \beta_n f_n(X) + \epsilon
\end{split}
\end{equation*}
\sphinxAtStartPar
where:
\begin{itemize}
\item {} 
\sphinxAtStartPar
\( \beta_0, \beta_1, \dots, \beta_n \) are the regression coefficients.

\item {} 
\sphinxAtStartPar
\( f_i(X) \) are \sphinxstylestrong{functions of the input features} (which can be non\sphinxhyphen{}linear transformations like squares or interaction terms).

\item {} 
\sphinxAtStartPar
\( \epsilon \) is the error term.

\end{itemize}


\subsection{Example: Still a Linear Model}
\label{\detokenize{02:example-still-a-linear-model}}
\sphinxAtStartPar
Even if we introduce \sphinxstylestrong{non\sphinxhyphen{}linear transformations} of \(X\), as long as the model is \sphinxstylestrong{linear in the parameters} \( \beta_i \), it is still considered a \sphinxstylestrong{linear model}. For example:
\begin{equation*}
\begin{split}
y = \beta_0 + \beta_1 X + \beta_2 X^2 + \epsilon
\end{split}
\end{equation*}
\sphinxAtStartPar
This model includes a quadratic term (\(X^2\)), but is still a \sphinxstylestrong{linear regression model} because it remains a linear combination of the parameters \(\beta_0, \beta_1, \beta_2\).


\subsection{What Would NOT Be a Linear Model?}
\label{\detokenize{02:what-would-not-be-a-linear-model}}
\sphinxAtStartPar
A regression model becomes \sphinxstylestrong{non\sphinxhyphen{}linear in parameters} if the parameters appear in a non\sphinxhyphen{}linear way. For example:
\begin{equation*}
\begin{split}
y = \beta_0 + e^{\beta_1 X} + \epsilon
\end{split}
\end{equation*}
\sphinxAtStartPar
or
\begin{equation*}
\begin{split}
y = \frac{1}{\beta_0 + \beta_1 X} + \epsilon
\end{split}
\end{equation*}
\sphinxAtStartPar
These models are \sphinxstylestrong{non\sphinxhyphen{}linear in parameters} because the relationship between \(y\) and the \(\beta\) coefficients is no longer a simple linear combination.


\section{Why Use Polynomial Features in Linear Regression?}
\label{\detokenize{02:why-use-polynomial-features-in-linear-regression}}
\sphinxAtStartPar
Since linear regression allows us to use \sphinxstylestrong{non\sphinxhyphen{}linear transformations of \(X\)} while keeping the model linear in parameters, we can \sphinxstylestrong{extend linear regression} to capture more complex relationships.


\subsection{Polynomial Feature Expansion}
\label{\detokenize{02:polynomial-feature-expansion}}
\sphinxAtStartPar
Given original features \( X_1, X_2, X_3 \), we can introduce:
\begin{itemize}
\item {} 
\sphinxAtStartPar
\sphinxstylestrong{Squared terms}: \( X_1^2, X_2^2, X_3^2 \) → captures curvature.

\item {} 
\sphinxAtStartPar
\sphinxstylestrong{Interaction terms}: \( X_1X_2, X_1X_3, X_2X_3 \) → captures relationships between variables.

\end{itemize}

\sphinxAtStartPar
By expanding the features in this way, we allow a linear regression model to \sphinxstylestrong{approximate non\sphinxhyphen{}linear patterns} in the data, making it more flexible while maintaining its interpretability.


\subsection{Example}
\label{\detokenize{02:example}}
\sphinxAtStartPar
A second\sphinxhyphen{}degree polynomial regression model:
\begin{equation*}
\begin{split}
y = \beta_0 + \beta_1 X_1 + \beta_2 X_2 + \beta_3 X_3 + \beta_4 X_1^2 + \beta_5 X_2^2 + \beta_6 X_3^2 + \beta_7 X_1X_2 + \beta_8 X_1X_3 + \beta_9 X_2X_3 + \epsilon
\end{split}
\end{equation*}
\sphinxAtStartPar
This is still \sphinxstylestrong{a linear model in parameters} but now captures non\sphinxhyphen{}linear relationships in \(X\).


\subsection{How to Implement Polynomial Features in Python}
\label{\detokenize{02:how-to-implement-polynomial-features-in-python}}
\sphinxAtStartPar
In Python, we can generate polynomial features using \sphinxcode{\sphinxupquote{sklearn.preprocessing.PolynomialFeatures}}.

\begin{sphinxuseclass}{cell}\begin{sphinxVerbatimInput}

\begin{sphinxuseclass}{cell_input}
\begin{sphinxVerbatim}[commandchars=\\\{\}]
\PYG{k+kn}{from}\PYG{+w}{ }\PYG{n+nn}{sklearn}\PYG{n+nn}{.}\PYG{n+nn}{preprocessing}\PYG{+w}{ }\PYG{k+kn}{import} \PYG{n}{PolynomialFeatures}

\PYG{n}{poly} \PYG{o}{=} \PYG{n}{PolynomialFeatures}\PYG{p}{(}\PYG{n}{degree}\PYG{o}{=}\PYG{l+m+mi}{2}\PYG{p}{,} \PYG{n}{include\PYGZus{}bias}\PYG{o}{=}\PYG{k+kc}{False}\PYG{p}{)}
\PYG{n}{X\PYGZus{}poly} \PYG{o}{=} \PYG{n}{poly}\PYG{o}{.}\PYG{n}{fit\PYGZus{}transform}\PYG{p}{(}\PYG{n}{df}\PYG{p}{[}\PYG{p}{[}\PYG{l+s+s1}{\PYGZsq{}}\PYG{l+s+s1}{X1}\PYG{l+s+s1}{\PYGZsq{}}\PYG{p}{,}\PYG{l+s+s1}{\PYGZsq{}}\PYG{l+s+s1}{X2}\PYG{l+s+s1}{\PYGZsq{}}\PYG{p}{,}\PYG{l+s+s1}{\PYGZsq{}}\PYG{l+s+s1}{X3}\PYG{l+s+s1}{\PYGZsq{}}\PYG{p}{]}\PYG{p}{]}\PYG{p}{)}
\PYG{n+nb}{print}\PYG{p}{(}\PYG{l+s+s2}{\PYGZdq{}}\PYG{l+s+s2}{Original shape:}\PYG{l+s+s2}{\PYGZdq{}}\PYG{p}{,} \PYG{n}{df}\PYG{p}{[}\PYG{p}{[}\PYG{l+s+s1}{\PYGZsq{}}\PYG{l+s+s1}{X1}\PYG{l+s+s1}{\PYGZsq{}}\PYG{p}{,}\PYG{l+s+s1}{\PYGZsq{}}\PYG{l+s+s1}{X2}\PYG{l+s+s1}{\PYGZsq{}}\PYG{p}{,}\PYG{l+s+s1}{\PYGZsq{}}\PYG{l+s+s1}{X3}\PYG{l+s+s1}{\PYGZsq{}}\PYG{p}{]}\PYG{p}{]}\PYG{o}{.}\PYG{n}{shape}\PYG{p}{)}
\PYG{n+nb}{print}\PYG{p}{(}\PYG{l+s+s2}{\PYGZdq{}}\PYG{l+s+s2}{Polynomial shape:}\PYG{l+s+s2}{\PYGZdq{}}\PYG{p}{,} \PYG{n}{X\PYGZus{}poly}\PYG{o}{.}\PYG{n}{shape}\PYG{p}{)}

\PYG{n}{X\PYGZus{}poly\PYGZus{}mat} \PYG{o}{=} \PYG{n}{np}\PYG{o}{.}\PYG{n}{column\PYGZus{}stack}\PYG{p}{(}\PYG{p}{[}\PYG{n}{np}\PYG{o}{.}\PYG{n}{ones}\PYG{p}{(}\PYG{n}{N}\PYG{p}{)}\PYG{p}{,} \PYG{n}{X\PYGZus{}poly}\PYG{p}{]}\PYG{p}{)}
\end{sphinxVerbatim}

\end{sphinxuseclass}\end{sphinxVerbatimInput}
\begin{sphinxVerbatimOutput}

\begin{sphinxuseclass}{cell_output}
\begin{sphinxVerbatim}[commandchars=\\\{\}]
Original shape: (200, 3)
Polynomial shape: (200, 9)
\end{sphinxVerbatim}

\end{sphinxuseclass}\end{sphinxVerbatimOutput}

\end{sphinxuseclass}
\sphinxAtStartPar
The OLS is obtained using the same closed\sphinxhyphen{}form formulation:

\begin{sphinxuseclass}{cell}\begin{sphinxVerbatimInput}

\begin{sphinxuseclass}{cell_input}
\begin{sphinxVerbatim}[commandchars=\\\{\}]
\PYG{n}{XTX\PYGZus{}poly} \PYG{o}{=} \PYG{n}{X\PYGZus{}poly\PYGZus{}mat}\PYG{o}{.}\PYG{n}{T} \PYG{o}{@} \PYG{n}{X\PYGZus{}poly\PYGZus{}mat}
\PYG{n}{XTX\PYGZus{}poly\PYGZus{}inv} \PYG{o}{=} \PYG{n}{np}\PYG{o}{.}\PYG{n}{linalg}\PYG{o}{.}\PYG{n}{inv}\PYG{p}{(}\PYG{n}{XTX\PYGZus{}poly}\PYG{p}{)}
\PYG{n}{XTy\PYGZus{}poly} \PYG{o}{=} \PYG{n}{X\PYGZus{}poly\PYGZus{}mat}\PYG{o}{.}\PYG{n}{T} \PYG{o}{@} \PYG{n}{y\PYGZus{}vec}
\PYG{n}{beta\PYGZus{}hat\PYGZus{}poly\PYGZus{}manual} \PYG{o}{=} \PYG{n}{XTX\PYGZus{}poly\PYGZus{}inv} \PYG{o}{@} \PYG{n}{XTy\PYGZus{}poly}
\PYG{n}{beta\PYGZus{}hat\PYGZus{}poly\PYGZus{}manual} \PYG{o}{=} \PYG{n}{beta\PYGZus{}hat\PYGZus{}poly\PYGZus{}manual}\PYG{o}{.}\PYG{n}{flatten}\PYG{p}{(}\PYG{p}{)}

\PYG{n+nb}{print}\PYG{p}{(}\PYG{l+s+s2}{\PYGZdq{}}\PYG{l+s+s2}{Manual OLS (polynomial) shape of beta:}\PYG{l+s+s2}{\PYGZdq{}}\PYG{p}{,} \PYG{n}{beta\PYGZus{}hat\PYGZus{}poly\PYGZus{}manual}\PYG{o}{.}\PYG{n}{shape}\PYG{p}{)}
\PYG{n}{pprint}\PYG{p}{(}\PYG{l+s+s2}{\PYGZdq{}}\PYG{l+s+se}{\PYGZbs{}\PYGZbs{}}\PYG{l+s+s2}{text}\PYG{l+s+s2}{\PYGZob{}}\PYG{l+s+s2}{beta}\PYG{l+s+s2}{\PYGZbs{}}\PYG{l+s+s2}{\PYGZus{}hat}\PYG{l+s+s2}{\PYGZbs{}}\PYG{l+s+s2}{\PYGZus{}poly}\PYG{l+s+s2}{\PYGZbs{}}\PYG{l+s+s2}{\PYGZus{}manual\PYGZcb{}:}\PYG{l+s+s2}{\PYGZdq{}}\PYG{p}{,} \PYG{n}{beta\PYGZus{}hat\PYGZus{}poly\PYGZus{}manual}\PYG{p}{)}
\end{sphinxVerbatim}

\end{sphinxuseclass}\end{sphinxVerbatimInput}
\begin{sphinxVerbatimOutput}

\begin{sphinxuseclass}{cell_output}
\begin{sphinxVerbatim}[commandchars=\\\{\}]
Manual OLS (polynomial) shape of beta: (10,)
\end{sphinxVerbatim}
\begin{equation*}
\begin{split}\displaystyle \text{beta\_hat\_poly\_manual}:\left[
\begin{array}{}
  5.0526 &  1.4773 & -1.9646 &  0.5315 &  0.0036 &  0.0054 &  0.0249 &  0.0037 &  0.0133 & -0.0026
\end{array}
\right]\end{split}
\end{equation*}
\end{sphinxuseclass}\end{sphinxVerbatimOutput}

\end{sphinxuseclass}
\sphinxAtStartPar
We can compare the results offered by \sphinxcode{\sphinxupquote{statsmodel}}:

\begin{sphinxuseclass}{cell}\begin{sphinxVerbatimInput}

\begin{sphinxuseclass}{cell_input}
\begin{sphinxVerbatim}[commandchars=\\\{\}]
\PYG{n}{X\PYGZus{}poly\PYGZus{}sm} \PYG{o}{=} \PYG{n}{sm}\PYG{o}{.}\PYG{n}{add\PYGZus{}constant}\PYG{p}{(}\PYG{n}{X\PYGZus{}poly}\PYG{p}{)}
\PYG{n}{model\PYGZus{}poly\PYGZus{}sm} \PYG{o}{=} \PYG{n}{sm}\PYG{o}{.}\PYG{n}{OLS}\PYG{p}{(}\PYG{n}{df}\PYG{p}{[}\PYG{l+s+s1}{\PYGZsq{}}\PYG{l+s+s1}{Y}\PYG{l+s+s1}{\PYGZsq{}}\PYG{p}{]}\PYG{p}{,} \PYG{n}{X\PYGZus{}poly\PYGZus{}sm}\PYG{p}{)}\PYG{o}{.}\PYG{n}{fit}\PYG{p}{(}\PYG{p}{)}
\PYG{n}{beta\PYGZus{}hat\PYGZus{}poly\PYGZus{}sm} \PYG{o}{=} \PYG{n}{model\PYGZus{}poly\PYGZus{}sm}\PYG{o}{.}\PYG{n}{params}\PYG{o}{.}\PYG{n}{values}

\PYG{n}{pprint}\PYG{p}{(}\PYG{l+s+s2}{\PYGZdq{}}\PYG{l+s+se}{\PYGZbs{}\PYGZbs{}}\PYG{l+s+s2}{text}\PYG{l+s+s2}{\PYGZob{}}\PYG{l+s+s2}{Statsmodels polynomial OLS\PYGZcb{}:}\PYG{l+s+s2}{\PYGZdq{}}\PYG{p}{,} \PYG{n}{beta\PYGZus{}hat\PYGZus{}poly\PYGZus{}sm}\PYG{p}{)}
\PYG{n}{pprint}\PYG{p}{(}\PYG{l+s+s2}{\PYGZdq{}}\PYG{l+s+se}{\PYGZbs{}\PYGZbs{}}\PYG{l+s+s2}{text}\PYG{l+s+s2}{\PYGZob{}}\PYG{l+s+s2}{Manual polynomial OLS\PYGZcb{}:}\PYG{l+s+s2}{\PYGZdq{}}\PYG{p}{,} \PYG{n}{beta\PYGZus{}hat\PYGZus{}poly\PYGZus{}manual}\PYG{p}{)}
\PYG{n+nb}{print}\PYG{p}{(}\PYG{l+s+s2}{\PYGZdq{}}\PYG{l+s+se}{\PYGZbs{}n}\PYG{l+s+s2}{Polynomial OLS summary:}\PYG{l+s+se}{\PYGZbs{}n}\PYG{l+s+s2}{\PYGZdq{}}\PYG{p}{,} \PYG{n}{model\PYGZus{}poly\PYGZus{}sm}\PYG{o}{.}\PYG{n}{summary}\PYG{p}{(}\PYG{p}{)}\PYG{p}{)}
\end{sphinxVerbatim}

\end{sphinxuseclass}\end{sphinxVerbatimInput}
\begin{sphinxVerbatimOutput}

\begin{sphinxuseclass}{cell_output}\begin{equation*}
\begin{split}\displaystyle \text{Statsmodels polynomial OLS}:\left[
\begin{array}{}
  5.0526 &  1.4773 & -1.9646 &  0.5315 &  0.0036 &  0.0054 &  0.0249 &  0.0037 &  0.0133 & -0.0026
\end{array}
\right]\end{split}
\end{equation*}\begin{equation*}
\begin{split}\displaystyle \text{Manual polynomial OLS}:\left[
\begin{array}{}
  5.0526 &  1.4773 & -1.9646 &  0.5315 &  0.0036 &  0.0054 &  0.0249 &  0.0037 &  0.0133 & -0.0026
\end{array}
\right]\end{split}
\end{equation*}
\begin{sphinxVerbatim}[commandchars=\\\{\}]
Polynomial OLS summary:
                             OLS Regression Results                            
==============================================================================
Dep. Variable:                      Y   R\PYGZhy{}squared:                       0.976
Model:                            OLS   Adj. R\PYGZhy{}squared:                  0.975
Method:                 Least Squares   F\PYGZhy{}statistic:                     872.4
Date:                Tue, 16 Sep 2025   Prob (F\PYGZhy{}statistic):          2.13e\PYGZhy{}149
Time:                        13:31:42   Log\PYGZhy{}Likelihood:                \PYGZhy{}413.46
No. Observations:                 200   AIC:                             846.9
Df Residuals:                     190   BIC:                             879.9
Df Model:                           9                                         
Covariance Type:            nonrobust                                         
==============================================================================
                 coef    std err          t      P\PYGZgt{}|t|      [0.025      0.975]
\PYGZhy{}\PYGZhy{}\PYGZhy{}\PYGZhy{}\PYGZhy{}\PYGZhy{}\PYGZhy{}\PYGZhy{}\PYGZhy{}\PYGZhy{}\PYGZhy{}\PYGZhy{}\PYGZhy{}\PYGZhy{}\PYGZhy{}\PYGZhy{}\PYGZhy{}\PYGZhy{}\PYGZhy{}\PYGZhy{}\PYGZhy{}\PYGZhy{}\PYGZhy{}\PYGZhy{}\PYGZhy{}\PYGZhy{}\PYGZhy{}\PYGZhy{}\PYGZhy{}\PYGZhy{}\PYGZhy{}\PYGZhy{}\PYGZhy{}\PYGZhy{}\PYGZhy{}\PYGZhy{}\PYGZhy{}\PYGZhy{}\PYGZhy{}\PYGZhy{}\PYGZhy{}\PYGZhy{}\PYGZhy{}\PYGZhy{}\PYGZhy{}\PYGZhy{}\PYGZhy{}\PYGZhy{}\PYGZhy{}\PYGZhy{}\PYGZhy{}\PYGZhy{}\PYGZhy{}\PYGZhy{}\PYGZhy{}\PYGZhy{}\PYGZhy{}\PYGZhy{}\PYGZhy{}\PYGZhy{}\PYGZhy{}\PYGZhy{}\PYGZhy{}\PYGZhy{}\PYGZhy{}\PYGZhy{}\PYGZhy{}\PYGZhy{}\PYGZhy{}\PYGZhy{}\PYGZhy{}\PYGZhy{}\PYGZhy{}\PYGZhy{}\PYGZhy{}\PYGZhy{}\PYGZhy{}\PYGZhy{}
const          5.0526      0.287     17.620      0.000       4.487       5.618
x1             1.4773      0.053     27.754      0.000       1.372       1.582
x2            \PYGZhy{}1.9646      0.024    \PYGZhy{}82.615      0.000      \PYGZhy{}2.012      \PYGZhy{}1.918
x3             0.5315      0.073      7.290      0.000       0.388       0.675
x4             0.0036      0.020      0.179      0.858      \PYGZhy{}0.036       0.043
x5             0.0054      0.009      0.571      0.569      \PYGZhy{}0.013       0.024
x6             0.0249      0.028      0.879      0.381      \PYGZhy{}0.031       0.081
x7             0.0037      0.005      0.803      0.423      \PYGZhy{}0.005       0.013
x8             0.0133      0.012      1.074      0.284      \PYGZhy{}0.011       0.038
x9            \PYGZhy{}0.0026      0.028     \PYGZhy{}0.095      0.924      \PYGZhy{}0.058       0.052
==============================================================================
Omnibus:                        0.925   Durbin\PYGZhy{}Watson:                   2.024
Prob(Omnibus):                  0.630   Jarque\PYGZhy{}Bera (JB):                1.006
Skew:                          \PYGZhy{}0.090   Prob(JB):                        0.605
Kurtosis:                       2.703   Cond. No.                         97.1
==============================================================================

Notes:
[1] Standard Errors assume that the covariance matrix of the errors is correctly specified.
\end{sphinxVerbatim}

\end{sphinxuseclass}\end{sphinxVerbatimOutput}

\end{sphinxuseclass}

\subsubsection{Residual analysis and QQ\sphinxhyphen{}plot}
\label{\detokenize{02:residual-analysis-and-qq-plot}}
\begin{sphinxuseclass}{cell}\begin{sphinxVerbatimInput}

\begin{sphinxuseclass}{cell_input}
\begin{sphinxVerbatim}[commandchars=\\\{\}]
\PYG{n}{residuals\PYGZus{}poly} \PYG{o}{=} \PYG{n}{model\PYGZus{}poly\PYGZus{}sm}\PYG{o}{.}\PYG{n}{resid}
\PYG{n}{fitted\PYGZus{}poly} \PYG{o}{=} \PYG{n}{model\PYGZus{}poly\PYGZus{}sm}\PYG{o}{.}\PYG{n}{fittedvalues}

\PYG{n}{plt}\PYG{o}{.}\PYG{n}{figure}\PYG{p}{(}\PYG{n}{figsize}\PYG{o}{=}\PYG{p}{(}\PYG{l+m+mi}{10}\PYG{p}{,} \PYG{l+m+mi}{6}\PYG{p}{)}\PYG{p}{)}
\PYG{n}{plt}\PYG{o}{.}\PYG{n}{scatter}\PYG{p}{(}\PYG{n}{fitted\PYGZus{}poly}\PYG{p}{,} \PYG{n}{residuals\PYGZus{}poly}\PYG{p}{,} \PYG{n}{alpha}\PYG{o}{=}\PYG{l+m+mf}{0.5}\PYG{p}{)}
\PYG{n}{plt}\PYG{o}{.}\PYG{n}{axhline}\PYG{p}{(}\PYG{n}{y}\PYG{o}{=}\PYG{l+m+mi}{0}\PYG{p}{,} \PYG{n}{color}\PYG{o}{=}\PYG{l+s+s1}{\PYGZsq{}}\PYG{l+s+s1}{green}\PYG{l+s+s1}{\PYGZsq{}}\PYG{p}{,} \PYG{n}{linestyle}\PYG{o}{=}\PYG{l+s+s1}{\PYGZsq{}}\PYG{l+s+s1}{\PYGZhy{}\PYGZhy{}}\PYG{l+s+s1}{\PYGZsq{}}\PYG{p}{)}
\PYG{n}{plt}\PYG{o}{.}\PYG{n}{xlabel}\PYG{p}{(}\PYG{l+s+s2}{\PYGZdq{}}\PYG{l+s+s2}{Fitted values (poly)}\PYG{l+s+s2}{\PYGZdq{}}\PYG{p}{)}
\PYG{n}{plt}\PYG{o}{.}\PYG{n}{ylabel}\PYG{p}{(}\PYG{l+s+s2}{\PYGZdq{}}\PYG{l+s+s2}{Residuals}\PYG{l+s+s2}{\PYGZdq{}}\PYG{p}{)}
\PYG{n}{plt}\PYG{o}{.}\PYG{n}{title}\PYG{p}{(}\PYG{l+s+s2}{\PYGZdq{}}\PYG{l+s+s2}{Residuals vs. Fitted (Polynomial Model)}\PYG{l+s+s2}{\PYGZdq{}}\PYG{p}{)}
\PYG{n}{plt}\PYG{o}{.}\PYG{n}{grid}\PYG{p}{(}\PYG{l+s+s2}{\PYGZdq{}}\PYG{l+s+s2}{on}\PYG{l+s+s2}{\PYGZdq{}}\PYG{p}{)}
\PYG{n}{plt}\PYG{o}{.}\PYG{n}{show}\PYG{p}{(}\PYG{p}{)}

\PYG{n}{plt}\PYG{o}{.}\PYG{n}{figure}\PYG{p}{(}\PYG{n}{figsize}\PYG{o}{=}\PYG{p}{(}\PYG{l+m+mi}{10}\PYG{p}{,} \PYG{l+m+mi}{6}\PYG{p}{)}\PYG{p}{)}
\PYG{n}{probplot}\PYG{p}{(}\PYG{n}{residuals\PYGZus{}poly}\PYG{p}{,} \PYG{n}{dist}\PYG{o}{=}\PYG{l+s+s2}{\PYGZdq{}}\PYG{l+s+s2}{norm}\PYG{l+s+s2}{\PYGZdq{}}\PYG{p}{,} \PYG{n}{plot}\PYG{o}{=}\PYG{n}{plt}\PYG{p}{)}
\PYG{n}{plt}\PYG{o}{.}\PYG{n}{title}\PYG{p}{(}\PYG{l+s+s2}{\PYGZdq{}}\PYG{l+s+s2}{Normal Q\PYGZhy{}Q Plot (Polynomial Model Residuals)}\PYG{l+s+s2}{\PYGZdq{}}\PYG{p}{)}
\PYG{n}{plt}\PYG{o}{.}\PYG{n}{grid}\PYG{p}{(}\PYG{l+s+s2}{\PYGZdq{}}\PYG{l+s+s2}{on}\PYG{l+s+s2}{\PYGZdq{}}\PYG{p}{)}
\PYG{n}{plt}\PYG{o}{.}\PYG{n}{show}\PYG{p}{(}\PYG{p}{)}
\end{sphinxVerbatim}

\end{sphinxuseclass}\end{sphinxVerbatimInput}
\begin{sphinxVerbatimOutput}

\begin{sphinxuseclass}{cell_output}
\noindent\sphinxincludegraphics{{058510b3c11a0fb724a1c6dfdfbdee47a0154a934032ba892466cc547793b11e}.png}

\noindent\sphinxincludegraphics{{63232262e329da9e0ce61297e428d201f3e8faa1f0c49f84f37b3be167099356}.png}

\end{sphinxuseclass}\end{sphinxVerbatimOutput}

\end{sphinxuseclass}
\sphinxAtStartPar
\sphinxstylestrong{Interpretation}:
\begin{itemize}
\item {} 
\sphinxAtStartPar
If the polynomial model captures more structure, we may see less pattern in the residuals plot.

\item {} 
\sphinxAtStartPar
Alternatively, if the true relationship was already linear, polynomial terms might just overfit or remain insignificant.

\end{itemize}


\section{RBF Expansion}
\label{\detokenize{02:rbf-expansion}}
\sphinxAtStartPar
Polynomial feature expansion is a powerful tool to introduce \sphinxstylestrong{non\sphinxhyphen{}linearity} into a linear regression model, but it has several limitations:
\begin{itemize}
\item {} 
\sphinxAtStartPar
\sphinxstylestrong{Feature explosion}: As the polynomial degree increases, the number of new features grows \sphinxstylestrong{exponentially}, making models computationally expensive and prone to overfitting.

\item {} 
\sphinxAtStartPar
\sphinxstylestrong{Global transformations}: Polynomial terms like \(X^2, X^3\) affect the entire input space, meaning that small changes in \(X\) have a large impact everywhere, which might not always be desirable.

\item {} 
\sphinxAtStartPar
\sphinxstylestrong{Poor handling of local variations}: Polynomial functions are not ideal when the relationship between \(X\) and \(y\) has local variations or is non\sphinxhyphen{}smooth.

\end{itemize}

\sphinxAtStartPar
This is where \sphinxstylestrong{Radial Basis Function (RBF) expansion} can be useful.


\subsection{What is RBF Expansion?}
\label{\detokenize{02:what-is-rbf-expansion}}
\sphinxAtStartPar
RBF expansion \sphinxstylestrong{maps the original features into a higher\sphinxhyphen{}dimensional space} using functions that respond \sphinxstylestrong{locally} to the input values.

\sphinxAtStartPar
Instead of using polynomial transformations like:
\begin{equation*}
\begin{split}
y = \beta_0 + \beta_1 X_1 + \beta_2 X_2 + \beta_3 X_3 + \beta_4 X_1^2 + \beta_5 X_2^2 + \beta_6 X_3^2 + \beta_7 X_1X_2 + ...
\end{split}
\end{equation*}
\sphinxAtStartPar
We use \sphinxstylestrong{random Fourier features} to approximate an RBF kernel, generating transformed variables like:
\begin{equation*}
\begin{split}
\Phi(X) = \cos(WX) + \sin(WX)
\end{split}
\end{equation*}
\sphinxAtStartPar
where \( W \) is a random weight matrix sampled from a Gaussian distribution.


\subsubsection{Why Use RBF Instead of Polynomial Expansion?}
\label{\detokenize{02:why-use-rbf-instead-of-polynomial-expansion}}\begin{itemize}
\item {} 
\sphinxAtStartPar
\sphinxstylestrong{Better for local variations}: Unlike polynomials, RBFs can capture \sphinxstylestrong{localized} structures in data.

\item {} 
\sphinxAtStartPar
\sphinxstylestrong{Fixed feature size}: The number of transformed features is independent of the input size, making it computationally efficient.

\item {} 
\sphinxAtStartPar
\sphinxstylestrong{Good for high\sphinxhyphen{}dimensional data}: Works well when the input space is \sphinxstylestrong{high\sphinxhyphen{}dimensional} because it avoids feature explosion.

\end{itemize}


\subsection{Implementing RBF Expansion with Random Fourier Features}
\label{\detokenize{02:implementing-rbf-expansion-with-random-fourier-features}}
\sphinxAtStartPar
We use \sphinxcode{\sphinxupquote{sklearn.kernel\_approximation.RBFSampler}} to generate \sphinxstylestrong{random Fourier features} that approximate the RBF kernel.

\begin{sphinxuseclass}{cell}\begin{sphinxVerbatimInput}

\begin{sphinxuseclass}{cell_input}
\begin{sphinxVerbatim}[commandchars=\\\{\}]
\PYG{k+kn}{from}\PYG{+w}{ }\PYG{n+nn}{sklearn}\PYG{n+nn}{.}\PYG{n+nn}{kernel\PYGZus{}approximation}\PYG{+w}{ }\PYG{k+kn}{import} \PYG{n}{RBFSampler}

\PYG{n}{rbf} \PYG{o}{=} \PYG{n}{RBFSampler}\PYG{p}{(}\PYG{n}{gamma}\PYG{o}{=}\PYG{l+m+mf}{0.05}\PYG{p}{,} \PYG{n}{n\PYGZus{}components}\PYG{o}{=}\PYG{l+m+mi}{50}\PYG{p}{,} \PYG{n}{random\PYGZus{}state}\PYG{o}{=}\PYG{l+m+mi}{123}\PYG{p}{)}
\PYG{n}{X\PYGZus{}rbf} \PYG{o}{=} \PYG{n}{rbf}\PYG{o}{.}\PYG{n}{fit\PYGZus{}transform}\PYG{p}{(}\PYG{n}{df}\PYG{p}{[}\PYG{p}{[}\PYG{l+s+s1}{\PYGZsq{}}\PYG{l+s+s1}{X1}\PYG{l+s+s1}{\PYGZsq{}}\PYG{p}{,}\PYG{l+s+s1}{\PYGZsq{}}\PYG{l+s+s1}{X2}\PYG{l+s+s1}{\PYGZsq{}}\PYG{p}{,}\PYG{l+s+s1}{\PYGZsq{}}\PYG{l+s+s1}{X3}\PYG{l+s+s1}{\PYGZsq{}}\PYG{p}{]}\PYG{p}{]}\PYG{p}{)}
\PYG{n+nb}{print}\PYG{p}{(}\PYG{l+s+s2}{\PYGZdq{}}\PYG{l+s+s2}{RBF\PYGZhy{}transformed shape:}\PYG{l+s+s2}{\PYGZdq{}}\PYG{p}{,} \PYG{n}{X\PYGZus{}rbf}\PYG{o}{.}\PYG{n}{shape}\PYG{p}{)}

\PYG{c+c1}{\PYGZsh{} We\PYGZsq{}ll do OLS with statsmodels for brevity:}
\PYG{n}{X\PYGZus{}rbf\PYGZus{}sm} \PYG{o}{=} \PYG{n}{sm}\PYG{o}{.}\PYG{n}{add\PYGZus{}constant}\PYG{p}{(}\PYG{n}{X\PYGZus{}rbf}\PYG{p}{)}
\PYG{n}{model\PYGZus{}rbf\PYGZus{}sm} \PYG{o}{=} \PYG{n}{sm}\PYG{o}{.}\PYG{n}{OLS}\PYG{p}{(}\PYG{n}{df}\PYG{p}{[}\PYG{l+s+s1}{\PYGZsq{}}\PYG{l+s+s1}{Y}\PYG{l+s+s1}{\PYGZsq{}}\PYG{p}{]}\PYG{p}{,} \PYG{n}{X\PYGZus{}rbf\PYGZus{}sm}\PYG{p}{)}\PYG{o}{.}\PYG{n}{fit}\PYG{p}{(}\PYG{p}{)}

\PYG{n+nb}{print}\PYG{p}{(}\PYG{n}{model\PYGZus{}rbf\PYGZus{}sm}\PYG{o}{.}\PYG{n}{summary}\PYG{p}{(}\PYG{p}{)}\PYG{p}{)}

\PYG{n}{resid\PYGZus{}rbf} \PYG{o}{=} \PYG{n}{model\PYGZus{}rbf\PYGZus{}sm}\PYG{o}{.}\PYG{n}{resid}
\PYG{n}{fitted\PYGZus{}rbf} \PYG{o}{=} \PYG{n}{model\PYGZus{}rbf\PYGZus{}sm}\PYG{o}{.}\PYG{n}{fittedvalues}

\PYG{n}{plt}\PYG{o}{.}\PYG{n}{figure}\PYG{p}{(}\PYG{n}{figsize}\PYG{o}{=}\PYG{p}{(}\PYG{l+m+mi}{10}\PYG{p}{,} \PYG{l+m+mi}{6}\PYG{p}{)}\PYG{p}{)}
\PYG{n}{plt}\PYG{o}{.}\PYG{n}{scatter}\PYG{p}{(}\PYG{n}{fitted\PYGZus{}rbf}\PYG{p}{,} \PYG{n}{resid\PYGZus{}rbf}\PYG{p}{,} \PYG{n}{alpha}\PYG{o}{=}\PYG{l+m+mf}{0.5}\PYG{p}{)}
\PYG{n}{plt}\PYG{o}{.}\PYG{n}{axhline}\PYG{p}{(}\PYG{n}{y}\PYG{o}{=}\PYG{l+m+mi}{0}\PYG{p}{,} \PYG{n}{color}\PYG{o}{=}\PYG{l+s+s1}{\PYGZsq{}}\PYG{l+s+s1}{green}\PYG{l+s+s1}{\PYGZsq{}}\PYG{p}{,} \PYG{n}{linestyle}\PYG{o}{=}\PYG{l+s+s1}{\PYGZsq{}}\PYG{l+s+s1}{\PYGZhy{}\PYGZhy{}}\PYG{l+s+s1}{\PYGZsq{}}\PYG{p}{)}
\PYG{n}{plt}\PYG{o}{.}\PYG{n}{xlabel}\PYG{p}{(}\PYG{l+s+s2}{\PYGZdq{}}\PYG{l+s+s2}{Fitted values (RBF)}\PYG{l+s+s2}{\PYGZdq{}}\PYG{p}{)}
\PYG{n}{plt}\PYG{o}{.}\PYG{n}{ylabel}\PYG{p}{(}\PYG{l+s+s2}{\PYGZdq{}}\PYG{l+s+s2}{Residuals}\PYG{l+s+s2}{\PYGZdq{}}\PYG{p}{)}
\PYG{n}{plt}\PYG{o}{.}\PYG{n}{title}\PYG{p}{(}\PYG{l+s+s2}{\PYGZdq{}}\PYG{l+s+s2}{Residuals vs. Fitted (RBF)}\PYG{l+s+s2}{\PYGZdq{}}\PYG{p}{)}
\PYG{n}{plt}\PYG{o}{.}\PYG{n}{grid}\PYG{p}{(}\PYG{l+s+s2}{\PYGZdq{}}\PYG{l+s+s2}{on}\PYG{l+s+s2}{\PYGZdq{}}\PYG{p}{)}
\PYG{n}{plt}\PYG{o}{.}\PYG{n}{show}\PYG{p}{(}\PYG{p}{)}
\end{sphinxVerbatim}

\end{sphinxuseclass}\end{sphinxVerbatimInput}
\begin{sphinxVerbatimOutput}

\begin{sphinxuseclass}{cell_output}
\begin{sphinxVerbatim}[commandchars=\\\{\}]
RBF\PYGZhy{}transformed shape: (200, 50)
                            OLS Regression Results                            
==============================================================================
Dep. Variable:                      Y   R\PYGZhy{}squared:                       0.980
Model:                            OLS   Adj. R\PYGZhy{}squared:                  0.973
Method:                 Least Squares   F\PYGZhy{}statistic:                     146.9
Date:                Tue, 16 Sep 2025   Prob (F\PYGZhy{}statistic):          5.70e\PYGZhy{}105
Time:                        13:31:42   Log\PYGZhy{}Likelihood:                \PYGZhy{}396.22
No. Observations:                 200   AIC:                             894.4
Df Residuals:                     149   BIC:                             1063.
Df Model:                          50                                         
Covariance Type:            nonrobust                                         
==============================================================================
                 coef    std err          t      P\PYGZgt{}|t|      [0.025      0.975]
\PYGZhy{}\PYGZhy{}\PYGZhy{}\PYGZhy{}\PYGZhy{}\PYGZhy{}\PYGZhy{}\PYGZhy{}\PYGZhy{}\PYGZhy{}\PYGZhy{}\PYGZhy{}\PYGZhy{}\PYGZhy{}\PYGZhy{}\PYGZhy{}\PYGZhy{}\PYGZhy{}\PYGZhy{}\PYGZhy{}\PYGZhy{}\PYGZhy{}\PYGZhy{}\PYGZhy{}\PYGZhy{}\PYGZhy{}\PYGZhy{}\PYGZhy{}\PYGZhy{}\PYGZhy{}\PYGZhy{}\PYGZhy{}\PYGZhy{}\PYGZhy{}\PYGZhy{}\PYGZhy{}\PYGZhy{}\PYGZhy{}\PYGZhy{}\PYGZhy{}\PYGZhy{}\PYGZhy{}\PYGZhy{}\PYGZhy{}\PYGZhy{}\PYGZhy{}\PYGZhy{}\PYGZhy{}\PYGZhy{}\PYGZhy{}\PYGZhy{}\PYGZhy{}\PYGZhy{}\PYGZhy{}\PYGZhy{}\PYGZhy{}\PYGZhy{}\PYGZhy{}\PYGZhy{}\PYGZhy{}\PYGZhy{}\PYGZhy{}\PYGZhy{}\PYGZhy{}\PYGZhy{}\PYGZhy{}\PYGZhy{}\PYGZhy{}\PYGZhy{}\PYGZhy{}\PYGZhy{}\PYGZhy{}\PYGZhy{}\PYGZhy{}\PYGZhy{}\PYGZhy{}\PYGZhy{}\PYGZhy{}
const        \PYGZhy{}43.4898     19.729     \PYGZhy{}2.204      0.029     \PYGZhy{}82.475      \PYGZhy{}4.505
x1             7.1364     11.512      0.620      0.536     \PYGZhy{}15.612      29.885
x2             8.9773      3.334      2.692      0.008       2.388      15.566
x3             2.6549      7.719      0.344      0.731     \PYGZhy{}12.598      17.908
x4             1.1046      1.872      0.590      0.556      \PYGZhy{}2.595       4.804
x5           137.0639     46.449      2.951      0.004      45.280     228.848
x6            16.5251     11.039      1.497      0.137      \PYGZhy{}5.288      38.338
x7            \PYGZhy{}3.4503      4.473     \PYGZhy{}0.771      0.442     \PYGZhy{}12.288       5.388
x8             2.9028      3.919      0.741      0.460      \PYGZhy{}4.842      10.648
x9            \PYGZhy{}0.5036      6.801     \PYGZhy{}0.074      0.941     \PYGZhy{}13.942      12.934
x10           16.0558     15.230      1.054      0.294     \PYGZhy{}14.040      46.151
x11           24.7962     12.423      1.996      0.048       0.249      49.344
x12           59.4056      7.430      7.995      0.000      44.723      74.088
x13           \PYGZhy{}6.8051     11.750     \PYGZhy{}0.579      0.563     \PYGZhy{}30.024      16.414
x14           12.4240     24.130      0.515      0.607     \PYGZhy{}35.258      60.106
x15          \PYGZhy{}19.3450      7.617     \PYGZhy{}2.540      0.012     \PYGZhy{}34.397      \PYGZhy{}4.293
x16         \PYGZhy{}266.9437     65.123     \PYGZhy{}4.099      0.000    \PYGZhy{}395.627    \PYGZhy{}138.260
x17            5.0042      4.000      1.251      0.213      \PYGZhy{}2.900      12.908
x18           11.1227      5.828      1.909      0.058      \PYGZhy{}0.393      22.639
x19           24.4491     13.325      1.835      0.069      \PYGZhy{}1.882      50.780
x20           \PYGZhy{}6.4866      3.810     \PYGZhy{}1.702      0.091     \PYGZhy{}14.016       1.042
x21         \PYGZhy{}136.7679     40.658     \PYGZhy{}3.364      0.001    \PYGZhy{}217.108     \PYGZhy{}56.428
x22           \PYGZhy{}0.5810      7.714     \PYGZhy{}0.075      0.940     \PYGZhy{}15.825      14.663
x23           29.9036     16.223      1.843      0.067      \PYGZhy{}2.153      61.960
x24           43.4672     31.745      1.369      0.173     \PYGZhy{}19.261     106.195
x25           \PYGZhy{}2.5586      1.587     \PYGZhy{}1.612      0.109      \PYGZhy{}5.695       0.578
x26           24.3295     50.242      0.484      0.629     \PYGZhy{}74.950     123.609
x27           24.6458     14.259      1.728      0.086      \PYGZhy{}3.530      52.821
x28           \PYGZhy{}9.9883     14.837     \PYGZhy{}0.673      0.502     \PYGZhy{}39.306      19.329
x29           98.4893     31.798      3.097      0.002      35.657     161.322
x30           \PYGZhy{}5.2150     10.521     \PYGZhy{}0.496      0.621     \PYGZhy{}26.004      15.574
x31           \PYGZhy{}2.7180     16.134     \PYGZhy{}0.168      0.866     \PYGZhy{}34.600      29.164
x32            2.4668      1.976      1.248      0.214      \PYGZhy{}1.438       6.372
x33            2.5672      2.217      1.158      0.249      \PYGZhy{}1.814       6.948
x34          200.9372     54.077      3.716      0.000      94.081     307.793
x35           42.1785     14.095      2.992      0.003      14.327      70.031
x36          \PYGZhy{}24.5715    128.680     \PYGZhy{}0.191      0.849    \PYGZhy{}278.845     229.702
x37          156.2145     31.866      4.902      0.000      93.246     219.183
x38           \PYGZhy{}7.9412     16.226     \PYGZhy{}0.489      0.625     \PYGZhy{}40.004      24.122
x39           \PYGZhy{}1.1294      1.594     \PYGZhy{}0.709      0.480      \PYGZhy{}4.279       2.020
x40         \PYGZhy{}225.6936     88.747     \PYGZhy{}2.543      0.012    \PYGZhy{}401.059     \PYGZhy{}50.328
x41          \PYGZhy{}12.5555     36.848     \PYGZhy{}0.341      0.734     \PYGZhy{}85.367      60.256
x42          \PYGZhy{}12.9924      9.099     \PYGZhy{}1.428      0.155     \PYGZhy{}30.972       4.987
x43           25.7314     13.007      1.978      0.050       0.029      51.433
x44            7.2633     15.228      0.477      0.634     \PYGZhy{}22.827      37.353
x45          \PYGZhy{}70.8756     20.960     \PYGZhy{}3.381      0.001    \PYGZhy{}112.293     \PYGZhy{}29.458
x46          \PYGZhy{}46.6215    172.329     \PYGZhy{}0.271      0.787    \PYGZhy{}387.147     293.904
x47            0.5652      3.121      0.181      0.857      \PYGZhy{}5.602       6.732
x48           \PYGZhy{}6.3771     10.449     \PYGZhy{}0.610      0.543     \PYGZhy{}27.025      14.271
x49          153.3197     46.325      3.310      0.001      61.781     244.858
x50            9.1846      5.433      1.691      0.093      \PYGZhy{}1.551      19.920
==============================================================================
Omnibus:                        0.557   Durbin\PYGZhy{}Watson:                   2.049
Prob(Omnibus):                  0.757   Jarque\PYGZhy{}Bera (JB):                0.682
Skew:                          \PYGZhy{}0.041   Prob(JB):                        0.711
Kurtosis:                       2.726   Cond. No.                     1.75e+03
==============================================================================

Notes:
[1] Standard Errors assume that the covariance matrix of the errors is correctly specified.
[2] The condition number is large, 1.75e+03. This might indicate that there are
strong multicollinearity or other numerical problems.
\end{sphinxVerbatim}

\noindent\sphinxincludegraphics{{780d93454acec8c3ab6c4f25370c0190e0db903a733cec5cb2002e60c7327e26}.png}

\end{sphinxuseclass}\end{sphinxVerbatimOutput}

\end{sphinxuseclass}

\subsubsection{Comments}
\label{\detokenize{02:id1}}\begin{itemize}
\item {} 
\sphinxAtStartPar
Each RBF feature is basically a sinusoidal basis function in random projection space.

\item {} 
\sphinxAtStartPar
We could also attempt a manual approach:
\begin{equation*}
\begin{split}
    \hat{\beta} = (X_{\mathrm{rbf}}^\top X_{\mathrm{rbf}})^{-1} X_{\mathrm{rbf}}^\top y
  \end{split}
\end{equation*}
\sphinxAtStartPar
if we want to replicate the normal equation style.

\item {} 
\sphinxAtStartPar
The dimensionality (50) can be large, so be mindful of potential numerical issues and overfitting.

\end{itemize}


\section{Ridge Regression}
\label{\detokenize{02:ridge-regression}}
\sphinxAtStartPar
For complex expansions (like polynomial of high degree or many RBF components), \sphinxstylestrong{overfitting} can happen.
Ridge regression adds an L2 penalty:
\begin{equation*}
\begin{split}
  \min_{\beta} \;\; \| Y - X \beta \|^2 + \alpha \| \beta \|^2.
\end{split}
\end{equation*}

\subsection{Ridge Closed\sphinxhyphen{}form}
\label{\detokenize{02:ridge-closed-form}}
\sphinxAtStartPar
Ridge also has a known closed\sphinxhyphen{}form:
\begin{equation*}
\begin{split}
  \hat{\beta}_{\mathrm{ridge}} = (X^\top X + \alpha I)^{-1} X^\top Y.
\end{split}
\end{equation*}
\sphinxAtStartPar
We’ll show a manual solution, then compare to scikit\sphinxhyphen{}learn. We’ll do it on the \sphinxstylestrong{polynomial features} for illustration.

\begin{sphinxuseclass}{cell}\begin{sphinxVerbatimInput}

\begin{sphinxuseclass}{cell_input}
\begin{sphinxVerbatim}[commandchars=\\\{\}]
\PYG{n}{alpha\PYGZus{}ridge} \PYG{o}{=} \PYG{l+m+mf}{10.0}
\PYG{n}{I\PYGZus{}p} \PYG{o}{=} \PYG{n}{np}\PYG{o}{.}\PYG{n}{eye}\PYG{p}{(}\PYG{n}{X\PYGZus{}poly\PYGZus{}mat}\PYG{o}{.}\PYG{n}{shape}\PYG{p}{[}\PYG{l+m+mi}{1}\PYG{p}{]}\PYG{p}{)}  \PYG{c+c1}{\PYGZsh{} dimension p x p, where p = number of poly features + 1 (intercept)}

\PYG{c+c1}{\PYGZsh{} Usually we do not regularize the intercept, so we might set I\PYGZus{}p[0,0] = 0.}
\PYG{c+c1}{\PYGZsh{} Let\PYGZsq{}s do that to avoid shrinking the intercept:}
\PYG{n}{I\PYGZus{}p}\PYG{p}{[}\PYG{l+m+mi}{0}\PYG{p}{,}\PYG{l+m+mi}{0}\PYG{p}{]} \PYG{o}{=} \PYG{l+m+mf}{0.0}

\PYG{n}{XTX\PYGZus{}ridge} \PYG{o}{=} \PYG{n}{X\PYGZus{}poly\PYGZus{}mat}\PYG{o}{.}\PYG{n}{T} \PYG{o}{@} \PYG{n}{X\PYGZus{}poly\PYGZus{}mat} \PYG{o}{+} \PYG{n}{alpha\PYGZus{}ridge}\PYG{o}{*}\PYG{n}{I\PYGZus{}p}
\PYG{n}{XTX\PYGZus{}ridge\PYGZus{}inv} \PYG{o}{=} \PYG{n}{np}\PYG{o}{.}\PYG{n}{linalg}\PYG{o}{.}\PYG{n}{inv}\PYG{p}{(}\PYG{n}{XTX\PYGZus{}ridge}\PYG{p}{)}
\PYG{n}{XTy\PYGZus{}ridge} \PYG{o}{=} \PYG{n}{X\PYGZus{}poly\PYGZus{}mat}\PYG{o}{.}\PYG{n}{T} \PYG{o}{@} \PYG{n}{y\PYGZus{}vec}
\PYG{n}{beta\PYGZus{}hat\PYGZus{}ridge\PYGZus{}manual} \PYG{o}{=} \PYG{n}{XTX\PYGZus{}ridge\PYGZus{}inv} \PYG{o}{@} \PYG{n}{XTy\PYGZus{}ridge}
\PYG{n}{beta\PYGZus{}hat\PYGZus{}ridge\PYGZus{}manual} \PYG{o}{=} \PYG{n}{beta\PYGZus{}hat\PYGZus{}ridge\PYGZus{}manual}\PYG{o}{.}\PYG{n}{flatten}\PYG{p}{(}\PYG{p}{)}

\PYG{n+nb}{print}\PYG{p}{(}\PYG{l+s+s2}{\PYGZdq{}}\PYG{l+s+s2}{Manual Ridge (polynomial) shape:}\PYG{l+s+s2}{\PYGZdq{}}\PYG{p}{,} \PYG{n}{beta\PYGZus{}hat\PYGZus{}ridge\PYGZus{}manual}\PYG{o}{.}\PYG{n}{shape}\PYG{p}{)}
\PYG{n}{pprint}\PYG{p}{(}\PYG{l+s+s2}{\PYGZdq{}}\PYG{l+s+se}{\PYGZbs{}\PYGZbs{}}\PYG{l+s+s2}{text}\PYG{l+s+s2}{\PYGZob{}}\PYG{l+s+s2}{Manual Ridge Coeffs\PYGZcb{}:}\PYG{l+s+s2}{\PYGZdq{}}\PYG{p}{,} \PYG{n}{beta\PYGZus{}hat\PYGZus{}ridge\PYGZus{}manual}\PYG{p}{)}
\end{sphinxVerbatim}

\end{sphinxuseclass}\end{sphinxVerbatimInput}
\begin{sphinxVerbatimOutput}

\begin{sphinxuseclass}{cell_output}
\begin{sphinxVerbatim}[commandchars=\\\{\}]
Manual Ridge (polynomial) shape: (10,)
\end{sphinxVerbatim}
\begin{equation*}
\begin{split}\displaystyle \text{Manual Ridge Coeffs}:\left[
\begin{array}{}
  5.0547 &  1.4667 & -1.9617 &  0.5238 &  0.0037 &  0.0051 &  0.0242 &  0.0037 &  0.0133 & -0.0038
\end{array}
\right]\end{split}
\end{equation*}
\end{sphinxuseclass}\end{sphinxVerbatimOutput}

\end{sphinxuseclass}
\begin{sphinxuseclass}{cell}\begin{sphinxVerbatimInput}

\begin{sphinxuseclass}{cell_input}
\begin{sphinxVerbatim}[commandchars=\\\{\}]
\PYG{k+kn}{from}\PYG{+w}{ }\PYG{n+nn}{sklearn}\PYG{n+nn}{.}\PYG{n+nn}{linear\PYGZus{}model}\PYG{+w}{ }\PYG{k+kn}{import} \PYG{n}{Ridge}

\PYG{n}{ridge\PYGZus{}sklearn} \PYG{o}{=} \PYG{n}{Ridge}\PYG{p}{(}\PYG{n}{alpha}\PYG{o}{=}\PYG{n}{alpha\PYGZus{}ridge}\PYG{p}{,} \PYG{n}{fit\PYGZus{}intercept}\PYG{o}{=}\PYG{k+kc}{False}\PYG{p}{)} 
\PYG{c+c1}{\PYGZsh{} \PYGZdq{}fit\PYGZus{}intercept=False\PYGZdq{} because we already included the intercept column in X\PYGZus{}poly\PYGZus{}mat}
\PYG{c+c1}{\PYGZsh{} We will manually set the intercept to match the first column or we can remove intercept from X\PYGZus{}poly\PYGZus{}mat.}

\PYG{n}{ridge\PYGZus{}sklearn}\PYG{o}{.}\PYG{n}{fit}\PYG{p}{(}\PYG{n}{X\PYGZus{}poly\PYGZus{}mat}\PYG{p}{,} \PYG{n}{df}\PYG{p}{[}\PYG{l+s+s1}{\PYGZsq{}}\PYG{l+s+s1}{Y}\PYG{l+s+s1}{\PYGZsq{}}\PYG{p}{]}\PYG{p}{)}
\PYG{n}{beta\PYGZus{}hat\PYGZus{}ridge\PYGZus{}sklearn} \PYG{o}{=} \PYG{n}{ridge\PYGZus{}sklearn}\PYG{o}{.}\PYG{n}{coef\PYGZus{}}

\PYG{n}{pprint}\PYG{p}{(}\PYG{l+s+s2}{\PYGZdq{}}\PYG{l+s+se}{\PYGZbs{}\PYGZbs{}}\PYG{l+s+s2}{text}\PYG{l+s+s2}{\PYGZob{}}\PYG{l+s+s2}{scikit\PYGZhy{}learn Ridge Coeffs\PYGZcb{}:}\PYG{l+s+s2}{\PYGZdq{}}\PYG{p}{,} \PYG{n}{beta\PYGZus{}hat\PYGZus{}ridge\PYGZus{}sklearn}\PYG{p}{)}
\PYG{n}{pprint}\PYG{p}{(}\PYG{l+s+s2}{\PYGZdq{}}\PYG{l+s+se}{\PYGZbs{}\PYGZbs{}}\PYG{l+s+s2}{text}\PYG{l+s+s2}{\PYGZob{}}\PYG{l+s+s2}{Difference (manual \PYGZhy{} sklearn)\PYGZcb{}:}\PYG{l+s+s2}{\PYGZdq{}}\PYG{p}{,} \PYG{p}{(}\PYG{n}{beta\PYGZus{}hat\PYGZus{}ridge\PYGZus{}manual} \PYG{o}{\PYGZhy{}} \PYG{n}{beta\PYGZus{}hat\PYGZus{}ridge\PYGZus{}sklearn}\PYG{p}{)}\PYG{p}{)}
\end{sphinxVerbatim}

\end{sphinxuseclass}\end{sphinxVerbatimInput}
\begin{sphinxVerbatimOutput}

\begin{sphinxuseclass}{cell_output}\begin{equation*}
\begin{split}\displaystyle \text{scikit-learn Ridge Coeffs}:\left[
\begin{array}{}
  4.1657 &  1.4733 & -1.9611 &  0.5235 &  0.0387 &  0.0073 &  0.0330 &  0.0122 &  0.0117 &  0.0322
\end{array}
\right]\end{split}
\end{equation*}\begin{equation*}
\begin{split}\displaystyle \text{Difference (manual - sklearn)}:\left[
\begin{array}{}
  0.8890 & -0.0067 & -0.0006 &  0.0003 & -0.0350 & -0.0022 & -0.0088 & -0.0085 &  0.0016 & -0.0360
\end{array}
\right]\end{split}
\end{equation*}
\end{sphinxuseclass}\end{sphinxVerbatimOutput}

\end{sphinxuseclass}
\sphinxAtStartPar
They should be nearly identical (tiny floating discrepancies may occur).


\subsection{Residual analysis for Ridge}
\label{\detokenize{02:residual-analysis-for-ridge}}
\begin{sphinxuseclass}{cell}\begin{sphinxVerbatimInput}

\begin{sphinxuseclass}{cell_input}
\begin{sphinxVerbatim}[commandchars=\\\{\}]
\PYG{n}{y\PYGZus{}pred\PYGZus{}ridge} \PYG{o}{=} \PYG{n}{X\PYGZus{}poly\PYGZus{}mat} \PYG{o}{@} \PYG{n}{beta\PYGZus{}hat\PYGZus{}ridge\PYGZus{}manual}
\PYG{n}{resid\PYGZus{}ridge} \PYG{o}{=} \PYG{n}{df}\PYG{p}{[}\PYG{l+s+s1}{\PYGZsq{}}\PYG{l+s+s1}{Y}\PYG{l+s+s1}{\PYGZsq{}}\PYG{p}{]} \PYG{o}{\PYGZhy{}} \PYG{n}{y\PYGZus{}pred\PYGZus{}ridge}

\PYG{n}{plt}\PYG{o}{.}\PYG{n}{figure}\PYG{p}{(}\PYG{n}{figsize}\PYG{o}{=}\PYG{p}{(}\PYG{l+m+mi}{10}\PYG{p}{,} \PYG{l+m+mi}{6}\PYG{p}{)}\PYG{p}{)}
\PYG{n}{plt}\PYG{o}{.}\PYG{n}{scatter}\PYG{p}{(}\PYG{n}{y\PYGZus{}pred\PYGZus{}ridge}\PYG{p}{,} \PYG{n}{resid\PYGZus{}ridge}\PYG{p}{,} \PYG{n}{alpha}\PYG{o}{=}\PYG{l+m+mf}{0.5}\PYG{p}{)}
\PYG{n}{plt}\PYG{o}{.}\PYG{n}{axhline}\PYG{p}{(}\PYG{n}{y}\PYG{o}{=}\PYG{l+m+mi}{0}\PYG{p}{,} \PYG{n}{color}\PYG{o}{=}\PYG{l+s+s1}{\PYGZsq{}}\PYG{l+s+s1}{g}\PYG{l+s+s1}{\PYGZsq{}}\PYG{p}{,} \PYG{n}{linestyle}\PYG{o}{=}\PYG{l+s+s1}{\PYGZsq{}}\PYG{l+s+s1}{\PYGZhy{}\PYGZhy{}}\PYG{l+s+s1}{\PYGZsq{}}\PYG{p}{)}
\PYG{n}{plt}\PYG{o}{.}\PYG{n}{xlabel}\PYG{p}{(}\PYG{l+s+s2}{\PYGZdq{}}\PYG{l+s+s2}{Fitted values (Ridge, polynomial)}\PYG{l+s+s2}{\PYGZdq{}}\PYG{p}{)}
\PYG{n}{plt}\PYG{o}{.}\PYG{n}{ylabel}\PYG{p}{(}\PYG{l+s+s2}{\PYGZdq{}}\PYG{l+s+s2}{Residuals}\PYG{l+s+s2}{\PYGZdq{}}\PYG{p}{)}
\PYG{n}{plt}\PYG{o}{.}\PYG{n}{title}\PYG{p}{(}\PYG{l+s+s2}{\PYGZdq{}}\PYG{l+s+s2}{Residuals vs. Fitted (Ridge)}\PYG{l+s+s2}{\PYGZdq{}}\PYG{p}{)}
\PYG{n}{plt}\PYG{o}{.}\PYG{n}{grid}\PYG{p}{(}\PYG{l+s+s2}{\PYGZdq{}}\PYG{l+s+s2}{on}\PYG{l+s+s2}{\PYGZdq{}}\PYG{p}{)}
\PYG{n}{plt}\PYG{o}{.}\PYG{n}{show}\PYG{p}{(}\PYG{p}{)}
\end{sphinxVerbatim}

\end{sphinxuseclass}\end{sphinxVerbatimInput}
\begin{sphinxVerbatimOutput}

\begin{sphinxuseclass}{cell_output}
\noindent\sphinxincludegraphics{{46cd873ca0b7cf4c0e09251a497c889a8edddc5011ec75ab3d6d35306dbee6b6}.png}

\end{sphinxuseclass}\end{sphinxVerbatimOutput}

\end{sphinxuseclass}
\sphinxAtStartPar
\sphinxstylestrong{Remarks}:
\begin{itemize}
\item {} 
\sphinxAtStartPar
As alpha\_ridge increases, the polynomial coefficients shrink more, typically reducing variance (but possibly increasing bias).

\item {} 
\sphinxAtStartPar
We can further use cross\sphinxhyphen{}validation to select alpha automatically or do a hyperparameter grid.

\end{itemize}


\section{Conclusion}
\label{\detokenize{02:conclusion}}
\sphinxAtStartPar
This chapter has provided a comprehensive, hands\sphinxhyphen{}on exploration of linear models—starting from the manual derivation of the Ordinary Least Squares (OLS) estimator to the application of regularization techniques and feature transformations. Through synthetic data, we validated the accuracy of our estimators, examined residuals for assumption\sphinxhyphen{}checking, and used diagnostic plots to evaluate model fit.

\sphinxAtStartPar
We have seen how linear regression, while conceptually straightforward, remains a powerful modeling tool when used thoughtfully. The incorporation of polynomial and radial basis function expansions extends the expressive capacity of linear models, enabling them to approximate complex, non\sphinxhyphen{}linear relationships without abandoning interpretability.

\sphinxAtStartPar
We also introduced Ridge regression, which provides a principled approach to combat overfitting and multicollinearity, especially in high\sphinxhyphen{}dimensional settings. This regularized extension highlights a key transition in statistical learning: balancing model flexibility with generalization performance.

\sphinxAtStartPar
Ultimately, this chapter emphasizes not only how to implement linear models, but how to critically assess their assumptions and limitations. These foundational skills will be essential as we move on to more sophisticated models in the chapters ahead.

\sphinxstepscope


\chapter{Data Preprocessing, Local Models and Tree\sphinxhyphen{}Based Models}
\label{\detokenize{03:data-preprocessing-local-models-and-tree-based-models}}\label{\detokenize{03::doc}}

\section{Introduction}
\label{\detokenize{03:introduction}}
\sphinxAtStartPar
Machine Learning, as you now know, is based on statistical analysis of data. Data comes in many kinds of formats, resolutions, annotations, … and therefore should be processed before performing analysis. This is the first content of this practical. We will focus on regression tasks, as already defined in the previous practicals, showing to commonly used strategies. The practical is structured around five key objectives:
\begin{enumerate}
\sphinxsetlistlabels{\arabic}{enumi}{enumii}{}{.}%
\item {} 
\sphinxAtStartPar
Understanding and applying essential data preprocessing techniques\\
Before training a machine learning model, the data must be properly prepared. We will cover the following common preprocessing steps:
\begin{itemize}
\item {} 
\sphinxAtStartPar
Handling missing values: strategies to impute or remove data with missing entries.

\item {} 
\sphinxAtStartPar
Feature scaling and normalization: ensuring features are on comparable scales to improve model performance.

\item {} 
\sphinxAtStartPar
Encoding categorical variables: transforming non\sphinxhyphen{}numeric data (e.g., categories) into numerical representations suitable for machine learning algorithms.

\end{itemize}

\item {} 
\sphinxAtStartPar
Implementing the K\sphinxhyphen{}Nearest Neighbors (KNN) algorithm from scratch\\
We will write a simple, “vanilla” version of the KNN algorithm for regression tasks. This implementation will help demystify how KNN works by reinforcing the idea of prediction based on similarity in feature space.

\item {} 
\sphinxAtStartPar
Training and evaluating tree\sphinxhyphen{}based regression models\\
We will work with two popular types of decision\sphinxhyphen{}tree\sphinxhyphen{}based models for regression:
\begin{itemize}
\item {} 
\sphinxAtStartPar
Regression trees: models that recursively split the data space into regions to predict a continuous output.

\item {} 
\sphinxAtStartPar
Random Forests: an ensemble method that combines multiple decision trees to improve predictive accuracy and reduce overfitting.

\end{itemize}

\item {} 
\sphinxAtStartPar
Interpreting feature importances in tree\sphinxhyphen{}based models\\
One advantage of tree\sphinxhyphen{}based models is their ability to estimate the relative importance of each input feature. We will learn how to extract and interpret these feature importances to gain insights into which variables are most influential in the prediction.

\item {} 
\sphinxAtStartPar
Using \sphinxcode{\sphinxupquote{scikit\sphinxhyphen{}learn}} Pipelines to streamline preprocessing and modeling\\
To ensure clean and reproducible workflows, we will introduce \sphinxcode{\sphinxupquote{scikit\sphinxhyphen{}learn}} pipelines, which allow us to:
\begin{itemize}
\item {} 
\sphinxAtStartPar
Chain together preprocessing steps and model training into a single object.

\item {} 
\sphinxAtStartPar
Simplify cross\sphinxhyphen{}validation and hyperparameter tuning.

\item {} 
\sphinxAtStartPar
Prevent data leakage by ensuring that preprocessing is properly applied to both training and test data.

\end{itemize}

\end{enumerate}

\sphinxAtStartPar
By the end of this practical, you should be able to preprocess real\sphinxhyphen{}world data appropriately, implement a basic regression model, use tree\sphinxhyphen{}based models effectively, and combine preprocessing and modeling steps into a coherent, automated pipeline using \sphinxcode{\sphinxupquote{scikit\sphinxhyphen{}learn}}.


\subsection{Important precision about the evaluation metric used in this practical}
\label{\detokenize{03:important-precision-about-the-evaluation-metric-used-in-this-practical}}
\sphinxAtStartPar
In this practical, we will use the \sphinxstyleemphasis{Mean Squared Error} (MSE) as a convenient shorthand for the estimated \sphinxstyleemphasis{Mean Integrated Squared Error} \((\widehat{\mathrm{MISE}})\). Indeed, in practice, we rarely know the true MISE and instead rely on estimates obtained via techniques like \(K\)\sphinxhyphen{}fold cross\sphinxhyphen{}validation or leave\sphinxhyphen{}one\sphinxhyphen{}out (LOO) cross\sphinxhyphen{}validation. So keep in mind that when we refer to the MSE in these contexts, what we are really computing is an estimation of \(\widehat{\mathrm{MISE}}\) from our data.


\subsection{Data preprocessing overview}
\label{\detokenize{03:data-preprocessing-overview}}
\sphinxAtStartPar
Data preprocessing is the foundation of any successful machine learning project. Many Machine Learning models see their performances degraded when trained on poorly preprocessed data. Key steps include:
\begin{itemize}
\item {} 
\sphinxAtStartPar
Handling Missing Values: Missing data can bias results if not addressed. This problem affects many models—such as linear models, support vector machines (SVMs), and neural networks—but it is particularly critical for distance\sphinxhyphen{}based algorithms like KNN where missing feature values can disrupt the calculation of distances, leading to inaccurate outcomes. Common strategies for handling missing values include:
\begin{itemize}
\item {} 
\sphinxAtStartPar
\sphinxstyleemphasis{Dropping} rows/columns with missing data.

\item {} 
\sphinxAtStartPar
\sphinxstyleemphasis{Imputing} missing data using simple statistics (mean, median) or more sophisticated methods like \sphinxcode{\sphinxupquote{KNNImputer}}.

\end{itemize}

\item {} 
\sphinxAtStartPar
Normalization / Scaling: Many algorithms assume features are on similar scales. Without proper scaling, features with larger numerical ranges may dominate those with smaller ranges, potentially biasing the model’s results. Common scaling strategies include:
\begin{itemize}
\item {} 
\sphinxAtStartPar
\sphinxstyleemphasis{Standardization}: Transform features to have zero mean and unit variance.

\item {} 
\sphinxAtStartPar
\sphinxstyleemphasis{Min\sphinxhyphen{}Max Scaling}: Scale features to a fixed \([0,1]\) range.

\end{itemize}

\item {} 
\sphinxAtStartPar
Dealing with Categorical Variables: Categorical variables must be converted to a numeric representation that preserves meaningful differences. Common strategies for dealing with categorical variables include:
\begin{itemize}
\item {} 
\sphinxAtStartPar
\sphinxstyleemphasis{Label Encoding}: Assign an integer to each category (useful if there is an ordinal relationship).

\item {} 
\sphinxAtStartPar
\sphinxstyleemphasis{One\sphinxhyphen{}Hot Encoding (Dummy Encoding)}: Create a new binary column for each category (common for nominal features). This approach can cause a significant increase in input dimensionality (see the curse of dimensionality).

\end{itemize}

\end{itemize}


\section{K\sphinxhyphen{}Nearest Neighbors}
\label{\detokenize{03:k-nearest-neighbors}}
\begin{sphinxShadowBox}
\sphinxstylesidebartitle{}

\sphinxAtStartPar
💡Did you know ? The \(K\)\sphinxhyphen{}Neighrest Neighbors algorithm was first proposed in Fix and Hodges {[}\hyperlink{cite.bibliography:id2}{FH51}{]}.
\end{sphinxShadowBox}

\sphinxAtStartPar
The \sphinxstyleemphasis{\(K\)\sphinxhyphen{}Nearest Neighbors} (KNN) algorithm is a local modeling approach that infers its prediction from the data points closest (in some metric) to a query. Suppose we have a training set of inputs \(\{x_i\}\) with associated labels or targets \(\{y_i\}\). For a new query point \(q\):
\begin{enumerate}
\sphinxsetlistlabels{\arabic}{enumi}{enumii}{}{.}%
\item {} 
\sphinxAtStartPar
Distance computation: Compute the distance (e.g., Euclidean) between \(q\) and each training example \(x_i\).

\item {} 
\sphinxAtStartPar
Ranking: Sort or rank the training points by increasing distance to \(q\).

\item {} 
\sphinxAtStartPar
Neighborhood selection: Identify the subset of the \(K\) nearest neighbors \(\{x_{[1]}, \dots, x_{[K]}\}\) where \(y_{[i]}\) denotes the label of neighbor \(x_{[i]}\).

\item {} 
\sphinxAtStartPar
Regression or classification: KNN can be viewed as a local estimator of the conditional expectation \(\mathbb{E}[y|x=q]\).

\end{enumerate}

\sphinxAtStartPar
For regression tasks, the simplest approach is to average the targets of the \(K\) nearest neighbors:
\begin{equation*}
\begin{split}
\hat{y}(q) = \frac{1}{K}\sum_{i=1}^{K} y_{[i]}.
\end{split}
\end{equation*}
\sphinxAtStartPar
Alternatively, a local linear model may be used:
\begin{equation*}
\begin{split}
\hat{y}(q) = \hat{a}^T q + \hat{b},
\end{split}
\end{equation*}
\sphinxAtStartPar
where \(\hat{a}\) and \(\hat{b}\) are estimated (locally) by least\sphinxhyphen{}squares on the neighbors.

\sphinxAtStartPar
For classification, one can estimate the conditional probability of belonging to a specific class (e.g., class “1”) by the proportion of neighbors labeled “1”:
\begin{equation*}
\begin{split}
\hat{p}_1(q) = \frac{1}{K}\sum_{i=1}^{K} y_{[i]}.
\end{split}
\end{equation*}
\sphinxAtStartPar
A threshold (e.g., 0.5) can then be used to decide the class label.

\sphinxAtStartPar
The key hyperparameter here is \(K\). A \sphinxstyleemphasis{smaller} \(K\) reduces bias but increases variance, while a \sphinxstyleemphasis{larger} \(K\) produces a smoother model at the risk of higher bias. KNN’s performance also heavily depends on the distance metric, especially when features have different scales or include categorical variables (which may require specialized encodings or distance definitions).


\section{Regression Trees}
\label{\detokenize{03:regression-trees}}
\sphinxAtStartPar
Regression trees rely on a tree\sphinxhyphen{}based structure of \sphinxstyleemphasis{internal nodes} (where decisions are made) and \sphinxstyleemphasis{terminal nodes} (leaves), which partition the input space into mutually exclusive regions, each associated with a simple (often constant) local model. Its construction begins with a tree growing phase: starting from the root node that contains all data, we recursively choose a split that best reduces the empirical error. Specifically, for a node \(t\) containing \(N(t)\) samples, we define
\begin{equation*}
\begin{split}
R_{\mathrm{emp}}(t) \;=\; \min_{\alpha_t} \sum_{i=1}^{N(t)} L\bigl(y_i,\;h_t(x_i,\;\alpha_t)\bigr),
\end{split}
\end{equation*}
\sphinxAtStartPar
where \(L\) is typically the squared error \((y - \hat{y})^2\) for regression. Given a potential split \(s\) dividing \(t\) into children \(t_l\) and \(t_r\), we consider the decrease in empirical error:
\begin{equation*}
\begin{split}
\Delta E(s,\;t) \;=\; R_{\mathrm{emp}}(t)\;-\;\bigl(R_{\mathrm{emp}}(t_l)\;+\;R_{\mathrm{emp}}(t_r)\bigr).
\end{split}
\end{equation*}
\sphinxAtStartPar
We choose the split \(s^*\) maximizing \(\Delta E\), partition the dataset accordingly, and repeat the procedure recursively until no further improvement is found or a stopping criterion is met. This exhaustive splitting often yields a very large tree that overfits the data.

\sphinxAtStartPar
To address overfitting, a cost\sphinxhyphen{}complexity pruning procedure is commonly used. Introducing a complexity parameter \(\lambda \ge 0\) that penalizes the number of leaf nodes, we define the cost\sphinxhyphen{}complexity measure:
\begin{equation*}
\begin{split}
R_{\lambda}(T) = R_{\mathrm{emp}}(T) + \lambda \, |T|,
\end{split}
\end{equation*}
\sphinxAtStartPar
where \(|T|\) is the number of terminal nodes (leaves) in tree \(T\).
\begin{itemize}
\item {} 
\sphinxAtStartPar
By gradually increasing \(\lambda\), we obtain a \sphinxstyleemphasis{sequence} of subtrees:

\end{itemize}
\begin{equation*}
\begin{split}
T_{\max} \supset T_{L-1} \supset \dots \supset T_2 \supset T_1.
\end{split}
\end{equation*}
\sphinxAtStartPar
Each has fewer leaves. We consider all admissible subtrees \(T_t \subset T\) of the large tree and compute the smallest
\begin{equation*}
\begin{split}
\lambda_t = \frac{R_{\mathrm{emp}}(T) - R_{\mathrm{emp}}(T_t)}{|T| - |T_t|}
\end{split}
\end{equation*}
\sphinxAtStartPar
that yields a lower cost. We select the subtree that minimizes \(R_{\lambda}(T)\) (the best balances between empirical error and model complexity).
The final tree structure is typically chosen via cross\sphinxhyphen{}validation or held\sphinxhyphen{}out validation to find the best \(\lambda\).
A more visual representation of tree\sphinxhyphen{}based models, an animated version of decision trees (classification) can be found \sphinxhref{http://www.r2d3.us/visual-intro-to-machine-learning-part-1/}{here}.

\begin{sphinxShadowBox}
\sphinxstylesidebartitle{}

\sphinxAtStartPar
Random Forest have been first presented in Ho {[}\hyperlink{cite.bibliography:id3}{Ho98}{]}.
\end{sphinxShadowBox}


\section{Random Forests}
\label{\detokenize{03:random-forests}}
\sphinxAtStartPar
A Random Forest (RF), is an ensemble learning technique designed to reduce the variance of decision trees by combining two main ideas:
\begin{enumerate}
\sphinxsetlistlabels{\arabic}{enumi}{enumii}{}{.}%
\item {} 
\sphinxAtStartPar
Bootstrap Sampling (Bagging): We generate \(B\) \sphinxstyleemphasis{bootstrap} datasets, each by sampling (with replacement) from the original training data.

\item {} 
\sphinxAtStartPar
Random Feature Selection (Feature Bagging): At each split, a \sphinxstyleemphasis{random subset} of \(n' < n\) features is chosen, and the best split is found \sphinxstyleemphasis{only among those \(n'\) features}.

\end{enumerate}

\sphinxAtStartPar
Hence, each tree \(h_b(\mathbf{x}, \alpha_b)\) is built from a \sphinxstyleemphasis{different} resampled dataset and uses only a random subset of features at each split. As a result, the trees are more decorrelated and often are \sphinxstyleemphasis{not pruned} heavily.

\sphinxAtStartPar
Formally, for a regression task:
\begin{equation*}
\begin{split}
  h_{\mathrm{rf}}(\mathbf{x}) = \frac{1}{B} \sum_{b=1}^{B} h_b(\mathbf{x}, \alpha_b),
\end{split}
\end{equation*}
\sphinxAtStartPar
i.e., the average of all tree predictions. (In classification, we take the majority vote.)


\subsection{Why does this reduce variance?}
\label{\detokenize{03:why-does-this-reduce-variance}}
\sphinxAtStartPar
Suppose each tree \(h_b\) has variance \(\sigma^2\) and they have pairwise correlation \(\rho\). The variance of the RF predictor \(h_{\mathrm{rf}}\) can be approximated as:
\begin{equation*}
\begin{split}
  \mathrm{Var}[h_{\mathrm{rf}}] \approx \frac{(1 - \rho)\,\sigma^2}{B} + \rho\,\sigma^2,
\end{split}
\end{equation*}
\sphinxAtStartPar
which shows that increasing \(B\) (the number of trees) reduces the first term, while lowering the correlation \(\rho\) among trees reduces overall variance. Random feature selection helps reduce \(\rho\), because each tree sees a different subset of features.


\subsection{Main Hyperparameters}
\label{\detokenize{03:main-hyperparameters}}\begin{enumerate}
\sphinxsetlistlabels{\arabic}{enumi}{enumii}{}{.}%
\item {} 
\sphinxAtStartPar
\(B\): the number of trees in the forest.

\item {} 
\sphinxAtStartPar
\(n'\): the number of randomly selected features at each split (often \(\sqrt{n}\) by default in classification).

\item {} 
\sphinxAtStartPar
Tree parameters: e.g., maximum depth, minimum samples per leaf, etc.

\end{enumerate}

\sphinxAtStartPar
In practice, random forests typically provide strong performance out of the box and are less sensitive to hyperparameter tuning compared to single trees or more complex models. They also allow computing a useful estimate of feature importance by measuring, for each variable, how much it contributes to the cost function improvement across all splits in all trees.


\section{Feature Importances}
\label{\detokenize{03:feature-importances}}
\sphinxAtStartPar
Most \sphinxcode{\sphinxupquote{scikit\sphinxhyphen{}learn}} models allows you to estimate how each feature contribues to the final prediction. For example, for tree\sphinxhyphen{}based models, both \sphinxcode{\sphinxupquote{DecisionTreeRegressor}} and \sphinxcode{\sphinxupquote{RandomForestRegressor}} in \sphinxcode{\sphinxupquote{scikit\sphinxhyphen{}learn}} store a \sphinxcode{\sphinxupquote{feature\_importances\_}} attribute after fitting. This gives an estimate of how much each feature contributed to reducing the split criterion (e.g., MSE).


\section{\sphinxstyleliteralintitle{\sphinxupquote{Scikit\sphinxhyphen{}learn}} Pipelines}
\label{\detokenize{03:scikit-learn-pipelines}}
\sphinxAtStartPar
Pipelines in \sphinxcode{\sphinxupquote{scikit\sphinxhyphen{}learn}} let us combine data preprocessing (e.g., imputation, scaling) and model training (e.g., a random forest) into a single workflow object. Below is a simple example:

\begin{sphinxuseclass}{cell}\begin{sphinxVerbatimInput}

\begin{sphinxuseclass}{cell_input}
\begin{sphinxVerbatim}[commandchars=\\\{\}]
\PYG{k+kn}{from}\PYG{+w}{ }\PYG{n+nn}{sklearn}\PYG{n+nn}{.}\PYG{n+nn}{pipeline}\PYG{+w}{ }\PYG{k+kn}{import} \PYG{n}{Pipeline}
\PYG{k+kn}{from}\PYG{+w}{ }\PYG{n+nn}{sklearn}\PYG{n+nn}{.}\PYG{n+nn}{preprocessing}\PYG{+w}{ }\PYG{k+kn}{import} \PYG{n}{StandardScaler}
\PYG{k+kn}{from}\PYG{+w}{ }\PYG{n+nn}{sklearn}\PYG{n+nn}{.}\PYG{n+nn}{linear\PYGZus{}model}\PYG{+w}{ }\PYG{k+kn}{import} \PYG{n}{LinearRegression}
\PYG{k+kn}{from}\PYG{+w}{ }\PYG{n+nn}{sklearn}\PYG{n+nn}{.}\PYG{n+nn}{impute}\PYG{+w}{ }\PYG{k+kn}{import} \PYG{n}{SimpleImputer}

\PYG{n}{example\PYGZus{}pipeline} \PYG{o}{=} \PYG{n}{Pipeline}\PYG{p}{(}\PYG{p}{[}
    \PYG{p}{(}\PYG{l+s+s2}{\PYGZdq{}}\PYG{l+s+s2}{imputer}\PYG{l+s+s2}{\PYGZdq{}}\PYG{p}{,} \PYG{n}{SimpleImputer}\PYG{p}{(}\PYG{n}{strategy}\PYG{o}{=}\PYG{l+s+s2}{\PYGZdq{}}\PYG{l+s+s2}{mean}\PYG{l+s+s2}{\PYGZdq{}}\PYG{p}{)}\PYG{p}{)}\PYG{p}{,}
    \PYG{p}{(}\PYG{l+s+s2}{\PYGZdq{}}\PYG{l+s+s2}{scaler}\PYG{l+s+s2}{\PYGZdq{}}\PYG{p}{,} \PYG{n}{StandardScaler}\PYG{p}{(}\PYG{p}{)}\PYG{p}{)}\PYG{p}{,}
    \PYG{p}{(}\PYG{l+s+s2}{\PYGZdq{}}\PYG{l+s+s2}{regressor}\PYG{l+s+s2}{\PYGZdq{}}\PYG{p}{,} \PYG{n}{LinearRegression}\PYG{p}{(}\PYG{p}{)}\PYG{p}{)}
\PYG{p}{]}\PYG{p}{)}
\end{sphinxVerbatim}

\end{sphinxuseclass}\end{sphinxVerbatimInput}

\end{sphinxuseclass}
\begin{sphinxuseclass}{cell}\begin{sphinxVerbatimInput}

\begin{sphinxuseclass}{cell_input}
\begin{sphinxVerbatim}[commandchars=\\\{\}]
\PYG{k+kn}{from}\PYG{+w}{ }\PYG{n+nn}{sklearn}\PYG{n+nn}{.}\PYG{n+nn}{base}\PYG{+w}{ }\PYG{k+kn}{import} \PYG{n}{BaseEstimator}\PYG{p}{,} \PYG{n}{TransformerMixin}
\PYG{k+kn}{from}\PYG{+w}{ }\PYG{n+nn}{sklearn}\PYG{n+nn}{.}\PYG{n+nn}{ensemble}\PYG{+w}{ }\PYG{k+kn}{import} \PYG{n}{RandomForestRegressor}

\PYG{k}{class}\PYG{+w}{ }\PYG{n+nc}{AddConstantTransformer}\PYG{p}{(}\PYG{n}{BaseEstimator}\PYG{p}{,} \PYG{n}{TransformerMixin}\PYG{p}{)}\PYG{p}{:}
\PYG{+w}{    }\PYG{l+s+sd}{\PYGZdq{}\PYGZdq{}\PYGZdq{}}
\PYG{l+s+sd}{    Custom transformer that adds a constant value to all features.}
\PYG{l+s+sd}{    \PYGZdq{}\PYGZdq{}\PYGZdq{}}
    \PYG{k}{def}\PYG{+w}{ }\PYG{n+nf+fm}{\PYGZus{}\PYGZus{}init\PYGZus{}\PYGZus{}}\PYG{p}{(}\PYG{n+nb+bp}{self}\PYG{p}{,} \PYG{n}{constant}\PYG{o}{=}\PYG{l+m+mf}{0.0}\PYG{p}{)}\PYG{p}{:}
        \PYG{n+nb+bp}{self}\PYG{o}{.}\PYG{n}{constant} \PYG{o}{=} \PYG{n}{constant}

    \PYG{k}{def}\PYG{+w}{ }\PYG{n+nf}{fit}\PYG{p}{(}\PYG{n+nb+bp}{self}\PYG{p}{,} \PYG{n}{X}\PYG{p}{,} \PYG{n}{y}\PYG{o}{=}\PYG{k+kc}{None}\PYG{p}{)}\PYG{p}{:}
        \PYG{k}{return} \PYG{n+nb+bp}{self}

    \PYG{k}{def}\PYG{+w}{ }\PYG{n+nf}{transform}\PYG{p}{(}\PYG{n+nb+bp}{self}\PYG{p}{,} \PYG{n}{X}\PYG{p}{,} \PYG{n}{y}\PYG{o}{=}\PYG{k+kc}{None}\PYG{p}{)}\PYG{p}{:}
        \PYG{k}{return} \PYG{n}{X} \PYG{o}{+} \PYG{n+nb+bp}{self}\PYG{o}{.}\PYG{n}{constant}

\PYG{n}{custom\PYGZus{}pipeline} \PYG{o}{=} \PYG{n}{Pipeline}\PYG{p}{(}\PYG{p}{[}
    \PYG{p}{(}\PYG{l+s+s2}{\PYGZdq{}}\PYG{l+s+s2}{imputer}\PYG{l+s+s2}{\PYGZdq{}}\PYG{p}{,} \PYG{n}{SimpleImputer}\PYG{p}{(}\PYG{n}{strategy}\PYG{o}{=}\PYG{l+s+s2}{\PYGZdq{}}\PYG{l+s+s2}{mean}\PYG{l+s+s2}{\PYGZdq{}}\PYG{p}{)}\PYG{p}{)}\PYG{p}{,}
    \PYG{p}{(}\PYG{l+s+s2}{\PYGZdq{}}\PYG{l+s+s2}{custom\PYGZus{}add\PYGZus{}constant}\PYG{l+s+s2}{\PYGZdq{}}\PYG{p}{,} \PYG{n}{AddConstantTransformer}\PYG{p}{(}\PYG{n}{constant}\PYG{o}{=}\PYG{l+m+mf}{1.0}\PYG{p}{)}\PYG{p}{)}\PYG{p}{,}
    \PYG{p}{(}\PYG{l+s+s2}{\PYGZdq{}}\PYG{l+s+s2}{scaler}\PYG{l+s+s2}{\PYGZdq{}}\PYG{p}{,} \PYG{n}{StandardScaler}\PYG{p}{(}\PYG{p}{)}\PYG{p}{)}\PYG{p}{,}
    \PYG{p}{(}\PYG{l+s+s2}{\PYGZdq{}}\PYG{l+s+s2}{regressor}\PYG{l+s+s2}{\PYGZdq{}}\PYG{p}{,} \PYG{n}{RandomForestRegressor}\PYG{p}{(}\PYG{n}{n\PYGZus{}estimators}\PYG{o}{=}\PYG{l+m+mi}{10}\PYG{p}{,} \PYG{n}{random\PYGZus{}state}\PYG{o}{=}\PYG{l+m+mi}{42}\PYG{p}{)}\PYG{p}{)}
\PYG{p}{]}\PYG{p}{)}
\end{sphinxVerbatim}

\end{sphinxuseclass}\end{sphinxVerbatimInput}

\end{sphinxuseclass}

\section{Exercises}
\label{\detokenize{03:exercises}}

\subsection{California Housing}
\label{\detokenize{03:california-housing}}\begin{enumerate}
\sphinxsetlistlabels{\arabic}{enumi}{enumii}{}{.}%
\item {} 
\sphinxAtStartPar
Load the well\sphinxhyphen{}known \sphinxstyleemphasis{California Housing} regression dataset using the following function:

\end{enumerate}

\begin{sphinxVerbatim}[commandchars=\\\{\}]
\PYG{k+kn}{from}\PYG{+w}{ }\PYG{n+nn}{sklearn}\PYG{n+nn}{.}\PYG{n+nn}{datasets}\PYG{+w}{ }\PYG{k+kn}{import} \PYG{n}{fetch\PYGZus{}california\PYGZus{}housing}
\end{sphinxVerbatim}
\begin{enumerate}
\sphinxsetlistlabels{\arabic}{enumi}{enumii}{}{.}%
\setcounter{enumi}{1}
\item {} 
\sphinxAtStartPar
Explore the dataset (distribution, summary statistics).

\item {} 
\sphinxAtStartPar
Check for missing values. If the dataset does not contain missing values, artificially introduce some for the purpose of this exercice. Then, handle missing values by:
\begin{itemize}
\item {} 
\sphinxAtStartPar
Dropping them (simple approach).

\item {} 
\sphinxAtStartPar
Imputing them (mean or median).

\item {} 
\sphinxAtStartPar
Using a more sophisticated method (e.g., \sphinxcode{\sphinxupquote{KNNImputer}}).

\end{itemize}

\item {} 
\sphinxAtStartPar
Normalize the features (using e.g., \sphinxcode{\sphinxupquote{StandardScaler}}). Remember that, as all processing steps, scaling is fit on the training set (not on the entire dataset) to avoid biasing the test set!

\item {} 
\sphinxAtStartPar
Preserve original vs. processed versions of the data for comparison.

\end{enumerate}


\subsection{Solution}
\label{\detokenize{03:solution}}

\subsubsection{Dataset and exploration}
\label{\detokenize{03:dataset-and-exploration}}
\sphinxAtStartPar
For this exercise, we will use the California Housing dataset from \sphinxcode{\sphinxupquote{sklearn.datasets.fetch\_california\_housing}}. It contains features related to California housing data, and the target is the average house value.

\sphinxAtStartPar
Alternatively, if you would like test your solution on other datasets, \sphinxcode{\sphinxupquote{scikit\sphinxhyphen{}learn}} provides other built\sphinxhyphen{}in datasets, such as:
\begin{itemize}
\item {} 
\sphinxAtStartPar
Diabetes (a small dataset with 10 features for a regression task on disease progression).

\item {} 
\sphinxAtStartPar
Linnerud (exercise/physiological data, a small multi\sphinxhyphen{}output regression problem).
If you want a dataset with even more features or complexity, you could consider Ames Housing or Kaggle House Price, both of which have many categorical and numeric features but may require additional steps to download or preprocess. You can also explore the UCI Machine Learning Repository for a wide range of tabular datasets.

\end{itemize}

\begin{sphinxuseclass}{cell}\begin{sphinxVerbatimInput}

\begin{sphinxuseclass}{cell_input}
\begin{sphinxVerbatim}[commandchars=\\\{\}]
\PYG{k+kn}{import}\PYG{+w}{ }\PYG{n+nn}{numpy}\PYG{+w}{ }\PYG{k}{as}\PYG{+w}{ }\PYG{n+nn}{np}
\PYG{k+kn}{import}\PYG{+w}{ }\PYG{n+nn}{pandas}\PYG{+w}{ }\PYG{k}{as}\PYG{+w}{ }\PYG{n+nn}{pd}
\PYG{k+kn}{from}\PYG{+w}{ }\PYG{n+nn}{sklearn}\PYG{n+nn}{.}\PYG{n+nn}{datasets}\PYG{+w}{ }\PYG{k+kn}{import} \PYG{n}{fetch\PYGZus{}california\PYGZus{}housing}

\PYG{n}{cal\PYGZus{}housing} \PYG{o}{=} \PYG{n}{fetch\PYGZus{}california\PYGZus{}housing}\PYG{p}{(}\PYG{p}{)}
\PYG{n}{X\PYGZus{}original} \PYG{o}{=} \PYG{n}{pd}\PYG{o}{.}\PYG{n}{DataFrame}\PYG{p}{(}\PYG{n}{cal\PYGZus{}housing}\PYG{o}{.}\PYG{n}{data}\PYG{p}{,} \PYG{n}{columns}\PYG{o}{=}\PYG{n}{cal\PYGZus{}housing}\PYG{o}{.}\PYG{n}{feature\PYGZus{}names}\PYG{p}{)}
\PYG{n}{y\PYGZus{}original} \PYG{o}{=} \PYG{n}{pd}\PYG{o}{.}\PYG{n}{Series}\PYG{p}{(}\PYG{n}{cal\PYGZus{}housing}\PYG{o}{.}\PYG{n}{target}\PYG{p}{,} \PYG{n}{name}\PYG{o}{=}\PYG{l+s+s1}{\PYGZsq{}}\PYG{l+s+s1}{Target}\PYG{l+s+s1}{\PYGZsq{}}\PYG{p}{)}

\PYG{n+nb}{print}\PYG{p}{(}\PYG{l+s+s2}{\PYGZdq{}}\PYG{l+s+s2}{Shape of the feature matrix:}\PYG{l+s+s2}{\PYGZdq{}}\PYG{p}{,} \PYG{n}{X\PYGZus{}original}\PYG{o}{.}\PYG{n}{shape}\PYG{p}{)}
\PYG{n+nb}{print}\PYG{p}{(}\PYG{l+s+s2}{\PYGZdq{}}\PYG{l+s+s2}{Shape of the target vector:}\PYG{l+s+s2}{\PYGZdq{}}\PYG{p}{,} \PYG{n}{y\PYGZus{}original}\PYG{o}{.}\PYG{n}{shape}\PYG{p}{)}

\PYG{n}{X\PYGZus{}original}\PYG{o}{.}\PYG{n}{info}\PYG{p}{(}\PYG{p}{)}
\PYG{n}{display}\PYG{p}{(}\PYG{n}{X\PYGZus{}original}\PYG{o}{.}\PYG{n}{describe}\PYG{p}{(}\PYG{p}{)}\PYG{p}{)}
\end{sphinxVerbatim}

\end{sphinxuseclass}\end{sphinxVerbatimInput}
\begin{sphinxVerbatimOutput}

\begin{sphinxuseclass}{cell_output}
\begin{sphinxVerbatim}[commandchars=\\\{\}]
Shape of the feature matrix: (20640, 8)
Shape of the target vector: (20640,)
\PYGZlt{}class \PYGZsq{}pandas.core.frame.DataFrame\PYGZsq{}\PYGZgt{}
RangeIndex: 20640 entries, 0 to 20639
Data columns (total 8 columns):
 \PYGZsh{}   Column      Non\PYGZhy{}Null Count  Dtype  
\PYGZhy{}\PYGZhy{}\PYGZhy{}  \PYGZhy{}\PYGZhy{}\PYGZhy{}\PYGZhy{}\PYGZhy{}\PYGZhy{}      \PYGZhy{}\PYGZhy{}\PYGZhy{}\PYGZhy{}\PYGZhy{}\PYGZhy{}\PYGZhy{}\PYGZhy{}\PYGZhy{}\PYGZhy{}\PYGZhy{}\PYGZhy{}\PYGZhy{}\PYGZhy{}  \PYGZhy{}\PYGZhy{}\PYGZhy{}\PYGZhy{}\PYGZhy{}  
 0   MedInc      20640 non\PYGZhy{}null  float64
 1   HouseAge    20640 non\PYGZhy{}null  float64
 2   AveRooms    20640 non\PYGZhy{}null  float64
 3   AveBedrms   20640 non\PYGZhy{}null  float64
 4   Population  20640 non\PYGZhy{}null  float64
 5   AveOccup    20640 non\PYGZhy{}null  float64
 6   Latitude    20640 non\PYGZhy{}null  float64
 7   Longitude   20640 non\PYGZhy{}null  float64
dtypes: float64(8)
memory usage: 1.3 MB
\end{sphinxVerbatim}

\begin{sphinxVerbatim}[commandchars=\\\{\}]
             MedInc      HouseAge      AveRooms     AveBedrms    Population  \PYGZbs{}
count  20640.000000  20640.000000  20640.000000  20640.000000  20640.000000   
mean       3.870671     28.639486      5.429000      1.096675   1425.476744   
std        1.899822     12.585558      2.474173      0.473911   1132.462122   
min        0.499900      1.000000      0.846154      0.333333      3.000000   
25\PYGZpc{}        2.563400     18.000000      4.440716      1.006079    787.000000   
50\PYGZpc{}        3.534800     29.000000      5.229129      1.048780   1166.000000   
75\PYGZpc{}        4.743250     37.000000      6.052381      1.099526   1725.000000   
max       15.000100     52.000000    141.909091     34.066667  35682.000000   

           AveOccup      Latitude     Longitude  
count  20640.000000  20640.000000  20640.000000  
mean       3.070655     35.631861   \PYGZhy{}119.569704  
std       10.386050      2.135952      2.003532  
min        0.692308     32.540000   \PYGZhy{}124.350000  
25\PYGZpc{}        2.429741     33.930000   \PYGZhy{}121.800000  
50\PYGZpc{}        2.818116     34.260000   \PYGZhy{}118.490000  
75\PYGZpc{}        3.282261     37.710000   \PYGZhy{}118.010000  
max     1243.333333     41.950000   \PYGZhy{}114.310000  
\end{sphinxVerbatim}

\end{sphinxuseclass}\end{sphinxVerbatimOutput}

\end{sphinxuseclass}

\subsubsection{Handling Missing Values}
\label{\detokenize{03:handling-missing-values}}
\sphinxAtStartPar
Although the California Housing dataset typically does not contain missing values, let’s artificially introduce some missing values for demonstration. We will then show three approaches:
\begin{enumerate}
\sphinxsetlistlabels{\arabic}{enumi}{enumii}{}{.}%
\item {} 
\sphinxAtStartPar
Dropping rows with missing values.

\item {} 
\sphinxAtStartPar
Imputing with mean or median.

\item {} 
\sphinxAtStartPar
\sphinxcode{\sphinxupquote{KNNImputer}} or another advanced method.

\end{enumerate}


\paragraph{An advanced method: \sphinxstyleliteralintitle{\sphinxupquote{KNNImputer}}}
\label{\detokenize{03:an-advanced-method-knnimputer}}
\sphinxAtStartPar
\sphinxcode{\sphinxupquote{KNNImputer}} uses a K\sphinxhyphen{}Nearest Neighbors approach to fill missing entries: for each row that has missing features, it finds the nearest \sphinxstyleemphasis{k} rows (in feature space) that do not have missing values in those columns, and uses their values (e.g., average) to impute the missing ones. This can often be more accurate than a simple mean or median if the data has local structure—rows that are similar in some features are likely also similar in the others.

\begin{sphinxuseclass}{cell}\begin{sphinxVerbatimInput}

\begin{sphinxuseclass}{cell_input}
\begin{sphinxVerbatim}[commandchars=\\\{\}]
\PYG{k+kn}{from}\PYG{+w}{ }\PYG{n+nn}{sklearn}\PYG{n+nn}{.}\PYG{n+nn}{model\PYGZus{}selection}\PYG{+w}{ }\PYG{k+kn}{import} \PYG{n}{train\PYGZus{}test\PYGZus{}split}
\PYG{k+kn}{from}\PYG{+w}{ }\PYG{n+nn}{sklearn}\PYG{n+nn}{.}\PYG{n+nn}{impute}\PYG{+w}{ }\PYG{k+kn}{import} \PYG{n}{KNNImputer}

\PYG{n}{missing\PYGZus{}mask} \PYG{o}{=} \PYG{n}{np}\PYG{o}{.}\PYG{n}{random}\PYG{o}{.}\PYG{n}{rand}\PYG{p}{(}\PYG{o}{*}\PYG{n}{X\PYGZus{}original}\PYG{o}{.}\PYG{n}{shape}\PYG{p}{)} \PYG{o}{\PYGZlt{}} \PYG{l+m+mf}{0.01}
\PYG{n}{X\PYGZus{}missing} \PYG{o}{=} \PYG{n}{X\PYGZus{}original}\PYG{o}{.}\PYG{n}{copy}\PYG{p}{(}\PYG{p}{)}
\PYG{n}{X\PYGZus{}missing}\PYG{p}{[}\PYG{n}{missing\PYGZus{}mask}\PYG{p}{]} \PYG{o}{=} \PYG{n}{np}\PYG{o}{.}\PYG{n}{nan}

\PYG{n+nb}{print}\PYG{p}{(}\PYG{l+s+s2}{\PYGZdq{}}\PYG{l+s+s2}{\PYGZhy{}\PYGZhy{}\PYGZhy{} Number of missing values (artificial) \PYGZhy{}\PYGZhy{}\PYGZhy{}}\PYG{l+s+s2}{\PYGZdq{}}\PYG{p}{)}
\PYG{n+nb}{print}\PYG{p}{(}\PYG{n}{X\PYGZus{}missing}\PYG{o}{.}\PYG{n}{isnull}\PYG{p}{(}\PYG{p}{)}\PYG{o}{.}\PYG{n}{sum}\PYG{p}{(}\PYG{p}{)}\PYG{p}{)}

\PYG{n}{X\PYGZus{}train\PYGZus{}full}\PYG{p}{,} \PYG{n}{X\PYGZus{}test\PYGZus{}full}\PYG{p}{,} \PYG{n}{y\PYGZus{}train\PYGZus{}full}\PYG{p}{,} \PYG{n}{y\PYGZus{}test\PYGZus{}full} \PYG{o}{=} \PYG{n}{train\PYGZus{}test\PYGZus{}split}\PYG{p}{(}
    \PYG{n}{X\PYGZus{}missing}\PYG{p}{,} \PYG{n}{y\PYGZus{}original}\PYG{p}{,} \PYG{n}{test\PYGZus{}size}\PYG{o}{=}\PYG{l+m+mf}{0.2}\PYG{p}{,} \PYG{n}{random\PYGZus{}state}\PYG{o}{=}\PYG{l+m+mi}{42}
\PYG{p}{)}

\PYG{n}{X\PYGZus{}train\PYGZus{}drop} \PYG{o}{=} \PYG{n}{X\PYGZus{}train\PYGZus{}full}\PYG{o}{.}\PYG{n}{dropna}\PYG{p}{(}\PYG{n}{axis}\PYG{o}{=}\PYG{l+m+mi}{0}\PYG{p}{)}
\PYG{n}{y\PYGZus{}train\PYGZus{}drop} \PYG{o}{=} \PYG{n}{y\PYGZus{}train\PYGZus{}full}\PYG{o}{.}\PYG{n}{loc}\PYG{p}{[}\PYG{n}{X\PYGZus{}train\PYGZus{}drop}\PYG{o}{.}\PYG{n}{index}\PYG{p}{]}
\PYG{n}{X\PYGZus{}test\PYGZus{}drop} \PYG{o}{=} \PYG{n}{X\PYGZus{}test\PYGZus{}full}\PYG{o}{.}\PYG{n}{dropna}\PYG{p}{(}\PYG{n}{axis}\PYG{o}{=}\PYG{l+m+mi}{0}\PYG{p}{)}
\PYG{n}{y\PYGZus{}test\PYGZus{}drop} \PYG{o}{=} \PYG{n}{y\PYGZus{}test\PYGZus{}full}\PYG{o}{.}\PYG{n}{loc}\PYG{p}{[}\PYG{n}{X\PYGZus{}test\PYGZus{}drop}\PYG{o}{.}\PYG{n}{index}\PYG{p}{]}

\PYG{n+nb}{print}\PYG{p}{(}\PYG{l+s+s2}{\PYGZdq{}}\PYG{l+s+s2}{\PYGZhy{}\PYGZhy{}\PYGZhy{} After Dropping Missing Values \PYGZhy{}\PYGZhy{}\PYGZhy{}}\PYG{l+s+s2}{\PYGZdq{}}\PYG{p}{)}
\PYG{n+nb}{print}\PYG{p}{(}\PYG{l+s+s2}{\PYGZdq{}}\PYG{l+s+s2}{Training set size:}\PYG{l+s+s2}{\PYGZdq{}}\PYG{p}{,} \PYG{n}{X\PYGZus{}train\PYGZus{}drop}\PYG{o}{.}\PYG{n}{shape}\PYG{p}{)}
\PYG{n+nb}{print}\PYG{p}{(}\PYG{l+s+s2}{\PYGZdq{}}\PYG{l+s+s2}{Test set size:}\PYG{l+s+s2}{\PYGZdq{}}\PYG{p}{,} \PYG{n}{X\PYGZus{}test\PYGZus{}drop}\PYG{o}{.}\PYG{n}{shape}\PYG{p}{)}

\PYG{k+kn}{from}\PYG{+w}{ }\PYG{n+nn}{sklearn}\PYG{n+nn}{.}\PYG{n+nn}{impute}\PYG{+w}{ }\PYG{k+kn}{import} \PYG{n}{SimpleImputer}

\PYG{n}{mean\PYGZus{}imputer} \PYG{o}{=} \PYG{n}{SimpleImputer}\PYG{p}{(}\PYG{n}{strategy}\PYG{o}{=}\PYG{l+s+s1}{\PYGZsq{}}\PYG{l+s+s1}{mean}\PYG{l+s+s1}{\PYGZsq{}}\PYG{p}{)}
\PYG{n}{X\PYGZus{}train\PYGZus{}mean} \PYG{o}{=} \PYG{n}{mean\PYGZus{}imputer}\PYG{o}{.}\PYG{n}{fit\PYGZus{}transform}\PYG{p}{(}\PYG{n}{X\PYGZus{}train\PYGZus{}full}\PYG{p}{)}
\PYG{n}{X\PYGZus{}test\PYGZus{}mean} \PYG{o}{=} \PYG{n}{mean\PYGZus{}imputer}\PYG{o}{.}\PYG{n}{transform}\PYG{p}{(}\PYG{n}{X\PYGZus{}test\PYGZus{}full}\PYG{p}{)}

\PYG{n}{median\PYGZus{}imputer} \PYG{o}{=} \PYG{n}{SimpleImputer}\PYG{p}{(}\PYG{n}{strategy}\PYG{o}{=}\PYG{l+s+s1}{\PYGZsq{}}\PYG{l+s+s1}{median}\PYG{l+s+s1}{\PYGZsq{}}\PYG{p}{)}
\PYG{n}{X\PYGZus{}train\PYGZus{}median} \PYG{o}{=} \PYG{n}{median\PYGZus{}imputer}\PYG{o}{.}\PYG{n}{fit\PYGZus{}transform}\PYG{p}{(}\PYG{n}{X\PYGZus{}train\PYGZus{}full}\PYG{p}{)}
\PYG{n}{X\PYGZus{}test\PYGZus{}median} \PYG{o}{=} \PYG{n}{median\PYGZus{}imputer}\PYG{o}{.}\PYG{n}{transform}\PYG{p}{(}\PYG{n}{X\PYGZus{}test\PYGZus{}full}\PYG{p}{)}

\PYG{n}{knn\PYGZus{}imputer} \PYG{o}{=} \PYG{n}{KNNImputer}\PYG{p}{(}\PYG{n}{n\PYGZus{}neighbors}\PYG{o}{=}\PYG{l+m+mi}{3}\PYG{p}{)}
\PYG{n}{X\PYGZus{}train\PYGZus{}knn} \PYG{o}{=} \PYG{n}{knn\PYGZus{}imputer}\PYG{o}{.}\PYG{n}{fit\PYGZus{}transform}\PYG{p}{(}\PYG{n}{X\PYGZus{}train\PYGZus{}full}\PYG{p}{)}
\PYG{n}{X\PYGZus{}test\PYGZus{}knn} \PYG{o}{=} \PYG{n}{knn\PYGZus{}imputer}\PYG{o}{.}\PYG{n}{transform}\PYG{p}{(}\PYG{n}{X\PYGZus{}test\PYGZus{}full}\PYG{p}{)}
\end{sphinxVerbatim}

\end{sphinxuseclass}\end{sphinxVerbatimInput}
\begin{sphinxVerbatimOutput}

\begin{sphinxuseclass}{cell_output}
\begin{sphinxVerbatim}[commandchars=\\\{\}]
\PYGZhy{}\PYGZhy{}\PYGZhy{} Number of missing values (artificial) \PYGZhy{}\PYGZhy{}\PYGZhy{}
MedInc        195
HouseAge      201
AveRooms      205
AveBedrms     219
Population    191
AveOccup      206
Latitude      177
Longitude     223
dtype: int64
\PYGZhy{}\PYGZhy{}\PYGZhy{} After Dropping Missing Values \PYGZhy{}\PYGZhy{}\PYGZhy{}
Training set size: (15284, 8)
Test set size: (3798, 8)
\end{sphinxVerbatim}

\end{sphinxuseclass}\end{sphinxVerbatimOutput}

\end{sphinxuseclass}

\subsubsection{Dealing with Categorical Values}
\label{\detokenize{03:dealing-with-categorical-values}}
\sphinxAtStartPar
In this dataset, we have only numerical features. However, real\sphinxhyphen{}world datasets often include categorical variables. Two popular encoding strategies are:
\begin{enumerate}
\sphinxsetlistlabels{\arabic}{enumi}{enumii}{}{.}%
\item {} 
\sphinxAtStartPar
Label Encoding: Each category is assigned an integer.

\item {} 
\sphinxAtStartPar
One\sphinxhyphen{}Hot Encoding (Dummy Encoding): Create a new binary column for each category.

\end{enumerate}

\sphinxAtStartPar
Below is a small example using a synthetic dataframe containing a categorical column.

\begin{sphinxuseclass}{cell}\begin{sphinxVerbatimInput}

\begin{sphinxuseclass}{cell_input}
\begin{sphinxVerbatim}[commandchars=\\\{\}]
\PYG{k+kn}{from}\PYG{+w}{ }\PYG{n+nn}{sklearn}\PYG{n+nn}{.}\PYG{n+nn}{preprocessing}\PYG{+w}{ }\PYG{k+kn}{import} \PYG{n}{LabelEncoder}

\PYG{n}{example\PYGZus{}df} \PYG{o}{=} \PYG{n}{pd}\PYG{o}{.}\PYG{n}{DataFrame}\PYG{p}{(}\PYG{p}{\PYGZob{}}
    \PYG{l+s+s2}{\PYGZdq{}}\PYG{l+s+s2}{Size}\PYG{l+s+s2}{\PYGZdq{}}\PYG{p}{:} \PYG{p}{[}\PYG{l+s+s2}{\PYGZdq{}}\PYG{l+s+s2}{Small}\PYG{l+s+s2}{\PYGZdq{}}\PYG{p}{,} \PYG{l+s+s2}{\PYGZdq{}}\PYG{l+s+s2}{Medium}\PYG{l+s+s2}{\PYGZdq{}}\PYG{p}{,} \PYG{l+s+s2}{\PYGZdq{}}\PYG{l+s+s2}{Large}\PYG{l+s+s2}{\PYGZdq{}}\PYG{p}{,} \PYG{l+s+s2}{\PYGZdq{}}\PYG{l+s+s2}{Medium}\PYG{l+s+s2}{\PYGZdq{}}\PYG{p}{,} \PYG{l+s+s2}{\PYGZdq{}}\PYG{l+s+s2}{Small}\PYG{l+s+s2}{\PYGZdq{}}\PYG{p}{]}\PYG{p}{,}
    \PYG{l+s+s2}{\PYGZdq{}}\PYG{l+s+s2}{Weight}\PYG{l+s+s2}{\PYGZdq{}}\PYG{p}{:} \PYG{p}{[}\PYG{l+m+mf}{1.0}\PYG{p}{,} \PYG{l+m+mf}{2.3}\PYG{p}{,} \PYG{l+m+mf}{5.5}\PYG{p}{,} \PYG{l+m+mf}{2.0}\PYG{p}{,} \PYG{l+m+mf}{1.1}\PYG{p}{]}
\PYG{p}{\PYGZcb{}}\PYG{p}{)}

\PYG{n}{label\PYGZus{}encoder} \PYG{o}{=} \PYG{n}{LabelEncoder}\PYG{p}{(}\PYG{p}{)}
\PYG{n}{example\PYGZus{}df}\PYG{p}{[}\PYG{l+s+s2}{\PYGZdq{}}\PYG{l+s+s2}{Size\PYGZus{}label}\PYG{l+s+s2}{\PYGZdq{}}\PYG{p}{]} \PYG{o}{=} \PYG{n}{label\PYGZus{}encoder}\PYG{o}{.}\PYG{n}{fit\PYGZus{}transform}\PYG{p}{(}\PYG{n}{example\PYGZus{}df}\PYG{p}{[}\PYG{l+s+s2}{\PYGZdq{}}\PYG{l+s+s2}{Size}\PYG{l+s+s2}{\PYGZdq{}}\PYG{p}{]}\PYG{p}{)}

\PYG{n}{one\PYGZus{}hot\PYGZus{}df} \PYG{o}{=} \PYG{n}{pd}\PYG{o}{.}\PYG{n}{get\PYGZus{}dummies}\PYG{p}{(}\PYG{n}{example\PYGZus{}df}\PYG{p}{,} \PYG{n}{columns}\PYG{o}{=}\PYG{p}{[}\PYG{l+s+s2}{\PYGZdq{}}\PYG{l+s+s2}{Size}\PYG{l+s+s2}{\PYGZdq{}}\PYG{p}{]}\PYG{p}{,} \PYG{n}{prefix}\PYG{o}{=}\PYG{l+s+s2}{\PYGZdq{}}\PYG{l+s+s2}{Size}\PYG{l+s+s2}{\PYGZdq{}}\PYG{p}{)}
\end{sphinxVerbatim}

\end{sphinxuseclass}\end{sphinxVerbatimInput}

\end{sphinxuseclass}
\sphinxAtStartPar
Original DataFrame with Label Encoding:

\begin{sphinxuseclass}{cell}\begin{sphinxVerbatimInput}

\begin{sphinxuseclass}{cell_input}
\begin{sphinxVerbatim}[commandchars=\\\{\}]
\PYG{n}{display}\PYG{p}{(}\PYG{n}{example\PYGZus{}df}\PYG{p}{)}
\end{sphinxVerbatim}

\end{sphinxuseclass}\end{sphinxVerbatimInput}
\begin{sphinxVerbatimOutput}

\begin{sphinxuseclass}{cell_output}
\begin{sphinxVerbatim}[commandchars=\\\{\}]
     Size  Weight  Size\PYGZus{}label
0   Small     1.0           2
1  Medium     2.3           1
2   Large     5.5           0
3  Medium     2.0           1
4   Small     1.1           2
\end{sphinxVerbatim}

\end{sphinxuseclass}\end{sphinxVerbatimOutput}

\end{sphinxuseclass}
\sphinxAtStartPar
DataFrame with One\sphinxhyphen{}Hot Encoding:

\begin{sphinxuseclass}{cell}\begin{sphinxVerbatimInput}

\begin{sphinxuseclass}{cell_input}
\begin{sphinxVerbatim}[commandchars=\\\{\}]
\PYG{n}{display}\PYG{p}{(}\PYG{n}{one\PYGZus{}hot\PYGZus{}df}\PYG{p}{)}
\end{sphinxVerbatim}

\end{sphinxuseclass}\end{sphinxVerbatimInput}
\begin{sphinxVerbatimOutput}

\begin{sphinxuseclass}{cell_output}
\begin{sphinxVerbatim}[commandchars=\\\{\}]
   Weight  Size\PYGZus{}label  Size\PYGZus{}Large  Size\PYGZus{}Medium  Size\PYGZus{}Small
0     1.0           2       False        False        True
1     2.3           1       False         True       False
2     5.5           0        True        False       False
3     2.0           1       False         True       False
4     1.1           2       False        False        True
\end{sphinxVerbatim}

\end{sphinxuseclass}\end{sphinxVerbatimOutput}

\end{sphinxuseclass}
\sphinxAtStartPar
We’ll stick to numeric data for the main practical, but this is how you’d handle categorical features.


\subsubsection{Normalization/Scaling}
\label{\detokenize{03:normalization-scaling}}
\sphinxAtStartPar
Scaling is critical for distance\sphinxhyphen{}based methods. We will use \sphinxcode{\sphinxupquote{StandardScaler}}, which transforms each feature to have zero mean and unit variance, but note that \sphinxcode{\sphinxupquote{MinMaxScaler}} is also a standard.
\begin{quote}

\sphinxAtStartPar
\sphinxstylestrong{Remember}: \sphinxstyleemphasis{Always fit the scaler on the training set, then apply the same transformation to the test set. This prevents data leakage and preserves the integrity of the test data.}
\end{quote}

\begin{sphinxuseclass}{cell}\begin{sphinxVerbatimInput}

\begin{sphinxuseclass}{cell_input}
\begin{sphinxVerbatim}[commandchars=\\\{\}]
\PYG{k+kn}{from}\PYG{+w}{ }\PYG{n+nn}{sklearn}\PYG{n+nn}{.}\PYG{n+nn}{preprocessing}\PYG{+w}{ }\PYG{k+kn}{import} \PYG{n}{StandardScaler}

\PYG{n}{scaler} \PYG{o}{=} \PYG{n}{StandardScaler}\PYG{p}{(}\PYG{p}{)}
\PYG{n}{X\PYGZus{}train\PYGZus{}scaled} \PYG{o}{=} \PYG{n}{scaler}\PYG{o}{.}\PYG{n}{fit\PYGZus{}transform}\PYG{p}{(}\PYG{n}{X\PYGZus{}train\PYGZus{}mean}\PYG{p}{)}
\PYG{n}{X\PYGZus{}test\PYGZus{}scaled} \PYG{o}{=} \PYG{n}{scaler}\PYG{o}{.}\PYG{n}{transform}\PYG{p}{(}\PYG{n}{X\PYGZus{}test\PYGZus{}mean}\PYG{p}{)}

\PYG{n}{X\PYGZus{}train\PYGZus{}final} \PYG{o}{=} \PYG{n}{X\PYGZus{}train\PYGZus{}scaled}
\PYG{n}{X\PYGZus{}test\PYGZus{}final} \PYG{o}{=} \PYG{n}{X\PYGZus{}test\PYGZus{}scaled}
\PYG{n}{y\PYGZus{}train\PYGZus{}final} \PYG{o}{=} \PYG{n}{y\PYGZus{}train\PYGZus{}full}
\PYG{n}{y\PYGZus{}test\PYGZus{}final} \PYG{o}{=} \PYG{n}{y\PYGZus{}test\PYGZus{}full}
\end{sphinxVerbatim}

\end{sphinxuseclass}\end{sphinxVerbatimInput}

\end{sphinxuseclass}

\subsection{Vanilla KNN Implementation}
\label{\detokenize{03:vanilla-knn-implementation}}\begin{enumerate}
\sphinxsetlistlabels{\arabic}{enumi}{enumii}{}{.}%
\item {} 
\sphinxAtStartPar
Split your data into training and test sets (we should have already done that in the previous exercise).

\item {} 
\sphinxAtStartPar
Implement a vanilla KNN regressor from scratch (no \sphinxcode{\sphinxupquote{scikit\sphinxhyphen{}learn}} for this part).

\item {} 
\sphinxAtStartPar
Choose a value of \(k\) (e.g., 3, 5, 7).

\item {} 
\sphinxAtStartPar
Compare predictions with the true values using Mean Squared Error (MSE).

\end{enumerate}


\subsubsection{Solution}
\label{\detokenize{03:id3}}
\begin{sphinxShadowBox}
\sphinxstylesidebartitle{}

\sphinxAtStartPar
On this graph, we clearly see that there is an optimal value for the number of neighbors to use: a value too small increases the variance of the predictions while a value too high increases the bias. The optimal value, of course, depends on the problem definition and on the training data.
\end{sphinxShadowBox}

\begin{sphinxuseclass}{cell}\begin{sphinxVerbatimInput}

\begin{sphinxuseclass}{cell_input}
\begin{sphinxVerbatim}[commandchars=\\\{\}]
\PYG{k+kn}{import}\PYG{+w}{ }\PYG{n+nn}{matplotlib}\PYG{n+nn}{.}\PYG{n+nn}{pyplot}\PYG{+w}{ }\PYG{k}{as}\PYG{+w}{ }\PYG{n+nn}{plt}
\PYG{n}{plt}\PYG{o}{.}\PYG{n}{style}\PYG{o}{.}\PYG{n}{use}\PYG{p}{(}\PYG{p}{[}\PYG{l+s+s1}{\PYGZsq{}}\PYG{l+s+s1}{seaborn\PYGZhy{}v0\PYGZus{}8\PYGZhy{}muted}\PYG{l+s+s1}{\PYGZsq{}}\PYG{p}{,} \PYG{l+s+s1}{\PYGZsq{}}\PYG{l+s+s1}{practicals.mplstyle}\PYG{l+s+s1}{\PYGZsq{}}\PYG{p}{]}\PYG{p}{)}
\PYG{k+kn}{from}\PYG{+w}{ }\PYG{n+nn}{sklearn}\PYG{n+nn}{.}\PYG{n+nn}{metrics}\PYG{+w}{ }\PYG{k+kn}{import} \PYG{n}{mean\PYGZus{}squared\PYGZus{}error}
\PYG{k+kn}{from}\PYG{+w}{ }\PYG{n+nn}{tqdm}\PYG{n+nn}{.}\PYG{n+nn}{auto}\PYG{+w}{ }\PYG{k+kn}{import} \PYG{n}{tqdm}

\PYG{k}{def}\PYG{+w}{ }\PYG{n+nf}{euclidean\PYGZus{}distance}\PYG{p}{(}\PYG{n}{x1}\PYG{p}{,} \PYG{n}{x2}\PYG{p}{)}\PYG{p}{:}
    \PYG{k}{return} \PYG{n}{np}\PYG{o}{.}\PYG{n}{sqrt}\PYG{p}{(}\PYG{n}{np}\PYG{o}{.}\PYG{n}{sum}\PYG{p}{(}\PYG{p}{(}\PYG{n}{x1} \PYG{o}{\PYGZhy{}} \PYG{n}{x2}\PYG{p}{)} \PYG{o}{*}\PYG{o}{*} \PYG{l+m+mi}{2}\PYG{p}{)}\PYG{p}{)}

\PYG{k}{def}\PYG{+w}{ }\PYG{n+nf}{knn\PYGZus{}regressor}\PYG{p}{(}\PYG{n}{X\PYGZus{}train}\PYG{p}{,} \PYG{n}{y\PYGZus{}train}\PYG{p}{,} \PYG{n}{X\PYGZus{}test}\PYG{p}{,} \PYG{n}{k}\PYG{o}{=}\PYG{l+m+mi}{3}\PYG{p}{)}\PYG{p}{:}
\PYG{+w}{    }\PYG{l+s+sd}{\PYGZdq{}\PYGZdq{}\PYGZdq{}}
\PYG{l+s+sd}{    A simple KNN regressor (manual implementation).}
\PYG{l+s+sd}{    \PYGZhy{} X\PYGZus{}train: training features (numpy array)}
\PYG{l+s+sd}{    \PYGZhy{} y\PYGZus{}train: training targets (numpy array)}
\PYG{l+s+sd}{    \PYGZhy{} X\PYGZus{}test: test features (numpy array)}
\PYG{l+s+sd}{    \PYGZhy{} k: number of neighbors to consider}
\PYG{l+s+sd}{    \PYGZdq{}\PYGZdq{}\PYGZdq{}}
    \PYG{n}{predictions} \PYG{o}{=} \PYG{p}{[}\PYG{p}{]}
    \PYG{k}{for} \PYG{n}{test\PYGZus{}point} \PYG{o+ow}{in} \PYG{n}{tqdm}\PYG{p}{(}\PYG{n}{X\PYGZus{}test}\PYG{p}{)}\PYG{p}{:}
        \PYG{n}{distances} \PYG{o}{=} \PYG{p}{[}\PYG{n}{euclidean\PYGZus{}distance}\PYG{p}{(}\PYG{n}{test\PYGZus{}point}\PYG{p}{,} \PYG{n}{x\PYGZus{}train}\PYG{p}{)} \PYG{k}{for} \PYG{n}{x\PYGZus{}train} \PYG{o+ow}{in} \PYG{n}{X\PYGZus{}train}\PYG{p}{]}
        \PYG{n}{k\PYGZus{}indices} \PYG{o}{=} \PYG{n}{np}\PYG{o}{.}\PYG{n}{argsort}\PYG{p}{(}\PYG{n}{distances}\PYG{p}{)}\PYG{p}{[}\PYG{p}{:}\PYG{n}{k}\PYG{p}{]}
        \PYG{n}{prediction} \PYG{o}{=} \PYG{n}{np}\PYG{o}{.}\PYG{n}{mean}\PYG{p}{(}\PYG{n}{y\PYGZus{}train}\PYG{p}{[}\PYG{n}{k\PYGZus{}indices}\PYG{p}{]}\PYG{p}{)}
        \PYG{n}{predictions}\PYG{o}{.}\PYG{n}{append}\PYG{p}{(}\PYG{n}{prediction}\PYG{p}{)}
    \PYG{k}{return} \PYG{n}{np}\PYG{o}{.}\PYG{n}{array}\PYG{p}{(}\PYG{n}{predictions}\PYG{p}{)}

\PYG{k+kn}{from}\PYG{+w}{ }\PYG{n+nn}{scipy}\PYG{n+nn}{.}\PYG{n+nn}{spatial}\PYG{+w}{ }\PYG{k+kn}{import} \PYG{n}{distance}

\PYG{k}{def}\PYG{+w}{ }\PYG{n+nf}{knn\PYGZus{}regressor\PYGZus{}fast}\PYG{p}{(}\PYG{n}{X\PYGZus{}train}\PYG{p}{,} \PYG{n}{y\PYGZus{}train}\PYG{p}{,} \PYG{n}{X\PYGZus{}test}\PYG{p}{,} \PYG{n}{k}\PYG{o}{=}\PYG{l+m+mi}{3}\PYG{p}{)}\PYG{p}{:}
\PYG{+w}{    }\PYG{l+s+sd}{\PYGZdq{}\PYGZdq{}\PYGZdq{}}
\PYG{l+s+sd}{    A faster KNN regressor using vectorized distance computation.}
\PYG{l+s+sd}{    \PYGZdq{}\PYGZdq{}\PYGZdq{}}
    \PYG{n}{dists} \PYG{o}{=} \PYG{n}{distance}\PYG{o}{.}\PYG{n}{cdist}\PYG{p}{(}\PYG{n}{X\PYGZus{}test}\PYG{p}{,} \PYG{n}{X\PYGZus{}train}\PYG{p}{,} \PYG{n}{metric}\PYG{o}{=}\PYG{l+s+s1}{\PYGZsq{}}\PYG{l+s+s1}{euclidean}\PYG{l+s+s1}{\PYGZsq{}}\PYG{p}{)}
    \PYG{n}{k\PYGZus{}indices} \PYG{o}{=} \PYG{n}{np}\PYG{o}{.}\PYG{n}{argpartition}\PYG{p}{(}\PYG{n}{dists}\PYG{p}{,} \PYG{n}{kth}\PYG{o}{=}\PYG{n}{k}\PYG{p}{,} \PYG{n}{axis}\PYG{o}{=}\PYG{l+m+mi}{1}\PYG{p}{)}\PYG{p}{[}\PYG{p}{:}\PYG{p}{,} \PYG{p}{:}\PYG{n}{k}\PYG{p}{]}
    \PYG{n}{predictions} \PYG{o}{=} \PYG{n}{np}\PYG{o}{.}\PYG{n}{mean}\PYG{p}{(}\PYG{n}{y\PYGZus{}train}\PYG{p}{[}\PYG{n}{k\PYGZus{}indices}\PYG{p}{]}\PYG{p}{,} \PYG{n}{axis}\PYG{o}{=}\PYG{l+m+mi}{1}\PYG{p}{)}
    \PYG{k}{return} \PYG{n}{predictions}

\PYG{n}{mses\PYGZus{}knn} \PYG{o}{=} \PYG{p}{[}\PYG{p}{]}
\PYG{k}{for} \PYG{n}{k} \PYG{o+ow}{in} \PYG{n}{tqdm}\PYG{p}{(}\PYG{n+nb}{range}\PYG{p}{(}\PYG{l+m+mi}{1}\PYG{p}{,} \PYG{l+m+mi}{51}\PYG{p}{,} \PYG{l+m+mi}{2}\PYG{p}{)}\PYG{p}{)}\PYG{p}{:}
  \PYG{n}{y\PYGZus{}pred\PYGZus{}knn} \PYG{o}{=} \PYG{n}{knn\PYGZus{}regressor\PYGZus{}fast}\PYG{p}{(}\PYG{n}{X\PYGZus{}train\PYGZus{}final}\PYG{p}{,} \PYG{n}{y\PYGZus{}train\PYGZus{}final}\PYG{o}{.}\PYG{n}{values}\PYG{p}{,} \PYG{n}{X\PYGZus{}test\PYGZus{}final}\PYG{p}{,} \PYG{n}{k}\PYG{o}{=}\PYG{n}{k}\PYG{p}{)}
  \PYG{n}{mses\PYGZus{}knn}\PYG{o}{.}\PYG{n}{append}\PYG{p}{(}\PYG{n}{mean\PYGZus{}squared\PYGZus{}error}\PYG{p}{(}\PYG{n}{y\PYGZus{}test\PYGZus{}final}\PYG{p}{,} \PYG{n}{y\PYGZus{}pred\PYGZus{}knn}\PYG{p}{)}\PYG{p}{)}

\PYG{n}{plt}\PYG{o}{.}\PYG{n}{figure}\PYG{p}{(}\PYG{n}{figsize}\PYG{o}{=}\PYG{p}{(}\PYG{l+m+mi}{10}\PYG{p}{,} \PYG{l+m+mi}{6}\PYG{p}{)}\PYG{p}{)}
\PYG{n}{plt}\PYG{o}{.}\PYG{n}{plot}\PYG{p}{(}\PYG{n+nb}{list}\PYG{p}{(}\PYG{n+nb}{range}\PYG{p}{(}\PYG{l+m+mi}{1}\PYG{p}{,} \PYG{l+m+mi}{51}\PYG{p}{,} \PYG{l+m+mi}{2}\PYG{p}{)}\PYG{p}{)}\PYG{p}{,} \PYG{n}{mses\PYGZus{}knn}\PYG{p}{,} \PYG{n}{marker}\PYG{o}{=}\PYG{l+s+s1}{\PYGZsq{}}\PYG{l+s+s1}{o}\PYG{l+s+s1}{\PYGZsq{}}\PYG{p}{)}
\PYG{n}{plt}\PYG{o}{.}\PYG{n}{title}\PYG{p}{(}\PYG{l+s+s2}{\PYGZdq{}}\PYG{l+s+s2}{MSE vs K (KNN)}\PYG{l+s+s2}{\PYGZdq{}}\PYG{p}{)}
\PYG{n}{plt}\PYG{o}{.}\PYG{n}{xlabel}\PYG{p}{(}\PYG{l+s+s2}{\PYGZdq{}}\PYG{l+s+s2}{K}\PYG{l+s+s2}{\PYGZdq{}}\PYG{p}{)}
\PYG{n}{plt}\PYG{o}{.}\PYG{n}{ylabel}\PYG{p}{(}\PYG{l+s+s2}{\PYGZdq{}}\PYG{l+s+s2}{MSE}\PYG{l+s+s2}{\PYGZdq{}}\PYG{p}{)}
\PYG{n}{plt}\PYG{o}{.}\PYG{n}{grid}\PYG{p}{(}\PYG{l+s+s2}{\PYGZdq{}}\PYG{l+s+s2}{on}\PYG{l+s+s2}{\PYGZdq{}}\PYG{p}{)}
\PYG{n}{plt}\PYG{o}{.}\PYG{n}{show}\PYG{p}{(}\PYG{p}{)}
\end{sphinxVerbatim}

\end{sphinxuseclass}\end{sphinxVerbatimInput}
\begin{sphinxVerbatimOutput}

\begin{sphinxuseclass}{cell_output}
\begin{sphinxVerbatim}[commandchars=\\\{\}]
  0\PYGZpc{}|          | 0/25 [00:00\PYGZlt{}?, ?it/s]
\end{sphinxVerbatim}

\noindent\sphinxincludegraphics{{bd03f00d3bb200eb71e6f9599e8bbaf2231030fb46841e810d48d3cbe6009a70}.png}

\end{sphinxuseclass}\end{sphinxVerbatimOutput}

\end{sphinxuseclass}

\subsection{Regression Trees and Random Forests}
\label{\detokenize{03:regression-trees-and-random-forests}}\begin{enumerate}
\sphinxsetlistlabels{\arabic}{enumi}{enumii}{}{.}%
\item {} 
\sphinxAtStartPar
Train a regression tree on the training set.

\item {} 
\sphinxAtStartPar
Train a random forest on the training set.

\item {} 
\sphinxAtStartPar
Evaluate both models using MSE.

\item {} 
\sphinxAtStartPar
Compare their performance with your KNN model.

\item {} 
\sphinxAtStartPar
Determine how modifying the \sphinxcode{\sphinxupquote{max\_depth}} (for the regression trees) and \sphinxcode{\sphinxupquote{n\_estimators}} (for the random forests) impacts the performances of the model. Plot these results graphically.

\item {} 
\sphinxAtStartPar
Plot the features importance of the best model with the optimal parameters

\end{enumerate}


\subsection{Solution}
\label{\detokenize{03:id4}}

\subsubsection{Implementation with \sphinxstyleliteralintitle{\sphinxupquote{scikit\sphinxhyphen{}learn}}}
\label{\detokenize{03:implementation-with-scikit-learn}}
\begin{sphinxuseclass}{cell}\begin{sphinxVerbatimInput}

\begin{sphinxuseclass}{cell_input}
\begin{sphinxVerbatim}[commandchars=\\\{\}]
\PYG{k+kn}{from}\PYG{+w}{ }\PYG{n+nn}{sklearn}\PYG{n+nn}{.}\PYG{n+nn}{tree}\PYG{+w}{ }\PYG{k+kn}{import} \PYG{n}{DecisionTreeRegressor}
\PYG{k+kn}{from}\PYG{+w}{ }\PYG{n+nn}{sklearn}\PYG{n+nn}{.}\PYG{n+nn}{ensemble}\PYG{+w}{ }\PYG{k+kn}{import} \PYG{n}{RandomForestRegressor}

\PYG{n}{tree\PYGZus{}reg} \PYG{o}{=} \PYG{n}{DecisionTreeRegressor}\PYG{p}{(}\PYG{n}{max\PYGZus{}depth}\PYG{o}{=}\PYG{l+m+mi}{9}\PYG{p}{,} \PYG{n}{random\PYGZus{}state}\PYG{o}{=}\PYG{l+m+mi}{42}\PYG{p}{)}
\PYG{n}{tree\PYGZus{}reg}\PYG{o}{.}\PYG{n}{fit}\PYG{p}{(}\PYG{n}{X\PYGZus{}train\PYGZus{}final}\PYG{p}{,} \PYG{n}{y\PYGZus{}train\PYGZus{}final}\PYG{p}{)}
\PYG{n}{y\PYGZus{}pred\PYGZus{}tree} \PYG{o}{=} \PYG{n}{tree\PYGZus{}reg}\PYG{o}{.}\PYG{n}{predict}\PYG{p}{(}\PYG{n}{X\PYGZus{}test\PYGZus{}final}\PYG{p}{)}
\PYG{n}{mse\PYGZus{}tree} \PYG{o}{=} \PYG{n}{mean\PYGZus{}squared\PYGZus{}error}\PYG{p}{(}\PYG{n}{y\PYGZus{}test\PYGZus{}final}\PYG{p}{,} \PYG{n}{y\PYGZus{}pred\PYGZus{}tree}\PYG{p}{)}

\PYG{n}{forest\PYGZus{}reg} \PYG{o}{=} \PYG{n}{RandomForestRegressor}\PYG{p}{(}\PYG{n}{n\PYGZus{}estimators}\PYG{o}{=}\PYG{l+m+mi}{60}\PYG{p}{,} \PYG{n}{random\PYGZus{}state}\PYG{o}{=}\PYG{l+m+mi}{42}\PYG{p}{)}
\PYG{n}{forest\PYGZus{}reg}\PYG{o}{.}\PYG{n}{fit}\PYG{p}{(}\PYG{n}{X\PYGZus{}train\PYGZus{}final}\PYG{p}{,} \PYG{n}{y\PYGZus{}train\PYGZus{}final}\PYG{p}{)}
\PYG{n}{y\PYGZus{}pred\PYGZus{}forest} \PYG{o}{=} \PYG{n}{forest\PYGZus{}reg}\PYG{o}{.}\PYG{n}{predict}\PYG{p}{(}\PYG{n}{X\PYGZus{}test\PYGZus{}final}\PYG{p}{)}
\PYG{n}{mse\PYGZus{}forest} \PYG{o}{=} \PYG{n}{mean\PYGZus{}squared\PYGZus{}error}\PYG{p}{(}\PYG{n}{y\PYGZus{}test\PYGZus{}final}\PYG{p}{,} \PYG{n}{y\PYGZus{}pred\PYGZus{}forest}\PYG{p}{)}

\PYG{n+nb}{print}\PYG{p}{(}\PYG{l+s+s2}{\PYGZdq{}}\PYG{l+s+s2}{Comparison of MSE:}\PYG{l+s+s2}{\PYGZdq{}}\PYG{p}{)}
\PYG{n+nb}{print}\PYG{p}{(}\PYG{l+s+sa}{f}\PYG{l+s+s2}{\PYGZdq{}}\PYG{l+s+s2}{KNN (k=13): }\PYG{l+s+si}{\PYGZob{}}\PYG{n}{mean\PYGZus{}squared\PYGZus{}error}\PYG{p}{(}\PYG{n}{y\PYGZus{}test\PYGZus{}final}\PYG{p}{,}\PYG{+w}{ }\PYG{n}{knn\PYGZus{}regressor\PYGZus{}fast}\PYG{p}{(}\PYG{n}{X\PYGZus{}train\PYGZus{}final}\PYG{p}{,}\PYG{+w}{ }\PYG{n}{y\PYGZus{}train\PYGZus{}final}\PYG{o}{.}\PYG{n}{values}\PYG{p}{,}\PYG{+w}{ }\PYG{n}{X\PYGZus{}test\PYGZus{}final}\PYG{p}{,}\PYG{+w}{ }\PYG{n}{k}\PYG{o}{=}\PYG{l+m+mi}{13}\PYG{p}{)}\PYG{p}{)}\PYG{l+s+si}{:}\PYG{l+s+s2}{.4f}\PYG{l+s+si}{\PYGZcb{}}\PYG{l+s+s2}{\PYGZdq{}}\PYG{p}{)}
\PYG{n+nb}{print}\PYG{p}{(}\PYG{l+s+sa}{f}\PYG{l+s+s2}{\PYGZdq{}}\PYG{l+s+s2}{Decision Tree: }\PYG{l+s+si}{\PYGZob{}}\PYG{n}{mse\PYGZus{}tree}\PYG{l+s+si}{:}\PYG{l+s+s2}{.4f}\PYG{l+s+si}{\PYGZcb{}}\PYG{l+s+s2}{\PYGZdq{}}\PYG{p}{)}
\PYG{n+nb}{print}\PYG{p}{(}\PYG{l+s+sa}{f}\PYG{l+s+s2}{\PYGZdq{}}\PYG{l+s+s2}{Random Forest: }\PYG{l+s+si}{\PYGZob{}}\PYG{n}{mse\PYGZus{}forest}\PYG{l+s+si}{:}\PYG{l+s+s2}{.4f}\PYG{l+s+si}{\PYGZcb{}}\PYG{l+s+s2}{\PYGZdq{}}\PYG{p}{)}
\end{sphinxVerbatim}

\end{sphinxuseclass}\end{sphinxVerbatimInput}
\begin{sphinxVerbatimOutput}

\begin{sphinxuseclass}{cell_output}
\begin{sphinxVerbatim}[commandchars=\\\{\}]
Comparison of MSE:
\end{sphinxVerbatim}

\begin{sphinxVerbatim}[commandchars=\\\{\}]
KNN (k=13): 0.4419
Decision Tree: 0.4417
Random Forest: 0.2693
\end{sphinxVerbatim}

\end{sphinxuseclass}\end{sphinxVerbatimOutput}

\end{sphinxuseclass}

\paragraph{Exploring Key Parameters}
\label{\detokenize{03:exploring-key-parameters}}

\subparagraph{Decision Tree: \sphinxstyleliteralintitle{\sphinxupquote{max\_depth}}}
\label{\detokenize{03:decision-tree-max-depth}}
\sphinxAtStartPar
We can examine how \sphinxcode{\sphinxupquote{max\_depth}} (the maximum number of splits from the root to the leaf) affects performance.

\begin{sphinxuseclass}{cell}\begin{sphinxVerbatimInput}

\begin{sphinxuseclass}{cell_input}
\begin{sphinxVerbatim}[commandchars=\\\{\}]
\PYG{n}{depths} \PYG{o}{=} \PYG{n+nb}{list}\PYG{p}{(}\PYG{n+nb}{range}\PYG{p}{(}\PYG{l+m+mi}{2}\PYG{p}{,} \PYG{l+m+mi}{21}\PYG{p}{)}\PYG{p}{)} \PYG{o}{+} \PYG{p}{[}\PYG{k+kc}{None}\PYG{p}{]}
\PYG{n}{mse\PYGZus{}values\PYGZus{}tree} \PYG{o}{=} \PYG{p}{[}\PYG{p}{]}

\PYG{k}{for} \PYG{n}{d} \PYG{o+ow}{in} \PYG{n}{depths}\PYG{p}{:}
    \PYG{n}{dt} \PYG{o}{=} \PYG{n}{DecisionTreeRegressor}\PYG{p}{(}\PYG{n}{max\PYGZus{}depth}\PYG{o}{=}\PYG{n}{d}\PYG{p}{,} \PYG{n}{random\PYGZus{}state}\PYG{o}{=}\PYG{l+m+mi}{42}\PYG{p}{)}
    \PYG{n}{dt}\PYG{o}{.}\PYG{n}{fit}\PYG{p}{(}\PYG{n}{X\PYGZus{}train\PYGZus{}final}\PYG{p}{,} \PYG{n}{y\PYGZus{}train\PYGZus{}final}\PYG{p}{)}
    \PYG{n}{y\PYGZus{}pred\PYGZus{}dt} \PYG{o}{=} \PYG{n}{dt}\PYG{o}{.}\PYG{n}{predict}\PYG{p}{(}\PYG{n}{X\PYGZus{}test\PYGZus{}final}\PYG{p}{)}
    \PYG{n}{mse\PYGZus{}values\PYGZus{}tree}\PYG{o}{.}\PYG{n}{append}\PYG{p}{(}\PYG{n}{mean\PYGZus{}squared\PYGZus{}error}\PYG{p}{(}\PYG{n}{y\PYGZus{}test\PYGZus{}final}\PYG{p}{,} \PYG{n}{y\PYGZus{}pred\PYGZus{}dt}\PYG{p}{)}\PYG{p}{)}

\PYG{n}{plt}\PYG{o}{.}\PYG{n}{figure}\PYG{p}{(}\PYG{n}{figsize}\PYG{o}{=}\PYG{p}{(}\PYG{l+m+mi}{10}\PYG{p}{,} \PYG{l+m+mi}{6}\PYG{p}{)}\PYG{p}{)}
\PYG{n}{plt}\PYG{o}{.}\PYG{n}{plot}\PYG{p}{(}\PYG{p}{[}\PYG{n+nb}{str}\PYG{p}{(}\PYG{n}{d}\PYG{p}{)} \PYG{k}{for} \PYG{n}{d} \PYG{o+ow}{in} \PYG{n}{depths}\PYG{p}{]}\PYG{p}{,} \PYG{n}{mse\PYGZus{}values\PYGZus{}tree}\PYG{p}{,} \PYG{n}{marker}\PYG{o}{=}\PYG{l+s+s1}{\PYGZsq{}}\PYG{l+s+s1}{o}\PYG{l+s+s1}{\PYGZsq{}}\PYG{p}{)}
\PYG{n}{plt}\PYG{o}{.}\PYG{n}{title}\PYG{p}{(}\PYG{l+s+s2}{\PYGZdq{}}\PYG{l+s+s2}{MSE vs max\PYGZus{}depth (Decision Tree)}\PYG{l+s+s2}{\PYGZdq{}}\PYG{p}{)}
\PYG{n}{plt}\PYG{o}{.}\PYG{n}{xlabel}\PYG{p}{(}\PYG{l+s+s2}{\PYGZdq{}}\PYG{l+s+s2}{max\PYGZus{}depth}\PYG{l+s+s2}{\PYGZdq{}}\PYG{p}{)}
\PYG{n}{plt}\PYG{o}{.}\PYG{n}{ylabel}\PYG{p}{(}\PYG{l+s+s2}{\PYGZdq{}}\PYG{l+s+s2}{MSE}\PYG{l+s+s2}{\PYGZdq{}}\PYG{p}{)}
\PYG{n}{plt}\PYG{o}{.}\PYG{n}{grid}\PYG{p}{(}\PYG{l+s+s2}{\PYGZdq{}}\PYG{l+s+s2}{on}\PYG{l+s+s2}{\PYGZdq{}}\PYG{p}{)}
\PYG{n}{plt}\PYG{o}{.}\PYG{n}{show}\PYG{p}{(}\PYG{p}{)}
\end{sphinxVerbatim}

\end{sphinxuseclass}\end{sphinxVerbatimInput}
\begin{sphinxVerbatimOutput}

\begin{sphinxuseclass}{cell_output}
\noindent\sphinxincludegraphics{{508b01aca098ca5b4df94fa8fc83044de632e8f88d65cef139469c8311346e8d}.png}

\end{sphinxuseclass}\end{sphinxVerbatimOutput}

\end{sphinxuseclass}

\paragraph{Random Forest: \sphinxstyleliteralintitle{\sphinxupquote{n\_estimators}}}
\label{\detokenize{03:random-forest-n-estimators}}
\sphinxAtStartPar
Similarly, we can see how \sphinxcode{\sphinxupquote{n\_estimators}} (the number of trees in the forest) impacts the MSE.

\begin{sphinxuseclass}{cell}\begin{sphinxVerbatimInput}

\begin{sphinxuseclass}{cell_input}
\begin{sphinxVerbatim}[commandchars=\\\{\}]
\PYG{n}{estimators\PYGZus{}range} \PYG{o}{=} \PYG{n+nb}{list}\PYG{p}{(}\PYG{n+nb}{range}\PYG{p}{(}\PYG{l+m+mi}{1}\PYG{p}{,} \PYG{l+m+mi}{10}\PYG{p}{)}\PYG{p}{)}\PYG{o}{+}\PYG{n+nb}{list}\PYG{p}{(}\PYG{n+nb}{range}\PYG{p}{(}\PYG{l+m+mi}{10}\PYG{p}{,} \PYG{l+m+mi}{50}\PYG{p}{,} \PYG{l+m+mi}{2}\PYG{p}{)}\PYG{p}{)}\PYG{o}{+}\PYG{n+nb}{list}\PYG{p}{(}\PYG{n+nb}{range}\PYG{p}{(}\PYG{l+m+mi}{50}\PYG{p}{,} \PYG{l+m+mi}{100}\PYG{p}{,} \PYG{l+m+mi}{5}\PYG{p}{)}\PYG{p}{)}
\PYG{n}{mse\PYGZus{}values\PYGZus{}rf} \PYG{o}{=} \PYG{p}{[}\PYG{p}{]}

\PYG{k}{for} \PYG{n}{n} \PYG{o+ow}{in} \PYG{n}{tqdm}\PYG{p}{(}\PYG{n}{estimators\PYGZus{}range}\PYG{p}{)}\PYG{p}{:}
    \PYG{n}{rf} \PYG{o}{=} \PYG{n}{RandomForestRegressor}\PYG{p}{(}\PYG{n}{n\PYGZus{}estimators}\PYG{o}{=}\PYG{n}{n}\PYG{p}{,} \PYG{n}{random\PYGZus{}state}\PYG{o}{=}\PYG{l+m+mi}{42}\PYG{p}{)}
    \PYG{n}{rf}\PYG{o}{.}\PYG{n}{fit}\PYG{p}{(}\PYG{n}{X\PYGZus{}train\PYGZus{}final}\PYG{p}{,} \PYG{n}{y\PYGZus{}train\PYGZus{}final}\PYG{p}{)}
    \PYG{n}{y\PYGZus{}pred\PYGZus{}rf} \PYG{o}{=} \PYG{n}{rf}\PYG{o}{.}\PYG{n}{predict}\PYG{p}{(}\PYG{n}{X\PYGZus{}test\PYGZus{}final}\PYG{p}{)}
    \PYG{n}{mse\PYGZus{}values\PYGZus{}rf}\PYG{o}{.}\PYG{n}{append}\PYG{p}{(}\PYG{n}{mean\PYGZus{}squared\PYGZus{}error}\PYG{p}{(}\PYG{n}{y\PYGZus{}test\PYGZus{}final}\PYG{p}{,} \PYG{n}{y\PYGZus{}pred\PYGZus{}rf}\PYG{p}{)}\PYG{p}{)}

\PYG{n}{plt}\PYG{o}{.}\PYG{n}{figure}\PYG{p}{(}\PYG{n}{figsize}\PYG{o}{=}\PYG{p}{(}\PYG{l+m+mi}{10}\PYG{p}{,} \PYG{l+m+mi}{6}\PYG{p}{)}\PYG{p}{)}
\PYG{n}{plt}\PYG{o}{.}\PYG{n}{plot}\PYG{p}{(}\PYG{n}{estimators\PYGZus{}range}\PYG{p}{,} \PYG{n}{mse\PYGZus{}values\PYGZus{}rf}\PYG{p}{,} \PYG{n}{marker}\PYG{o}{=}\PYG{l+s+s1}{\PYGZsq{}}\PYG{l+s+s1}{o}\PYG{l+s+s1}{\PYGZsq{}}\PYG{p}{,} \PYG{n}{color}\PYG{o}{=}\PYG{l+s+s1}{\PYGZsq{}}\PYG{l+s+s1}{green}\PYG{l+s+s1}{\PYGZsq{}}\PYG{p}{)}
\PYG{n}{plt}\PYG{o}{.}\PYG{n}{title}\PYG{p}{(}\PYG{l+s+s2}{\PYGZdq{}}\PYG{l+s+s2}{MSE vs n\PYGZus{}estimators (Random Forest)}\PYG{l+s+s2}{\PYGZdq{}}\PYG{p}{)}
\PYG{n}{plt}\PYG{o}{.}\PYG{n}{xlabel}\PYG{p}{(}\PYG{l+s+s2}{\PYGZdq{}}\PYG{l+s+s2}{n\PYGZus{}estimators}\PYG{l+s+s2}{\PYGZdq{}}\PYG{p}{)}
\PYG{n}{plt}\PYG{o}{.}\PYG{n}{ylabel}\PYG{p}{(}\PYG{l+s+s2}{\PYGZdq{}}\PYG{l+s+s2}{MSE}\PYG{l+s+s2}{\PYGZdq{}}\PYG{p}{)}
\PYG{n}{plt}\PYG{o}{.}\PYG{n}{grid}\PYG{p}{(}\PYG{l+s+s2}{\PYGZdq{}}\PYG{l+s+s2}{on}\PYG{l+s+s2}{\PYGZdq{}}\PYG{p}{)}
\PYG{n}{plt}\PYG{o}{.}\PYG{n}{show}\PYG{p}{(}\PYG{p}{)}
\end{sphinxVerbatim}

\end{sphinxuseclass}\end{sphinxVerbatimInput}
\begin{sphinxVerbatimOutput}

\begin{sphinxuseclass}{cell_output}
\begin{sphinxVerbatim}[commandchars=\\\{\}]
  0\PYGZpc{}|          | 0/39 [00:00\PYGZlt{}?, ?it/s]
\end{sphinxVerbatim}

\noindent\sphinxincludegraphics{{33c88984e5a48bf8d4b4db84ee4276024270e5dbf0f860bc95680d6ee2d155b2}.png}

\end{sphinxuseclass}\end{sphinxVerbatimOutput}

\end{sphinxuseclass}

\subsubsection{Feature importance with the best models}
\label{\detokenize{03:feature-importance-with-the-best-models}}
\sphinxAtStartPar
Based on the experiments above, let’s define user\sphinxhyphen{}chosen parameters for the \sphinxstyleemphasis{best} (or at least improved) Decision Tree and Random Forest, then compute and display their feature importances.

\begin{sphinxShadowBox}
\sphinxstylesidebartitle{}

\sphinxAtStartPar
Feel free to update \sphinxcode{\sphinxupquote{max\_depth}} and \sphinxcode{\sphinxupquote{n\_estimators}} according to the results you obtained above. Later, we will see how this can be used as a \sphinxstyleemphasis{feature selection} method.
\end{sphinxShadowBox}

\begin{sphinxuseclass}{cell}\begin{sphinxVerbatimInput}

\begin{sphinxuseclass}{cell_input}
\begin{sphinxVerbatim}[commandchars=\\\{\}]
\PYG{n}{max\PYGZus{}depth} \PYG{o}{=} \PYG{l+m+mi}{9}
\PYG{n}{n\PYGZus{}estimators} \PYG{o}{=} \PYG{l+m+mi}{60} 

\PYG{n}{best\PYGZus{}tree} \PYG{o}{=} \PYG{n}{DecisionTreeRegressor}\PYG{p}{(}\PYG{n}{max\PYGZus{}depth}\PYG{o}{=}\PYG{n}{max\PYGZus{}depth}\PYG{p}{,} \PYG{n}{random\PYGZus{}state}\PYG{o}{=}\PYG{l+m+mi}{42}\PYG{p}{)}
\PYG{n}{best\PYGZus{}tree}\PYG{o}{.}\PYG{n}{fit}\PYG{p}{(}\PYG{n}{X\PYGZus{}train\PYGZus{}final}\PYG{p}{,} \PYG{n}{y\PYGZus{}train\PYGZus{}final}\PYG{p}{)}

\PYG{n}{best\PYGZus{}forest} \PYG{o}{=} \PYG{n}{RandomForestRegressor}\PYG{p}{(}\PYG{n}{n\PYGZus{}estimators}\PYG{o}{=}\PYG{n}{n\PYGZus{}estimators}\PYG{p}{,} \PYG{n}{random\PYGZus{}state}\PYG{o}{=}\PYG{l+m+mi}{42}\PYG{p}{)}
\PYG{n}{best\PYGZus{}forest}\PYG{o}{.}\PYG{n}{fit}\PYG{p}{(}\PYG{n}{X\PYGZus{}train\PYGZus{}final}\PYG{p}{,} \PYG{n}{y\PYGZus{}train\PYGZus{}final}\PYG{p}{)}

\PYG{n}{feature\PYGZus{}names} \PYG{o}{=} \PYG{n}{X\PYGZus{}original}\PYG{o}{.}\PYG{n}{columns}

\PYG{n}{tree\PYGZus{}importances} \PYG{o}{=} \PYG{n}{best\PYGZus{}tree}\PYG{o}{.}\PYG{n}{feature\PYGZus{}importances\PYGZus{}}
\PYG{n}{forest\PYGZus{}importances} \PYG{o}{=} \PYG{n}{best\PYGZus{}forest}\PYG{o}{.}\PYG{n}{feature\PYGZus{}importances\PYGZus{}}

\PYG{n}{tree\PYGZus{}sorted\PYGZus{}idx} \PYG{o}{=} \PYG{n}{np}\PYG{o}{.}\PYG{n}{argsort}\PYG{p}{(}\PYG{n}{tree\PYGZus{}importances}\PYG{p}{)}\PYG{p}{[}\PYG{p}{:}\PYG{p}{:}\PYG{o}{\PYGZhy{}}\PYG{l+m+mi}{1}\PYG{p}{]}
\PYG{n}{forest\PYGZus{}sorted\PYGZus{}idx} \PYG{o}{=} \PYG{n}{np}\PYG{o}{.}\PYG{n}{argsort}\PYG{p}{(}\PYG{n}{forest\PYGZus{}importances}\PYG{p}{)}\PYG{p}{[}\PYG{p}{:}\PYG{p}{:}\PYG{o}{\PYGZhy{}}\PYG{l+m+mi}{1}\PYG{p}{]}

\PYG{n}{plt}\PYG{o}{.}\PYG{n}{figure}\PYG{p}{(}\PYG{n}{figsize}\PYG{o}{=}\PYG{p}{(}\PYG{l+m+mi}{12}\PYG{p}{,} \PYG{l+m+mi}{6}\PYG{p}{)}\PYG{p}{)}
\PYG{n}{plt}\PYG{o}{.}\PYG{n}{subplot}\PYG{p}{(}\PYG{l+m+mi}{1}\PYG{p}{,} \PYG{l+m+mi}{2}\PYG{p}{,} \PYG{l+m+mi}{1}\PYG{p}{)}
\PYG{n}{plt}\PYG{o}{.}\PYG{n}{title}\PYG{p}{(}\PYG{l+s+s2}{\PYGZdq{}}\PYG{l+s+s2}{Decision Tree Feature Importances}\PYG{l+s+s2}{\PYGZdq{}}\PYG{p}{)}
\PYG{n}{plt}\PYG{o}{.}\PYG{n}{bar}\PYG{p}{(}\PYG{n+nb}{range}\PYG{p}{(}\PYG{n+nb}{len}\PYG{p}{(}\PYG{n}{tree\PYGZus{}importances}\PYG{p}{)}\PYG{p}{)}\PYG{p}{,} \PYG{n}{tree\PYGZus{}importances}\PYG{p}{[}\PYG{n}{tree\PYGZus{}sorted\PYGZus{}idx}\PYG{p}{]}\PYG{p}{,} \PYG{n}{align}\PYG{o}{=}\PYG{l+s+s1}{\PYGZsq{}}\PYG{l+s+s1}{center}\PYG{l+s+s1}{\PYGZsq{}}\PYG{p}{)}
\PYG{n}{plt}\PYG{o}{.}\PYG{n}{xticks}\PYG{p}{(}\PYG{n+nb}{range}\PYG{p}{(}\PYG{n+nb}{len}\PYG{p}{(}\PYG{n}{tree\PYGZus{}importances}\PYG{p}{)}\PYG{p}{)}\PYG{p}{,} \PYG{n}{feature\PYGZus{}names}\PYG{p}{[}\PYG{n}{tree\PYGZus{}sorted\PYGZus{}idx}\PYG{p}{]}\PYG{p}{,} \PYG{n}{rotation}\PYG{o}{=}\PYG{l+m+mi}{45}\PYG{p}{)}
\PYG{n}{plt}\PYG{o}{.}\PYG{n}{grid}\PYG{p}{(}\PYG{l+s+s2}{\PYGZdq{}}\PYG{l+s+s2}{on}\PYG{l+s+s2}{\PYGZdq{}}\PYG{p}{)}
\PYG{n}{plt}\PYG{o}{.}\PYG{n}{subplot}\PYG{p}{(}\PYG{l+m+mi}{1}\PYG{p}{,} \PYG{l+m+mi}{2}\PYG{p}{,} \PYG{l+m+mi}{2}\PYG{p}{)}
\PYG{n}{plt}\PYG{o}{.}\PYG{n}{title}\PYG{p}{(}\PYG{l+s+s2}{\PYGZdq{}}\PYG{l+s+s2}{Random Forest Feature Importances}\PYG{l+s+s2}{\PYGZdq{}}\PYG{p}{)}
\PYG{n}{plt}\PYG{o}{.}\PYG{n}{bar}\PYG{p}{(}\PYG{n+nb}{range}\PYG{p}{(}\PYG{n+nb}{len}\PYG{p}{(}\PYG{n}{forest\PYGZus{}importances}\PYG{p}{)}\PYG{p}{)}\PYG{p}{,} \PYG{n}{forest\PYGZus{}importances}\PYG{p}{[}\PYG{n}{forest\PYGZus{}sorted\PYGZus{}idx}\PYG{p}{]}\PYG{p}{,} \PYG{n}{align}\PYG{o}{=}\PYG{l+s+s1}{\PYGZsq{}}\PYG{l+s+s1}{center}\PYG{l+s+s1}{\PYGZsq{}}\PYG{p}{,} \PYG{n}{color}\PYG{o}{=}\PYG{l+s+s1}{\PYGZsq{}}\PYG{l+s+s1}{green}\PYG{l+s+s1}{\PYGZsq{}}\PYG{p}{)}
\PYG{n}{plt}\PYG{o}{.}\PYG{n}{xticks}\PYG{p}{(}\PYG{n+nb}{range}\PYG{p}{(}\PYG{n+nb}{len}\PYG{p}{(}\PYG{n}{forest\PYGZus{}importances}\PYG{p}{)}\PYG{p}{)}\PYG{p}{,} \PYG{n}{feature\PYGZus{}names}\PYG{p}{[}\PYG{n}{forest\PYGZus{}sorted\PYGZus{}idx}\PYG{p}{]}\PYG{p}{,} \PYG{n}{rotation}\PYG{o}{=}\PYG{l+m+mi}{45}\PYG{p}{)}
\PYG{n}{plt}\PYG{o}{.}\PYG{n}{grid}\PYG{p}{(}\PYG{l+s+s2}{\PYGZdq{}}\PYG{l+s+s2}{on}\PYG{l+s+s2}{\PYGZdq{}}\PYG{p}{)}
\PYG{n}{plt}\PYG{o}{.}\PYG{n}{tight\PYGZus{}layout}\PYG{p}{(}\PYG{p}{)}
\PYG{n}{plt}\PYG{o}{.}\PYG{n}{show}\PYG{p}{(}\PYG{p}{)}
\end{sphinxVerbatim}

\end{sphinxuseclass}\end{sphinxVerbatimInput}
\begin{sphinxVerbatimOutput}

\begin{sphinxuseclass}{cell_output}
\noindent\sphinxincludegraphics{{9842af6cdf010e1e4c19acbccb07afac23fd30e6fe0ca1ccf93374efe72d64c2}.png}

\end{sphinxuseclass}\end{sphinxVerbatimOutput}

\end{sphinxuseclass}

\subsection{Building a Pipeline}
\label{\detokenize{03:building-a-pipeline}}\begin{enumerate}
\sphinxsetlistlabels{\arabic}{enumi}{enumii}{}{.}%
\item {} 
\sphinxAtStartPar
Define a pipeline that applies the preprocessing steps (e.g., imputation, scaling) and the best\sphinxhyphen{}performing model (from Exercise 3).

\item {} 
\sphinxAtStartPar
Fit this pipeline on the training data and evaluate on the test data.

\item {} 
\sphinxAtStartPar
Observe how \sphinxcode{\sphinxupquote{scikit\sphinxhyphen{}learn}} handles transformations during \sphinxcode{\sphinxupquote{fit}} and \sphinxcode{\sphinxupquote{predict}}.

\end{enumerate}


\subsubsection{Solution}
\label{\detokenize{03:id5}}
\sphinxAtStartPar
Below, we define a pipeline that:
\begin{enumerate}
\sphinxsetlistlabels{\arabic}{enumi}{enumii}{}{.}%
\item {} 
\sphinxAtStartPar
Imputes missing values with mean.

\item {} 
\sphinxAtStartPar
Scales the features.

\item {} 
\sphinxAtStartPar
Trains a random forest regressor.

\end{enumerate}

\begin{sphinxuseclass}{cell}\begin{sphinxVerbatimInput}

\begin{sphinxuseclass}{cell_input}
\begin{sphinxVerbatim}[commandchars=\\\{\}]
\PYG{k+kn}{from}\PYG{+w}{ }\PYG{n+nn}{sklearn}\PYG{n+nn}{.}\PYG{n+nn}{pipeline}\PYG{+w}{ }\PYG{k+kn}{import} \PYG{n}{Pipeline}
\PYG{k+kn}{from}\PYG{+w}{ }\PYG{n+nn}{sklearn}\PYG{n+nn}{.}\PYG{n+nn}{impute}\PYG{+w}{ }\PYG{k+kn}{import} \PYG{n}{SimpleImputer}
\PYG{k+kn}{from}\PYG{+w}{ }\PYG{n+nn}{sklearn}\PYG{n+nn}{.}\PYG{n+nn}{preprocessing}\PYG{+w}{ }\PYG{k+kn}{import} \PYG{n}{StandardScaler}
\PYG{k+kn}{from}\PYG{+w}{ }\PYG{n+nn}{sklearn}\PYG{n+nn}{.}\PYG{n+nn}{metrics}\PYG{+w}{ }\PYG{k+kn}{import} \PYG{n}{mean\PYGZus{}squared\PYGZus{}error}

\PYG{n}{pipeline\PYGZus{}rf} \PYG{o}{=} \PYG{n}{Pipeline}\PYG{p}{(}\PYG{p}{[}
    \PYG{p}{(}\PYG{l+s+s2}{\PYGZdq{}}\PYG{l+s+s2}{imputer}\PYG{l+s+s2}{\PYGZdq{}}\PYG{p}{,} \PYG{n}{SimpleImputer}\PYG{p}{(}\PYG{n}{strategy}\PYG{o}{=}\PYG{l+s+s2}{\PYGZdq{}}\PYG{l+s+s2}{mean}\PYG{l+s+s2}{\PYGZdq{}}\PYG{p}{)}\PYG{p}{)}\PYG{p}{,}
    \PYG{p}{(}\PYG{l+s+s2}{\PYGZdq{}}\PYG{l+s+s2}{scaler}\PYG{l+s+s2}{\PYGZdq{}}\PYG{p}{,} \PYG{n}{StandardScaler}\PYG{p}{(}\PYG{p}{)}\PYG{p}{)}\PYG{p}{,}
    \PYG{p}{(}\PYG{l+s+s2}{\PYGZdq{}}\PYG{l+s+s2}{regressor}\PYG{l+s+s2}{\PYGZdq{}}\PYG{p}{,} \PYG{n}{RandomForestRegressor}\PYG{p}{(}\PYG{n}{n\PYGZus{}estimators}\PYG{o}{=}\PYG{l+m+mi}{60}\PYG{p}{,} \PYG{n}{random\PYGZus{}state}\PYG{o}{=}\PYG{l+m+mi}{42}\PYG{p}{)}\PYG{p}{)}
\PYG{p}{]}\PYG{p}{)}

\PYG{n}{pipeline\PYGZus{}rf}\PYG{o}{.}\PYG{n}{fit}\PYG{p}{(}\PYG{n}{X\PYGZus{}train\PYGZus{}full}\PYG{p}{,} \PYG{n}{y\PYGZus{}train\PYGZus{}full}\PYG{p}{)}
\PYG{n}{y\PYGZus{}pred\PYGZus{}pipe} \PYG{o}{=} \PYG{n}{pipeline\PYGZus{}rf}\PYG{o}{.}\PYG{n}{predict}\PYG{p}{(}\PYG{n}{X\PYGZus{}test\PYGZus{}full}\PYG{p}{)}
\PYG{n}{mse\PYGZus{}pipeline\PYGZus{}rf} \PYG{o}{=} \PYG{n}{mean\PYGZus{}squared\PYGZus{}error}\PYG{p}{(}\PYG{n}{y\PYGZus{}test\PYGZus{}full}\PYG{p}{,} \PYG{n}{y\PYGZus{}pred\PYGZus{}pipe}\PYG{p}{)}
\PYG{n+nb}{print}\PYG{p}{(}\PYG{l+s+sa}{f}\PYG{l+s+s2}{\PYGZdq{}}\PYG{l+s+s2}{Pipeline (Random Forest) MSE: }\PYG{l+s+si}{\PYGZob{}}\PYG{n}{mse\PYGZus{}pipeline\PYGZus{}rf}\PYG{l+s+si}{:}\PYG{l+s+s2}{.4f}\PYG{l+s+si}{\PYGZcb{}}\PYG{l+s+s2}{\PYGZdq{}}\PYG{p}{)}
\end{sphinxVerbatim}

\end{sphinxuseclass}\end{sphinxVerbatimInput}
\begin{sphinxVerbatimOutput}

\begin{sphinxuseclass}{cell_output}
\begin{sphinxVerbatim}[commandchars=\\\{\}]
Pipeline (Random Forest) MSE: 0.2693
\end{sphinxVerbatim}

\end{sphinxuseclass}\end{sphinxVerbatimOutput}

\end{sphinxuseclass}

\subsection{Additional exercise: bias–variance trade\sphinxhyphen{}off in regression with decision trees}
\label{\detokenize{03:additional-exercise-biasvariance-trade-off-in-regression-with-decision-trees}}
\sphinxAtStartPar
In this exercise, you will investigate the bias–variance trade\sphinxhyphen{}off for decision tree regressors of varying complexity. Specifically, you will perform simulations to estimate how the depth of a decision tree influences its performance.
\begin{enumerate}
\sphinxsetlistlabels{\arabic}{enumi}{enumii}{}{.}%
\item {} 
\sphinxAtStartPar
Select a nonlinear, moderately complex target function e.g.:

\end{enumerate}
\begin{equation*}
\begin{split}f(x) = \sin(4\pi x) + w\end{split}
\end{equation*}
\sphinxAtStartPar
or
\begin{equation*}
\begin{split}f(x) = \sin(2\pi x_1 x_2 x_3) + w\end{split}
\end{equation*}\begin{enumerate}
\sphinxsetlistlabels{\arabic}{enumi}{enumii}{}{.}%
\setcounter{enumi}{1}
\item {} 
\sphinxAtStartPar
Evaluate three decision tree regressors, each with a different complexity level:
\begin{itemize}
\item {} 
\sphinxAtStartPar
Shallow tree: depth = 2

\item {} 
\sphinxAtStartPar
Medium\sphinxhyphen{}depth tree: depth = 6

\item {} 
\sphinxAtStartPar
Deep tree: depth = 20

\end{itemize}

\item {} 
\sphinxAtStartPar
Generate training datasets composed of \(n_{\text{train}} = 400\) random samples:

\end{enumerate}
\begin{equation*}
\begin{split}
X_{\text{train}} \sim U(0,1), \quad y_{\text{train}} = f(X_{\text{train}}) + w
\end{split}
\end{equation*}
\sphinxAtStartPar
Here, noise \(w\) follows a Gaussian distribution with mean 0 and standard deviation \(0.3\):
\begin{equation*}
\begin{split}
w \sim \mathcal{N}(0, 0.3^2)
\end{split}
\end{equation*}\begin{enumerate}
\sphinxsetlistlabels{\arabic}{enumi}{enumii}{}{.}%
\setcounter{enumi}{4}
\item {} 
\sphinxAtStartPar
Repeat the entire training and evaluation process \(n_{\text{runs}} = 100\) times. For each run, compute and store the mean squared error (MSE) using 10\sphinxhyphen{}fold cross\sphinxhyphen{}validation on the training set.

\item {} 
\sphinxAtStartPar
Evaluate your models on a test set \(X_{\text{test}}\) evenly spaced in the interval \([0,1]\) (e.g., 400 points).
\begin{itemize}
\item {} 
\sphinxAtStartPar
For each test point, compute predictions across all 100 simulation runs.

\item {} 
\sphinxAtStartPar
Using these predictions, compute
\begin{itemize}
\item {} 
\sphinxAtStartPar
MSE (CV): The average cross\sphinxhyphen{}validation mean squared error (which we will call \(\widehat{\text{MISE}}_{\text{kfold}}\)) obtained across the 100 runs.

\item {} 
\sphinxAtStartPar
Bias\(^2\): Based on how far the \sphinxstyleemphasis{average prediction} per sample is from the sample’s true label.

\item {} 
\sphinxAtStartPar
Variance: Based on the variability of predictions per sample around that sample’s mean prediction.

\item {} 
\sphinxAtStartPar
MSE (true): the approximation \(
\text{MSE} \approx \text{Bias}^2 + \text{Variance}.\)

\end{itemize}

\end{itemize}

\item {} 
\sphinxAtStartPar
Discussion\\
Analyze your results to illustrate how model complexity (tree depth) affects the balance between bias and variance:
\begin{itemize}
\item {} 
\sphinxAtStartPar
Identify which depth leads to underfitting (high bias, low variance).

\item {} 
\sphinxAtStartPar
Identify which depth leads to overfitting (low bias, high variance).

\end{itemize}

\end{enumerate}


\subsubsection{Solution}
\label{\detokenize{03:id6}}
\begin{sphinxuseclass}{cell}\begin{sphinxVerbatimInput}

\begin{sphinxuseclass}{cell_input}
\begin{sphinxVerbatim}[commandchars=\\\{\}]
\PYG{k+kn}{import}\PYG{+w}{ }\PYG{n+nn}{pandas}\PYG{+w}{ }\PYG{k}{as}\PYG{+w}{ }\PYG{n+nn}{pd}
\PYG{k+kn}{from}\PYG{+w}{ }\PYG{n+nn}{sklearn}\PYG{n+nn}{.}\PYG{n+nn}{model\PYGZus{}selection}\PYG{+w}{ }\PYG{k+kn}{import} \PYG{n}{cross\PYGZus{}val\PYGZus{}score}\PYG{p}{,} \PYG{n}{KFold}

\PYG{n}{n\PYGZus{}train} \PYG{o}{=} \PYG{l+m+mi}{400}
\PYG{n}{n\PYGZus{}runs} \PYG{o}{=} \PYG{l+m+mi}{100}
\PYG{n}{noise\PYGZus{}std} \PYG{o}{=} \PYG{l+m+mf}{0.3}
\PYG{n}{depths} \PYG{o}{=} \PYG{p}{[}\PYG{l+m+mi}{2}\PYG{p}{,} \PYG{l+m+mi}{6}\PYG{p}{,} \PYG{l+m+mi}{20}\PYG{p}{]}
\PYG{n}{K} \PYG{o}{=} \PYG{l+m+mi}{10}

\PYG{k}{def}\PYG{+w}{ }\PYG{n+nf}{f\PYGZus{}func}\PYG{p}{(}\PYG{n}{x}\PYG{p}{)}\PYG{p}{:}
    \PYG{k}{return} \PYG{n}{np}\PYG{o}{.}\PYG{n}{sin}\PYG{p}{(}\PYG{l+m+mi}{4} \PYG{o}{*} \PYG{n}{np}\PYG{o}{.}\PYG{n}{pi} \PYG{o}{*} \PYG{n}{x}\PYG{p}{)}\PYG{o}{.}\PYG{n}{ravel}\PYG{p}{(}\PYG{p}{)}

\PYG{n}{X\PYGZus{}test} \PYG{o}{=} \PYG{n}{np}\PYG{o}{.}\PYG{n}{linspace}\PYG{p}{(}\PYG{l+m+mi}{0}\PYG{p}{,} \PYG{l+m+mi}{1}\PYG{p}{,} \PYG{l+m+mi}{400}\PYG{p}{)}\PYG{p}{[}\PYG{p}{:}\PYG{p}{,} \PYG{k+kc}{None}\PYG{p}{]}
\PYG{n}{y\PYGZus{}true} \PYG{o}{=} \PYG{n}{f\PYGZus{}func}\PYG{p}{(}\PYG{n}{X\PYGZus{}test}\PYG{p}{)}

\PYG{n}{predictions} \PYG{o}{=} \PYG{p}{\PYGZob{}}\PYG{n}{depth}\PYG{p}{:} \PYG{n}{np}\PYG{o}{.}\PYG{n}{zeros}\PYG{p}{(}\PYG{p}{(}\PYG{n}{n\PYGZus{}runs}\PYG{p}{,} \PYG{n}{X\PYGZus{}test}\PYG{o}{.}\PYG{n}{shape}\PYG{p}{[}\PYG{l+m+mi}{0}\PYG{p}{]}\PYG{p}{)}\PYG{p}{)} \PYG{k}{for} \PYG{n}{depth} \PYG{o+ow}{in} \PYG{n}{depths}\PYG{p}{\PYGZcb{}}
\PYG{n}{cv\PYGZus{}mse\PYGZus{}values} \PYG{o}{=} \PYG{p}{\PYGZob{}}\PYG{n}{depth}\PYG{p}{:} \PYG{p}{[}\PYG{p}{]} \PYG{k}{for} \PYG{n}{depth} \PYG{o+ow}{in} \PYG{n}{depths}\PYG{p}{\PYGZcb{}}

\PYG{k}{for} \PYG{n}{i} \PYG{o+ow}{in} \PYG{n+nb}{range}\PYG{p}{(}\PYG{n}{n\PYGZus{}runs}\PYG{p}{)}\PYG{p}{:}
    \PYG{n}{X\PYGZus{}train} \PYG{o}{=} \PYG{n}{np}\PYG{o}{.}\PYG{n}{random}\PYG{o}{.}\PYG{n}{rand}\PYG{p}{(}\PYG{n}{n\PYGZus{}train}\PYG{p}{,} \PYG{l+m+mi}{1}\PYG{p}{)}
    \PYG{n}{y\PYGZus{}train} \PYG{o}{=} \PYG{n}{f\PYGZus{}func}\PYG{p}{(}\PYG{n}{X\PYGZus{}train}\PYG{p}{)} \PYG{o}{+} \PYG{n}{np}\PYG{o}{.}\PYG{n}{random}\PYG{o}{.}\PYG{n}{normal}\PYG{p}{(}\PYG{l+m+mi}{0}\PYG{p}{,} \PYG{n}{noise\PYGZus{}std}\PYG{p}{,} \PYG{n}{size}\PYG{o}{=}\PYG{n}{n\PYGZus{}train}\PYG{p}{)}

    \PYG{k}{for} \PYG{n}{depth} \PYG{o+ow}{in} \PYG{n}{depths}\PYG{p}{:}
        \PYG{n}{model} \PYG{o}{=} \PYG{n}{DecisionTreeRegressor}\PYG{p}{(}\PYG{n}{max\PYGZus{}depth}\PYG{o}{=}\PYG{n}{depth}\PYG{p}{)}
        \PYG{n}{model}\PYG{o}{.}\PYG{n}{fit}\PYG{p}{(}\PYG{n}{X\PYGZus{}train}\PYG{p}{,} \PYG{n}{y\PYGZus{}train}\PYG{p}{)}

        \PYG{n}{predictions}\PYG{p}{[}\PYG{n}{depth}\PYG{p}{]}\PYG{p}{[}\PYG{n}{i}\PYG{p}{,} \PYG{p}{:}\PYG{p}{]} \PYG{o}{=} \PYG{n}{model}\PYG{o}{.}\PYG{n}{predict}\PYG{p}{(}\PYG{n}{X\PYGZus{}test}\PYG{p}{)}
        
        \PYG{n}{cv\PYGZus{}scores} \PYG{o}{=} \PYG{n}{cross\PYGZus{}val\PYGZus{}score}\PYG{p}{(}
            \PYG{n}{DecisionTreeRegressor}\PYG{p}{(}\PYG{n}{max\PYGZus{}depth}\PYG{o}{=}\PYG{n}{depth}\PYG{p}{)}\PYG{p}{,}
            \PYG{n}{X\PYGZus{}train}\PYG{p}{,} \PYG{n}{y\PYGZus{}train}\PYG{p}{,}
            \PYG{n}{scoring}\PYG{o}{=}\PYG{l+s+s1}{\PYGZsq{}}\PYG{l+s+s1}{neg\PYGZus{}mean\PYGZus{}squared\PYGZus{}error}\PYG{l+s+s1}{\PYGZsq{}}\PYG{p}{,}
            \PYG{n}{cv}\PYG{o}{=}\PYG{n}{KFold}\PYG{p}{(}\PYG{n}{n\PYGZus{}splits}\PYG{o}{=}\PYG{n}{K}\PYG{p}{,} \PYG{n}{shuffle}\PYG{o}{=}\PYG{k+kc}{True}\PYG{p}{,} \PYG{n}{random\PYGZus{}state}\PYG{o}{=}\PYG{l+m+mi}{1}\PYG{p}{)}
        \PYG{p}{)}
        \PYG{n}{cv\PYGZus{}mse} \PYG{o}{=} \PYG{o}{\PYGZhy{}}\PYG{n}{cv\PYGZus{}scores}\PYG{o}{.}\PYG{n}{mean}\PYG{p}{(}\PYG{p}{)}
        \PYG{n}{cv\PYGZus{}mse\PYGZus{}values}\PYG{p}{[}\PYG{n}{depth}\PYG{p}{]}\PYG{o}{.}\PYG{n}{append}\PYG{p}{(}\PYG{n}{cv\PYGZus{}mse}\PYG{p}{)}

\PYG{n}{noise\PYGZus{}var} \PYG{o}{=} \PYG{n}{noise\PYGZus{}std}\PYG{o}{*}\PYG{o}{*}\PYG{l+m+mi}{2}
\PYG{n}{metrics} \PYG{o}{=} \PYG{p}{\PYGZob{}}\PYG{p}{\PYGZcb{}}

\PYG{k}{for} \PYG{n}{depth} \PYG{o+ow}{in} \PYG{n}{depths}\PYG{p}{:}
    \PYG{n}{mean\PYGZus{}pred} \PYG{o}{=} \PYG{n}{np}\PYG{o}{.}\PYG{n}{mean}\PYG{p}{(}\PYG{n}{predictions}\PYG{p}{[}\PYG{n}{depth}\PYG{p}{]}\PYG{p}{,} \PYG{n}{axis}\PYG{o}{=}\PYG{l+m+mi}{0}\PYG{p}{)}
    \PYG{n}{bias\PYGZus{}sq} \PYG{o}{=} \PYG{n}{np}\PYG{o}{.}\PYG{n}{mean}\PYG{p}{(}\PYG{p}{(}\PYG{n}{mean\PYGZus{}pred} \PYG{o}{\PYGZhy{}} \PYG{n}{y\PYGZus{}true}\PYG{p}{)}\PYG{o}{*}\PYG{o}{*}\PYG{l+m+mi}{2}\PYG{p}{)}
    \PYG{n}{variance} \PYG{o}{=} \PYG{n}{np}\PYG{o}{.}\PYG{n}{mean}\PYG{p}{(}\PYG{n}{np}\PYG{o}{.}\PYG{n}{var}\PYG{p}{(}\PYG{n}{predictions}\PYG{p}{[}\PYG{n}{depth}\PYG{p}{]}\PYG{p}{,} \PYG{n}{axis}\PYG{o}{=}\PYG{l+m+mi}{0}\PYG{p}{)}\PYG{p}{)}
    \PYG{n}{true\PYGZus{}mse} \PYG{o}{=} \PYG{n}{bias\PYGZus{}sq} \PYG{o}{+} \PYG{n}{variance} \PYG{o}{+} \PYG{n}{noise\PYGZus{}var}
    \PYG{n}{cv\PYGZus{}mse\PYGZus{}mean} \PYG{o}{=} \PYG{n}{np}\PYG{o}{.}\PYG{n}{mean}\PYG{p}{(}\PYG{n}{cv\PYGZus{}mse\PYGZus{}values}\PYG{p}{[}\PYG{n}{depth}\PYG{p}{]}\PYG{p}{)}

    \PYG{n}{metrics}\PYG{p}{[}\PYG{n}{depth}\PYG{p}{]} \PYG{o}{=} \PYG{p}{\PYGZob{}}
        \PYG{l+s+s2}{\PYGZdq{}}\PYG{l+s+s2}{Bias\PYGZca{}2}\PYG{l+s+s2}{\PYGZdq{}}\PYG{p}{:} \PYG{n}{bias\PYGZus{}sq}\PYG{p}{,}
        \PYG{l+s+s2}{\PYGZdq{}}\PYG{l+s+s2}{Variance}\PYG{l+s+s2}{\PYGZdq{}}\PYG{p}{:} \PYG{n}{variance}\PYG{p}{,}
        \PYG{l+s+s2}{\PYGZdq{}}\PYG{l+s+s2}{True MSE}\PYG{l+s+s2}{\PYGZdq{}}\PYG{p}{:} \PYG{n}{true\PYGZus{}mse}\PYG{p}{,}
        \PYG{l+s+s2}{\PYGZdq{}}\PYG{l+s+s2}{CV MSE}\PYG{l+s+s2}{\PYGZdq{}}\PYG{p}{:} \PYG{n}{cv\PYGZus{}mse\PYGZus{}mean}
    \PYG{p}{\PYGZcb{}}

\PYG{n}{metrics\PYGZus{}df} \PYG{o}{=} \PYG{n}{pd}\PYG{o}{.}\PYG{n}{DataFrame}\PYG{p}{(}\PYG{n}{metrics}\PYG{p}{)}\PYG{o}{.}\PYG{n}{T}
\PYG{n}{metrics\PYGZus{}df}\PYG{o}{.}\PYG{n}{index}\PYG{o}{.}\PYG{n}{name} \PYG{o}{=} \PYG{l+s+s2}{\PYGZdq{}}\PYG{l+s+s2}{Tree Depth}\PYG{l+s+s2}{\PYGZdq{}}
\PYG{n}{metrics\PYGZus{}df} \PYG{o}{=} \PYG{n}{metrics\PYGZus{}df}\PYG{o}{.}\PYG{n}{round}\PYG{p}{(}\PYG{l+m+mi}{4}\PYG{p}{)}
\PYG{n+nb}{print}\PYG{p}{(}\PYG{n}{metrics\PYGZus{}df}\PYG{p}{)}
\end{sphinxVerbatim}

\end{sphinxuseclass}\end{sphinxVerbatimInput}
\begin{sphinxVerbatimOutput}

\begin{sphinxuseclass}{cell_output}
\begin{sphinxVerbatim}[commandchars=\\\{\}]
            Bias\PYGZca{}2  Variance  True MSE  CV MSE
Tree Depth                                    
2           0.1781    0.1099    0.3780  0.3800
6           0.0010    0.0315    0.1225  0.1242
20          0.0009    0.0883    0.1792  0.1788
\end{sphinxVerbatim}

\end{sphinxuseclass}\end{sphinxVerbatimOutput}

\end{sphinxuseclass}
\sphinxstepscope


\chapter{Neural Networks for Regression}
\label{\detokenize{04:neural-networks-for-regression}}\label{\detokenize{04::doc}}

\section{Introduction}
\label{\detokenize{04:introduction}}
\sphinxAtStartPar
This practical has not as intention to reintroduce entirely the notion of Neural Network. We redirect the student to Chapter 8 of \sphinxhref{http://researchgate.net/publication/242692234\_Statistical\_foundations\_of\_machine\_learning\_the\_handbook}{the handbook} for a formal introduction to those tools.

\sphinxAtStartPar
However, we aim at giving a practical sight on what Neural Networks actually are. Concretely, Neural Networks form a family of approximations of functions. Generally speaking, most Neural Networks make the hypothesis of non\sphinxhyphen{}linearity in the relationship between the inputs and the outputs. The design of the Neural Network involves, non\sphinxhyphen{}exhaustively, the definition of the input and output shape, the size and the relationships between the hidden layers, and the non\sphinxhyphen{}linear activation functions for each of the layers.

\begin{sphinxShadowBox}
\sphinxstylesidebartitle{}

\sphinxAtStartPar
One main difference between \sphinxstyleemphasis{tensors} and \sphinxstyleemphasis{matrices}, mathematically speaking, is the fact that the tensors do not depend on the chosen basis, while matrices do. Matrices can serve as good representations of tensors. Note that tensors can be seen members of vector spaces. However, this is not important for the covered subject, nor for the practical tensorial algebra that is used in Neural Networks design and optimization.
\end{sphinxShadowBox}

\sphinxAtStartPar
In most modern Neural Networks frameworks, the design of a Neural Network also implies the definition of the optimization algorithm and the definition of some hyperparameters. In most of the existing libraries, the data needs to be transformed to \sphinxstyleemphasis{tensors}, which, despite the fact that they have another mathematical definition, behave really closely to the \sphinxcode{\sphinxupquote{NumPy}} arrays that we have presented earlier. Modern computing architectures take into account the usual operations computed on tensors in order to optimize them, which, like it was the case for the arrays and the GPU’s (Graphical Processing Units), gives birth to the TPU’s (Tensor Processing Units).

\sphinxAtStartPar
In this practical, we will present one of the most used framework for Neural Networks design and optimization, \sphinxcode{\sphinxupquote{PyTorch}}, to introduce two famous Neural Network families: the \sphinxstyleemphasis{Multi\sphinxhyphen{}Layers Perceptrons} (MLP), where layers are densely connected with the next one, and \sphinxstyleemphasis{Convolutional Neural Networks} (CNN), which take advantage of spatial (or temporal) proximity of the data’s components.

\begin{sphinxShadowBox}
\sphinxstylesidebartitle{}

\sphinxAtStartPar
As usual, this line can be skipped if the required libraries are already installed on your system.
\end{sphinxShadowBox}

\begin{sphinxuseclass}{cell}\begin{sphinxVerbatimInput}

\begin{sphinxuseclass}{cell_input}
\begin{sphinxVerbatim}[commandchars=\\\{\}]
\PYG{c+ch}{\PYGZsh{}!pip install torch torchsummary}
\end{sphinxVerbatim}

\end{sphinxuseclass}\end{sphinxVerbatimInput}

\end{sphinxuseclass}
\sphinxAtStartPar
Let us define some basic imports and a \sphinxcode{\sphinxupquote{pprint}} function.

\begin{sphinxuseclass}{cell}\begin{sphinxVerbatimInput}

\begin{sphinxuseclass}{cell_input}
\begin{sphinxVerbatim}[commandchars=\\\{\}]
\PYG{k+kn}{import}\PYG{+w}{ }\PYG{n+nn}{numpy}\PYG{+w}{ }\PYG{k}{as}\PYG{+w}{ }\PYG{n+nn}{np}
\PYG{k+kn}{import}\PYG{+w}{ }\PYG{n+nn}{torch}
\PYG{k+kn}{from}\PYG{+w}{ }\PYG{n+nn}{IPython}\PYG{n+nn}{.}\PYG{n+nn}{display}\PYG{+w}{ }\PYG{k+kn}{import} \PYG{n}{display}\PYG{p}{,} \PYG{n}{Math}
\PYG{k+kn}{from}\PYG{+w}{ }\PYG{n+nn}{numpyarray\PYGZus{}to\PYGZus{}latex}\PYG{n+nn}{.}\PYG{n+nn}{jupyter}\PYG{+w}{ }\PYG{k+kn}{import} \PYG{n}{to\PYGZus{}ltx}
\PYG{k+kn}{import}\PYG{+w}{ }\PYG{n+nn}{matplotlib}\PYG{n+nn}{.}\PYG{n+nn}{pyplot}\PYG{+w}{ }\PYG{k}{as}\PYG{+w}{ }\PYG{n+nn}{plt}
\PYG{n}{plt}\PYG{o}{.}\PYG{n}{style}\PYG{o}{.}\PYG{n}{use}\PYG{p}{(}\PYG{p}{[}\PYG{l+s+s1}{\PYGZsq{}}\PYG{l+s+s1}{seaborn\PYGZhy{}v0\PYGZus{}8\PYGZhy{}muted}\PYG{l+s+s1}{\PYGZsq{}}\PYG{p}{,} \PYG{l+s+s1}{\PYGZsq{}}\PYG{l+s+s1}{practicals.mplstyle}\PYG{l+s+s1}{\PYGZsq{}}\PYG{p}{]}\PYG{p}{)}
\PYG{k}{def}\PYG{+w}{ }\PYG{n+nf}{pprint}\PYG{p}{(}\PYG{o}{*}\PYG{n}{args}\PYG{p}{)}\PYG{p}{:}
    \PYG{n}{res} \PYG{o}{=} \PYG{l+s+s2}{\PYGZdq{}}\PYG{l+s+s2}{\PYGZdq{}}
    \PYG{k}{for} \PYG{n}{i} \PYG{o+ow}{in} \PYG{n}{args}\PYG{p}{:}
        \PYG{k}{if} \PYG{n+nb}{type}\PYG{p}{(}\PYG{n}{i}\PYG{p}{)} \PYG{o}{==} \PYG{n}{np}\PYG{o}{.}\PYG{n}{ndarray}\PYG{p}{:}
            \PYG{n}{res} \PYG{o}{+}\PYG{o}{=} \PYG{n}{to\PYGZus{}ltx}\PYG{p}{(}\PYG{n}{i}\PYG{p}{,} \PYG{n}{brackets}\PYG{o}{=}\PYG{l+s+s1}{\PYGZsq{}}\PYG{l+s+s1}{[]}\PYG{l+s+s1}{\PYGZsq{}}\PYG{p}{)}
        \PYG{k}{elif} \PYG{n+nb}{type}\PYG{p}{(}\PYG{n}{i}\PYG{p}{)} \PYG{o}{==} \PYG{n}{torch}\PYG{o}{.}\PYG{n}{Tensor}\PYG{p}{:}
            \PYG{n}{res} \PYG{o}{+}\PYG{o}{=} \PYG{n}{to\PYGZus{}ltx}\PYG{p}{(}\PYG{n}{i}\PYG{o}{.}\PYG{n}{detach}\PYG{p}{(}\PYG{p}{)}\PYG{o}{.}\PYG{n}{cpu}\PYG{p}{(}\PYG{p}{)}\PYG{o}{.}\PYG{n}{numpy}\PYG{p}{(}\PYG{p}{)}\PYG{p}{,} \PYG{n}{brackets}\PYG{o}{=}\PYG{l+s+s1}{\PYGZsq{}}\PYG{l+s+s1}{[]}\PYG{l+s+s1}{\PYGZsq{}}\PYG{p}{,}\PYG{p}{)}
        \PYG{k}{elif} \PYG{n+nb}{type}\PYG{p}{(}\PYG{n}{i}\PYG{p}{)} \PYG{o}{==} \PYG{n+nb}{str}\PYG{p}{:}
            \PYG{n}{res} \PYG{o}{+}\PYG{o}{=} \PYG{n}{i}
    \PYG{n}{display}\PYG{p}{(}\PYG{n}{Math}\PYG{p}{(}\PYG{n}{res}\PYG{p}{)}\PYG{p}{)}
\end{sphinxVerbatim}

\end{sphinxuseclass}\end{sphinxVerbatimInput}

\end{sphinxuseclass}

\subsection{Data Preprocessing for Neural Networks}
\label{\detokenize{04:data-preprocessing-for-neural-networks}}

\subsubsection{Normalization and Standardization}
\label{\detokenize{04:normalization-and-standardization}}
\sphinxAtStartPar
Gradient\sphinxhyphen{}based optimization of neural networks often becomes unstable if different features vary widely in scale. Hence, \sphinxstylestrong{Sec. 8.1} of the handbook (and earlier practicals) reiterate that scaling inputs to zero mean and unit variance—or to a small bounded range—is essential. This ensures the gradients have more balanced magnitudes across parameters, which helps the optimizer converge more reliably. Typically, normalization and standardization follow those steps:
\begin{enumerate}
\sphinxsetlistlabels{\arabic}{enumi}{enumii}{}{.}%
\item {} 
\sphinxAtStartPar
Compute the empirical mean and standard deviation of each input dimension:

\end{enumerate}
\begin{equation*}
\begin{split}
\mu_j = \frac{1}{N} \sum_{i=1}^N x_{ij}, \quad
\sigma_j = \sqrt{\frac{1}{N}\sum_{i=1}^N (x_{ij}-\mu_j)^2}.
\end{split}
\end{equation*}\begin{enumerate}
\sphinxsetlistlabels{\arabic}{enumi}{enumii}{}{.}%
\setcounter{enumi}{1}
\item {} 
\sphinxAtStartPar
Transform (standardise) the data as

\end{enumerate}
\begin{equation*}
\begin{split}
x_{ij}^{(std)} = \frac{x_{ij} - \mu_j}{\sigma_j}.
\end{split}
\end{equation*}\begin{enumerate}
\sphinxsetlistlabels{\arabic}{enumi}{enumii}{}{.}%
\setcounter{enumi}{2}
\item {} 
\sphinxAtStartPar
Use the same \(\mu_j\) and \(\sigma_j\) for the validation and test sets to maintain consistency.

\end{enumerate}


\subsubsection{Input Data Shaping}
\label{\detokenize{04:input-data-shaping}}\begin{itemize}
\item {} 
\sphinxAtStartPar
\sphinxstylestrong{For MLPs}: Usually arrange training data into a matrix of shape \((\text{batch}, n)\), where \(n\) is the number of input features.

\item {} 
\sphinxAtStartPar
\sphinxstylestrong{For 2D\sphinxhyphen{}CNNs}: Reshape data into a 4D tensor of shape \((\text{batch}, \text{channels}, \text{height}, \text{width})\). Convolutional layers exploit the spatial dimensions for images or other grid\sphinxhyphen{}like inputs.

\end{itemize}


\section{Feed\sphinxhyphen{}Forward Neural Networks (MLPs)}
\label{\detokenize{04:feed-forward-neural-networks-mlps}}

\subsection{Architecture and Notation}
\label{\detokenize{04:architecture-and-notation}}
\sphinxAtStartPar
An MLP (multi\sphinxhyphen{}layer perceptron, or feed\sphinxhyphen{}forward network) with \(L\) layers (including hidden and output layers) maps an input vector \(\mathbf{x}\in\mathbb{R}^n\) to an output (for regression, typically \(\hat{y}\in\mathbb{R}\) or \(\mathbb{R}^m\)) through successive affine transformations and nonlinear activations. In the notation of \sphinxstylestrong{Eqs. \((8.1.2)\)–\((8.1.5)\)}:
\begin{enumerate}
\sphinxsetlistlabels{\arabic}{enumi}{enumii}{}{.}%
\item {} 
\sphinxAtStartPar
\sphinxstylestrong{Layer 1}:

\end{enumerate}
\begin{equation*}
\begin{split}
A^{(1)} = Z^{(0)} W^{(1)} + b^{(1)},\quad Z^{(1)} = g^{(1)}(A^{(1)}),
\end{split}
\end{equation*}
\sphinxAtStartPar
where \(Z^{(0)} = \mathbf{x}\) is the input and \(W^{(1)}\) is an \((n \times h_1)\) weight matrix if the first hidden layer has \(h_1\) units.
\begin{enumerate}
\sphinxsetlistlabels{\arabic}{enumi}{enumii}{}{.}%
\setcounter{enumi}{1}
\item {} 
\sphinxAtStartPar
\sphinxstylestrong{Layer 2} (or subsequent hidden layers):

\end{enumerate}
\begin{equation*}
\begin{split}
A^{(2)} = Z^{(1)} W^{(2)} + b^{(2)},\quad Z^{(2)} = g^{(2)}(A^{(2)}),
\end{split}
\end{equation*}
\sphinxAtStartPar
continuing similarly up to layer \(L-1\).
\begin{enumerate}
\sphinxsetlistlabels{\arabic}{enumi}{enumii}{}{.}%
\setcounter{enumi}{2}
\item {} 
\sphinxAtStartPar
\sphinxstylestrong{Output Layer} (\(l=L\)):
For regression with a linear output,

\end{enumerate}
\begin{equation*}
\begin{split}
A^{(L)} = Z^{(L-1)} W^{(L)} + b^{(L)},\quad Z^{(L)} = A^{(L)}.
\end{split}
\end{equation*}
\sphinxAtStartPar
If the output dimension is 1, \(W^{(L)}\) is \((h_{L-1} \times 1)\). For vector outputs, adjust accordingly.

\sphinxAtStartPar
An example of representation of a MLP made of layers of size (16, 10, 6, 1), made with \sphinxhref{https://alexlenail.me/NN-SVG/index.html}{this tool}, is given here:

\sphinxAtStartPar
\sphinxincludegraphics{{mlp_example}.jpg}


\section{Convolutional Neural Networks (CNNs)}
\label{\detokenize{04:convolutional-neural-networks-cnns}}
\sphinxAtStartPar
Though the practical focuses on regression, CNNs (p. 198–199) illustrate deep architectures that handle grid\sphinxhyphen{}structured data (e.g., images). Main ideas:
\begin{enumerate}
\sphinxsetlistlabels{\arabic}{enumi}{enumii}{}{.}%
\item {} 
\sphinxAtStartPar
\sphinxstylestrong{Local Connectivity}: Each neuron in a convolutional layer only connects to a local patch of the input.

\item {} 
\sphinxAtStartPar
\sphinxstylestrong{Shared Weights}: A small kernel (filter) is “convolved” across the entire image or feature map.

\item {} 
\sphinxAtStartPar
\sphinxstylestrong{Translation Equivariance}: The same filter is applied at all spatial locations.

\end{enumerate}


\subsection{Convolution Operation}
\label{\detokenize{04:convolution-operation}}
\sphinxAtStartPar
A single 2D filter:
\begin{equation*}
\begin{split}
y_{ij} = \sum_{m=1}^{k_h} \sum_{n=1}^{k_w} W_{mn} x_{(i + m - 1, j + n - 1)} + b,
\end{split}
\end{equation*}
\sphinxAtStartPar
where \((k_h, k_w)\) is the kernel size. ReLU or another activation is then applied.


\subsection{Pooling Layers}
\label{\detokenize{04:pooling-layers}}
\sphinxAtStartPar
Often, a \sphinxstylestrong{pooling} layer follows one or more convolutional layers to reduce the spatial resolution. For instance, \sphinxstylestrong{max pooling} with a \(2\times2\) window and stride 2 halves the height and width.

\sphinxAtStartPar
An example of represention of a CNN is given here:

\sphinxAtStartPar
\sphinxincludegraphics{{cnn_example}.jpg}


\section{Activation Functions}
\label{\detokenize{04:activation-functions}}
\sphinxAtStartPar
Nonlinearities \(g^{(l)}(\cdot)\) are critical. Common examples in the course:
\begin{itemize}
\item {} 
\sphinxAtStartPar
\sphinxstylestrong{Sigmoid}: \(\sigma(z)=1/(1+e^{-z})\), which saturates for large \(|z|\).

\item {} 
\sphinxAtStartPar
\sphinxstylestrong{Tanh}: \(\tanh(z)\in(-1,1)\), similar saturation behavior.

\item {} 
\sphinxAtStartPar
\sphinxstylestrong{ReLU}: \(\max(0,z)\), which alleviates vanishing gradients on the positive side and is widely used in deep networks.

\end{itemize}

\sphinxAtStartPar
Without a nonlinear \(g^{(l)}\), the stacked linear transformations collapse to an effective single linear mapping. Nonlinearity ensures the network can approximate complex functions (Universal Approximation property).


\section{Cost function and matrix\sphinxhyphen{}based Back\sphinxhyphen{}Propagation}
\label{\detokenize{04:cost-function-and-matrix-based-back-propagation}}

\subsection{Sum of Squared Errors (SSE)}
\label{\detokenize{04:sum-of-squared-errors-sse}}
\sphinxAtStartPar
For regression tasks, the \sphinxstyleemphasis{empirical sum of squared errors}:
\begin{equation*}
\begin{split}
\mathrm{SSE}_{\mathrm{emp}}(\alpha_N) = \sum_{i=1}^N (y_i - \hat{y}(x_i))^2,
\end{split}
\end{equation*}
\sphinxAtStartPar
where \(\hat{y}(x_i)\) is the network’s prediction for input \(x_i\), and \(y_i\) is the true target. The parameter vector \(\alpha_N\) includes \(\{W^{(l)}, b^{(l)}\}\) for \(l=1,\ldots,L\).


\subsection{Gradient Descent Update}
\label{\detokenize{04:gradient-descent-update}}
\sphinxAtStartPar
Using the notation from \sphinxstylestrong{Algorithm 2 (p. 194)}, we iteratively update each parameter in the direction that reduces the SSE cost:
\begin{equation*}
\begin{split}
\alpha_N(k+1) = \alpha_N(k) - \eta\, \frac{\partial\,\mathrm{SSE}_{\mathrm{emp}}}{\partial\,\alpha_N}(k),
\end{split}
\end{equation*}
\sphinxAtStartPar
where \(\eta\) is the learning rate.


\subsection{Back\sphinxhyphen{}Propagation (Algorithm 2)}
\label{\detokenize{04:back-propagation-algorithm-2}}
\sphinxAtStartPar
Back\sphinxhyphen{}propagation is the application of the chain rule to compute \(\nabla \mathrm{SSE}_{\mathrm{emp}}\) with respect to each weight matrix \(W^{(l)}\) and bias \(b^{(l)}\).
\begin{enumerate}
\sphinxsetlistlabels{\arabic}{enumi}{enumii}{}{.}%
\item {} 
\sphinxAtStartPar
\sphinxstylestrong{Output\sphinxhyphen{}Layer Delta}:

\end{enumerate}
\begin{equation*}
\begin{split}
\delta^{(L)} = \frac{\partial\,\mathrm{SSE}_{\mathrm{emp}}}{\partial\,A^{(L)}} \odot g'^{(L)}(A^{(L)}).
\end{split}
\end{equation*}\begin{itemize}
\item {} 
\sphinxAtStartPar
In linear output regression, \(g'^{(L)}(\cdot) = 1\), and the cost derivative is \(\delta^{(L)} = (\hat{y}-y).\)

\end{itemize}
\begin{enumerate}
\sphinxsetlistlabels{\arabic}{enumi}{enumii}{}{.}%
\setcounter{enumi}{1}
\item {} 
\sphinxAtStartPar
\sphinxstylestrong{Backward Recursion}:

\end{enumerate}
\begin{equation*}
\begin{split}
\delta^{(l)} = \bigl(W^{(l+1)} \delta^{(l+1)}\bigr) \odot g'^{(l)}(A^{(l)}), \quad l = L-1,\ldots,1.
\end{split}
\end{equation*}
\sphinxAtStartPar
This “error” \(\delta^{(l)}\) is used to compute partial derivatives for \(W^{(l)}\) and \(b^{(l)}\).
\begin{enumerate}
\sphinxsetlistlabels{\arabic}{enumi}{enumii}{}{.}%
\setcounter{enumi}{2}
\item {} 
\sphinxAtStartPar
\sphinxstylestrong{Gradient w.r.t. the weights}:

\end{enumerate}
\begin{equation*}
\begin{split}
\frac{\partial\,\mathrm{SSE}_{\mathrm{emp}}}{\partial\,W^{(l)}} = (Z^{(l-1)})^\top \,\delta^{(l)},\quad
\frac{\partial\,\mathrm{SSE}_{\mathrm{emp}}}{\partial\,b^{(l)}} = \sum_{i=1}^{N} \delta^{(l)}_i.
\end{split}
\end{equation*}
\sphinxAtStartPar
using the matrix\sphinxhyphen{}based form.


\section{Example 1: Back\sphinxhyphen{}Propagation in a MLP}
\label{\detokenize{04:example-1-back-propagation-in-a-mlp}}
\begin{sphinxuseclass}{cell}\begin{sphinxVerbatimInput}

\begin{sphinxuseclass}{cell_input}
\begin{sphinxVerbatim}[commandchars=\\\{\}]
\PYG{k+kn}{import}\PYG{+w}{ }\PYG{n+nn}{torch}\PYG{n+nn}{.}\PYG{n+nn}{nn}\PYG{+w}{ }\PYG{k}{as}\PYG{+w}{ }\PYG{n+nn}{nn}

\PYG{k}{class}\PYG{+w}{ }\PYG{n+nc}{Simple2LayerMLP}\PYG{p}{(}\PYG{n}{nn}\PYG{o}{.}\PYG{n}{Module}\PYG{p}{)}\PYG{p}{:}
    \PYG{k}{def}\PYG{+w}{ }\PYG{n+nf+fm}{\PYGZus{}\PYGZus{}init\PYGZus{}\PYGZus{}}\PYG{p}{(}\PYG{n+nb+bp}{self}\PYG{p}{,} \PYG{n}{input\PYGZus{}dim}\PYG{p}{,} \PYG{n}{hidden\PYGZus{}dim}\PYG{p}{,} \PYG{n}{output\PYGZus{}dim}\PYG{o}{=}\PYG{l+m+mi}{1}\PYG{p}{)}\PYG{p}{:}
        \PYG{n+nb}{super}\PYG{p}{(}\PYG{p}{)}\PYG{o}{.}\PYG{n+nf+fm}{\PYGZus{}\PYGZus{}init\PYGZus{}\PYGZus{}}\PYG{p}{(}\PYG{p}{)}
        \PYG{n+nb+bp}{self}\PYG{o}{.}\PYG{n}{fc1} \PYG{o}{=} \PYG{n}{nn}\PYG{o}{.}\PYG{n}{Linear}\PYG{p}{(}\PYG{n}{input\PYGZus{}dim}\PYG{p}{,} \PYG{n}{hidden\PYGZus{}dim}\PYG{p}{)}
        \PYG{n+nb+bp}{self}\PYG{o}{.}\PYG{n}{fc2} \PYG{o}{=} \PYG{n}{nn}\PYG{o}{.}\PYG{n}{Linear}\PYG{p}{(}\PYG{n}{hidden\PYGZus{}dim}\PYG{p}{,} \PYG{n}{output\PYGZus{}dim}\PYG{p}{)}

    \PYG{k}{def}\PYG{+w}{ }\PYG{n+nf}{forward}\PYG{p}{(}\PYG{n+nb+bp}{self}\PYG{p}{,} \PYG{n}{x}\PYG{p}{)}\PYG{p}{:}
        \PYG{n}{x} \PYG{o}{=} \PYG{n+nb+bp}{self}\PYG{o}{.}\PYG{n}{fc1}\PYG{p}{(}\PYG{n}{x}\PYG{p}{)}
        \PYG{n}{x} \PYG{o}{=} \PYG{n}{torch}\PYG{o}{.}\PYG{n}{relu}\PYG{p}{(}\PYG{n}{x}\PYG{p}{)}
        \PYG{n}{x} \PYG{o}{=} \PYG{n+nb+bp}{self}\PYG{o}{.}\PYG{n}{fc2}\PYG{p}{(}\PYG{n}{x}\PYG{p}{)}
        \PYG{k}{return} \PYG{n}{x}

\PYG{n}{model} \PYG{o}{=} \PYG{n}{Simple2LayerMLP}\PYG{p}{(}\PYG{n}{input\PYGZus{}dim}\PYG{o}{=}\PYG{l+m+mi}{10}\PYG{p}{,} \PYG{n}{hidden\PYGZus{}dim}\PYG{o}{=}\PYG{l+m+mi}{15}\PYG{p}{,} \PYG{n}{output\PYGZus{}dim}\PYG{o}{=}\PYG{l+m+mi}{1}\PYG{p}{)}
\PYG{n}{criterion} \PYG{o}{=} \PYG{n}{nn}\PYG{o}{.}\PYG{n}{MSELoss}\PYG{p}{(}\PYG{n}{reduction}\PYG{o}{=}\PYG{l+s+s1}{\PYGZsq{}}\PYG{l+s+s1}{sum}\PYG{l+s+s1}{\PYGZsq{}}\PYG{p}{)}
\PYG{n}{eta} \PYG{o}{=} \PYG{l+m+mf}{1e\PYGZhy{}2}

\PYG{n}{X\PYGZus{}sample} \PYG{o}{=} \PYG{n}{torch}\PYG{o}{.}\PYG{n}{randn}\PYG{p}{(}\PYG{l+m+mi}{5}\PYG{p}{,} \PYG{l+m+mi}{10}\PYG{p}{)}
\PYG{n}{y\PYGZus{}sample} \PYG{o}{=} \PYG{n}{torch}\PYG{o}{.}\PYG{n}{randn}\PYG{p}{(}\PYG{l+m+mi}{5}\PYG{p}{,} \PYG{l+m+mi}{1}\PYG{p}{)}

\PYG{n}{y\PYGZus{}pred} \PYG{o}{=} \PYG{n}{model}\PYG{p}{(}\PYG{n}{X\PYGZus{}sample}\PYG{p}{)}
\PYG{n}{loss} \PYG{o}{=} \PYG{n}{criterion}\PYG{p}{(}\PYG{n}{y\PYGZus{}pred}\PYG{p}{,} \PYG{n}{y\PYGZus{}sample}\PYG{p}{)}

\PYG{n}{loss}\PYG{o}{.}\PYG{n}{backward}\PYG{p}{(}\PYG{p}{)}

\PYG{k}{with} \PYG{n}{torch}\PYG{o}{.}\PYG{n}{no\PYGZus{}grad}\PYG{p}{(}\PYG{p}{)}\PYG{p}{:}
    \PYG{k}{for} \PYG{n}{param} \PYG{o+ow}{in} \PYG{n}{model}\PYG{o}{.}\PYG{n}{parameters}\PYG{p}{(}\PYG{p}{)}\PYG{p}{:}
        \PYG{n}{param} \PYG{o}{\PYGZhy{}}\PYG{o}{=} \PYG{n}{eta} \PYG{o}{*} \PYG{n}{param}\PYG{o}{.}\PYG{n}{grad}
        \PYG{n}{param}\PYG{o}{.}\PYG{n}{grad}\PYG{o}{.}\PYG{n}{zero\PYGZus{}}\PYG{p}{(}\PYG{p}{)}
\end{sphinxVerbatim}

\end{sphinxuseclass}\end{sphinxVerbatimInput}

\end{sphinxuseclass}

\section{Overfitting and Cross\sphinxhyphen{}Validation}
\label{\detokenize{04:overfitting-and-cross-validation}}

\subsection{Overfitting Phenomenon}
\label{\detokenize{04:overfitting-phenomenon}}
\sphinxAtStartPar
In the handbook’s example on \sphinxstylestrong{p. 197}, a network with 15 hidden units achieves extremely low training SSE but high test SSE, illustrating \sphinxstylestrong{overfitting}. Overfitting arises when a model fits training data “too well,” capturing noise rather than the underlying trend.


\subsection{Detection and Mitigation}
\label{\detokenize{04:detection-and-mitigation}}\begin{itemize}
\item {} 
\sphinxAtStartPar
\sphinxstylestrong{Validation Set}: Split data into training and validation. Monitor the validation error (e.g., MSE). If it increases while training error decreases, the model is overfitting.

\item {} 
\sphinxAtStartPar
\sphinxstylestrong{Cross\sphinxhyphen{}Validation}: Alternatively, use \(k\)\sphinxhyphen{}fold cross\sphinxhyphen{}validation.

\item {} 
\sphinxAtStartPar
\sphinxstylestrong{Regularization}:
\begin{itemize}
\item {} 
\sphinxAtStartPar
\sphinxstylestrong{Dropout}: Temporarily “drop” a fraction of hidden units during training.

\item {} 
\sphinxAtStartPar
\sphinxstylestrong{Weight Decay (L2)}: Penalizes large weights.

\item {} 
\sphinxAtStartPar
\sphinxstylestrong{Early Stopping}: Halt training when validation performance stops improving.

\end{itemize}

\end{itemize}


\section{Model Persistence: Saving and Loading Models in PyTorch}
\label{\detokenize{04:model-persistence-saving-and-loading-models-in-pytorch}}
\sphinxAtStartPar
For safe and reproducible experiments, one typically needs to \sphinxstylestrong{save}:
\begin{enumerate}
\sphinxsetlistlabels{\arabic}{enumi}{enumii}{}{.}%
\item {} 
\sphinxAtStartPar
\sphinxstylestrong{Model Weights (state\_dict)}: The learned parameters \(\{W^{(l)},b^{(l)}\}\).

\item {} 
\sphinxAtStartPar
\sphinxstylestrong{Optimizer State (if resuming)}: Internal states like momentum buffers in gradient\sphinxhyphen{}based optimizers.

\end{enumerate}

\sphinxAtStartPar
A typical code pattern:

\begin{sphinxuseclass}{cell}\begin{sphinxVerbatimInput}

\begin{sphinxuseclass}{cell_input}
\begin{sphinxVerbatim}[commandchars=\\\{\}]
\PYG{k+kn}{import}\PYG{+w}{ }\PYG{n+nn}{torch}\PYG{n+nn}{.}\PYG{n+nn}{optim}\PYG{+w}{ }\PYG{k}{as}\PYG{+w}{ }\PYG{n+nn}{optim}

\PYG{n}{model} \PYG{o}{=} \PYG{n}{Simple2LayerMLP}\PYG{p}{(}\PYG{l+m+mi}{10}\PYG{p}{,} \PYG{l+m+mi}{15}\PYG{p}{,} \PYG{l+m+mi}{1}\PYG{p}{)}
\PYG{n}{optimizer} \PYG{o}{=} \PYG{n}{optim}\PYG{o}{.}\PYG{n}{SGD}\PYG{p}{(}\PYG{n}{model}\PYG{o}{.}\PYG{n}{parameters}\PYG{p}{(}\PYG{p}{)}\PYG{p}{,} \PYG{n}{lr}\PYG{o}{=}\PYG{l+m+mf}{0.01}\PYG{p}{)}

\PYG{n}{torch}\PYG{o}{.}\PYG{n}{save}\PYG{p}{(}\PYG{p}{\PYGZob{}}
    \PYG{l+s+s1}{\PYGZsq{}}\PYG{l+s+s1}{model\PYGZus{}state\PYGZus{}dict}\PYG{l+s+s1}{\PYGZsq{}}\PYG{p}{:} \PYG{n}{model}\PYG{o}{.}\PYG{n}{state\PYGZus{}dict}\PYG{p}{(}\PYG{p}{)}\PYG{p}{,}
    \PYG{l+s+s1}{\PYGZsq{}}\PYG{l+s+s1}{optimizer\PYGZus{}state\PYGZus{}dict}\PYG{l+s+s1}{\PYGZsq{}}\PYG{p}{:} \PYG{n}{optimizer}\PYG{o}{.}\PYG{n}{state\PYGZus{}dict}\PYG{p}{(}\PYG{p}{)}\PYG{p}{,}
    \PYG{l+s+s1}{\PYGZsq{}}\PYG{l+s+s1}{epoch}\PYG{l+s+s1}{\PYGZsq{}}\PYG{p}{:} \PYG{l+m+mi}{10}
\PYG{p}{\PYGZcb{}}\PYG{p}{,} \PYG{l+s+s2}{\PYGZdq{}}\PYG{l+s+s2}{checkpoint.pth}\PYG{l+s+s2}{\PYGZdq{}}\PYG{p}{)}

\PYG{n}{checkpoint} \PYG{o}{=} \PYG{n}{torch}\PYG{o}{.}\PYG{n}{load}\PYG{p}{(}\PYG{l+s+s2}{\PYGZdq{}}\PYG{l+s+s2}{checkpoint.pth}\PYG{l+s+s2}{\PYGZdq{}}\PYG{p}{,} \PYG{n}{weights\PYGZus{}only}\PYG{o}{=}\PYG{k+kc}{False}\PYG{p}{)}
\PYG{n}{model}\PYG{o}{.}\PYG{n}{load\PYGZus{}state\PYGZus{}dict}\PYG{p}{(}\PYG{n}{checkpoint}\PYG{p}{[}\PYG{l+s+s1}{\PYGZsq{}}\PYG{l+s+s1}{model\PYGZus{}state\PYGZus{}dict}\PYG{l+s+s1}{\PYGZsq{}}\PYG{p}{]}\PYG{p}{)}
\PYG{n}{optimizer}\PYG{o}{.}\PYG{n}{load\PYGZus{}state\PYGZus{}dict}\PYG{p}{(}\PYG{n}{checkpoint}\PYG{p}{[}\PYG{l+s+s1}{\PYGZsq{}}\PYG{l+s+s1}{optimizer\PYGZus{}state\PYGZus{}dict}\PYG{l+s+s1}{\PYGZsq{}}\PYG{p}{]}\PYG{p}{)}
\PYG{n}{start\PYGZus{}epoch} \PYG{o}{=} \PYG{n}{checkpoint}\PYG{p}{[}\PYG{l+s+s1}{\PYGZsq{}}\PYG{l+s+s1}{epoch}\PYG{l+s+s1}{\PYGZsq{}}\PYG{p}{]} \PYG{o}{+} \PYG{l+m+mi}{1}

\PYG{c+c1}{\PYGZsh{} For inference, only model\PYGZus{}state\PYGZus{}dict is strictly required.}
\end{sphinxVerbatim}

\end{sphinxuseclass}\end{sphinxVerbatimInput}

\end{sphinxuseclass}

\section{Example 2: CNN for Regression}
\label{\detokenize{04:example-2-cnn-for-regression}}
\sphinxAtStartPar
Below is a more detailed illustration of an implementation of a CNN in PyTorch for regression highlighting many of the points presented above:

\begin{sphinxuseclass}{cell}\begin{sphinxVerbatimInput}

\begin{sphinxuseclass}{cell_input}
\begin{sphinxVerbatim}[commandchars=\\\{\}]
\PYG{k+kn}{import}\PYG{+w}{ }\PYG{n+nn}{torch}\PYG{n+nn}{.}\PYG{n+nn}{nn}\PYG{n+nn}{.}\PYG{n+nn}{functional}\PYG{+w}{ }\PYG{k}{as}\PYG{+w}{ }\PYG{n+nn}{F}
\PYG{k+kn}{from}\PYG{+w}{ }\PYG{n+nn}{torch}\PYG{n+nn}{.}\PYG{n+nn}{utils}\PYG{n+nn}{.}\PYG{n+nn}{data}\PYG{+w}{ }\PYG{k+kn}{import} \PYG{n}{Dataset}\PYG{p}{,} \PYG{n}{DataLoader}\PYG{p}{,} \PYG{n}{TensorDataset}

\PYG{k}{class}\PYG{+w}{ }\PYG{n+nc}{SimpleCNNRegression}\PYG{p}{(}\PYG{n}{nn}\PYG{o}{.}\PYG{n}{Module}\PYG{p}{)}\PYG{p}{:}
    \PYG{k}{def}\PYG{+w}{ }\PYG{n+nf+fm}{\PYGZus{}\PYGZus{}init\PYGZus{}\PYGZus{}}\PYG{p}{(}\PYG{n+nb+bp}{self}\PYG{p}{)}\PYG{p}{:}
        \PYG{n+nb}{super}\PYG{p}{(}\PYG{p}{)}\PYG{o}{.}\PYG{n+nf+fm}{\PYGZus{}\PYGZus{}init\PYGZus{}\PYGZus{}}\PYG{p}{(}\PYG{p}{)}
        
        \PYG{n+nb+bp}{self}\PYG{o}{.}\PYG{n}{conv1} \PYG{o}{=} \PYG{n}{nn}\PYG{o}{.}\PYG{n}{Conv2d}\PYG{p}{(}\PYG{n}{in\PYGZus{}channels}\PYG{o}{=}\PYG{l+m+mi}{1}\PYG{p}{,} \PYG{n}{out\PYGZus{}channels}\PYG{o}{=}\PYG{l+m+mi}{8}\PYG{p}{,} \PYG{n}{kernel\PYGZus{}size}\PYG{o}{=}\PYG{l+m+mi}{3}\PYG{p}{,} \PYG{n}{padding}\PYG{o}{=}\PYG{l+m+mi}{1}\PYG{p}{)}
        \PYG{n+nb+bp}{self}\PYG{o}{.}\PYG{n}{conv2} \PYG{o}{=} \PYG{n}{nn}\PYG{o}{.}\PYG{n}{Conv2d}\PYG{p}{(}\PYG{n}{in\PYGZus{}channels}\PYG{o}{=}\PYG{l+m+mi}{8}\PYG{p}{,} \PYG{n}{out\PYGZus{}channels}\PYG{o}{=}\PYG{l+m+mi}{16}\PYG{p}{,} \PYG{n}{kernel\PYGZus{}size}\PYG{o}{=}\PYG{l+m+mi}{3}\PYG{p}{,} \PYG{n}{padding}\PYG{o}{=}\PYG{l+m+mi}{1}\PYG{p}{)}
        \PYG{n+nb+bp}{self}\PYG{o}{.}\PYG{n}{fc1}   \PYG{o}{=} \PYG{n}{nn}\PYG{o}{.}\PYG{n}{Linear}\PYG{p}{(}\PYG{l+m+mi}{16} \PYG{o}{*} \PYG{l+m+mi}{7} \PYG{o}{*} \PYG{l+m+mi}{7}\PYG{p}{,} \PYG{l+m+mi}{64}\PYG{p}{)}  
        \PYG{n+nb+bp}{self}\PYG{o}{.}\PYG{n}{fc2}   \PYG{o}{=} \PYG{n}{nn}\PYG{o}{.}\PYG{n}{Linear}\PYG{p}{(}\PYG{l+m+mi}{64}\PYG{p}{,} \PYG{l+m+mi}{1}\PYG{p}{)}
        
        \PYG{n+nb+bp}{self}\PYG{o}{.}\PYG{n}{dropout} \PYG{o}{=} \PYG{n}{nn}\PYG{o}{.}\PYG{n}{Dropout}\PYG{p}{(}\PYG{n}{p}\PYG{o}{=}\PYG{l+m+mf}{0.25}\PYG{p}{)}
        \PYG{n+nb+bp}{self}\PYG{o}{.}\PYG{n}{pool}   \PYG{o}{=} \PYG{n}{nn}\PYG{o}{.}\PYG{n}{MaxPool2d}\PYG{p}{(}\PYG{n}{kernel\PYGZus{}size}\PYG{o}{=}\PYG{l+m+mi}{2}\PYG{p}{,} \PYG{n}{stride}\PYG{o}{=}\PYG{l+m+mi}{2}\PYG{p}{)}

    \PYG{k}{def}\PYG{+w}{ }\PYG{n+nf}{forward}\PYG{p}{(}\PYG{n+nb+bp}{self}\PYG{p}{,} \PYG{n}{x}\PYG{p}{)}\PYG{p}{:}
        \PYG{c+c1}{\PYGZsh{} First conv: conv1 \PYGZhy{}\PYGZgt{} ReLU \PYGZhy{}\PYGZgt{} pool =\PYGZgt{} shape: (N, 8, 14, 14)}
        \PYG{n}{x} \PYG{o}{=} \PYG{n+nb+bp}{self}\PYG{o}{.}\PYG{n}{conv1}\PYG{p}{(}\PYG{n}{x}\PYG{p}{)}        
        \PYG{n}{x} \PYG{o}{=} \PYG{n}{F}\PYG{o}{.}\PYG{n}{relu}\PYG{p}{(}\PYG{n}{x}\PYG{p}{)}
        \PYG{n}{x} \PYG{o}{=} \PYG{n+nb+bp}{self}\PYG{o}{.}\PYG{n}{pool}\PYG{p}{(}\PYG{n}{x}\PYG{p}{)}
        
        \PYG{c+c1}{\PYGZsh{} Second conv: conv2 \PYGZhy{}\PYGZgt{} ReLU \PYGZhy{}\PYGZgt{} pool =\PYGZgt{} shape: (N, 16, 7, 7)}
        \PYG{n}{x} \PYG{o}{=} \PYG{n+nb+bp}{self}\PYG{o}{.}\PYG{n}{conv2}\PYG{p}{(}\PYG{n}{x}\PYG{p}{)}        
        \PYG{n}{x} \PYG{o}{=} \PYG{n}{F}\PYG{o}{.}\PYG{n}{relu}\PYG{p}{(}\PYG{n}{x}\PYG{p}{)}
        \PYG{n}{x} \PYG{o}{=} \PYG{n+nb+bp}{self}\PYG{o}{.}\PYG{n}{pool}\PYG{p}{(}\PYG{n}{x}\PYG{p}{)}
        
        \PYG{c+c1}{\PYGZsh{} Flatten =\PYGZgt{} shape: (N, 16*7*7 = 784)}
        \PYG{n}{x} \PYG{o}{=} \PYG{n}{x}\PYG{o}{.}\PYG{n}{view}\PYG{p}{(}\PYG{n}{x}\PYG{o}{.}\PYG{n}{size}\PYG{p}{(}\PYG{l+m+mi}{0}\PYG{p}{)}\PYG{p}{,} \PYG{o}{\PYGZhy{}}\PYG{l+m+mi}{1}\PYG{p}{)}
        
        \PYG{c+c1}{\PYGZsh{} First fully connected =\PYGZgt{} ReLU =\PYGZgt{} Dropout =\PYGZgt{} shape: (N, 64)}
        \PYG{n}{x} \PYG{o}{=} \PYG{n+nb+bp}{self}\PYG{o}{.}\PYG{n}{fc1}\PYG{p}{(}\PYG{n}{x}\PYG{p}{)}
        \PYG{n}{x} \PYG{o}{=} \PYG{n}{F}\PYG{o}{.}\PYG{n}{relu}\PYG{p}{(}\PYG{n}{x}\PYG{p}{)}
        \PYG{n}{x} \PYG{o}{=} \PYG{n+nb+bp}{self}\PYG{o}{.}\PYG{n}{dropout}\PYG{p}{(}\PYG{n}{x}\PYG{p}{)}
        
        \PYG{c+c1}{\PYGZsh{} Final fully connected =\PYGZgt{} shape: (N, 1)}
        \PYG{n}{x} \PYG{o}{=} \PYG{n+nb+bp}{self}\PYG{o}{.}\PYG{n}{fc2}\PYG{p}{(}\PYG{n}{x}\PYG{p}{)}
        \PYG{k}{return} \PYG{n}{x}

\PYG{n}{N\PYGZus{}train} \PYG{o}{=} \PYG{l+m+mi}{1000}
\PYG{n}{N\PYGZus{}val}   \PYG{o}{=} \PYG{l+m+mi}{200}
\PYG{n}{X\PYGZus{}train} \PYG{o}{=} \PYG{n}{torch}\PYG{o}{.}\PYG{n}{randn}\PYG{p}{(}\PYG{n}{N\PYGZus{}train}\PYG{p}{,} \PYG{l+m+mi}{1}\PYG{p}{,} \PYG{l+m+mi}{28}\PYG{p}{,} \PYG{l+m+mi}{28}\PYG{p}{)}
\PYG{n}{y\PYGZus{}train} \PYG{o}{=} \PYG{n}{torch}\PYG{o}{.}\PYG{n}{randn}\PYG{p}{(}\PYG{n}{N\PYGZus{}train}\PYG{p}{,} \PYG{l+m+mi}{1}\PYG{p}{)} 
\PYG{n}{X\PYGZus{}val}   \PYG{o}{=} \PYG{n}{torch}\PYG{o}{.}\PYG{n}{randn}\PYG{p}{(}\PYG{n}{N\PYGZus{}val}\PYG{p}{,} \PYG{l+m+mi}{1}\PYG{p}{,} \PYG{l+m+mi}{28}\PYG{p}{,} \PYG{l+m+mi}{28}\PYG{p}{)}
\PYG{n}{y\PYGZus{}val}   \PYG{o}{=} \PYG{n}{torch}\PYG{o}{.}\PYG{n}{randn}\PYG{p}{(}\PYG{n}{N\PYGZus{}val}\PYG{p}{,} \PYG{l+m+mi}{1}\PYG{p}{)}

\PYG{n}{train\PYGZus{}dataset} \PYG{o}{=} \PYG{n}{TensorDataset}\PYG{p}{(}\PYG{n}{X\PYGZus{}train}\PYG{p}{,} \PYG{n}{y\PYGZus{}train}\PYG{p}{)}
\PYG{n}{val\PYGZus{}dataset}   \PYG{o}{=} \PYG{n}{TensorDataset}\PYG{p}{(}\PYG{n}{X\PYGZus{}val}\PYG{p}{,} \PYG{n}{y\PYGZus{}val}\PYG{p}{)}
\PYG{n}{batch\PYGZus{}size}    \PYG{o}{=} \PYG{l+m+mi}{50}
\PYG{n}{train\PYGZus{}loader}  \PYG{o}{=} \PYG{n}{DataLoader}\PYG{p}{(}\PYG{n}{train\PYGZus{}dataset}\PYG{p}{,} \PYG{n}{batch\PYGZus{}size}\PYG{o}{=}\PYG{n}{batch\PYGZus{}size}\PYG{p}{,} \PYG{n}{shuffle}\PYG{o}{=}\PYG{k+kc}{True}\PYG{p}{)}
\PYG{n}{val\PYGZus{}loader}    \PYG{o}{=} \PYG{n}{DataLoader}\PYG{p}{(}\PYG{n}{val\PYGZus{}dataset}\PYG{p}{,}   \PYG{n}{batch\PYGZus{}size}\PYG{o}{=}\PYG{n}{batch\PYGZus{}size}\PYG{p}{,} \PYG{n}{shuffle}\PYG{o}{=}\PYG{k+kc}{False}\PYG{p}{)}

\PYG{n}{model} \PYG{o}{=} \PYG{n}{SimpleCNNRegression}\PYG{p}{(}\PYG{p}{)}
\PYG{n}{criterion} \PYG{o}{=} \PYG{n}{nn}\PYG{o}{.}\PYG{n}{MSELoss}\PYG{p}{(}\PYG{p}{)}

\PYG{n}{optimizer} \PYG{o}{=} \PYG{n}{optim}\PYG{o}{.}\PYG{n}{Adam}\PYG{p}{(}\PYG{n}{model}\PYG{o}{.}\PYG{n}{parameters}\PYG{p}{(}\PYG{p}{)}\PYG{p}{,} \PYG{n}{lr}\PYG{o}{=}\PYG{l+m+mf}{0.001}\PYG{p}{,} \PYG{n}{weight\PYGZus{}decay}\PYG{o}{=}\PYG{l+m+mf}{1e\PYGZhy{}4}\PYG{p}{)}

\PYG{n}{max\PYGZus{}epochs}         \PYG{o}{=} \PYG{l+m+mi}{20}
\PYG{n}{patience}           \PYG{o}{=} \PYG{l+m+mi}{3} 
\PYG{n}{best\PYGZus{}val\PYGZus{}loss}      \PYG{o}{=} \PYG{n+nb}{float}\PYG{p}{(}\PYG{l+s+s1}{\PYGZsq{}}\PYG{l+s+s1}{inf}\PYG{l+s+s1}{\PYGZsq{}}\PYG{p}{)}
\PYG{n}{epochs\PYGZus{}no\PYGZus{}improve}  \PYG{o}{=} \PYG{l+m+mi}{0}
\PYG{n}{best\PYGZus{}model\PYGZus{}path}    \PYG{o}{=} \PYG{l+s+s2}{\PYGZdq{}}\PYG{l+s+s2}{best\PYGZus{}model.pth}\PYG{l+s+s2}{\PYGZdq{}}

\PYG{k}{for} \PYG{n}{epoch} \PYG{o+ow}{in} \PYG{n+nb}{range}\PYG{p}{(}\PYG{n}{max\PYGZus{}epochs}\PYG{p}{)}\PYG{p}{:}
    \PYG{n}{model}\PYG{o}{.}\PYG{n}{train}\PYG{p}{(}\PYG{p}{)}
    \PYG{n}{running\PYGZus{}train\PYGZus{}loss} \PYG{o}{=} \PYG{l+m+mf}{0.0}
    \PYG{k}{for} \PYG{n}{X\PYGZus{}batch}\PYG{p}{,} \PYG{n}{y\PYGZus{}batch} \PYG{o+ow}{in} \PYG{n}{train\PYGZus{}loader}\PYG{p}{:}
        \PYG{n}{optimizer}\PYG{o}{.}\PYG{n}{zero\PYGZus{}grad}\PYG{p}{(}\PYG{p}{)}
        \PYG{n}{preds} \PYG{o}{=} \PYG{n}{model}\PYG{p}{(}\PYG{n}{X\PYGZus{}batch}\PYG{p}{)}
        \PYG{n}{loss} \PYG{o}{=} \PYG{n}{criterion}\PYG{p}{(}\PYG{n}{preds}\PYG{p}{,} \PYG{n}{y\PYGZus{}batch}\PYG{p}{)}
        \PYG{n}{loss}\PYG{o}{.}\PYG{n}{backward}\PYG{p}{(}\PYG{p}{)}
        \PYG{n}{optimizer}\PYG{o}{.}\PYG{n}{step}\PYG{p}{(}\PYG{p}{)}
        \PYG{n}{running\PYGZus{}train\PYGZus{}loss} \PYG{o}{+}\PYG{o}{=} \PYG{n}{loss}\PYG{o}{.}\PYG{n}{item}\PYG{p}{(}\PYG{p}{)}
    \PYG{n}{avg\PYGZus{}train\PYGZus{}loss} \PYG{o}{=} \PYG{n}{running\PYGZus{}train\PYGZus{}loss} \PYG{o}{/} \PYG{n+nb}{len}\PYG{p}{(}\PYG{n}{train\PYGZus{}loader}\PYG{p}{)}
    
    \PYG{n}{model}\PYG{o}{.}\PYG{n}{eval}\PYG{p}{(}\PYG{p}{)}
    \PYG{n}{running\PYGZus{}val\PYGZus{}loss} \PYG{o}{=} \PYG{l+m+mf}{0.0}
    \PYG{k}{with} \PYG{n}{torch}\PYG{o}{.}\PYG{n}{no\PYGZus{}grad}\PYG{p}{(}\PYG{p}{)}\PYG{p}{:}
        \PYG{k}{for} \PYG{n}{Xv\PYGZus{}batch}\PYG{p}{,} \PYG{n}{yv\PYGZus{}batch} \PYG{o+ow}{in} \PYG{n}{val\PYGZus{}loader}\PYG{p}{:}
            \PYG{n}{val\PYGZus{}preds} \PYG{o}{=} \PYG{n}{model}\PYG{p}{(}\PYG{n}{Xv\PYGZus{}batch}\PYG{p}{)}
            \PYG{n}{val\PYGZus{}loss} \PYG{o}{=} \PYG{n}{criterion}\PYG{p}{(}\PYG{n}{val\PYGZus{}preds}\PYG{p}{,} \PYG{n}{yv\PYGZus{}batch}\PYG{p}{)}
            \PYG{n}{running\PYGZus{}val\PYGZus{}loss} \PYG{o}{+}\PYG{o}{=} \PYG{n}{val\PYGZus{}loss}\PYG{o}{.}\PYG{n}{item}\PYG{p}{(}\PYG{p}{)}
    
    \PYG{n}{avg\PYGZus{}val\PYGZus{}loss} \PYG{o}{=} \PYG{n}{running\PYGZus{}val\PYGZus{}loss} \PYG{o}{/} \PYG{n+nb}{len}\PYG{p}{(}\PYG{n}{val\PYGZus{}loader}\PYG{p}{)}
    
    \PYG{n+nb}{print}\PYG{p}{(}\PYG{l+s+sa}{f}\PYG{l+s+s2}{\PYGZdq{}}\PYG{l+s+s2}{Epoch [}\PYG{l+s+si}{\PYGZob{}}\PYG{n}{epoch}\PYG{o}{+}\PYG{l+m+mi}{1}\PYG{l+s+si}{\PYGZcb{}}\PYG{l+s+s2}{/}\PYG{l+s+si}{\PYGZob{}}\PYG{n}{max\PYGZus{}epochs}\PYG{l+s+si}{\PYGZcb{}}\PYG{l+s+s2}{]  }\PYG{l+s+s2}{\PYGZdq{}}
          \PYG{l+s+sa}{f}\PYG{l+s+s2}{\PYGZdq{}}\PYG{l+s+s2}{Train MSE: }\PYG{l+s+si}{\PYGZob{}}\PYG{n}{avg\PYGZus{}train\PYGZus{}loss}\PYG{l+s+si}{:}\PYG{l+s+s2}{.4f}\PYG{l+s+si}{\PYGZcb{}}\PYG{l+s+s2}{  }\PYG{l+s+s2}{\PYGZdq{}}
          \PYG{l+s+sa}{f}\PYG{l+s+s2}{\PYGZdq{}}\PYG{l+s+s2}{Val MSE: }\PYG{l+s+si}{\PYGZob{}}\PYG{n}{avg\PYGZus{}val\PYGZus{}loss}\PYG{l+s+si}{:}\PYG{l+s+s2}{.4f}\PYG{l+s+si}{\PYGZcb{}}\PYG{l+s+s2}{\PYGZdq{}}\PYG{p}{)}
    
    \PYG{k}{if} \PYG{n}{avg\PYGZus{}val\PYGZus{}loss} \PYG{o}{\PYGZlt{}} \PYG{n}{best\PYGZus{}val\PYGZus{}loss}\PYG{p}{:}
        \PYG{n}{best\PYGZus{}val\PYGZus{}loss} \PYG{o}{=} \PYG{n}{avg\PYGZus{}val\PYGZus{}loss}
        \PYG{n}{epochs\PYGZus{}no\PYGZus{}improve} \PYG{o}{=} \PYG{l+m+mi}{0}
        \PYG{n}{torch}\PYG{o}{.}\PYG{n}{save}\PYG{p}{(}\PYG{n}{model}\PYG{o}{.}\PYG{n}{state\PYGZus{}dict}\PYG{p}{(}\PYG{p}{)}\PYG{p}{,} \PYG{n}{best\PYGZus{}model\PYGZus{}path}\PYG{p}{)}
        \PYG{n+nb}{print}\PYG{p}{(}\PYG{l+s+s2}{\PYGZdq{}}\PYG{l+s+s2}{  Best model saved (improved validation MSE).}\PYG{l+s+s2}{\PYGZdq{}}\PYG{p}{)}
    \PYG{k}{else}\PYG{p}{:}
        \PYG{n}{epochs\PYGZus{}no\PYGZus{}improve} \PYG{o}{+}\PYG{o}{=} \PYG{l+m+mi}{1}
        \PYG{k}{if} \PYG{n}{epochs\PYGZus{}no\PYGZus{}improve} \PYG{o}{\PYGZgt{}}\PYG{o}{=} \PYG{n}{patience}\PYG{p}{:}
            \PYG{n+nb}{print}\PYG{p}{(}\PYG{l+s+s2}{\PYGZdq{}}\PYG{l+s+s2}{Early stopping triggered (no improvement).}\PYG{l+s+s2}{\PYGZdq{}}\PYG{p}{)}
            \PYG{k}{break}

\PYG{n+nb}{print}\PYG{p}{(}\PYG{l+s+s2}{\PYGZdq{}}\PYG{l+s+se}{\PYGZbs{}n}\PYG{l+s+s2}{Training complete. Loading best model for final checks.}\PYG{l+s+s2}{\PYGZdq{}}\PYG{p}{)}
\PYG{n}{best\PYGZus{}model} \PYG{o}{=} \PYG{n}{SimpleCNNRegression}\PYG{p}{(}\PYG{p}{)}
\PYG{n}{best\PYGZus{}model}\PYG{o}{.}\PYG{n}{load\PYGZus{}state\PYGZus{}dict}\PYG{p}{(}\PYG{n}{torch}\PYG{o}{.}\PYG{n}{load}\PYG{p}{(}\PYG{n}{best\PYGZus{}model\PYGZus{}path}\PYG{p}{,} \PYG{n}{weights\PYGZus{}only}\PYG{o}{=}\PYG{k+kc}{False}\PYG{p}{)}\PYG{p}{)}
\PYG{n}{best\PYGZus{}model}\PYG{o}{.}\PYG{n}{eval}\PYG{p}{(}\PYG{p}{)}

\PYG{n}{final\PYGZus{}val\PYGZus{}loss} \PYG{o}{=} \PYG{l+m+mf}{0.0}
\PYG{k}{with} \PYG{n}{torch}\PYG{o}{.}\PYG{n}{no\PYGZus{}grad}\PYG{p}{(}\PYG{p}{)}\PYG{p}{:}
    \PYG{k}{for} \PYG{n}{Xv\PYGZus{}batch}\PYG{p}{,} \PYG{n}{yv\PYGZus{}batch} \PYG{o+ow}{in} \PYG{n}{val\PYGZus{}loader}\PYG{p}{:}
        \PYG{n}{val\PYGZus{}preds} \PYG{o}{=} \PYG{n}{best\PYGZus{}model}\PYG{p}{(}\PYG{n}{Xv\PYGZus{}batch}\PYG{p}{)}
        \PYG{n}{val\PYGZus{}loss} \PYG{o}{=} \PYG{n}{criterion}\PYG{p}{(}\PYG{n}{val\PYGZus{}preds}\PYG{p}{,} \PYG{n}{yv\PYGZus{}batch}\PYG{p}{)}
        \PYG{n}{final\PYGZus{}val\PYGZus{}loss} \PYG{o}{+}\PYG{o}{=} \PYG{n}{val\PYGZus{}loss}\PYG{o}{.}\PYG{n}{item}\PYG{p}{(}\PYG{p}{)}
\PYG{n+nb}{print}\PYG{p}{(}\PYG{l+s+sa}{f}\PYG{l+s+s2}{\PYGZdq{}}\PYG{l+s+s2}{Final validation MSE (best model): }\PYG{l+s+si}{\PYGZob{}}\PYG{n}{final\PYGZus{}val\PYGZus{}loss}\PYG{+w}{ }\PYG{o}{/}\PYG{+w}{ }\PYG{n+nb}{len}\PYG{p}{(}\PYG{n}{val\PYGZus{}loader}\PYG{p}{)}\PYG{l+s+si}{:}\PYG{l+s+s2}{.4f}\PYG{l+s+si}{\PYGZcb{}}\PYG{l+s+s2}{\PYGZdq{}}\PYG{p}{)}
\end{sphinxVerbatim}

\end{sphinxuseclass}\end{sphinxVerbatimInput}
\begin{sphinxVerbatimOutput}

\begin{sphinxuseclass}{cell_output}
\begin{sphinxVerbatim}[commandchars=\\\{\}]
Epoch [1/20]  Train MSE: 1.1341  Val MSE: 0.8154
  Best model saved (improved validation MSE).
Epoch [2/20]  Train MSE: 1.1007  Val MSE: 0.8157
\end{sphinxVerbatim}

\begin{sphinxVerbatim}[commandchars=\\\{\}]
Epoch [3/20]  Train MSE: 1.0860  Val MSE: 0.8227
Epoch [4/20]  Train MSE: 1.0493  Val MSE: 0.8578
Early stopping triggered (no improvement).

Training complete. Loading best model for final checks.
Final validation MSE (best model): 0.8154
\end{sphinxVerbatim}

\end{sphinxuseclass}\end{sphinxVerbatimOutput}

\end{sphinxuseclass}

\section{Introduction to PyTorch tensors}
\label{\detokenize{04:introduction-to-pytorch-tensors}}
\sphinxAtStartPar
\sphinxcode{\sphinxupquote{PyTorch}} is a popular open\sphinxhyphen{}source deep learning framework, and \sphinxstylestrong{tensors} are its core data structure. In PyTorch, tensors are the central data abstraction – a specialized data structure very similar to arrays or matrices, used to represent data (inputs, outputs, model parameters) in deep learning models. Tensors are conceptually like NumPy’s multi\sphinxhyphen{}dimensional arrays, but with two key advantages: they can run on GPUs (for accelerated computing), and they support automatic differentiation (through PyTorch’s autograd engine) for computing gradients. These features make tensors fundamental for building and training neural networks.

\sphinxAtStartPar
In this introduction, we’ll explore PyTorch tensors from the ground up. We assume you are comfortable with Python and NumPy, but completely new to PyTorch. We’ll cover what tensors are and how they compare to NumPy arrays, how to create and manipulate them, basic operations, how PyTorch’s autograd works for automatic differentiation, and some additional important tensor functionalities like broadcasting, using GPUs, and in\sphinxhyphen{}place operations. Along the way, you’ll find hands\sphinxhyphen{}on code examples you can run and modify to solidify your understanding.

\sphinxAtStartPar
Let’s get started by importing PyTorch (the \sphinxcode{\sphinxupquote{torch}} library) and seeing how to create tensors in various ways.

\sphinxAtStartPar
PyTorch (\sphinxcode{\sphinxupquote{torch}}) provides the \sphinxcode{\sphinxupquote{Tensor}} class, which, as mentioned, is similar to NumPy’s \sphinxcode{\sphinxupquote{ndarray}}s but with GPU support and gradient computation. If you’re familiar with NumPy arrays, you’ll find the tensor API quite intuitive. In fact, many operations (creation, indexing, math, etc.) have a very similar syntax.

\sphinxAtStartPar
One important difference is that PyTorch operations can be automatically differentiated (if required), which is essential for tasks like training neural networks.


\subsection{Creating Tensors}
\label{\detokenize{04:creating-tensors}}
\sphinxAtStartPar
There are many ways to create a tensor in PyTorch. Here are some of the most common methods:
\begin{enumerate}
\sphinxsetlistlabels{\arabic}{enumi}{enumii}{}{.}%
\item {} 
\sphinxAtStartPar
\sphinxstylestrong{From Python data}: You can create a tensor directly from a Python list or sequence of numbers using \sphinxcode{\sphinxupquote{torch.tensor()}}. PyTorch will automatically infer the data type (dtype) if not specified.

\item {} 
\sphinxAtStartPar
\sphinxstylestrong{From a NumPy array}: PyTorch can convert NumPy \sphinxcode{\sphinxupquote{ndarray}}s into tensors with \sphinxcode{\sphinxupquote{torch.from\_numpy()}}, and vice versa using the Tensor’s \sphinxcode{\sphinxupquote{.numpy()}} method. This makes it easy to move data between NumPy and PyTorch. (Notably, if the NumPy array is on CPU, the tensor and the array will share the same memory buffer, so changing one will also change the other.)

\item {} 
\sphinxAtStartPar
\sphinxstylestrong{Using factory functions}: PyTorch provides factory methods to create tensors with specific values or distributions. For example, \sphinxcode{\sphinxupquote{torch.zeros(shape)}} creates a tensor filled with zeros, \sphinxcode{\sphinxupquote{torch.ones(shape)}} creates one filled with ones, and \sphinxcode{\sphinxupquote{torch.rand(shape)}} creates one with random values uniformly sampled between 0 and 1. You can also create tensors based on an existing tensor’s properties with functions like \sphinxcode{\sphinxupquote{torch.ones\_like(existing\_tensor)}} or \sphinxcode{\sphinxupquote{torch.rand\_like(existing\_tensor)}}.

\item {} 
\sphinxAtStartPar
\sphinxstylestrong{Uninitialized tensors}: If you need an uninitialized tensor (just allocating memory without setting values), you can use \sphinxcode{\sphinxupquote{torch.empty(shape)}}.

\item {} 
\sphinxAtStartPar
\sphinxstylestrong{Specifying data types and devices}: All the above functions accept a \sphinxcode{\sphinxupquote{dtype}} parameter to specify the tensor’s data type (e.g. \sphinxcode{\sphinxupquote{torch.float32}}, \sphinxcode{\sphinxupquote{torch.int64}}, etc.), and a \sphinxcode{\sphinxupquote{device}} parameter to specify whether the tensor should reside on the CPU or GPU. By default, PyTorch tensors are created with \sphinxcode{\sphinxupquote{dtype=torch.float32}} (for floating point numbers) and on the CPU.

\end{enumerate}

\sphinxAtStartPar
Let’s go through some examples of creating tensors.

\begin{sphinxuseclass}{cell}\begin{sphinxVerbatimInput}

\begin{sphinxuseclass}{cell_input}
\begin{sphinxVerbatim}[commandchars=\\\{\}]
\PYG{c+c1}{\PYGZsh{} Creating a tensor from a Python list}
\PYG{n}{data\PYGZus{}list} \PYG{o}{=} \PYG{p}{[}\PYG{p}{[}\PYG{l+m+mi}{1}\PYG{p}{,} \PYG{l+m+mi}{2}\PYG{p}{,} \PYG{l+m+mi}{3}\PYG{p}{]}\PYG{p}{,} \PYG{p}{[}\PYG{l+m+mi}{4}\PYG{p}{,} \PYG{l+m+mi}{5}\PYG{p}{,} \PYG{l+m+mi}{6}\PYG{p}{]}\PYG{p}{]}
\PYG{n}{x\PYGZus{}data} \PYG{o}{=} \PYG{n}{torch}\PYG{o}{.}\PYG{n}{tensor}\PYG{p}{(}\PYG{n}{data\PYGZus{}list}\PYG{p}{)}
\PYG{n}{pprint}\PYG{p}{(}\PYG{l+s+s2}{\PYGZdq{}}\PYG{l+s+se}{\PYGZbs{}\PYGZbs{}}\PYG{l+s+s2}{text}\PYG{l+s+s2}{\PYGZob{}}\PYG{l+s+s2}{Tensor from list:\PYGZcb{}}\PYG{l+s+s2}{\PYGZdq{}}\PYG{p}{,} \PYG{n}{x\PYGZus{}data}\PYG{p}{)}

\PYG{c+c1}{\PYGZsh{} Creating a tensor from a NumPy array}
\PYG{n}{np\PYGZus{}array} \PYG{o}{=} \PYG{n}{np}\PYG{o}{.}\PYG{n}{array}\PYG{p}{(}\PYG{n}{data\PYGZus{}list}\PYG{p}{)}
\PYG{n}{x\PYGZus{}np} \PYG{o}{=} \PYG{n}{torch}\PYG{o}{.}\PYG{n}{from\PYGZus{}numpy}\PYG{p}{(}\PYG{n}{np\PYGZus{}array}\PYG{p}{)}
\PYG{n}{pprint}\PYG{p}{(}\PYG{l+s+s2}{\PYGZdq{}}\PYG{l+s+se}{\PYGZbs{}\PYGZbs{}}\PYG{l+s+s2}{text}\PYG{l+s+s2}{\PYGZob{}}\PYG{l+s+s2}{Tensor from NumPy array:\PYGZcb{}}\PYG{l+s+s2}{\PYGZdq{}}\PYG{p}{,} \PYG{n}{x\PYGZus{}np}\PYG{p}{)}

\PYG{c+c1}{\PYGZsh{} Verify that modifying the NumPy array changes the tensor (shared memory)}
\PYG{n}{np\PYGZus{}array}\PYG{p}{[}\PYG{l+m+mi}{0}\PYG{p}{,} \PYG{l+m+mi}{0}\PYG{p}{]} \PYG{o}{=} \PYG{l+m+mi}{99}
\PYG{n}{pprint}\PYG{p}{(}\PYG{l+s+s2}{\PYGZdq{}}\PYG{l+s+se}{\PYGZbs{}\PYGZbs{}}\PYG{l+s+s2}{text}\PYG{l+s+s2}{\PYGZob{}}\PYG{l+s+s2}{After modifying the NumPy array, the Tensor becomes:\PYGZcb{}}\PYG{l+s+s2}{\PYGZdq{}}\PYG{p}{,} \PYG{n}{x\PYGZus{}np}\PYG{p}{)}

\PYG{c+c1}{\PYGZsh{} Creating tensors using factory functions}
\PYG{n}{shape} \PYG{o}{=} \PYG{p}{(}\PYG{l+m+mi}{2}\PYG{p}{,} \PYG{l+m+mi}{3}\PYG{p}{)}
\PYG{n}{x\PYGZus{}zeros} \PYG{o}{=} \PYG{n}{torch}\PYG{o}{.}\PYG{n}{zeros}\PYG{p}{(}\PYG{n}{shape}\PYG{p}{)}
\PYG{n}{x\PYGZus{}ones} \PYG{o}{=} \PYG{n}{torch}\PYG{o}{.}\PYG{n}{ones}\PYG{p}{(}\PYG{n}{shape}\PYG{p}{)}
\PYG{n}{x\PYGZus{}rand} \PYG{o}{=} \PYG{n}{torch}\PYG{o}{.}\PYG{n}{rand}\PYG{p}{(}\PYG{n}{shape}\PYG{p}{)}
\PYG{n}{pprint}\PYG{p}{(}\PYG{l+s+s2}{\PYGZdq{}}\PYG{l+s+se}{\PYGZbs{}\PYGZbs{}}\PYG{l+s+s2}{text}\PYG{l+s+s2}{\PYGZob{}}\PYG{l+s+s2}{Zero Tensor:\PYGZcb{}}\PYG{l+s+s2}{\PYGZdq{}}\PYG{p}{,} \PYG{n}{x\PYGZus{}zeros}\PYG{p}{)}
\PYG{n}{pprint}\PYG{p}{(}\PYG{l+s+s2}{\PYGZdq{}}\PYG{l+s+se}{\PYGZbs{}\PYGZbs{}}\PYG{l+s+s2}{text}\PYG{l+s+s2}{\PYGZob{}}\PYG{l+s+s2}{Ones Tensor:\PYGZcb{}}\PYG{l+s+s2}{\PYGZdq{}}\PYG{p}{,} \PYG{n}{x\PYGZus{}ones}\PYG{p}{)}
\PYG{n}{pprint}\PYG{p}{(}\PYG{l+s+s2}{\PYGZdq{}}\PYG{l+s+se}{\PYGZbs{}\PYGZbs{}}\PYG{l+s+s2}{text}\PYG{l+s+s2}{\PYGZob{}}\PYG{l+s+s2}{Random Tensor:\PYGZcb{}}\PYG{l+s+s2}{\PYGZdq{}}\PYG{p}{,} \PYG{n}{x\PYGZus{}rand}\PYG{p}{)}

\PYG{c+c1}{\PYGZsh{} Creating tensors based on existing ones}
\PYG{n}{x\PYGZus{}like} \PYG{o}{=} \PYG{n}{torch}\PYG{o}{.}\PYG{n}{ones\PYGZus{}like}\PYG{p}{(}\PYG{n}{x\PYGZus{}rand}\PYG{p}{)}        \PYG{c+c1}{\PYGZsh{} tensor of ones with same shape as x\PYGZus{}rand}
\PYG{n}{x\PYGZus{}like\PYGZus{}float} \PYG{o}{=} \PYG{n}{torch}\PYG{o}{.}\PYG{n}{rand\PYGZus{}like}\PYG{p}{(}\PYG{n}{x\PYGZus{}data}\PYG{p}{,} \PYG{n}{dtype}\PYG{o}{=}\PYG{n}{torch}\PYG{o}{.}\PYG{n}{float32}\PYG{p}{)}  \PYG{c+c1}{\PYGZsh{} random tensor with same shape as x\PYGZus{}data, but forced float type}
\PYG{n}{pprint}\PYG{p}{(}\PYG{l+s+s2}{\PYGZdq{}}\PYG{l+s+se}{\PYGZbs{}\PYGZbs{}}\PYG{l+s+s2}{text}\PYG{l+s+s2}{\PYGZob{}}\PYG{l+s+s2}{Ones\PYGZhy{}like Tensor:\PYGZcb{}}\PYG{l+s+s2}{\PYGZdq{}}\PYG{p}{,} \PYG{n}{x\PYGZus{}like}\PYG{p}{)}
\PYG{n}{pprint}\PYG{p}{(}\PYG{l+s+s2}{\PYGZdq{}}\PYG{l+s+se}{\PYGZbs{}\PYGZbs{}}\PYG{l+s+s2}{text}\PYG{l+s+s2}{\PYGZob{}}\PYG{l+s+s2}{Rand\PYGZhy{}like Tensor with float32 dtype:\PYGZcb{}}\PYG{l+s+s2}{\PYGZdq{}}\PYG{p}{,} \PYG{n}{x\PYGZus{}like\PYGZus{}float}\PYG{p}{)}

\PYG{c+c1}{\PYGZsh{} Checking tensor attributes}
\PYG{n+nb}{print}\PYG{p}{(}\PYG{l+s+s2}{\PYGZdq{}}\PYG{l+s+s2}{Tensor shape:}\PYG{l+s+s2}{\PYGZdq{}}\PYG{p}{,} \PYG{n}{x\PYGZus{}data}\PYG{o}{.}\PYG{n}{shape}\PYG{p}{)}
\PYG{n+nb}{print}\PYG{p}{(}\PYG{l+s+s2}{\PYGZdq{}}\PYG{l+s+s2}{Tensor data type:}\PYG{l+s+s2}{\PYGZdq{}}\PYG{p}{,} \PYG{n}{x\PYGZus{}data}\PYG{o}{.}\PYG{n}{dtype}\PYG{p}{)}
\PYG{n+nb}{print}\PYG{p}{(}\PYG{l+s+s2}{\PYGZdq{}}\PYG{l+s+s2}{Tensor device:}\PYG{l+s+s2}{\PYGZdq{}}\PYG{p}{,} \PYG{n}{x\PYGZus{}data}\PYG{o}{.}\PYG{n}{device}\PYG{p}{)}
\end{sphinxVerbatim}

\end{sphinxuseclass}\end{sphinxVerbatimInput}
\begin{sphinxVerbatimOutput}

\begin{sphinxuseclass}{cell_output}\begin{equation*}
\begin{split}\displaystyle \text{Tensor from list:}\left[
\begin{array}{}
  1.0000 &  2.0000 &  3.0000\\
  4.0000 &  5.0000 &  6.0000
\end{array}
\right]\end{split}
\end{equation*}\begin{equation*}
\begin{split}\displaystyle \text{Tensor from NumPy array:}\left[
\begin{array}{}
  1.0000 &  2.0000 &  3.0000\\
  4.0000 &  5.0000 &  6.0000
\end{array}
\right]\end{split}
\end{equation*}\begin{equation*}
\begin{split}\displaystyle \text{After modifying the NumPy array, the Tensor becomes:}\left[
\begin{array}{}
  99.0000 &  2.0000 &  3.0000\\
  4.0000 &  5.0000 &  6.0000
\end{array}
\right]\end{split}
\end{equation*}\begin{equation*}
\begin{split}\displaystyle \text{Zero Tensor:}\left[
\begin{array}{}
  0.0000 &  0.0000 &  0.0000\\
  0.0000 &  0.0000 &  0.0000
\end{array}
\right]\end{split}
\end{equation*}\begin{equation*}
\begin{split}\displaystyle \text{Ones Tensor:}\left[
\begin{array}{}
  1.0000 &  1.0000 &  1.0000\\
  1.0000 &  1.0000 &  1.0000
\end{array}
\right]\end{split}
\end{equation*}\begin{equation*}
\begin{split}\displaystyle \text{Random Tensor:}\left[
\begin{array}{}
  0.1394 &  0.5029 &  0.2952\\
  0.5932 &  0.1593 &  0.1439
\end{array}
\right]\end{split}
\end{equation*}\begin{equation*}
\begin{split}\displaystyle \text{Ones-like Tensor:}\left[
\begin{array}{}
  1.0000 &  1.0000 &  1.0000\\
  1.0000 &  1.0000 &  1.0000
\end{array}
\right]\end{split}
\end{equation*}\begin{equation*}
\begin{split}\displaystyle \text{Rand-like Tensor with float32 dtype:}\left[
\begin{array}{}
  0.9453 &  0.5497 &  0.7615\\
  0.4744 &  0.6142 &  0.6309
\end{array}
\right]\end{split}
\end{equation*}
\begin{sphinxVerbatim}[commandchars=\\\{\}]
Tensor shape: torch.Size([2, 3])
Tensor data type: torch.int64
Tensor device: cpu
\end{sphinxVerbatim}

\end{sphinxuseclass}\end{sphinxVerbatimOutput}

\end{sphinxuseclass}

\subsection{Basic Tensor Operations}
\label{\detokenize{04:basic-tensor-operations}}
\sphinxAtStartPar
Once you have tensors, PyTorch provides a rich set of operations to manipulate them. In fact, there are over 100 tensor operations available, including arithmetic, linear algebra, slicing, reshaping, and more. Many of these mirror NumPy operations in syntax and function. Let’s go through some fundamental operations.


\subsubsection{Arithmetic operations}
\label{\detokenize{04:arithmetic-operations}}
\sphinxAtStartPar
Arithmetic on tensors is typically element\sphinxhyphen{}wise (just like with NumPy arrays). You can use Python arithmetic operators (\sphinxcode{\sphinxupquote{+}}, \sphinxcode{\sphinxupquote{\sphinxhyphen{}}}, \sphinxcode{\sphinxupquote{*}}, \sphinxcode{\sphinxupquote{/}}) or the equivalent \sphinxcode{\sphinxupquote{torch}} functions (\sphinxcode{\sphinxupquote{torch.add}}, \sphinxcode{\sphinxupquote{torch.sub}}, etc.). For example:

\begin{sphinxuseclass}{cell}\begin{sphinxVerbatimInput}

\begin{sphinxuseclass}{cell_input}
\begin{sphinxVerbatim}[commandchars=\\\{\}]
\PYG{c+c1}{\PYGZsh{} Element\PYGZhy{}wise arithmetic operations}
\PYG{n}{a} \PYG{o}{=} \PYG{n}{torch}\PYG{o}{.}\PYG{n}{tensor}\PYG{p}{(}\PYG{p}{[}\PYG{l+m+mf}{1.0}\PYG{p}{,} \PYG{l+m+mf}{2.0}\PYG{p}{,} \PYG{l+m+mf}{3.0}\PYG{p}{]}\PYG{p}{)}
\PYG{n}{b} \PYG{o}{=} \PYG{n}{torch}\PYG{o}{.}\PYG{n}{tensor}\PYG{p}{(}\PYG{p}{[}\PYG{l+m+mf}{4.0}\PYG{p}{,} \PYG{l+m+mf}{5.0}\PYG{p}{,} \PYG{l+m+mf}{6.0}\PYG{p}{]}\PYG{p}{)}

\PYG{n}{pprint}\PYG{p}{(}\PYG{l+s+s2}{\PYGZdq{}}\PYG{l+s+s2}{a=}\PYG{l+s+s2}{\PYGZdq{}}\PYG{p}{,} \PYG{n}{a}\PYG{p}{)}
\PYG{n}{pprint}\PYG{p}{(}\PYG{l+s+s2}{\PYGZdq{}}\PYG{l+s+s2}{b=}\PYG{l+s+s2}{\PYGZdq{}}\PYG{p}{,} \PYG{n}{b}\PYG{p}{)}
\PYG{n}{pprint}\PYG{p}{(}\PYG{l+s+s2}{\PYGZdq{}}\PYG{l+s+s2}{a + b =}\PYG{l+s+s2}{\PYGZdq{}}\PYG{p}{,} \PYG{n}{a} \PYG{o}{+} \PYG{n}{b}\PYG{p}{)}        \PYG{c+c1}{\PYGZsh{} element\PYGZhy{}wise addition}
\PYG{n}{pprint}\PYG{p}{(}\PYG{l+s+s2}{\PYGZdq{}}\PYG{l+s+s2}{a }\PYG{l+s+se}{\PYGZbs{}\PYGZbs{}}\PYG{l+s+s2}{times b =}\PYG{l+s+s2}{\PYGZdq{}}\PYG{p}{,} \PYG{n}{a} \PYG{o}{*} \PYG{n}{b}\PYG{p}{)}        \PYG{c+c1}{\PYGZsh{} element\PYGZhy{}wise multiplication}
\PYG{n}{pprint}\PYG{p}{(}\PYG{l+s+s2}{\PYGZdq{}}\PYG{l+s+s2}{a \PYGZhy{} b =}\PYG{l+s+s2}{\PYGZdq{}}\PYG{p}{,} \PYG{n}{a} \PYG{o}{\PYGZhy{}} \PYG{n}{b}\PYG{p}{)}        \PYG{c+c1}{\PYGZsh{} element\PYGZhy{}wise subtraction}
\PYG{n}{pprint}\PYG{p}{(}\PYG{l+s+s2}{\PYGZdq{}}\PYG{l+s+se}{\PYGZbs{}\PYGZbs{}}\PYG{l+s+s2}{frac}\PYG{l+s+si}{\PYGZob{}a\PYGZcb{}}\PYG{l+s+si}{\PYGZob{}b\PYGZcb{}}\PYG{l+s+s2}{ =}\PYG{l+s+s2}{\PYGZdq{}}\PYG{p}{,} \PYG{n}{a} \PYG{o}{/} \PYG{n}{b}\PYG{p}{)}        \PYG{c+c1}{\PYGZsh{} element\PYGZhy{}wise division}
\end{sphinxVerbatim}

\end{sphinxuseclass}\end{sphinxVerbatimInput}
\begin{sphinxVerbatimOutput}

\begin{sphinxuseclass}{cell_output}\begin{equation*}
\begin{split}\displaystyle a=\left[
\begin{array}{}
  1.0000 &  2.0000 &  3.0000
\end{array}
\right]\end{split}
\end{equation*}\begin{equation*}
\begin{split}\displaystyle b=\left[
\begin{array}{}
  4.0000 &  5.0000 &  6.0000
\end{array}
\right]\end{split}
\end{equation*}\begin{equation*}
\begin{split}\displaystyle a + b =\left[
\begin{array}{}
  5.0000 &  7.0000 &  9.0000
\end{array}
\right]\end{split}
\end{equation*}\begin{equation*}
\begin{split}\displaystyle a \times b =\left[
\begin{array}{}
  4.0000 &  10.0000 &  18.0000
\end{array}
\right]\end{split}
\end{equation*}\begin{equation*}
\begin{split}\displaystyle a - b =\left[
\begin{array}{}
 -3.0000 & -3.0000 & -3.0000
\end{array}
\right]\end{split}
\end{equation*}\begin{equation*}
\begin{split}\displaystyle \frac{a}{b} =\left[
\begin{array}{}
  0.2500 &  0.4000 &  0.5000
\end{array}
\right]\end{split}
\end{equation*}
\end{sphinxuseclass}\end{sphinxVerbatimOutput}

\end{sphinxuseclass}
\sphinxAtStartPar
The results should be as expected:
\begin{itemize}
\item {} 
\sphinxAtStartPar
\sphinxcode{\sphinxupquote{a + b}} gives \sphinxcode{\sphinxupquote{{[}5.0, 7.0, 9.0{]}}}

\item {} 
\sphinxAtStartPar
\sphinxcode{\sphinxupquote{a * b}} gives \sphinxcode{\sphinxupquote{{[}4.0, 10.0, 18.0{]}}}

\item {} 
\sphinxAtStartPar
etc.

\end{itemize}

\sphinxAtStartPar
If you want to do \sphinxstylestrong{matrix multiplication} (dot products, etc.), you should use the \sphinxcode{\sphinxupquote{@}} operator or \sphinxcode{\sphinxupquote{torch.matmul}}. For example, \sphinxcode{\sphinxupquote{A @ B}} does matrix multiply if \sphinxcode{\sphinxupquote{A}} and \sphinxcode{\sphinxupquote{B}} are 2\sphinxhyphen{}D matrices (or higher\sphinxhyphen{}dim, following broadcast batch rules). For 1\sphinxhyphen{}D vectors, \sphinxcode{\sphinxupquote{a @ b}} will give a scalar dot product. Let’s see a quick example:

\begin{sphinxuseclass}{cell}\begin{sphinxVerbatimInput}

\begin{sphinxuseclass}{cell_input}
\begin{sphinxVerbatim}[commandchars=\\\{\}]
\PYG{c+c1}{\PYGZsh{} Matrix multiplication vs elementwise}
\PYG{n}{A} \PYG{o}{=} \PYG{n}{torch}\PYG{o}{.}\PYG{n}{tensor}\PYG{p}{(}\PYG{p}{[}\PYG{p}{[}\PYG{l+m+mi}{1}\PYG{p}{,} \PYG{l+m+mi}{2}\PYG{p}{]}\PYG{p}{,}
                  \PYG{p}{[}\PYG{l+m+mi}{3}\PYG{p}{,} \PYG{l+m+mi}{4}\PYG{p}{]}\PYG{p}{]}\PYG{p}{)}
\PYG{n}{B} \PYG{o}{=} \PYG{n}{torch}\PYG{o}{.}\PYG{n}{tensor}\PYG{p}{(}\PYG{p}{[}\PYG{p}{[}\PYG{l+m+mi}{5}\PYG{p}{,} \PYG{l+m+mi}{6}\PYG{p}{]}\PYG{p}{,}
                  \PYG{p}{[}\PYG{l+m+mi}{7}\PYG{p}{,} \PYG{l+m+mi}{8}\PYG{p}{]}\PYG{p}{]}\PYG{p}{)}
\PYG{n}{pprint}\PYG{p}{(}\PYG{l+s+s2}{\PYGZdq{}}\PYG{l+s+s2}{A=}\PYG{l+s+s2}{\PYGZdq{}}\PYG{p}{,} \PYG{n}{A}\PYG{p}{)}
\PYG{n}{pprint}\PYG{p}{(}\PYG{l+s+s2}{\PYGZdq{}}\PYG{l+s+s2}{B=}\PYG{l+s+s2}{\PYGZdq{}}\PYG{p}{,} \PYG{n}{B}\PYG{p}{)}
\PYG{n}{pprint}\PYG{p}{(}\PYG{l+s+s2}{\PYGZdq{}}\PYG{l+s+se}{\PYGZbs{}\PYGZbs{}}\PYG{l+s+s2}{text}\PYG{l+s+s2}{\PYGZob{}}\PYG{l+s+s2}{Element\PYGZhy{}wise multiply:\PYGZcb{}}\PYG{l+s+s2}{\PYGZdq{}}\PYG{p}{,} \PYG{n}{A} \PYG{o}{*} \PYG{n}{B}\PYG{p}{)}   \PYG{c+c1}{\PYGZsh{} same shape \PYGZhy{}\PYGZgt{} elementwise product}
\PYG{n}{pprint}\PYG{p}{(}\PYG{l+s+s2}{\PYGZdq{}}\PYG{l+s+se}{\PYGZbs{}\PYGZbs{}}\PYG{l+s+s2}{text}\PYG{l+s+s2}{\PYGZob{}}\PYG{l+s+s2}{Matrix multiply:\PYGZcb{}}\PYG{l+s+s2}{\PYGZdq{}}\PYG{p}{,} \PYG{n}{A} \PYG{o}{@} \PYG{n}{B}\PYG{p}{)}         \PYG{c+c1}{\PYGZsh{} 2x2 @ 2x2 \PYGZhy{}\PYGZgt{} 2x2 matrix result}
\end{sphinxVerbatim}

\end{sphinxuseclass}\end{sphinxVerbatimInput}
\begin{sphinxVerbatimOutput}

\begin{sphinxuseclass}{cell_output}\begin{equation*}
\begin{split}\displaystyle A=\left[
\begin{array}{}
  1.0000 &  2.0000\\
  3.0000 &  4.0000
\end{array}
\right]\end{split}
\end{equation*}\begin{equation*}
\begin{split}\displaystyle B=\left[
\begin{array}{}
  5.0000 &  6.0000\\
  7.0000 &  8.0000
\end{array}
\right]\end{split}
\end{equation*}\begin{equation*}
\begin{split}\displaystyle \text{Element-wise multiply:}\left[
\begin{array}{}
  5.0000 &  12.0000\\
  21.0000 &  32.0000
\end{array}
\right]\end{split}
\end{equation*}\begin{equation*}
\begin{split}\displaystyle \text{Matrix multiply:}\left[
\begin{array}{}
  19.0000 &  22.0000\\
  43.0000 &  50.0000
\end{array}
\right]\end{split}
\end{equation*}
\end{sphinxuseclass}\end{sphinxVerbatimOutput}

\end{sphinxuseclass}
\sphinxAtStartPar
Here, \sphinxcode{\sphinxupquote{A * B}} will do element\sphinxhyphen{}wise multiplication (resulting in \(\begin{bmatrix}5&12\\21&32\end{bmatrix}\)), whereas \sphinxcode{\sphinxupquote{A @ B}} will compute the matrix product (result \(\begin{bmatrix}19&22\\43&50\end{bmatrix}\)).


\subsubsection{Indexing and slicing}
\label{\detokenize{04:indexing-and-slicing}}
\sphinxAtStartPar
You can access and modify parts of tensors using indexing and slicing, just like with NumPy. PyTorch uses 0\sphinxhyphen{}based indexing. You can use indices, ranges, colons, and even Boolean masks or index tensors (though advanced indexing is outside our scope here). Some examples:

\begin{sphinxuseclass}{cell}\begin{sphinxVerbatimInput}

\begin{sphinxuseclass}{cell_input}
\begin{sphinxVerbatim}[commandchars=\\\{\}]
\PYG{c+c1}{\PYGZsh{} Create a 2\PYGZhy{}D tensor (matrix) for demonstration}
\PYG{n}{x} \PYG{o}{=} \PYG{n}{torch}\PYG{o}{.}\PYG{n}{tensor}\PYG{p}{(}\PYG{p}{[}\PYG{p}{[}\PYG{l+m+mi}{10}\PYG{p}{,} \PYG{l+m+mi}{11}\PYG{p}{,} \PYG{l+m+mi}{12}\PYG{p}{]}\PYG{p}{,}
                  \PYG{p}{[}\PYG{l+m+mi}{20}\PYG{p}{,} \PYG{l+m+mi}{21}\PYG{p}{,} \PYG{l+m+mi}{22}\PYG{p}{]}\PYG{p}{,}
                  \PYG{p}{[}\PYG{l+m+mi}{30}\PYG{p}{,} \PYG{l+m+mi}{31}\PYG{p}{,} \PYG{l+m+mi}{32}\PYG{p}{]}\PYG{p}{]}\PYG{p}{)}
\PYG{n}{pprint}\PYG{p}{(}\PYG{l+s+s2}{\PYGZdq{}}\PYG{l+s+se}{\PYGZbs{}\PYGZbs{}}\PYG{l+s+s2}{text}\PYG{l+s+s2}{\PYGZob{}}\PYG{l+s+s2}{Original tensor: \PYGZcb{}}\PYG{l+s+s2}{\PYGZdq{}}\PYG{p}{,} \PYG{n}{x}\PYG{p}{)}

\PYG{c+c1}{\PYGZsh{} Indexing a single element}
\PYG{n}{pprint}\PYG{p}{(}\PYG{l+s+s2}{\PYGZdq{}}\PYG{l+s+se}{\PYGZbs{}\PYGZbs{}}\PYG{l+s+s2}{text}\PYG{l+s+s2}{\PYGZob{}}\PYG{l+s+s2}{Element at (0, 1): \PYGZcb{}}\PYG{l+s+s2}{\PYGZdq{}}\PYG{p}{,} \PYG{n+nb}{str}\PYG{p}{(}\PYG{n}{x}\PYG{p}{[}\PYG{l+m+mi}{0}\PYG{p}{,} \PYG{l+m+mi}{1}\PYG{p}{]}\PYG{o}{.}\PYG{n}{item}\PYG{p}{(}\PYG{p}{)}\PYG{p}{)}\PYG{p}{)}   \PYG{c+c1}{\PYGZsh{} .item() to get Python scalar}

\PYG{c+c1}{\PYGZsh{} Slicing rows and columns}
\PYG{n}{pprint}\PYG{p}{(}\PYG{l+s+s2}{\PYGZdq{}}\PYG{l+s+se}{\PYGZbs{}\PYGZbs{}}\PYG{l+s+s2}{text}\PYG{l+s+s2}{\PYGZob{}}\PYG{l+s+s2}{First row: \PYGZcb{}}\PYG{l+s+s2}{\PYGZdq{}}\PYG{p}{,} \PYG{n}{x}\PYG{p}{[}\PYG{l+m+mi}{0}\PYG{p}{]}\PYG{p}{)}            \PYG{c+c1}{\PYGZsh{} equivalent to x[0, :]}
\PYG{n}{pprint}\PYG{p}{(}\PYG{l+s+s2}{\PYGZdq{}}\PYG{l+s+se}{\PYGZbs{}\PYGZbs{}}\PYG{l+s+s2}{text}\PYG{l+s+s2}{\PYGZob{}}\PYG{l+s+s2}{Last column: \PYGZcb{}}\PYG{l+s+s2}{\PYGZdq{}}\PYG{p}{,} \PYG{n}{x}\PYG{p}{[}\PYG{p}{:}\PYG{p}{,} \PYG{o}{\PYGZhy{}}\PYG{l+m+mi}{1}\PYG{p}{]}\PYG{p}{)}       \PYG{c+c1}{\PYGZsh{} all rows, last column}

\PYG{c+c1}{\PYGZsh{} Slice a submatrix (e.g., top\PYGZhy{}left 2x2 block)}
\PYG{n}{sub\PYGZus{}x} \PYG{o}{=} \PYG{n}{x}\PYG{p}{[}\PYG{l+m+mi}{0}\PYG{p}{:}\PYG{l+m+mi}{2}\PYG{p}{,} \PYG{l+m+mi}{0}\PYG{p}{:}\PYG{l+m+mi}{2}\PYG{p}{]}   \PYG{c+c1}{\PYGZsh{} rows 0\PYGZhy{}1 and cols 0\PYGZhy{}1}
\PYG{n}{pprint}\PYG{p}{(}\PYG{l+s+s2}{\PYGZdq{}}\PYG{l+s+se}{\PYGZbs{}\PYGZbs{}}\PYG{l+s+s2}{text}\PYG{l+s+s2}{\PYGZob{}}\PYG{l+s+s2}{Top\PYGZhy{}left 2x2 sub\PYGZhy{}tensor: \PYGZcb{}}\PYG{l+s+s2}{\PYGZdq{}}\PYG{p}{,} \PYG{n}{sub\PYGZus{}x}\PYG{p}{)}

\PYG{c+c1}{\PYGZsh{} Modify part of the tensor}
\PYG{n}{x}\PYG{p}{[}\PYG{l+m+mi}{1}\PYG{p}{:}\PYG{p}{,} \PYG{l+m+mi}{1}\PYG{p}{:}\PYG{p}{]} \PYG{o}{=} \PYG{n}{torch}\PYG{o}{.}\PYG{n}{tensor}\PYG{p}{(}\PYG{p}{[}\PYG{p}{[} \PYG{o}{\PYGZhy{}}\PYG{l+m+mi}{1}\PYG{p}{,} \PYG{o}{\PYGZhy{}}\PYG{l+m+mi}{2}\PYG{p}{]}\PYG{p}{,}
                          \PYG{p}{[} \PYG{o}{\PYGZhy{}}\PYG{l+m+mi}{3}\PYG{p}{,} \PYG{o}{\PYGZhy{}}\PYG{l+m+mi}{4}\PYG{p}{]}\PYG{p}{]}\PYG{p}{)}   \PYG{c+c1}{\PYGZsh{} set bottom\PYGZhy{}right 2x2 block to new values}
\PYG{n}{pprint}\PYG{p}{(}\PYG{l+s+s2}{\PYGZdq{}}\PYG{l+s+se}{\PYGZbs{}\PYGZbs{}}\PYG{l+s+s2}{text}\PYG{l+s+s2}{\PYGZob{}}\PYG{l+s+s2}{Tensor after modification: \PYGZcb{}}\PYG{l+s+s2}{\PYGZdq{}}\PYG{p}{,} \PYG{n}{x}\PYG{p}{)}
\end{sphinxVerbatim}

\end{sphinxuseclass}\end{sphinxVerbatimInput}
\begin{sphinxVerbatimOutput}

\begin{sphinxuseclass}{cell_output}\begin{equation*}
\begin{split}\displaystyle \text{Original tensor: }\left[
\begin{array}{}
  10.0000 &  11.0000 &  12.0000\\
  20.0000 &  21.0000 &  22.0000\\
  30.0000 &  31.0000 &  32.0000
\end{array}
\right]\end{split}
\end{equation*}\begin{equation*}
\begin{split}\displaystyle \text{Element at (0, 1): }11\end{split}
\end{equation*}\begin{equation*}
\begin{split}\displaystyle \text{First row: }\left[
\begin{array}{}
  10.0000 &  11.0000 &  12.0000
\end{array}
\right]\end{split}
\end{equation*}\begin{equation*}
\begin{split}\displaystyle \text{Last column: }\left[
\begin{array}{}
  12.0000 &  22.0000 &  32.0000
\end{array}
\right]\end{split}
\end{equation*}\begin{equation*}
\begin{split}\displaystyle \text{Top-left 2x2 sub-tensor: }\left[
\begin{array}{}
  10.0000 &  11.0000\\
  20.0000 &  21.0000
\end{array}
\right]\end{split}
\end{equation*}\begin{equation*}
\begin{split}\displaystyle \text{Tensor after modification: }\left[
\begin{array}{}
  10.0000 &  11.0000 &  12.0000\\
  20.0000 & -1.0000 & -2.0000\\
  30.0000 & -3.0000 & -4.0000
\end{array}
\right]\end{split}
\end{equation*}
\end{sphinxuseclass}\end{sphinxVerbatimOutput}

\end{sphinxuseclass}

\paragraph{In this code:}
\label{\detokenize{04:in-this-code}}\begin{itemize}
\item {} 
\sphinxAtStartPar
We created a 3×3 tensor \sphinxcode{\sphinxupquote{x}} with some distinct values.

\item {} 
\sphinxAtStartPar
We accessed a single element \sphinxcode{\sphinxupquote{x{[}0, 1{]}}} (which should be \sphinxcode{\sphinxupquote{11}} in the original tensor). Using \sphinxcode{\sphinxupquote{.item()}} converts the 0\sphinxhyphen{}dim tensor to a Python number.

\item {} 
\sphinxAtStartPar
We took the first row \sphinxcode{\sphinxupquote{x{[}0{]}}} (which yields \sphinxcode{\sphinxupquote{{[}10, 11, 12{]}}}) and the last column \sphinxcode{\sphinxupquote{x{[}:, \sphinxhyphen{}1{]}}} (which yields \sphinxcode{\sphinxupquote{{[}12, 22, 32{]}}}).

\item {} 
\sphinxAtStartPar
We sliced a 2×2 sub\sphinxhyphen{}tensor (the upper\sphinxhyphen{}left corner).

\item {} 
\sphinxAtStartPar
We then modified a sub\sphinxhyphen{}section of the tensor (\sphinxcode{\sphinxupquote{x{[}1:, 1:{]}}}) by assigning a new tensor of the same shape. The original tensor \sphinxcode{\sphinxupquote{x}} is changed in place.

\end{itemize}

\sphinxAtStartPar
Note: Slicing a tensor (e.g., \sphinxcode{\sphinxupquote{y = x{[}0:2, 0:2{]}}}) will result in a \sphinxstylestrong{view} of the original tensor whenever possible (i.e., it doesn’t copy the data). So modifying \sphinxcode{\sphinxupquote{y}} would also modify \sphinxcode{\sphinxupquote{x}}. If you need to explicitly copy a tensor or sub\sphinxhyphen{}tensor, you can use \sphinxcode{\sphinxupquote{.clone()}}.


\subsubsection{Reshaping and concatenation}
\label{\detokenize{04:reshaping-and-concatenation}}
\sphinxAtStartPar
Often you’ll want to change the shape or dimensionality of a tensor without changing its data. This is done with operations like \sphinxcode{\sphinxupquote{tensor.view()}} or \sphinxcode{\sphinxupquote{tensor.reshape()}}. Both do similar things (though \sphinxcode{\sphinxupquote{.view()}} only works on contiguous tensors). For simplicity, we’ll use \sphinxcode{\sphinxupquote{reshape}} here.

\sphinxAtStartPar
You can also join tensors together (concatenate) along a given dimension using \sphinxcode{\sphinxupquote{torch.cat}} or stack them with \sphinxcode{\sphinxupquote{torch.stack}} (which adds a new dimension). Examples:

\begin{sphinxuseclass}{cell}\begin{sphinxVerbatimInput}

\begin{sphinxuseclass}{cell_input}
\begin{sphinxVerbatim}[commandchars=\\\{\}]
\PYG{c+c1}{\PYGZsh{} Reshape (view) a tensor}
\PYG{n}{t} \PYG{o}{=} \PYG{n}{torch}\PYG{o}{.}\PYG{n}{arange}\PYG{p}{(}\PYG{l+m+mi}{16}\PYG{p}{)}        \PYG{c+c1}{\PYGZsh{} 1\PYGZhy{}D tensor [0, 1, 2, ..., 15]}
\PYG{n}{pprint}\PYG{p}{(}\PYG{l+s+s2}{\PYGZdq{}}\PYG{l+s+s2}{t=}\PYG{l+s+s2}{\PYGZdq{}}\PYG{p}{,} \PYG{n}{t}\PYG{p}{)}

\PYG{n}{t\PYGZus{}matrix} \PYG{o}{=} \PYG{n}{t}\PYG{o}{.}\PYG{n}{reshape}\PYG{p}{(}\PYG{l+m+mi}{4}\PYG{p}{,} \PYG{l+m+mi}{4}\PYG{p}{)}  \PYG{c+c1}{\PYGZsh{} reshape to 4x4 matrix}
\PYG{n}{pprint}\PYG{p}{(}\PYG{l+s+s2}{\PYGZdq{}}\PYG{l+s+se}{\PYGZbs{}\PYGZbs{}}\PYG{l+s+s2}{text}\PYG{l+s+s2}{\PYGZob{}}\PYG{l+s+s2}{t reshaped to 4x4: \PYGZcb{}}\PYG{l+s+s2}{\PYGZdq{}}\PYG{p}{,} \PYG{n}{t\PYGZus{}matrix}\PYG{p}{)}

\PYG{c+c1}{\PYGZsh{} Transpose the matrix (swap dimensions)}
\PYG{n}{pprint}\PYG{p}{(}\PYG{l+s+s2}{\PYGZdq{}}\PYG{l+s+s2}{t\PYGZca{}T:}\PYG{l+s+s2}{\PYGZbs{}}\PYG{l+s+s2}{ }\PYG{l+s+s2}{\PYGZdq{}}\PYG{p}{,} \PYG{n}{t\PYGZus{}matrix}\PYG{o}{.}\PYG{n}{T}\PYG{p}{)}
\end{sphinxVerbatim}

\end{sphinxuseclass}\end{sphinxVerbatimInput}
\begin{sphinxVerbatimOutput}

\begin{sphinxuseclass}{cell_output}\begin{equation*}
\begin{split}\displaystyle t=\left[
\begin{array}{}
  0.0000 &  1.0000 &  2.0000 &  3.0000 &  4.0000 &  5.0000 &  6.0000 &  7.0000 &  8.0000 &  9.0000 &  10.0000 &  11.0000 &  12.0000 &  13.0000 &  14.0000 &  15.0000
\end{array}
\right]\end{split}
\end{equation*}\begin{equation*}
\begin{split}\displaystyle \text{t reshaped to 4x4: }\left[
\begin{array}{}
  0.0000 &  1.0000 &  2.0000 &  3.0000\\
  4.0000 &  5.0000 &  6.0000 &  7.0000\\
  8.0000 &  9.0000 &  10.0000 &  11.0000\\
  12.0000 &  13.0000 &  14.0000 &  15.0000
\end{array}
\right]\end{split}
\end{equation*}\begin{equation*}
\begin{split}\displaystyle t^T:\ \left[
\begin{array}{}
  0.0000 &  4.0000 &  8.0000 &  12.0000\\
  1.0000 &  5.0000 &  9.0000 &  13.0000\\
  2.0000 &  6.0000 &  10.0000 &  14.0000\\
  3.0000 &  7.0000 &  11.0000 &  15.0000
\end{array}
\right]\end{split}
\end{equation*}
\end{sphinxuseclass}\end{sphinxVerbatimOutput}

\end{sphinxuseclass}
\begin{sphinxuseclass}{cell}\begin{sphinxVerbatimInput}

\begin{sphinxuseclass}{cell_input}
\begin{sphinxVerbatim}[commandchars=\\\{\}]
\PYG{c+c1}{\PYGZsh{} Concatenate two tensors}
\PYG{n}{t1} \PYG{o}{=} \PYG{n}{torch}\PYG{o}{.}\PYG{n}{tensor}\PYG{p}{(}\PYG{p}{[}\PYG{p}{[}\PYG{l+m+mi}{1}\PYG{p}{,} \PYG{l+m+mi}{2}\PYG{p}{,} \PYG{l+m+mi}{3}\PYG{p}{]}\PYG{p}{,}
                   \PYG{p}{[}\PYG{l+m+mi}{4}\PYG{p}{,} \PYG{l+m+mi}{5}\PYG{p}{,} \PYG{l+m+mi}{6}\PYG{p}{]}\PYG{p}{]}\PYG{p}{)}
\PYG{n}{t2} \PYG{o}{=} \PYG{n}{torch}\PYG{o}{.}\PYG{n}{tensor}\PYG{p}{(}\PYG{p}{[}\PYG{p}{[}\PYG{l+m+mi}{7}\PYG{p}{,} \PYG{l+m+mi}{8}\PYG{p}{,} \PYG{l+m+mi}{9}\PYG{p}{]}\PYG{p}{,}
                   \PYG{p}{[}\PYG{l+m+mi}{10}\PYG{p}{,} \PYG{l+m+mi}{11}\PYG{p}{,} \PYG{l+m+mi}{12}\PYG{p}{]}\PYG{p}{]}\PYG{p}{)}
\PYG{n}{t\PYGZus{}cat0} \PYG{o}{=} \PYG{n}{torch}\PYG{o}{.}\PYG{n}{cat}\PYG{p}{(}\PYG{p}{[}\PYG{n}{t1}\PYG{p}{,} \PYG{n}{t2}\PYG{p}{]}\PYG{p}{,} \PYG{n}{dim}\PYG{o}{=}\PYG{l+m+mi}{0}\PYG{p}{)}  \PYG{c+c1}{\PYGZsh{} concatenate along rows (dim=0)}
\PYG{n}{t\PYGZus{}cat1} \PYG{o}{=} \PYG{n}{torch}\PYG{o}{.}\PYG{n}{cat}\PYG{p}{(}\PYG{p}{[}\PYG{n}{t1}\PYG{p}{,} \PYG{n}{t2}\PYG{p}{]}\PYG{p}{,} \PYG{n}{dim}\PYG{o}{=}\PYG{l+m+mi}{1}\PYG{p}{)}  \PYG{c+c1}{\PYGZsh{} concatenate along columns (dim=1)}

\PYG{n}{pprint}\PYG{p}{(}\PYG{l+s+s2}{\PYGZdq{}}\PYG{l+s+s2}{t\PYGZus{}1=}\PYG{l+s+s2}{\PYGZdq{}}\PYG{p}{,} \PYG{n}{t1}\PYG{p}{)}
\PYG{n}{pprint}\PYG{p}{(}\PYG{l+s+s2}{\PYGZdq{}}\PYG{l+s+s2}{t\PYGZus{}2:=}\PYG{l+s+s2}{\PYGZdq{}}\PYG{p}{,} \PYG{n}{t2}\PYG{p}{)}
\PYG{n}{pprint}\PYG{p}{(}\PYG{l+s+s2}{\PYGZdq{}}\PYG{l+s+s2}{t\PYGZus{}1}\PYG{l+s+se}{\PYGZbs{}\PYGZbs{}}\PYG{l+s+s2}{text}\PYG{l+s+s2}{\PYGZob{}}\PYG{l+s+s2}{ and \PYGZcb{}t\PYGZus{}2}\PYG{l+s+se}{\PYGZbs{}\PYGZbs{}}\PYG{l+s+s2}{text}\PYG{l+s+s2}{\PYGZob{}}\PYG{l+s+s2}{  concatenated along dim=0: \PYGZcb{}}\PYG{l+s+s2}{\PYGZdq{}}\PYG{p}{,} \PYG{n}{t\PYGZus{}cat0}\PYG{p}{)}
\PYG{n}{pprint}\PYG{p}{(}\PYG{l+s+s2}{\PYGZdq{}}\PYG{l+s+s2}{t\PYGZus{}1}\PYG{l+s+se}{\PYGZbs{}\PYGZbs{}}\PYG{l+s+s2}{text}\PYG{l+s+s2}{\PYGZob{}}\PYG{l+s+s2}{  and \PYGZcb{}t\PYGZus{}2}\PYG{l+s+se}{\PYGZbs{}\PYGZbs{}}\PYG{l+s+s2}{text}\PYG{l+s+s2}{\PYGZob{}}\PYG{l+s+s2}{  concatenated along dim=1: \PYGZcb{}}\PYG{l+s+s2}{\PYGZdq{}}\PYG{p}{,} \PYG{n}{t\PYGZus{}cat1}\PYG{p}{)}
\end{sphinxVerbatim}

\end{sphinxuseclass}\end{sphinxVerbatimInput}
\begin{sphinxVerbatimOutput}

\begin{sphinxuseclass}{cell_output}\begin{equation*}
\begin{split}\displaystyle t_1=\left[
\begin{array}{}
  1.0000 &  2.0000 &  3.0000\\
  4.0000 &  5.0000 &  6.0000
\end{array}
\right]\end{split}
\end{equation*}\begin{equation*}
\begin{split}\displaystyle t_2:=\left[
\begin{array}{}
  7.0000 &  8.0000 &  9.0000\\
  10.0000 &  11.0000 &  12.0000
\end{array}
\right]\end{split}
\end{equation*}\begin{equation*}
\begin{split}\displaystyle t_1\text{ and }t_2\text{  concatenated along dim=0: }\left[
\begin{array}{}
  1.0000 &  2.0000 &  3.0000\\
  4.0000 &  5.0000 &  6.0000\\
  7.0000 &  8.0000 &  9.0000\\
  10.0000 &  11.0000 &  12.0000
\end{array}
\right]\end{split}
\end{equation*}\begin{equation*}
\begin{split}\displaystyle t_1\text{  and }t_2\text{  concatenated along dim=1: }\left[
\begin{array}{}
  1.0000 &  2.0000 &  3.0000 &  7.0000 &  8.0000 &  9.0000\\
  4.0000 &  5.0000 &  6.0000 &  10.0000 &  11.0000 &  12.0000
\end{array}
\right]\end{split}
\end{equation*}
\end{sphinxuseclass}\end{sphinxVerbatimOutput}

\end{sphinxuseclass}
\sphinxAtStartPar
Here:
\begin{itemize}
\item {} 
\sphinxAtStartPar
We created a 1\sphinxhyphen{}D tensor \sphinxcode{\sphinxupquote{t}} with 16 sequential elements. We reshaped it into a 4×4 tensor \sphinxcode{\sphinxupquote{t\_matrix}}.

\item {} 
\sphinxAtStartPar
We transposed \sphinxcode{\sphinxupquote{t\_matrix}} with \sphinxcode{\sphinxupquote{.T}}.

\item {} 
\sphinxAtStartPar
We took two 2×3 tensors \sphinxcode{\sphinxupquote{t1}} and \sphinxcode{\sphinxupquote{t2}} and concatenated them. Along \sphinxcode{\sphinxupquote{dim=0}} (rows), we got a 4×3 tensor by stacking \sphinxcode{\sphinxupquote{t2}} below \sphinxcode{\sphinxupquote{t1}}. Along \sphinxcode{\sphinxupquote{dim=1}} (columns), we got a 2×6 tensor by stacking \sphinxcode{\sphinxupquote{t2}} to the right of \sphinxcode{\sphinxupquote{t1}}.

\end{itemize}

\sphinxAtStartPar
These basic operations (arithmetic, indexing, reshaping, concatenation, etc.) allow you to manipulate tensor data in flexible ways. PyTorch’s API is designed to be highly compatible with NumPy’s where it makes sense.


\subsection{Autograd and Gradients}
\label{\detokenize{04:autograd-and-gradients}}
\sphinxAtStartPar
One of the most powerful features of PyTorch is its \sphinxstylestrong{autograd} system, which allows automatic computation of gradients for tensor operations. This is what enables training of neural networks using gradient descent: you define a forward computation, and PyTorch can compute the backward gradients for you.

\sphinxAtStartPar
In PyTorch, the \sphinxcode{\sphinxupquote{torch.autograd}} engine tracks all operations on tensors that have \sphinxcode{\sphinxupquote{requires\_grad=True}}. By default, tensors do \sphinxstylestrong{not} track gradients. You need to specify which tensors require gradient computation (usually the learnable parameters of your model, or any input for which you want to compute a derivative) by setting \sphinxcode{\sphinxupquote{requires\_grad=True}} when creating the tensor or by calling \sphinxcode{\sphinxupquote{.requires\_grad\_()}} on an existing tensor.

\sphinxAtStartPar
When you perform operations on such tensors, PyTorch builds a computational graph behind the scenes: nodes are tensors, and edges are functions that produce output tensors from input tensors. Once you have a result (typically a scalar loss value), you can call \sphinxcode{\sphinxupquote{.backward()}} on it, and PyTorch will automatically traverse the graph to compute the gradient of that result with respect to every tensor that has \sphinxcode{\sphinxupquote{requires\_grad=True}} (using the chain rule of calculus). Let’s see a simple example to illustrate this:

\begin{sphinxuseclass}{cell}\begin{sphinxVerbatimInput}

\begin{sphinxuseclass}{cell_input}
\begin{sphinxVerbatim}[commandchars=\\\{\}]
\PYG{c+c1}{\PYGZsh{} A simple autograd example: y = f(x) where f(x) = x\PYGZca{}2 + 2x + 1}
\PYG{n}{x} \PYG{o}{=} \PYG{n}{torch}\PYG{o}{.}\PYG{n}{tensor}\PYG{p}{(}\PYG{l+m+mf}{3.0}\PYG{p}{,} \PYG{n}{requires\PYGZus{}grad}\PYG{o}{=}\PYG{k+kc}{True}\PYG{p}{)}   \PYG{c+c1}{\PYGZsh{} define a tensor x with gradient tracking}
\PYG{n+nb}{print}\PYG{p}{(}\PYG{l+s+s2}{\PYGZdq{}}\PYG{l+s+s2}{x:}\PYG{l+s+s2}{\PYGZdq{}}\PYG{p}{,} \PYG{n}{x}\PYG{p}{)}

\PYG{c+c1}{\PYGZsh{} Define a function of x}
\PYG{n}{y} \PYG{o}{=} \PYG{n}{x}\PYG{o}{*}\PYG{o}{*}\PYG{l+m+mi}{2} \PYG{o}{+} \PYG{l+m+mi}{2}\PYG{o}{*}\PYG{n}{x} \PYG{o}{+} \PYG{l+m+mi}{1}
\PYG{n+nb}{print}\PYG{p}{(}\PYG{l+s+s2}{\PYGZdq{}}\PYG{l+s+s2}{y:}\PYG{l+s+s2}{\PYGZdq{}}\PYG{p}{,} \PYG{n}{y}\PYG{p}{)}

\PYG{c+c1}{\PYGZsh{} Compute the gradient dy/dx by calling backward on y}
\PYG{n}{y}\PYG{o}{.}\PYG{n}{backward}\PYG{p}{(}\PYG{p}{)}    \PYG{c+c1}{\PYGZsh{} computes gradient of y w.r.t. x}
\PYG{n+nb}{print}\PYG{p}{(}\PYG{l+s+s2}{\PYGZdq{}}\PYG{l+s+s2}{dy/dx at x=3:}\PYG{l+s+s2}{\PYGZdq{}}\PYG{p}{,} \PYG{n}{x}\PYG{o}{.}\PYG{n}{grad}\PYG{p}{)}
\end{sphinxVerbatim}

\end{sphinxuseclass}\end{sphinxVerbatimInput}
\begin{sphinxVerbatimOutput}

\begin{sphinxuseclass}{cell_output}
\begin{sphinxVerbatim}[commandchars=\\\{\}]
x: tensor(3., requires\PYGZus{}grad=True)
y: tensor(16., grad\PYGZus{}fn=\PYGZlt{}AddBackward0\PYGZgt{})
dy/dx at x=3: tensor(8.)
\end{sphinxVerbatim}

\end{sphinxuseclass}\end{sphinxVerbatimOutput}

\end{sphinxuseclass}
\sphinxAtStartPar
In this example:
\begin{itemize}
\item {} 
\sphinxAtStartPar
We created a tensor \sphinxcode{\sphinxupquote{x}} with value \sphinxcode{\sphinxupquote{3.0}} and \sphinxcode{\sphinxupquote{requires\_grad=True}}. This means PyTorch will track operations on \sphinxcode{\sphinxupquote{x}}.

\item {} 
\sphinxAtStartPar
We computed \sphinxcode{\sphinxupquote{y = x**2 + 2*x + 1}}. Because \sphinxcode{\sphinxupquote{x}} requires grad, \sphinxcode{\sphinxupquote{y}} will also by default require grad and have a grad function associated.

\item {} 
\sphinxAtStartPar
We called \sphinxcode{\sphinxupquote{y.backward()}}. Since \sphinxcode{\sphinxupquote{y}} is a scalar (single value), we can call \sphinxcode{\sphinxupquote{backward}} directly. After this call, the gradient of \sphinxcode{\sphinxupquote{y}} with respect to \sphinxcode{\sphinxupquote{x}} is computed and stored in \sphinxcode{\sphinxupquote{x.grad}}.

\item {} 
\sphinxAtStartPar
We printed \sphinxcode{\sphinxupquote{x.grad}}. The derivative \( y = x^2 + 2x + 1 \) is \( dy/dx = 2x + 2 \). At \( x = 3 \), that is \sphinxcode{\sphinxupquote{8}}.

\end{itemize}

\sphinxAtStartPar
\sphinxstylestrong{Important notes about autograd}:
\begin{enumerate}
\sphinxsetlistlabels{\arabic}{enumi}{enumii}{}{.}%
\item {} 
\sphinxAtStartPar
You can disable gradient tracking by wrapping code in \sphinxcode{\sphinxupquote{with torch.no\_grad():}} or by using \sphinxcode{\sphinxupquote{.detach()}} on a tensor to get a new tensor that shares the same data but with no grad tracking.

\item {} 
\sphinxAtStartPar
If \sphinxcode{\sphinxupquote{y}} is not a scalar, you need to specify a gradient argument in \sphinxcode{\sphinxupquote{y.backward(gradient=...)}} which is the tensor of same shape as \sphinxcode{\sphinxupquote{y}} containing the gradient of the final result w.r.t. each element of \sphinxcode{\sphinxupquote{y}}.

\item {} 
\sphinxAtStartPar
By default, gradients are \sphinxstylestrong{accumulated} into the \sphinxcode{\sphinxupquote{.grad}} property. So typically you want to zero out gradients (e.g., \sphinxcode{\sphinxupquote{x.grad.zero\_()}}) between independent backward passes.

\end{enumerate}

\sphinxAtStartPar
Autograd is what powers training routines. If you had a multi\sphinxhyphen{}parameter function (like a neural network with many weights), you would set \sphinxcode{\sphinxupquote{requires\_grad=True}} for all those parameters. After computing a loss (scalar), you’d call \sphinxcode{\sphinxupquote{loss.backward()}}, and each parameter tensor’s \sphinxcode{\sphinxupquote{.grad}} will be filled with its gradient.

\sphinxAtStartPar
To reinforce understanding, let’s do a slightly more complex autograd example with two inputs:

\begin{sphinxuseclass}{cell}\begin{sphinxVerbatimInput}

\begin{sphinxuseclass}{cell_input}
\begin{sphinxVerbatim}[commandchars=\\\{\}]
\PYG{c+c1}{\PYGZsh{} Autograd with two inputs: z = x * y + y\PYGZca{}2}
\PYG{n}{x} \PYG{o}{=} \PYG{n}{torch}\PYG{o}{.}\PYG{n}{tensor}\PYG{p}{(}\PYG{l+m+mf}{2.0}\PYG{p}{,} \PYG{n}{requires\PYGZus{}grad}\PYG{o}{=}\PYG{k+kc}{True}\PYG{p}{)}
\PYG{n}{y} \PYG{o}{=} \PYG{n}{torch}\PYG{o}{.}\PYG{n}{tensor}\PYG{p}{(}\PYG{o}{\PYGZhy{}}\PYG{l+m+mf}{3.0}\PYG{p}{,} \PYG{n}{requires\PYGZus{}grad}\PYG{o}{=}\PYG{k+kc}{True}\PYG{p}{)}

\PYG{n}{z} \PYG{o}{=} \PYG{n}{x} \PYG{o}{*} \PYG{n}{y} \PYG{o}{+} \PYG{n}{y}\PYG{o}{*}\PYG{o}{*}\PYG{l+m+mi}{2}   \PYG{c+c1}{\PYGZsh{} z is a function of both x and y}
\PYG{n+nb}{print}\PYG{p}{(}\PYG{l+s+s2}{\PYGZdq{}}\PYG{l+s+s2}{z:}\PYG{l+s+s2}{\PYGZdq{}}\PYG{p}{,} \PYG{n}{z}\PYG{p}{)}

\PYG{n}{z}\PYG{o}{.}\PYG{n}{backward}\PYG{p}{(}\PYG{p}{)}  \PYG{c+c1}{\PYGZsh{} Compute gradients dz/dx and dz/dy}
\PYG{n+nb}{print}\PYG{p}{(}\PYG{l+s+s2}{\PYGZdq{}}\PYG{l+s+s2}{dz/dx:}\PYG{l+s+s2}{\PYGZdq{}}\PYG{p}{,} \PYG{n}{x}\PYG{o}{.}\PYG{n}{grad}\PYG{p}{)}   \PYG{c+c1}{\PYGZsh{} should be y}
\PYG{n+nb}{print}\PYG{p}{(}\PYG{l+s+s2}{\PYGZdq{}}\PYG{l+s+s2}{dz/dy:}\PYG{l+s+s2}{\PYGZdq{}}\PYG{p}{,} \PYG{n}{y}\PYG{o}{.}\PYG{n}{grad}\PYG{p}{)}   \PYG{c+c1}{\PYGZsh{} should be x + 2y}
\end{sphinxVerbatim}

\end{sphinxuseclass}\end{sphinxVerbatimInput}
\begin{sphinxVerbatimOutput}

\begin{sphinxuseclass}{cell_output}
\begin{sphinxVerbatim}[commandchars=\\\{\}]
z: tensor(3., grad\PYGZus{}fn=\PYGZlt{}AddBackward0\PYGZgt{})
dz/dx: tensor(\PYGZhy{}3.)
dz/dy: tensor(\PYGZhy{}4.)
\end{sphinxVerbatim}

\end{sphinxuseclass}\end{sphinxVerbatimOutput}

\end{sphinxuseclass}
\sphinxAtStartPar
Here \( z = x \times y + y^2 \). We expect \( dz/dx = y \) and \( dz/dy = x + 2y \). With \( x = 2.0 \) and \( y = -3.0 \), those values are \( dz/dx = -3 \) and \( dz/dy = 2 + 2(-3) = -4 \). The autograd should produce the same results in \sphinxcode{\sphinxupquote{x.grad}} and \sphinxcode{\sphinxupquote{y.grad}}.


\subsection{Advanced Tensor Functions}
\label{\detokenize{04:advanced-tensor-functions}}

\subsubsection{Broadcasting}
\label{\detokenize{04:broadcasting}}
\sphinxAtStartPar
Like NumPy, PyTorch supports \sphinxstylestrong{broadcasting} – a mechanism that allows arithmetic operations on tensors of different shapes by automatically expanding one of them to match the shape of the other. Broadcasting rules in PyTorch are the same as NumPy’s. In short, two tensors are compatible for an elementwise operation if their shapes are equal or align in such a way that one of them can be expanded (a dimension of size 1 can be expanded to match the other’s size, and missing dimensions are treated as size 1).

\sphinxAtStartPar
For example, if you have a tensor of shape (4, 3) and another of shape (3,), you can add them because the second tensor can be treated as if it were (1, 3) and then expanded to (4, 3) to match the first tensor’s shape.

\sphinxAtStartPar
In code, we can demonstrate broadcasting with a simple example:

\begin{sphinxuseclass}{cell}\begin{sphinxVerbatimInput}

\begin{sphinxuseclass}{cell_input}
\begin{sphinxVerbatim}[commandchars=\\\{\}]
\PYG{c+c1}{\PYGZsh{} Broadcasting example: add a vector to each row of a matrix}
\PYG{n}{mat} \PYG{o}{=} \PYG{n}{torch}\PYG{o}{.}\PYG{n}{tensor}\PYG{p}{(}\PYG{p}{[}\PYG{p}{[}\PYG{l+m+mi}{0}\PYG{p}{,} \PYG{l+m+mi}{0}\PYG{p}{,} \PYG{l+m+mi}{0}\PYG{p}{]}\PYG{p}{,}
                    \PYG{p}{[}\PYG{l+m+mi}{10}\PYG{p}{,} \PYG{l+m+mi}{10}\PYG{p}{,} \PYG{l+m+mi}{10}\PYG{p}{]}\PYG{p}{,}
                    \PYG{p}{[}\PYG{l+m+mi}{20}\PYG{p}{,} \PYG{l+m+mi}{20}\PYG{p}{,} \PYG{l+m+mi}{20}\PYG{p}{]}\PYG{p}{,}
                    \PYG{p}{[}\PYG{l+m+mi}{30}\PYG{p}{,} \PYG{l+m+mi}{30}\PYG{p}{,} \PYG{l+m+mi}{30}\PYG{p}{]}\PYG{p}{]}\PYG{p}{)}   \PYG{c+c1}{\PYGZsh{} shape (4,3)}
\PYG{n}{vec} \PYG{o}{=} \PYG{n}{torch}\PYG{o}{.}\PYG{n}{tensor}\PYG{p}{(}\PYG{p}{[}\PYG{l+m+mi}{1}\PYG{p}{,} \PYG{l+m+mi}{2}\PYG{p}{,} \PYG{l+m+mi}{3}\PYG{p}{]}\PYG{p}{)}        \PYG{c+c1}{\PYGZsh{} shape (3,)}

\PYG{n}{result} \PYG{o}{=} \PYG{n}{mat} \PYG{o}{+} \PYG{n}{vec}  \PYG{c+c1}{\PYGZsh{} vec is broadcast to shape (4,3)}
\PYG{n}{pprint}\PYG{p}{(}\PYG{l+s+s2}{\PYGZdq{}}\PYG{l+s+s2}{M=}\PYG{l+s+s2}{\PYGZdq{}}\PYG{p}{,} \PYG{n}{mat}\PYG{p}{)}
\PYG{n}{pprint}\PYG{p}{(}\PYG{l+s+s2}{\PYGZdq{}}\PYG{l+s+s2}{V=}\PYG{l+s+s2}{\PYGZdq{}}\PYG{p}{,} \PYG{n}{vec}\PYG{p}{)}
\PYG{n}{pprint}\PYG{p}{(}\PYG{l+s+s2}{\PYGZdq{}}\PYG{l+s+s2}{M+V}\PYG{l+s+s2}{\PYGZdq{}}\PYG{p}{,} \PYG{n}{result}\PYG{p}{)}
\end{sphinxVerbatim}

\end{sphinxuseclass}\end{sphinxVerbatimInput}
\begin{sphinxVerbatimOutput}

\begin{sphinxuseclass}{cell_output}\begin{equation*}
\begin{split}\displaystyle M=\left[
\begin{array}{}
  0.0000 &  0.0000 &  0.0000\\
  10.0000 &  10.0000 &  10.0000\\
  20.0000 &  20.0000 &  20.0000\\
  30.0000 &  30.0000 &  30.0000
\end{array}
\right]\end{split}
\end{equation*}\begin{equation*}
\begin{split}\displaystyle V=\left[
\begin{array}{}
  1.0000 &  2.0000 &  3.0000
\end{array}
\right]\end{split}
\end{equation*}\begin{equation*}
\begin{split}\displaystyle M+V\left[
\begin{array}{}
  1.0000 &  2.0000 &  3.0000\\
  11.0000 &  12.0000 &  13.0000\\
  21.0000 &  22.0000 &  23.0000\\
  31.0000 &  32.0000 &  33.0000
\end{array}
\right]\end{split}
\end{equation*}
\end{sphinxuseclass}\end{sphinxVerbatimOutput}

\end{sphinxuseclass}
\sphinxAtStartPar
Here, \sphinxcode{\sphinxupquote{vec}} is added to each row of \sphinxcode{\sphinxupquote{mat}}. The expected result would be each row of \sphinxcode{\sphinxupquote{mat}} incremented by \sphinxcode{\sphinxupquote{{[}1, 2, 3{]}}}. Broadcasting is a powerful feature that makes your code more concise and often more efficient.


\subsubsection{Device Management (CPU vs GPU)}
\label{\detokenize{04:device-management-cpu-vs-gpu}}
\sphinxAtStartPar
One of PyTorch’s strengths is the ability to perform computations on a \sphinxstylestrong{GPU} for speed. Tensors can reside on either the CPU or the GPU. By default, tensors are created on the CPU.
You can check where a tensor is located by looking at its \sphinxcode{\sphinxupquote{.device}} attribute. To leverage a GPU (if you have one and PyTorch is installed with CUDA support), you need to explicitly move tensors to the GPU. There are a few ways to do this:
\begin{enumerate}
\sphinxsetlistlabels{\arabic}{enumi}{enumii}{}{.}%
\item {} 
\sphinxAtStartPar
Using the \sphinxcode{\sphinxupquote{.to(device)}} method on a tensor, where \sphinxcode{\sphinxupquote{device}} is something like \sphinxcode{\sphinxupquote{torch.device("cuda")}} or the shorthand string \sphinxcode{\sphinxupquote{"cuda"}}.

\item {} 
\sphinxAtStartPar
Using \sphinxcode{\sphinxupquote{.cuda()}} method on the tensor.

\item {} 
\sphinxAtStartPar
Creating the tensor directly on the GPU by specifying \sphinxcode{\sphinxupquote{device=torch.device("cuda")}} (or \sphinxcode{\sphinxupquote{device="cuda"}}) in the creation function.

\end{enumerate}

\sphinxAtStartPar
It’s common to write code that is device\sphinxhyphen{}agnostic (runs on CPU if no GPU is available, otherwise uses the GPU). For example:

\begin{sphinxuseclass}{cell}\begin{sphinxVerbatimInput}

\begin{sphinxuseclass}{cell_input}
\begin{sphinxVerbatim}[commandchars=\\\{\}]
\PYG{c+c1}{\PYGZsh{} Device\PYGZhy{}agnostic code: choose CPU or CUDA}
\PYG{n}{device} \PYG{o}{=} \PYG{n}{torch}\PYG{o}{.}\PYG{n}{device}\PYG{p}{(}\PYG{l+s+s2}{\PYGZdq{}}\PYG{l+s+s2}{cuda}\PYG{l+s+s2}{\PYGZdq{}} \PYG{k}{if} \PYG{n}{torch}\PYG{o}{.}\PYG{n}{cuda}\PYG{o}{.}\PYG{n}{is\PYGZus{}available}\PYG{p}{(}\PYG{p}{)} \PYG{k}{else} \PYG{p}{(}\PYG{l+s+s2}{\PYGZdq{}}\PYG{l+s+s2}{mps}\PYG{l+s+s2}{\PYGZdq{}} \PYG{k}{if} \PYG{n}{torch}\PYG{o}{.}\PYG{n}{backends}\PYG{o}{.}\PYG{n}{mps}\PYG{o}{.}\PYG{n}{is\PYGZus{}available}\PYG{p}{(}\PYG{p}{)} \PYG{k}{else} \PYG{l+s+s2}{\PYGZdq{}}\PYG{l+s+s2}{cpu}\PYG{l+s+s2}{\PYGZdq{}}\PYG{p}{)}\PYG{p}{)}
\PYG{n+nb}{print}\PYG{p}{(}\PYG{l+s+s2}{\PYGZdq{}}\PYG{l+s+s2}{Using device:}\PYG{l+s+s2}{\PYGZdq{}}\PYG{p}{,} \PYG{n}{device}\PYG{p}{)}

\PYG{c+c1}{\PYGZsh{} Create a tensor on the chosen device}
\PYG{n}{x} \PYG{o}{=} \PYG{n}{torch}\PYG{o}{.}\PYG{n}{tensor}\PYG{p}{(}\PYG{p}{[}\PYG{l+m+mf}{1.0}\PYG{p}{,} \PYG{l+m+mf}{2.0}\PYG{p}{,} \PYG{l+m+mf}{3.0}\PYG{p}{]}\PYG{p}{,} \PYG{n}{device}\PYG{o}{=}\PYG{n}{device}\PYG{p}{)}
\PYG{n+nb}{print}\PYG{p}{(}\PYG{l+s+s2}{\PYGZdq{}}\PYG{l+s+s2}{x device:}\PYG{l+s+s2}{\PYGZdq{}}\PYG{p}{,} \PYG{n}{x}\PYG{o}{.}\PYG{n}{device}\PYG{p}{)}

\PYG{c+c1}{\PYGZsh{} Alternatively, move an existing tensor to the device}
\PYG{n}{y} \PYG{o}{=} \PYG{n}{torch}\PYG{o}{.}\PYG{n}{ones}\PYG{p}{(}\PYG{l+m+mi}{3}\PYG{p}{)}            \PYG{c+c1}{\PYGZsh{} defaults to CPU}
\PYG{n}{y} \PYG{o}{=} \PYG{n}{y}\PYG{o}{.}\PYG{n}{to}\PYG{p}{(}\PYG{n}{device}\PYG{p}{)}             \PYG{c+c1}{\PYGZsh{} move to GPU if available}
\PYG{n+nb}{print}\PYG{p}{(}\PYG{l+s+s2}{\PYGZdq{}}\PYG{l+s+s2}{y device:}\PYG{l+s+s2}{\PYGZdq{}}\PYG{p}{,} \PYG{n}{y}\PYG{o}{.}\PYG{n}{device}\PYG{p}{)}
\end{sphinxVerbatim}

\end{sphinxuseclass}\end{sphinxVerbatimInput}
\begin{sphinxVerbatimOutput}

\begin{sphinxuseclass}{cell_output}
\begin{sphinxVerbatim}[commandchars=\\\{\}]
Using device: mps
\end{sphinxVerbatim}

\begin{sphinxVerbatim}[commandchars=\\\{\}]
x device: mps:0
y device: mps:0
\end{sphinxVerbatim}

\end{sphinxuseclass}\end{sphinxVerbatimOutput}

\end{sphinxuseclass}
\sphinxAtStartPar
If you actually run this on a system with a GPU and CUDA installed, you should see device printed as \sphinxcode{\sphinxupquote{cuda:0}} and the tensor devices as such. On a CPU\sphinxhyphen{}only system, it will say \sphinxcode{\sphinxupquote{cpu}}.

\sphinxAtStartPar
A quick verification of CPU vs GPU performance can be done by timing a large matrix multiplication. (Note: if you actually run this, avoid printing the giant result; just measure the time.)

\begin{sphinxuseclass}{cell}\begin{sphinxVerbatimInput}

\begin{sphinxuseclass}{cell_input}
\begin{sphinxVerbatim}[commandchars=\\\{\}]
\PYG{k+kn}{import}\PYG{+w}{ }\PYG{n+nn}{time}

\PYG{c+c1}{\PYGZsh{} Large tensor operation on CPU vs GPU (if available)}
\PYG{n}{large\PYGZus{}cpu} \PYG{o}{=} \PYG{n}{torch}\PYG{o}{.}\PYG{n}{randn}\PYG{p}{(}\PYG{l+m+mi}{10000}\PYG{p}{,} \PYG{l+m+mi}{10000}\PYG{p}{)}  \PYG{c+c1}{\PYGZsh{} 10000x10000 matrix on CPU}
\PYG{n}{start} \PYG{o}{=} \PYG{n}{time}\PYG{o}{.}\PYG{n}{time}\PYG{p}{(}\PYG{p}{)}
\PYG{n}{res\PYGZus{}cpu} \PYG{o}{=} \PYG{n}{large\PYGZus{}cpu} \PYG{o}{@} \PYG{n}{large\PYGZus{}cpu}       \PYG{c+c1}{\PYGZsh{} matrix multiply on CPU}
\PYG{n}{end} \PYG{o}{=} \PYG{n}{time}\PYG{o}{.}\PYG{n}{time}\PYG{p}{(}\PYG{p}{)}
\PYG{n+nb}{print}\PYG{p}{(}\PYG{l+s+s2}{\PYGZdq{}}\PYG{l+s+s2}{CPU computation time:}\PYG{l+s+s2}{\PYGZdq{}}\PYG{p}{,} \PYG{n}{end} \PYG{o}{\PYGZhy{}} \PYG{n}{start}\PYG{p}{,} \PYG{l+s+s2}{\PYGZdq{}}\PYG{l+s+s2}{seconds}\PYG{l+s+s2}{\PYGZdq{}}\PYG{p}{)}

\PYG{k}{if} \PYG{n}{torch}\PYG{o}{.}\PYG{n}{cuda}\PYG{o}{.}\PYG{n}{is\PYGZus{}available}\PYG{p}{(}\PYG{p}{)}\PYG{p}{:}
    \PYG{n}{large\PYGZus{}gpu} \PYG{o}{=} \PYG{n}{large\PYGZus{}cpu}\PYG{o}{.}\PYG{n}{to}\PYG{p}{(}\PYG{l+s+s2}{\PYGZdq{}}\PYG{l+s+s2}{cuda}\PYG{l+s+s2}{\PYGZdq{}}\PYG{p}{)}
    \PYG{n}{torch}\PYG{o}{.}\PYG{n}{cuda}\PYG{o}{.}\PYG{n}{synchronize}\PYG{p}{(}\PYG{p}{)}  \PYG{c+c1}{\PYGZsh{} ensure data is transferred}
    \PYG{n}{start} \PYG{o}{=} \PYG{n}{time}\PYG{o}{.}\PYG{n}{time}\PYG{p}{(}\PYG{p}{)}
    \PYG{n}{res\PYGZus{}gpu} \PYG{o}{=} \PYG{n}{large\PYGZus{}gpu} \PYG{o}{@} \PYG{n}{large\PYGZus{}gpu}   \PYG{c+c1}{\PYGZsh{} matrix multiply on GPU}
    \PYG{n}{torch}\PYG{o}{.}\PYG{n}{cuda}\PYG{o}{.}\PYG{n}{synchronize}\PYG{p}{(}\PYG{p}{)}  \PYG{c+c1}{\PYGZsh{} wait for GPU to finish computation}
    \PYG{n}{end} \PYG{o}{=} \PYG{n}{time}\PYG{o}{.}\PYG{n}{time}\PYG{p}{(}\PYG{p}{)}
    \PYG{n+nb}{print}\PYG{p}{(}\PYG{l+s+s2}{\PYGZdq{}}\PYG{l+s+s2}{GPU computation time:}\PYG{l+s+s2}{\PYGZdq{}}\PYG{p}{,} \PYG{n}{end} \PYG{o}{\PYGZhy{}} \PYG{n}{start}\PYG{p}{,} \PYG{l+s+s2}{\PYGZdq{}}\PYG{l+s+s2}{seconds}\PYG{l+s+s2}{\PYGZdq{}}\PYG{p}{)}

\PYG{k}{if} \PYG{n}{torch}\PYG{o}{.}\PYG{n}{backends}\PYG{o}{.}\PYG{n}{mps}\PYG{o}{.}\PYG{n}{is\PYGZus{}available}\PYG{p}{(}\PYG{p}{)}\PYG{p}{:}
    \PYG{n}{large\PYGZus{}gpu} \PYG{o}{=} \PYG{n}{large\PYGZus{}cpu}\PYG{o}{.}\PYG{n}{to}\PYG{p}{(}\PYG{l+s+s2}{\PYGZdq{}}\PYG{l+s+s2}{mps}\PYG{l+s+s2}{\PYGZdq{}}\PYG{p}{)}
    \PYG{n}{start} \PYG{o}{=} \PYG{n}{time}\PYG{o}{.}\PYG{n}{time}\PYG{p}{(}\PYG{p}{)}
    \PYG{n}{res\PYGZus{}gpu} \PYG{o}{=} \PYG{n}{large\PYGZus{}gpu} \PYG{o}{@} \PYG{n}{large\PYGZus{}gpu}
    \PYG{n}{end} \PYG{o}{=} \PYG{n}{time}\PYG{o}{.}\PYG{n}{time}\PYG{p}{(}\PYG{p}{)}
    \PYG{n+nb}{print}\PYG{p}{(}\PYG{l+s+s2}{\PYGZdq{}}\PYG{l+s+s2}{MPS computation time:}\PYG{l+s+s2}{\PYGZdq{}}\PYG{p}{,} \PYG{n}{end} \PYG{o}{\PYGZhy{}} \PYG{n}{start}\PYG{p}{,} \PYG{l+s+s2}{\PYGZdq{}}\PYG{l+s+s2}{seconds}\PYG{l+s+s2}{\PYGZdq{}}\PYG{p}{)}
\end{sphinxVerbatim}

\end{sphinxuseclass}\end{sphinxVerbatimInput}
\begin{sphinxVerbatimOutput}

\begin{sphinxuseclass}{cell_output}
\begin{sphinxVerbatim}[commandchars=\\\{\}]
CPU computation time: 1.7131898403167725 seconds
\end{sphinxVerbatim}

\begin{sphinxVerbatim}[commandchars=\\\{\}]
MPS computation time: 0.21483182907104492 seconds
\end{sphinxVerbatim}

\end{sphinxuseclass}\end{sphinxVerbatimOutput}

\end{sphinxuseclass}
\sphinxAtStartPar
This snippet will likely show that the GPU computation is significantly faster for large matrix operations (depending on hardware). Remember to move your \sphinxstylestrong{model parameters} and \sphinxstylestrong{data} to the same device before doing operations.

\sphinxAtStartPar
For example, the code below moves a tensor to GPU if available:

\begin{sphinxuseclass}{cell}\begin{sphinxVerbatimInput}

\begin{sphinxuseclass}{cell_input}
\begin{sphinxVerbatim}[commandchars=\\\{\}]
\PYG{n}{tensor} \PYG{o}{=} \PYG{n}{torch}\PYG{o}{.}\PYG{n}{rand}\PYG{p}{(}\PYG{l+m+mi}{3}\PYG{p}{)}
\PYG{n+nb}{print}\PYG{p}{(}\PYG{l+s+s2}{\PYGZdq{}}\PYG{l+s+s2}{Before moving, device:}\PYG{l+s+s2}{\PYGZdq{}}\PYG{p}{,} \PYG{n}{tensor}\PYG{o}{.}\PYG{n}{device}\PYG{p}{)}
\PYG{k}{if} \PYG{n}{torch}\PYG{o}{.}\PYG{n}{cuda}\PYG{o}{.}\PYG{n}{is\PYGZus{}available}\PYG{p}{(}\PYG{p}{)}\PYG{p}{:}
    \PYG{n}{tensor} \PYG{o}{=} \PYG{n}{tensor}\PYG{o}{.}\PYG{n}{to}\PYG{p}{(}\PYG{l+s+s1}{\PYGZsq{}}\PYG{l+s+s1}{cuda}\PYG{l+s+s1}{\PYGZsq{}}\PYG{p}{)}
    \PYG{n+nb}{print}\PYG{p}{(}\PYG{l+s+s2}{\PYGZdq{}}\PYG{l+s+s2}{After moving, device:}\PYG{l+s+s2}{\PYGZdq{}}\PYG{p}{,} \PYG{n}{tensor}\PYG{o}{.}\PYG{n}{device}\PYG{p}{)}
\PYG{k}{if} \PYG{n}{torch}\PYG{o}{.}\PYG{n}{backends}\PYG{o}{.}\PYG{n}{mps}\PYG{o}{.}\PYG{n}{is\PYGZus{}available}\PYG{p}{(}\PYG{p}{)}\PYG{p}{:}
    \PYG{n}{tensor} \PYG{o}{=} \PYG{n}{tensor}\PYG{o}{.}\PYG{n}{to}\PYG{p}{(}\PYG{l+s+s1}{\PYGZsq{}}\PYG{l+s+s1}{mps}\PYG{l+s+s1}{\PYGZsq{}}\PYG{p}{)}
    \PYG{n+nb}{print}\PYG{p}{(}\PYG{l+s+s2}{\PYGZdq{}}\PYG{l+s+s2}{After moving, device:}\PYG{l+s+s2}{\PYGZdq{}}\PYG{p}{,} \PYG{n}{tensor}\PYG{o}{.}\PYG{n}{device}\PYG{p}{)}
\end{sphinxVerbatim}

\end{sphinxuseclass}\end{sphinxVerbatimInput}
\begin{sphinxVerbatimOutput}

\begin{sphinxuseclass}{cell_output}
\begin{sphinxVerbatim}[commandchars=\\\{\}]
Before moving, device: cpu
\end{sphinxVerbatim}

\begin{sphinxVerbatim}[commandchars=\\\{\}]
After moving, device: mps:0
\end{sphinxVerbatim}

\end{sphinxuseclass}\end{sphinxVerbatimOutput}

\end{sphinxuseclass}
\sphinxAtStartPar
Running this will print:

\begin{sphinxVerbatim}[commandchars=\\\{\}]
\PYG{n}{Before} \PYG{n}{moving}\PYG{p}{,} \PYG{n}{device}\PYG{p}{:} \PYG{n}{cpu}
\PYG{n}{After} \PYG{n}{moving}\PYG{p}{,} \PYG{n}{device}\PYG{p}{:} \PYG{n}{cuda}\PYG{p}{:}\PYG{l+m+mi}{0}
\end{sphinxVerbatim}

\sphinxAtStartPar
if you have a GPU, or

\begin{sphinxVerbatim}[commandchars=\\\{\}]
\PYG{n}{Before} \PYG{n}{moving}\PYG{p}{,} \PYG{n}{device}\PYG{p}{:} \PYG{n}{cpu}
\PYG{n}{After} \PYG{n}{moving}\PYG{p}{,} \PYG{n}{device}\PYG{p}{:} \PYG{n}{mps}\PYG{p}{:}\PYG{l+m+mi}{0}
\end{sphinxVerbatim}

\sphinxAtStartPar
on M\sphinxhyphen{}chip Apple devices.

\begin{sphinxuseclass}{cell}\begin{sphinxVerbatimInput}

\begin{sphinxuseclass}{cell_input}
\begin{sphinxVerbatim}[commandchars=\\\{\}]
\PYG{n}{torch}\PYG{o}{.}\PYG{n}{set\PYGZus{}default\PYGZus{}device}\PYG{p}{(}\PYG{n}{device}\PYG{p}{)}
\end{sphinxVerbatim}

\end{sphinxuseclass}\end{sphinxVerbatimInput}

\end{sphinxuseclass}

\subsubsection{In\sphinxhyphen{}place Operations}
\label{\detokenize{04:in-place-operations}}
\sphinxAtStartPar
PyTorch differentiates between operations that create new tensors and those that modify an existing tensor \sphinxstylestrong{in\sphinxhyphen{}place}. In\sphinxhyphen{}place operations are identified by an underscore suffix in the method name. For example, \sphinxcode{\sphinxupquote{tensor.add\_(5)}} will add 5 to every element of the tensor \sphinxstylestrong{in\sphinxhyphen{}place} (mutating the tensor), whereas \sphinxcode{\sphinxupquote{tensor.add(5)}} returns a new tensor with the result and leaves the original unchanged. All functions ending in \sphinxcode{\sphinxupquote{\_}} modify the tensor in\sphinxhyphen{}place.

\sphinxAtStartPar
Why use in\sphinxhyphen{}place ops? They can save memory and sometimes a bit of compute, because you don’t allocate a new tensor for the result. However, be careful: in\sphinxhyphen{}place operations can sometimes interfere with autograd (if you modify values required to compute gradients). As a rule of thumb, avoid using in\sphinxhyphen{}place ops on tensors that require grad unless you know what you’re doing.

\sphinxAtStartPar
Let’s demonstrate the difference between an in\sphinxhyphen{}place operation and its out\sphinxhyphen{}of\sphinxhyphen{}place counterpart:

\begin{sphinxuseclass}{cell}\begin{sphinxVerbatimInput}

\begin{sphinxuseclass}{cell_input}
\begin{sphinxVerbatim}[commandchars=\\\{\}]
\PYG{n}{a} \PYG{o}{=} \PYG{n}{torch}\PYG{o}{.}\PYG{n}{ones}\PYG{p}{(}\PYG{l+m+mi}{5}\PYG{p}{)}
\PYG{n}{pprint}\PYG{p}{(}\PYG{l+s+s2}{\PYGZdq{}}\PYG{l+s+s2}{a=}\PYG{l+s+s2}{\PYGZdq{}}\PYG{p}{,} \PYG{n}{a}\PYG{p}{)}

\PYG{c+c1}{\PYGZsh{} Out\PYGZhy{}of\PYGZhy{}place addition}
\PYG{n}{b} \PYG{o}{=} \PYG{n}{a}\PYG{o}{.}\PYG{n}{add}\PYG{p}{(}\PYG{l+m+mi}{5}\PYG{p}{)}
\PYG{n+nb}{print}\PYG{p}{(}\PYG{l+s+s2}{\PYGZdq{}}\PYG{l+s+s2}{After b = a.add(5):}\PYG{l+s+s2}{\PYGZdq{}}\PYG{p}{)}
\PYG{n}{pprint}\PYG{p}{(}\PYG{l+s+s2}{\PYGZdq{}}\PYG{l+s+s2}{a=}\PYG{l+s+s2}{\PYGZdq{}}\PYG{p}{,} \PYG{n}{a}\PYG{p}{)}  \PYG{c+c1}{\PYGZsh{} a should remain unchanged}
\PYG{n}{pprint}\PYG{p}{(}\PYG{l+s+s2}{\PYGZdq{}}\PYG{l+s+s2}{b=}\PYG{l+s+s2}{\PYGZdq{}}\PYG{p}{,} \PYG{n}{b}\PYG{p}{)}  \PYG{c+c1}{\PYGZsh{} b is the new tensor}

\PYG{c+c1}{\PYGZsh{} In\PYGZhy{}place addition}
\PYG{n}{a}\PYG{o}{.}\PYG{n}{add\PYGZus{}}\PYG{p}{(}\PYG{l+m+mi}{5}\PYG{p}{)}
\PYG{n+nb}{print}\PYG{p}{(}\PYG{l+s+s2}{\PYGZdq{}}\PYG{l+s+se}{\PYGZbs{}n}\PYG{l+s+s2}{After a.add\PYGZus{}(5):}\PYG{l+s+s2}{\PYGZdq{}}\PYG{p}{)}
\PYG{n}{pprint}\PYG{p}{(}\PYG{l+s+s2}{\PYGZdq{}}\PYG{l+s+s2}{a=}\PYG{l+s+s2}{\PYGZdq{}}\PYG{p}{,} \PYG{n}{a}\PYG{p}{)}  \PYG{c+c1}{\PYGZsh{} a is now changed in\PYGZhy{}place}
\end{sphinxVerbatim}

\end{sphinxuseclass}\end{sphinxVerbatimInput}
\begin{sphinxVerbatimOutput}

\begin{sphinxuseclass}{cell_output}\begin{equation*}
\begin{split}\displaystyle a=\left[
\begin{array}{}
  1.0000 &  1.0000 &  1.0000 &  1.0000 &  1.0000
\end{array}
\right]\end{split}
\end{equation*}
\begin{sphinxVerbatim}[commandchars=\\\{\}]
After b = a.add(5):
\end{sphinxVerbatim}
\begin{equation*}
\begin{split}\displaystyle a=\left[
\begin{array}{}
  1.0000 &  1.0000 &  1.0000 &  1.0000 &  1.0000
\end{array}
\right]\end{split}
\end{equation*}\begin{equation*}
\begin{split}\displaystyle b=\left[
\begin{array}{}
  6.0000 &  6.0000 &  6.0000 &  6.0000 &  6.0000
\end{array}
\right]\end{split}
\end{equation*}
\begin{sphinxVerbatim}[commandchars=\\\{\}]
After a.add\PYGZus{}(5):
\end{sphinxVerbatim}
\begin{equation*}
\begin{split}\displaystyle a=\left[
\begin{array}{}
  6.0000 &  6.0000 &  6.0000 &  6.0000 &  6.0000
\end{array}
\right]\end{split}
\end{equation*}
\end{sphinxuseclass}\end{sphinxVerbatimOutput}

\end{sphinxuseclass}
\sphinxAtStartPar
Output should show:
\begin{enumerate}
\sphinxsetlistlabels{\arabic}{enumi}{enumii}{}{.}%
\item {} 
\sphinxAtStartPar
Initially, \sphinxcode{\sphinxupquote{a}} is \sphinxcode{\sphinxupquote{{[}1, 1, 1, 1, 1{]}}}.

\item {} 
\sphinxAtStartPar
After \sphinxcode{\sphinxupquote{b = a.add(5)}}, \sphinxcode{\sphinxupquote{a}} stays \sphinxcode{\sphinxupquote{{[}1, 1, 1, 1, 1{]}}} but \sphinxcode{\sphinxupquote{b}} is \sphinxcode{\sphinxupquote{{[}6, 6, 6, 6, 6{]}}}.

\item {} 
\sphinxAtStartPar
After \sphinxcode{\sphinxupquote{a.add\_(5)}}, now \sphinxcode{\sphinxupquote{a}} has become \sphinxcode{\sphinxupquote{{[}6, 6, 6, 6, 6{]}}} in\sphinxhyphen{}place.

\end{enumerate}

\sphinxAtStartPar
Another example with multiplication:

\begin{sphinxuseclass}{cell}\begin{sphinxVerbatimInput}

\begin{sphinxuseclass}{cell_input}
\begin{sphinxVerbatim}[commandchars=\\\{\}]
\PYG{n}{x} \PYG{o}{=} \PYG{n}{torch}\PYG{o}{.}\PYG{n}{ones}\PYG{p}{(}\PYG{l+m+mi}{3}\PYG{p}{)}
\PYG{n}{y} \PYG{o}{=} \PYG{n}{x}\PYG{o}{.}\PYG{n}{mul\PYGZus{}}\PYG{p}{(}\PYG{l+m+mi}{2}\PYG{p}{)}   \PYG{c+c1}{\PYGZsh{} in\PYGZhy{}place multiply x by 2}
\PYG{n}{pprint}\PYG{p}{(}\PYG{l+s+s2}{\PYGZdq{}}\PYG{l+s+s2}{x=}\PYG{l+s+s2}{\PYGZdq{}}\PYG{p}{,} \PYG{n}{x}\PYG{p}{)}
\PYG{n}{pprint}\PYG{p}{(}\PYG{l+s+s2}{\PYGZdq{}}\PYG{l+s+s2}{y=}\PYG{l+s+s2}{\PYGZdq{}}\PYG{p}{,} \PYG{n}{y}\PYG{p}{)}

\PYG{n}{pprint}\PYG{p}{(}\PYG{l+s+s2}{\PYGZdq{}}\PYG{l+s+se}{\PYGZbs{}\PYGZbs{}}\PYG{l+s+s2}{text}\PYG{l+s+s2}{\PYGZob{}}\PYG{l+s+s2}{x after in\PYGZhy{}place mul}\PYG{l+s+s2}{\PYGZbs{}}\PYG{l+s+s2}{\PYGZus{}: \PYGZcb{}}\PYG{l+s+s2}{\PYGZdq{}}\PYG{p}{,} \PYG{n}{x}\PYG{p}{)}
\PYG{n}{pprint}\PYG{p}{(}\PYG{l+s+s2}{\PYGZdq{}}\PYG{l+s+se}{\PYGZbs{}\PYGZbs{}}\PYG{l+s+s2}{text}\PYG{l+s+s2}{\PYGZob{}}\PYG{l+s+s2}{y (result of x.mul}\PYG{l+s+s2}{\PYGZbs{}}\PYG{l+s+s2}{\PYGZus{}(2)): \PYGZcb{}}\PYG{l+s+s2}{\PYGZdq{}}\PYG{p}{,} \PYG{n}{y}\PYG{p}{)}

\PYG{c+c1}{\PYGZsh{} Now x is [2,2,2], let\PYGZsq{}s do out\PYGZhy{}of\PYGZhy{}place multiply}
\PYG{n}{z} \PYG{o}{=} \PYG{n}{x}\PYG{o}{.}\PYG{n}{mul}\PYG{p}{(}\PYG{l+m+mi}{3}\PYG{p}{)}
\PYG{n}{pprint}\PYG{p}{(}\PYG{l+s+s2}{\PYGZdq{}}\PYG{l+s+se}{\PYGZbs{}\PYGZbs{}}\PYG{l+s+s2}{text}\PYG{l+s+s2}{\PYGZob{}}\PYG{l+s+s2}{x after out\PYGZhy{}of\PYGZhy{}place mul(3): \PYGZcb{}}\PYG{l+s+s2}{\PYGZdq{}}\PYG{p}{,} \PYG{n}{x}\PYG{p}{)}  \PYG{c+c1}{\PYGZsh{} x remains [2,2,2]}
\PYG{n}{pprint}\PYG{p}{(}\PYG{l+s+s2}{\PYGZdq{}}\PYG{l+s+se}{\PYGZbs{}\PYGZbs{}}\PYG{l+s+s2}{text}\PYG{l+s+s2}{\PYGZob{}}\PYG{l+s+s2}{z (result of x.mul(3)): \PYGZcb{}}\PYG{l+s+s2}{\PYGZdq{}}\PYG{p}{,} \PYG{n}{z}\PYG{p}{)}       \PYG{c+c1}{\PYGZsh{} z is [6,6,6]}
\end{sphinxVerbatim}

\end{sphinxuseclass}\end{sphinxVerbatimInput}
\begin{sphinxVerbatimOutput}

\begin{sphinxuseclass}{cell_output}\begin{equation*}
\begin{split}\displaystyle x=\left[
\begin{array}{}
  2.0000 &  2.0000 &  2.0000
\end{array}
\right]\end{split}
\end{equation*}\begin{equation*}
\begin{split}\displaystyle y=\left[
\begin{array}{}
  2.0000 &  2.0000 &  2.0000
\end{array}
\right]\end{split}
\end{equation*}\begin{equation*}
\begin{split}\displaystyle \text{x after in-place mul\_: }\left[
\begin{array}{}
  2.0000 &  2.0000 &  2.0000
\end{array}
\right]\end{split}
\end{equation*}\begin{equation*}
\begin{split}\displaystyle \text{y (result of x.mul\_(2)): }\left[
\begin{array}{}
  2.0000 &  2.0000 &  2.0000
\end{array}
\right]\end{split}
\end{equation*}\begin{equation*}
\begin{split}\displaystyle \text{x after out-of-place mul(3): }\left[
\begin{array}{}
  2.0000 &  2.0000 &  2.0000
\end{array}
\right]\end{split}
\end{equation*}\begin{equation*}
\begin{split}\displaystyle \text{z (result of x.mul(3)): }\left[
\begin{array}{}
  6.0000 &  6.0000 &  6.0000
\end{array}
\right]\end{split}
\end{equation*}
\end{sphinxuseclass}\end{sphinxVerbatimOutput}

\end{sphinxuseclass}
\sphinxAtStartPar
This shows that after the in\sphinxhyphen{}place multiply, \sphinxcode{\sphinxupquote{x}} (and \sphinxcode{\sphinxupquote{y}}) have doubled their values, while the out\sphinxhyphen{}of\sphinxhyphen{}place multiply by 3 produces a new tensor \sphinxcode{\sphinxupquote{z}} and leaves \sphinxcode{\sphinxupquote{x}} the same.


\section{Exercises}
\label{\detokenize{04:exercises}}

\subsection{Comparison of symbolic and automatic differentiation in a simple feedforward neural network}
\label{\detokenize{04:comparison-of-symbolic-and-automatic-differentiation-in-a-simple-feedforward-neural-network}}
\sphinxAtStartPar
In this exercise, we will illustrate how \sphinxstylestrong{automatic differentiation} (as implemented in PyTorch) matches with the result of \sphinxstylestrong{symbolic differentiation} for a simple two\sphinxhyphen{}hidden\sphinxhyphen{}neuron feedforward neural network. You should:
\begin{enumerate}
\sphinxsetlistlabels{\arabic}{enumi}{enumii}{}{.}%
\item {} 
\sphinxAtStartPar
\sphinxstylestrong{Define} a simple feedforward neural network (FNN) with two hidden neurons \(z_1\) and \(z_2\) using the values and notation from the figure below.

\item {} 
\sphinxAtStartPar
\sphinxstylestrong{Compute} the network output and the loss function (a squared error).

\item {} 
\sphinxAtStartPar
\sphinxstylestrong{Derive} the gradient of the loss function with respect to each weight \sphinxstylestrong{symbolically}.

\item {} 
\sphinxAtStartPar
\sphinxstylestrong{Use} PyTorch’s \sphinxstylestrong{automatic differentiation} to compute the same gradients \sphinxstylestrong{numerically}.

\item {} 
\sphinxAtStartPar
\sphinxstylestrong{Compare} the two results to verify they match.

\end{enumerate}

\sphinxAtStartPar
The feedforward neural network architecture that we will use in this exercise is the following:




\section{Problem setup}
\label{\detokenize{04:problem-setup}}
\sphinxAtStartPar
Let’s consider the following simplified network:
\begin{itemize}
\item {} 
\sphinxAtStartPar
\sphinxstylestrong{Input}: \(x\), a single scalar input.

\item {} 
\sphinxAtStartPar
\sphinxstylestrong{Hidden layer}: two neurons \(z_1\) and \(z_2\)

\item {} 
\sphinxAtStartPar
\sphinxstylestrong{Output}: \(\hat{y}\)

\end{itemize}

\sphinxAtStartPar
For a single input \(x\), the hidden neurons are given by:
\begin{equation*}
\begin{split}\begin{aligned}
\text{Neuron } z_1: & \quad a_1^{(1)} = w_{11}^{(1)} \cdot x, && \quad z_1 = g(a_1^{(1)}) = \frac{1}{1 + e^{-a_1^{(1)}}}, \\
\text{Neuron } z_2: & \quad a_2^{(1)} = w_{12}^{(1)} \cdot x, && \quad z_2 = g(a_2^{(1)}) = \frac{1}{1 + e^{-a_2^{(1)}}}.
\end{aligned}
\end{split}
\end{equation*}
\sphinxAtStartPar
Then the output neuron combines these two hidden neurons:
\begin{equation*}
\begin{split}\begin{aligned}
a_1^{(2)} &= w_{11}^{(2)} z_1 + w_{21}^{(2)} z_2, \\
\hat{y} &= g(a_1^{(2)}) = \frac{1}{1 + e^{-a_1^{(2)}}}.
\end{aligned}
\end{split}
\end{equation*}
\sphinxAtStartPar
We consider a \sphinxstylestrong{loss} (squared error w.r.t. a true label \(y\)):
\begin{equation*}
\begin{split}L = (y - \hat{y})^2.\end{split}
\end{equation*}
\sphinxAtStartPar
Our goal is to \sphinxstylestrong{differentiate} \(L\) with respect to the parameters:
\begin{equation*}
\begin{split}\begin{aligned}
w_{11}^{(1)},\quad w_{12}^{(1)},\quad w_{11}^{(2)},\quad w_{21}^{(2)}.
\end{aligned}
\end{split}
\end{equation*}
\sphinxAtStartPar
We will compute these derivatives in two ways:
\begin{enumerate}
\sphinxsetlistlabels{\arabic}{enumi}{enumii}{}{.}%
\item {} 
\sphinxAtStartPar
\sphinxstylestrong{Symbolic Differentiation} (analytical formulas)

\item {} 
\sphinxAtStartPar
\sphinxstylestrong{Automatic Differentiation} (via PyTorch \sphinxcode{\sphinxupquote{backward()}})

\end{enumerate}

\sphinxAtStartPar
We will show that these two results match.


\section{PyTorch implementation}
\label{\detokenize{04:pytorch-implementation}}

\subsection{Define the activation function and its derivative}
\label{\detokenize{04:define-the-activation-function-and-its-derivative}}
\sphinxAtStartPar
We will use the logistic sigmoid \(g(a) = 1 / (1 + e^{-a})\) as the activation function. Its derivative with respect to \(a\) is:
\begin{equation*}
\begin{split}\frac{d}{da} g(a) = g(a) \bigl(1 - g(a)\bigr).\end{split}
\end{equation*}
\sphinxAtStartPar
However, for clarity, we can directly implement it as given in the example code:

\begin{sphinxuseclass}{cell}\begin{sphinxVerbatimInput}

\begin{sphinxuseclass}{cell_input}
\begin{sphinxVerbatim}[commandchars=\\\{\}]
\PYG{k}{def}\PYG{+w}{ }\PYG{n+nf}{g}\PYG{p}{(}\PYG{n}{a}\PYG{p}{)}\PYG{p}{:}
\PYG{+w}{    }\PYG{l+s+sd}{\PYGZdq{}\PYGZdq{}\PYGZdq{}}
\PYG{l+s+sd}{    Sigmoid function g(a) = 1 / (1 + e\PYGZca{}\PYGZhy{}a)}
\PYG{l+s+sd}{    \PYGZdq{}\PYGZdq{}\PYGZdq{}}
    \PYG{k}{return} \PYG{l+m+mf}{1.0} \PYG{o}{/} \PYG{p}{(}\PYG{l+m+mf}{1.0} \PYG{o}{+} \PYG{n}{torch}\PYG{o}{.}\PYG{n}{exp}\PYG{p}{(}\PYG{o}{\PYGZhy{}}\PYG{n}{a}\PYG{p}{)}\PYG{p}{)}

\PYG{k}{def}\PYG{+w}{ }\PYG{n+nf}{gp}\PYG{p}{(}\PYG{n}{a}\PYG{p}{)}\PYG{p}{:}
\PYG{+w}{    }\PYG{l+s+sd}{\PYGZdq{}\PYGZdq{}\PYGZdq{}}
\PYG{l+s+sd}{    Derivative of the sigmoid function g\PYGZsq{}(a)}
\PYG{l+s+sd}{    = exp(\PYGZhy{}a) / (1 + exp(\PYGZhy{}a))\PYGZca{}2}
\PYG{l+s+sd}{    \PYGZdq{}\PYGZdq{}\PYGZdq{}}
    \PYG{k}{return} \PYG{n}{torch}\PYG{o}{.}\PYG{n}{exp}\PYG{p}{(}\PYG{o}{\PYGZhy{}}\PYG{n}{a}\PYG{p}{)} \PYG{o}{/} \PYG{p}{(}\PYG{l+m+mi}{1} \PYG{o}{+} \PYG{n}{torch}\PYG{o}{.}\PYG{n}{exp}\PYG{p}{(}\PYG{o}{\PYGZhy{}}\PYG{n}{a}\PYG{p}{)}\PYG{p}{)}\PYG{o}{*}\PYG{o}{*}\PYG{l+m+mi}{2}
\end{sphinxVerbatim}

\end{sphinxuseclass}\end{sphinxVerbatimInput}

\end{sphinxuseclass}

\subsection{Define the working variables}
\label{\detokenize{04:define-the-working-variables}}
\sphinxAtStartPar
We will define:
\begin{itemize}
\item {} 
\sphinxAtStartPar
A target value \(y\)

\item {} 
\sphinxAtStartPar
An input \(x\)

\item {} 
\sphinxAtStartPar
Four weights we wish to learn/differentiate:
\begin{itemize}
\item {} 
\sphinxAtStartPar
\(w_{11}^{(1)}\) and \(w_{12}^{(1)}\) for the hidden layer

\item {} 
\sphinxAtStartPar
\(w_{11}^{(2)}\) and \(w_{21}^{(2)}\) for the output layer

\end{itemize}

\end{itemize}

\sphinxAtStartPar
By setting \sphinxcode{\sphinxupquote{requires\_grad=True}}, PyTorch will keep track of operations involving these tensors, allowing automatic differentiation via the \sphinxcode{\sphinxupquote{backward()}} call.

\begin{sphinxuseclass}{cell}\begin{sphinxVerbatimInput}

\begin{sphinxuseclass}{cell_input}
\begin{sphinxVerbatim}[commandchars=\\\{\}]
\PYG{c+c1}{\PYGZsh{} Input}
\PYG{n}{x} \PYG{o}{=} \PYG{n}{torch}\PYG{o}{.}\PYG{n}{tensor}\PYG{p}{(}\PYG{p}{[}\PYG{o}{\PYGZhy{}}\PYG{l+m+mf}{1.7}\PYG{p}{]}\PYG{p}{)}  \PYG{c+c1}{\PYGZsh{} A single scalar input}

\PYG{c+c1}{\PYGZsh{} True target}
\PYG{n}{y} \PYG{o}{=} \PYG{n}{torch}\PYG{o}{.}\PYG{n}{tensor}\PYG{p}{(}\PYG{p}{[}\PYG{o}{\PYGZhy{}}\PYG{l+m+mf}{4.5}\PYG{p}{]}\PYG{p}{)}  \PYG{c+c1}{\PYGZsh{} A single scalar target}

\PYG{c+c1}{\PYGZsh{} Hidden layer weights}
\PYG{n}{w11\PYGZus{}1} \PYG{o}{=} \PYG{n}{torch}\PYG{o}{.}\PYG{n}{tensor}\PYG{p}{(}\PYG{p}{[}\PYG{l+m+mf}{0.13}\PYG{p}{]}\PYG{p}{,} \PYG{n}{requires\PYGZus{}grad}\PYG{o}{=}\PYG{k+kc}{True}\PYG{p}{)}  \PYG{c+c1}{\PYGZsh{} w\PYGZus{}\PYGZob{}11\PYGZcb{}\PYGZca{}\PYGZob{}(1)\PYGZcb{}}
\PYG{n}{w12\PYGZus{}1} \PYG{o}{=} \PYG{n}{torch}\PYG{o}{.}\PYG{n}{tensor}\PYG{p}{(}\PYG{p}{[}\PYG{o}{\PYGZhy{}}\PYG{l+m+mf}{2.0}\PYG{p}{]}\PYG{p}{,} \PYG{n}{requires\PYGZus{}grad}\PYG{o}{=}\PYG{k+kc}{True}\PYG{p}{)}  \PYG{c+c1}{\PYGZsh{} w\PYGZus{}\PYGZob{}12\PYGZcb{}\PYGZca{}\PYGZob{}(1)\PYGZcb{}}

\PYG{c+c1}{\PYGZsh{} Output layer weights}
\PYG{n}{w11\PYGZus{}2} \PYG{o}{=} \PYG{n}{torch}\PYG{o}{.}\PYG{n}{tensor}\PYG{p}{(}\PYG{p}{[}\PYG{l+m+mf}{1.0}\PYG{p}{]}\PYG{p}{,} \PYG{n}{requires\PYGZus{}grad}\PYG{o}{=}\PYG{k+kc}{True}\PYG{p}{)}   \PYG{c+c1}{\PYGZsh{} w\PYGZus{}\PYGZob{}11\PYGZcb{}\PYGZca{}\PYGZob{}(2)\PYGZcb{}}
\PYG{n}{w21\PYGZus{}2} \PYG{o}{=} \PYG{n}{torch}\PYG{o}{.}\PYG{n}{tensor}\PYG{p}{(}\PYG{p}{[}\PYG{l+m+mf}{0.5}\PYG{p}{]}\PYG{p}{,} \PYG{n}{requires\PYGZus{}grad}\PYG{o}{=}\PYG{k+kc}{True}\PYG{p}{)}   \PYG{c+c1}{\PYGZsh{} w\PYGZus{}\PYGZob{}21\PYGZcb{}\PYGZca{}\PYGZob{}(2)\PYGZcb{}}
\end{sphinxVerbatim}

\end{sphinxuseclass}\end{sphinxVerbatimInput}

\end{sphinxuseclass}

\subsection{Forward pass}
\label{\detokenize{04:forward-pass}}
\sphinxAtStartPar
We compute:
\begin{equation*}
\begin{split}a_1^{(1)} = w_{11}^{(1)} x,\quad
a_2^{(1)} = w_{12}^{(1)} x,\quad
z_1 = g(a_1^{(1)}),\quad z_2 = g(a_2^{(1)}),\quad
a_1^{(2)} = w_{11}^{(2)} z_1 + w_{21}^{(2)} z_2,\quad
\hat{y} = g(a_1^{(2)})
\end{split}
\end{equation*}
\sphinxAtStartPar
and finally the loss \(L = (y - \hat{y})^2\).

\begin{sphinxuseclass}{cell}\begin{sphinxVerbatimInput}

\begin{sphinxuseclass}{cell_input}
\begin{sphinxVerbatim}[commandchars=\\\{\}]
\PYG{c+c1}{\PYGZsh{} Hidden neuron 1}
\PYG{n}{a1\PYGZus{}1} \PYG{o}{=} \PYG{n}{w11\PYGZus{}1} \PYG{o}{*} \PYG{n}{x}
\PYG{n}{z1}   \PYG{o}{=} \PYG{n}{g}\PYG{p}{(}\PYG{n}{a1\PYGZus{}1}\PYG{p}{)}

\PYG{c+c1}{\PYGZsh{} Hidden neuron 2}
\PYG{n}{a2\PYGZus{}1} \PYG{o}{=} \PYG{n}{w12\PYGZus{}1} \PYG{o}{*} \PYG{n}{x}
\PYG{n}{z2}   \PYG{o}{=} \PYG{n}{g}\PYG{p}{(}\PYG{n}{a2\PYGZus{}1}\PYG{p}{)}

\PYG{c+c1}{\PYGZsh{} Output neuron}
\PYG{n}{a1\PYGZus{}2} \PYG{o}{=} \PYG{n}{w11\PYGZus{}2} \PYG{o}{*} \PYG{n}{z1} \PYG{o}{+} \PYG{n}{w21\PYGZus{}2} \PYG{o}{*} \PYG{n}{z2}
\PYG{n}{yhat} \PYG{o}{=} \PYG{n}{g}\PYG{p}{(}\PYG{n}{a1\PYGZus{}2}\PYG{p}{)}

\PYG{c+c1}{\PYGZsh{} Loss}
\PYG{n}{L} \PYG{o}{=} \PYG{p}{(}\PYG{n}{y} \PYG{o}{\PYGZhy{}} \PYG{n}{yhat}\PYG{p}{)}\PYG{o}{*}\PYG{o}{*}\PYG{l+m+mi}{2}
\end{sphinxVerbatim}

\end{sphinxuseclass}\end{sphinxVerbatimInput}

\end{sphinxuseclass}

\subsection{Symbolic gradients}
\label{\detokenize{04:symbolic-gradients}}
\sphinxAtStartPar
From the slides and from our direct symbolic differentiation, we have:
\begin{equation*}
\begin{split}\begin{aligned}
\frac{\partial L}{\partial w_{11}^{(1)}} 
&= -2 \,(y - \hat{y})\; g'(a_1^{(2)})\; w_{11}^{(2)}\; g'(a_1^{(1)})\; x,\\
\frac{\partial L}{\partial w_{12}^{(1)}} 
&= -2 \,(y - \hat{y})\; g'(a_1^{(2)})\; w_{21}^{(2)}\; g'(a_2^{(1)})\; x,\\
\frac{\partial L}{\partial w_{11}^{(2)}} 
&= -2 \,(y - \hat{y})\; g'(a_1^{(2)})\; z_1,\\
\frac{\partial L}{\partial w_{21}^{(2)}} 
&= -2 \,(y - \hat{y})\; g'(a_1^{(2)})\; z_2.
\end{aligned}\end{split}
\end{equation*}
\sphinxAtStartPar
Let’s compute these values explicitly in PyTorch (still using the same \sphinxcode{\sphinxupquote{gp}} function for \(g'(a)\)).

\begin{sphinxuseclass}{cell}\begin{sphinxVerbatimInput}

\begin{sphinxuseclass}{cell_input}
\begin{sphinxVerbatim}[commandchars=\\\{\}]
\PYG{c+c1}{\PYGZsh{} Symbolic gradients}
\PYG{n}{Sgrad\PYGZus{}L\PYGZus{}w11\PYGZus{}1} \PYG{o}{=} \PYG{o}{\PYGZhy{}}\PYG{l+m+mi}{2} \PYG{o}{*} \PYG{p}{(}\PYG{n}{y} \PYG{o}{\PYGZhy{}} \PYG{n}{yhat}\PYG{p}{)} \PYG{o}{*} \PYG{n}{gp}\PYG{p}{(}\PYG{n}{a1\PYGZus{}2}\PYG{p}{)} \PYG{o}{*} \PYG{n}{w11\PYGZus{}2} \PYG{o}{*} \PYG{n}{gp}\PYG{p}{(}\PYG{n}{a1\PYGZus{}1}\PYG{p}{)} \PYG{o}{*} \PYG{n}{x}
\PYG{n}{Sgrad\PYGZus{}L\PYGZus{}w12\PYGZus{}1} \PYG{o}{=} \PYG{o}{\PYGZhy{}}\PYG{l+m+mi}{2} \PYG{o}{*} \PYG{p}{(}\PYG{n}{y} \PYG{o}{\PYGZhy{}} \PYG{n}{yhat}\PYG{p}{)} \PYG{o}{*} \PYG{n}{gp}\PYG{p}{(}\PYG{n}{a1\PYGZus{}2}\PYG{p}{)} \PYG{o}{*} \PYG{n}{w21\PYGZus{}2} \PYG{o}{*} \PYG{n}{gp}\PYG{p}{(}\PYG{n}{a2\PYGZus{}1}\PYG{p}{)} \PYG{o}{*} \PYG{n}{x}
\PYG{n}{Sgrad\PYGZus{}L\PYGZus{}w11\PYGZus{}2} \PYG{o}{=} \PYG{o}{\PYGZhy{}}\PYG{l+m+mi}{2} \PYG{o}{*} \PYG{p}{(}\PYG{n}{y} \PYG{o}{\PYGZhy{}} \PYG{n}{yhat}\PYG{p}{)} \PYG{o}{*} \PYG{n}{gp}\PYG{p}{(}\PYG{n}{a1\PYGZus{}2}\PYG{p}{)} \PYG{o}{*} \PYG{n}{z1}
\PYG{n}{Sgrad\PYGZus{}L\PYGZus{}w21\PYGZus{}2} \PYG{o}{=} \PYG{o}{\PYGZhy{}}\PYG{l+m+mi}{2} \PYG{o}{*} \PYG{p}{(}\PYG{n}{y} \PYG{o}{\PYGZhy{}} \PYG{n}{yhat}\PYG{p}{)} \PYG{o}{*} \PYG{n}{gp}\PYG{p}{(}\PYG{n}{a1\PYGZus{}2}\PYG{p}{)} \PYG{o}{*} \PYG{n}{z2}

\PYG{n}{Sgrad\PYGZus{}L\PYGZus{}w11\PYGZus{}1}\PYG{p}{,} \PYG{n}{Sgrad\PYGZus{}L\PYGZus{}w12\PYGZus{}1}\PYG{p}{,} \PYG{n}{Sgrad\PYGZus{}L\PYGZus{}w11\PYGZus{}2}\PYG{p}{,} \PYG{n}{Sgrad\PYGZus{}L\PYGZus{}w21\PYGZus{}2}
\end{sphinxVerbatim}

\end{sphinxuseclass}\end{sphinxVerbatimInput}
\begin{sphinxVerbatimOutput}

\begin{sphinxuseclass}{cell_output}
\begin{sphinxVerbatim}[commandchars=\\\{\}]
(tensor([\PYGZhy{}0.8892], device=\PYGZsq{}mps:0\PYGZsq{}, grad\PYGZus{}fn=\PYGZlt{}MulBackward0\PYGZgt{}),
 tensor([\PYGZhy{}0.0563], device=\PYGZsq{}mps:0\PYGZsq{}, grad\PYGZus{}fn=\PYGZlt{}MulBackward0\PYGZgt{}),
 tensor([0.9424], device=\PYGZsq{}mps:0\PYGZsq{}, grad\PYGZus{}fn=\PYGZlt{}MulBackward0\PYGZgt{}),
 tensor([2.0494], device=\PYGZsq{}mps:0\PYGZsq{}, grad\PYGZus{}fn=\PYGZlt{}MulBackward0\PYGZgt{}))
\end{sphinxVerbatim}

\end{sphinxuseclass}\end{sphinxVerbatimOutput}

\end{sphinxuseclass}

\subsection{Automatic differentiation}
\label{\detokenize{04:automatic-differentiation}}
\sphinxAtStartPar
We can then rely on PyTorch’s \sphinxstylestrong{autograd} engine to compute the gradients for us. Simply call:

\begin{sphinxVerbatim}[commandchars=\\\{\}]
\PYG{n}{L}\PYG{o}{.}\PYG{n}{backward}\PYG{p}{(}\PYG{p}{)}
\end{sphinxVerbatim}

\sphinxAtStartPar
and then retrieve the \sphinxcode{\sphinxupquote{.grad}} attribute of each parameter of interest.

\begin{sphinxuseclass}{cell}\begin{sphinxVerbatimInput}

\begin{sphinxuseclass}{cell_input}
\begin{sphinxVerbatim}[commandchars=\\\{\}]
\PYG{c+c1}{\PYGZsh{} Reset gradients (in case this cell is run multiple times)}
\PYG{n}{w11\PYGZus{}1}\PYG{o}{.}\PYG{n}{grad} \PYG{o}{=} \PYG{k+kc}{None}
\PYG{n}{w12\PYGZus{}1}\PYG{o}{.}\PYG{n}{grad} \PYG{o}{=} \PYG{k+kc}{None}
\PYG{n}{w11\PYGZus{}2}\PYG{o}{.}\PYG{n}{grad} \PYG{o}{=} \PYG{k+kc}{None}
\PYG{n}{w21\PYGZus{}2}\PYG{o}{.}\PYG{n}{grad} \PYG{o}{=} \PYG{k+kc}{None}

\PYG{c+c1}{\PYGZsh{} Perform backprop}
\PYG{n}{L}\PYG{o}{.}\PYG{n}{backward}\PYG{p}{(}\PYG{p}{)}

\PYG{c+c1}{\PYGZsh{} Automatic gradients}
\PYG{n}{grad\PYGZus{}L\PYGZus{}w11\PYGZus{}1} \PYG{o}{=} \PYG{n}{w11\PYGZus{}1}\PYG{o}{.}\PYG{n}{grad}
\PYG{n}{grad\PYGZus{}L\PYGZus{}w12\PYGZus{}1} \PYG{o}{=} \PYG{n}{w12\PYGZus{}1}\PYG{o}{.}\PYG{n}{grad}
\PYG{n}{grad\PYGZus{}L\PYGZus{}w11\PYGZus{}2} \PYG{o}{=} \PYG{n}{w11\PYGZus{}2}\PYG{o}{.}\PYG{n}{grad}
\PYG{n}{grad\PYGZus{}L\PYGZus{}w21\PYGZus{}2} \PYG{o}{=} \PYG{n}{w21\PYGZus{}2}\PYG{o}{.}\PYG{n}{grad}

\PYG{n}{grad\PYGZus{}L\PYGZus{}w11\PYGZus{}1}\PYG{p}{,} \PYG{n}{grad\PYGZus{}L\PYGZus{}w12\PYGZus{}1}\PYG{p}{,} \PYG{n}{grad\PYGZus{}L\PYGZus{}w11\PYGZus{}2}\PYG{p}{,} \PYG{n}{grad\PYGZus{}L\PYGZus{}w21\PYGZus{}2}
\end{sphinxVerbatim}

\end{sphinxuseclass}\end{sphinxVerbatimInput}
\begin{sphinxVerbatimOutput}

\begin{sphinxuseclass}{cell_output}
\begin{sphinxVerbatim}[commandchars=\\\{\}]
(tensor([\PYGZhy{}0.8892], device=\PYGZsq{}mps:0\PYGZsq{}),
 tensor([\PYGZhy{}0.0563], device=\PYGZsq{}mps:0\PYGZsq{}),
 tensor([0.9424], device=\PYGZsq{}mps:0\PYGZsq{}),
 tensor([2.0494], device=\PYGZsq{}mps:0\PYGZsq{}))
\end{sphinxVerbatim}

\end{sphinxuseclass}\end{sphinxVerbatimOutput}

\end{sphinxuseclass}

\subsection{Comparison of symbolic and automatic differentiation}
\label{\detokenize{04:comparison-of-symbolic-and-automatic-differentiation}}
\sphinxAtStartPar
Finally, let’s print them side by side to verify they match (or are very close, up to floating\sphinxhyphen{}point rounding).

\begin{sphinxuseclass}{cell}\begin{sphinxVerbatimInput}

\begin{sphinxuseclass}{cell_input}
\begin{sphinxVerbatim}[commandchars=\\\{\}]
\PYG{n+nb}{print}\PYG{p}{(}\PYG{l+s+s2}{\PYGZdq{}}\PYG{l+s+s2}{Loss L =}\PYG{l+s+s2}{\PYGZdq{}}\PYG{p}{,} \PYG{n}{L}\PYG{o}{.}\PYG{n}{item}\PYG{p}{(}\PYG{p}{)}\PYG{p}{)}

\PYG{n+nb}{print}\PYG{p}{(}\PYG{l+s+s2}{\PYGZdq{}}\PYG{l+s+se}{\PYGZbs{}n}\PYG{l+s+s2}{\PYGZhy{}\PYGZhy{}\PYGZhy{} dL/dw11\PYGZca{}(2) \PYGZhy{}\PYGZhy{}\PYGZhy{}}\PYG{l+s+s2}{\PYGZdq{}}\PYG{p}{)}
\PYG{n+nb}{print}\PYG{p}{(}\PYG{l+s+s2}{\PYGZdq{}}\PYG{l+s+s2}{Auto\PYGZhy{}grad:}\PYG{l+s+s2}{\PYGZdq{}}\PYG{p}{,} \PYG{n}{grad\PYGZus{}L\PYGZus{}w11\PYGZus{}2}\PYG{o}{.}\PYG{n}{item}\PYG{p}{(}\PYG{p}{)}\PYG{p}{,} \PYG{l+s+s2}{\PYGZdq{}}\PYG{l+s+se}{\PYGZbs{}t}\PYG{l+s+s2}{ Symbolic:}\PYG{l+s+s2}{\PYGZdq{}}\PYG{p}{,} \PYG{n}{Sgrad\PYGZus{}L\PYGZus{}w11\PYGZus{}2}\PYG{o}{.}\PYG{n}{item}\PYG{p}{(}\PYG{p}{)}\PYG{p}{)}

\PYG{n+nb}{print}\PYG{p}{(}\PYG{l+s+s2}{\PYGZdq{}}\PYG{l+s+se}{\PYGZbs{}n}\PYG{l+s+s2}{\PYGZhy{}\PYGZhy{}\PYGZhy{} dL/dw21\PYGZca{}(2) \PYGZhy{}\PYGZhy{}\PYGZhy{}}\PYG{l+s+s2}{\PYGZdq{}}\PYG{p}{)}
\PYG{n+nb}{print}\PYG{p}{(}\PYG{l+s+s2}{\PYGZdq{}}\PYG{l+s+s2}{Auto\PYGZhy{}grad:}\PYG{l+s+s2}{\PYGZdq{}}\PYG{p}{,} \PYG{n}{grad\PYGZus{}L\PYGZus{}w21\PYGZus{}2}\PYG{o}{.}\PYG{n}{item}\PYG{p}{(}\PYG{p}{)}\PYG{p}{,} \PYG{l+s+s2}{\PYGZdq{}}\PYG{l+s+se}{\PYGZbs{}t}\PYG{l+s+s2}{ Symbolic:}\PYG{l+s+s2}{\PYGZdq{}}\PYG{p}{,} \PYG{n}{Sgrad\PYGZus{}L\PYGZus{}w21\PYGZus{}2}\PYG{o}{.}\PYG{n}{item}\PYG{p}{(}\PYG{p}{)}\PYG{p}{)}

\PYG{n+nb}{print}\PYG{p}{(}\PYG{l+s+s2}{\PYGZdq{}}\PYG{l+s+se}{\PYGZbs{}n}\PYG{l+s+s2}{\PYGZhy{}\PYGZhy{}\PYGZhy{} dL/dw11\PYGZca{}(1) \PYGZhy{}\PYGZhy{}\PYGZhy{}}\PYG{l+s+s2}{\PYGZdq{}}\PYG{p}{)}
\PYG{n+nb}{print}\PYG{p}{(}\PYG{l+s+s2}{\PYGZdq{}}\PYG{l+s+s2}{Auto\PYGZhy{}grad:}\PYG{l+s+s2}{\PYGZdq{}}\PYG{p}{,} \PYG{n}{grad\PYGZus{}L\PYGZus{}w11\PYGZus{}1}\PYG{o}{.}\PYG{n}{item}\PYG{p}{(}\PYG{p}{)}\PYG{p}{,} \PYG{l+s+s2}{\PYGZdq{}}\PYG{l+s+se}{\PYGZbs{}t}\PYG{l+s+s2}{ Symbolic:}\PYG{l+s+s2}{\PYGZdq{}}\PYG{p}{,} \PYG{n}{Sgrad\PYGZus{}L\PYGZus{}w11\PYGZus{}1}\PYG{o}{.}\PYG{n}{item}\PYG{p}{(}\PYG{p}{)}\PYG{p}{)}

\PYG{n+nb}{print}\PYG{p}{(}\PYG{l+s+s2}{\PYGZdq{}}\PYG{l+s+se}{\PYGZbs{}n}\PYG{l+s+s2}{\PYGZhy{}\PYGZhy{}\PYGZhy{} dL/dw12\PYGZca{}(1) \PYGZhy{}\PYGZhy{}\PYGZhy{}}\PYG{l+s+s2}{\PYGZdq{}}\PYG{p}{)}
\PYG{n+nb}{print}\PYG{p}{(}\PYG{l+s+s2}{\PYGZdq{}}\PYG{l+s+s2}{Auto\PYGZhy{}grad:}\PYG{l+s+s2}{\PYGZdq{}}\PYG{p}{,} \PYG{n}{grad\PYGZus{}L\PYGZus{}w12\PYGZus{}1}\PYG{o}{.}\PYG{n}{item}\PYG{p}{(}\PYG{p}{)}\PYG{p}{,} \PYG{l+s+s2}{\PYGZdq{}}\PYG{l+s+se}{\PYGZbs{}t}\PYG{l+s+s2}{ Symbolic:}\PYG{l+s+s2}{\PYGZdq{}}\PYG{p}{,} \PYG{n}{Sgrad\PYGZus{}L\PYGZus{}w12\PYGZus{}1}\PYG{o}{.}\PYG{n}{item}\PYG{p}{(}\PYG{p}{)}\PYG{p}{)}
\end{sphinxVerbatim}

\end{sphinxuseclass}\end{sphinxVerbatimInput}
\begin{sphinxVerbatimOutput}

\begin{sphinxuseclass}{cell_output}
\begin{sphinxVerbatim}[commandchars=\\\{\}]
Loss L = 27.21538734436035

\PYGZhy{}\PYGZhy{}\PYGZhy{} dL/dw11\PYGZca{}(2) \PYGZhy{}\PYGZhy{}\PYGZhy{}
Auto\PYGZhy{}grad: 0.942385196685791 	 Symbolic: 0.9423852562904358

\PYGZhy{}\PYGZhy{}\PYGZhy{} dL/dw21\PYGZca{}(2) \PYGZhy{}\PYGZhy{}\PYGZhy{}
Auto\PYGZhy{}grad: 2.049447774887085 	 Symbolic: 2.049448013305664

\PYGZhy{}\PYGZhy{}\PYGZhy{} dL/dw11\PYGZca{}(1) \PYGZhy{}\PYGZhy{}\PYGZhy{}
Auto\PYGZhy{}grad: \PYGZhy{}0.8891823887825012 	 Symbolic: \PYGZhy{}0.8891823887825012

\PYGZhy{}\PYGZhy{}\PYGZhy{} dL/dw12\PYGZca{}(1) \PYGZhy{}\PYGZhy{}\PYGZhy{}
Auto\PYGZhy{}grad: \PYGZhy{}0.05625968053936958 	 Symbolic: \PYGZhy{}0.05625968798995018
\end{sphinxVerbatim}

\end{sphinxuseclass}\end{sphinxVerbatimOutput}

\end{sphinxuseclass}
\sphinxAtStartPar
In this exercise, we demonstrated that for our simple 2\sphinxhyphen{}hidden\sphinxhyphen{}neuron feedforward network:
\begin{enumerate}
\sphinxsetlistlabels{\arabic}{enumi}{enumii}{}{.}%
\item {} 
\sphinxAtStartPar
\sphinxstylestrong{Symbolic differentiation} (i.e., deriving the gradient expressions by hand)

\item {} 
\sphinxAtStartPar
\sphinxstylestrong{Automatic differentiation} (i.e., via PyTorch’s \sphinxcode{\sphinxupquote{backward()}})

\end{enumerate}

\sphinxAtStartPar
yield the \sphinxstylestrong{same} results, up to numerical precision. This is a powerful illustration of why automatic differentiation is so valuable: it \sphinxstylestrong{automates} what would otherwise be a lengthy (and error\sphinxhyphen{}prone) symbolic derivation process.

\sphinxAtStartPar
Feel free to change the values of \(x\), \(y\), and the weights \(w\) to see how the gradients match up in other situations.
You can now expand this example to more complex network architectures or different activation and loss functions.


\subsection{Facial keypoints detection with PyTorch FCN and CNN architectures}
\label{\detokenize{04:facial-keypoints-detection-with-pytorch-fcn-and-cnn-architectures}}
\sphinxAtStartPar
Facial keypoint detection is the task of identifying important landmarks on a human face (e.g., corners of the eyes, nose tip, mouth corners) in an image. These keypoints are crucial for many applications in computer vision:
\begin{itemize}
\item {} 
\sphinxAtStartPar
Aligning faces for recognition or verification.

\item {} 
\sphinxAtStartPar
Tracking facial expressions for emotion detection.

\item {} 
\sphinxAtStartPar
Applying augmented reality filters or effects accurately on a face.

\end{itemize}

\sphinxAtStartPar
In this exercise, we will build fully connected neural networks and convolutional neural networks to detect \sphinxstylestrong{15 keypoints} on \sphinxstylestrong{96x96 grayscale face images}. The network will take an image as input and output the \( (x,y) \) coordinates of the \sphinxstylestrong{15 facial keypoints}.

\sphinxAtStartPar
Source: \sphinxhref{https://danielnouri.org/notes/2014/12/17/using-convolutional-neural-nets-to-detect-facial-keypoints-tutorial/}{danielnouri.org}


\subsection{Dataset overview}
\label{\detokenize{04:dataset-overview}}
\sphinxAtStartPar
We will use the \sphinxstylestrong{Kaggle Facial Keypoints Detection} dataset (\sphinxhref{https://www.kaggle.com/competitions/facial-keypoints-detection}{Kaggle challenge webpage}). This dataset consists of a training set of \sphinxstylestrong{7,049} face images (96x96 pixels, grayscale). Each image has up to \sphinxstylestrong{15 keypoint coordinates} labeled (30 values: \( x \) and \( y \) for each of 15 facial landmarks). For example, the keypoints include the centers of the eyes, the tip of the nose, and the corners of the mouth.


\subsubsection{Data format}
\label{\detokenize{04:data-format}}
\sphinxAtStartPar
The training data is provided as a CSV file (\sphinxcode{\sphinxupquote{training.csv}}). Each row contains:
\begin{itemize}
\item {} 
\sphinxAtStartPar
\sphinxstylestrong{30 columns} of keypoint coordinates: \sphinxcode{\sphinxupquote{left\_eye\_center\_x}}, \sphinxcode{\sphinxupquote{left\_eye\_center\_y}}, \sphinxcode{\sphinxupquote{right\_eye\_center\_x}}, … etc. (15 keypoints \(\times\) 2). If a keypoint was not labeled for a given image, its value is missing.

\item {} 
\sphinxAtStartPar
An \sphinxstylestrong{Image} column containing 96x96 pixel intensity values (9216 numbers) in a single string, separated by spaces.

\end{itemize}

\sphinxAtStartPar
If you downloaded this exercise from the Université Virtuelle, you should have a \sphinxcode{\sphinxupquote{training.csv}} file in the \sphinxcode{\sphinxupquote{data/}} directory. Now, let’s load the training data using pandas and examine its structure.


\subsubsection{Missing values}
\label{\detokenize{04:missing-values}}
\sphinxAtStartPar
Not all images have all 15 keypoints annotated. In fact, only about 30\% of the images have all 15 keypoints; the rest have some missing keypoint labels. Nearly all images have at least a few keypoints (e.g., eyes, nose, mouth) labeled, but many did not have the more peripheral points. In this tutorial, for simplicity, we will use only images with \sphinxstylestrong{complete keypoint data} (all 15 points). This means we’ll discard images with any missing keypoint values. (In practice, one could use data augmentation or multi\sphinxhyphen{}stage training to leverage partially labeled images, but we will not cover that here.)

\sphinxAtStartPar
In this exercise, you should:
\begin{enumerate}
\sphinxsetlistlabels{\arabic}{enumi}{enumii}{}{.}%
\item {} 
\sphinxAtStartPar
Explore the dataset

\item {} 
\sphinxAtStartPar
Preprocess the dataset to use it with neural networks including:
\begin{itemize}
\item {} 
\sphinxAtStartPar
handling missing values

\item {} 
\sphinxAtStartPar
converting the image data from strings to numeric arrays

\item {} 
\sphinxAtStartPar
normalizing the pixel values

\item {} 
\sphinxAtStartPar
reshaping the images

\item {} 
\sphinxAtStartPar
preparing the keypoint coordinates labels

\end{itemize}

\item {} 
\sphinxAtStartPar
Implement a fully connected neural network and compute the validated MSE

\item {} 
\sphinxAtStartPar
Implement a convolutional neural network and compute the validated MSE

\item {} 
\sphinxAtStartPar
Visualize the predictions of both models with respect to the true keypoints coordinates on images

\end{enumerate}

\begin{sphinxuseclass}{cell}\begin{sphinxVerbatimInput}

\begin{sphinxuseclass}{cell_input}
\begin{sphinxVerbatim}[commandchars=\\\{\}]
\PYG{k+kn}{import}\PYG{+w}{ }\PYG{n+nn}{pandas}\PYG{+w}{ }\PYG{k}{as}\PYG{+w}{ }\PYG{n+nn}{pd}

\PYG{c+c1}{\PYGZsh{} Load the training data}
\PYG{n}{df} \PYG{o}{=} \PYG{n}{pd}\PYG{o}{.}\PYG{n}{read\PYGZus{}csv}\PYG{p}{(}\PYG{l+s+s1}{\PYGZsq{}}\PYG{l+s+s1}{facial\PYGZhy{}keypoints\PYGZhy{}detection/training.csv}\PYG{l+s+s1}{\PYGZsq{}}\PYG{p}{)}
\PYG{n+nb}{print}\PYG{p}{(}\PYG{l+s+s2}{\PYGZdq{}}\PYG{l+s+s2}{Training data shape:}\PYG{l+s+s2}{\PYGZdq{}}\PYG{p}{,} \PYG{n}{df}\PYG{o}{.}\PYG{n}{shape}\PYG{p}{)}
\PYG{n}{df}\PYG{o}{.}\PYG{n}{head}\PYG{p}{(}\PYG{l+m+mi}{5}\PYG{p}{)}  \PYG{c+c1}{\PYGZsh{} show first 5 rows as an example}
\end{sphinxVerbatim}

\end{sphinxuseclass}\end{sphinxVerbatimInput}
\begin{sphinxVerbatimOutput}

\begin{sphinxuseclass}{cell_output}
\begin{sphinxVerbatim}[commandchars=\\\{\}]
Training data shape: (7049, 31)
\end{sphinxVerbatim}

\begin{sphinxVerbatim}[commandchars=\\\{\}]
   left\PYGZus{}eye\PYGZus{}center\PYGZus{}x  left\PYGZus{}eye\PYGZus{}center\PYGZus{}y  right\PYGZus{}eye\PYGZus{}center\PYGZus{}x  \PYGZbs{}
0          66.033564          39.002274           30.227008   
1          64.332936          34.970077           29.949277   
2          65.057053          34.909642           30.903789   
3          65.225739          37.261774           32.023096   
4          66.725301          39.621261           32.244810   

   right\PYGZus{}eye\PYGZus{}center\PYGZus{}y  left\PYGZus{}eye\PYGZus{}inner\PYGZus{}corner\PYGZus{}x  left\PYGZus{}eye\PYGZus{}inner\PYGZus{}corner\PYGZus{}y  \PYGZbs{}
0           36.421678                59.582075                39.647423   
1           33.448715                58.856170                35.274349   
2           34.909642                59.412000                36.320968   
3           37.261774                60.003339                39.127179   
4           38.042032                58.565890                39.621261   

   left\PYGZus{}eye\PYGZus{}outer\PYGZus{}corner\PYGZus{}x  left\PYGZus{}eye\PYGZus{}outer\PYGZus{}corner\PYGZus{}y  right\PYGZus{}eye\PYGZus{}inner\PYGZus{}corner\PYGZus{}x  \PYGZbs{}
0                73.130346                39.969997                 36.356571   
1                70.722723                36.187166                 36.034723   
2                70.984421                36.320968                 37.678105   
3                72.314713                38.380967                 37.618643   
4                72.515926                39.884466                 36.982380   

   right\PYGZus{}eye\PYGZus{}inner\PYGZus{}corner\PYGZus{}y  ...  nose\PYGZus{}tip\PYGZus{}y  mouth\PYGZus{}left\PYGZus{}corner\PYGZus{}x  \PYGZbs{}
0                 37.389402  ...   57.066803            61.195308   
1                 34.361532  ...   55.660936            56.421447   
2                 36.320968  ...   53.538947            60.822947   
3                 38.754115  ...   54.166539            65.598887   
4                 39.094852  ...   64.889521            60.671411   

   mouth\PYGZus{}left\PYGZus{}corner\PYGZus{}y  mouth\PYGZus{}right\PYGZus{}corner\PYGZus{}x  mouth\PYGZus{}right\PYGZus{}corner\PYGZus{}y  \PYGZbs{}
0            79.970165             28.614496             77.388992   
1            76.352000             35.122383             76.047660   
2            73.014316             33.726316             72.732000   
3            72.703722             37.245496             74.195478   
4            77.523239             31.191755             76.997301   

   mouth\PYGZus{}center\PYGZus{}top\PYGZus{}lip\PYGZus{}x  mouth\PYGZus{}center\PYGZus{}top\PYGZus{}lip\PYGZus{}y  mouth\PYGZus{}center\PYGZus{}bottom\PYGZus{}lip\PYGZus{}x  \PYGZbs{}
0               43.312602               72.935459                  43.130707   
1               46.684596               70.266553                  45.467915   
2               47.274947               70.191789                  47.274947   
3               50.303165               70.091687                  51.561183   
4               44.962748               73.707387                  44.227141   

   mouth\PYGZus{}center\PYGZus{}bottom\PYGZus{}lip\PYGZus{}y  \PYGZbs{}
0                  84.485774   
1                  85.480170   
2                  78.659368   
3                  78.268383   
4                  86.871166   

                                               Image  
0  238 236 237 238 240 240 239 241 241 243 240 23...  
1  219 215 204 196 204 211 212 200 180 168 178 19...  
2  144 142 159 180 188 188 184 180 167 132 84 59 ...  
3  193 192 193 194 194 194 193 192 168 111 50 12 ...  
4  147 148 160 196 215 214 216 217 219 220 206 18...  

[5 rows x 31 columns]
\end{sphinxVerbatim}

\end{sphinxuseclass}\end{sphinxVerbatimOutput}

\end{sphinxuseclass}
\sphinxAtStartPar
This will show the DataFrame with columns for each keypoint and the “Image” column containing pixel values. Let’s check how many values are present in each column and how many are missing:

\begin{sphinxuseclass}{cell}\begin{sphinxVerbatimInput}

\begin{sphinxuseclass}{cell_input}
\begin{sphinxVerbatim}[commandchars=\\\{\}]
\PYG{c+c1}{\PYGZsh{} Count non\PYGZhy{}null values in each column}
\PYG{n+nb}{print}\PYG{p}{(}\PYG{n}{df}\PYG{o}{.}\PYG{n}{count}\PYG{p}{(}\PYG{p}{)}\PYG{p}{)}
\end{sphinxVerbatim}

\end{sphinxuseclass}\end{sphinxVerbatimInput}
\begin{sphinxVerbatimOutput}

\begin{sphinxuseclass}{cell_output}
\begin{sphinxVerbatim}[commandchars=\\\{\}]
left\PYGZus{}eye\PYGZus{}center\PYGZus{}x            7039
left\PYGZus{}eye\PYGZus{}center\PYGZus{}y            7039
right\PYGZus{}eye\PYGZus{}center\PYGZus{}x           7036
right\PYGZus{}eye\PYGZus{}center\PYGZus{}y           7036
left\PYGZus{}eye\PYGZus{}inner\PYGZus{}corner\PYGZus{}x      2271
left\PYGZus{}eye\PYGZus{}inner\PYGZus{}corner\PYGZus{}y      2271
left\PYGZus{}eye\PYGZus{}outer\PYGZus{}corner\PYGZus{}x      2267
left\PYGZus{}eye\PYGZus{}outer\PYGZus{}corner\PYGZus{}y      2267
right\PYGZus{}eye\PYGZus{}inner\PYGZus{}corner\PYGZus{}x     2268
right\PYGZus{}eye\PYGZus{}inner\PYGZus{}corner\PYGZus{}y     2268
right\PYGZus{}eye\PYGZus{}outer\PYGZus{}corner\PYGZus{}x     2268
right\PYGZus{}eye\PYGZus{}outer\PYGZus{}corner\PYGZus{}y     2268
left\PYGZus{}eyebrow\PYGZus{}inner\PYGZus{}end\PYGZus{}x     2270
left\PYGZus{}eyebrow\PYGZus{}inner\PYGZus{}end\PYGZus{}y     2270
left\PYGZus{}eyebrow\PYGZus{}outer\PYGZus{}end\PYGZus{}x     2225
left\PYGZus{}eyebrow\PYGZus{}outer\PYGZus{}end\PYGZus{}y     2225
right\PYGZus{}eyebrow\PYGZus{}inner\PYGZus{}end\PYGZus{}x    2270
right\PYGZus{}eyebrow\PYGZus{}inner\PYGZus{}end\PYGZus{}y    2270
right\PYGZus{}eyebrow\PYGZus{}outer\PYGZus{}end\PYGZus{}x    2236
right\PYGZus{}eyebrow\PYGZus{}outer\PYGZus{}end\PYGZus{}y    2236
nose\PYGZus{}tip\PYGZus{}x                   7049
nose\PYGZus{}tip\PYGZus{}y                   7049
mouth\PYGZus{}left\PYGZus{}corner\PYGZus{}x          2269
mouth\PYGZus{}left\PYGZus{}corner\PYGZus{}y          2269
mouth\PYGZus{}right\PYGZus{}corner\PYGZus{}x         2270
mouth\PYGZus{}right\PYGZus{}corner\PYGZus{}y         2270
mouth\PYGZus{}center\PYGZus{}top\PYGZus{}lip\PYGZus{}x       2275
mouth\PYGZus{}center\PYGZus{}top\PYGZus{}lip\PYGZus{}y       2275
mouth\PYGZus{}center\PYGZus{}bottom\PYGZus{}lip\PYGZus{}x    7016
mouth\PYGZus{}center\PYGZus{}bottom\PYGZus{}lip\PYGZus{}y    7016
Image                        7049
dtype: int64
\end{sphinxVerbatim}

\end{sphinxuseclass}\end{sphinxVerbatimOutput}

\end{sphinxuseclass}
\sphinxAtStartPar
From this, we can see which keypoint columns have missing values. We expect some columns to have 7049 values (no missing data for those keypoints) and others to have fewer (indicating missing labels). If a column has fewer than 7049 non\sphinxhyphen{}null entries, the difference is the number of missing values.


\section{Data exploration}
\label{\detokenize{04:data-exploration}}
\sphinxAtStartPar
Before training a model, it’s important to understand the dataset. We will:
\begin{enumerate}
\sphinxsetlistlabels{\arabic}{enumi}{enumii}{}{.}%
\item {} 
\sphinxAtStartPar
Identify how many keypoint annotations are missing.

\item {} 
\sphinxAtStartPar
Visualize a few images with their annotated keypoints to get an intuition of the task.

\end{enumerate}


\subsection{Missing Values}
\label{\detokenize{04:id1}}
\sphinxAtStartPar
We already counted non\sphinxhyphen{}null values. Let’s quantify missing data more directly:

\begin{sphinxuseclass}{cell}\begin{sphinxVerbatimInput}

\begin{sphinxuseclass}{cell_input}
\begin{sphinxVerbatim}[commandchars=\\\{\}]
\PYG{c+c1}{\PYGZsh{} Check total missing values per column}
\PYG{n}{missing\PYGZus{}counts} \PYG{o}{=} \PYG{n}{df}\PYG{o}{.}\PYG{n}{isnull}\PYG{p}{(}\PYG{p}{)}\PYG{o}{.}\PYG{n}{sum}\PYG{p}{(}\PYG{p}{)}
\PYG{n+nb}{print}\PYG{p}{(}\PYG{l+s+s2}{\PYGZdq{}}\PYG{l+s+s2}{Missing values per column:}\PYG{l+s+s2}{\PYGZdq{}}\PYG{p}{)}
\PYG{n+nb}{print}\PYG{p}{(}\PYG{n}{missing\PYGZus{}counts}\PYG{p}{)}
\end{sphinxVerbatim}

\end{sphinxuseclass}\end{sphinxVerbatimInput}
\begin{sphinxVerbatimOutput}

\begin{sphinxuseclass}{cell_output}
\begin{sphinxVerbatim}[commandchars=\\\{\}]
Missing values per column:
left\PYGZus{}eye\PYGZus{}center\PYGZus{}x              10
left\PYGZus{}eye\PYGZus{}center\PYGZus{}y              10
right\PYGZus{}eye\PYGZus{}center\PYGZus{}x             13
right\PYGZus{}eye\PYGZus{}center\PYGZus{}y             13
left\PYGZus{}eye\PYGZus{}inner\PYGZus{}corner\PYGZus{}x      4778
left\PYGZus{}eye\PYGZus{}inner\PYGZus{}corner\PYGZus{}y      4778
left\PYGZus{}eye\PYGZus{}outer\PYGZus{}corner\PYGZus{}x      4782
left\PYGZus{}eye\PYGZus{}outer\PYGZus{}corner\PYGZus{}y      4782
right\PYGZus{}eye\PYGZus{}inner\PYGZus{}corner\PYGZus{}x     4781
right\PYGZus{}eye\PYGZus{}inner\PYGZus{}corner\PYGZus{}y     4781
right\PYGZus{}eye\PYGZus{}outer\PYGZus{}corner\PYGZus{}x     4781
right\PYGZus{}eye\PYGZus{}outer\PYGZus{}corner\PYGZus{}y     4781
left\PYGZus{}eyebrow\PYGZus{}inner\PYGZus{}end\PYGZus{}x     4779
left\PYGZus{}eyebrow\PYGZus{}inner\PYGZus{}end\PYGZus{}y     4779
left\PYGZus{}eyebrow\PYGZus{}outer\PYGZus{}end\PYGZus{}x     4824
left\PYGZus{}eyebrow\PYGZus{}outer\PYGZus{}end\PYGZus{}y     4824
right\PYGZus{}eyebrow\PYGZus{}inner\PYGZus{}end\PYGZus{}x    4779
right\PYGZus{}eyebrow\PYGZus{}inner\PYGZus{}end\PYGZus{}y    4779
right\PYGZus{}eyebrow\PYGZus{}outer\PYGZus{}end\PYGZus{}x    4813
right\PYGZus{}eyebrow\PYGZus{}outer\PYGZus{}end\PYGZus{}y    4813
nose\PYGZus{}tip\PYGZus{}x                      0
nose\PYGZus{}tip\PYGZus{}y                      0
mouth\PYGZus{}left\PYGZus{}corner\PYGZus{}x          4780
mouth\PYGZus{}left\PYGZus{}corner\PYGZus{}y          4780
mouth\PYGZus{}right\PYGZus{}corner\PYGZus{}x         4779
mouth\PYGZus{}right\PYGZus{}corner\PYGZus{}y         4779
mouth\PYGZus{}center\PYGZus{}top\PYGZus{}lip\PYGZus{}x       4774
mouth\PYGZus{}center\PYGZus{}top\PYGZus{}lip\PYGZus{}y       4774
mouth\PYGZus{}center\PYGZus{}bottom\PYGZus{}lip\PYGZus{}x      33
mouth\PYGZus{}center\PYGZus{}bottom\PYGZus{}lip\PYGZus{}y      33
Image                           0
dtype: int64
\end{sphinxVerbatim}

\end{sphinxuseclass}\end{sphinxVerbatimOutput}

\end{sphinxuseclass}
\sphinxAtStartPar
You should see that some keypoint columns (e.g., \sphinxcode{\sphinxupquote{left\_eye\_center\_x}}) have very few missing values, while others (e.g., some mouth or eyebrow points) have thousands of missing entries. In fact, only 2,140 images have \sphinxstylestrong{all} 15 keypoints labeled, meaning the remainder have some missing points. We will handle this in preprocessing by removing incomplete rows.


\subsection{Sample visualization}
\label{\detokenize{04:sample-visualization}}
\sphinxAtStartPar
Let’s visualize some of the images with their keypoints to understand the data. We will pick a few examples (ensuring they have no missing keypoints) and plot them in a grid.

\begin{sphinxuseclass}{cell}\begin{sphinxVerbatimInput}

\begin{sphinxuseclass}{cell_input}
\begin{sphinxVerbatim}[commandchars=\\\{\}]
\PYG{n}{df\PYGZus{}complete} \PYG{o}{=} \PYG{n}{df}\PYG{o}{.}\PYG{n}{dropna}\PYG{p}{(}\PYG{p}{)}\PYG{o}{.}\PYG{n}{reset\PYGZus{}index}\PYG{p}{(}\PYG{n}{drop}\PYG{o}{=}\PYG{k+kc}{True}\PYG{p}{)}
\PYG{n+nb}{print}\PYG{p}{(}\PYG{l+s+s2}{\PYGZdq{}}\PYG{l+s+s2}{Number of complete cases (no missing keypoints):}\PYG{l+s+s2}{\PYGZdq{}}\PYG{p}{,} \PYG{n+nb}{len}\PYG{p}{(}\PYG{n}{df\PYGZus{}complete}\PYG{p}{)}\PYG{p}{)}

\PYG{n}{sample\PYGZus{}images} \PYG{o}{=} \PYG{n}{df\PYGZus{}complete}\PYG{p}{[}\PYG{l+s+s1}{\PYGZsq{}}\PYG{l+s+s1}{Image}\PYG{l+s+s1}{\PYGZsq{}}\PYG{p}{]}\PYG{o}{.}\PYG{n}{iloc}\PYG{p}{[}\PYG{p}{:}\PYG{l+m+mi}{9}\PYG{p}{]}\PYG{o}{.}\PYG{n}{values}
\PYG{n}{sample\PYGZus{}images} \PYG{o}{=} \PYG{p}{[}\PYG{n}{np}\PYG{o}{.}\PYG{n}{fromstring}\PYG{p}{(}\PYG{n}{img\PYGZus{}str}\PYG{p}{,} \PYG{n}{sep}\PYG{o}{=}\PYG{l+s+s1}{\PYGZsq{}}\PYG{l+s+s1}{ }\PYG{l+s+s1}{\PYGZsq{}}\PYG{p}{)}\PYG{o}{.}\PYG{n}{reshape}\PYG{p}{(}\PYG{l+m+mi}{96}\PYG{p}{,} \PYG{l+m+mi}{96}\PYG{p}{)} \PYG{k}{for} \PYG{n}{img\PYGZus{}str} \PYG{o+ow}{in} \PYG{n}{sample\PYGZus{}images}\PYG{p}{]}
\PYG{n}{sample\PYGZus{}keypoints} \PYG{o}{=} \PYG{n}{df\PYGZus{}complete}\PYG{o}{.}\PYG{n}{iloc}\PYG{p}{[}\PYG{p}{:}\PYG{l+m+mi}{9}\PYG{p}{,} \PYG{p}{:}\PYG{o}{\PYGZhy{}}\PYG{l+m+mi}{1}\PYG{p}{]}\PYG{o}{.}\PYG{n}{values}  \PYG{c+c1}{\PYGZsh{} first 30 columns are keypoints}

\PYG{n}{fig}\PYG{p}{,} \PYG{n}{axes} \PYG{o}{=} \PYG{n}{plt}\PYG{o}{.}\PYG{n}{subplots}\PYG{p}{(}\PYG{l+m+mi}{3}\PYG{p}{,} \PYG{l+m+mi}{3}\PYG{p}{,} \PYG{n}{figsize}\PYG{o}{=}\PYG{p}{(}\PYG{l+m+mi}{10}\PYG{p}{,} \PYG{l+m+mi}{10}\PYG{p}{)}\PYG{p}{)}
\PYG{n}{axes} \PYG{o}{=} \PYG{n}{axes}\PYG{o}{.}\PYG{n}{flatten}\PYG{p}{(}\PYG{p}{)}
\PYG{k}{for} \PYG{n}{i} \PYG{o+ow}{in} \PYG{n+nb}{range}\PYG{p}{(}\PYG{l+m+mi}{9}\PYG{p}{)}\PYG{p}{:}
    \PYG{n}{axes}\PYG{p}{[}\PYG{n}{i}\PYG{p}{]}\PYG{o}{.}\PYG{n}{imshow}\PYG{p}{(}\PYG{n}{sample\PYGZus{}images}\PYG{p}{[}\PYG{n}{i}\PYG{p}{]}\PYG{p}{,} \PYG{n}{cmap}\PYG{o}{=}\PYG{l+s+s1}{\PYGZsq{}}\PYG{l+s+s1}{gray}\PYG{l+s+s1}{\PYGZsq{}}\PYG{p}{)}
    \PYG{c+c1}{\PYGZsh{} Plot keypoints as red dots}
    \PYG{n}{kp} \PYG{o}{=} \PYG{n}{sample\PYGZus{}keypoints}\PYG{p}{[}\PYG{n}{i}\PYG{p}{]}
    \PYG{n}{axes}\PYG{p}{[}\PYG{n}{i}\PYG{p}{]}\PYG{o}{.}\PYG{n}{scatter}\PYG{p}{(}\PYG{n}{kp}\PYG{p}{[}\PYG{l+m+mi}{0}\PYG{p}{:}\PYG{p}{:}\PYG{l+m+mi}{2}\PYG{p}{]}\PYG{p}{,} \PYG{n}{kp}\PYG{p}{[}\PYG{l+m+mi}{1}\PYG{p}{:}\PYG{p}{:}\PYG{l+m+mi}{2}\PYG{p}{]}\PYG{p}{,} \PYG{n}{c}\PYG{o}{=}\PYG{l+s+s1}{\PYGZsq{}}\PYG{l+s+s1}{red}\PYG{l+s+s1}{\PYGZsq{}}\PYG{p}{,} \PYG{n}{s}\PYG{o}{=}\PYG{l+m+mi}{15}\PYG{p}{)}
    \PYG{n}{axes}\PYG{p}{[}\PYG{n}{i}\PYG{p}{]}\PYG{o}{.}\PYG{n}{axis}\PYG{p}{(}\PYG{l+s+s1}{\PYGZsq{}}\PYG{l+s+s1}{off}\PYG{l+s+s1}{\PYGZsq{}}\PYG{p}{)}
\PYG{n}{plt}\PYG{o}{.}\PYG{n}{tight\PYGZus{}layout}\PYG{p}{(}\PYG{p}{)}
\PYG{n}{plt}\PYG{o}{.}\PYG{n}{show}\PYG{p}{(}\PYG{p}{)}
\end{sphinxVerbatim}

\end{sphinxuseclass}\end{sphinxVerbatimInput}
\begin{sphinxVerbatimOutput}

\begin{sphinxuseclass}{cell_output}
\begin{sphinxVerbatim}[commandchars=\\\{\}]
Number of complete cases (no missing keypoints): 2140
\end{sphinxVerbatim}

\noindent\sphinxincludegraphics{{4db5c26022f0cc6fd9047647a3e6d2c1fba2cbb555a021c9d72e10108c03dcd3}.png}

\end{sphinxuseclass}\end{sphinxVerbatimOutput}

\end{sphinxuseclass}
\sphinxAtStartPar
We see 9 face images with red dots marking the locations of eyes, nose, mouth, etc. This visual check helps verify that the keypoints line up with facial features (for instance, dots around the eyes, nose, and mouth).


\section{Data preprocessing}
\label{\detokenize{04:data-preprocessing}}
\sphinxAtStartPar
Next, we need to prepare the data for training a PyTorch model. This involves:
\begin{enumerate}
\sphinxsetlistlabels{\arabic}{enumi}{enumii}{}{.}%
\item {} 
\sphinxAtStartPar
\sphinxstylestrong{Handling missing keypoint values}: We’ll simplify by using only the images that have all 15 keypoints labeled.

\item {} 
\sphinxAtStartPar
\sphinxstylestrong{Converting the image data} from strings to numeric arrays.

\item {} 
\sphinxAtStartPar
\sphinxstylestrong{Normalizing} the pixel values (scaling them to \([0,1]\)) by dividing by 255.

\item {} 
\sphinxAtStartPar
\sphinxstylestrong{Reshaping} the images to the proper format (\( (N, 1, 96, 96) \) for a single\sphinxhyphen{}channel image).

\item {} 
\sphinxAtStartPar
\sphinxstylestrong{Preparing the keypoint coordinates} array of shape (\( N, 30 \)). Each sample has 30 target values: 15 \( x \)\sphinxhyphen{}coordinates and 15 \( y \)\sphinxhyphen{}coordinates.

\end{enumerate}

\sphinxAtStartPar
We encourage you to try at home to use more sophisticated data augmentation strategies or partial labeling approaches, but we’ll keep things simple here.

\begin{sphinxuseclass}{cell}\begin{sphinxVerbatimInput}

\begin{sphinxuseclass}{cell_input}
\begin{sphinxVerbatim}[commandchars=\\\{\}]
\PYG{c+c1}{\PYGZsh{} 1. Drop rows with missing values}
\PYG{n}{df} \PYG{o}{=} \PYG{n}{df}\PYG{o}{.}\PYG{n}{dropna}\PYG{p}{(}\PYG{p}{)}\PYG{o}{.}\PYG{n}{reset\PYGZus{}index}\PYG{p}{(}\PYG{n}{drop}\PYG{o}{=}\PYG{k+kc}{True}\PYG{p}{)}
\PYG{n+nb}{print}\PYG{p}{(}\PYG{l+s+s2}{\PYGZdq{}}\PYG{l+s+s2}{Remaining training samples after dropping missing:}\PYG{l+s+s2}{\PYGZdq{}}\PYG{p}{,} \PYG{n+nb}{len}\PYG{p}{(}\PYG{n}{df}\PYG{p}{)}\PYG{p}{)}

\PYG{c+c1}{\PYGZsh{} 2. Convert the \PYGZsq{}Image\PYGZsq{} column from strings to numpy arrays of shape (96,96)}
\PYG{n}{df}\PYG{p}{[}\PYG{l+s+s1}{\PYGZsq{}}\PYG{l+s+s1}{Image}\PYG{l+s+s1}{\PYGZsq{}}\PYG{p}{]} \PYG{o}{=} \PYG{n}{df}\PYG{p}{[}\PYG{l+s+s1}{\PYGZsq{}}\PYG{l+s+s1}{Image}\PYG{l+s+s1}{\PYGZsq{}}\PYG{p}{]}\PYG{o}{.}\PYG{n}{apply}\PYG{p}{(}\PYG{k}{lambda} \PYG{n}{img}\PYG{p}{:} \PYG{n}{np}\PYG{o}{.}\PYG{n}{fromstring}\PYG{p}{(}\PYG{n}{img}\PYG{p}{,} \PYG{n}{sep}\PYG{o}{=}\PYG{l+s+s1}{\PYGZsq{}}\PYG{l+s+s1}{ }\PYG{l+s+s1}{\PYGZsq{}}\PYG{p}{,} \PYG{n}{dtype}\PYG{o}{=}\PYG{n}{np}\PYG{o}{.}\PYG{n}{float32}\PYG{p}{)}\PYG{p}{)}

\PYG{c+c1}{\PYGZsh{} 3. Normalize pixel values to [0,1]}
\PYG{n}{df}\PYG{p}{[}\PYG{l+s+s1}{\PYGZsq{}}\PYG{l+s+s1}{Image}\PYG{l+s+s1}{\PYGZsq{}}\PYG{p}{]} \PYG{o}{=} \PYG{n}{df}\PYG{p}{[}\PYG{l+s+s1}{\PYGZsq{}}\PYG{l+s+s1}{Image}\PYG{l+s+s1}{\PYGZsq{}}\PYG{p}{]}\PYG{o}{.}\PYG{n}{apply}\PYG{p}{(}\PYG{k}{lambda} \PYG{n}{img}\PYG{p}{:} \PYG{n}{img}\PYG{o}{/}\PYG{l+m+mf}{255.0}\PYG{p}{)}

\PYG{c+c1}{\PYGZsh{} 4. Convert list of image arrays to a 4D numpy array for model input}
\PYG{n}{X\PYGZus{}fcn} \PYG{o}{=} \PYG{n}{np}\PYG{o}{.}\PYG{n}{vstack}\PYG{p}{(}\PYG{n}{df}\PYG{p}{[}\PYG{l+s+s1}{\PYGZsq{}}\PYG{l+s+s1}{Image}\PYG{l+s+s1}{\PYGZsq{}}\PYG{p}{]}\PYG{o}{.}\PYG{n}{values}\PYG{p}{)}  \PYG{c+c1}{\PYGZsh{} shape (N, 9216)}
\PYG{n}{X\PYGZus{}cnn} \PYG{o}{=} \PYG{n}{np}\PYG{o}{.}\PYG{n}{vstack}\PYG{p}{(}\PYG{n}{df}\PYG{p}{[}\PYG{l+s+s1}{\PYGZsq{}}\PYG{l+s+s1}{Image}\PYG{l+s+s1}{\PYGZsq{}}\PYG{p}{]}\PYG{o}{.}\PYG{n}{values}\PYG{p}{)}\PYG{o}{.}\PYG{n}{reshape}\PYG{p}{(}\PYG{o}{\PYGZhy{}}\PYG{l+m+mi}{1}\PYG{p}{,} \PYG{l+m+mi}{1}\PYG{p}{,} \PYG{l+m+mi}{96}\PYG{p}{,} \PYG{l+m+mi}{96}\PYG{p}{)}\PYG{o}{.}\PYG{n}{astype}\PYG{p}{(}\PYG{n}{np}\PYG{o}{.}\PYG{n}{float32}\PYG{p}{)}  \PYG{c+c1}{\PYGZsh{} (N, 1, 96, 96)}

\PYG{c+c1}{\PYGZsh{} 5. Prepare keypoint targets as a numpy array}
\PYG{n}{y} \PYG{o}{=} \PYG{n}{df}\PYG{p}{[}\PYG{n}{df}\PYG{o}{.}\PYG{n}{columns}\PYG{p}{[}\PYG{p}{:}\PYG{o}{\PYGZhy{}}\PYG{l+m+mi}{1}\PYG{p}{]}\PYG{p}{]}\PYG{o}{.}\PYG{n}{values}  \PYG{c+c1}{\PYGZsh{} all columns except the \PYGZsq{}Image\PYGZsq{} column}
\PYG{n}{y} \PYG{o}{=} \PYG{n}{y}\PYG{o}{.}\PYG{n}{astype}\PYG{p}{(}\PYG{n}{np}\PYG{o}{.}\PYG{n}{float32}\PYG{p}{)}
\PYG{n+nb}{print}\PYG{p}{(}\PYG{l+s+s2}{\PYGZdq{}}\PYG{l+s+s2}{X\PYGZus{}fcn shape:}\PYG{l+s+s2}{\PYGZdq{}}\PYG{p}{,} \PYG{n}{X\PYGZus{}fcn}\PYG{o}{.}\PYG{n}{shape}\PYG{p}{,} \PYG{l+s+s2}{\PYGZdq{}}\PYG{l+s+s2}{y shape:}\PYG{l+s+s2}{\PYGZdq{}}\PYG{p}{,} \PYG{n}{y}\PYG{o}{.}\PYG{n}{shape}\PYG{p}{)}
\PYG{n+nb}{print}\PYG{p}{(}\PYG{l+s+s2}{\PYGZdq{}}\PYG{l+s+s2}{X\PYGZus{}cnn shape:}\PYG{l+s+s2}{\PYGZdq{}}\PYG{p}{,} \PYG{n}{X\PYGZus{}cnn}\PYG{o}{.}\PYG{n}{shape}\PYG{p}{,} \PYG{l+s+s2}{\PYGZdq{}}\PYG{l+s+s2}{y shape:}\PYG{l+s+s2}{\PYGZdq{}}\PYG{p}{,} \PYG{n}{y}\PYG{o}{.}\PYG{n}{shape}\PYG{p}{)}
\end{sphinxVerbatim}

\end{sphinxuseclass}\end{sphinxVerbatimInput}
\begin{sphinxVerbatimOutput}

\begin{sphinxuseclass}{cell_output}
\begin{sphinxVerbatim}[commandchars=\\\{\}]
Remaining training samples after dropping missing: 2140
\end{sphinxVerbatim}

\begin{sphinxVerbatim}[commandchars=\\\{\}]
X\PYGZus{}fcn shape: (2140, 9216) y shape: (2140, 30)
X\PYGZus{}cnn shape: (2140, 1, 96, 96) y shape: (2140, 30)
\end{sphinxVerbatim}

\end{sphinxuseclass}\end{sphinxVerbatimOutput}

\end{sphinxuseclass}

\section{Fully connected neural network (FCN)}
\label{\detokenize{04:fully-connected-neural-network-fcn}}

\subsection{FCN architecture}
\label{\detokenize{04:fcn-architecture}}
\sphinxAtStartPar
We create a \sphinxstylestrong{Fully Connected Neural Network (FCN)} using PyTorch to regress the 30 keypoint coordinates \((x, y)\) directly from flattened grayscale images (of size \(96 \times 96 = 9216\) pixels). Feel free to adjust the hidden layer sizes and dropout probabilities based on your experimentation.

\sphinxAtStartPar
The architecture of our FCN is as follows:
\begin{itemize}
\item {} 
\sphinxAtStartPar
\sphinxstylestrong{Fully Connected Layer 1:}
\begin{itemize}
\item {} 
\sphinxAtStartPar
Input: 9216 features (flattened image)

\item {} 
\sphinxAtStartPar
Output: 1024 neurons

\item {} 
\sphinxAtStartPar
Activation: ReLU

\item {} 
\sphinxAtStartPar
Dropout: 0.3 probability (to reduce overfitting)

\end{itemize}

\item {} 
\sphinxAtStartPar
\sphinxstylestrong{Fully Connected Layer 2:}
\begin{itemize}
\item {} 
\sphinxAtStartPar
Input: 1024 neurons

\item {} 
\sphinxAtStartPar
Output: 512 neurons

\item {} 
\sphinxAtStartPar
Activation: ReLU

\item {} 
\sphinxAtStartPar
Dropout: 0.3 probability

\end{itemize}

\item {} 
\sphinxAtStartPar
\sphinxstylestrong{Output Layer:}
\begin{itemize}
\item {} 
\sphinxAtStartPar
Input: 512 neurons

\item {} 
\sphinxAtStartPar
Output: 30 neurons (coordinates for 15 keypoints)

\item {} 
\sphinxAtStartPar
Linear activation (no final nonlinearity, since this is a regression output)

\end{itemize}

\end{itemize}

\sphinxAtStartPar
This FCN is trained using the \sphinxstylestrong{Mean Squared Error (MSE)} loss function to predict the \((x, y)\) coordinates of 15 facial keypoints from the flattened images.

\begin{sphinxuseclass}{cell}\begin{sphinxVerbatimInput}

\begin{sphinxuseclass}{cell_input}
\begin{sphinxVerbatim}[commandchars=\\\{\}]
\PYG{k}{class}\PYG{+w}{ }\PYG{n+nc}{KeypointFCN}\PYG{p}{(}\PYG{n}{nn}\PYG{o}{.}\PYG{n}{Module}\PYG{p}{)}\PYG{p}{:}
    \PYG{k}{def}\PYG{+w}{ }\PYG{n+nf+fm}{\PYGZus{}\PYGZus{}init\PYGZus{}\PYGZus{}}\PYG{p}{(}\PYG{n+nb+bp}{self}\PYG{p}{)}\PYG{p}{:}
        \PYG{n+nb}{super}\PYG{p}{(}\PYG{n}{KeypointFCN}\PYG{p}{,} \PYG{n+nb+bp}{self}\PYG{p}{)}\PYG{o}{.}\PYG{n+nf+fm}{\PYGZus{}\PYGZus{}init\PYGZus{}\PYGZus{}}\PYG{p}{(}\PYG{p}{)}
        \PYG{c+c1}{\PYGZsh{} We have 96*96 = 9216 input features per image}
        \PYG{n+nb+bp}{self}\PYG{o}{.}\PYG{n}{fc1} \PYG{o}{=} \PYG{n}{nn}\PYG{o}{.}\PYG{n}{Linear}\PYG{p}{(}\PYG{l+m+mi}{96}\PYG{o}{*}\PYG{l+m+mi}{96}\PYG{p}{,} \PYG{l+m+mi}{1024}\PYG{p}{)}
        \PYG{n+nb+bp}{self}\PYG{o}{.}\PYG{n}{dropout1} \PYG{o}{=} \PYG{n}{nn}\PYG{o}{.}\PYG{n}{Dropout}\PYG{p}{(}\PYG{l+m+mf}{0.3}\PYG{p}{)}
        \PYG{n+nb+bp}{self}\PYG{o}{.}\PYG{n}{fc2} \PYG{o}{=} \PYG{n}{nn}\PYG{o}{.}\PYG{n}{Linear}\PYG{p}{(}\PYG{l+m+mi}{1024}\PYG{p}{,} \PYG{l+m+mi}{512}\PYG{p}{)}
        \PYG{n+nb+bp}{self}\PYG{o}{.}\PYG{n}{dropout2} \PYG{o}{=} \PYG{n}{nn}\PYG{o}{.}\PYG{n}{Dropout}\PYG{p}{(}\PYG{l+m+mf}{0.3}\PYG{p}{)}
        \PYG{n+nb+bp}{self}\PYG{o}{.}\PYG{n}{fc3} \PYG{o}{=} \PYG{n}{nn}\PYG{o}{.}\PYG{n}{Linear}\PYG{p}{(}\PYG{l+m+mi}{512}\PYG{p}{,} \PYG{l+m+mi}{30}\PYG{p}{)}  \PYG{c+c1}{\PYGZsh{} 30 = (x,y) for 15 keypoints}

    \PYG{k}{def}\PYG{+w}{ }\PYG{n+nf}{forward}\PYG{p}{(}\PYG{n+nb+bp}{self}\PYG{p}{,} \PYG{n}{x}\PYG{p}{)}\PYG{p}{:}
        \PYG{c+c1}{\PYGZsh{} x shape: (batch\PYGZus{}size, 9216) }
        \PYG{c+c1}{\PYGZsh{} (If x is not already flattened, do x = x.view(x.size(0), \PYGZhy{}1))}
        \PYG{n}{x} \PYG{o}{=} \PYG{n}{F}\PYG{o}{.}\PYG{n}{relu}\PYG{p}{(}\PYG{n+nb+bp}{self}\PYG{o}{.}\PYG{n}{fc1}\PYG{p}{(}\PYG{n}{x}\PYG{p}{)}\PYG{p}{)}
        \PYG{n}{x} \PYG{o}{=} \PYG{n+nb+bp}{self}\PYG{o}{.}\PYG{n}{dropout1}\PYG{p}{(}\PYG{n}{x}\PYG{p}{)}
        \PYG{n}{x} \PYG{o}{=} \PYG{n}{F}\PYG{o}{.}\PYG{n}{relu}\PYG{p}{(}\PYG{n+nb+bp}{self}\PYG{o}{.}\PYG{n}{fc2}\PYG{p}{(}\PYG{n}{x}\PYG{p}{)}\PYG{p}{)}
        \PYG{n}{x} \PYG{o}{=} \PYG{n+nb+bp}{self}\PYG{o}{.}\PYG{n}{dropout2}\PYG{p}{(}\PYG{n}{x}\PYG{p}{)}
        \PYG{n}{x} \PYG{o}{=} \PYG{n+nb+bp}{self}\PYG{o}{.}\PYG{n}{fc3}\PYG{p}{(}\PYG{n}{x}\PYG{p}{)}
        \PYG{k}{return} \PYG{n}{x}
\end{sphinxVerbatim}

\end{sphinxuseclass}\end{sphinxVerbatimInput}

\end{sphinxuseclass}

\subsection{Split the Flattened Dataset for Training/Validation}
\label{\detokenize{04:split-the-flattened-dataset-for-training-validation}}
\sphinxAtStartPar
We already have \(X_{fcn}\) of shape \((N, 9216)\) and \(y\) of shape \((N, 30)\). We split them into training and validation sets:

\begin{sphinxuseclass}{cell}\begin{sphinxVerbatimInput}

\begin{sphinxuseclass}{cell_input}
\begin{sphinxVerbatim}[commandchars=\\\{\}]
\PYG{k+kn}{from}\PYG{+w}{ }\PYG{n+nn}{sklearn}\PYG{n+nn}{.}\PYG{n+nn}{model\PYGZus{}selection}\PYG{+w}{ }\PYG{k+kn}{import} \PYG{n}{train\PYGZus{}test\PYGZus{}split}

\PYG{c+c1}{\PYGZsh{} 90/10 train\PYGZhy{}validation split using X\PYGZus{}fcn}
\PYG{n}{X\PYGZus{}train\PYGZus{}fcn}\PYG{p}{,} \PYG{n}{X\PYGZus{}val\PYGZus{}fcn}\PYG{p}{,} \PYG{n}{y\PYGZus{}train\PYGZus{}fcn}\PYG{p}{,} \PYG{n}{y\PYGZus{}val\PYGZus{}fcn} \PYG{o}{=} \PYG{n}{train\PYGZus{}test\PYGZus{}split}\PYG{p}{(}
    \PYG{n}{X\PYGZus{}fcn}\PYG{p}{,} \PYG{n}{y}\PYG{p}{,} \PYG{n}{test\PYGZus{}size}\PYG{o}{=}\PYG{l+m+mf}{0.1}\PYG{p}{,} \PYG{n}{random\PYGZus{}state}\PYG{o}{=}\PYG{l+m+mi}{42}
\PYG{p}{)}

\PYG{n+nb}{print}\PYG{p}{(}\PYG{l+s+s2}{\PYGZdq{}}\PYG{l+s+s2}{Training samples (FCN):}\PYG{l+s+s2}{\PYGZdq{}}\PYG{p}{,} \PYG{n+nb}{len}\PYG{p}{(}\PYG{n}{X\PYGZus{}train\PYGZus{}fcn}\PYG{p}{)}\PYG{p}{,} 
      \PYG{l+s+s2}{\PYGZdq{}}\PYG{l+s+s2}{Validation samples (FCN):}\PYG{l+s+s2}{\PYGZdq{}}\PYG{p}{,} \PYG{n+nb}{len}\PYG{p}{(}\PYG{n}{X\PYGZus{}val\PYGZus{}fcn}\PYG{p}{)}\PYG{p}{)}
\end{sphinxVerbatim}

\end{sphinxuseclass}\end{sphinxVerbatimInput}
\begin{sphinxVerbatimOutput}

\begin{sphinxuseclass}{cell_output}
\begin{sphinxVerbatim}[commandchars=\\\{\}]
Training samples (FCN): 1926 Validation samples (FCN): 214
\end{sphinxVerbatim}

\end{sphinxuseclass}\end{sphinxVerbatimOutput}

\end{sphinxuseclass}

\subsection{Create PyTorch Datasets and DataLoaders}
\label{\detokenize{04:create-pytorch-datasets-and-dataloaders}}
\sphinxAtStartPar
We’ll create TensorDatasets from the NumPy arrays, then wrap them in DataLoaders for batching and shuffling. Notice that we \sphinxstylestrong{did not} reshape the input for the FCN version.

\begin{sphinxuseclass}{cell}\begin{sphinxVerbatimInput}

\begin{sphinxuseclass}{cell_input}
\begin{sphinxVerbatim}[commandchars=\\\{\}]
\PYG{n}{train\PYGZus{}dataset\PYGZus{}fcn} \PYG{o}{=} \PYG{n}{TensorDataset}\PYG{p}{(}
    \PYG{n}{torch}\PYG{o}{.}\PYG{n}{from\PYGZus{}numpy}\PYG{p}{(}\PYG{n}{X\PYGZus{}train\PYGZus{}fcn}\PYG{p}{)}\PYG{o}{.}\PYG{n}{to}\PYG{p}{(}\PYG{n}{device}\PYG{p}{)}\PYG{p}{,}  \PYG{c+c1}{\PYGZsh{} shape: (batch\PYGZus{}size, 9216)}
    \PYG{n}{torch}\PYG{o}{.}\PYG{n}{from\PYGZus{}numpy}\PYG{p}{(}\PYG{n}{y\PYGZus{}train\PYGZus{}fcn}\PYG{p}{)}\PYG{o}{.}\PYG{n}{to}\PYG{p}{(}\PYG{n}{device}\PYG{p}{)}   \PYG{c+c1}{\PYGZsh{} shape: (batch\PYGZus{}size, 30)}
\PYG{p}{)}
\PYG{n}{val\PYGZus{}dataset\PYGZus{}fcn} \PYG{o}{=} \PYG{n}{TensorDataset}\PYG{p}{(}
    \PYG{n}{torch}\PYG{o}{.}\PYG{n}{from\PYGZus{}numpy}\PYG{p}{(}\PYG{n}{X\PYGZus{}val\PYGZus{}fcn}\PYG{p}{)}\PYG{o}{.}\PYG{n}{to}\PYG{p}{(}\PYG{n}{device}\PYG{p}{)}\PYG{p}{,}
    \PYG{n}{torch}\PYG{o}{.}\PYG{n}{from\PYGZus{}numpy}\PYG{p}{(}\PYG{n}{y\PYGZus{}val\PYGZus{}fcn}\PYG{p}{)}\PYG{o}{.}\PYG{n}{to}\PYG{p}{(}\PYG{n}{device}\PYG{p}{)}
\PYG{p}{)}

\PYG{n}{g} \PYG{o}{=} \PYG{n}{torch}\PYG{o}{.}\PYG{n}{Generator}\PYG{p}{(}\PYG{n}{device}\PYG{o}{=}\PYG{n}{device}\PYG{p}{)}
\PYG{n}{train\PYGZus{}loader\PYGZus{}fcn} \PYG{o}{=} \PYG{n}{DataLoader}\PYG{p}{(}\PYG{n}{train\PYGZus{}dataset\PYGZus{}fcn}\PYG{p}{,} \PYG{n}{batch\PYGZus{}size}\PYG{o}{=}\PYG{l+m+mi}{64}\PYG{p}{,} \PYG{n}{shuffle}\PYG{o}{=}\PYG{k+kc}{True}\PYG{p}{,} \PYG{n}{generator}\PYG{o}{=}\PYG{n}{g}\PYG{p}{)}
\PYG{n}{val\PYGZus{}loader\PYGZus{}fcn} \PYG{o}{=} \PYG{n}{DataLoader}\PYG{p}{(}\PYG{n}{val\PYGZus{}dataset\PYGZus{}fcn}\PYG{p}{,} \PYG{n}{batch\PYGZus{}size}\PYG{o}{=}\PYG{l+m+mi}{64}\PYG{p}{,} \PYG{n}{shuffle}\PYG{o}{=}\PYG{k+kc}{False}\PYG{p}{)}
\end{sphinxVerbatim}

\end{sphinxuseclass}\end{sphinxVerbatimInput}

\end{sphinxuseclass}

\subsection{Instantiate the FCN Model, define criterion and the optimizer}
\label{\detokenize{04:instantiate-the-fcn-model-define-criterion-and-the-optimizer}}
\sphinxAtStartPar
We use an MSE loss since we’re doing keypoint regression. We use the Adam optimiser in this exercise but you can experiment with any optimizer you like (Adam, SGD, etc.):

\begin{sphinxuseclass}{cell}\begin{sphinxVerbatimInput}

\begin{sphinxuseclass}{cell_input}
\begin{sphinxVerbatim}[commandchars=\\\{\}]
\PYG{k+kn}{from}\PYG{+w}{ }\PYG{n+nn}{torchsummary}\PYG{+w}{ }\PYG{k+kn}{import} \PYG{n}{summary}
\PYG{n}{fcn\PYGZus{}model} \PYG{o}{=} \PYG{n}{KeypointFCN}\PYG{p}{(}\PYG{p}{)}\PYG{o}{.}\PYG{n}{to}\PYG{p}{(}\PYG{n}{device}\PYG{p}{)}
\PYG{n}{summary}\PYG{p}{(}\PYG{n}{fcn\PYGZus{}model}\PYG{o}{.}\PYG{n}{cpu}\PYG{p}{(}\PYG{p}{)}\PYG{p}{,} \PYG{n}{X\PYGZus{}train\PYGZus{}fcn}\PYG{p}{[}\PYG{l+m+mi}{0}\PYG{p}{]}\PYG{o}{.}\PYG{n}{shape}\PYG{p}{)}
\PYG{n}{fcn\PYGZus{}model} \PYG{o}{=} \PYG{n}{fcn\PYGZus{}model}\PYG{o}{.}\PYG{n}{to}\PYG{p}{(}\PYG{n}{device}\PYG{p}{)}
\PYG{n}{criterion\PYGZus{}fcn} \PYG{o}{=} \PYG{n}{nn}\PYG{o}{.}\PYG{n}{MSELoss}\PYG{p}{(}\PYG{p}{)}  
\PYG{n}{optimizer\PYGZus{}fcn} \PYG{o}{=} \PYG{n}{optim}\PYG{o}{.}\PYG{n}{Adam}\PYG{p}{(}\PYG{n}{fcn\PYGZus{}model}\PYG{o}{.}\PYG{n}{parameters}\PYG{p}{(}\PYG{p}{)}\PYG{p}{,} \PYG{n}{lr}\PYG{o}{=}\PYG{l+m+mf}{0.001}\PYG{p}{)}
\end{sphinxVerbatim}

\end{sphinxuseclass}\end{sphinxVerbatimInput}
\begin{sphinxVerbatimOutput}

\begin{sphinxuseclass}{cell_output}
\begin{sphinxVerbatim}[commandchars=\\\{\}]
\PYGZhy{}\PYGZhy{}\PYGZhy{}\PYGZhy{}\PYGZhy{}\PYGZhy{}\PYGZhy{}\PYGZhy{}\PYGZhy{}\PYGZhy{}\PYGZhy{}\PYGZhy{}\PYGZhy{}\PYGZhy{}\PYGZhy{}\PYGZhy{}\PYGZhy{}\PYGZhy{}\PYGZhy{}\PYGZhy{}\PYGZhy{}\PYGZhy{}\PYGZhy{}\PYGZhy{}\PYGZhy{}\PYGZhy{}\PYGZhy{}\PYGZhy{}\PYGZhy{}\PYGZhy{}\PYGZhy{}\PYGZhy{}\PYGZhy{}\PYGZhy{}\PYGZhy{}\PYGZhy{}\PYGZhy{}\PYGZhy{}\PYGZhy{}\PYGZhy{}\PYGZhy{}\PYGZhy{}\PYGZhy{}\PYGZhy{}\PYGZhy{}\PYGZhy{}\PYGZhy{}\PYGZhy{}\PYGZhy{}\PYGZhy{}\PYGZhy{}\PYGZhy{}\PYGZhy{}\PYGZhy{}\PYGZhy{}\PYGZhy{}\PYGZhy{}\PYGZhy{}\PYGZhy{}\PYGZhy{}\PYGZhy{}\PYGZhy{}\PYGZhy{}\PYGZhy{}
        Layer (type)               Output Shape         Param \PYGZsh{}
================================================================
            Linear\PYGZhy{}1                 [\PYGZhy{}1, 1024]       9,438,208
           Dropout\PYGZhy{}2                 [\PYGZhy{}1, 1024]               0
            Linear\PYGZhy{}3                  [\PYGZhy{}1, 512]         524,800
           Dropout\PYGZhy{}4                  [\PYGZhy{}1, 512]               0
            Linear\PYGZhy{}5                   [\PYGZhy{}1, 30]          15,390
================================================================
Total params: 9,978,398
Trainable params: 9,978,398
Non\PYGZhy{}trainable params: 0
\PYGZhy{}\PYGZhy{}\PYGZhy{}\PYGZhy{}\PYGZhy{}\PYGZhy{}\PYGZhy{}\PYGZhy{}\PYGZhy{}\PYGZhy{}\PYGZhy{}\PYGZhy{}\PYGZhy{}\PYGZhy{}\PYGZhy{}\PYGZhy{}\PYGZhy{}\PYGZhy{}\PYGZhy{}\PYGZhy{}\PYGZhy{}\PYGZhy{}\PYGZhy{}\PYGZhy{}\PYGZhy{}\PYGZhy{}\PYGZhy{}\PYGZhy{}\PYGZhy{}\PYGZhy{}\PYGZhy{}\PYGZhy{}\PYGZhy{}\PYGZhy{}\PYGZhy{}\PYGZhy{}\PYGZhy{}\PYGZhy{}\PYGZhy{}\PYGZhy{}\PYGZhy{}\PYGZhy{}\PYGZhy{}\PYGZhy{}\PYGZhy{}\PYGZhy{}\PYGZhy{}\PYGZhy{}\PYGZhy{}\PYGZhy{}\PYGZhy{}\PYGZhy{}\PYGZhy{}\PYGZhy{}\PYGZhy{}\PYGZhy{}\PYGZhy{}\PYGZhy{}\PYGZhy{}\PYGZhy{}\PYGZhy{}\PYGZhy{}\PYGZhy{}\PYGZhy{}
Input size (MB): 0.04
Forward/backward pass size (MB): 0.02
Params size (MB): 38.06
Estimated Total Size (MB): 38.12
\PYGZhy{}\PYGZhy{}\PYGZhy{}\PYGZhy{}\PYGZhy{}\PYGZhy{}\PYGZhy{}\PYGZhy{}\PYGZhy{}\PYGZhy{}\PYGZhy{}\PYGZhy{}\PYGZhy{}\PYGZhy{}\PYGZhy{}\PYGZhy{}\PYGZhy{}\PYGZhy{}\PYGZhy{}\PYGZhy{}\PYGZhy{}\PYGZhy{}\PYGZhy{}\PYGZhy{}\PYGZhy{}\PYGZhy{}\PYGZhy{}\PYGZhy{}\PYGZhy{}\PYGZhy{}\PYGZhy{}\PYGZhy{}\PYGZhy{}\PYGZhy{}\PYGZhy{}\PYGZhy{}\PYGZhy{}\PYGZhy{}\PYGZhy{}\PYGZhy{}\PYGZhy{}\PYGZhy{}\PYGZhy{}\PYGZhy{}\PYGZhy{}\PYGZhy{}\PYGZhy{}\PYGZhy{}\PYGZhy{}\PYGZhy{}\PYGZhy{}\PYGZhy{}\PYGZhy{}\PYGZhy{}\PYGZhy{}\PYGZhy{}\PYGZhy{}\PYGZhy{}\PYGZhy{}\PYGZhy{}\PYGZhy{}\PYGZhy{}\PYGZhy{}\PYGZhy{}
\end{sphinxVerbatim}

\end{sphinxuseclass}\end{sphinxVerbatimOutput}

\end{sphinxuseclass}

\subsection{Training with ongoing evaluation of the validation error}
\label{\detokenize{04:training-with-ongoing-evaluation-of-the-validation-error}}
\begin{sphinxuseclass}{cell}\begin{sphinxVerbatimInput}

\begin{sphinxuseclass}{cell_input}
\begin{sphinxVerbatim}[commandchars=\\\{\}]
\PYG{k+kn}{from}\PYG{+w}{ }\PYG{n+nn}{tqdm}\PYG{n+nn}{.}\PYG{n+nn}{auto}\PYG{+w}{ }\PYG{k+kn}{import} \PYG{n}{tqdm}
\PYG{n}{epochs} \PYG{o}{=} \PYG{l+m+mi}{20}
\PYG{k}{for} \PYG{n}{epoch} \PYG{o+ow}{in} \PYG{p}{(}\PYG{n}{pbar}\PYG{o}{:=}\PYG{n}{tqdm}\PYG{p}{(}\PYG{n+nb}{range}\PYG{p}{(}\PYG{l+m+mi}{1}\PYG{p}{,} \PYG{n}{epochs}\PYG{o}{+}\PYG{l+m+mi}{1}\PYG{p}{)}\PYG{p}{)}\PYG{p}{)}\PYG{p}{:}
    \PYG{c+c1}{\PYGZsh{} \PYGZhy{}\PYGZhy{}\PYGZhy{} Training \PYGZhy{}\PYGZhy{}\PYGZhy{}}
    \PYG{n}{fcn\PYGZus{}model}\PYG{o}{.}\PYG{n}{train}\PYG{p}{(}\PYG{p}{)}
    \PYG{n}{train\PYGZus{}loss} \PYG{o}{=} \PYG{l+m+mf}{0.0}
    \PYG{k}{for} \PYG{n}{batch\PYGZus{}images}\PYG{p}{,} \PYG{n}{batch\PYGZus{}keypoints} \PYG{o+ow}{in} \PYG{n}{train\PYGZus{}loader\PYGZus{}fcn}\PYG{p}{:}
        \PYG{n}{outputs} \PYG{o}{=} \PYG{n}{fcn\PYGZus{}model}\PYG{p}{(}\PYG{n}{batch\PYGZus{}images}\PYG{p}{)}
        \PYG{n}{loss} \PYG{o}{=} \PYG{n}{criterion\PYGZus{}fcn}\PYG{p}{(}\PYG{n}{outputs}\PYG{p}{,} \PYG{n}{batch\PYGZus{}keypoints}\PYG{p}{)}

        \PYG{n}{optimizer\PYGZus{}fcn}\PYG{o}{.}\PYG{n}{zero\PYGZus{}grad}\PYG{p}{(}\PYG{p}{)}
        \PYG{n}{loss}\PYG{o}{.}\PYG{n}{backward}\PYG{p}{(}\PYG{p}{)}
        \PYG{n}{optimizer\PYGZus{}fcn}\PYG{o}{.}\PYG{n}{step}\PYG{p}{(}\PYG{p}{)}

        \PYG{n}{train\PYGZus{}loss} \PYG{o}{+}\PYG{o}{=} \PYG{n}{loss}\PYG{o}{.}\PYG{n}{item}\PYG{p}{(}\PYG{p}{)} \PYG{o}{*} \PYG{n}{batch\PYGZus{}images}\PYG{o}{.}\PYG{n}{size}\PYG{p}{(}\PYG{l+m+mi}{0}\PYG{p}{)}
    \PYG{n}{train\PYGZus{}loss} \PYG{o}{/}\PYG{o}{=} \PYG{n+nb}{len}\PYG{p}{(}\PYG{n}{train\PYGZus{}loader\PYGZus{}fcn}\PYG{o}{.}\PYG{n}{dataset}\PYG{p}{)}

    \PYG{c+c1}{\PYGZsh{} \PYGZhy{}\PYGZhy{}\PYGZhy{} Validation \PYGZhy{}\PYGZhy{}\PYGZhy{}}
    \PYG{n}{fcn\PYGZus{}model}\PYG{o}{.}\PYG{n}{eval}\PYG{p}{(}\PYG{p}{)}
    \PYG{n}{val\PYGZus{}loss} \PYG{o}{=} \PYG{l+m+mf}{0.0}
    \PYG{k}{with} \PYG{n}{torch}\PYG{o}{.}\PYG{n}{no\PYGZus{}grad}\PYG{p}{(}\PYG{p}{)}\PYG{p}{:}
        \PYG{k}{for} \PYG{n}{batch\PYGZus{}images}\PYG{p}{,} \PYG{n}{batch\PYGZus{}keypoints} \PYG{o+ow}{in} \PYG{n}{val\PYGZus{}loader\PYGZus{}fcn}\PYG{p}{:}
            \PYG{n}{preds} \PYG{o}{=} \PYG{n}{fcn\PYGZus{}model}\PYG{p}{(}\PYG{n}{batch\PYGZus{}images}\PYG{p}{)}
            \PYG{n}{loss} \PYG{o}{=} \PYG{n}{criterion\PYGZus{}fcn}\PYG{p}{(}\PYG{n}{preds}\PYG{p}{,} \PYG{n}{batch\PYGZus{}keypoints}\PYG{p}{)}
            \PYG{n}{val\PYGZus{}loss} \PYG{o}{+}\PYG{o}{=} \PYG{n}{loss}\PYG{o}{.}\PYG{n}{item}\PYG{p}{(}\PYG{p}{)} \PYG{o}{*} \PYG{n}{batch\PYGZus{}images}\PYG{o}{.}\PYG{n}{size}\PYG{p}{(}\PYG{l+m+mi}{0}\PYG{p}{)}
    \PYG{n}{val\PYGZus{}loss} \PYG{o}{/}\PYG{o}{=} \PYG{n+nb}{len}\PYG{p}{(}\PYG{n}{val\PYGZus{}loader\PYGZus{}fcn}\PYG{o}{.}\PYG{n}{dataset}\PYG{p}{)}

    \PYG{n}{pbar}\PYG{o}{.}\PYG{n}{set\PYGZus{}description}\PYG{p}{(}\PYG{l+s+sa}{f}\PYG{l+s+s2}{\PYGZdq{}}\PYG{l+s+s2}{Epoch }\PYG{l+s+si}{\PYGZob{}}\PYG{n}{epoch}\PYG{l+s+si}{\PYGZcb{}}\PYG{l+s+s2}{/}\PYG{l+s+si}{\PYGZob{}}\PYG{n}{epochs}\PYG{l+s+si}{\PYGZcb{}}\PYG{l+s+s2}{ \PYGZhy{} Train Loss: }\PYG{l+s+si}{\PYGZob{}}\PYG{n}{train\PYGZus{}loss}\PYG{l+s+si}{:}\PYG{l+s+s2}{.4f}\PYG{l+s+si}{\PYGZcb{}}\PYG{l+s+s2}{ \PYGZhy{} Val Loss: }\PYG{l+s+si}{\PYGZob{}}\PYG{n}{val\PYGZus{}loss}\PYG{l+s+si}{:}\PYG{l+s+s2}{.4f}\PYG{l+s+si}{\PYGZcb{}}\PYG{l+s+s2}{\PYGZdq{}}\PYG{p}{)}
\end{sphinxVerbatim}

\end{sphinxuseclass}\end{sphinxVerbatimInput}
\begin{sphinxVerbatimOutput}

\begin{sphinxuseclass}{cell_output}
\begin{sphinxVerbatim}[commandchars=\\\{\}]
  0\PYGZpc{}|          | 0/20 [00:00\PYGZlt{}?, ?it/s]
\end{sphinxVerbatim}

\end{sphinxuseclass}\end{sphinxVerbatimOutput}

\end{sphinxuseclass}

\subsection{Evaluate on validation set}
\label{\detokenize{04:evaluate-on-validation-set}}
\sphinxAtStartPar
After the training completes, we can compute the validation set predictions and compute appropriate metrics to estimate the performance of the fully connected neural network model:

\begin{sphinxuseclass}{cell}\begin{sphinxVerbatimInput}

\begin{sphinxuseclass}{cell_input}
\begin{sphinxVerbatim}[commandchars=\\\{\}]
\PYG{n}{fcn\PYGZus{}model}\PYG{o}{.}\PYG{n}{eval}\PYG{p}{(}\PYG{p}{)}
\PYG{k}{with} \PYG{n}{torch}\PYG{o}{.}\PYG{n}{no\PYGZus{}grad}\PYG{p}{(}\PYG{p}{)}\PYG{p}{:}
    \PYG{n}{val\PYGZus{}preds\PYGZus{}fcn} \PYG{o}{=} \PYG{n}{fcn\PYGZus{}model}\PYG{p}{(}\PYG{n}{torch}\PYG{o}{.}\PYG{n}{Tensor}\PYG{p}{(}\PYG{n}{X\PYGZus{}val\PYGZus{}fcn}\PYG{p}{)}\PYG{o}{.}\PYG{n}{to}\PYG{p}{(}\PYG{n}{device}\PYG{p}{)}\PYG{p}{)}\PYG{o}{.}\PYG{n}{cpu}\PYG{p}{(}\PYG{p}{)}\PYG{o}{.}\PYG{n}{numpy}\PYG{p}{(}\PYG{p}{)}
    
\PYG{n}{val\PYGZus{}truth\PYGZus{}fcn} \PYG{o}{=} \PYG{n}{y\PYGZus{}val\PYGZus{}fcn}
\PYG{n}{mse\PYGZus{}val\PYGZus{}fcn} \PYG{o}{=} \PYG{n}{np}\PYG{o}{.}\PYG{n}{mean}\PYG{p}{(}\PYG{p}{(}\PYG{n}{val\PYGZus{}preds\PYGZus{}fcn} \PYG{o}{\PYGZhy{}} \PYG{n}{val\PYGZus{}truth\PYGZus{}fcn}\PYG{p}{)}\PYG{o}{*}\PYG{o}{*}\PYG{l+m+mi}{2}\PYG{p}{)}
\PYG{n}{rmse\PYGZus{}val\PYGZus{}fcn} \PYG{o}{=} \PYG{n}{np}\PYG{o}{.}\PYG{n}{sqrt}\PYG{p}{(}\PYG{n}{mse\PYGZus{}val\PYGZus{}fcn}\PYG{p}{)}

\PYG{n+nb}{print}\PYG{p}{(}\PYG{l+s+sa}{f}\PYG{l+s+s2}{\PYGZdq{}}\PYG{l+s+s2}{Validation RMSE (FCN): }\PYG{l+s+si}{\PYGZob{}}\PYG{n}{rmse\PYGZus{}val\PYGZus{}fcn}\PYG{l+s+si}{:}\PYG{l+s+s2}{.3f}\PYG{l+s+si}{\PYGZcb{}}\PYG{l+s+s2}{ pixels}\PYG{l+s+s2}{\PYGZdq{}}\PYG{p}{)}
\end{sphinxVerbatim}

\end{sphinxuseclass}\end{sphinxVerbatimInput}
\begin{sphinxVerbatimOutput}

\begin{sphinxuseclass}{cell_output}
\begin{sphinxVerbatim}[commandchars=\\\{\}]
Validation RMSE (FCN): 10.461 pixels
\end{sphinxVerbatim}

\end{sphinxuseclass}\end{sphinxVerbatimOutput}

\end{sphinxuseclass}

\section{Convolutional neural network}
\label{\detokenize{04:convolutional-neural-network}}
\sphinxAtStartPar
We will create a Convolutional Neural Network (CNN) to \sphinxstylestrong{regress} the 30 keypoint coordinates from the input image. CNNs are highly effective in image analysis as they can automatically learn spatial hierarchies of features from images – from low\sphinxhyphen{}level edges to high\sphinxhyphen{}level shapes – which makes them well\sphinxhyphen{}suited for detecting patterns like facial features.

\sphinxAtStartPar
The architecture of our CNN will be as follows:
\begin{enumerate}
\sphinxsetlistlabels{\arabic}{enumi}{enumii}{}{.}%
\item {} 
\sphinxAtStartPar
\sphinxstylestrong{Conv Layer 1}: 1 input channel (grayscale) \sphinxhyphen{}> 32 filters, kernel size 3x3, activation ReLU, followed by Max Pooling (2x2).

\item {} 
\sphinxAtStartPar
\sphinxstylestrong{Conv Layer 2}: 32 \sphinxhyphen{}> 64 filters, 3x3, ReLU, followed by Max Pooling (2x2).

\item {} 
\sphinxAtStartPar
\sphinxstylestrong{Conv Layer 3}: 64 \sphinxhyphen{}> 128 filters, 3x3, ReLU, followed by Max Pooling (2x2).

\item {} 
\sphinxAtStartPar
\sphinxstylestrong{Flatten}: Flatten the feature maps into a vector.

\item {} 
\sphinxAtStartPar
\sphinxstylestrong{Fully Connected 1}: 256 neurons, ReLU, with Dropout (to reduce overfitting).

\item {} 
\sphinxAtStartPar
\sphinxstylestrong{Fully Connected 2}: 128 neurons, ReLU, with Dropout.

\item {} 
\sphinxAtStartPar
\sphinxstylestrong{Output Layer}: 30 outputs (the \( x,y \) coordinates for 15 keypoints), linear activation (no final nonlinearity, since this is a regression output).

\end{enumerate}

\sphinxAtStartPar
Dropout layers will help regularize the network by randomly dropping a fraction of neurons during training, which can improve generalization. We expect the CNN to learn to detect facial features through the convolutional layers (which capture local patterns like edges, corners of eyes/mouth, etc.), and the fully connected layers will combine these to predict the precise coordinates.

\begin{sphinxuseclass}{cell}\begin{sphinxVerbatimInput}

\begin{sphinxuseclass}{cell_input}
\begin{sphinxVerbatim}[commandchars=\\\{\}]
\PYG{k}{class}\PYG{+w}{ }\PYG{n+nc}{KeypointCNN}\PYG{p}{(}\PYG{n}{nn}\PYG{o}{.}\PYG{n}{Module}\PYG{p}{)}\PYG{p}{:}
    \PYG{k}{def}\PYG{+w}{ }\PYG{n+nf+fm}{\PYGZus{}\PYGZus{}init\PYGZus{}\PYGZus{}}\PYG{p}{(}\PYG{n+nb+bp}{self}\PYG{p}{)}\PYG{p}{:}
        \PYG{n+nb}{super}\PYG{p}{(}\PYG{n}{KeypointCNN}\PYG{p}{,} \PYG{n+nb+bp}{self}\PYG{p}{)}\PYG{o}{.}\PYG{n+nf+fm}{\PYGZus{}\PYGZus{}init\PYGZus{}\PYGZus{}}\PYG{p}{(}\PYG{p}{)}
        \PYG{c+c1}{\PYGZsh{} Convolutional layers}
        \PYG{n+nb+bp}{self}\PYG{o}{.}\PYG{n}{conv1} \PYG{o}{=} \PYG{n}{nn}\PYG{o}{.}\PYG{n}{Conv2d}\PYG{p}{(}\PYG{l+m+mi}{1}\PYG{p}{,} \PYG{l+m+mi}{32}\PYG{p}{,} \PYG{n}{kernel\PYGZus{}size}\PYG{o}{=}\PYG{l+m+mi}{3}\PYG{p}{)}   \PYG{c+c1}{\PYGZsh{} 96x96 \PYGZhy{}\PYGZgt{} 94x94 (then 47x47 after pool)}
        \PYG{n+nb+bp}{self}\PYG{o}{.}\PYG{n}{conv2} \PYG{o}{=} \PYG{n}{nn}\PYG{o}{.}\PYG{n}{Conv2d}\PYG{p}{(}\PYG{l+m+mi}{32}\PYG{p}{,} \PYG{l+m+mi}{64}\PYG{p}{,} \PYG{n}{kernel\PYGZus{}size}\PYG{o}{=}\PYG{l+m+mi}{3}\PYG{p}{)}  \PYG{c+c1}{\PYGZsh{} 47x47 \PYGZhy{}\PYGZgt{} 45x45 (then 22x22 after pool)}
        \PYG{n+nb+bp}{self}\PYG{o}{.}\PYG{n}{conv3} \PYG{o}{=} \PYG{n}{nn}\PYG{o}{.}\PYG{n}{Conv2d}\PYG{p}{(}\PYG{l+m+mi}{64}\PYG{p}{,} \PYG{l+m+mi}{128}\PYG{p}{,} \PYG{n}{kernel\PYGZus{}size}\PYG{o}{=}\PYG{l+m+mi}{3}\PYG{p}{)} \PYG{c+c1}{\PYGZsh{} 22x22 \PYGZhy{}\PYGZgt{} 20x20 (then 10x10 after pool)}
        \PYG{n+nb+bp}{self}\PYG{o}{.}\PYG{n}{pool} \PYG{o}{=} \PYG{n}{nn}\PYG{o}{.}\PYG{n}{MaxPool2d}\PYG{p}{(}\PYG{l+m+mi}{2}\PYG{p}{,} \PYG{l+m+mi}{2}\PYG{p}{)}

        \PYG{c+c1}{\PYGZsh{} Dropout layers}
        \PYG{n+nb+bp}{self}\PYG{o}{.}\PYG{n}{dropout1} \PYG{o}{=} \PYG{n}{nn}\PYG{o}{.}\PYG{n}{Dropout}\PYG{p}{(}\PYG{l+m+mf}{0.1}\PYG{p}{)}
        \PYG{n+nb+bp}{self}\PYG{o}{.}\PYG{n}{dropout2} \PYG{o}{=} \PYG{n}{nn}\PYG{o}{.}\PYG{n}{Dropout}\PYG{p}{(}\PYG{l+m+mf}{0.2}\PYG{p}{)}

        \PYG{c+c1}{\PYGZsh{} Fully connected layers}
        \PYG{n+nb+bp}{self}\PYG{o}{.}\PYG{n}{fc1} \PYG{o}{=} \PYG{n}{nn}\PYG{o}{.}\PYG{n}{Linear}\PYG{p}{(}\PYG{l+m+mi}{128} \PYG{o}{*} \PYG{l+m+mi}{10} \PYG{o}{*} \PYG{l+m+mi}{10}\PYG{p}{,} \PYG{l+m+mi}{256}\PYG{p}{)}  \PYG{c+c1}{\PYGZsh{} 128 feature maps * 10 * 10}
        \PYG{n+nb+bp}{self}\PYG{o}{.}\PYG{n}{fc2} \PYG{o}{=} \PYG{n}{nn}\PYG{o}{.}\PYG{n}{Linear}\PYG{p}{(}\PYG{l+m+mi}{256}\PYG{p}{,} \PYG{l+m+mi}{128}\PYG{p}{)}
        \PYG{n+nb+bp}{self}\PYG{o}{.}\PYG{n}{fc3} \PYG{o}{=} \PYG{n}{nn}\PYG{o}{.}\PYG{n}{Linear}\PYG{p}{(}\PYG{l+m+mi}{128}\PYG{p}{,} \PYG{l+m+mi}{30}\PYG{p}{)}  \PYG{c+c1}{\PYGZsh{} 30 outputs (x,y for 15 keypoints)}
        
    \PYG{k}{def}\PYG{+w}{ }\PYG{n+nf}{forward}\PYG{p}{(}\PYG{n+nb+bp}{self}\PYG{p}{,} \PYG{n}{x}\PYG{p}{)}\PYG{p}{:}
        \PYG{c+c1}{\PYGZsh{} Three conv layers with ReLU and pooling}
        \PYG{n}{x} \PYG{o}{=} \PYG{n+nb+bp}{self}\PYG{o}{.}\PYG{n}{pool}\PYG{p}{(}\PYG{n}{F}\PYG{o}{.}\PYG{n}{relu}\PYG{p}{(}\PYG{n+nb+bp}{self}\PYG{o}{.}\PYG{n}{conv1}\PYG{p}{(}\PYG{n}{x}\PYG{p}{)}\PYG{p}{)}\PYG{p}{)}
        \PYG{n}{x} \PYG{o}{=} \PYG{n+nb+bp}{self}\PYG{o}{.}\PYG{n}{pool}\PYG{p}{(}\PYG{n}{F}\PYG{o}{.}\PYG{n}{relu}\PYG{p}{(}\PYG{n+nb+bp}{self}\PYG{o}{.}\PYG{n}{conv2}\PYG{p}{(}\PYG{n}{x}\PYG{p}{)}\PYG{p}{)}\PYG{p}{)}
        \PYG{n}{x} \PYG{o}{=} \PYG{n+nb+bp}{self}\PYG{o}{.}\PYG{n}{pool}\PYG{p}{(}\PYG{n}{F}\PYG{o}{.}\PYG{n}{relu}\PYG{p}{(}\PYG{n+nb+bp}{self}\PYG{o}{.}\PYG{n}{conv3}\PYG{p}{(}\PYG{n}{x}\PYG{p}{)}\PYG{p}{)}\PYG{p}{)}

        \PYG{c+c1}{\PYGZsh{} Flatten}
        \PYG{n}{x} \PYG{o}{=} \PYG{n}{x}\PYG{o}{.}\PYG{n}{view}\PYG{p}{(}\PYG{n}{x}\PYG{o}{.}\PYG{n}{size}\PYG{p}{(}\PYG{l+m+mi}{0}\PYG{p}{)}\PYG{p}{,} \PYG{o}{\PYGZhy{}}\PYG{l+m+mi}{1}\PYG{p}{)}

        \PYG{c+c1}{\PYGZsh{} Fully connected layers with dropout and ReLU}
        \PYG{n}{x} \PYG{o}{=} \PYG{n}{F}\PYG{o}{.}\PYG{n}{relu}\PYG{p}{(}\PYG{n+nb+bp}{self}\PYG{o}{.}\PYG{n}{fc1}\PYG{p}{(}\PYG{n}{x}\PYG{p}{)}\PYG{p}{)}
        \PYG{n}{x} \PYG{o}{=} \PYG{n+nb+bp}{self}\PYG{o}{.}\PYG{n}{dropout1}\PYG{p}{(}\PYG{n}{x}\PYG{p}{)}
        \PYG{n}{x} \PYG{o}{=} \PYG{n}{F}\PYG{o}{.}\PYG{n}{relu}\PYG{p}{(}\PYG{n+nb+bp}{self}\PYG{o}{.}\PYG{n}{fc2}\PYG{p}{(}\PYG{n}{x}\PYG{p}{)}\PYG{p}{)}
        \PYG{n}{x} \PYG{o}{=} \PYG{n+nb+bp}{self}\PYG{o}{.}\PYG{n}{dropout2}\PYG{p}{(}\PYG{n}{x}\PYG{p}{)}
        \PYG{n}{x} \PYG{o}{=} \PYG{n+nb+bp}{self}\PYG{o}{.}\PYG{n}{fc3}\PYG{p}{(}\PYG{n}{x}\PYG{p}{)}
        \PYG{k}{return} \PYG{n}{x}
\end{sphinxVerbatim}

\end{sphinxuseclass}\end{sphinxVerbatimInput}

\end{sphinxuseclass}

\section{Training the model}
\label{\detokenize{04:training-the-model}}
\sphinxAtStartPar
Now, let’s train the model on our training data and monitor its performance on the validation set.
\begin{itemize}
\item {} 
\sphinxAtStartPar
\sphinxstylestrong{Data Loaders}: We’ll use PyTorch \sphinxcode{\sphinxupquote{DataLoader}} to batch and shuffle the data for training and to batch the validation data.

\item {} 
\sphinxAtStartPar
\sphinxstylestrong{Loss Function}: This is a regression problem (predicting continuous coordinate values), so a common choice is \sphinxstylestrong{Mean Squared Error} (\sphinxcode{\sphinxupquote{nn.MSELoss}}) in PyTorch.

\item {} 
\sphinxAtStartPar
\sphinxstylestrong{Optimizer}: We’ll use \sphinxstylestrong{Adam} for efficient stochastic gradient descent.

\item {} 
\sphinxAtStartPar
\sphinxstylestrong{Training Loop}: We’ll train for a number of epochs (iterations over the whole training set). In each epoch:
\begin{enumerate}
\sphinxsetlistlabels{\arabic}{enumi}{enumii}{}{.}%
\item {} 
\sphinxAtStartPar
Set the model to training mode.

\item {} 
\sphinxAtStartPar
Loop over mini\sphinxhyphen{}batches from the training loader.

\item {} 
\sphinxAtStartPar
For each batch, do a forward pass to get predictions, compute the loss, do a backward pass to compute gradients, and update weights.

\item {} 
\sphinxAtStartPar
Evaluate the model on the validation set after each epoch to track performance.

\end{enumerate}

\end{itemize}

\begin{sphinxuseclass}{cell}\begin{sphinxVerbatimInput}

\begin{sphinxuseclass}{cell_input}
\begin{sphinxVerbatim}[commandchars=\\\{\}]
\PYG{k+kn}{from}\PYG{+w}{ }\PYG{n+nn}{sklearn}\PYG{n+nn}{.}\PYG{n+nn}{model\PYGZus{}selection}\PYG{+w}{ }\PYG{k+kn}{import} \PYG{n}{train\PYGZus{}test\PYGZus{}split}

\PYG{c+c1}{\PYGZsh{} We\PYGZsq{}ll do a 90/10 train\PYGZhy{}validation split}
\PYG{n}{X\PYGZus{}train\PYGZus{}cnn}\PYG{p}{,} \PYG{n}{X\PYGZus{}val\PYGZus{}cnn}\PYG{p}{,} \PYG{n}{y\PYGZus{}train\PYGZus{}cnn}\PYG{p}{,} \PYG{n}{y\PYGZus{}val\PYGZus{}cnn} \PYG{o}{=} \PYG{n}{train\PYGZus{}test\PYGZus{}split}\PYG{p}{(}\PYG{n}{X\PYGZus{}cnn}\PYG{p}{,} \PYG{n}{y}\PYG{p}{,} \PYG{n}{test\PYGZus{}size}\PYG{o}{=}\PYG{l+m+mf}{0.1}\PYG{p}{,} \PYG{n}{random\PYGZus{}state}\PYG{o}{=}\PYG{l+m+mi}{42}\PYG{p}{)}
\PYG{n+nb}{print}\PYG{p}{(}\PYG{l+s+s2}{\PYGZdq{}}\PYG{l+s+s2}{Training samples:}\PYG{l+s+s2}{\PYGZdq{}}\PYG{p}{,} \PYG{n+nb}{len}\PYG{p}{(}\PYG{n}{X\PYGZus{}train\PYGZus{}cnn}\PYG{p}{)}\PYG{p}{,} \PYG{l+s+s2}{\PYGZdq{}}\PYG{l+s+s2}{Validation samples:}\PYG{l+s+s2}{\PYGZdq{}}\PYG{p}{,} \PYG{n+nb}{len}\PYG{p}{(}\PYG{n}{X\PYGZus{}val\PYGZus{}cnn}\PYG{p}{)}\PYG{p}{)}
\end{sphinxVerbatim}

\end{sphinxuseclass}\end{sphinxVerbatimInput}
\begin{sphinxVerbatimOutput}

\begin{sphinxuseclass}{cell_output}
\begin{sphinxVerbatim}[commandchars=\\\{\}]
Training samples: 1926 Validation samples: 214
\end{sphinxVerbatim}

\end{sphinxuseclass}\end{sphinxVerbatimOutput}

\end{sphinxuseclass}
\begin{sphinxuseclass}{cell}\begin{sphinxVerbatimInput}

\begin{sphinxuseclass}{cell_input}
\begin{sphinxVerbatim}[commandchars=\\\{\}]
\PYG{c+c1}{\PYGZsh{} Create TensorDatasets for train and validation}
\PYG{n}{train\PYGZus{}dataset} \PYG{o}{=} \PYG{n}{TensorDataset}\PYG{p}{(}\PYG{n}{torch}\PYG{o}{.}\PYG{n}{from\PYGZus{}numpy}\PYG{p}{(}\PYG{n}{X\PYGZus{}train\PYGZus{}cnn}\PYG{p}{)}\PYG{o}{.}\PYG{n}{to}\PYG{p}{(}\PYG{n}{device}\PYG{p}{)}\PYG{p}{,} \PYG{n}{torch}\PYG{o}{.}\PYG{n}{from\PYGZus{}numpy}\PYG{p}{(}\PYG{n}{y\PYGZus{}train\PYGZus{}cnn}\PYG{p}{)}\PYG{o}{.}\PYG{n}{to}\PYG{p}{(}\PYG{n}{device}\PYG{p}{)}\PYG{p}{)}
\PYG{n}{val\PYGZus{}dataset} \PYG{o}{=} \PYG{n}{TensorDataset}\PYG{p}{(}\PYG{n}{torch}\PYG{o}{.}\PYG{n}{from\PYGZus{}numpy}\PYG{p}{(}\PYG{n}{X\PYGZus{}val\PYGZus{}cnn}\PYG{p}{)}\PYG{o}{.}\PYG{n}{to}\PYG{p}{(}\PYG{n}{device}\PYG{p}{)}\PYG{p}{,} \PYG{n}{torch}\PYG{o}{.}\PYG{n}{from\PYGZus{}numpy}\PYG{p}{(}\PYG{n}{y\PYGZus{}val\PYGZus{}cnn}\PYG{p}{)}\PYG{o}{.}\PYG{n}{to}\PYG{p}{(}\PYG{n}{device}\PYG{p}{)}\PYG{p}{)}

\PYG{c+c1}{\PYGZsh{} DataLoader for batching}
\PYG{n}{g} \PYG{o}{=} \PYG{n}{torch}\PYG{o}{.}\PYG{n}{Generator}\PYG{p}{(}\PYG{n}{device}\PYG{o}{=}\PYG{n}{device}\PYG{p}{)}
\PYG{n}{train\PYGZus{}loader} \PYG{o}{=} \PYG{n}{DataLoader}\PYG{p}{(}\PYG{n}{train\PYGZus{}dataset}\PYG{p}{,} \PYG{n}{batch\PYGZus{}size}\PYG{o}{=}\PYG{l+m+mi}{64}\PYG{p}{,} \PYG{n}{shuffle}\PYG{o}{=}\PYG{k+kc}{True}\PYG{p}{,} \PYG{n}{generator}\PYG{o}{=}\PYG{n}{g}\PYG{p}{)}
\PYG{n}{val\PYGZus{}loader} \PYG{o}{=} \PYG{n}{DataLoader}\PYG{p}{(}\PYG{n}{val\PYGZus{}dataset}\PYG{p}{,} \PYG{n}{batch\PYGZus{}size}\PYG{o}{=}\PYG{l+m+mi}{64}\PYG{p}{,} \PYG{n}{shuffle}\PYG{o}{=}\PYG{k+kc}{False}\PYG{p}{)}
\end{sphinxVerbatim}

\end{sphinxuseclass}\end{sphinxVerbatimInput}

\end{sphinxuseclass}
\begin{sphinxuseclass}{cell}\begin{sphinxVerbatimInput}

\begin{sphinxuseclass}{cell_input}
\begin{sphinxVerbatim}[commandchars=\\\{\}]
\PYG{n}{cnn\PYGZus{}model} \PYG{o}{=} \PYG{n}{KeypointCNN}\PYG{p}{(}\PYG{p}{)}
\PYG{n}{summary}\PYG{p}{(}\PYG{n}{cnn\PYGZus{}model}\PYG{o}{.}\PYG{n}{cpu}\PYG{p}{(}\PYG{p}{)}\PYG{p}{,} \PYG{n}{X\PYGZus{}train\PYGZus{}cnn}\PYG{p}{[}\PYG{l+m+mi}{0}\PYG{p}{]}\PYG{o}{.}\PYG{n}{shape}\PYG{p}{)}
\PYG{n}{cnn\PYGZus{}model} \PYG{o}{=} \PYG{n}{cnn\PYGZus{}model}\PYG{o}{.}\PYG{n}{to}\PYG{p}{(}\PYG{n}{device}\PYG{p}{)}
\end{sphinxVerbatim}

\end{sphinxuseclass}\end{sphinxVerbatimInput}
\begin{sphinxVerbatimOutput}

\begin{sphinxuseclass}{cell_output}
\begin{sphinxVerbatim}[commandchars=\\\{\}]
\PYGZhy{}\PYGZhy{}\PYGZhy{}\PYGZhy{}\PYGZhy{}\PYGZhy{}\PYGZhy{}\PYGZhy{}\PYGZhy{}\PYGZhy{}\PYGZhy{}\PYGZhy{}\PYGZhy{}\PYGZhy{}\PYGZhy{}\PYGZhy{}\PYGZhy{}\PYGZhy{}\PYGZhy{}\PYGZhy{}\PYGZhy{}\PYGZhy{}\PYGZhy{}\PYGZhy{}\PYGZhy{}\PYGZhy{}\PYGZhy{}\PYGZhy{}\PYGZhy{}\PYGZhy{}\PYGZhy{}\PYGZhy{}\PYGZhy{}\PYGZhy{}\PYGZhy{}\PYGZhy{}\PYGZhy{}\PYGZhy{}\PYGZhy{}\PYGZhy{}\PYGZhy{}\PYGZhy{}\PYGZhy{}\PYGZhy{}\PYGZhy{}\PYGZhy{}\PYGZhy{}\PYGZhy{}\PYGZhy{}\PYGZhy{}\PYGZhy{}\PYGZhy{}\PYGZhy{}\PYGZhy{}\PYGZhy{}\PYGZhy{}\PYGZhy{}\PYGZhy{}\PYGZhy{}\PYGZhy{}\PYGZhy{}\PYGZhy{}\PYGZhy{}\PYGZhy{}
        Layer (type)               Output Shape         Param \PYGZsh{}
================================================================
            Conv2d\PYGZhy{}1           [\PYGZhy{}1, 32, 94, 94]             320
         MaxPool2d\PYGZhy{}2           [\PYGZhy{}1, 32, 47, 47]               0
            Conv2d\PYGZhy{}3           [\PYGZhy{}1, 64, 45, 45]          18,496
         MaxPool2d\PYGZhy{}4           [\PYGZhy{}1, 64, 22, 22]               0
            Conv2d\PYGZhy{}5          [\PYGZhy{}1, 128, 20, 20]          73,856
         MaxPool2d\PYGZhy{}6          [\PYGZhy{}1, 128, 10, 10]               0
            Linear\PYGZhy{}7                  [\PYGZhy{}1, 256]       3,277,056
           Dropout\PYGZhy{}8                  [\PYGZhy{}1, 256]               0
            Linear\PYGZhy{}9                  [\PYGZhy{}1, 128]          32,896
          Dropout\PYGZhy{}10                  [\PYGZhy{}1, 128]               0
           Linear\PYGZhy{}11                   [\PYGZhy{}1, 30]           3,870
================================================================
Total params: 3,406,494
Trainable params: 3,406,494
Non\PYGZhy{}trainable params: 0
\PYGZhy{}\PYGZhy{}\PYGZhy{}\PYGZhy{}\PYGZhy{}\PYGZhy{}\PYGZhy{}\PYGZhy{}\PYGZhy{}\PYGZhy{}\PYGZhy{}\PYGZhy{}\PYGZhy{}\PYGZhy{}\PYGZhy{}\PYGZhy{}\PYGZhy{}\PYGZhy{}\PYGZhy{}\PYGZhy{}\PYGZhy{}\PYGZhy{}\PYGZhy{}\PYGZhy{}\PYGZhy{}\PYGZhy{}\PYGZhy{}\PYGZhy{}\PYGZhy{}\PYGZhy{}\PYGZhy{}\PYGZhy{}\PYGZhy{}\PYGZhy{}\PYGZhy{}\PYGZhy{}\PYGZhy{}\PYGZhy{}\PYGZhy{}\PYGZhy{}\PYGZhy{}\PYGZhy{}\PYGZhy{}\PYGZhy{}\PYGZhy{}\PYGZhy{}\PYGZhy{}\PYGZhy{}\PYGZhy{}\PYGZhy{}\PYGZhy{}\PYGZhy{}\PYGZhy{}\PYGZhy{}\PYGZhy{}\PYGZhy{}\PYGZhy{}\PYGZhy{}\PYGZhy{}\PYGZhy{}\PYGZhy{}\PYGZhy{}\PYGZhy{}\PYGZhy{}
Input size (MB): 0.04
Forward/backward pass size (MB): 4.42
Params size (MB): 12.99
Estimated Total Size (MB): 17.45
\PYGZhy{}\PYGZhy{}\PYGZhy{}\PYGZhy{}\PYGZhy{}\PYGZhy{}\PYGZhy{}\PYGZhy{}\PYGZhy{}\PYGZhy{}\PYGZhy{}\PYGZhy{}\PYGZhy{}\PYGZhy{}\PYGZhy{}\PYGZhy{}\PYGZhy{}\PYGZhy{}\PYGZhy{}\PYGZhy{}\PYGZhy{}\PYGZhy{}\PYGZhy{}\PYGZhy{}\PYGZhy{}\PYGZhy{}\PYGZhy{}\PYGZhy{}\PYGZhy{}\PYGZhy{}\PYGZhy{}\PYGZhy{}\PYGZhy{}\PYGZhy{}\PYGZhy{}\PYGZhy{}\PYGZhy{}\PYGZhy{}\PYGZhy{}\PYGZhy{}\PYGZhy{}\PYGZhy{}\PYGZhy{}\PYGZhy{}\PYGZhy{}\PYGZhy{}\PYGZhy{}\PYGZhy{}\PYGZhy{}\PYGZhy{}\PYGZhy{}\PYGZhy{}\PYGZhy{}\PYGZhy{}\PYGZhy{}\PYGZhy{}\PYGZhy{}\PYGZhy{}\PYGZhy{}\PYGZhy{}\PYGZhy{}\PYGZhy{}\PYGZhy{}\PYGZhy{}
\end{sphinxVerbatim}

\end{sphinxuseclass}\end{sphinxVerbatimOutput}

\end{sphinxuseclass}
\begin{sphinxuseclass}{cell}\begin{sphinxVerbatimInput}

\begin{sphinxuseclass}{cell_input}
\begin{sphinxVerbatim}[commandchars=\\\{\}]
\PYG{n}{criterion} \PYG{o}{=} \PYG{n}{nn}\PYG{o}{.}\PYG{n}{MSELoss}\PYG{p}{(}\PYG{p}{)}  \PYG{c+c1}{\PYGZsh{} Mean Squared Error loss}
\PYG{n}{optimizer} \PYG{o}{=} \PYG{n}{optim}\PYG{o}{.}\PYG{n}{Adam}\PYG{p}{(}\PYG{n}{cnn\PYGZus{}model}\PYG{o}{.}\PYG{n}{parameters}\PYG{p}{(}\PYG{p}{)}\PYG{p}{,} \PYG{n}{lr}\PYG{o}{=}\PYG{l+m+mf}{0.001}\PYG{p}{)}

\PYG{n}{epochs} \PYG{o}{=} \PYG{l+m+mi}{20}
\PYG{k}{for} \PYG{n}{epoch} \PYG{o+ow}{in} \PYG{p}{(}\PYG{n}{pbar}\PYG{o}{:=}\PYG{n}{tqdm}\PYG{p}{(}\PYG{n+nb}{range}\PYG{p}{(}\PYG{l+m+mi}{1}\PYG{p}{,} \PYG{n}{epochs}\PYG{o}{+}\PYG{l+m+mi}{1}\PYG{p}{)}\PYG{p}{)}\PYG{p}{)}\PYG{p}{:}
    \PYG{n}{cnn\PYGZus{}model}\PYG{o}{.}\PYG{n}{train}\PYG{p}{(}\PYG{p}{)}
    \PYG{n}{train\PYGZus{}loss} \PYG{o}{=} \PYG{l+m+mf}{0.0}
    \PYG{k}{for} \PYG{n}{batch\PYGZus{}images}\PYG{p}{,} \PYG{n}{batch\PYGZus{}keypoints} \PYG{o+ow}{in} \PYG{n}{train\PYGZus{}loader}\PYG{p}{:}
        \PYG{n}{images} \PYG{o}{=} \PYG{n}{batch\PYGZus{}images}
        \PYG{n}{keypoints} \PYG{o}{=} \PYG{n}{batch\PYGZus{}keypoints}

        \PYG{n}{outputs} \PYG{o}{=} \PYG{n}{cnn\PYGZus{}model}\PYG{p}{(}\PYG{n}{images}\PYG{p}{)}
        \PYG{n}{loss} \PYG{o}{=} \PYG{n}{criterion}\PYG{p}{(}\PYG{n}{outputs}\PYG{p}{,} \PYG{n}{keypoints}\PYG{p}{)}

        \PYG{n}{optimizer}\PYG{o}{.}\PYG{n}{zero\PYGZus{}grad}\PYG{p}{(}\PYG{p}{)}
        \PYG{n}{loss}\PYG{o}{.}\PYG{n}{backward}\PYG{p}{(}\PYG{p}{)}
        \PYG{n}{optimizer}\PYG{o}{.}\PYG{n}{step}\PYG{p}{(}\PYG{p}{)}

        \PYG{n}{train\PYGZus{}loss} \PYG{o}{+}\PYG{o}{=} \PYG{n}{loss}\PYG{o}{.}\PYG{n}{item}\PYG{p}{(}\PYG{p}{)} \PYG{o}{*} \PYG{n}{images}\PYG{o}{.}\PYG{n}{size}\PYG{p}{(}\PYG{l+m+mi}{0}\PYG{p}{)}
    \PYG{n}{train\PYGZus{}loss} \PYG{o}{/}\PYG{o}{=} \PYG{n+nb}{len}\PYG{p}{(}\PYG{n}{train\PYGZus{}loader}\PYG{o}{.}\PYG{n}{dataset}\PYG{p}{)}

    \PYG{c+c1}{\PYGZsh{} Validation phase}
    \PYG{n}{cnn\PYGZus{}model}\PYG{o}{.}\PYG{n}{eval}\PYG{p}{(}\PYG{p}{)}
    \PYG{n}{val\PYGZus{}loss} \PYG{o}{=} \PYG{l+m+mf}{0.0}
    \PYG{k}{with} \PYG{n}{torch}\PYG{o}{.}\PYG{n}{no\PYGZus{}grad}\PYG{p}{(}\PYG{p}{)}\PYG{p}{:}
        \PYG{k}{for} \PYG{n}{batch\PYGZus{}images}\PYG{p}{,} \PYG{n}{batch\PYGZus{}keypoints} \PYG{o+ow}{in} \PYG{n}{val\PYGZus{}loader}\PYG{p}{:}
            \PYG{n}{images} \PYG{o}{=} \PYG{n}{batch\PYGZus{}images}
            \PYG{n}{keypoints} \PYG{o}{=} \PYG{n}{batch\PYGZus{}keypoints}
            \PYG{n}{preds} \PYG{o}{=} \PYG{n}{cnn\PYGZus{}model}\PYG{p}{(}\PYG{n}{images}\PYG{p}{)}
            \PYG{n}{loss} \PYG{o}{=} \PYG{n}{criterion}\PYG{p}{(}\PYG{n}{preds}\PYG{p}{,} \PYG{n}{keypoints}\PYG{p}{)}
            \PYG{n}{val\PYGZus{}loss} \PYG{o}{+}\PYG{o}{=} \PYG{n}{loss}\PYG{o}{.}\PYG{n}{item}\PYG{p}{(}\PYG{p}{)} \PYG{o}{*} \PYG{n}{images}\PYG{o}{.}\PYG{n}{size}\PYG{p}{(}\PYG{l+m+mi}{0}\PYG{p}{)}
    \PYG{n}{val\PYGZus{}loss} \PYG{o}{/}\PYG{o}{=} \PYG{n+nb}{len}\PYG{p}{(}\PYG{n}{val\PYGZus{}loader}\PYG{o}{.}\PYG{n}{dataset}\PYG{p}{)}

    \PYG{n}{pbar}\PYG{o}{.}\PYG{n}{set\PYGZus{}description}\PYG{p}{(}\PYG{l+s+sa}{f}\PYG{l+s+s2}{\PYGZdq{}}\PYG{l+s+s2}{Epoch }\PYG{l+s+si}{\PYGZob{}}\PYG{n}{epoch}\PYG{l+s+si}{\PYGZcb{}}\PYG{l+s+s2}{/}\PYG{l+s+si}{\PYGZob{}}\PYG{n}{epochs}\PYG{l+s+si}{\PYGZcb{}}\PYG{l+s+s2}{ \PYGZhy{} Training Loss: }\PYG{l+s+si}{\PYGZob{}}\PYG{n}{train\PYGZus{}loss}\PYG{l+s+si}{:}\PYG{l+s+s2}{.4f}\PYG{l+s+si}{\PYGZcb{}}\PYG{l+s+s2}{ \PYGZhy{} Validation Loss: }\PYG{l+s+si}{\PYGZob{}}\PYG{n}{val\PYGZus{}loss}\PYG{l+s+si}{:}\PYG{l+s+s2}{.4f}\PYG{l+s+si}{\PYGZcb{}}\PYG{l+s+s2}{\PYGZdq{}}\PYG{p}{)}
\end{sphinxVerbatim}

\end{sphinxuseclass}\end{sphinxVerbatimInput}
\begin{sphinxVerbatimOutput}

\begin{sphinxuseclass}{cell_output}
\begin{sphinxVerbatim}[commandchars=\\\{\}]
  0\PYGZpc{}|          | 0/20 [00:00\PYGZlt{}?, ?it/s]
\end{sphinxVerbatim}

\end{sphinxuseclass}\end{sphinxVerbatimOutput}

\end{sphinxuseclass}

\section{Evaluation and results}
\label{\detokenize{04:evaluation-and-results}}
\sphinxAtStartPar
After training, we’ll evaluate the model’s performance. One common metric for this task is \sphinxstylestrong{Root Mean Squared Error (RMSE)} of the keypoint predictions. RMSE is the square root of MSE, giving an error in the same units as the coordinates (pixels). We will calculate RMSE on the validation set.

\sphinxAtStartPar
For instance, an RMSE of 4.0 would mean on average the predictions are about 4 pixels away from the true keypoint positions.

\begin{sphinxuseclass}{cell}\begin{sphinxVerbatimInput}

\begin{sphinxuseclass}{cell_input}
\begin{sphinxVerbatim}[commandchars=\\\{\}]
\PYG{c+c1}{\PYGZsh{} Switch to evaluation mode}
\PYG{n}{cnn\PYGZus{}model}\PYG{o}{.}\PYG{n}{eval}\PYG{p}{(}\PYG{p}{)}

\PYG{c+c1}{\PYGZsh{} Predict on the entire validation set}
\PYG{k}{with} \PYG{n}{torch}\PYG{o}{.}\PYG{n}{no\PYGZus{}grad}\PYG{p}{(}\PYG{p}{)}\PYG{p}{:}
    \PYG{n}{val\PYGZus{}preds} \PYG{o}{=} \PYG{n}{cnn\PYGZus{}model}\PYG{p}{(}\PYG{n}{torch}\PYG{o}{.}\PYG{n}{Tensor}\PYG{p}{(}\PYG{n}{X\PYGZus{}val\PYGZus{}cnn}\PYG{p}{)}\PYG{o}{.}\PYG{n}{to}\PYG{p}{(}\PYG{n}{device}\PYG{p}{)}\PYG{p}{)}\PYG{o}{.}\PYG{n}{cpu}\PYG{p}{(}\PYG{p}{)}\PYG{o}{.}\PYG{n}{numpy}\PYG{p}{(}\PYG{p}{)}

\PYG{c+c1}{\PYGZsh{} Compute RMSE on validation set}
\PYG{n}{val\PYGZus{}truth} \PYG{o}{=} \PYG{n}{y\PYGZus{}val\PYGZus{}cnn}  \PYG{c+c1}{\PYGZsh{} actual keypoints}
\PYG{n}{mse\PYGZus{}val} \PYG{o}{=} \PYG{n}{np}\PYG{o}{.}\PYG{n}{mean}\PYG{p}{(}\PYG{p}{(}\PYG{n}{val\PYGZus{}preds} \PYG{o}{\PYGZhy{}} \PYG{n}{val\PYGZus{}truth}\PYG{p}{)}\PYG{o}{*}\PYG{o}{*}\PYG{l+m+mi}{2}\PYG{p}{)}
\PYG{n}{rmse\PYGZus{}val} \PYG{o}{=} \PYG{n}{np}\PYG{o}{.}\PYG{n}{sqrt}\PYG{p}{(}\PYG{n}{mse\PYGZus{}val}\PYG{p}{)}
\PYG{n+nb}{print}\PYG{p}{(}\PYG{l+s+sa}{f}\PYG{l+s+s2}{\PYGZdq{}}\PYG{l+s+s2}{Validation RMSE: }\PYG{l+s+si}{\PYGZob{}}\PYG{n}{rmse\PYGZus{}val}\PYG{l+s+si}{:}\PYG{l+s+s2}{.3f}\PYG{l+s+si}{\PYGZcb{}}\PYG{l+s+s2}{ pixels}\PYG{l+s+s2}{\PYGZdq{}}\PYG{p}{)}
\end{sphinxVerbatim}

\end{sphinxuseclass}\end{sphinxVerbatimInput}
\begin{sphinxVerbatimOutput}

\begin{sphinxuseclass}{cell_output}
\begin{sphinxVerbatim}[commandchars=\\\{\}]
Validation RMSE: 3.167 pixels
\end{sphinxVerbatim}

\end{sphinxuseclass}\end{sphinxVerbatimOutput}

\end{sphinxuseclass}

\subsection{Visualizing predictions}
\label{\detokenize{04:visualizing-predictions}}
\sphinxAtStartPar
Finally, let’s visualize some predictions to see how the model is performing. We’ll take a few images from the validation set (which the model hasn’t seen during training) and plot the image with:
\begin{itemize}
\item {} 
\sphinxAtStartPar
The \sphinxstylestrong{ground truth} keypoints (in green)

\item {} 
\sphinxAtStartPar
The \sphinxstylestrong{predicted} keypoints (in red)

\end{itemize}

\begin{sphinxuseclass}{cell}\begin{sphinxVerbatimInput}

\begin{sphinxuseclass}{cell_input}
\begin{sphinxVerbatim}[commandchars=\\\{\}]
\PYG{c+c1}{\PYGZsh{} Choose a few samples from validation to visualize}
\PYG{n}{num\PYGZus{}samples\PYGZus{}to\PYGZus{}show} \PYG{o}{=} \PYG{l+m+mi}{5}
\PYG{n}{indices} \PYG{o}{=} \PYG{n}{np}\PYG{o}{.}\PYG{n}{random}\PYG{o}{.}\PYG{n}{choice}\PYG{p}{(}\PYG{n+nb}{len}\PYG{p}{(}\PYG{n}{X\PYGZus{}val}\PYG{p}{)}\PYG{p}{,} \PYG{n}{size}\PYG{o}{=}\PYG{n}{num\PYGZus{}samples\PYGZus{}to\PYGZus{}show}\PYG{p}{,} \PYG{n}{replace}\PYG{o}{=}\PYG{k+kc}{False}\PYG{p}{)}

\PYG{n}{fig}\PYG{p}{,} \PYG{n}{axes} \PYG{o}{=} \PYG{n}{plt}\PYG{o}{.}\PYG{n}{subplots}\PYG{p}{(}\PYG{l+m+mi}{1}\PYG{p}{,} \PYG{n}{num\PYGZus{}samples\PYGZus{}to\PYGZus{}show}\PYG{p}{,} \PYG{n}{figsize}\PYG{o}{=}\PYG{p}{(}\PYG{l+m+mi}{10}\PYG{p}{,} \PYG{l+m+mi}{6}\PYG{p}{)}\PYG{p}{)}
\PYG{k}{for} \PYG{n}{i}\PYG{p}{,} \PYG{n}{idx} \PYG{o+ow}{in} \PYG{n+nb}{enumerate}\PYG{p}{(}\PYG{n}{indices}\PYG{p}{)}\PYG{p}{:}
    \PYG{n}{img} \PYG{o}{=} \PYG{n}{X\PYGZus{}val\PYGZus{}fcn}\PYG{p}{[}\PYG{n}{idx}\PYG{p}{]}\PYG{o}{.}\PYG{n}{reshape}\PYG{p}{(}\PYG{l+m+mi}{96}\PYG{p}{,} \PYG{l+m+mi}{96}\PYG{p}{)}
    \PYG{n}{true\PYGZus{}kp} \PYG{o}{=} \PYG{n}{y\PYGZus{}val\PYGZus{}fcn}\PYG{p}{[}\PYG{n}{idx}\PYG{p}{]}
    \PYG{n}{cnn\PYGZus{}pred\PYGZus{}kp} \PYG{o}{=} \PYG{n}{val\PYGZus{}preds}\PYG{p}{[}\PYG{n}{idx}\PYG{p}{]}
    \PYG{n}{fcn\PYGZus{}pred\PYGZus{}kp} \PYG{o}{=} \PYG{n}{fcn\PYGZus{}model}\PYG{p}{(}\PYG{n}{torch}\PYG{o}{.}\PYG{n}{tensor}\PYG{p}{(}\PYG{n}{X\PYGZus{}val\PYGZus{}fcn}\PYG{p}{[}\PYG{n}{idx}\PYG{p}{]}\PYG{o}{.}\PYG{n}{reshape}\PYG{p}{(}\PYG{l+m+mi}{1}\PYG{p}{,} \PYG{l+m+mi}{96}\PYG{o}{*}\PYG{l+m+mi}{96}\PYG{p}{)}\PYG{p}{)}\PYG{o}{.}\PYG{n}{to}\PYG{p}{(}\PYG{n}{device}\PYG{p}{)}\PYG{p}{)}\PYG{o}{.}\PYG{n}{cpu}\PYG{p}{(}\PYG{p}{)}\PYG{o}{.}\PYG{n}{detach}\PYG{p}{(}\PYG{p}{)}\PYG{o}{.}\PYG{n}{numpy}\PYG{p}{(}\PYG{p}{)}\PYG{p}{[}\PYG{l+m+mi}{0}\PYG{p}{]}

    \PYG{n}{axes}\PYG{p}{[}\PYG{n}{i}\PYG{p}{]}\PYG{o}{.}\PYG{n}{imshow}\PYG{p}{(}\PYG{n}{img}\PYG{p}{,} \PYG{n}{cmap}\PYG{o}{=}\PYG{l+s+s1}{\PYGZsq{}}\PYG{l+s+s1}{gray}\PYG{l+s+s1}{\PYGZsq{}}\PYG{p}{)}
    \PYG{c+c1}{\PYGZsh{} Plot true keypoints in green}
    \PYG{n}{axes}\PYG{p}{[}\PYG{n}{i}\PYG{p}{]}\PYG{o}{.}\PYG{n}{scatter}\PYG{p}{(}\PYG{n}{true\PYGZus{}kp}\PYG{p}{[}\PYG{l+m+mi}{0}\PYG{p}{:}\PYG{p}{:}\PYG{l+m+mi}{2}\PYG{p}{]}\PYG{p}{,} \PYG{n}{true\PYGZus{}kp}\PYG{p}{[}\PYG{l+m+mi}{1}\PYG{p}{:}\PYG{p}{:}\PYG{l+m+mi}{2}\PYG{p}{]}\PYG{p}{,} \PYG{n}{c}\PYG{o}{=}\PYG{l+s+s1}{\PYGZsq{}}\PYG{l+s+s1}{green}\PYG{l+s+s1}{\PYGZsq{}}\PYG{p}{,} \PYG{n}{s}\PYG{o}{=}\PYG{l+m+mi}{20}\PYG{p}{,} \PYG{n}{label}\PYG{o}{=}\PYG{l+s+s1}{\PYGZsq{}}\PYG{l+s+s1}{True}\PYG{l+s+s1}{\PYGZsq{}}\PYG{p}{)}
    \PYG{c+c1}{\PYGZsh{} Plot predicted keypoints in red}
    \PYG{n}{axes}\PYG{p}{[}\PYG{n}{i}\PYG{p}{]}\PYG{o}{.}\PYG{n}{scatter}\PYG{p}{(}\PYG{n}{fcn\PYGZus{}pred\PYGZus{}kp}\PYG{p}{[}\PYG{l+m+mi}{0}\PYG{p}{:}\PYG{p}{:}\PYG{l+m+mi}{2}\PYG{p}{]}\PYG{p}{,} \PYG{n}{fcn\PYGZus{}pred\PYGZus{}kp}\PYG{p}{[}\PYG{l+m+mi}{1}\PYG{p}{:}\PYG{p}{:}\PYG{l+m+mi}{2}\PYG{p}{]}\PYG{p}{,} \PYG{n}{c}\PYG{o}{=}\PYG{l+s+s1}{\PYGZsq{}}\PYG{l+s+s1}{red}\PYG{l+s+s1}{\PYGZsq{}}\PYG{p}{,} \PYG{n}{s}\PYG{o}{=}\PYG{l+m+mi}{20}\PYG{p}{,} \PYG{n}{marker}\PYG{o}{=}\PYG{l+s+s1}{\PYGZsq{}}\PYG{l+s+s1}{x}\PYG{l+s+s1}{\PYGZsq{}}\PYG{p}{,} \PYG{n}{label}\PYG{o}{=}\PYG{l+s+s1}{\PYGZsq{}}\PYG{l+s+s1}{Pred}\PYG{l+s+s1}{\PYGZsq{}}\PYG{p}{)}
    \PYG{c+c1}{\PYGZsh{} Plot predicted keypoints in red}
    \PYG{n}{axes}\PYG{p}{[}\PYG{n}{i}\PYG{p}{]}\PYG{o}{.}\PYG{n}{scatter}\PYG{p}{(}\PYG{n}{cnn\PYGZus{}pred\PYGZus{}kp}\PYG{p}{[}\PYG{l+m+mi}{0}\PYG{p}{:}\PYG{p}{:}\PYG{l+m+mi}{2}\PYG{p}{]}\PYG{p}{,} \PYG{n}{cnn\PYGZus{}pred\PYGZus{}kp}\PYG{p}{[}\PYG{l+m+mi}{1}\PYG{p}{:}\PYG{p}{:}\PYG{l+m+mi}{2}\PYG{p}{]}\PYG{p}{,} \PYG{n}{c}\PYG{o}{=}\PYG{l+s+s1}{\PYGZsq{}}\PYG{l+s+s1}{blue}\PYG{l+s+s1}{\PYGZsq{}}\PYG{p}{,} \PYG{n}{s}\PYG{o}{=}\PYG{l+m+mi}{20}\PYG{p}{,} \PYG{n}{marker}\PYG{o}{=}\PYG{l+s+s1}{\PYGZsq{}}\PYG{l+s+s1}{x}\PYG{l+s+s1}{\PYGZsq{}}\PYG{p}{,} \PYG{n}{label}\PYG{o}{=}\PYG{l+s+s1}{\PYGZsq{}}\PYG{l+s+s1}{Pred}\PYG{l+s+s1}{\PYGZsq{}}\PYG{p}{)}
    \PYG{n}{axes}\PYG{p}{[}\PYG{n}{i}\PYG{p}{]}\PYG{o}{.}\PYG{n}{axis}\PYG{p}{(}\PYG{l+s+s1}{\PYGZsq{}}\PYG{l+s+s1}{off}\PYG{l+s+s1}{\PYGZsq{}}\PYG{p}{)}

\PYG{n}{plt}\PYG{o}{.}\PYG{n}{tight\PYGZus{}layout}\PYG{p}{(}\PYG{p}{)}
\PYG{n}{plt}\PYG{o}{.}\PYG{n}{show}\PYG{p}{(}\PYG{p}{)}
\end{sphinxVerbatim}

\end{sphinxuseclass}\end{sphinxVerbatimInput}
\begin{sphinxVerbatimOutput}

\begin{sphinxuseclass}{cell_output}
\noindent\sphinxincludegraphics{{896322e73b9fb09d1fb112805f4be3e6ffd416b6eb46cc915c21b1f8d2d892ca}.png}

\end{sphinxuseclass}\end{sphinxVerbatimOutput}

\end{sphinxuseclass}
\sphinxAtStartPar
In the displayed images, \sphinxstylestrong{green} dots should mark the actual keypoint positions and \sphinxstylestrong{red} x’s mark the model’s predicted positions. Ideally, they will be close for most points. You might notice the model does well on prominent features like the eye centers or nose tip, but could be less accurate on some others, especially if the network hasn’t fully converged or if data is limited. A deeper network or additional data augmentation could further improve performance.

\sphinxstepscope


\chapter{Ensembles of models and feature selection}
\label{\detokenize{05:ensembles-of-models-and-feature-selection}}\label{\detokenize{05::doc}}
\sphinxAtStartPar
Let us define some basic imports and a \sphinxcode{\sphinxupquote{pprint}} function.

\begin{sphinxuseclass}{cell}\begin{sphinxVerbatimInput}

\begin{sphinxuseclass}{cell_input}
\begin{sphinxVerbatim}[commandchars=\\\{\}]
\PYG{k+kn}{import}\PYG{+w}{ }\PYG{n+nn}{numpy}\PYG{+w}{ }\PYG{k}{as}\PYG{+w}{ }\PYG{n+nn}{np}
\PYG{k+kn}{import}\PYG{+w}{ }\PYG{n+nn}{pandas}\PYG{+w}{ }\PYG{k}{as}\PYG{+w}{ }\PYG{n+nn}{pd}
\PYG{k+kn}{import}\PYG{+w}{ }\PYG{n+nn}{sklearn}
\PYG{k+kn}{import}\PYG{+w}{ }\PYG{n+nn}{warnings}
\PYG{n}{warnings}\PYG{o}{.}\PYG{n}{filterwarnings}\PYG{p}{(}\PYG{l+s+s1}{\PYGZsq{}}\PYG{l+s+s1}{ignore}\PYG{l+s+s1}{\PYGZsq{}}\PYG{p}{)}
\PYG{k+kn}{from}\PYG{+w}{ }\PYG{n+nn}{tqdm}\PYG{n+nn}{.}\PYG{n+nn}{auto}\PYG{+w}{ }\PYG{k+kn}{import} \PYG{n}{tqdm}
\PYG{k+kn}{from}\PYG{+w}{ }\PYG{n+nn}{IPython}\PYG{n+nn}{.}\PYG{n+nn}{display}\PYG{+w}{ }\PYG{k+kn}{import} \PYG{n}{display}\PYG{p}{,} \PYG{n}{Math}
\PYG{k+kn}{from}\PYG{+w}{ }\PYG{n+nn}{numpyarray\PYGZus{}to\PYGZus{}latex}\PYG{n+nn}{.}\PYG{n+nn}{jupyter}\PYG{+w}{ }\PYG{k+kn}{import} \PYG{n}{to\PYGZus{}ltx}
\PYG{k+kn}{import}\PYG{+w}{ }\PYG{n+nn}{matplotlib}\PYG{n+nn}{.}\PYG{n+nn}{pyplot}\PYG{+w}{ }\PYG{k}{as}\PYG{+w}{ }\PYG{n+nn}{plt}
\PYG{n}{plt}\PYG{o}{.}\PYG{n}{style}\PYG{o}{.}\PYG{n}{use}\PYG{p}{(}\PYG{p}{[}\PYG{l+s+s1}{\PYGZsq{}}\PYG{l+s+s1}{seaborn\PYGZhy{}v0\PYGZus{}8\PYGZhy{}muted}\PYG{l+s+s1}{\PYGZsq{}}\PYG{p}{,} \PYG{l+s+s1}{\PYGZsq{}}\PYG{l+s+s1}{practicals.mplstyle}\PYG{l+s+s1}{\PYGZsq{}}\PYG{p}{]}\PYG{p}{)}
\PYG{k+kn}{import}\PYG{+w}{ }\PYG{n+nn}{seaborn}\PYG{+w}{ }\PYG{k}{as}\PYG{+w}{ }\PYG{n+nn}{sns}
\PYG{k}{def}\PYG{+w}{ }\PYG{n+nf}{pprint}\PYG{p}{(}\PYG{o}{*}\PYG{n}{args}\PYG{p}{)}\PYG{p}{:}
    \PYG{n}{res} \PYG{o}{=} \PYG{l+s+s2}{\PYGZdq{}}\PYG{l+s+s2}{\PYGZdq{}}
    \PYG{k}{for} \PYG{n}{i} \PYG{o+ow}{in} \PYG{n}{args}\PYG{p}{:}
        \PYG{k}{if} \PYG{n+nb}{type}\PYG{p}{(}\PYG{n}{i}\PYG{p}{)} \PYG{o}{==} \PYG{n}{np}\PYG{o}{.}\PYG{n}{ndarray}\PYG{p}{:}
            \PYG{n}{res} \PYG{o}{+}\PYG{o}{=} \PYG{n}{to\PYGZus{}ltx}\PYG{p}{(}\PYG{n}{i}\PYG{p}{,} \PYG{n}{brackets}\PYG{o}{=}\PYG{l+s+s1}{\PYGZsq{}}\PYG{l+s+s1}{[]}\PYG{l+s+s1}{\PYGZsq{}}\PYG{p}{)}
        \PYG{k}{elif} \PYG{n+nb}{type}\PYG{p}{(}\PYG{n}{i}\PYG{p}{)} \PYG{o}{==} \PYG{n+nb}{str}\PYG{p}{:}
            \PYG{n}{res} \PYG{o}{+}\PYG{o}{=} \PYG{n}{i}
    \PYG{n}{display}\PYG{p}{(}\PYG{n}{Math}\PYG{p}{(}\PYG{n}{res}\PYG{p}{)}\PYG{p}{)}
\end{sphinxVerbatim}

\end{sphinxuseclass}\end{sphinxVerbatimInput}

\end{sphinxuseclass}

\section{Reminder:  Supervised learning}
\label{\detokenize{05:reminder-supervised-learning}}
\sphinxAtStartPar
Supervised learning aims to model the (unknown) stochastic dependence between a set of \(n\) inputs \(x\) (also called features) and an output \({\mathbf y}\) on the basis of a training set \(D_N\) of size \(N\). Supervised learning tasks are decomposed into \sphinxstyleemphasis{regression} and \sphinxstyleemphasis{classification} tasks.


\subsection{Regression}
\label{\detokenize{05:regression}}
\sphinxAtStartPar
In regression the stochastic dependence is given by
\begin{equation*}
\begin{split} {\mathbf y} = f(x)+\mathbf{w}\end{split}
\end{equation*}
\sphinxAtStartPar
where:
\begin{itemize}
\item {} 
\sphinxAtStartPar
\(\mathbf{y} \in \mathbb{R}\) represents the output variable (also called target)

\item {} 
\sphinxAtStartPar
\(x \in \mathbb{R}^n\) represents the vector of inputs (also called features)

\item {} 
\sphinxAtStartPar
\(\mathbf{w}\) denotes the noise, and is typically assumed that \(E[\mathbf{w}]=0\) and the variance
\(\text{Var}[\mathbf{w}]=\sigma^2_{\mathbf{w}}\) is constant

\item {} 
\sphinxAtStartPar
\(f(x)=E[{\mathbf y} | x]\) is the (unknown) mapping between input and outputs, also known as the \sphinxstyleemphasis{regression function}.

\end{itemize}

\sphinxAtStartPar
In regression, the goal of learning is to return an estimator
\begin{equation*}
\begin{split}h(x,\alpha)\end{split}
\end{equation*}
\sphinxAtStartPar
of the regression function \(f(x)\), where \(\alpha\) denotes the set of parameters of the model \(h\).


\subsection{Error measurement}
\label{\detokenize{05:error-measurement}}
\sphinxAtStartPar
We already have used the \sphinxstyleemphasis{mean squared error} as an evaluation tool for the quality of the predictions of a model. Unfortunately, this tool does not take into account how “variable” the expected outputs are, and therefore two models could become undistinguishable when using the MSE. We can normalize the MSE by dividing it by the (measured) variance of the output data. This leads to the \sphinxstyleemphasis{normalized mean squared error}:
\begin{equation*}
\begin{split}\text{NMSE} = \frac{\sum(y-\hat{y})^2}{N\sigma^2_y}\end{split}
\end{equation*}
\sphinxAtStartPar
We can see that the most naive model which outputs the mean of the observed outputs should give:
\begin{equation*}
\begin{split}\frac{\sum(y-\mu_y)^2}{N\sigma^2_y} = \frac{\sigma^2_y}{\sigma^2_y} = 1\end{split}
\end{equation*}
\sphinxAtStartPar
Any model that has a NMSE greater than \(1\) is therefore performing worse than the most naive model.

\begin{sphinxuseclass}{cell}\begin{sphinxVerbatimInput}

\begin{sphinxuseclass}{cell_input}
\begin{sphinxVerbatim}[commandchars=\\\{\}]
\PYG{k+kn}{from}\PYG{+w}{ }\PYG{n+nn}{sklearn}\PYG{n+nn}{.}\PYG{n+nn}{linear\PYGZus{}model}\PYG{+w}{ }\PYG{k+kn}{import} \PYG{n}{LinearRegression}
\PYG{n}{x} \PYG{o}{=} \PYG{n}{np}\PYG{o}{.}\PYG{n}{linspace}\PYG{p}{(}\PYG{o}{\PYGZhy{}}\PYG{l+m+mi}{2}\PYG{p}{,} \PYG{l+m+mi}{2}\PYG{p}{,} \PYG{l+m+mi}{1000}\PYG{p}{)}\PYG{o}{.}\PYG{n}{reshape}\PYG{p}{(}\PYG{l+m+mi}{1000}\PYG{p}{,}\PYG{o}{\PYGZhy{}}\PYG{l+m+mi}{1}\PYG{p}{)}
\PYG{n}{B1}\PYG{p}{,} \PYG{n}{B2} \PYG{o}{=} \PYG{n}{np}\PYG{o}{.}\PYG{n}{random}\PYG{o}{.}\PYG{n}{normal}\PYG{p}{(}\PYG{l+m+mi}{0}\PYG{p}{,} \PYG{l+m+mi}{1}\PYG{p}{,} \PYG{n}{size}\PYG{o}{=}\PYG{p}{(}\PYG{l+m+mi}{1000}\PYG{p}{,}\PYG{l+m+mi}{1}\PYG{p}{)}\PYG{p}{)}\PYG{p}{,} \PYG{n}{np}\PYG{o}{.}\PYG{n}{random}\PYG{o}{.}\PYG{n}{normal}\PYG{p}{(}\PYG{l+m+mi}{0}\PYG{p}{,} \PYG{l+m+mi}{10}\PYG{p}{,} \PYG{n}{size}\PYG{o}{=}\PYG{p}{(}\PYG{l+m+mi}{1000}\PYG{p}{,}\PYG{l+m+mi}{1}\PYG{p}{)}\PYG{p}{)}
\PYG{n}{A} \PYG{o}{=} \PYG{l+m+mi}{2}
\PYG{n}{y1} \PYG{o}{=} \PYG{n}{A}\PYG{o}{*}\PYG{n}{x}\PYG{o}{+}\PYG{n}{B1}
\PYG{n}{y2} \PYG{o}{=} \PYG{n}{A}\PYG{o}{*}\PYG{n}{x}\PYG{o}{+}\PYG{n}{B2}

\PYG{n}{model1} \PYG{o}{=} \PYG{n}{LinearRegression}\PYG{p}{(}\PYG{p}{)}\PYG{o}{.}\PYG{n}{fit}\PYG{p}{(}\PYG{n}{x}\PYG{p}{,} \PYG{n}{y1}\PYG{p}{)}
\PYG{n}{model2} \PYG{o}{=} \PYG{n}{LinearRegression}\PYG{p}{(}\PYG{p}{)}\PYG{o}{.}\PYG{n}{fit}\PYG{p}{(}\PYG{n}{x}\PYG{p}{,} \PYG{n}{y2}\PYG{p}{)}

\PYG{n}{plt}\PYG{o}{.}\PYG{n}{figure}\PYG{p}{(}\PYG{n}{figsize}\PYG{o}{=}\PYG{p}{(}\PYG{l+m+mi}{10}\PYG{p}{,} \PYG{l+m+mi}{6}\PYG{p}{)}\PYG{p}{)}
\PYG{n}{plt}\PYG{o}{.}\PYG{n}{scatter}\PYG{p}{(}\PYG{n}{x}\PYG{p}{,} \PYG{n}{y1}\PYG{p}{,} \PYG{n}{marker}\PYG{o}{=}\PYG{l+s+s1}{\PYGZsq{}}\PYG{l+s+s1}{.}\PYG{l+s+s1}{\PYGZsq{}}\PYG{p}{,} \PYG{n}{label}\PYG{o}{=}\PYG{l+s+s1}{\PYGZsq{}}\PYG{l+s+s1}{y1}\PYG{l+s+s1}{\PYGZsq{}}\PYG{p}{)}
\PYG{n}{plt}\PYG{o}{.}\PYG{n}{plot}\PYG{p}{(}\PYG{n}{x}\PYG{p}{,} \PYG{n}{model1}\PYG{o}{.}\PYG{n}{predict}\PYG{p}{(}\PYG{n}{x}\PYG{p}{)}\PYG{p}{,} \PYG{n}{color}\PYG{o}{=}\PYG{l+s+s1}{\PYGZsq{}}\PYG{l+s+s1}{red}\PYG{l+s+s1}{\PYGZsq{}}\PYG{p}{,} \PYG{n}{label}\PYG{o}{=}\PYG{l+s+s1}{\PYGZsq{}}\PYG{l+s+s1}{model1}\PYG{l+s+s1}{\PYGZsq{}}\PYG{p}{)}
\PYG{n}{plt}\PYG{o}{.}\PYG{n}{scatter}\PYG{p}{(}\PYG{n}{x}\PYG{p}{,} \PYG{n}{y2}\PYG{p}{,} \PYG{n}{marker}\PYG{o}{=}\PYG{l+s+s1}{\PYGZsq{}}\PYG{l+s+s1}{.}\PYG{l+s+s1}{\PYGZsq{}}\PYG{p}{,} \PYG{n}{label}\PYG{o}{=}\PYG{l+s+s1}{\PYGZsq{}}\PYG{l+s+s1}{y2}\PYG{l+s+s1}{\PYGZsq{}}\PYG{p}{)}
\PYG{n}{plt}\PYG{o}{.}\PYG{n}{plot}\PYG{p}{(}\PYG{n}{x}\PYG{p}{,} \PYG{n}{model2}\PYG{o}{.}\PYG{n}{predict}\PYG{p}{(}\PYG{n}{x}\PYG{p}{)}\PYG{p}{,} \PYG{n}{color}\PYG{o}{=}\PYG{l+s+s1}{\PYGZsq{}}\PYG{l+s+s1}{lightgreen}\PYG{l+s+s1}{\PYGZsq{}}\PYG{p}{,} \PYG{n}{label}\PYG{o}{=}\PYG{l+s+s1}{\PYGZsq{}}\PYG{l+s+s1}{model2}\PYG{l+s+s1}{\PYGZsq{}}\PYG{p}{)}
\PYG{n}{plt}\PYG{o}{.}\PYG{n}{title}\PYG{p}{(}\PYG{l+s+s2}{\PYGZdq{}}\PYG{l+s+s2}{Two linear models on data with different variances}\PYG{l+s+s2}{\PYGZdq{}}\PYG{p}{)}
\PYG{n}{plt}\PYG{o}{.}\PYG{n}{legend}\PYG{p}{(}\PYG{p}{)}
\PYG{n}{plt}\PYG{o}{.}\PYG{n}{grid}\PYG{p}{(}\PYG{l+s+s2}{\PYGZdq{}}\PYG{l+s+s2}{on}\PYG{l+s+s2}{\PYGZdq{}}\PYG{p}{)}
\PYG{n}{plt}\PYG{o}{.}\PYG{n}{show}\PYG{p}{(}\PYG{p}{)}
\end{sphinxVerbatim}

\end{sphinxuseclass}\end{sphinxVerbatimInput}
\begin{sphinxVerbatimOutput}

\begin{sphinxuseclass}{cell_output}
\noindent\sphinxincludegraphics{{2b28417cb3777dd44b57925fa60db1cc5b7e324df3b526e7642cc8a76a94c065}.png}

\end{sphinxuseclass}\end{sphinxVerbatimOutput}

\end{sphinxuseclass}
\begin{sphinxuseclass}{cell}\begin{sphinxVerbatimInput}

\begin{sphinxuseclass}{cell_input}
\begin{sphinxVerbatim}[commandchars=\\\{\}]
\PYG{k}{def}\PYG{+w}{ }\PYG{n+nf}{mse}\PYG{p}{(}\PYG{n}{y}\PYG{p}{,} \PYG{n}{y\PYGZus{}hat}\PYG{p}{)}\PYG{p}{:}
    \PYG{k}{return} \PYG{n}{np}\PYG{o}{.}\PYG{n}{mean}\PYG{p}{(}\PYG{p}{(}\PYG{n}{y}\PYG{o}{\PYGZhy{}}\PYG{n}{y\PYGZus{}hat}\PYG{p}{)}\PYG{o}{*}\PYG{o}{*}\PYG{l+m+mi}{2}\PYG{p}{)}

\PYG{k}{def}\PYG{+w}{ }\PYG{n+nf}{nmse}\PYG{p}{(}\PYG{n}{y}\PYG{p}{,} \PYG{n}{y\PYGZus{}hat}\PYG{p}{)}\PYG{p}{:}
    \PYG{k}{return} \PYG{n}{np}\PYG{o}{.}\PYG{n}{mean}\PYG{p}{(}\PYG{p}{(}\PYG{n}{y}\PYG{o}{\PYGZhy{}}\PYG{n}{y\PYGZus{}hat}\PYG{p}{)}\PYG{o}{*}\PYG{o}{*}\PYG{l+m+mi}{2}\PYG{p}{)}\PYG{o}{/}\PYG{n}{np}\PYG{o}{.}\PYG{n}{std}\PYG{p}{(}\PYG{n}{y}\PYG{p}{)}\PYG{o}{*}\PYG{o}{*}\PYG{l+m+mi}{2}
\end{sphinxVerbatim}

\end{sphinxuseclass}\end{sphinxVerbatimInput}

\end{sphinxuseclass}
\begin{sphinxuseclass}{cell}\begin{sphinxVerbatimInput}

\begin{sphinxuseclass}{cell_input}
\begin{sphinxVerbatim}[commandchars=\\\{\}]
\PYG{n+nb}{print}\PYG{p}{(}\PYG{l+s+sa}{f}\PYG{l+s+s2}{\PYGZdq{}}\PYG{l+s+s2}{MSEs: }\PYG{l+s+si}{\PYGZob{}}\PYG{n}{mse}\PYG{p}{(}\PYG{n}{y1}\PYG{p}{,}\PYG{+w}{ }\PYG{n}{model1}\PYG{o}{.}\PYG{n}{predict}\PYG{p}{(}\PYG{n}{x}\PYG{p}{)}\PYG{p}{)}\PYG{l+s+si}{\PYGZcb{}}\PYG{l+s+s2}{, }\PYG{l+s+si}{\PYGZob{}}\PYG{n}{mse}\PYG{p}{(}\PYG{n}{y2}\PYG{p}{,}\PYG{+w}{ }\PYG{n}{model2}\PYG{o}{.}\PYG{n}{predict}\PYG{p}{(}\PYG{n}{x}\PYG{p}{)}\PYG{p}{)}\PYG{l+s+si}{\PYGZcb{}}\PYG{l+s+s2}{\PYGZdq{}}\PYG{p}{)}
\PYG{n+nb}{print}\PYG{p}{(}\PYG{l+s+sa}{f}\PYG{l+s+s2}{\PYGZdq{}}\PYG{l+s+s2}{NMSEs: }\PYG{l+s+si}{\PYGZob{}}\PYG{n}{nmse}\PYG{p}{(}\PYG{n}{y1}\PYG{p}{,}\PYG{+w}{ }\PYG{n}{model1}\PYG{o}{.}\PYG{n}{predict}\PYG{p}{(}\PYG{n}{x}\PYG{p}{)}\PYG{p}{)}\PYG{l+s+si}{\PYGZcb{}}\PYG{l+s+s2}{, }\PYG{l+s+si}{\PYGZob{}}\PYG{n}{nmse}\PYG{p}{(}\PYG{n}{y2}\PYG{p}{,}\PYG{+w}{ }\PYG{n}{model2}\PYG{o}{.}\PYG{n}{predict}\PYG{p}{(}\PYG{n}{x}\PYG{p}{)}\PYG{p}{)}\PYG{l+s+si}{\PYGZcb{}}\PYG{l+s+s2}{\PYGZdq{}}\PYG{p}{)}
\end{sphinxVerbatim}

\end{sphinxuseclass}\end{sphinxVerbatimInput}
\begin{sphinxVerbatimOutput}

\begin{sphinxuseclass}{cell_output}
\begin{sphinxVerbatim}[commandchars=\\\{\}]
MSEs: 1.051874601439384, 102.696829717452
NMSEs: 0.16885724586364886, 0.963320410110144
\end{sphinxVerbatim}

\end{sphinxuseclass}\end{sphinxVerbatimOutput}

\end{sphinxuseclass}

\section{Feature selection}
\label{\detokenize{05:feature-selection}}
\sphinxAtStartPar
Feature selection and ensembles of models are two techniques which can be used to improve the accuracy of predictions.

\sphinxAtStartPar
Feature selection aims at reducing the dimensionality of the problem, and is useful when input variables contain redundant or irrelevant (noisy) information. Benefits are twofold: it decreases the training time by simplifying the problem, and it decreases the complexity of the predictive model. This in turn usually improves the prediction accuracy, since high\sphinxhyphen{}dimensionality makes predictive models more prone to overfitting, and estimates of parameters more variant.

\sphinxAtStartPar
There are three main approaches to feature selection:
\begin{itemize}
\item {} 
\sphinxAtStartPar
\sphinxstylestrong{Filter methods:}
These methods rely solely on the data and their intrinsic properties, without considering the impact of the selected features on the learning algorithm performance. For this reason, they are often used as preprocessing techniques.

\item {} 
\sphinxAtStartPar
\sphinxstylestrong{Wrapper methods:}
These methods assess subsets of variables according to their usefulness to a given predictor. The feature selection is performed using an evaluation function that includes the predictive performance of the considered learning algorithm as a selection criterion.

\item {} 
\sphinxAtStartPar
\sphinxstylestrong{Embedded methods:}
These methods are specific to given learning machines, and usually built\sphinxhyphen{}in in the learning procedure (e.g. random forest, regularization based techniques).

\end{itemize}

\sphinxAtStartPar
Ensembles of models consist in building several predictive models using resampled subsets of the original training set. The method works particularly well for predictive models with high variance (for example, decision trees or neural networks). The average prediction of the resulting models usually strongly decreases the variance component of the error, and as a consequence improves the prediction accuracy.

\sphinxAtStartPar
In this session, we will use a synthetic dataset, made of 30 continuous variables. 24 of them will be useless, and 6 will serve in the output. The output will be a function of the useful input variables.


\subsection{Data overview and preprocessing}
\label{\detokenize{05:data-overview-and-preprocessing}}
\begin{sphinxuseclass}{cell}\begin{sphinxVerbatimInput}

\begin{sphinxuseclass}{cell_input}
\begin{sphinxVerbatim}[commandchars=\\\{\}]
\PYG{k}{def}\PYG{+w}{ }\PYG{n+nf}{compute}\PYG{p}{(}\PYG{n}{data}\PYG{p}{)}\PYG{p}{:}
    \PYG{n}{data}\PYG{p}{[}\PYG{l+s+s1}{\PYGZsq{}}\PYG{l+s+s1}{y}\PYG{l+s+s1}{\PYGZsq{}}\PYG{p}{]} \PYG{o}{=} \PYG{n}{np}\PYG{o}{.}\PYG{n}{zeros}\PYG{p}{(}\PYG{n}{data}\PYG{o}{.}\PYG{n}{shape}\PYG{p}{[}\PYG{l+m+mi}{0}\PYG{p}{]}\PYG{p}{)}
    \PYG{k}{for} \PYG{n}{i} \PYG{o+ow}{in} \PYG{n+nb}{range}\PYG{p}{(}\PYG{l+m+mi}{6}\PYG{p}{)}\PYG{p}{:}
        \PYG{k}{if} \PYG{n}{i}\PYG{o}{\PYGZpc{}}\PYG{k}{3} == 0:
            \PYG{n}{data}\PYG{p}{[}\PYG{l+s+s1}{\PYGZsq{}}\PYG{l+s+s1}{y}\PYG{l+s+s1}{\PYGZsq{}}\PYG{p}{]} \PYG{o}{+}\PYG{o}{=} \PYG{n}{np}\PYG{o}{.}\PYG{n}{random}\PYG{o}{.}\PYG{n}{uniform}\PYG{p}{(}\PYG{o}{\PYGZhy{}}\PYG{l+m+mi}{1}\PYG{p}{,} \PYG{l+m+mi}{1}\PYG{p}{)} \PYG{o}{*} \PYG{n}{data}\PYG{p}{[}\PYG{l+s+sa}{f}\PYG{l+s+s1}{\PYGZsq{}}\PYG{l+s+s1}{x\PYGZus{}useful\PYGZus{}}\PYG{l+s+si}{\PYGZob{}}\PYG{n}{i}\PYG{o}{+}\PYG{l+m+mi}{1}\PYG{l+s+si}{\PYGZcb{}}\PYG{l+s+s1}{\PYGZsq{}}\PYG{p}{]} \PYG{o}{+} \PYG{n}{np}\PYG{o}{.}\PYG{n}{random}\PYG{o}{.}\PYG{n}{uniform}\PYG{p}{(}\PYG{o}{\PYGZhy{}}\PYG{l+m+mi}{5}\PYG{p}{,} \PYG{l+m+mi}{5}\PYG{p}{)}
        \PYG{k}{elif} \PYG{n}{i}\PYG{o}{\PYGZpc{}}\PYG{k}{3} == 1:
            \PYG{n}{data}\PYG{p}{[}\PYG{l+s+s1}{\PYGZsq{}}\PYG{l+s+s1}{y}\PYG{l+s+s1}{\PYGZsq{}}\PYG{p}{]} \PYG{o}{+}\PYG{o}{=} \PYG{n}{np}\PYG{o}{.}\PYG{n}{random}\PYG{o}{.}\PYG{n}{uniform}\PYG{p}{(}\PYG{o}{\PYGZhy{}}\PYG{l+m+mi}{1}\PYG{p}{,} \PYG{l+m+mi}{1}\PYG{p}{)} \PYG{o}{*} \PYG{n}{data}\PYG{p}{[}\PYG{l+s+sa}{f}\PYG{l+s+s1}{\PYGZsq{}}\PYG{l+s+s1}{x\PYGZus{}useful\PYGZus{}}\PYG{l+s+si}{\PYGZob{}}\PYG{n}{i}\PYG{o}{+}\PYG{l+m+mi}{1}\PYG{l+s+si}{\PYGZcb{}}\PYG{l+s+s1}{\PYGZsq{}}\PYG{p}{]}\PYG{o}{*}\PYG{o}{*}\PYG{l+m+mi}{2} \PYG{o}{+} \PYG{n}{np}\PYG{o}{.}\PYG{n}{random}\PYG{o}{.}\PYG{n}{uniform}\PYG{p}{(}\PYG{o}{\PYGZhy{}}\PYG{l+m+mi}{5}\PYG{p}{,} \PYG{l+m+mi}{5}\PYG{p}{)}
        \PYG{k}{elif} \PYG{n}{i}\PYG{o}{\PYGZpc{}}\PYG{k}{3} == 2:
            \PYG{n}{data}\PYG{p}{[}\PYG{l+s+s1}{\PYGZsq{}}\PYG{l+s+s1}{y}\PYG{l+s+s1}{\PYGZsq{}}\PYG{p}{]} \PYG{o}{+}\PYG{o}{=} \PYG{n}{np}\PYG{o}{.}\PYG{n}{random}\PYG{o}{.}\PYG{n}{uniform}\PYG{p}{(}\PYG{o}{\PYGZhy{}}\PYG{l+m+mi}{1}\PYG{p}{,} \PYG{l+m+mi}{1}\PYG{p}{)} \PYG{o}{*} \PYG{n}{np}\PYG{o}{.}\PYG{n}{sin}\PYG{p}{(}\PYG{n}{data}\PYG{p}{[}\PYG{l+s+sa}{f}\PYG{l+s+s1}{\PYGZsq{}}\PYG{l+s+s1}{x\PYGZus{}useful\PYGZus{}}\PYG{l+s+si}{\PYGZob{}}\PYG{n}{i}\PYG{o}{+}\PYG{l+m+mi}{1}\PYG{l+s+si}{\PYGZcb{}}\PYG{l+s+s1}{\PYGZsq{}}\PYG{p}{]}\PYG{p}{)} \PYG{o}{+} \PYG{n}{np}\PYG{o}{.}\PYG{n}{random}\PYG{o}{.}\PYG{n}{uniform}\PYG{p}{(}\PYG{o}{\PYGZhy{}}\PYG{l+m+mi}{5}\PYG{p}{,} \PYG{l+m+mi}{5}\PYG{p}{)}
    \PYG{k}{return} \PYG{n}{data}
\end{sphinxVerbatim}

\end{sphinxuseclass}\end{sphinxVerbatimInput}

\end{sphinxuseclass}
\begin{sphinxuseclass}{cell}\begin{sphinxVerbatimInput}

\begin{sphinxuseclass}{cell_input}
\begin{sphinxVerbatim}[commandchars=\\\{\}]
\PYG{n}{n\PYGZus{}x\PYGZus{}useful} \PYG{o}{=} \PYG{l+m+mi}{6}
\PYG{n}{n\PYGZus{}x\PYGZus{}useless} \PYG{o}{=} \PYG{l+m+mi}{24}
\PYG{n}{n\PYGZus{}samples} \PYG{o}{=} \PYG{l+m+mi}{1000}
\PYG{n}{d} \PYG{o}{=} \PYG{p}{\PYGZob{}}\PYG{p}{\PYGZcb{}}
\PYG{n}{np}\PYG{o}{.}\PYG{n}{random}\PYG{o}{.}\PYG{n}{seed}\PYG{p}{(}\PYG{l+m+mi}{0}\PYG{p}{)}

\PYG{k}{for} \PYG{n}{i} \PYG{o+ow}{in} \PYG{n+nb}{range}\PYG{p}{(}\PYG{n}{n\PYGZus{}x\PYGZus{}useful}\PYG{p}{)}\PYG{p}{:}
    \PYG{n}{m}\PYG{p}{,} \PYG{n}{s} \PYG{o}{=} \PYG{n}{np}\PYG{o}{.}\PYG{n}{random}\PYG{o}{.}\PYG{n}{uniform}\PYG{p}{(}\PYG{o}{\PYGZhy{}}\PYG{l+m+mi}{1}\PYG{p}{,} \PYG{l+m+mi}{1}\PYG{p}{)}\PYG{p}{,} \PYG{n}{np}\PYG{o}{.}\PYG{n}{random}\PYG{o}{.}\PYG{n}{uniform}\PYG{p}{(}\PYG{l+m+mf}{0.5}\PYG{p}{,} \PYG{l+m+mf}{2.5}\PYG{p}{)}
    \PYG{n}{d}\PYG{p}{[}\PYG{l+s+sa}{f}\PYG{l+s+s1}{\PYGZsq{}}\PYG{l+s+s1}{x\PYGZus{}useful\PYGZus{}}\PYG{l+s+si}{\PYGZob{}}\PYG{n}{i}\PYG{o}{+}\PYG{l+m+mi}{1}\PYG{l+s+si}{\PYGZcb{}}\PYG{l+s+s1}{\PYGZsq{}}\PYG{p}{]} \PYG{o}{=} \PYG{n}{np}\PYG{o}{.}\PYG{n}{random}\PYG{o}{.}\PYG{n}{normal}\PYG{p}{(}\PYG{n}{m}\PYG{p}{,} \PYG{n}{s}\PYG{p}{,} \PYG{n}{size}\PYG{o}{=}\PYG{p}{(}\PYG{n}{n\PYGZus{}samples}\PYG{p}{,}\PYG{p}{)}\PYG{p}{)}
\PYG{k}{for} \PYG{n}{i} \PYG{o+ow}{in} \PYG{n+nb}{range}\PYG{p}{(}\PYG{n}{n\PYGZus{}x\PYGZus{}useless}\PYG{p}{)}\PYG{p}{:}
    \PYG{n}{m}\PYG{p}{,} \PYG{n}{s} \PYG{o}{=} \PYG{n}{np}\PYG{o}{.}\PYG{n}{random}\PYG{o}{.}\PYG{n}{uniform}\PYG{p}{(}\PYG{o}{\PYGZhy{}}\PYG{l+m+mi}{1}\PYG{p}{,} \PYG{l+m+mi}{1}\PYG{p}{)}\PYG{p}{,} \PYG{n}{np}\PYG{o}{.}\PYG{n}{random}\PYG{o}{.}\PYG{n}{uniform}\PYG{p}{(}\PYG{l+m+mf}{0.5}\PYG{p}{,} \PYG{l+m+mf}{2.5}\PYG{p}{)}
    \PYG{n}{d}\PYG{p}{[}\PYG{l+s+sa}{f}\PYG{l+s+s1}{\PYGZsq{}}\PYG{l+s+s1}{x\PYGZus{}useless\PYGZus{}}\PYG{l+s+si}{\PYGZob{}}\PYG{n}{i}\PYG{o}{+}\PYG{l+m+mi}{1}\PYG{l+s+si}{\PYGZcb{}}\PYG{l+s+s1}{\PYGZsq{}}\PYG{p}{]} \PYG{o}{=} \PYG{n}{np}\PYG{o}{.}\PYG{n}{random}\PYG{o}{.}\PYG{n}{normal}\PYG{p}{(}\PYG{n}{m}\PYG{p}{,} \PYG{n}{s}\PYG{p}{,} \PYG{n}{size}\PYG{o}{=}\PYG{p}{(}\PYG{n}{n\PYGZus{}samples}\PYG{p}{,}\PYG{p}{)}\PYG{p}{)}

\PYG{n}{data} \PYG{o}{=} \PYG{n}{pd}\PYG{o}{.}\PYG{n}{DataFrame}\PYG{p}{(}\PYG{n}{d}\PYG{p}{)}

\PYG{n}{data} \PYG{o}{=} \PYG{n}{compute}\PYG{p}{(}\PYG{n}{data}\PYG{p}{)}
\end{sphinxVerbatim}

\end{sphinxuseclass}\end{sphinxVerbatimInput}

\end{sphinxuseclass}
\begin{sphinxuseclass}{cell}\begin{sphinxVerbatimInput}

\begin{sphinxuseclass}{cell_input}
\begin{sphinxVerbatim}[commandchars=\\\{\}]
\PYG{n}{data}\PYG{o}{.}\PYG{n}{describe}\PYG{p}{(}\PYG{n}{include}\PYG{o}{=}\PYG{l+s+s1}{\PYGZsq{}}\PYG{l+s+s1}{all}\PYG{l+s+s1}{\PYGZsq{}}\PYG{p}{)}
\end{sphinxVerbatim}

\end{sphinxuseclass}\end{sphinxVerbatimInput}
\begin{sphinxVerbatimOutput}

\begin{sphinxuseclass}{cell_output}
\begin{sphinxVerbatim}[commandchars=\\\{\}]
        x\PYGZus{}useful\PYGZus{}1   x\PYGZus{}useful\PYGZus{}2   x\PYGZus{}useful\PYGZus{}3   x\PYGZus{}useful\PYGZus{}4   x\PYGZus{}useful\PYGZus{}5  \PYGZbs{}
count  1000.000000  1000.000000  1000.000000  1000.000000  1000.000000   
mean      0.008883     0.787868    \PYGZhy{}0.862362    \PYGZhy{}0.220739     0.900991   
std       1.904117     2.362233     1.466222     0.881817     1.535395   
min      \PYGZhy{}5.782583    \PYGZhy{}6.520188    \PYGZhy{}5.587836    \PYGZhy{}3.444757    \PYGZhy{}3.751236   
25\PYGZpc{}      \PYGZhy{}1.250588    \PYGZhy{}0.839990    \PYGZhy{}1.910921    \PYGZhy{}0.821390    \PYGZhy{}0.111143   
50\PYGZpc{}      \PYGZhy{}0.014389     0.822298    \PYGZhy{}0.868373    \PYGZhy{}0.229933     0.883910   
75\PYGZpc{}       1.269272     2.284124     0.175776     0.372973     1.957939   
max       5.424227     8.481613     3.712036     3.085683     5.435128   

        x\PYGZus{}useful\PYGZus{}6  x\PYGZus{}useless\PYGZus{}1  x\PYGZus{}useless\PYGZus{}2  x\PYGZus{}useless\PYGZus{}3  x\PYGZus{}useless\PYGZus{}4  ...  \PYGZbs{}
count  1000.000000  1000.000000  1000.000000  1000.000000  1000.000000  ...   
mean      0.048531     0.848947     0.019645     0.655122    \PYGZhy{}0.189598  ...   
std       0.512844     1.362313     1.729953     1.665836     0.771923  ...   
min      \PYGZhy{}1.380297    \PYGZhy{}3.305645    \PYGZhy{}5.469208    \PYGZhy{}4.426479    \PYGZhy{}2.355184  ...   
25\PYGZpc{}      \PYGZhy{}0.293415    \PYGZhy{}0.091018    \PYGZhy{}1.090598    \PYGZhy{}0.432532    \PYGZhy{}0.738497  ...   
50\PYGZpc{}       0.050555     0.825699     0.014043     0.655699    \PYGZhy{}0.199494  ...   
75\PYGZpc{}       0.386908     1.786537     1.128925     1.744190     0.350852  ...   
max       1.539755     4.825280     5.912388     5.919659     2.026439  ...   

       x\PYGZus{}useless\PYGZus{}16  x\PYGZus{}useless\PYGZus{}17  x\PYGZus{}useless\PYGZus{}18  x\PYGZus{}useless\PYGZus{}19  x\PYGZus{}useless\PYGZus{}20  \PYGZbs{}
count   1000.000000   1000.000000   1000.000000   1000.000000   1000.000000   
mean      \PYGZhy{}0.421467      0.821257     \PYGZhy{}0.011411      0.441756     \PYGZhy{}0.958489   
std        0.922338      0.862039      2.287231      1.402679      1.452495   
min       \PYGZhy{}3.443265     \PYGZhy{}1.763982     \PYGZhy{}7.985982     \PYGZhy{}4.070035     \PYGZhy{}5.422766   
25\PYGZpc{}       \PYGZhy{}1.025653      0.250196     \PYGZhy{}1.464619     \PYGZhy{}0.527079     \PYGZhy{}1.967275   
50\PYGZpc{}       \PYGZhy{}0.423784      0.844540      0.021261      0.423711     \PYGZhy{}1.010335   
75\PYGZpc{}        0.211144      1.361641      1.522601      1.368873      0.002406   
max        3.010346      4.035696      7.689106      5.682900      3.504355   

       x\PYGZus{}useless\PYGZus{}21  x\PYGZus{}useless\PYGZus{}22  x\PYGZus{}useless\PYGZus{}23  x\PYGZus{}useless\PYGZus{}24            y  
count   1000.000000   1000.000000   1000.000000   1000.000000  1000.000000  
mean      \PYGZhy{}0.384471     \PYGZhy{}0.152625     \PYGZhy{}0.518022      0.051116    \PYGZhy{}7.379040  
std        0.787796      0.610938      2.501076      1.985946     1.890504  
min       \PYGZhy{}2.947481     \PYGZhy{}1.917751    \PYGZhy{}11.958636     \PYGZhy{}6.453296   \PYGZhy{}13.989287  
25\PYGZpc{}       \PYGZhy{}0.890462     \PYGZhy{}0.564771     \PYGZhy{}2.155428     \PYGZhy{}1.329528    \PYGZhy{}8.514752  
50\PYGZpc{}       \PYGZhy{}0.403122     \PYGZhy{}0.147460     \PYGZhy{}0.412912      0.054168    \PYGZhy{}7.414389  
75\PYGZpc{}        0.173736      0.252462      1.202683      1.373122    \PYGZhy{}6.287151  
max        2.146540      1.745356      8.160750      5.685927     2.202998  

[8 rows x 31 columns]
\end{sphinxVerbatim}

\end{sphinxuseclass}\end{sphinxVerbatimOutput}

\end{sphinxuseclass}

\subsection{Input and output variables}
\label{\detokenize{05:input-and-output-variables}}
\sphinxAtStartPar
The output variable (Y) is \sphinxcode{\sphinxupquote{y}}, and all other variables (X) are considered as inputs.

\begin{sphinxuseclass}{cell}\begin{sphinxVerbatimInput}

\begin{sphinxuseclass}{cell_input}
\begin{sphinxVerbatim}[commandchars=\\\{\}]
\PYG{n}{Y} \PYG{o}{=} \PYG{n}{data}\PYG{p}{[}\PYG{l+s+s1}{\PYGZsq{}}\PYG{l+s+s1}{y}\PYG{l+s+s1}{\PYGZsq{}}\PYG{p}{]}\PYG{o}{.}\PYG{n}{values}
\PYG{n}{X} \PYG{o}{=} \PYG{n}{data}\PYG{o}{.}\PYG{n}{drop}\PYG{p}{(}\PYG{n}{columns}\PYG{o}{=}\PYG{p}{[}\PYG{l+s+s1}{\PYGZsq{}}\PYG{l+s+s1}{y}\PYG{l+s+s1}{\PYGZsq{}}\PYG{p}{]}\PYG{p}{)}

\PYG{n}{N} \PYG{o}{=} \PYG{n}{X}\PYG{o}{.}\PYG{n}{shape}\PYG{p}{[}\PYG{l+m+mi}{0}\PYG{p}{]}
\PYG{n}{n} \PYG{o}{=} \PYG{n}{X}\PYG{o}{.}\PYG{n}{shape}\PYG{p}{[}\PYG{l+m+mi}{1}\PYG{p}{]}

\PYG{n}{plt}\PYG{o}{.}\PYG{n}{figure}\PYG{p}{(}\PYG{n}{figsize}\PYG{o}{=}\PYG{p}{(}\PYG{l+m+mi}{10}\PYG{p}{,} \PYG{l+m+mi}{6}\PYG{p}{)}\PYG{p}{)}
\PYG{n}{plt}\PYG{o}{.}\PYG{n}{hist}\PYG{p}{(}\PYG{n}{Y}\PYG{p}{,} \PYG{n}{bins}\PYG{o}{=}\PYG{l+m+mi}{50}\PYG{p}{)}
\PYG{n}{plt}\PYG{o}{.}\PYG{n}{title}\PYG{p}{(}\PYG{l+s+s2}{\PYGZdq{}}\PYG{l+s+s2}{Distribution of y}\PYG{l+s+s2}{\PYGZdq{}}\PYG{p}{)}
\PYG{n}{plt}\PYG{o}{.}\PYG{n}{grid}\PYG{p}{(}\PYG{l+s+s2}{\PYGZdq{}}\PYG{l+s+s2}{on}\PYG{l+s+s2}{\PYGZdq{}}\PYG{p}{)}
\PYG{n}{plt}\PYG{o}{.}\PYG{n}{show}\PYG{p}{(}\PYG{p}{)}

\PYG{n+nb}{print}\PYG{p}{(}\PYG{l+s+s2}{\PYGZdq{}}\PYG{l+s+s2}{Mean of Y:}\PYG{l+s+s2}{\PYGZdq{}}\PYG{p}{,} \PYG{n}{np}\PYG{o}{.}\PYG{n}{mean}\PYG{p}{(}\PYG{n}{Y}\PYG{p}{)}\PYG{p}{)}
\PYG{n+nb}{print}\PYG{p}{(}\PYG{l+s+s2}{\PYGZdq{}}\PYG{l+s+s2}{Variance of Y:}\PYG{l+s+s2}{\PYGZdq{}}\PYG{p}{,} \PYG{n}{np}\PYG{o}{.}\PYG{n}{var}\PYG{p}{(}\PYG{n}{Y}\PYG{p}{)}\PYG{p}{)}
\end{sphinxVerbatim}

\end{sphinxuseclass}\end{sphinxVerbatimInput}
\begin{sphinxVerbatimOutput}

\begin{sphinxuseclass}{cell_output}
\noindent\sphinxincludegraphics{{46179c5858e03263e6978f21dcad1199965f6554b3d0abfb06e31966e92726d6}.png}

\begin{sphinxVerbatim}[commandchars=\\\{\}]
Mean of Y: \PYGZhy{}7.379040070934719
Variance of Y: 3.5704323990926374
\end{sphinxVerbatim}

\end{sphinxuseclass}\end{sphinxVerbatimOutput}

\end{sphinxuseclass}

\subsection{Modelling with linear and decision tree models}
\label{\detokenize{05:modelling-with-linear-and-decision-tree-models}}

\subsubsection{Linear model}
\label{\detokenize{05:linear-model}}\begin{itemize}
\item {} 
\sphinxAtStartPar
Create a linear model \(h(x)\) for predicting the IMDB score on the basis of the input variables, and compute its empirical (or training) mean square error.

\end{itemize}

\begin{sphinxuseclass}{cell}\begin{sphinxVerbatimInput}

\begin{sphinxuseclass}{cell_input}
\begin{sphinxVerbatim}[commandchars=\\\{\}]
\PYG{k+kn}{from}\PYG{+w}{ }\PYG{n+nn}{sklearn}\PYG{n+nn}{.}\PYG{n+nn}{linear\PYGZus{}model}\PYG{+w}{ }\PYG{k+kn}{import} \PYG{n}{LinearRegression}

\PYG{n}{model} \PYG{o}{=} \PYG{n}{LinearRegression}\PYG{p}{(}\PYG{p}{)}
\PYG{n}{model}\PYG{o}{.}\PYG{n}{fit}\PYG{p}{(}\PYG{n}{X}\PYG{p}{,} \PYG{n}{Y}\PYG{p}{)}
\PYG{n}{Y\PYGZus{}hat} \PYG{o}{=} \PYG{n}{model}\PYG{o}{.}\PYG{n}{predict}\PYG{p}{(}\PYG{n}{X}\PYG{p}{)}

\PYG{n}{empirical\PYGZus{}error} \PYG{o}{=} \PYG{n}{nmse}\PYG{p}{(}\PYG{n}{Y}\PYG{p}{,} \PYG{n}{Y\PYGZus{}hat}\PYG{p}{)}
\PYG{n+nb}{print}\PYG{p}{(}\PYG{l+s+sa}{f}\PYG{l+s+s2}{\PYGZdq{}}\PYG{l+s+s2}{Empirical error: }\PYG{l+s+si}{\PYGZob{}}\PYG{n}{empirical\PYGZus{}error}\PYG{l+s+si}{:}\PYG{l+s+s2}{.5f}\PYG{l+s+si}{\PYGZcb{}}\PYG{l+s+s2}{\PYGZdq{}}\PYG{p}{)}
\end{sphinxVerbatim}

\end{sphinxuseclass}\end{sphinxVerbatimInput}
\begin{sphinxVerbatimOutput}

\begin{sphinxuseclass}{cell_output}
\begin{sphinxVerbatim}[commandchars=\\\{\}]
Empirical error: 0.55375
\end{sphinxVerbatim}

\end{sphinxuseclass}\end{sphinxVerbatimOutput}

\end{sphinxuseclass}\begin{itemize}
\item {} 
\sphinxAtStartPar
Which input variables are statistically correlated with the output?

\end{itemize}

\sphinxAtStartPar
In Python, we can check the coefficients of the linear model.

\begin{sphinxuseclass}{cell}\begin{sphinxVerbatimInput}

\begin{sphinxuseclass}{cell_input}
\begin{sphinxVerbatim}[commandchars=\\\{\}]
\PYG{n}{coefs} \PYG{o}{=} \PYG{n}{pd}\PYG{o}{.}\PYG{n}{Series}\PYG{p}{(}\PYG{n}{model}\PYG{o}{.}\PYG{n}{coef\PYGZus{}}\PYG{p}{,} \PYG{n}{index}\PYG{o}{=}\PYG{n}{X}\PYG{o}{.}\PYG{n}{columns}\PYG{p}{)}
\PYG{n}{plt}\PYG{o}{.}\PYG{n}{figure}\PYG{p}{(}\PYG{n}{figsize}\PYG{o}{=}\PYG{p}{(}\PYG{l+m+mi}{15}\PYG{p}{,} \PYG{l+m+mi}{6}\PYG{p}{)}\PYG{p}{)}
\PYG{n}{sns}\PYG{o}{.}\PYG{n}{barplot}\PYG{p}{(}\PYG{n}{coefs}\PYG{p}{)}
\PYG{n}{plt}\PYG{o}{.}\PYG{n}{title}\PYG{p}{(}\PYG{l+s+s2}{\PYGZdq{}}\PYG{l+s+s2}{Coefficient of each variable in the linear model}\PYG{l+s+s2}{\PYGZdq{}}\PYG{p}{)}
\PYG{n}{plt}\PYG{o}{.}\PYG{n}{xticks}\PYG{p}{(}\PYG{n}{rotation}\PYG{o}{=}\PYG{l+m+mi}{45}\PYG{p}{,} \PYG{n}{ha}\PYG{o}{=}\PYG{l+s+s1}{\PYGZsq{}}\PYG{l+s+s1}{right}\PYG{l+s+s1}{\PYGZsq{}}\PYG{p}{)}
\PYG{n}{plt}\PYG{o}{.}\PYG{n}{grid}\PYG{p}{(}\PYG{l+s+s2}{\PYGZdq{}}\PYG{l+s+s2}{on}\PYG{l+s+s2}{\PYGZdq{}}\PYG{p}{)}
\PYG{n}{plt}\PYG{o}{.}\PYG{n}{show}\PYG{p}{(}\PYG{p}{)}
\end{sphinxVerbatim}

\end{sphinxuseclass}\end{sphinxVerbatimInput}
\begin{sphinxVerbatimOutput}

\begin{sphinxuseclass}{cell_output}
\noindent\sphinxincludegraphics{{5ae216c4470e2ab3bb3e30f81f0d9b747e4e67dea7bb641e67c5c56baa838322}.png}

\end{sphinxuseclass}\end{sphinxVerbatimOutput}

\end{sphinxuseclass}\begin{itemize}
\item {} 
\sphinxAtStartPar
Compute the validation error with a 10\sphinxhyphen{}fold cross\sphinxhyphen{}validation

\end{itemize}

\begin{sphinxuseclass}{cell}\begin{sphinxVerbatimInput}

\begin{sphinxuseclass}{cell_input}
\begin{sphinxVerbatim}[commandchars=\\\{\}]
\PYG{k+kn}{from}\PYG{+w}{ }\PYG{n+nn}{sklearn}\PYG{n+nn}{.}\PYG{n+nn}{model\PYGZus{}selection}\PYG{+w}{ }\PYG{k+kn}{import} \PYG{n}{KFold}

\PYG{n}{kf} \PYG{o}{=} \PYG{n}{KFold}\PYG{p}{(}\PYG{n}{n\PYGZus{}splits}\PYG{o}{=}\PYG{l+m+mi}{10}\PYG{p}{,} \PYG{n}{shuffle}\PYG{o}{=}\PYG{k+kc}{True}\PYG{p}{,} \PYG{n}{random\PYGZus{}state}\PYG{o}{=}\PYG{l+m+mi}{3}\PYG{p}{)}
\PYG{n}{CV\PYGZus{}err\PYGZus{}lm\PYGZus{}single\PYGZus{}model} \PYG{o}{=} \PYG{p}{[}\PYG{p}{]}

\PYG{k}{for} \PYG{n}{train\PYGZus{}index}\PYG{p}{,} \PYG{n}{test\PYGZus{}index} \PYG{o+ow}{in} \PYG{n}{kf}\PYG{o}{.}\PYG{n}{split}\PYG{p}{(}\PYG{n}{X}\PYG{p}{)}\PYG{p}{:}
    \PYG{n}{X\PYGZus{}tr}\PYG{p}{,} \PYG{n}{X\PYGZus{}ts} \PYG{o}{=} \PYG{n}{X}\PYG{o}{.}\PYG{n}{iloc}\PYG{p}{[}\PYG{n}{train\PYGZus{}index}\PYG{p}{]}\PYG{p}{,} \PYG{n}{X}\PYG{o}{.}\PYG{n}{iloc}\PYG{p}{[}\PYG{n}{test\PYGZus{}index}\PYG{p}{]}
    \PYG{n}{Y\PYGZus{}tr}\PYG{p}{,} \PYG{n}{Y\PYGZus{}ts} \PYG{o}{=} \PYG{n}{Y}\PYG{p}{[}\PYG{n}{train\PYGZus{}index}\PYG{p}{]}\PYG{p}{,} \PYG{n}{Y}\PYG{p}{[}\PYG{n}{test\PYGZus{}index}\PYG{p}{]}
    
    \PYG{n}{model} \PYG{o}{=} \PYG{n}{LinearRegression}\PYG{p}{(}\PYG{p}{)}
    \PYG{n}{model}\PYG{o}{.}\PYG{n}{fit}\PYG{p}{(}\PYG{n}{X\PYGZus{}tr}\PYG{p}{,} \PYG{n}{Y\PYGZus{}tr}\PYG{p}{)}
    \PYG{n}{Y\PYGZus{}hat\PYGZus{}ts} \PYG{o}{=} \PYG{n}{model}\PYG{o}{.}\PYG{n}{predict}\PYG{p}{(}\PYG{n}{X\PYGZus{}ts}\PYG{p}{)}
    \PYG{n}{CV\PYGZus{}err\PYGZus{}lm\PYGZus{}single\PYGZus{}model}\PYG{o}{.}\PYG{n}{append}\PYG{p}{(}\PYG{n}{nmse}\PYG{p}{(}\PYG{n}{Y\PYGZus{}ts}\PYG{p}{,} \PYG{n}{Y\PYGZus{}hat\PYGZus{}ts}\PYG{p}{)}\PYG{p}{)}

\PYG{n+nb}{print}\PYG{p}{(}\PYG{l+s+sa}{f}\PYG{l+s+s2}{\PYGZdq{}}\PYG{l+s+s2}{CV error: }\PYG{l+s+si}{\PYGZob{}}\PYG{n}{np}\PYG{o}{.}\PYG{n}{mean}\PYG{p}{(}\PYG{n}{CV\PYGZus{}err\PYGZus{}lm\PYGZus{}single\PYGZus{}model}\PYG{p}{)}\PYG{l+s+si}{:}\PYG{l+s+s2}{.5f}\PYG{l+s+si}{\PYGZcb{}}\PYG{l+s+s2}{, std dev: }\PYG{l+s+si}{\PYGZob{}}\PYG{n}{np}\PYG{o}{.}\PYG{n}{std}\PYG{p}{(}\PYG{n}{CV\PYGZus{}err\PYGZus{}lm\PYGZus{}single\PYGZus{}model}\PYG{p}{)}\PYG{l+s+si}{:}\PYG{l+s+s2}{.5f}\PYG{l+s+si}{\PYGZcb{}}\PYG{l+s+s2}{\PYGZdq{}}\PYG{p}{)}
\end{sphinxVerbatim}

\end{sphinxuseclass}\end{sphinxVerbatimInput}
\begin{sphinxVerbatimOutput}

\begin{sphinxuseclass}{cell_output}
\begin{sphinxVerbatim}[commandchars=\\\{\}]
CV error: 0.61060, std dev: 0.08403
\end{sphinxVerbatim}

\end{sphinxuseclass}\end{sphinxVerbatimOutput}

\end{sphinxuseclass}

\subsubsection{Decision tree}
\label{\detokenize{05:decision-tree}}\begin{itemize}
\item {} 
\sphinxAtStartPar
Modify the previous code to compute the empirical error using a decision tree model. Use sklearn’s \sphinxcode{\sphinxupquote{DecisionTreeRegressor}}.

\end{itemize}

\begin{sphinxuseclass}{cell}\begin{sphinxVerbatimInput}

\begin{sphinxuseclass}{cell_input}
\begin{sphinxVerbatim}[commandchars=\\\{\}]
\PYG{k+kn}{from}\PYG{+w}{ }\PYG{n+nn}{sklearn}\PYG{n+nn}{.}\PYG{n+nn}{tree}\PYG{+w}{ }\PYG{k+kn}{import} \PYG{n}{DecisionTreeRegressor}

\PYG{n}{model} \PYG{o}{=} \PYG{n}{DecisionTreeRegressor}\PYG{p}{(}\PYG{n}{random\PYGZus{}state}\PYG{o}{=}\PYG{l+m+mi}{3}\PYG{p}{)}
\PYG{n}{model}\PYG{o}{.}\PYG{n}{fit}\PYG{p}{(}\PYG{n}{X}\PYG{p}{,} \PYG{n}{Y}\PYG{p}{)}
\PYG{n}{Y\PYGZus{}hat} \PYG{o}{=} \PYG{n}{model}\PYG{o}{.}\PYG{n}{predict}\PYG{p}{(}\PYG{n}{X}\PYG{p}{)}

\PYG{n}{empirical\PYGZus{}error} \PYG{o}{=} \PYG{n}{nmse}\PYG{p}{(}\PYG{n}{Y}\PYG{p}{,} \PYG{n}{Y\PYGZus{}hat}\PYG{p}{)}
\PYG{n+nb}{print}\PYG{p}{(}\PYG{l+s+sa}{f}\PYG{l+s+s2}{\PYGZdq{}}\PYG{l+s+s2}{Empirical error: }\PYG{l+s+si}{\PYGZob{}}\PYG{n}{empirical\PYGZus{}error}\PYG{l+s+si}{:}\PYG{l+s+s2}{.5f}\PYG{l+s+si}{\PYGZcb{}}\PYG{l+s+s2}{\PYGZdq{}}\PYG{p}{)}
\end{sphinxVerbatim}

\end{sphinxuseclass}\end{sphinxVerbatimInput}
\begin{sphinxVerbatimOutput}

\begin{sphinxuseclass}{cell_output}
\begin{sphinxVerbatim}[commandchars=\\\{\}]
Empirical error: 0.00000
\end{sphinxVerbatim}

\end{sphinxuseclass}\end{sphinxVerbatimOutput}

\end{sphinxuseclass}\begin{itemize}
\item {} 
\sphinxAtStartPar
Plot the resulting tree is more complicated in Python due to size, but we can just visualize its structure or get the depth.

\end{itemize}

\begin{sphinxuseclass}{cell}\begin{sphinxVerbatimInput}

\begin{sphinxuseclass}{cell_input}
\begin{sphinxVerbatim}[commandchars=\\\{\}]
\PYG{n+nb}{print}\PYG{p}{(}\PYG{l+s+sa}{f}\PYG{l+s+s2}{\PYGZdq{}}\PYG{l+s+s2}{Depth: }\PYG{l+s+si}{\PYGZob{}}\PYG{n}{model}\PYG{o}{.}\PYG{n}{get\PYGZus{}depth}\PYG{p}{(}\PYG{p}{)}\PYG{l+s+si}{\PYGZcb{}}\PYG{l+s+s2}{, number of leaves: }\PYG{l+s+si}{\PYGZob{}}\PYG{n}{model}\PYG{o}{.}\PYG{n}{get\PYGZus{}n\PYGZus{}leaves}\PYG{p}{(}\PYG{p}{)}\PYG{l+s+si}{\PYGZcb{}}\PYG{l+s+s2}{\PYGZdq{}}\PYG{p}{)}
\end{sphinxVerbatim}

\end{sphinxuseclass}\end{sphinxVerbatimInput}
\begin{sphinxVerbatimOutput}

\begin{sphinxuseclass}{cell_output}
\begin{sphinxVerbatim}[commandchars=\\\{\}]
Depth: 22, number of leaves: 1000
\end{sphinxVerbatim}

\end{sphinxuseclass}\end{sphinxVerbatimOutput}

\end{sphinxuseclass}\begin{itemize}
\item {} 
\sphinxAtStartPar
What is the 10\sphinxhyphen{}fold cross\sphinxhyphen{}validation error using a decision tree model?

\end{itemize}

\begin{sphinxuseclass}{cell}\begin{sphinxVerbatimInput}

\begin{sphinxuseclass}{cell_input}
\begin{sphinxVerbatim}[commandchars=\\\{\}]
\PYG{n}{CV\PYGZus{}err\PYGZus{}rpart\PYGZus{}single\PYGZus{}model} \PYG{o}{=} \PYG{p}{[}\PYG{p}{]}

\PYG{k}{for} \PYG{n}{train\PYGZus{}index}\PYG{p}{,} \PYG{n}{test\PYGZus{}index} \PYG{o+ow}{in} \PYG{n}{kf}\PYG{o}{.}\PYG{n}{split}\PYG{p}{(}\PYG{n}{X}\PYG{p}{)}\PYG{p}{:}
    \PYG{n}{X\PYGZus{}tr}\PYG{p}{,} \PYG{n}{X\PYGZus{}ts} \PYG{o}{=} \PYG{n}{X}\PYG{o}{.}\PYG{n}{iloc}\PYG{p}{[}\PYG{n}{train\PYGZus{}index}\PYG{p}{]}\PYG{p}{,} \PYG{n}{X}\PYG{o}{.}\PYG{n}{iloc}\PYG{p}{[}\PYG{n}{test\PYGZus{}index}\PYG{p}{]}
    \PYG{n}{Y\PYGZus{}tr}\PYG{p}{,} \PYG{n}{Y\PYGZus{}ts} \PYG{o}{=} \PYG{n}{Y}\PYG{p}{[}\PYG{n}{train\PYGZus{}index}\PYG{p}{]}\PYG{p}{,} \PYG{n}{Y}\PYG{p}{[}\PYG{n}{test\PYGZus{}index}\PYG{p}{]}
    
    \PYG{n}{model} \PYG{o}{=} \PYG{n}{DecisionTreeRegressor}\PYG{p}{(}\PYG{n}{random\PYGZus{}state}\PYG{o}{=}\PYG{l+m+mi}{3}\PYG{p}{)}
    \PYG{n}{model}\PYG{o}{.}\PYG{n}{fit}\PYG{p}{(}\PYG{n}{X\PYGZus{}tr}\PYG{p}{,} \PYG{n}{Y\PYGZus{}tr}\PYG{p}{)}
    \PYG{n}{Y\PYGZus{}hat\PYGZus{}ts} \PYG{o}{=} \PYG{n}{model}\PYG{o}{.}\PYG{n}{predict}\PYG{p}{(}\PYG{n}{X\PYGZus{}ts}\PYG{p}{)}
    \PYG{n}{CV\PYGZus{}err\PYGZus{}rpart\PYGZus{}single\PYGZus{}model}\PYG{o}{.}\PYG{n}{append}\PYG{p}{(}\PYG{n}{nmse}\PYG{p}{(}\PYG{n}{Y\PYGZus{}ts}\PYG{p}{,} \PYG{n}{Y\PYGZus{}hat\PYGZus{}ts}\PYG{p}{)}\PYG{p}{)}

\PYG{n+nb}{print}\PYG{p}{(}\PYG{l+s+sa}{f}\PYG{l+s+s2}{\PYGZdq{}}\PYG{l+s+s2}{CV error: }\PYG{l+s+si}{\PYGZob{}}\PYG{n}{np}\PYG{o}{.}\PYG{n}{mean}\PYG{p}{(}\PYG{n}{CV\PYGZus{}err\PYGZus{}rpart\PYGZus{}single\PYGZus{}model}\PYG{p}{)}\PYG{l+s+si}{:}\PYG{l+s+s2}{.5f}\PYG{l+s+si}{\PYGZcb{}}\PYG{l+s+s2}{, std dev: }\PYG{l+s+si}{\PYGZob{}}\PYG{n}{np}\PYG{o}{.}\PYG{n}{std}\PYG{p}{(}\PYG{n}{CV\PYGZus{}err\PYGZus{}rpart\PYGZus{}single\PYGZus{}model}\PYG{p}{)}\PYG{l+s+si}{:}\PYG{l+s+s2}{.5f}\PYG{l+s+si}{\PYGZcb{}}\PYG{l+s+s2}{\PYGZdq{}}\PYG{p}{)}
\end{sphinxVerbatim}

\end{sphinxuseclass}\end{sphinxVerbatimInput}
\begin{sphinxVerbatimOutput}

\begin{sphinxuseclass}{cell_output}
\begin{sphinxVerbatim}[commandchars=\\\{\}]
CV error: 0.33924, std dev: 0.08885
\end{sphinxVerbatim}

\end{sphinxuseclass}\end{sphinxVerbatimOutput}

\end{sphinxuseclass}
\sphinxAtStartPar
Why is the result so different using the 10\sphinxhyphen{}fold cross\sphinxhyphen{}validation?


\subsubsection{Ridge regression with LOO\sphinxhyphen{}CV}
\label{\detokenize{05:ridge-regression-with-loo-cv}}
\sphinxAtStartPar
Recall: Ridge regression looks like a linear regression but includes a penalty term for large coefficients:
\begin{equation*}
\begin{split}
  \min_{\beta} \;\; \| Y - X \beta \|^2 + \lambda \| \beta \|^2.
\end{split}
\end{equation*}
\sphinxAtStartPar
Ridge also has a known closed\sphinxhyphen{}form:
\begin{equation*}
\begin{split}
  \hat{\beta}_{\mathrm{ridge}} = (X^\top X + \lambda I)^{-1} X^\top Y.
\end{split}
\end{equation*}\begin{itemize}
\item {} 
\sphinxAtStartPar
Leave\sphinxhyphen{}one out cross\sphinxhyphen{}validation is a method that consists in evaluating a given model by training it on the entire training set, except for one sample used for validation. By performing this validation, eliminating each of the samples at a time, we get an estimation of the quality of the model.

\item {} 
\sphinxAtStartPar
Then, the coefficients can give an idea of the importance of each of the features in the prediction task.

\end{itemize}

\sphinxAtStartPar
Here, we limit the dataset size to \(50\) items, to avoid too long computations.

\begin{sphinxuseclass}{cell}\begin{sphinxVerbatimInput}

\begin{sphinxuseclass}{cell_input}
\begin{sphinxVerbatim}[commandchars=\\\{\}]
\PYG{n}{X\PYGZus{}np} \PYG{o}{=} \PYG{n}{X}\PYG{o}{.}\PYG{n}{to\PYGZus{}numpy}\PYG{p}{(}\PYG{p}{)}\PYG{p}{[}\PYG{p}{:}\PYG{l+m+mi}{50}\PYG{p}{]}

\PYG{n}{N} \PYG{o}{=} \PYG{n}{X\PYGZus{}np}\PYG{o}{.}\PYG{n}{shape}\PYG{p}{[}\PYG{l+m+mi}{0}\PYG{p}{]}
\PYG{n}{lambdas} \PYG{o}{=} \PYG{n}{np}\PYG{o}{.}\PYG{n}{linspace}\PYG{p}{(}\PYG{l+m+mf}{0.01}\PYG{p}{,} \PYG{l+m+mf}{0.5}\PYG{p}{,} \PYG{l+m+mi}{10}\PYG{p}{)}
\PYG{n}{E} \PYG{o}{=} \PYG{n}{np}\PYG{o}{.}\PYG{n}{empty}\PYG{p}{(}\PYG{p}{(}\PYG{n}{N}\PYG{p}{,} \PYG{n+nb}{len}\PYG{p}{(}\PYG{n}{lambdas}\PYG{p}{)}\PYG{p}{)}\PYG{p}{)}  
\PYG{n}{E}\PYG{p}{[}\PYG{p}{:}\PYG{p}{]} \PYG{o}{=} \PYG{n}{np}\PYG{o}{.}\PYG{n}{nan}

\PYG{k}{for} \PYG{n}{i} \PYG{o+ow}{in} \PYG{n}{tqdm}\PYG{p}{(}\PYG{n+nb}{range}\PYG{p}{(}\PYG{n}{N}\PYG{p}{)}\PYG{p}{)}\PYG{p}{:}
    \PYG{n}{Xtr} \PYG{o}{=} \PYG{n}{np}\PYG{o}{.}\PYG{n}{column\PYGZus{}stack}\PYG{p}{(}\PYG{p}{(}\PYG{n}{np}\PYG{o}{.}\PYG{n}{ones}\PYG{p}{(}\PYG{n}{N} \PYG{o}{\PYGZhy{}} \PYG{l+m+mi}{1}\PYG{p}{)}\PYG{p}{,} \PYG{n}{np}\PYG{o}{.}\PYG{n}{delete}\PYG{p}{(}\PYG{n}{X\PYGZus{}np}\PYG{p}{,} \PYG{n}{i}\PYG{p}{,} \PYG{n}{axis}\PYG{o}{=}\PYG{l+m+mi}{0}\PYG{p}{)}\PYG{p}{)}\PYG{p}{)}
    \PYG{n}{Ytr} \PYG{o}{=} \PYG{n}{np}\PYG{o}{.}\PYG{n}{delete}\PYG{p}{(}\PYG{n}{Y}\PYG{p}{[}\PYG{p}{:}\PYG{l+m+mi}{50}\PYG{p}{]}\PYG{p}{,} \PYG{n}{i}\PYG{p}{)}
    
    \PYG{n}{Xts} \PYG{o}{=} \PYG{n}{np}\PYG{o}{.}\PYG{n}{concatenate}\PYG{p}{(}\PYG{p}{(}\PYG{p}{[}\PYG{l+m+mi}{1}\PYG{p}{]}\PYG{p}{,} \PYG{n}{X\PYGZus{}np}\PYG{p}{[}\PYG{n}{i}\PYG{p}{,}\PYG{p}{]}\PYG{p}{)}\PYG{p}{)}
    
    \PYG{n}{cnt} \PYG{o}{=} \PYG{l+m+mi}{0}
    \PYG{k}{for} \PYG{n}{l} \PYG{o+ow}{in} \PYG{n}{lambdas}\PYG{p}{:}
        \PYG{n}{A} \PYG{o}{=} \PYG{n}{Xtr}\PYG{o}{.}\PYG{n}{T} \PYG{o}{@} \PYG{n}{Xtr} \PYG{o}{+} \PYG{n}{l} \PYG{o}{*} \PYG{n}{np}\PYG{o}{.}\PYG{n}{eye}\PYG{p}{(}\PYG{n}{n} \PYG{o}{+} \PYG{l+m+mi}{1}\PYG{p}{)}
        \PYG{n}{b} \PYG{o}{=} \PYG{n}{Xtr}\PYG{o}{.}\PYG{n}{T} \PYG{o}{@} \PYG{n}{Ytr}
        \PYG{n}{betahat} \PYG{o}{=} \PYG{n}{np}\PYG{o}{.}\PYG{n}{linalg}\PYG{o}{.}\PYG{n}{inv}\PYG{p}{(}\PYG{n}{A}\PYG{p}{)} \PYG{o}{@} \PYG{n}{b}
        \PYG{n}{Yhati} \PYG{o}{=} \PYG{n}{Xts} \PYG{o}{@} \PYG{n}{betahat}
        \PYG{n}{E}\PYG{p}{[}\PYG{n}{i}\PYG{p}{,} \PYG{n}{cnt}\PYG{p}{]} \PYG{o}{=} \PYG{p}{(}\PYG{p}{(}\PYG{n}{Y}\PYG{p}{[}\PYG{n}{i}\PYG{p}{]} \PYG{o}{\PYGZhy{}} \PYG{n}{Yhati}\PYG{p}{)} \PYG{o}{*}\PYG{o}{*} \PYG{l+m+mi}{2}\PYG{p}{)}
        \PYG{n}{cnt} \PYG{o}{+}\PYG{o}{=} \PYG{l+m+mi}{1}

\PYG{n}{mseloo} \PYG{o}{=} \PYG{n}{np}\PYG{o}{.}\PYG{n}{mean}\PYG{p}{(}\PYG{n}{E}\PYG{p}{,} \PYG{n}{axis}\PYG{o}{=}\PYG{l+m+mi}{0}\PYG{p}{)}\PYG{o}{/}\PYG{n}{np}\PYG{o}{.}\PYG{n}{std}\PYG{p}{(}\PYG{n}{Y}\PYG{p}{)}\PYG{o}{*}\PYG{o}{*}\PYG{l+m+mi}{2}

\PYG{n}{bestLambda} \PYG{o}{=} \PYG{n}{lambdas}\PYG{p}{[}\PYG{n}{np}\PYG{o}{.}\PYG{n}{argmin}\PYG{p}{(}\PYG{n}{np}\PYG{o}{.}\PYG{n}{mean}\PYG{p}{(}\PYG{n}{E}\PYG{p}{,} \PYG{n}{axis}\PYG{o}{=}\PYG{l+m+mi}{0}\PYG{p}{)}\PYG{p}{)}\PYG{p}{]}
\PYG{n+nb}{print}\PYG{p}{(}\PYG{l+s+sa}{f}\PYG{l+s+s2}{\PYGZdq{}}\PYG{l+s+s2}{Best lambda: }\PYG{l+s+si}{\PYGZob{}}\PYG{n}{bestLambda}\PYG{l+s+si}{\PYGZcb{}}\PYG{l+s+s2}{\PYGZdq{}}\PYG{p}{)}

\PYG{n}{XX} \PYG{o}{=} \PYG{n}{np}\PYG{o}{.}\PYG{n}{column\PYGZus{}stack}\PYG{p}{(}\PYG{p}{(}\PYG{n}{np}\PYG{o}{.}\PYG{n}{ones}\PYG{p}{(}\PYG{n}{N}\PYG{p}{)}\PYG{p}{,} \PYG{n}{X\PYGZus{}np}\PYG{p}{)}\PYG{p}{)}
\PYG{n}{A\PYGZus{}final} \PYG{o}{=} \PYG{n}{np}\PYG{o}{.}\PYG{n}{dot}\PYG{p}{(}\PYG{n}{XX}\PYG{o}{.}\PYG{n}{T}\PYG{p}{,} \PYG{n}{XX}\PYG{p}{)} \PYG{o}{+} \PYG{n}{bestLambda} \PYG{o}{*} \PYG{n}{np}\PYG{o}{.}\PYG{n}{eye}\PYG{p}{(}\PYG{n}{n} \PYG{o}{+} \PYG{l+m+mi}{1}\PYG{p}{)}
\PYG{n}{b\PYGZus{}final} \PYG{o}{=} \PYG{n}{np}\PYG{o}{.}\PYG{n}{dot}\PYG{p}{(}\PYG{n}{XX}\PYG{o}{.}\PYG{n}{T}\PYG{p}{,} \PYG{n}{Y}\PYG{p}{[}\PYG{p}{:}\PYG{l+m+mi}{50}\PYG{p}{]}\PYG{p}{)}
\PYG{n}{betahat} \PYG{o}{=} \PYG{n}{np}\PYG{o}{.}\PYG{n}{linalg}\PYG{o}{.}\PYG{n}{inv}\PYG{p}{(}\PYG{n}{A\PYGZus{}final}\PYG{p}{)} \PYG{o}{@} \PYG{n}{XX}\PYG{o}{.}\PYG{n}{T} \PYG{o}{@} \PYG{n}{Y}\PYG{p}{[}\PYG{p}{:}\PYG{l+m+mi}{50}\PYG{p}{]}

\PYG{n}{abs\PYGZus{}betahat} \PYG{o}{=} \PYG{n}{np}\PYG{o}{.}\PYG{n}{abs}\PYG{p}{(}\PYG{n}{betahat}\PYG{p}{)}
\PYG{n}{sorted\PYGZus{}indices} \PYG{o}{=} \PYG{n}{np}\PYG{o}{.}\PYG{n}{argsort}\PYG{p}{(}\PYG{o}{\PYGZhy{}}\PYG{n}{abs\PYGZus{}betahat}\PYG{p}{)}
\PYG{n}{top4} \PYG{o}{=} \PYG{n}{sorted\PYGZus{}indices}\PYG{p}{[}\PYG{p}{:}\PYG{l+m+mi}{4}\PYG{p}{]} 
\PYG{n+nb}{print}\PYG{p}{(}\PYG{l+s+sa}{f}\PYG{l+s+s2}{\PYGZdq{}}\PYG{l+s+s2}{Greatests coefficients indexes: }\PYG{l+s+si}{\PYGZob{}}\PYG{n}{top4}\PYG{l+s+si}{\PYGZcb{}}\PYG{l+s+s2}{\PYGZdq{}}\PYG{p}{)}
\PYG{n+nb}{print}\PYG{p}{(}\PYG{l+s+sa}{f}\PYG{l+s+s2}{\PYGZdq{}}\PYG{l+s+s2}{Greatests coefficients: }\PYG{l+s+si}{\PYGZob{}}\PYG{n}{betahat}\PYG{p}{[}\PYG{n}{top4}\PYG{p}{]}\PYG{l+s+si}{\PYGZcb{}}\PYG{l+s+s2}{\PYGZdq{}}\PYG{p}{)}
\PYG{n}{plt}\PYG{o}{.}\PYG{n}{figure}\PYG{p}{(}\PYG{n}{figsize}\PYG{o}{=}\PYG{p}{(}\PYG{l+m+mi}{10}\PYG{p}{,} \PYG{l+m+mi}{6}\PYG{p}{)}\PYG{p}{)}
\PYG{n}{plt}\PYG{o}{.}\PYG{n}{plot}\PYG{p}{(}\PYG{n}{lambdas}\PYG{p}{,} \PYG{n}{mseloo}\PYG{p}{)}
\PYG{n}{plt}\PYG{o}{.}\PYG{n}{title}\PYG{p}{(}\PYG{l+s+s1}{\PYGZsq{}}\PYG{l+s+s1}{NMSE.loo vs Lambda}\PYG{l+s+s1}{\PYGZsq{}}\PYG{p}{)}
\PYG{n}{plt}\PYG{o}{.}\PYG{n}{grid}\PYG{p}{(}\PYG{l+s+s2}{\PYGZdq{}}\PYG{l+s+s2}{on}\PYG{l+s+s2}{\PYGZdq{}}\PYG{p}{)}
\PYG{n}{plt}\PYG{o}{.}\PYG{n}{show}\PYG{p}{(}\PYG{p}{)}
\end{sphinxVerbatim}

\end{sphinxuseclass}\end{sphinxVerbatimInput}
\begin{sphinxVerbatimOutput}

\begin{sphinxuseclass}{cell_output}
\begin{sphinxVerbatim}[commandchars=\\\{\}]
  0\PYGZpc{}|          | 0/50 [00:00\PYGZlt{}?, ?it/s]
\end{sphinxVerbatim}

\begin{sphinxVerbatim}[commandchars=\\\{\}]
Best lambda: 0.2822222222222222
Greatests coefficients indexes: [ 0  4 27 23]
Greatests coefficients: [\PYGZhy{}7.1367055   0.89720319 \PYGZhy{}0.34672033 \PYGZhy{}0.33732339]
\end{sphinxVerbatim}

\noindent\sphinxincludegraphics{{611fc69240f893f9a286157acc9f3d6b94c5024bc334e216f09d7a631a418c21}.png}

\end{sphinxuseclass}\end{sphinxVerbatimOutput}

\end{sphinxuseclass}

\subsubsection{Selection bias}
\label{\detokenize{05:selection-bias}}
\sphinxAtStartPar
Take care that if the features are selected using the error on the training set as an importance criterium leads to a flaw: the features may be over\sphinxhyphen{}fittingly selected and the performances then do not reflect on the testing performances dataset.

\begin{sphinxuseclass}{cell}\begin{sphinxVerbatimInput}

\begin{sphinxuseclass}{cell_input}
\begin{sphinxVerbatim}[commandchars=\\\{\}]
\PYG{k}{def}\PYG{+w}{ }\PYG{n+nf}{pred}\PYG{p}{(}\PYG{n}{X\PYGZus{}train}\PYG{p}{,} \PYG{n}{Y\PYGZus{}train}\PYG{p}{,} \PYG{n}{X\PYGZus{}test}\PYG{p}{)}\PYG{p}{:}
   \PYG{n}{reg} \PYG{o}{=} \PYG{n}{sklearn}\PYG{o}{.}\PYG{n}{tree}\PYG{o}{.}\PYG{n}{DecisionTreeRegressor}\PYG{p}{(}\PYG{p}{)}\PYG{o}{.}\PYG{n}{fit}\PYG{p}{(}\PYG{n}{X\PYGZus{}train}\PYG{p}{,} \PYG{n}{Y\PYGZus{}train}\PYG{p}{)}
   \PYG{k}{return} \PYG{n}{reg}\PYG{o}{.}\PYG{n}{predict}\PYG{p}{(}\PYG{n}{X\PYGZus{}test}\PYG{p}{)}

\PYG{n}{X\PYGZus{}tr} \PYG{o}{=} \PYG{n}{X}\PYG{o}{.}\PYG{n}{to\PYGZus{}numpy}\PYG{p}{(}\PYG{p}{)}\PYG{p}{[}\PYG{p}{:}\PYG{l+m+mi}{200}\PYG{p}{]}
\PYG{n}{Y\PYGZus{}tr} \PYG{o}{=} \PYG{n}{Y}\PYG{p}{[}\PYG{p}{:}\PYG{l+m+mi}{200}\PYG{p}{]}
\PYG{n}{X\PYGZus{}ts} \PYG{o}{=} \PYG{n}{X}\PYG{o}{.}\PYG{n}{to\PYGZus{}numpy}\PYG{p}{(}\PYG{p}{)}\PYG{p}{[}\PYG{l+m+mi}{200}\PYG{p}{:}\PYG{l+m+mi}{250}\PYG{p}{]}
\PYG{n}{Y\PYGZus{}ts} \PYG{o}{=} \PYG{n}{Y}\PYG{p}{[}\PYG{l+m+mi}{200}\PYG{p}{:}\PYG{l+m+mi}{250}\PYG{p}{]}

\PYG{n}{N\PYGZus{}tr} \PYG{o}{=} \PYG{n}{X\PYGZus{}tr}\PYG{o}{.}\PYG{n}{shape}\PYG{p}{[}\PYG{l+m+mi}{0}\PYG{p}{]}

\PYG{n}{fset} \PYG{o}{=} \PYG{p}{[}\PYG{p}{]}
\PYG{n}{bestr} \PYG{o}{=} \PYG{p}{[}\PYG{p}{]}
\PYG{n}{bests} \PYG{o}{=} \PYG{p}{[}\PYG{p}{]}

\PYG{k}{for} \PYG{n}{it} \PYG{o+ow}{in} \PYG{n}{tqdm}\PYG{p}{(}\PYG{n+nb}{range}\PYG{p}{(}\PYG{l+m+mi}{15}\PYG{p}{)}\PYG{p}{)}\PYG{p}{:}
    \PYG{n}{NMSEf} \PYG{o}{=} \PYG{p}{\PYGZob{}}\PYG{p}{\PYGZcb{}}
    \PYG{n}{candidates} \PYG{o}{=} \PYG{p}{[}\PYG{n}{f} \PYG{k}{for} \PYG{n}{f} \PYG{o+ow}{in} \PYG{n+nb}{range}\PYG{p}{(}\PYG{n}{n}\PYG{p}{)} \PYG{k}{if} \PYG{n}{f} \PYG{o+ow}{not} \PYG{o+ow}{in} \PYG{n}{fset}\PYG{p}{]}
    \PYG{k}{for} \PYG{n}{f} \PYG{o+ow}{in} \PYG{n}{candidates}\PYG{p}{:}
        \PYG{n}{e} \PYG{o}{=} \PYG{n}{np}\PYG{o}{.}\PYG{n}{empty}\PYG{p}{(}\PYG{n}{N\PYGZus{}tr}\PYG{p}{)}
        \PYG{k}{for} \PYG{n}{i} \PYG{o+ow}{in} \PYG{n+nb}{range}\PYG{p}{(}\PYG{n}{N\PYGZus{}tr}\PYG{p}{)}\PYG{p}{:}
            \PYG{n}{X\PYGZus{}train\PYGZus{}all} \PYG{o}{=} \PYG{n}{np}\PYG{o}{.}\PYG{n}{delete}\PYG{p}{(}\PYG{n}{X\PYGZus{}tr}\PYG{p}{,} \PYG{n}{i}\PYG{p}{,} \PYG{n}{axis}\PYG{o}{=}\PYG{l+m+mi}{0}\PYG{p}{)}
            \PYG{n}{Y\PYGZus{}train\PYGZus{}all} \PYG{o}{=} \PYG{n}{np}\PYG{o}{.}\PYG{n}{delete}\PYG{p}{(}\PYG{n}{Y\PYGZus{}tr}\PYG{p}{,} \PYG{n}{i}\PYG{p}{)} 
            \PYG{n}{cols} \PYG{o}{=} \PYG{p}{[}\PYG{n}{j} \PYG{o}{\PYGZhy{}} \PYG{l+m+mi}{1} \PYG{k}{for} \PYG{n}{j} \PYG{o+ow}{in} \PYG{p}{(}\PYG{n}{fset} \PYG{o}{+} \PYG{p}{[}\PYG{n}{f}\PYG{p}{]}\PYG{p}{)}\PYG{p}{]}
            \PYG{n}{X\PYGZus{}train} \PYG{o}{=} \PYG{n}{X\PYGZus{}train\PYGZus{}all}\PYG{p}{[}\PYG{p}{:}\PYG{p}{,} \PYG{n}{cols}\PYG{p}{]}
            \PYG{n}{X\PYGZus{}test\PYGZus{}sample} \PYG{o}{=} \PYG{n}{X\PYGZus{}tr}\PYG{p}{[}\PYG{n}{i}\PYG{p}{,} \PYG{n}{cols}\PYG{p}{]}
            \PYG{n}{Yhati} \PYG{o}{=} \PYG{n}{pred}\PYG{p}{(}\PYG{n}{X\PYGZus{}train}\PYG{p}{,} \PYG{n}{Y\PYGZus{}train\PYGZus{}all}\PYG{p}{,} \PYG{n}{X\PYGZus{}test\PYGZus{}sample}\PYG{o}{.}\PYG{n}{reshape}\PYG{p}{(}\PYG{l+m+mi}{1}\PYG{p}{,} \PYG{o}{\PYGZhy{}}\PYG{l+m+mi}{1}\PYG{p}{)}\PYG{p}{)}
            \PYG{n}{e}\PYG{p}{[}\PYG{n}{i}\PYG{p}{]} \PYG{o}{=} \PYG{p}{(}\PYG{n}{Y}\PYG{p}{[}\PYG{n}{i}\PYG{p}{]} \PYG{o}{\PYGZhy{}} \PYG{n}{Yhati}\PYG{p}{)}\PYG{o}{*}\PYG{o}{*}\PYG{l+m+mi}{2}\PYG{o}{/}\PYG{n}{np}\PYG{o}{.}\PYG{n}{std}\PYG{p}{(}\PYG{n}{Y\PYGZus{}tr}\PYG{p}{)}\PYG{o}{*}\PYG{o}{*}\PYG{l+m+mi}{2}
        \PYG{n}{NMSEf}\PYG{p}{[}\PYG{n}{f}\PYG{p}{]} \PYG{o}{=} \PYG{n}{np}\PYG{o}{.}\PYG{n}{mean}\PYG{p}{(}\PYG{n}{e}\PYG{p}{)}
    \PYG{n}{f\PYGZus{}best} \PYG{o}{=} \PYG{n+nb}{min}\PYG{p}{(}\PYG{n}{NMSEf}\PYG{p}{,} \PYG{n}{key}\PYG{o}{=}\PYG{n}{NMSEf}\PYG{o}{.}\PYG{n}{get}\PYG{p}{)}
    \PYG{n}{fset}\PYG{o}{.}\PYG{n}{append}\PYG{p}{(}\PYG{n}{f\PYGZus{}best}\PYG{p}{)}
    \PYG{n}{bestr}\PYG{o}{.}\PYG{n}{append}\PYG{p}{(}\PYG{n}{NMSEf}\PYG{p}{[}\PYG{n}{f\PYGZus{}best}\PYG{p}{]}\PYG{p}{)}
    \PYG{n}{cols} \PYG{o}{=} \PYG{p}{[}\PYG{n}{j} \PYG{o}{\PYGZhy{}} \PYG{l+m+mi}{1} \PYG{k}{for} \PYG{n}{j} \PYG{o+ow}{in} \PYG{n}{fset}\PYG{p}{]}
    \PYG{n}{Y\PYGZus{}hats} \PYG{o}{=} \PYG{n}{pred}\PYG{p}{(}\PYG{n}{X\PYGZus{}tr}\PYG{p}{[}\PYG{p}{:}\PYG{p}{,} \PYG{n}{cols}\PYG{p}{]}\PYG{p}{,} \PYG{n}{Y\PYGZus{}tr}\PYG{p}{,} \PYG{n}{X\PYGZus{}ts}\PYG{p}{[}\PYG{p}{:}\PYG{p}{,} \PYG{n}{cols}\PYG{p}{]}\PYG{p}{)}
    \PYG{n}{nmse\PYGZus{}test} \PYG{o}{=} \PYG{n}{nmse}\PYG{p}{(}\PYG{n}{Y\PYGZus{}hats}\PYG{p}{,} \PYG{n}{Y\PYGZus{}ts}\PYG{p}{)}
    \PYG{n}{bests}\PYG{o}{.}\PYG{n}{append}\PYG{p}{(}\PYG{n}{nmse\PYGZus{}test}\PYG{p}{)} 
    
\PYG{n}{pprint}\PYG{p}{(}\PYG{n}{fset}\PYG{p}{)}
\PYG{n}{pprint}\PYG{p}{(}\PYG{n}{np}\PYG{o}{.}\PYG{n}{array}\PYG{p}{(}\PYG{n}{bests}\PYG{p}{)} \PYG{o}{\PYGZhy{}} \PYG{n}{np}\PYG{o}{.}\PYG{n}{array}\PYG{p}{(}\PYG{n}{bestr}\PYG{p}{)}\PYG{p}{)}

\PYG{n}{plt}\PYG{o}{.}\PYG{n}{figure}\PYG{p}{(}\PYG{n}{figsize}\PYG{o}{=}\PYG{p}{(}\PYG{l+m+mi}{10}\PYG{p}{,} \PYG{l+m+mi}{6}\PYG{p}{)}\PYG{p}{)}
\PYG{n}{plt}\PYG{o}{.}\PYG{n}{plot}\PYG{p}{(}\PYG{n}{bests}\PYG{p}{,} \PYG{n}{color}\PYG{o}{=}\PYG{l+s+s2}{\PYGZdq{}}\PYG{l+s+s2}{red}\PYG{l+s+s2}{\PYGZdq{}}\PYG{p}{,} \PYG{n}{label}\PYG{o}{=}\PYG{l+s+s2}{\PYGZdq{}}\PYG{l+s+s2}{Test NMSE}\PYG{l+s+s2}{\PYGZdq{}}\PYG{p}{)}
\PYG{n}{plt}\PYG{o}{.}\PYG{n}{plot}\PYG{p}{(}\PYG{n}{bestr}\PYG{p}{,} \PYG{n}{linestyle}\PYG{o}{=}\PYG{l+s+s2}{\PYGZdq{}}\PYG{l+s+s2}{\PYGZhy{}\PYGZhy{}}\PYG{l+s+s2}{\PYGZdq{}}\PYG{p}{,} \PYG{n}{linewidth}\PYG{o}{=}\PYG{l+m+mi}{2}\PYG{p}{,} \PYG{n}{color}\PYG{o}{=}\PYG{l+s+s2}{\PYGZdq{}}\PYG{l+s+s2}{black}\PYG{l+s+s2}{\PYGZdq{}}\PYG{p}{,} \PYG{n}{label}\PYG{o}{=}\PYG{l+s+s2}{\PYGZdq{}}\PYG{l+s+s2}{Internal NMSE LOO}\PYG{l+s+s2}{\PYGZdq{}}\PYG{p}{)}
\PYG{n}{plt}\PYG{o}{.}\PYG{n}{ylim}\PYG{p}{(}\PYG{l+m+mi}{0}\PYG{p}{,} \PYG{n+nb}{max}\PYG{p}{(}\PYG{n}{bests}\PYG{p}{)}\PYG{p}{)}
\PYG{n}{plt}\PYG{o}{.}\PYG{n}{title}\PYG{p}{(}\PYG{l+s+s2}{\PYGZdq{}}\PYG{l+s+s2}{Selection bias and feature selection}\PYG{l+s+s2}{\PYGZdq{}}\PYG{p}{)}
\PYG{n}{plt}\PYG{o}{.}\PYG{n}{xlabel}\PYG{p}{(}\PYG{l+s+s2}{\PYGZdq{}}\PYG{l+s+s2}{Feature set size}\PYG{l+s+s2}{\PYGZdq{}}\PYG{p}{)}
\PYG{n}{plt}\PYG{o}{.}\PYG{n}{ylabel}\PYG{p}{(}\PYG{l+s+s2}{\PYGZdq{}}\PYG{l+s+s2}{\PYGZdq{}}\PYG{p}{)}
\PYG{n}{plt}\PYG{o}{.}\PYG{n}{legend}\PYG{p}{(}\PYG{n}{loc}\PYG{o}{=}\PYG{l+s+s2}{\PYGZdq{}}\PYG{l+s+s2}{upper left}\PYG{l+s+s2}{\PYGZdq{}}\PYG{p}{)}
\PYG{n}{plt}\PYG{o}{.}\PYG{n}{grid}\PYG{p}{(}\PYG{l+s+s2}{\PYGZdq{}}\PYG{l+s+s2}{on}\PYG{l+s+s2}{\PYGZdq{}}\PYG{p}{)}
\PYG{n}{plt}\PYG{o}{.}\PYG{n}{show}\PYG{p}{(}\PYG{p}{)}
\end{sphinxVerbatim}

\end{sphinxuseclass}\end{sphinxVerbatimInput}
\begin{sphinxVerbatimOutput}

\begin{sphinxuseclass}{cell_output}
\begin{sphinxVerbatim}[commandchars=\\\{\}]
  0\PYGZpc{}|          | 0/15 [00:00\PYGZlt{}?, ?it/s]
\end{sphinxVerbatim}
\begin{equation*}
\begin{split}\displaystyle \end{split}
\end{equation*}\begin{equation*}
\begin{split}\displaystyle \left[
\begin{array}{}
  0.3802 &  0.3628 & -0.0859 &  0.0500 & -0.0248 &  0.1694 &  0.2362 &  0.0735 &  0.2980 &  0.1232 &  0.2064 &  0.3436 &  0.3270 &  0.1484 &  0.2785
\end{array}
\right]\end{split}
\end{equation*}
\noindent\sphinxincludegraphics{{0522b1508ca0426b2a6aceddaa1031cfef55ed9c166bafc83d59f14d59cabe73}.png}

\end{sphinxuseclass}\end{sphinxVerbatimOutput}

\end{sphinxuseclass}

\subsection{Ensemble of models}
\label{\detokenize{05:ensemble-of-models}}
\sphinxAtStartPar
Let us now create an ensemble of R=20 linear models to make predictions.
\begin{itemize}
\item {} 
\sphinxAtStartPar
Use a linear model as the base model.

\end{itemize}

\begin{sphinxuseclass}{cell}\begin{sphinxVerbatimInput}

\begin{sphinxuseclass}{cell_input}
\begin{sphinxVerbatim}[commandchars=\\\{\}]
\PYG{n}{R} \PYG{o}{=} \PYG{l+m+mi}{20}
\PYG{n}{CV\PYGZus{}err\PYGZus{}lm\PYGZus{}ensemble\PYGZus{}model} \PYG{o}{=} \PYG{p}{[}\PYG{p}{]}

\PYG{k}{for} \PYG{n}{train\PYGZus{}index}\PYG{p}{,} \PYG{n}{test\PYGZus{}index} \PYG{o+ow}{in} \PYG{n}{kf}\PYG{o}{.}\PYG{n}{split}\PYG{p}{(}\PYG{n}{X}\PYG{p}{)}\PYG{p}{:}
    \PYG{n}{X\PYGZus{}tr}\PYG{p}{,} \PYG{n}{X\PYGZus{}ts} \PYG{o}{=} \PYG{n}{X}\PYG{o}{.}\PYG{n}{iloc}\PYG{p}{[}\PYG{n}{train\PYGZus{}index}\PYG{p}{]}\PYG{p}{,} \PYG{n}{X}\PYG{o}{.}\PYG{n}{iloc}\PYG{p}{[}\PYG{n}{test\PYGZus{}index}\PYG{p}{]}
    \PYG{n}{Y\PYGZus{}tr}\PYG{p}{,} \PYG{n}{Y\PYGZus{}ts} \PYG{o}{=} \PYG{n}{Y}\PYG{p}{[}\PYG{n}{train\PYGZus{}index}\PYG{p}{]}\PYG{p}{,} \PYG{n}{Y}\PYG{p}{[}\PYG{n}{test\PYGZus{}index}\PYG{p}{]}
    
    \PYG{n}{Y\PYGZus{}hat\PYGZus{}ts\PYGZus{}ensemble} \PYG{o}{=} \PYG{n}{np}\PYG{o}{.}\PYG{n}{zeros}\PYG{p}{(}\PYG{p}{(}\PYG{n}{X\PYGZus{}ts}\PYG{o}{.}\PYG{n}{shape}\PYG{p}{[}\PYG{l+m+mi}{0}\PYG{p}{]}\PYG{p}{,} \PYG{n}{R}\PYG{p}{)}\PYG{p}{)}
    \PYG{k}{for} \PYG{n}{r} \PYG{o+ow}{in} \PYG{n+nb}{range}\PYG{p}{(}\PYG{n}{R}\PYG{p}{)}\PYG{p}{:}
        \PYG{n}{idx\PYGZus{}tr\PYGZus{}resample} \PYG{o}{=} \PYG{n}{np}\PYG{o}{.}\PYG{n}{random}\PYG{o}{.}\PYG{n}{choice}\PYG{p}{(}\PYG{n}{train\PYGZus{}index}\PYG{p}{,} \PYG{n}{size}\PYG{o}{=}\PYG{n+nb}{len}\PYG{p}{(}\PYG{n}{train\PYGZus{}index}\PYG{p}{)}\PYG{p}{,} \PYG{n}{replace}\PYG{o}{=}\PYG{k+kc}{True}\PYG{p}{)}
        \PYG{n}{X\PYGZus{}tr\PYGZus{}res} \PYG{o}{=} \PYG{n}{X}\PYG{o}{.}\PYG{n}{iloc}\PYG{p}{[}\PYG{n}{idx\PYGZus{}tr\PYGZus{}resample}\PYG{p}{]}
        \PYG{n}{Y\PYGZus{}tr\PYGZus{}res} \PYG{o}{=} \PYG{n}{Y}\PYG{p}{[}\PYG{n}{idx\PYGZus{}tr\PYGZus{}resample}\PYG{p}{]}
        
        \PYG{n}{model} \PYG{o}{=} \PYG{n}{LinearRegression}\PYG{p}{(}\PYG{p}{)}
        \PYG{n}{model}\PYG{o}{.}\PYG{n}{fit}\PYG{p}{(}\PYG{n}{X\PYGZus{}tr\PYGZus{}res}\PYG{p}{,} \PYG{n}{Y\PYGZus{}tr\PYGZus{}res}\PYG{p}{)}
        \PYG{n}{Y\PYGZus{}hat\PYGZus{}ts\PYGZus{}ensemble}\PYG{p}{[}\PYG{p}{:}\PYG{p}{,} \PYG{n}{r}\PYG{p}{]} \PYG{o}{=} \PYG{n}{model}\PYG{o}{.}\PYG{n}{predict}\PYG{p}{(}\PYG{n}{X\PYGZus{}ts}\PYG{p}{)}
    
    \PYG{n}{Y\PYGZus{}hat\PYGZus{}ts} \PYG{o}{=} \PYG{n}{np}\PYG{o}{.}\PYG{n}{mean}\PYG{p}{(}\PYG{n}{Y\PYGZus{}hat\PYGZus{}ts\PYGZus{}ensemble}\PYG{p}{,} \PYG{n}{axis}\PYG{o}{=}\PYG{l+m+mi}{1}\PYG{p}{)}
    \PYG{n}{CV\PYGZus{}err\PYGZus{}lm\PYGZus{}ensemble\PYGZus{}model}\PYG{o}{.}\PYG{n}{append}\PYG{p}{(}\PYG{n}{nmse}\PYG{p}{(}\PYG{n}{Y\PYGZus{}ts}\PYG{p}{,} \PYG{n}{Y\PYGZus{}hat\PYGZus{}ts}\PYG{p}{)}\PYG{p}{)}

\PYG{n+nb}{print}\PYG{p}{(}\PYG{l+s+sa}{f}\PYG{l+s+s2}{\PYGZdq{}}\PYG{l+s+s2}{CV error: }\PYG{l+s+si}{\PYGZob{}}\PYG{n}{np}\PYG{o}{.}\PYG{n}{mean}\PYG{p}{(}\PYG{n}{CV\PYGZus{}err\PYGZus{}lm\PYGZus{}ensemble\PYGZus{}model}\PYG{p}{)}\PYG{l+s+si}{:}\PYG{l+s+s2}{.5f}\PYG{l+s+si}{\PYGZcb{}}\PYG{l+s+s2}{, std dev: }\PYG{l+s+si}{\PYGZob{}}\PYG{n}{np}\PYG{o}{.}\PYG{n}{std}\PYG{p}{(}\PYG{n}{CV\PYGZus{}err\PYGZus{}lm\PYGZus{}ensemble\PYGZus{}model}\PYG{p}{)}\PYG{l+s+si}{:}\PYG{l+s+s2}{.5f}\PYG{l+s+si}{\PYGZcb{}}\PYG{l+s+s2}{\PYGZdq{}}\PYG{p}{)}
\PYG{n+nb}{print}\PYG{p}{(}\PYG{l+s+s2}{\PYGZdq{}}\PYG{l+s+s2}{Is ensemble error lower than single model?}\PYG{l+s+s2}{\PYGZdq{}}\PYG{p}{,} \PYG{n}{np}\PYG{o}{.}\PYG{n}{mean}\PYG{p}{(}\PYG{n}{CV\PYGZus{}err\PYGZus{}lm\PYGZus{}ensemble\PYGZus{}model}\PYG{p}{)} \PYG{o}{\PYGZlt{}} \PYG{n}{np}\PYG{o}{.}\PYG{n}{mean}\PYG{p}{(}\PYG{n}{CV\PYGZus{}err\PYGZus{}lm\PYGZus{}single\PYGZus{}model}\PYG{p}{)}\PYG{p}{)}
\end{sphinxVerbatim}

\end{sphinxuseclass}\end{sphinxVerbatimInput}
\begin{sphinxVerbatimOutput}

\begin{sphinxuseclass}{cell_output}
\begin{sphinxVerbatim}[commandchars=\\\{\}]
CV error: 0.60855, std dev: 0.08307
Is ensemble error lower than single model? True
\end{sphinxVerbatim}

\end{sphinxuseclass}\end{sphinxVerbatimOutput}

\end{sphinxuseclass}\begin{itemize}
\item {} 
\sphinxAtStartPar
Use a decision tree as the base model. Is the CV error lower?

\end{itemize}

\begin{sphinxuseclass}{cell}\begin{sphinxVerbatimInput}

\begin{sphinxuseclass}{cell_input}
\begin{sphinxVerbatim}[commandchars=\\\{\}]
\PYG{n}{R} \PYG{o}{=} \PYG{l+m+mi}{20}
\PYG{n}{CV\PYGZus{}err\PYGZus{}rpart\PYGZus{}ensemble\PYGZus{}model} \PYG{o}{=} \PYG{p}{[}\PYG{p}{]}

\PYG{k}{for} \PYG{n}{train\PYGZus{}index}\PYG{p}{,} \PYG{n}{test\PYGZus{}index} \PYG{o+ow}{in} \PYG{n}{kf}\PYG{o}{.}\PYG{n}{split}\PYG{p}{(}\PYG{n}{X}\PYG{p}{)}\PYG{p}{:}
    \PYG{n}{X\PYGZus{}tr}\PYG{p}{,} \PYG{n}{X\PYGZus{}ts} \PYG{o}{=} \PYG{n}{X}\PYG{o}{.}\PYG{n}{iloc}\PYG{p}{[}\PYG{n}{train\PYGZus{}index}\PYG{p}{]}\PYG{p}{,} \PYG{n}{X}\PYG{o}{.}\PYG{n}{iloc}\PYG{p}{[}\PYG{n}{test\PYGZus{}index}\PYG{p}{]}
    \PYG{n}{Y\PYGZus{}tr}\PYG{p}{,} \PYG{n}{Y\PYGZus{}ts} \PYG{o}{=} \PYG{n}{Y}\PYG{p}{[}\PYG{n}{train\PYGZus{}index}\PYG{p}{]}\PYG{p}{,} \PYG{n}{Y}\PYG{p}{[}\PYG{n}{test\PYGZus{}index}\PYG{p}{]}
    
    \PYG{n}{Y\PYGZus{}hat\PYGZus{}ts\PYGZus{}ensemble} \PYG{o}{=} \PYG{n}{np}\PYG{o}{.}\PYG{n}{zeros}\PYG{p}{(}\PYG{p}{(}\PYG{n}{X\PYGZus{}ts}\PYG{o}{.}\PYG{n}{shape}\PYG{p}{[}\PYG{l+m+mi}{0}\PYG{p}{]}\PYG{p}{,} \PYG{n}{R}\PYG{p}{)}\PYG{p}{)}
    \PYG{k}{for} \PYG{n}{r} \PYG{o+ow}{in} \PYG{n+nb}{range}\PYG{p}{(}\PYG{n}{R}\PYG{p}{)}\PYG{p}{:}
        \PYG{n}{idx\PYGZus{}tr\PYGZus{}resample} \PYG{o}{=} \PYG{n}{np}\PYG{o}{.}\PYG{n}{random}\PYG{o}{.}\PYG{n}{choice}\PYG{p}{(}\PYG{n}{train\PYGZus{}index}\PYG{p}{,} \PYG{n}{size}\PYG{o}{=}\PYG{n+nb}{len}\PYG{p}{(}\PYG{n}{train\PYGZus{}index}\PYG{p}{)}\PYG{p}{,} \PYG{n}{replace}\PYG{o}{=}\PYG{k+kc}{True}\PYG{p}{)}
        \PYG{n}{X\PYGZus{}tr\PYGZus{}res} \PYG{o}{=} \PYG{n}{X}\PYG{o}{.}\PYG{n}{iloc}\PYG{p}{[}\PYG{n}{idx\PYGZus{}tr\PYGZus{}resample}\PYG{p}{]}
        \PYG{n}{Y\PYGZus{}tr\PYGZus{}res} \PYG{o}{=} \PYG{n}{Y}\PYG{p}{[}\PYG{n}{idx\PYGZus{}tr\PYGZus{}resample}\PYG{p}{]}
        
        \PYG{n}{model} \PYG{o}{=} \PYG{n}{DecisionTreeRegressor}\PYG{p}{(}\PYG{n}{random\PYGZus{}state}\PYG{o}{=}\PYG{n}{r}\PYG{p}{)}
        \PYG{n}{model}\PYG{o}{.}\PYG{n}{fit}\PYG{p}{(}\PYG{n}{X\PYGZus{}tr\PYGZus{}res}\PYG{p}{,} \PYG{n}{Y\PYGZus{}tr\PYGZus{}res}\PYG{p}{)}
        \PYG{n}{Y\PYGZus{}hat\PYGZus{}ts\PYGZus{}ensemble}\PYG{p}{[}\PYG{p}{:}\PYG{p}{,} \PYG{n}{r}\PYG{p}{]} \PYG{o}{=} \PYG{n}{model}\PYG{o}{.}\PYG{n}{predict}\PYG{p}{(}\PYG{n}{X\PYGZus{}ts}\PYG{p}{)}
    
    \PYG{n}{Y\PYGZus{}hat\PYGZus{}ts} \PYG{o}{=} \PYG{n}{np}\PYG{o}{.}\PYG{n}{mean}\PYG{p}{(}\PYG{n}{Y\PYGZus{}hat\PYGZus{}ts\PYGZus{}ensemble}\PYG{p}{,} \PYG{n}{axis}\PYG{o}{=}\PYG{l+m+mi}{1}\PYG{p}{)}
    \PYG{n}{CV\PYGZus{}err\PYGZus{}rpart\PYGZus{}ensemble\PYGZus{}model}\PYG{o}{.}\PYG{n}{append}\PYG{p}{(}\PYG{n}{nmse}\PYG{p}{(}\PYG{n}{Y\PYGZus{}ts}\PYG{p}{,} \PYG{n}{Y\PYGZus{}hat\PYGZus{}ts}\PYG{p}{)}\PYG{p}{)}

\PYG{n+nb}{print}\PYG{p}{(}\PYG{l+s+sa}{f}\PYG{l+s+s2}{\PYGZdq{}}\PYG{l+s+s2}{CV error: }\PYG{l+s+si}{\PYGZob{}}\PYG{n}{np}\PYG{o}{.}\PYG{n}{mean}\PYG{p}{(}\PYG{n}{CV\PYGZus{}err\PYGZus{}rpart\PYGZus{}ensemble\PYGZus{}model}\PYG{p}{)}\PYG{l+s+si}{:}\PYG{l+s+s2}{.5f}\PYG{l+s+si}{\PYGZcb{}}\PYG{l+s+s2}{, std dev: }\PYG{l+s+si}{\PYGZob{}}\PYG{n}{np}\PYG{o}{.}\PYG{n}{std}\PYG{p}{(}\PYG{n}{CV\PYGZus{}err\PYGZus{}rpart\PYGZus{}ensemble\PYGZus{}model}\PYG{p}{)}\PYG{l+s+si}{:}\PYG{l+s+s2}{.5f}\PYG{l+s+si}{\PYGZcb{}}\PYG{l+s+s2}{\PYGZdq{}}\PYG{p}{)}
\PYG{n+nb}{print}\PYG{p}{(}\PYG{l+s+s2}{\PYGZdq{}}\PYG{l+s+s2}{Is ensemble error lower than single model?}\PYG{l+s+s2}{\PYGZdq{}}\PYG{p}{,} \PYG{n}{np}\PYG{o}{.}\PYG{n}{mean}\PYG{p}{(}\PYG{n}{CV\PYGZus{}err\PYGZus{}rpart\PYGZus{}ensemble\PYGZus{}model}\PYG{p}{)} \PYG{o}{\PYGZlt{}} \PYG{n}{np}\PYG{o}{.}\PYG{n}{mean}\PYG{p}{(}\PYG{n}{CV\PYGZus{}err\PYGZus{}rpart\PYGZus{}single\PYGZus{}model}\PYG{p}{)}\PYG{p}{)}
\end{sphinxVerbatim}

\end{sphinxuseclass}\end{sphinxVerbatimInput}
\begin{sphinxVerbatimOutput}

\begin{sphinxuseclass}{cell_output}
\begin{sphinxVerbatim}[commandchars=\\\{\}]
CV error: 0.15982, std dev: 0.04733
Is ensemble error lower than single model? True
\end{sphinxVerbatim}

\end{sphinxuseclass}\end{sphinxVerbatimOutput}

\end{sphinxuseclass}

\section{Feature selection}
\label{\detokenize{05:id1}}

\subsection{Filter methods}
\label{\detokenize{05:filter-methods}}

\subsubsection{Correlation with the output}
\label{\detokenize{05:correlation-with-the-output}}
\sphinxAtStartPar
The following code performs feature selection by keeping the most correlated variables with the output. Compare the results for linear models and decision trees. What are the smallest CV errors, and how many features were needed?

\begin{sphinxuseclass}{cell}\begin{sphinxVerbatimInput}

\begin{sphinxuseclass}{cell_input}
\begin{sphinxVerbatim}[commandchars=\\\{\}]
\PYG{n}{correlations} \PYG{o}{=} \PYG{n}{np}\PYG{o}{.}\PYG{n}{abs}\PYG{p}{(}\PYG{n}{X}\PYG{o}{.}\PYG{n}{corrwith}\PYG{p}{(}\PYG{n}{pd}\PYG{o}{.}\PYG{n}{Series}\PYG{p}{(}\PYG{n}{Y}\PYG{p}{)}\PYG{p}{)}\PYG{p}{)}
\PYG{n}{ranking\PYGZus{}corr\PYGZus{}idx} \PYG{o}{=} \PYG{n}{correlations}\PYG{o}{.}\PYG{n}{sort\PYGZus{}values}\PYG{p}{(}\PYG{n}{ascending}\PYG{o}{=}\PYG{k+kc}{False}\PYG{p}{)}\PYG{o}{.}\PYG{n}{index}

\PYG{n}{CV\PYGZus{}err} \PYG{o}{=} \PYG{n}{np}\PYG{o}{.}\PYG{n}{zeros}\PYG{p}{(}\PYG{p}{(}\PYG{n}{n}\PYG{p}{,}\PYG{l+m+mi}{10}\PYG{p}{)}\PYG{p}{)}

\PYG{n}{fold\PYGZus{}id} \PYG{o}{=} \PYG{l+m+mi}{0}
\PYG{k}{for} \PYG{n}{train\PYGZus{}index}\PYG{p}{,} \PYG{n}{test\PYGZus{}index} \PYG{o+ow}{in} \PYG{n}{kf}\PYG{o}{.}\PYG{n}{split}\PYG{p}{(}\PYG{n}{X}\PYG{p}{)}\PYG{p}{:}
    \PYG{n}{X\PYGZus{}tr}\PYG{p}{,} \PYG{n}{X\PYGZus{}ts} \PYG{o}{=} \PYG{n}{X}\PYG{o}{.}\PYG{n}{iloc}\PYG{p}{[}\PYG{n}{train\PYGZus{}index}\PYG{p}{]}\PYG{p}{,} \PYG{n}{X}\PYG{o}{.}\PYG{n}{iloc}\PYG{p}{[}\PYG{n}{test\PYGZus{}index}\PYG{p}{]}
    \PYG{n}{Y\PYGZus{}tr}\PYG{p}{,} \PYG{n}{Y\PYGZus{}ts} \PYG{o}{=} \PYG{n}{Y}\PYG{p}{[}\PYG{n}{train\PYGZus{}index}\PYG{p}{]}\PYG{p}{,} \PYG{n}{Y}\PYG{p}{[}\PYG{n}{test\PYGZus{}index}\PYG{p}{]}
    
    \PYG{k}{for} \PYG{n}{nb\PYGZus{}features} \PYG{o+ow}{in} \PYG{n+nb}{range}\PYG{p}{(}\PYG{l+m+mi}{1}\PYG{p}{,} \PYG{n}{n}\PYG{o}{+}\PYG{l+m+mi}{1}\PYG{p}{)}\PYG{p}{:}
        \PYG{n}{selected\PYGZus{}features} \PYG{o}{=} \PYG{n}{ranking\PYGZus{}corr\PYGZus{}idx}\PYG{p}{[}\PYG{p}{:}\PYG{n}{nb\PYGZus{}features}\PYG{p}{]}
        \PYG{n}{model} \PYG{o}{=} \PYG{n}{LinearRegression}\PYG{p}{(}\PYG{p}{)}
        \PYG{n}{model}\PYG{o}{.}\PYG{n}{fit}\PYG{p}{(}\PYG{n}{X\PYGZus{}tr}\PYG{p}{[}\PYG{n}{selected\PYGZus{}features}\PYG{p}{]}\PYG{p}{,} \PYG{n}{Y\PYGZus{}tr}\PYG{p}{)}
        \PYG{n}{Y\PYGZus{}hat\PYGZus{}ts} \PYG{o}{=} \PYG{n}{model}\PYG{o}{.}\PYG{n}{predict}\PYG{p}{(}\PYG{n}{X\PYGZus{}ts}\PYG{p}{[}\PYG{n}{selected\PYGZus{}features}\PYG{p}{]}\PYG{p}{)}
        \PYG{n}{CV\PYGZus{}err}\PYG{p}{[}\PYG{n}{nb\PYGZus{}features}\PYG{o}{\PYGZhy{}}\PYG{l+m+mi}{1}\PYG{p}{,} \PYG{n}{fold\PYGZus{}id}\PYG{p}{]} \PYG{o}{=} \PYG{n}{nmse}\PYG{p}{(}\PYG{n}{Y\PYGZus{}ts}\PYG{p}{,} \PYG{n}{Y\PYGZus{}hat\PYGZus{}ts}\PYG{p}{)}
    \PYG{n}{fold\PYGZus{}id} \PYG{o}{+}\PYG{o}{=} \PYG{l+m+mi}{1}

\PYG{n}{x} \PYG{o}{=} \PYG{n+nb}{list}\PYG{p}{(}\PYG{n+nb}{range}\PYG{p}{(}\PYG{l+m+mi}{1}\PYG{p}{,} \PYG{n}{n}\PYG{o}{+}\PYG{l+m+mi}{1}\PYG{p}{)}\PYG{p}{)}
\PYG{n}{errors} \PYG{o}{=} \PYG{n}{np}\PYG{o}{.}\PYG{n}{array}\PYG{p}{(}\PYG{p}{[}\PYG{n}{np}\PYG{o}{.}\PYG{n}{mean}\PYG{p}{(}\PYG{n}{CV\PYGZus{}err}\PYG{p}{[}\PYG{n}{i}\PYG{p}{,} \PYG{p}{:}\PYG{p}{]}\PYG{p}{)} \PYG{k}{for} \PYG{n}{i} \PYG{o+ow}{in} \PYG{n+nb}{range}\PYG{p}{(}\PYG{n}{n}\PYG{p}{)}\PYG{p}{]}\PYG{p}{)}
\PYG{n}{stds} \PYG{o}{=} \PYG{n}{np}\PYG{o}{.}\PYG{n}{array}\PYG{p}{(}\PYG{p}{[}\PYG{n}{np}\PYG{o}{.}\PYG{n}{std}\PYG{p}{(}\PYG{n}{CV\PYGZus{}err}\PYG{p}{[}\PYG{n}{i}\PYG{p}{,} \PYG{p}{:}\PYG{p}{]}\PYG{p}{)} \PYG{k}{for} \PYG{n}{i} \PYG{o+ow}{in} \PYG{n+nb}{range}\PYG{p}{(}\PYG{n}{n}\PYG{p}{)}\PYG{p}{]}\PYG{p}{)}

\PYG{n}{fig}\PYG{p}{,} \PYG{n}{ax} \PYG{o}{=} \PYG{n}{plt}\PYG{o}{.}\PYG{n}{subplots}\PYG{p}{(}\PYG{n}{figsize}\PYG{o}{=}\PYG{p}{(}\PYG{l+m+mi}{10}\PYG{p}{,} \PYG{l+m+mi}{6}\PYG{p}{)}\PYG{p}{)}
\PYG{n}{ax}\PYG{o}{.}\PYG{n}{plot}\PYG{p}{(}\PYG{n}{x}\PYG{p}{,} \PYG{n}{errors}\PYG{p}{,} \PYG{n}{label}\PYG{o}{=}\PYG{l+s+s2}{\PYGZdq{}}\PYG{l+s+s2}{Average NMSE}\PYG{l+s+s2}{\PYGZdq{}}\PYG{p}{)}
\PYG{n}{ax}\PYG{o}{.}\PYG{n}{fill\PYGZus{}between}\PYG{p}{(}\PYG{n}{x}\PYG{p}{,} \PYG{n}{errors} \PYG{o}{\PYGZhy{}} \PYG{n}{stds}\PYG{p}{,} \PYG{n}{errors} \PYG{o}{+} \PYG{n}{stds}\PYG{p}{,} \PYG{n}{alpha}\PYG{o}{=}\PYG{l+m+mf}{0.2}\PYG{p}{,} \PYG{n}{label}\PYG{o}{=}\PYG{l+s+s2}{\PYGZdq{}}\PYG{l+s+s2}{Standard deviation of NMSE}\PYG{l+s+s2}{\PYGZdq{}}\PYG{p}{)}
\PYG{n}{ax}\PYG{o}{.}\PYG{n}{set\PYGZus{}title}\PYG{p}{(}\PYG{l+s+s2}{\PYGZdq{}}\PYG{l+s+s2}{CV Error depending on the number of selected features}\PYG{l+s+s2}{\PYGZdq{}}\PYG{p}{)}
\PYG{n}{ax}\PYG{o}{.}\PYG{n}{set\PYGZus{}xlabel}\PYG{p}{(}\PYG{l+s+s2}{\PYGZdq{}}\PYG{l+s+s2}{Selected features}\PYG{l+s+s2}{\PYGZdq{}}\PYG{p}{)}
\PYG{n}{ax}\PYG{o}{.}\PYG{n}{set\PYGZus{}ylabel}\PYG{p}{(}\PYG{l+s+s2}{\PYGZdq{}}\PYG{l+s+s2}{NMSE}\PYG{l+s+s2}{\PYGZdq{}}\PYG{p}{)}
\PYG{n}{ax}\PYG{o}{.}\PYG{n}{set\PYGZus{}xticks}\PYG{p}{(}\PYG{n}{ticks}\PYG{o}{=}\PYG{n}{x}\PYG{p}{,} \PYG{n}{labels}\PYG{o}{=}\PYG{n}{ranking\PYGZus{}corr\PYGZus{}idx}\PYG{o}{.}\PYG{n}{tolist}\PYG{p}{(}\PYG{p}{)}\PYG{p}{,} \PYG{n}{rotation}\PYG{o}{=}\PYG{l+m+mi}{45}\PYG{p}{)}
\PYG{n}{plt}\PYG{o}{.}\PYG{n}{legend}\PYG{p}{(}\PYG{p}{)}
\PYG{n}{plt}\PYG{o}{.}\PYG{n}{grid}\PYG{p}{(}\PYG{l+s+s2}{\PYGZdq{}}\PYG{l+s+s2}{on}\PYG{l+s+s2}{\PYGZdq{}}\PYG{p}{)}
\PYG{n}{plt}\PYG{o}{.}\PYG{n}{show}\PYG{p}{(}\PYG{p}{)}
\end{sphinxVerbatim}

\end{sphinxuseclass}\end{sphinxVerbatimInput}
\begin{sphinxVerbatimOutput}

\begin{sphinxuseclass}{cell_output}
\noindent\sphinxincludegraphics{{98907362051b4018520e5dc605409838417a79e58d0703ecc2231fc318feb3de}.png}

\end{sphinxuseclass}\end{sphinxVerbatimOutput}

\end{sphinxuseclass}

\subsubsection{mRMR}
\label{\detokenize{05:mrmr}}
\sphinxAtStartPar
We will implement a simple mRMR feature selection. mRMR uses mutual information. We will approximate mutual information via the correlation\sphinxhyphen{}based formula provided.

\begin{sphinxuseclass}{cell}\begin{sphinxVerbatimInput}

\begin{sphinxuseclass}{cell_input}
\begin{sphinxVerbatim}[commandchars=\\\{\}]
\PYG{k}{def}\PYG{+w}{ }\PYG{n+nf}{mutual\PYGZus{}info\PYGZus{}corr}\PYG{p}{(}\PYG{n}{X}\PYG{p}{,} \PYG{n}{Y}\PYG{p}{)}\PYG{p}{:}
    \PYG{n}{c} \PYG{o}{=} \PYG{n}{np}\PYG{o}{.}\PYG{n}{corrcoef}\PYG{p}{(}\PYG{n}{X}\PYG{p}{,} \PYG{n}{Y}\PYG{p}{)}\PYG{p}{[}\PYG{l+m+mi}{0}\PYG{p}{,}\PYG{l+m+mi}{1}\PYG{p}{]}
    \PYG{c+c1}{\PYGZsh{} Avoid invalid value if correlation == 1 or == \PYGZhy{}1}
    \PYG{k}{if} \PYG{n+nb}{abs}\PYG{p}{(}\PYG{n}{c}\PYG{p}{)}\PYG{o}{==}\PYG{l+m+mi}{1}\PYG{p}{:}
        \PYG{n}{c} \PYG{o}{=} \PYG{l+m+mf}{0.999999}
    \PYG{k}{return} \PYG{o}{\PYGZhy{}}\PYG{l+m+mf}{0.5} \PYG{o}{*} \PYG{n}{np}\PYG{o}{.}\PYG{n}{log}\PYG{p}{(}\PYG{l+m+mi}{1} \PYG{o}{\PYGZhy{}} \PYG{n}{c}\PYG{o}{*}\PYG{o}{*}\PYG{l+m+mi}{2}\PYG{p}{)}

\PYG{k}{def}\PYG{+w}{ }\PYG{n+nf}{compute\PYGZus{}mi\PYGZus{}vector}\PYG{p}{(}\PYG{n}{X\PYGZus{}tr}\PYG{p}{,} \PYG{n}{Y\PYGZus{}tr}\PYG{p}{)}\PYG{p}{:}
    \PYG{n}{mis} \PYG{o}{=} \PYG{p}{[}\PYG{p}{]}
    \PYG{k}{for} \PYG{n}{col} \PYG{o+ow}{in} \PYG{n}{X\PYGZus{}tr}\PYG{o}{.}\PYG{n}{columns}\PYG{p}{:}
        \PYG{n}{mi} \PYG{o}{=} \PYG{n}{mutual\PYGZus{}info\PYGZus{}corr}\PYG{p}{(}\PYG{n}{X\PYGZus{}tr}\PYG{p}{[}\PYG{n}{col}\PYG{p}{]}\PYG{o}{.}\PYG{n}{values}\PYG{p}{,} \PYG{n}{Y\PYGZus{}tr}\PYG{p}{)}
        \PYG{n}{mis}\PYG{o}{.}\PYG{n}{append}\PYG{p}{(}\PYG{n}{mi}\PYG{p}{)}
    \PYG{k}{return} \PYG{n}{np}\PYG{o}{.}\PYG{n}{array}\PYG{p}{(}\PYG{n}{mis}\PYG{p}{)}

\PYG{n}{CV\PYGZus{}err} \PYG{o}{=} \PYG{n}{np}\PYG{o}{.}\PYG{n}{zeros}\PYG{p}{(}\PYG{p}{(}\PYG{n}{n}\PYG{p}{,}\PYG{l+m+mi}{10}\PYG{p}{)}\PYG{p}{)}

\PYG{n}{fold\PYGZus{}id} \PYG{o}{=} \PYG{l+m+mi}{0}
\PYG{k}{for} \PYG{n}{train\PYGZus{}index}\PYG{p}{,} \PYG{n}{test\PYGZus{}index} \PYG{o+ow}{in} \PYG{n}{kf}\PYG{o}{.}\PYG{n}{split}\PYG{p}{(}\PYG{n}{X}\PYG{p}{)}\PYG{p}{:}
    \PYG{n}{X\PYGZus{}tr}\PYG{p}{,} \PYG{n}{X\PYGZus{}ts} \PYG{o}{=} \PYG{n}{X}\PYG{o}{.}\PYG{n}{iloc}\PYG{p}{[}\PYG{n}{train\PYGZus{}index}\PYG{p}{]}\PYG{p}{,} \PYG{n}{X}\PYG{o}{.}\PYG{n}{iloc}\PYG{p}{[}\PYG{n}{test\PYGZus{}index}\PYG{p}{]}
    \PYG{n}{Y\PYGZus{}tr}\PYG{p}{,} \PYG{n}{Y\PYGZus{}ts} \PYG{o}{=} \PYG{n}{Y}\PYG{p}{[}\PYG{n}{train\PYGZus{}index}\PYG{p}{]}\PYG{p}{,} \PYG{n}{Y}\PYG{p}{[}\PYG{n}{test\PYGZus{}index}\PYG{p}{]}
    
    \PYG{n}{mutual\PYGZus{}info\PYGZus{}values} \PYG{o}{=} \PYG{n}{compute\PYGZus{}mi\PYGZus{}vector}\PYG{p}{(}\PYG{n}{X\PYGZus{}tr}\PYG{p}{,} \PYG{n}{Y\PYGZus{}tr}\PYG{p}{)}
    \PYG{n}{selected} \PYG{o}{=} \PYG{p}{[}\PYG{p}{]}
    \PYG{n}{candidates} \PYG{o}{=} \PYG{n+nb}{list}\PYG{p}{(}\PYG{n+nb}{range}\PYG{p}{(}\PYG{n}{n}\PYG{p}{)}\PYG{p}{)}
    
    \PYG{k}{for} \PYG{n}{j} \PYG{o+ow}{in} \PYG{n+nb}{range}\PYG{p}{(}\PYG{n}{n}\PYG{p}{)}\PYG{p}{:}
        \PYG{n}{redundancy\PYGZus{}score} \PYG{o}{=} \PYG{n}{np}\PYG{o}{.}\PYG{n}{zeros}\PYG{p}{(}\PYG{n+nb}{len}\PYG{p}{(}\PYG{n}{candidates}\PYG{p}{)}\PYG{p}{)}
        \PYG{k}{if} \PYG{n+nb}{len}\PYG{p}{(}\PYG{n}{selected}\PYG{p}{)}\PYG{o}{\PYGZgt{}}\PYG{l+m+mi}{0}\PYG{p}{:}
            \PYG{c+c1}{\PYGZsh{} Compute pairwise mi between selected and candidates}
            \PYG{n}{mi\PYGZus{}sc} \PYG{o}{=} \PYG{p}{[}\PYG{p}{]}
            \PYG{k}{for} \PYG{n}{cidx} \PYG{o+ow}{in} \PYG{n}{candidates}\PYG{p}{:}
                \PYG{n}{col\PYGZus{}c} \PYG{o}{=} \PYG{n}{X\PYGZus{}tr}\PYG{o}{.}\PYG{n}{iloc}\PYG{p}{[}\PYG{p}{:}\PYG{p}{,} \PYG{n}{cidx}\PYG{p}{]}
                \PYG{n}{mis\PYGZus{}c} \PYG{o}{=} \PYG{p}{[}\PYG{p}{]}
                \PYG{k}{for} \PYG{n}{sidx} \PYG{o+ow}{in} \PYG{n}{selected}\PYG{p}{:}
                    \PYG{n}{col\PYGZus{}s} \PYG{o}{=} \PYG{n}{X\PYGZus{}tr}\PYG{o}{.}\PYG{n}{iloc}\PYG{p}{[}\PYG{p}{:}\PYG{p}{,} \PYG{n}{sidx}\PYG{p}{]}
                    \PYG{c+c1}{\PYGZsh{} Compute mutual info between col\PYGZus{}s and col\PYGZus{}c}
                    \PYG{n}{cc} \PYG{o}{=} \PYG{n}{np}\PYG{o}{.}\PYG{n}{corrcoef}\PYG{p}{(}\PYG{n}{col\PYGZus{}s}\PYG{p}{,} \PYG{n}{col\PYGZus{}c}\PYG{p}{)}\PYG{p}{[}\PYG{l+m+mi}{0}\PYG{p}{,}\PYG{l+m+mi}{1}\PYG{p}{]}
                    \PYG{k}{if} \PYG{n+nb}{abs}\PYG{p}{(}\PYG{n}{cc}\PYG{p}{)}\PYG{o}{==}\PYG{l+m+mi}{1}\PYG{p}{:}
                        \PYG{n}{cc}\PYG{o}{=}\PYG{l+m+mf}{0.999999}
                    \PYG{n}{mis\PYGZus{}c}\PYG{o}{.}\PYG{n}{append}\PYG{p}{(}\PYG{o}{\PYGZhy{}}\PYG{l+m+mf}{0.5}\PYG{o}{*}\PYG{n}{np}\PYG{o}{.}\PYG{n}{log}\PYG{p}{(}\PYG{l+m+mi}{1}\PYG{o}{\PYGZhy{}}\PYG{n}{cc}\PYG{o}{*}\PYG{o}{*}\PYG{l+m+mi}{2}\PYG{p}{)}\PYG{p}{)}
                \PYG{n}{redundancy\PYGZus{}score}\PYG{p}{[}\PYG{n}{candidates}\PYG{o}{.}\PYG{n}{index}\PYG{p}{(}\PYG{n}{cidx}\PYG{p}{)}\PYG{p}{]} \PYG{o}{=} \PYG{n}{np}\PYG{o}{.}\PYG{n}{mean}\PYG{p}{(}\PYG{n}{mis\PYGZus{}c}\PYG{p}{)}
        \PYG{n}{mRMR\PYGZus{}score} \PYG{o}{=} \PYG{n}{mutual\PYGZus{}info\PYGZus{}values}\PYG{p}{[}\PYG{n}{candidates}\PYG{p}{]} \PYG{o}{\PYGZhy{}} \PYG{n}{redundancy\PYGZus{}score}
        \PYG{n}{best\PYGZus{}idx} \PYG{o}{=} \PYG{n}{candidates}\PYG{p}{[}\PYG{n}{np}\PYG{o}{.}\PYG{n}{argmax}\PYG{p}{(}\PYG{n}{mRMR\PYGZus{}score}\PYG{p}{)}\PYG{p}{]}
        \PYG{n}{selected}\PYG{o}{.}\PYG{n}{append}\PYG{p}{(}\PYG{n}{best\PYGZus{}idx}\PYG{p}{)}
        \PYG{n}{candidates}\PYG{o}{.}\PYG{n}{remove}\PYG{p}{(}\PYG{n}{best\PYGZus{}idx}\PYG{p}{)}
    
    \PYG{c+c1}{\PYGZsh{} selected is the ranking}
    \PYG{k}{for} \PYG{n}{nb\PYGZus{}features} \PYG{o+ow}{in} \PYG{n+nb}{range}\PYG{p}{(}\PYG{l+m+mi}{1}\PYG{p}{,} \PYG{n}{n}\PYG{o}{+}\PYG{l+m+mi}{1}\PYG{p}{)}\PYG{p}{:}
        \PYG{n}{features\PYGZus{}to\PYGZus{}use} \PYG{o}{=} \PYG{p}{[}\PYG{n}{X}\PYG{o}{.}\PYG{n}{columns}\PYG{p}{[}\PYG{n}{i}\PYG{p}{]} \PYG{k}{for} \PYG{n}{i} \PYG{o+ow}{in} \PYG{n}{selected}\PYG{p}{[}\PYG{p}{:}\PYG{n}{nb\PYGZus{}features}\PYG{p}{]}\PYG{p}{]}
        \PYG{n}{model} \PYG{o}{=} \PYG{n}{LinearRegression}\PYG{p}{(}\PYG{p}{)}
        \PYG{n}{model}\PYG{o}{.}\PYG{n}{fit}\PYG{p}{(}\PYG{n}{X\PYGZus{}tr}\PYG{p}{[}\PYG{n}{features\PYGZus{}to\PYGZus{}use}\PYG{p}{]}\PYG{p}{,} \PYG{n}{Y\PYGZus{}tr}\PYG{p}{)}
        \PYG{n}{Y\PYGZus{}hat\PYGZus{}ts} \PYG{o}{=} \PYG{n}{model}\PYG{o}{.}\PYG{n}{predict}\PYG{p}{(}\PYG{n}{X\PYGZus{}ts}\PYG{p}{[}\PYG{n}{features\PYGZus{}to\PYGZus{}use}\PYG{p}{]}\PYG{p}{)}
        \PYG{n}{CV\PYGZus{}err}\PYG{p}{[}\PYG{n}{nb\PYGZus{}features}\PYG{o}{\PYGZhy{}}\PYG{l+m+mi}{1}\PYG{p}{,} \PYG{n}{fold\PYGZus{}id}\PYG{p}{]} \PYG{o}{=} \PYG{n}{nmse}\PYG{p}{(}\PYG{n}{Y\PYGZus{}ts}\PYG{p}{,} \PYG{n}{Y\PYGZus{}hat\PYGZus{}ts}\PYG{p}{)}
    \PYG{n}{fold\PYGZus{}id} \PYG{o}{+}\PYG{o}{=} \PYG{l+m+mi}{1}

\PYG{n}{x} \PYG{o}{=} \PYG{n+nb}{list}\PYG{p}{(}\PYG{n+nb}{range}\PYG{p}{(}\PYG{l+m+mi}{1}\PYG{p}{,} \PYG{n}{n}\PYG{o}{+}\PYG{l+m+mi}{1}\PYG{p}{)}\PYG{p}{)}
\PYG{n}{errors} \PYG{o}{=} \PYG{n}{np}\PYG{o}{.}\PYG{n}{array}\PYG{p}{(}\PYG{p}{[}\PYG{n}{np}\PYG{o}{.}\PYG{n}{mean}\PYG{p}{(}\PYG{n}{CV\PYGZus{}err}\PYG{p}{[}\PYG{n}{i}\PYG{p}{,} \PYG{p}{:}\PYG{p}{]}\PYG{p}{)} \PYG{k}{for} \PYG{n}{i} \PYG{o+ow}{in} \PYG{n+nb}{range}\PYG{p}{(}\PYG{n}{n}\PYG{p}{)}\PYG{p}{]}\PYG{p}{)}
\PYG{n}{stds} \PYG{o}{=} \PYG{n}{np}\PYG{o}{.}\PYG{n}{array}\PYG{p}{(}\PYG{p}{[}\PYG{n}{np}\PYG{o}{.}\PYG{n}{std}\PYG{p}{(}\PYG{n}{CV\PYGZus{}err}\PYG{p}{[}\PYG{n}{i}\PYG{p}{,} \PYG{p}{:}\PYG{p}{]}\PYG{p}{)} \PYG{k}{for} \PYG{n}{i} \PYG{o+ow}{in} \PYG{n+nb}{range}\PYG{p}{(}\PYG{n}{n}\PYG{p}{)}\PYG{p}{]}\PYG{p}{)}

\PYG{n}{fig}\PYG{p}{,} \PYG{n}{ax} \PYG{o}{=} \PYG{n}{plt}\PYG{o}{.}\PYG{n}{subplots}\PYG{p}{(}\PYG{n}{figsize}\PYG{o}{=}\PYG{p}{(}\PYG{l+m+mi}{10}\PYG{p}{,} \PYG{l+m+mi}{6}\PYG{p}{)}\PYG{p}{)}
\PYG{n}{ax}\PYG{o}{.}\PYG{n}{plot}\PYG{p}{(}\PYG{n}{x}\PYG{p}{,} \PYG{n}{errors}\PYG{p}{,} \PYG{n}{label}\PYG{o}{=}\PYG{l+s+s2}{\PYGZdq{}}\PYG{l+s+s2}{Average NMSE}\PYG{l+s+s2}{\PYGZdq{}}\PYG{p}{)}
\PYG{n}{ax}\PYG{o}{.}\PYG{n}{fill\PYGZus{}between}\PYG{p}{(}\PYG{n}{x}\PYG{p}{,} \PYG{n}{errors} \PYG{o}{\PYGZhy{}} \PYG{n}{stds}\PYG{p}{,} \PYG{n}{errors} \PYG{o}{+} \PYG{n}{stds}\PYG{p}{,} \PYG{n}{alpha}\PYG{o}{=}\PYG{l+m+mf}{0.2}\PYG{p}{,} \PYG{n}{label}\PYG{o}{=}\PYG{l+s+s2}{\PYGZdq{}}\PYG{l+s+s2}{Standard deviation of NMSE}\PYG{l+s+s2}{\PYGZdq{}}\PYG{p}{)}
\PYG{n}{ax}\PYG{o}{.}\PYG{n}{set\PYGZus{}title}\PYG{p}{(}\PYG{l+s+s2}{\PYGZdq{}}\PYG{l+s+s2}{CV Error depending on the number of selected features}\PYG{l+s+s2}{\PYGZdq{}}\PYG{p}{)}
\PYG{n}{ax}\PYG{o}{.}\PYG{n}{set\PYGZus{}xlabel}\PYG{p}{(}\PYG{l+s+s2}{\PYGZdq{}}\PYG{l+s+s2}{Selected features}\PYG{l+s+s2}{\PYGZdq{}}\PYG{p}{)}
\PYG{n}{ax}\PYG{o}{.}\PYG{n}{set\PYGZus{}ylabel}\PYG{p}{(}\PYG{l+s+s2}{\PYGZdq{}}\PYG{l+s+s2}{NMSE}\PYG{l+s+s2}{\PYGZdq{}}\PYG{p}{)}
\PYG{n}{ax}\PYG{o}{.}\PYG{n}{set\PYGZus{}xticks}\PYG{p}{(}\PYG{n}{ticks}\PYG{o}{=}\PYG{n}{x}\PYG{p}{,} \PYG{n}{labels}\PYG{o}{=}\PYG{n}{selected}\PYG{p}{,} \PYG{n}{rotation}\PYG{o}{=}\PYG{l+m+mi}{45}\PYG{p}{)}
\PYG{n}{plt}\PYG{o}{.}\PYG{n}{legend}\PYG{p}{(}\PYG{p}{)}
\PYG{n}{plt}\PYG{o}{.}\PYG{n}{grid}\PYG{p}{(}\PYG{l+s+s2}{\PYGZdq{}}\PYG{l+s+s2}{on}\PYG{l+s+s2}{\PYGZdq{}}\PYG{p}{)}
\PYG{n}{plt}\PYG{o}{.}\PYG{n}{show}\PYG{p}{(}\PYG{p}{)}
\end{sphinxVerbatim}

\end{sphinxuseclass}\end{sphinxVerbatimInput}
\begin{sphinxVerbatimOutput}

\begin{sphinxuseclass}{cell_output}
\noindent\sphinxincludegraphics{{c1d7660fbfad5f63b7764a8b817a8e5b1642a10abe23c746a118c6ce87ea6634}.png}

\end{sphinxuseclass}\end{sphinxVerbatimOutput}

\end{sphinxuseclass}

\subsubsection{PCA}
\label{\detokenize{05:pca}}
\sphinxAtStartPar
The following code performs features selection by first transforming the inputs using PCA, and then keeping the most relevant principal components in the model.

\begin{sphinxuseclass}{cell}\begin{sphinxVerbatimInput}

\begin{sphinxuseclass}{cell_input}
\begin{sphinxVerbatim}[commandchars=\\\{\}]
\PYG{k+kn}{from}\PYG{+w}{ }\PYG{n+nn}{sklearn}\PYG{n+nn}{.}\PYG{n+nn}{decomposition}\PYG{+w}{ }\PYG{k+kn}{import} \PYG{n}{PCA}

\PYG{n}{pca} \PYG{o}{=} \PYG{n}{PCA}\PYG{p}{(}\PYG{p}{)}
\PYG{n}{X\PYGZus{}pca} \PYG{o}{=} \PYG{n}{pca}\PYG{o}{.}\PYG{n}{fit\PYGZus{}transform}\PYG{p}{(}\PYG{n}{X}\PYG{p}{)}

\PYG{n}{CV\PYGZus{}err} \PYG{o}{=} \PYG{n}{np}\PYG{o}{.}\PYG{n}{zeros}\PYG{p}{(}\PYG{p}{(}\PYG{n}{n}\PYG{p}{,}\PYG{l+m+mi}{10}\PYG{p}{)}\PYG{p}{)}
\PYG{n}{fold\PYGZus{}id} \PYG{o}{=} \PYG{l+m+mi}{0}
\PYG{k}{for} \PYG{n}{train\PYGZus{}index}\PYG{p}{,} \PYG{n}{test\PYGZus{}index} \PYG{o+ow}{in} \PYG{n}{kf}\PYG{o}{.}\PYG{n}{split}\PYG{p}{(}\PYG{n}{X\PYGZus{}pca}\PYG{p}{)}\PYG{p}{:}
    \PYG{n}{X\PYGZus{}tr}\PYG{p}{,} \PYG{n}{X\PYGZus{}ts} \PYG{o}{=} \PYG{n}{X\PYGZus{}pca}\PYG{p}{[}\PYG{n}{train\PYGZus{}index}\PYG{p}{]}\PYG{p}{,} \PYG{n}{X\PYGZus{}pca}\PYG{p}{[}\PYG{n}{test\PYGZus{}index}\PYG{p}{]}
    \PYG{n}{Y\PYGZus{}tr}\PYG{p}{,} \PYG{n}{Y\PYGZus{}ts} \PYG{o}{=} \PYG{n}{Y}\PYG{p}{[}\PYG{n}{train\PYGZus{}index}\PYG{p}{]}\PYG{p}{,} \PYG{n}{Y}\PYG{p}{[}\PYG{n}{test\PYGZus{}index}\PYG{p}{]}
    
    \PYG{k}{for} \PYG{n}{nb\PYGZus{}components} \PYG{o+ow}{in} \PYG{n+nb}{range}\PYG{p}{(}\PYG{l+m+mi}{1}\PYG{p}{,} \PYG{n}{n}\PYG{o}{+}\PYG{l+m+mi}{1}\PYG{p}{)}\PYG{p}{:}
        \PYG{n}{model} \PYG{o}{=} \PYG{n}{LinearRegression}\PYG{p}{(}\PYG{p}{)}
        \PYG{n}{model}\PYG{o}{.}\PYG{n}{fit}\PYG{p}{(}\PYG{n}{X\PYGZus{}tr}\PYG{p}{[}\PYG{p}{:}\PYG{p}{,} \PYG{p}{:}\PYG{n}{nb\PYGZus{}components}\PYG{p}{]}\PYG{p}{,} \PYG{n}{Y\PYGZus{}tr}\PYG{p}{)}
        \PYG{n}{Y\PYGZus{}hat\PYGZus{}ts} \PYG{o}{=} \PYG{n}{model}\PYG{o}{.}\PYG{n}{predict}\PYG{p}{(}\PYG{n}{X\PYGZus{}ts}\PYG{p}{[}\PYG{p}{:}\PYG{p}{,} \PYG{p}{:}\PYG{n}{nb\PYGZus{}components}\PYG{p}{]}\PYG{p}{)}
        \PYG{n}{CV\PYGZus{}err}\PYG{p}{[}\PYG{n}{nb\PYGZus{}components}\PYG{o}{\PYGZhy{}}\PYG{l+m+mi}{1}\PYG{p}{,} \PYG{n}{fold\PYGZus{}id}\PYG{p}{]} \PYG{o}{=} \PYG{n}{nmse}\PYG{p}{(}\PYG{n}{Y\PYGZus{}ts}\PYG{p}{,} \PYG{n}{Y\PYGZus{}hat\PYGZus{}ts}\PYG{p}{)}
    \PYG{n}{fold\PYGZus{}id} \PYG{o}{+}\PYG{o}{=} \PYG{l+m+mi}{1}

\PYG{n}{x} \PYG{o}{=} \PYG{n+nb}{list}\PYG{p}{(}\PYG{n+nb}{range}\PYG{p}{(}\PYG{l+m+mi}{1}\PYG{p}{,} \PYG{n}{n}\PYG{o}{+}\PYG{l+m+mi}{1}\PYG{p}{)}\PYG{p}{)}
\PYG{n}{errors} \PYG{o}{=} \PYG{n}{np}\PYG{o}{.}\PYG{n}{array}\PYG{p}{(}\PYG{p}{[}\PYG{n}{np}\PYG{o}{.}\PYG{n}{mean}\PYG{p}{(}\PYG{n}{CV\PYGZus{}err}\PYG{p}{[}\PYG{n}{i}\PYG{p}{,} \PYG{p}{:}\PYG{p}{]}\PYG{p}{)} \PYG{k}{for} \PYG{n}{i} \PYG{o+ow}{in} \PYG{n+nb}{range}\PYG{p}{(}\PYG{n}{n}\PYG{p}{)}\PYG{p}{]}\PYG{p}{)}
\PYG{n}{stds} \PYG{o}{=} \PYG{n}{np}\PYG{o}{.}\PYG{n}{array}\PYG{p}{(}\PYG{p}{[}\PYG{n}{np}\PYG{o}{.}\PYG{n}{std}\PYG{p}{(}\PYG{n}{CV\PYGZus{}err}\PYG{p}{[}\PYG{n}{i}\PYG{p}{,} \PYG{p}{:}\PYG{p}{]}\PYG{p}{)} \PYG{k}{for} \PYG{n}{i} \PYG{o+ow}{in} \PYG{n+nb}{range}\PYG{p}{(}\PYG{n}{n}\PYG{p}{)}\PYG{p}{]}\PYG{p}{)}

\PYG{n}{fig}\PYG{p}{,} \PYG{n}{ax} \PYG{o}{=} \PYG{n}{plt}\PYG{o}{.}\PYG{n}{subplots}\PYG{p}{(}\PYG{n}{figsize}\PYG{o}{=}\PYG{p}{(}\PYG{l+m+mi}{10}\PYG{p}{,} \PYG{l+m+mi}{6}\PYG{p}{)}\PYG{p}{)}
\PYG{n}{ax}\PYG{o}{.}\PYG{n}{plot}\PYG{p}{(}\PYG{n}{x}\PYG{p}{,} \PYG{n}{errors}\PYG{p}{,} \PYG{n}{label}\PYG{o}{=}\PYG{l+s+s2}{\PYGZdq{}}\PYG{l+s+s2}{Average NMSE}\PYG{l+s+s2}{\PYGZdq{}}\PYG{p}{)}
\PYG{n}{ax}\PYG{o}{.}\PYG{n}{fill\PYGZus{}between}\PYG{p}{(}\PYG{n}{x}\PYG{p}{,} \PYG{n}{errors} \PYG{o}{\PYGZhy{}} \PYG{n}{stds}\PYG{p}{,} \PYG{n}{errors} \PYG{o}{+} \PYG{n}{stds}\PYG{p}{,} \PYG{n}{alpha}\PYG{o}{=}\PYG{l+m+mf}{0.2}\PYG{p}{,} \PYG{n}{label}\PYG{o}{=}\PYG{l+s+s2}{\PYGZdq{}}\PYG{l+s+s2}{Standard deviation of NMSE}\PYG{l+s+s2}{\PYGZdq{}}\PYG{p}{)}
\PYG{n}{ax}\PYG{o}{.}\PYG{n}{set\PYGZus{}title}\PYG{p}{(}\PYG{l+s+s2}{\PYGZdq{}}\PYG{l+s+s2}{CV Error depending on the number of selected features}\PYG{l+s+s2}{\PYGZdq{}}\PYG{p}{)}
\PYG{n}{ax}\PYG{o}{.}\PYG{n}{set\PYGZus{}xlabel}\PYG{p}{(}\PYG{l+s+s2}{\PYGZdq{}}\PYG{l+s+s2}{Number of principal components}\PYG{l+s+s2}{\PYGZdq{}}\PYG{p}{)}
\PYG{n}{ax}\PYG{o}{.}\PYG{n}{set\PYGZus{}ylabel}\PYG{p}{(}\PYG{l+s+s2}{\PYGZdq{}}\PYG{l+s+s2}{NMSE}\PYG{l+s+s2}{\PYGZdq{}}\PYG{p}{)}
\PYG{n}{plt}\PYG{o}{.}\PYG{n}{legend}\PYG{p}{(}\PYG{p}{)}
\PYG{n}{plt}\PYG{o}{.}\PYG{n}{grid}\PYG{p}{(}\PYG{l+s+s2}{\PYGZdq{}}\PYG{l+s+s2}{on}\PYG{l+s+s2}{\PYGZdq{}}\PYG{p}{)}
\PYG{n}{plt}\PYG{o}{.}\PYG{n}{show}\PYG{p}{(}\PYG{p}{)}
\end{sphinxVerbatim}

\end{sphinxuseclass}\end{sphinxVerbatimInput}
\begin{sphinxVerbatimOutput}

\begin{sphinxuseclass}{cell_output}
\noindent\sphinxincludegraphics{{57fc8e3f0df4194548ff60dbb595bbf4789c833bce50aa00d4e97baedba4e676}.png}

\end{sphinxuseclass}\end{sphinxVerbatimOutput}

\end{sphinxuseclass}

\subsection{Wrapper methods}
\label{\detokenize{05:wrapper-methods}}
\begin{sphinxuseclass}{cell}\begin{sphinxVerbatimInput}

\begin{sphinxuseclass}{cell_input}
\begin{sphinxVerbatim}[commandchars=\\\{\}]
\PYG{n}{selected} \PYG{o}{=} \PYG{p}{[}\PYG{p}{]}
\PYG{n}{CV\PYGZus{}errs} \PYG{o}{=} \PYG{p}{[}\PYG{p}{]}
\PYG{n}{std\PYGZus{}errs} \PYG{o}{=} \PYG{p}{[}\PYG{p}{]}
\PYG{k}{for} \PYG{n}{round\PYGZus{}i} \PYG{o+ow}{in} \PYG{n}{tqdm}\PYG{p}{(}\PYG{n+nb}{range}\PYG{p}{(}\PYG{n}{n}\PYG{p}{)}\PYG{p}{)}\PYG{p}{:}
    \PYG{n}{candidates} \PYG{o}{=} \PYG{n+nb}{list}\PYG{p}{(}\PYG{n+nb}{set}\PYG{p}{(}\PYG{n+nb}{range}\PYG{p}{(}\PYG{n}{n}\PYG{p}{)}\PYG{p}{)} \PYG{o}{\PYGZhy{}} \PYG{n+nb}{set}\PYG{p}{(}\PYG{n}{selected}\PYG{p}{)}\PYG{p}{)}
    \PYG{n}{CV\PYGZus{}err\PYGZus{}temp} \PYG{o}{=} \PYG{p}{[}\PYG{p}{]}
    \PYG{k}{for} \PYG{n}{c} \PYG{o+ow}{in} \PYG{n}{candidates}\PYG{p}{:}
        \PYG{n}{features\PYGZus{}to\PYGZus{}include} \PYG{o}{=} \PYG{n}{selected} \PYG{o}{+} \PYG{p}{[}\PYG{n}{c}\PYG{p}{]}
        \PYG{n}{fold\PYGZus{}errors} \PYG{o}{=} \PYG{p}{[}\PYG{p}{]}
        \PYG{k}{for} \PYG{n}{train\PYGZus{}index}\PYG{p}{,} \PYG{n}{test\PYGZus{}index} \PYG{o+ow}{in} \PYG{n}{kf}\PYG{o}{.}\PYG{n}{split}\PYG{p}{(}\PYG{n}{X}\PYG{p}{)}\PYG{p}{:}
            \PYG{n}{X\PYGZus{}tr}\PYG{p}{,} \PYG{n}{X\PYGZus{}ts} \PYG{o}{=} \PYG{n}{X}\PYG{o}{.}\PYG{n}{iloc}\PYG{p}{[}\PYG{n}{train\PYGZus{}index}\PYG{p}{,} \PYG{n}{features\PYGZus{}to\PYGZus{}include}\PYG{p}{]}\PYG{p}{,} \PYG{n}{X}\PYG{o}{.}\PYG{n}{iloc}\PYG{p}{[}\PYG{n}{test\PYGZus{}index}\PYG{p}{,} \PYG{n}{features\PYGZus{}to\PYGZus{}include}\PYG{p}{]}
            \PYG{n}{Y\PYGZus{}tr}\PYG{p}{,} \PYG{n}{Y\PYGZus{}ts} \PYG{o}{=} \PYG{n}{Y}\PYG{p}{[}\PYG{n}{train\PYGZus{}index}\PYG{p}{]}\PYG{p}{,} \PYG{n}{Y}\PYG{p}{[}\PYG{n}{test\PYGZus{}index}\PYG{p}{]}
            
            \PYG{n}{model} \PYG{o}{=} \PYG{n}{LinearRegression}\PYG{p}{(}\PYG{p}{)}
            \PYG{n}{model}\PYG{o}{.}\PYG{n}{fit}\PYG{p}{(}\PYG{n}{X\PYGZus{}tr}\PYG{p}{,} \PYG{n}{Y\PYGZus{}tr}\PYG{p}{)}
            \PYG{n}{Y\PYGZus{}hat\PYGZus{}ts} \PYG{o}{=} \PYG{n}{model}\PYG{o}{.}\PYG{n}{predict}\PYG{p}{(}\PYG{n}{X\PYGZus{}ts}\PYG{p}{)}
            \PYG{n}{fold\PYGZus{}errors}\PYG{o}{.}\PYG{n}{append}\PYG{p}{(}\PYG{n}{nmse}\PYG{p}{(}\PYG{n}{Y\PYGZus{}ts}\PYG{p}{,} \PYG{n}{Y\PYGZus{}hat\PYGZus{}ts}\PYG{p}{)}\PYG{p}{)}
        \PYG{n}{CV\PYGZus{}err\PYGZus{}temp}\PYG{o}{.}\PYG{n}{append}\PYG{p}{(}\PYG{n}{np}\PYG{o}{.}\PYG{n}{mean}\PYG{p}{(}\PYG{n}{fold\PYGZus{}errors}\PYG{p}{)}\PYG{p}{)}
    
    \PYG{n}{best\PYGZus{}candidate} \PYG{o}{=} \PYG{n}{candidates}\PYG{p}{[}\PYG{n}{np}\PYG{o}{.}\PYG{n}{argmin}\PYG{p}{(}\PYG{n}{CV\PYGZus{}err\PYGZus{}temp}\PYG{p}{)}\PYG{p}{]}
    \PYG{n}{selected}\PYG{o}{.}\PYG{n}{append}\PYG{p}{(}\PYG{n}{best\PYGZus{}candidate}\PYG{p}{)}
    \PYG{n}{CV\PYGZus{}errs}\PYG{o}{.}\PYG{n}{append}\PYG{p}{(}\PYG{n+nb}{min}\PYG{p}{(}\PYG{n}{CV\PYGZus{}err\PYGZus{}temp}\PYG{p}{)}\PYG{p}{)}
    \PYG{n}{std\PYGZus{}errs}\PYG{o}{.}\PYG{n}{append}\PYG{p}{(}\PYG{n}{np}\PYG{o}{.}\PYG{n}{std}\PYG{p}{(}\PYG{n}{fold\PYGZus{}errors}\PYG{p}{)}\PYG{p}{)}

\PYG{n}{x} \PYG{o}{=} \PYG{n+nb}{list}\PYG{p}{(}\PYG{n+nb}{range}\PYG{p}{(}\PYG{l+m+mi}{1}\PYG{p}{,} \PYG{n}{n}\PYG{o}{+}\PYG{l+m+mi}{1}\PYG{p}{)}\PYG{p}{)}
\PYG{n}{errors} \PYG{o}{=} \PYG{n}{np}\PYG{o}{.}\PYG{n}{array}\PYG{p}{(}\PYG{n}{CV\PYGZus{}errs}\PYG{p}{)}
\PYG{n}{stds} \PYG{o}{=} \PYG{n}{np}\PYG{o}{.}\PYG{n}{array}\PYG{p}{(}\PYG{n}{std\PYGZus{}errs}\PYG{p}{)}

\PYG{n}{fig}\PYG{p}{,} \PYG{n}{ax} \PYG{o}{=} \PYG{n}{plt}\PYG{o}{.}\PYG{n}{subplots}\PYG{p}{(}\PYG{n}{figsize}\PYG{o}{=}\PYG{p}{(}\PYG{l+m+mi}{10}\PYG{p}{,} \PYG{l+m+mi}{6}\PYG{p}{)}\PYG{p}{)}
\PYG{n}{ax}\PYG{o}{.}\PYG{n}{plot}\PYG{p}{(}\PYG{n}{x}\PYG{p}{,} \PYG{n}{errors}\PYG{p}{,} \PYG{n}{label}\PYG{o}{=}\PYG{l+s+s2}{\PYGZdq{}}\PYG{l+s+s2}{Average NMSE}\PYG{l+s+s2}{\PYGZdq{}}\PYG{p}{)}
\PYG{n}{ax}\PYG{o}{.}\PYG{n}{fill\PYGZus{}between}\PYG{p}{(}\PYG{n}{x}\PYG{p}{,} \PYG{n}{errors} \PYG{o}{\PYGZhy{}} \PYG{n}{stds}\PYG{p}{,} \PYG{n}{errors} \PYG{o}{+} \PYG{n}{stds}\PYG{p}{,} \PYG{n}{alpha}\PYG{o}{=}\PYG{l+m+mf}{0.2}\PYG{p}{,} \PYG{n}{label}\PYG{o}{=}\PYG{l+s+s2}{\PYGZdq{}}\PYG{l+s+s2}{Standard deviation of NMSE}\PYG{l+s+s2}{\PYGZdq{}}\PYG{p}{)}
\PYG{n}{ax}\PYG{o}{.}\PYG{n}{set\PYGZus{}title}\PYG{p}{(}\PYG{l+s+s2}{\PYGZdq{}}\PYG{l+s+s2}{CV Error depending on the number of selected features}\PYG{l+s+s2}{\PYGZdq{}}\PYG{p}{)}
\PYG{n}{ax}\PYG{o}{.}\PYG{n}{set\PYGZus{}xlabel}\PYG{p}{(}\PYG{l+s+s2}{\PYGZdq{}}\PYG{l+s+s2}{Selected features}\PYG{l+s+s2}{\PYGZdq{}}\PYG{p}{)}
\PYG{n}{ax}\PYG{o}{.}\PYG{n}{set\PYGZus{}ylabel}\PYG{p}{(}\PYG{l+s+s2}{\PYGZdq{}}\PYG{l+s+s2}{NMSE}\PYG{l+s+s2}{\PYGZdq{}}\PYG{p}{)}
\PYG{n}{ax}\PYG{o}{.}\PYG{n}{set\PYGZus{}xticks}\PYG{p}{(}\PYG{n}{ticks}\PYG{o}{=}\PYG{n}{x}\PYG{p}{,} \PYG{n}{labels}\PYG{o}{=}\PYG{n}{selected}\PYG{p}{,} \PYG{n}{rotation}\PYG{o}{=}\PYG{l+m+mi}{45}\PYG{p}{)}
\PYG{n}{plt}\PYG{o}{.}\PYG{n}{legend}\PYG{p}{(}\PYG{p}{)}
\PYG{n}{plt}\PYG{o}{.}\PYG{n}{grid}\PYG{p}{(}\PYG{l+s+s2}{\PYGZdq{}}\PYG{l+s+s2}{on}\PYG{l+s+s2}{\PYGZdq{}}\PYG{p}{)}
\PYG{n}{plt}\PYG{o}{.}\PYG{n}{show}\PYG{p}{(}\PYG{p}{)}
\end{sphinxVerbatim}

\end{sphinxuseclass}\end{sphinxVerbatimInput}
\begin{sphinxVerbatimOutput}

\begin{sphinxuseclass}{cell_output}
\begin{sphinxVerbatim}[commandchars=\\\{\}]
  0\PYGZpc{}|          | 0/30 [00:00\PYGZlt{}?, ?it/s]
\end{sphinxVerbatim}

\noindent\sphinxincludegraphics{{46b748cb00647c127ad0cf03914aa08d42f2b82e77666ef497a310bd8a194ba2}.png}

\end{sphinxuseclass}\end{sphinxVerbatimOutput}

\end{sphinxuseclass}

\section{Using other predictive models}
\label{\detokenize{05:using-other-predictive-models}}
\sphinxAtStartPar
We can try other models like SVM (SVR), Neural Networks (MLPRegressor), K\sphinxhyphen{}Nearest Neighbors (KNeighborsRegressor) and see if performance improves.
\begin{itemize}
\item {} 
\sphinxAtStartPar
Using SVR, show the error using all the features.

\end{itemize}

\begin{sphinxuseclass}{cell}\begin{sphinxVerbatimInput}

\begin{sphinxuseclass}{cell_input}
\begin{sphinxVerbatim}[commandchars=\\\{\}]
\PYG{k+kn}{from}\PYG{+w}{ }\PYG{n+nn}{sklearn}\PYG{n+nn}{.}\PYG{n+nn}{svm}\PYG{+w}{ }\PYG{k+kn}{import} \PYG{n}{SVR}

\PYG{n}{CV\PYGZus{}err\PYGZus{}svm\PYGZus{}single\PYGZus{}model} \PYG{o}{=} \PYG{p}{[}\PYG{p}{]}
\PYG{k}{for} \PYG{n}{train\PYGZus{}index}\PYG{p}{,} \PYG{n}{test\PYGZus{}index} \PYG{o+ow}{in} \PYG{n}{kf}\PYG{o}{.}\PYG{n}{split}\PYG{p}{(}\PYG{n}{X}\PYG{p}{)}\PYG{p}{:}
    \PYG{n}{X\PYGZus{}tr}\PYG{p}{,} \PYG{n}{X\PYGZus{}ts} \PYG{o}{=} \PYG{n}{X}\PYG{o}{.}\PYG{n}{iloc}\PYG{p}{[}\PYG{n}{train\PYGZus{}index}\PYG{p}{]}\PYG{p}{,} \PYG{n}{X}\PYG{o}{.}\PYG{n}{iloc}\PYG{p}{[}\PYG{n}{test\PYGZus{}index}\PYG{p}{]}
    \PYG{n}{Y\PYGZus{}tr}\PYG{p}{,} \PYG{n}{Y\PYGZus{}ts} \PYG{o}{=} \PYG{n}{Y}\PYG{p}{[}\PYG{n}{train\PYGZus{}index}\PYG{p}{]}\PYG{p}{,} \PYG{n}{Y}\PYG{p}{[}\PYG{n}{test\PYGZus{}index}\PYG{p}{]}
    
    \PYG{n}{model} \PYG{o}{=} \PYG{n}{SVR}\PYG{p}{(}\PYG{p}{)}
    \PYG{n}{model}\PYG{o}{.}\PYG{n}{fit}\PYG{p}{(}\PYG{n}{X\PYGZus{}tr}\PYG{p}{,} \PYG{n}{Y\PYGZus{}tr}\PYG{p}{)}
    \PYG{n}{Y\PYGZus{}hat\PYGZus{}ts} \PYG{o}{=} \PYG{n}{model}\PYG{o}{.}\PYG{n}{predict}\PYG{p}{(}\PYG{n}{X\PYGZus{}ts}\PYG{p}{)}
    \PYG{n}{CV\PYGZus{}err\PYGZus{}svm\PYGZus{}single\PYGZus{}model}\PYG{o}{.}\PYG{n}{append}\PYG{p}{(}\PYG{n}{nmse}\PYG{p}{(}\PYG{n}{Y\PYGZus{}ts}\PYG{p}{,} \PYG{n}{Y\PYGZus{}hat\PYGZus{}ts}\PYG{p}{)}\PYG{p}{)}

\PYG{n+nb}{print}\PYG{p}{(}\PYG{l+s+sa}{f}\PYG{l+s+s2}{\PYGZdq{}}\PYG{l+s+s2}{CV error: }\PYG{l+s+si}{\PYGZob{}}\PYG{n}{np}\PYG{o}{.}\PYG{n}{mean}\PYG{p}{(}\PYG{n}{CV\PYGZus{}err\PYGZus{}svm\PYGZus{}single\PYGZus{}model}\PYG{p}{)}\PYG{l+s+si}{:}\PYG{l+s+s2}{.5f}\PYG{l+s+si}{\PYGZcb{}}\PYG{l+s+s2}{, std dev: }\PYG{l+s+si}{\PYGZob{}}\PYG{n}{np}\PYG{o}{.}\PYG{n}{std}\PYG{p}{(}\PYG{n}{CV\PYGZus{}err\PYGZus{}svm\PYGZus{}single\PYGZus{}model}\PYG{p}{)}\PYG{l+s+si}{:}\PYG{l+s+s2}{.5f}\PYG{l+s+si}{\PYGZcb{}}\PYG{l+s+s2}{\PYGZdq{}}\PYG{p}{)}
\end{sphinxVerbatim}

\end{sphinxuseclass}\end{sphinxVerbatimInput}
\begin{sphinxVerbatimOutput}

\begin{sphinxuseclass}{cell_output}
\begin{sphinxVerbatim}[commandchars=\\\{\}]
CV error: 0.30621, std dev: 0.03799
\end{sphinxVerbatim}

\end{sphinxuseclass}\end{sphinxVerbatimOutput}

\end{sphinxuseclass}\begin{itemize}
\item {} 
\sphinxAtStartPar
Then, using cross\sphinxhyphen{}validation, use mRMR.

\end{itemize}

\begin{sphinxuseclass}{cell}\begin{sphinxVerbatimInput}

\begin{sphinxuseclass}{cell_input}
\begin{sphinxVerbatim}[commandchars=\\\{\}]
\PYG{c+c1}{\PYGZsh{} Feature selection with mRMR for SVM}

\PYG{n}{n\PYGZus{}variables} \PYG{o}{=} \PYG{l+m+mi}{10}

\PYG{n}{CV\PYGZus{}err\PYGZus{}svm\PYGZus{}single\PYGZus{}model\PYGZus{}fs} \PYG{o}{=} \PYG{n}{np}\PYG{o}{.}\PYG{n}{zeros}\PYG{p}{(}\PYG{p}{(}\PYG{n}{n\PYGZus{}variables}\PYG{p}{,}\PYG{l+m+mi}{10}\PYG{p}{)}\PYG{p}{)}
\PYG{n}{fold\PYGZus{}id} \PYG{o}{=} \PYG{l+m+mi}{0}

\PYG{k}{for} \PYG{n}{train\PYGZus{}index}\PYG{p}{,} \PYG{n}{test\PYGZus{}index} \PYG{o+ow}{in} \PYG{n}{kf}\PYG{o}{.}\PYG{n}{split}\PYG{p}{(}\PYG{n}{X}\PYG{p}{)}\PYG{p}{:}
    \PYG{n}{X\PYGZus{}tr}\PYG{p}{,} \PYG{n}{X\PYGZus{}ts} \PYG{o}{=} \PYG{n}{X}\PYG{o}{.}\PYG{n}{iloc}\PYG{p}{[}\PYG{n}{train\PYGZus{}index}\PYG{p}{]}\PYG{p}{,} \PYG{n}{X}\PYG{o}{.}\PYG{n}{iloc}\PYG{p}{[}\PYG{n}{test\PYGZus{}index}\PYG{p}{]}
    \PYG{n}{Y\PYGZus{}tr}\PYG{p}{,} \PYG{n}{Y\PYGZus{}ts} \PYG{o}{=} \PYG{n}{Y}\PYG{p}{[}\PYG{n}{train\PYGZus{}index}\PYG{p}{]}\PYG{p}{,} \PYG{n}{Y}\PYG{p}{[}\PYG{n}{test\PYGZus{}index}\PYG{p}{]}
    
    \PYG{n}{mutual\PYGZus{}info\PYGZus{}values} \PYG{o}{=} \PYG{n}{compute\PYGZus{}mi\PYGZus{}vector}\PYG{p}{(}\PYG{n}{X\PYGZus{}tr}\PYG{p}{,} \PYG{n}{Y\PYGZus{}tr}\PYG{p}{)}
    \PYG{n}{selected} \PYG{o}{=} \PYG{p}{[}\PYG{p}{]}
    \PYG{n}{candidates} \PYG{o}{=} \PYG{n+nb}{list}\PYG{p}{(}\PYG{n+nb}{range}\PYG{p}{(}\PYG{n}{n}\PYG{p}{)}\PYG{p}{)}
    
    \PYG{k}{for} \PYG{n}{j} \PYG{o+ow}{in} \PYG{n+nb}{range}\PYG{p}{(}\PYG{n}{n\PYGZus{}variables}\PYG{p}{)}\PYG{p}{:}
        \PYG{n}{redundancy\PYGZus{}score} \PYG{o}{=} \PYG{n}{np}\PYG{o}{.}\PYG{n}{zeros}\PYG{p}{(}\PYG{n+nb}{len}\PYG{p}{(}\PYG{n}{candidates}\PYG{p}{)}\PYG{p}{)}
        \PYG{k}{if} \PYG{n+nb}{len}\PYG{p}{(}\PYG{n}{selected}\PYG{p}{)}\PYG{o}{\PYGZgt{}}\PYG{l+m+mi}{0}\PYG{p}{:}
            \PYG{k}{for} \PYG{n}{ci}\PYG{p}{,} \PYG{n}{cidx} \PYG{o+ow}{in} \PYG{n+nb}{enumerate}\PYG{p}{(}\PYG{n}{candidates}\PYG{p}{)}\PYG{p}{:}
                \PYG{n}{col\PYGZus{}c} \PYG{o}{=} \PYG{n}{X\PYGZus{}tr}\PYG{o}{.}\PYG{n}{iloc}\PYG{p}{[}\PYG{p}{:}\PYG{p}{,} \PYG{n}{cidx}\PYG{p}{]}\PYG{o}{.}\PYG{n}{values}
                \PYG{n}{mis\PYGZus{}c} \PYG{o}{=} \PYG{p}{[}\PYG{p}{]}
                \PYG{k}{for} \PYG{n}{sidx} \PYG{o+ow}{in} \PYG{n}{selected}\PYG{p}{:}
                    \PYG{n}{col\PYGZus{}s} \PYG{o}{=} \PYG{n}{X\PYGZus{}tr}\PYG{o}{.}\PYG{n}{iloc}\PYG{p}{[}\PYG{p}{:}\PYG{p}{,} \PYG{n}{sidx}\PYG{p}{]}\PYG{o}{.}\PYG{n}{values}
                    \PYG{n}{cc} \PYG{o}{=} \PYG{n}{np}\PYG{o}{.}\PYG{n}{corrcoef}\PYG{p}{(}\PYG{n}{col\PYGZus{}s}\PYG{p}{,} \PYG{n}{col\PYGZus{}c}\PYG{p}{)}\PYG{p}{[}\PYG{l+m+mi}{0}\PYG{p}{,}\PYG{l+m+mi}{1}\PYG{p}{]}
                    \PYG{k}{if} \PYG{n+nb}{abs}\PYG{p}{(}\PYG{n}{cc}\PYG{p}{)}\PYG{o}{==}\PYG{l+m+mi}{1}\PYG{p}{:}
                        \PYG{n}{cc}\PYG{o}{=}\PYG{l+m+mf}{0.999999}
                    \PYG{n}{mis\PYGZus{}c}\PYG{o}{.}\PYG{n}{append}\PYG{p}{(}\PYG{o}{\PYGZhy{}}\PYG{l+m+mf}{0.5}\PYG{o}{*}\PYG{n}{np}\PYG{o}{.}\PYG{n}{log}\PYG{p}{(}\PYG{l+m+mi}{1}\PYG{o}{\PYGZhy{}}\PYG{n}{cc}\PYG{o}{*}\PYG{o}{*}\PYG{l+m+mi}{2}\PYG{p}{)}\PYG{p}{)}
                \PYG{n}{redundancy\PYGZus{}score}\PYG{p}{[}\PYG{n}{ci}\PYG{p}{]} \PYG{o}{=} \PYG{n}{np}\PYG{o}{.}\PYG{n}{mean}\PYG{p}{(}\PYG{n}{mis\PYGZus{}c}\PYG{p}{)}
        \PYG{n}{mRMR\PYGZus{}score} \PYG{o}{=} \PYG{n}{mutual\PYGZus{}info\PYGZus{}values}\PYG{p}{[}\PYG{n}{candidates}\PYG{p}{]} \PYG{o}{\PYGZhy{}} \PYG{n}{redundancy\PYGZus{}score}
        \PYG{n}{best\PYGZus{}idx} \PYG{o}{=} \PYG{n}{candidates}\PYG{p}{[}\PYG{n}{np}\PYG{o}{.}\PYG{n}{argmax}\PYG{p}{(}\PYG{n}{mRMR\PYGZus{}score}\PYG{p}{)}\PYG{p}{]}
        \PYG{n}{selected}\PYG{o}{.}\PYG{n}{append}\PYG{p}{(}\PYG{n}{best\PYGZus{}idx}\PYG{p}{)}
        \PYG{n}{candidates}\PYG{o}{.}\PYG{n}{remove}\PYG{p}{(}\PYG{n}{best\PYGZus{}idx}\PYG{p}{)}
        
    \PYG{c+c1}{\PYGZsh{} Evaluate performance with subsets of selected features}
    \PYG{k}{for} \PYG{n}{nb\PYGZus{}features} \PYG{o+ow}{in} \PYG{n+nb}{range}\PYG{p}{(}\PYG{l+m+mi}{1}\PYG{p}{,} \PYG{n}{n\PYGZus{}variables}\PYG{o}{+}\PYG{l+m+mi}{1}\PYG{p}{)}\PYG{p}{:}
        \PYG{n}{feats} \PYG{o}{=} \PYG{p}{[}\PYG{n}{X}\PYG{o}{.}\PYG{n}{columns}\PYG{p}{[}\PYG{n}{i}\PYG{p}{]} \PYG{k}{for} \PYG{n}{i} \PYG{o+ow}{in} \PYG{n}{selected}\PYG{p}{[}\PYG{p}{:}\PYG{n}{nb\PYGZus{}features}\PYG{p}{]}\PYG{p}{]}
        \PYG{n}{model} \PYG{o}{=} \PYG{n}{SVR}\PYG{p}{(}\PYG{p}{)}
        \PYG{n}{model}\PYG{o}{.}\PYG{n}{fit}\PYG{p}{(}\PYG{n}{X\PYGZus{}tr}\PYG{p}{[}\PYG{n}{feats}\PYG{p}{]}\PYG{p}{,} \PYG{n}{Y\PYGZus{}tr}\PYG{p}{)}
        \PYG{n}{Y\PYGZus{}hat\PYGZus{}ts} \PYG{o}{=} \PYG{n}{model}\PYG{o}{.}\PYG{n}{predict}\PYG{p}{(}\PYG{n}{X\PYGZus{}ts}\PYG{p}{[}\PYG{n}{feats}\PYG{p}{]}\PYG{p}{)}
        \PYG{n}{CV\PYGZus{}err\PYGZus{}svm\PYGZus{}single\PYGZus{}model\PYGZus{}fs}\PYG{p}{[}\PYG{n}{nb\PYGZus{}features}\PYG{o}{\PYGZhy{}}\PYG{l+m+mi}{1}\PYG{p}{,} \PYG{n}{fold\PYGZus{}id}\PYG{p}{]} \PYG{o}{=} \PYG{n}{nmse}\PYG{p}{(}\PYG{n}{Y\PYGZus{}ts}\PYG{p}{,} \PYG{n}{Y\PYGZus{}hat\PYGZus{}ts}\PYG{p}{)}
    \PYG{n}{fold\PYGZus{}id}\PYG{o}{+}\PYG{o}{=}\PYG{l+m+mi}{1}

\PYG{n}{x} \PYG{o}{=} \PYG{n+nb}{list}\PYG{p}{(}\PYG{n+nb}{range}\PYG{p}{(}\PYG{l+m+mi}{1}\PYG{p}{,} \PYG{n}{n\PYGZus{}variables}\PYG{o}{+}\PYG{l+m+mi}{1}\PYG{p}{)}\PYG{p}{)}
\PYG{n}{errors} \PYG{o}{=} \PYG{n}{np}\PYG{o}{.}\PYG{n}{array}\PYG{p}{(}\PYG{p}{[}\PYG{n}{np}\PYG{o}{.}\PYG{n}{mean}\PYG{p}{(}\PYG{n}{CV\PYGZus{}err\PYGZus{}svm\PYGZus{}single\PYGZus{}model\PYGZus{}fs}\PYG{p}{[}\PYG{n}{i}\PYG{p}{,} \PYG{p}{:}\PYG{p}{]}\PYG{p}{)} \PYG{k}{for} \PYG{n}{i} \PYG{o+ow}{in} \PYG{n+nb}{range}\PYG{p}{(}\PYG{n}{n\PYGZus{}variables}\PYG{p}{)}\PYG{p}{]}\PYG{p}{)}
\PYG{n}{stds} \PYG{o}{=} \PYG{n}{np}\PYG{o}{.}\PYG{n}{array}\PYG{p}{(}\PYG{p}{[}\PYG{n}{np}\PYG{o}{.}\PYG{n}{std}\PYG{p}{(}\PYG{n}{CV\PYGZus{}err\PYGZus{}svm\PYGZus{}single\PYGZus{}model\PYGZus{}fs}\PYG{p}{[}\PYG{n}{i}\PYG{p}{,} \PYG{p}{:}\PYG{p}{]}\PYG{p}{)} \PYG{k}{for} \PYG{n}{i} \PYG{o+ow}{in} \PYG{n+nb}{range}\PYG{p}{(}\PYG{n}{n\PYGZus{}variables}\PYG{p}{)}\PYG{p}{]}\PYG{p}{)}

\PYG{n}{fig}\PYG{p}{,} \PYG{n}{ax} \PYG{o}{=} \PYG{n}{plt}\PYG{o}{.}\PYG{n}{subplots}\PYG{p}{(}\PYG{n}{figsize}\PYG{o}{=}\PYG{p}{(}\PYG{l+m+mi}{10}\PYG{p}{,} \PYG{l+m+mi}{6}\PYG{p}{)}\PYG{p}{)}
\PYG{n}{ax}\PYG{o}{.}\PYG{n}{plot}\PYG{p}{(}\PYG{n}{x}\PYG{p}{,} \PYG{n}{errors}\PYG{p}{,} \PYG{n}{label}\PYG{o}{=}\PYG{l+s+s2}{\PYGZdq{}}\PYG{l+s+s2}{Average NMSE}\PYG{l+s+s2}{\PYGZdq{}}\PYG{p}{)}
\PYG{n}{ax}\PYG{o}{.}\PYG{n}{fill\PYGZus{}between}\PYG{p}{(}\PYG{n}{x}\PYG{p}{,} \PYG{n}{errors} \PYG{o}{\PYGZhy{}} \PYG{n}{stds}\PYG{p}{,} \PYG{n}{errors} \PYG{o}{+} \PYG{n}{stds}\PYG{p}{,} \PYG{n}{alpha}\PYG{o}{=}\PYG{l+m+mf}{0.2}\PYG{p}{,} \PYG{n}{label}\PYG{o}{=}\PYG{l+s+s2}{\PYGZdq{}}\PYG{l+s+s2}{Standard deviation of NMSE}\PYG{l+s+s2}{\PYGZdq{}}\PYG{p}{)}
\PYG{n}{ax}\PYG{o}{.}\PYG{n}{set\PYGZus{}title}\PYG{p}{(}\PYG{l+s+s2}{\PYGZdq{}}\PYG{l+s+s2}{CV Error depending on the number of selected features}\PYG{l+s+s2}{\PYGZdq{}}\PYG{p}{)}
\PYG{n}{ax}\PYG{o}{.}\PYG{n}{set\PYGZus{}xlabel}\PYG{p}{(}\PYG{l+s+s2}{\PYGZdq{}}\PYG{l+s+s2}{Selected features}\PYG{l+s+s2}{\PYGZdq{}}\PYG{p}{)}
\PYG{n}{ax}\PYG{o}{.}\PYG{n}{set\PYGZus{}ylabel}\PYG{p}{(}\PYG{l+s+s2}{\PYGZdq{}}\PYG{l+s+s2}{NMSE}\PYG{l+s+s2}{\PYGZdq{}}\PYG{p}{)}
\PYG{n}{ax}\PYG{o}{.}\PYG{n}{set\PYGZus{}xticks}\PYG{p}{(}\PYG{n}{ticks}\PYG{o}{=}\PYG{n}{x}\PYG{p}{,} \PYG{n}{labels}\PYG{o}{=}\PYG{n}{selected}\PYG{p}{,} \PYG{n}{rotation}\PYG{o}{=}\PYG{l+m+mi}{45}\PYG{p}{)}
\PYG{n}{plt}\PYG{o}{.}\PYG{n}{legend}\PYG{p}{(}\PYG{p}{)}
\PYG{n}{plt}\PYG{o}{.}\PYG{n}{grid}\PYG{p}{(}\PYG{l+s+s2}{\PYGZdq{}}\PYG{l+s+s2}{on}\PYG{l+s+s2}{\PYGZdq{}}\PYG{p}{)}
\PYG{n}{plt}\PYG{o}{.}\PYG{n}{show}\PYG{p}{(}\PYG{p}{)}
\end{sphinxVerbatim}

\end{sphinxuseclass}\end{sphinxVerbatimInput}
\begin{sphinxVerbatimOutput}

\begin{sphinxuseclass}{cell_output}
\noindent\sphinxincludegraphics{{5313b7bc7684e9c217e543fd1cbebf5ae4375e0d0953141054c967b428dbde01}.png}

\end{sphinxuseclass}\end{sphinxVerbatimOutput}

\end{sphinxuseclass}\begin{itemize}
\item {} 
\sphinxAtStartPar
Do the same for the two other models.

\end{itemize}

\begin{sphinxuseclass}{cell}\begin{sphinxVerbatimInput}

\begin{sphinxuseclass}{cell_input}
\begin{sphinxVerbatim}[commandchars=\\\{\}]
\PYG{k+kn}{from}\PYG{+w}{ }\PYG{n+nn}{sklearn}\PYG{n+nn}{.}\PYG{n+nn}{neural\PYGZus{}network}\PYG{+w}{ }\PYG{k+kn}{import} \PYG{n}{MLPRegressor}

\PYG{n}{CV\PYGZus{}err\PYGZus{}nnet\PYGZus{}single\PYGZus{}model} \PYG{o}{=} \PYG{p}{[}\PYG{p}{]}

\PYG{k}{for} \PYG{n}{train\PYGZus{}index}\PYG{p}{,} \PYG{n}{test\PYGZus{}index} \PYG{o+ow}{in} \PYG{n}{kf}\PYG{o}{.}\PYG{n}{split}\PYG{p}{(}\PYG{n}{X}\PYG{p}{)}\PYG{p}{:}
    \PYG{n}{X\PYGZus{}tr}\PYG{p}{,} \PYG{n}{X\PYGZus{}ts} \PYG{o}{=} \PYG{n}{X}\PYG{o}{.}\PYG{n}{iloc}\PYG{p}{[}\PYG{n}{train\PYGZus{}index}\PYG{p}{]}\PYG{p}{,} \PYG{n}{X}\PYG{o}{.}\PYG{n}{iloc}\PYG{p}{[}\PYG{n}{test\PYGZus{}index}\PYG{p}{]}
    \PYG{n}{Y\PYGZus{}tr}\PYG{p}{,} \PYG{n}{Y\PYGZus{}ts} \PYG{o}{=} \PYG{n}{Y}\PYG{p}{[}\PYG{n}{train\PYGZus{}index}\PYG{p}{]}\PYG{p}{,} \PYG{n}{Y}\PYG{p}{[}\PYG{n}{test\PYGZus{}index}\PYG{p}{]}

    \PYG{c+c1}{\PYGZsh{} Rescale input to 0\PYGZhy{}1}
    \PYG{n}{min\PYGZus{}X}\PYG{p}{,} \PYG{n}{max\PYGZus{}X} \PYG{o}{=} \PYG{n}{np}\PYG{o}{.}\PYG{n}{min}\PYG{p}{(}\PYG{n}{X\PYGZus{}tr}\PYG{p}{)}\PYG{p}{,} \PYG{n}{np}\PYG{o}{.}\PYG{n}{max}\PYG{p}{(}\PYG{n}{X\PYGZus{}tr}\PYG{p}{)}
    \PYG{n}{X\PYGZus{}tr} \PYG{o}{=} \PYG{p}{(}\PYG{n}{X\PYGZus{}tr}\PYG{o}{\PYGZhy{}}\PYG{n}{min\PYGZus{}X}\PYG{p}{)}\PYG{o}{/}\PYG{p}{(}\PYG{n}{max\PYGZus{}X}\PYG{o}{\PYGZhy{}}\PYG{n}{min\PYGZus{}X}\PYG{p}{)}
    \PYG{n}{X\PYGZus{}ts} \PYG{o}{=} \PYG{p}{(}\PYG{n}{X\PYGZus{}ts}\PYG{o}{\PYGZhy{}}\PYG{n}{min\PYGZus{}X}\PYG{p}{)}\PYG{o}{/}\PYG{p}{(}\PYG{n}{max\PYGZus{}X}\PYG{o}{\PYGZhy{}}\PYG{n}{min\PYGZus{}X}\PYG{p}{)}
    
    \PYG{c+c1}{\PYGZsh{} Rescale output to 0\PYGZhy{}1}
    \PYG{n}{min\PYGZus{}Y}\PYG{p}{,} \PYG{n}{max\PYGZus{}Y} \PYG{o}{=} \PYG{n}{np}\PYG{o}{.}\PYG{n}{min}\PYG{p}{(}\PYG{n}{Y\PYGZus{}tr}\PYG{p}{)}\PYG{p}{,} \PYG{n}{np}\PYG{o}{.}\PYG{n}{max}\PYG{p}{(}\PYG{n}{Y\PYGZus{}tr}\PYG{p}{)}
    \PYG{n}{Y\PYGZus{}tr\PYGZus{}rescale} \PYG{o}{=} \PYG{p}{(}\PYG{n}{Y\PYGZus{}tr}\PYG{o}{\PYGZhy{}}\PYG{n}{min\PYGZus{}Y}\PYG{p}{)}\PYG{o}{/}\PYG{p}{(}\PYG{n}{max\PYGZus{}Y}\PYG{o}{\PYGZhy{}}\PYG{n}{min\PYGZus{}Y}\PYG{p}{)}
    
    \PYG{n}{model} \PYG{o}{=} \PYG{n}{MLPRegressor}\PYG{p}{(}\PYG{n}{hidden\PYGZus{}layer\PYGZus{}sizes}\PYG{o}{=}\PYG{p}{(}\PYG{l+m+mi}{20}\PYG{p}{,}\PYG{p}{)}\PYG{p}{,} \PYG{n}{max\PYGZus{}iter}\PYG{o}{=}\PYG{l+m+mi}{10000}\PYG{p}{,} \PYG{n}{learning\PYGZus{}rate\PYGZus{}init}\PYG{o}{=}\PYG{l+m+mf}{0.05}\PYG{p}{,} \PYG{n}{random\PYGZus{}state}\PYG{o}{=}\PYG{l+m+mi}{3}\PYG{p}{)}
    \PYG{n}{model}\PYG{o}{.}\PYG{n}{fit}\PYG{p}{(}\PYG{n}{X\PYGZus{}tr}\PYG{p}{,} \PYG{n}{Y\PYGZus{}tr\PYGZus{}rescale}\PYG{p}{)}
    \PYG{n}{Y\PYGZus{}hat\PYGZus{}ts} \PYG{o}{=} \PYG{p}{(}\PYG{n}{model}\PYG{o}{.}\PYG{n}{predict}\PYG{p}{(}\PYG{n}{X\PYGZus{}ts}\PYG{p}{)}\PYG{o}{*}\PYG{p}{(}\PYG{n}{max\PYGZus{}Y}\PYG{o}{\PYGZhy{}}\PYG{n}{min\PYGZus{}Y}\PYG{p}{)}\PYG{p}{)}\PYG{o}{+}\PYG{n}{min\PYGZus{}Y}
    \PYG{n}{CV\PYGZus{}err\PYGZus{}nnet\PYGZus{}single\PYGZus{}model}\PYG{o}{.}\PYG{n}{append}\PYG{p}{(}\PYG{n}{nmse}\PYG{p}{(}\PYG{n}{Y\PYGZus{}ts}\PYG{p}{,} \PYG{n}{Y\PYGZus{}hat\PYGZus{}ts}\PYG{p}{)}\PYG{p}{)}

\PYG{n+nb}{print}\PYG{p}{(}\PYG{l+s+sa}{f}\PYG{l+s+s2}{\PYGZdq{}}\PYG{l+s+s2}{CV error: }\PYG{l+s+si}{\PYGZob{}}\PYG{n}{np}\PYG{o}{.}\PYG{n}{mean}\PYG{p}{(}\PYG{n}{CV\PYGZus{}err\PYGZus{}nnet\PYGZus{}single\PYGZus{}model}\PYG{p}{)}\PYG{l+s+si}{:}\PYG{l+s+s2}{.5f}\PYG{l+s+si}{\PYGZcb{}}\PYG{l+s+s2}{, std dev: }\PYG{l+s+si}{\PYGZob{}}\PYG{n}{np}\PYG{o}{.}\PYG{n}{std}\PYG{p}{(}\PYG{n}{CV\PYGZus{}err\PYGZus{}nnet\PYGZus{}single\PYGZus{}model}\PYG{p}{)}\PYG{l+s+si}{:}\PYG{l+s+s2}{.5f}\PYG{l+s+si}{\PYGZcb{}}\PYG{l+s+s2}{\PYGZdq{}}\PYG{p}{)}
\end{sphinxVerbatim}

\end{sphinxuseclass}\end{sphinxVerbatimInput}
\begin{sphinxVerbatimOutput}

\begin{sphinxuseclass}{cell_output}
\begin{sphinxVerbatim}[commandchars=\\\{\}]
CV error: 0.74304, std dev: 0.11244
\end{sphinxVerbatim}

\end{sphinxuseclass}\end{sphinxVerbatimOutput}

\end{sphinxuseclass}
\begin{sphinxuseclass}{cell}\begin{sphinxVerbatimInput}

\begin{sphinxuseclass}{cell_input}
\begin{sphinxVerbatim}[commandchars=\\\{\}]
\PYG{n}{CV\PYGZus{}err\PYGZus{}nnet\PYGZus{}single\PYGZus{}model\PYGZus{}fs} \PYG{o}{=} \PYG{n}{np}\PYG{o}{.}\PYG{n}{zeros}\PYG{p}{(}\PYG{p}{(}\PYG{n}{n\PYGZus{}variables}\PYG{p}{,}\PYG{l+m+mi}{10}\PYG{p}{)}\PYG{p}{)}
\PYG{n}{fold\PYGZus{}id}\PYG{o}{=}\PYG{l+m+mi}{0}
\PYG{k}{for} \PYG{n}{train\PYGZus{}index}\PYG{p}{,} \PYG{n}{test\PYGZus{}index} \PYG{o+ow}{in} \PYG{n}{kf}\PYG{o}{.}\PYG{n}{split}\PYG{p}{(}\PYG{n}{X}\PYG{p}{)}\PYG{p}{:}
    \PYG{n}{X\PYGZus{}tr}\PYG{p}{,} \PYG{n}{X\PYGZus{}ts} \PYG{o}{=} \PYG{n}{X}\PYG{o}{.}\PYG{n}{iloc}\PYG{p}{[}\PYG{n}{train\PYGZus{}index}\PYG{p}{]}\PYG{p}{,} \PYG{n}{X}\PYG{o}{.}\PYG{n}{iloc}\PYG{p}{[}\PYG{n}{test\PYGZus{}index}\PYG{p}{]}
    \PYG{n}{Y\PYGZus{}tr}\PYG{p}{,} \PYG{n}{Y\PYGZus{}ts} \PYG{o}{=} \PYG{n}{Y}\PYG{p}{[}\PYG{n}{train\PYGZus{}index}\PYG{p}{]}\PYG{p}{,} \PYG{n}{Y}\PYG{p}{[}\PYG{n}{test\PYGZus{}index}\PYG{p}{]}
    
     \PYG{c+c1}{\PYGZsh{} Rescale input to 0\PYGZhy{}1}
    \PYG{n}{min\PYGZus{}X}\PYG{p}{,} \PYG{n}{max\PYGZus{}X} \PYG{o}{=} \PYG{n}{np}\PYG{o}{.}\PYG{n}{min}\PYG{p}{(}\PYG{n}{X\PYGZus{}tr}\PYG{p}{)}\PYG{p}{,} \PYG{n}{np}\PYG{o}{.}\PYG{n}{max}\PYG{p}{(}\PYG{n}{X\PYGZus{}tr}\PYG{p}{)}
    \PYG{n}{X\PYGZus{}tr} \PYG{o}{=} \PYG{p}{(}\PYG{n}{X\PYGZus{}tr}\PYG{o}{\PYGZhy{}}\PYG{n}{min\PYGZus{}X}\PYG{p}{)}\PYG{o}{/}\PYG{p}{(}\PYG{n}{max\PYGZus{}X}\PYG{o}{\PYGZhy{}}\PYG{n}{min\PYGZus{}X}\PYG{p}{)}
    \PYG{n}{X\PYGZus{}ts} \PYG{o}{=} \PYG{p}{(}\PYG{n}{X\PYGZus{}ts}\PYG{o}{\PYGZhy{}}\PYG{n}{min\PYGZus{}X}\PYG{p}{)}\PYG{o}{/}\PYG{p}{(}\PYG{n}{max\PYGZus{}X}\PYG{o}{\PYGZhy{}}\PYG{n}{min\PYGZus{}X}\PYG{p}{)}
    
    \PYG{c+c1}{\PYGZsh{} Rescale output to 0\PYGZhy{}1}
    \PYG{n}{min\PYGZus{}Y}\PYG{p}{,} \PYG{n}{max\PYGZus{}Y} \PYG{o}{=} \PYG{n}{np}\PYG{o}{.}\PYG{n}{min}\PYG{p}{(}\PYG{n}{Y\PYGZus{}tr}\PYG{p}{)}\PYG{p}{,} \PYG{n}{np}\PYG{o}{.}\PYG{n}{max}\PYG{p}{(}\PYG{n}{Y\PYGZus{}tr}\PYG{p}{)}
    \PYG{n}{Y\PYGZus{}tr\PYGZus{}rescale} \PYG{o}{=} \PYG{p}{(}\PYG{n}{Y\PYGZus{}tr}\PYG{o}{\PYGZhy{}}\PYG{n}{min\PYGZus{}Y}\PYG{p}{)}\PYG{o}{/}\PYG{p}{(}\PYG{n}{max\PYGZus{}Y}\PYG{o}{\PYGZhy{}}\PYG{n}{min\PYGZus{}Y}\PYG{p}{)}
    
    \PYG{n}{mutual\PYGZus{}info\PYGZus{}values} \PYG{o}{=} \PYG{n}{compute\PYGZus{}mi\PYGZus{}vector}\PYG{p}{(}\PYG{n}{X\PYGZus{}tr}\PYG{p}{,} \PYG{n}{Y\PYGZus{}tr}\PYG{p}{)}
    \PYG{n}{selected} \PYG{o}{=} \PYG{p}{[}\PYG{p}{]}
    \PYG{n}{candidates} \PYG{o}{=} \PYG{n+nb}{list}\PYG{p}{(}\PYG{n+nb}{range}\PYG{p}{(}\PYG{n}{n}\PYG{p}{)}\PYG{p}{)}
    
    \PYG{k}{for} \PYG{n}{j} \PYG{o+ow}{in} \PYG{n+nb}{range}\PYG{p}{(}\PYG{n}{n\PYGZus{}variables}\PYG{p}{)}\PYG{p}{:}
        \PYG{n}{redundancy\PYGZus{}score} \PYG{o}{=} \PYG{n}{np}\PYG{o}{.}\PYG{n}{zeros}\PYG{p}{(}\PYG{n+nb}{len}\PYG{p}{(}\PYG{n}{candidates}\PYG{p}{)}\PYG{p}{)}
        \PYG{k}{if} \PYG{n+nb}{len}\PYG{p}{(}\PYG{n}{selected}\PYG{p}{)}\PYG{o}{\PYGZgt{}}\PYG{l+m+mi}{0}\PYG{p}{:}
            \PYG{k}{for} \PYG{n}{ci}\PYG{p}{,} \PYG{n}{cidx} \PYG{o+ow}{in} \PYG{n+nb}{enumerate}\PYG{p}{(}\PYG{n}{candidates}\PYG{p}{)}\PYG{p}{:}
                \PYG{n}{col\PYGZus{}c} \PYG{o}{=} \PYG{n}{X\PYGZus{}tr}\PYG{o}{.}\PYG{n}{iloc}\PYG{p}{[}\PYG{p}{:}\PYG{p}{,} \PYG{n}{cidx}\PYG{p}{]}\PYG{o}{.}\PYG{n}{values}
                \PYG{n}{mis\PYGZus{}c} \PYG{o}{=} \PYG{p}{[}\PYG{p}{]}
                \PYG{k}{for} \PYG{n}{sidx} \PYG{o+ow}{in} \PYG{n}{selected}\PYG{p}{:}
                    \PYG{n}{col\PYGZus{}s} \PYG{o}{=} \PYG{n}{X\PYGZus{}tr}\PYG{o}{.}\PYG{n}{iloc}\PYG{p}{[}\PYG{p}{:}\PYG{p}{,} \PYG{n}{sidx}\PYG{p}{]}\PYG{o}{.}\PYG{n}{values}
                    \PYG{n}{cc} \PYG{o}{=} \PYG{n}{np}\PYG{o}{.}\PYG{n}{corrcoef}\PYG{p}{(}\PYG{n}{col\PYGZus{}s}\PYG{p}{,} \PYG{n}{col\PYGZus{}c}\PYG{p}{)}\PYG{p}{[}\PYG{l+m+mi}{0}\PYG{p}{,}\PYG{l+m+mi}{1}\PYG{p}{]}
                    \PYG{k}{if} \PYG{n+nb}{abs}\PYG{p}{(}\PYG{n}{cc}\PYG{p}{)}\PYG{o}{==}\PYG{l+m+mi}{1}\PYG{p}{:}
                        \PYG{n}{cc}\PYG{o}{=}\PYG{l+m+mf}{0.999999}
                    \PYG{n}{mis\PYGZus{}c}\PYG{o}{.}\PYG{n}{append}\PYG{p}{(}\PYG{o}{\PYGZhy{}}\PYG{l+m+mf}{0.5}\PYG{o}{*}\PYG{n}{np}\PYG{o}{.}\PYG{n}{log}\PYG{p}{(}\PYG{l+m+mi}{1}\PYG{o}{\PYGZhy{}}\PYG{n}{cc}\PYG{o}{*}\PYG{o}{*}\PYG{l+m+mi}{2}\PYG{p}{)}\PYG{p}{)}
                \PYG{n}{redundancy\PYGZus{}score}\PYG{p}{[}\PYG{n}{ci}\PYG{p}{]} \PYG{o}{=} \PYG{n}{np}\PYG{o}{.}\PYG{n}{mean}\PYG{p}{(}\PYG{n}{mis\PYGZus{}c}\PYG{p}{)}
        \PYG{n}{mRMR\PYGZus{}score} \PYG{o}{=} \PYG{n}{mutual\PYGZus{}info\PYGZus{}values}\PYG{p}{[}\PYG{n}{candidates}\PYG{p}{]} \PYG{o}{\PYGZhy{}} \PYG{n}{redundancy\PYGZus{}score}
        \PYG{n}{best\PYGZus{}idx} \PYG{o}{=} \PYG{n}{candidates}\PYG{p}{[}\PYG{n}{np}\PYG{o}{.}\PYG{n}{argmax}\PYG{p}{(}\PYG{n}{mRMR\PYGZus{}score}\PYG{p}{)}\PYG{p}{]}
        \PYG{n}{selected}\PYG{o}{.}\PYG{n}{append}\PYG{p}{(}\PYG{n}{best\PYGZus{}idx}\PYG{p}{)}
        \PYG{n}{candidates}\PYG{o}{.}\PYG{n}{remove}\PYG{p}{(}\PYG{n}{best\PYGZus{}idx}\PYG{p}{)}
        
    \PYG{k}{for} \PYG{n}{nb\PYGZus{}features} \PYG{o+ow}{in} \PYG{n+nb}{range}\PYG{p}{(}\PYG{l+m+mi}{1}\PYG{p}{,} \PYG{n}{n\PYGZus{}variables}\PYG{o}{+}\PYG{l+m+mi}{1}\PYG{p}{)}\PYG{p}{:}
        \PYG{n}{feats} \PYG{o}{=} \PYG{p}{[}\PYG{n}{X}\PYG{o}{.}\PYG{n}{columns}\PYG{p}{[}\PYG{n}{i}\PYG{p}{]} \PYG{k}{for} \PYG{n}{i} \PYG{o+ow}{in} \PYG{n}{selected}\PYG{p}{[}\PYG{p}{:}\PYG{n}{nb\PYGZus{}features}\PYG{p}{]}\PYG{p}{]}
        \PYG{n}{Y\PYGZus{}tr\PYGZus{}rescale} \PYG{o}{=} \PYG{n}{Y\PYGZus{}tr}\PYG{o}{/}\PYG{l+m+mf}{10.0}
        \PYG{n}{model} \PYG{o}{=} \PYG{n}{MLPRegressor}\PYG{p}{(}\PYG{n}{hidden\PYGZus{}layer\PYGZus{}sizes}\PYG{o}{=}\PYG{p}{(}\PYG{l+m+mi}{20}\PYG{p}{,}\PYG{p}{)}\PYG{p}{,} \PYG{n}{max\PYGZus{}iter}\PYG{o}{=}\PYG{l+m+mi}{1000000}\PYG{p}{,} \PYG{n}{learning\PYGZus{}rate\PYGZus{}init}\PYG{o}{=}\PYG{l+m+mf}{0.005}\PYG{p}{,} \PYG{n}{random\PYGZus{}state}\PYG{o}{=}\PYG{l+m+mi}{3}\PYG{p}{)}
        \PYG{n}{model}\PYG{o}{.}\PYG{n}{fit}\PYG{p}{(}\PYG{n}{X\PYGZus{}tr}\PYG{p}{[}\PYG{n}{feats}\PYG{p}{]}\PYG{p}{,} \PYG{n}{Y\PYGZus{}tr\PYGZus{}rescale}\PYG{p}{)}
        \PYG{n}{Y\PYGZus{}hat\PYGZus{}ts} \PYG{o}{=} \PYG{n}{model}\PYG{o}{.}\PYG{n}{predict}\PYG{p}{(}\PYG{n}{X\PYGZus{}ts}\PYG{p}{[}\PYG{n}{feats}\PYG{p}{]}\PYG{p}{)}\PYG{o}{*}\PYG{l+m+mf}{10.0}
        \PYG{n}{CV\PYGZus{}err\PYGZus{}nnet\PYGZus{}single\PYGZus{}model\PYGZus{}fs}\PYG{p}{[}\PYG{n}{nb\PYGZus{}features}\PYG{o}{\PYGZhy{}}\PYG{l+m+mi}{1}\PYG{p}{,} \PYG{n}{fold\PYGZus{}id}\PYG{p}{]} \PYG{o}{=} \PYG{n}{nmse}\PYG{p}{(}\PYG{n}{Y\PYGZus{}ts}\PYG{p}{,} \PYG{n}{Y\PYGZus{}hat\PYGZus{}ts}\PYG{p}{)}
    \PYG{n}{fold\PYGZus{}id}\PYG{o}{+}\PYG{o}{=}\PYG{l+m+mi}{1}

\PYG{n}{x} \PYG{o}{=} \PYG{n+nb}{list}\PYG{p}{(}\PYG{n+nb}{range}\PYG{p}{(}\PYG{l+m+mi}{1}\PYG{p}{,} \PYG{n}{n\PYGZus{}variables}\PYG{o}{+}\PYG{l+m+mi}{1}\PYG{p}{)}\PYG{p}{)}
\PYG{n}{errors} \PYG{o}{=} \PYG{n}{np}\PYG{o}{.}\PYG{n}{array}\PYG{p}{(}\PYG{p}{[}\PYG{n}{np}\PYG{o}{.}\PYG{n}{mean}\PYG{p}{(}\PYG{n}{CV\PYGZus{}err\PYGZus{}nnet\PYGZus{}single\PYGZus{}model\PYGZus{}fs}\PYG{p}{[}\PYG{n}{i}\PYG{p}{,} \PYG{p}{:}\PYG{p}{]}\PYG{p}{)} \PYG{k}{for} \PYG{n}{i} \PYG{o+ow}{in} \PYG{n+nb}{range}\PYG{p}{(}\PYG{n}{n\PYGZus{}variables}\PYG{p}{)}\PYG{p}{]}\PYG{p}{)}
\PYG{n}{stds} \PYG{o}{=} \PYG{n}{np}\PYG{o}{.}\PYG{n}{array}\PYG{p}{(}\PYG{p}{[}\PYG{n}{np}\PYG{o}{.}\PYG{n}{std}\PYG{p}{(}\PYG{n}{CV\PYGZus{}err\PYGZus{}nnet\PYGZus{}single\PYGZus{}model\PYGZus{}fs}\PYG{p}{[}\PYG{n}{i}\PYG{p}{,} \PYG{p}{:}\PYG{p}{]}\PYG{p}{)} \PYG{k}{for} \PYG{n}{i} \PYG{o+ow}{in} \PYG{n+nb}{range}\PYG{p}{(}\PYG{n}{n\PYGZus{}variables}\PYG{p}{)}\PYG{p}{]}\PYG{p}{)}

\PYG{n}{fig}\PYG{p}{,} \PYG{n}{ax} \PYG{o}{=} \PYG{n}{plt}\PYG{o}{.}\PYG{n}{subplots}\PYG{p}{(}\PYG{n}{figsize}\PYG{o}{=}\PYG{p}{(}\PYG{l+m+mi}{10}\PYG{p}{,} \PYG{l+m+mi}{6}\PYG{p}{)}\PYG{p}{)}
\PYG{n}{ax}\PYG{o}{.}\PYG{n}{plot}\PYG{p}{(}\PYG{n}{x}\PYG{p}{,} \PYG{n}{errors}\PYG{p}{,} \PYG{n}{label}\PYG{o}{=}\PYG{l+s+s2}{\PYGZdq{}}\PYG{l+s+s2}{Average NMSE}\PYG{l+s+s2}{\PYGZdq{}}\PYG{p}{)}
\PYG{n}{ax}\PYG{o}{.}\PYG{n}{fill\PYGZus{}between}\PYG{p}{(}\PYG{n}{x}\PYG{p}{,} \PYG{n}{errors} \PYG{o}{\PYGZhy{}} \PYG{n}{stds}\PYG{p}{,} \PYG{n}{errors} \PYG{o}{+} \PYG{n}{stds}\PYG{p}{,} \PYG{n}{alpha}\PYG{o}{=}\PYG{l+m+mf}{0.2}\PYG{p}{,} \PYG{n}{label}\PYG{o}{=}\PYG{l+s+s2}{\PYGZdq{}}\PYG{l+s+s2}{Standard deviation of NMSE}\PYG{l+s+s2}{\PYGZdq{}}\PYG{p}{)}
\PYG{n}{ax}\PYG{o}{.}\PYG{n}{set\PYGZus{}title}\PYG{p}{(}\PYG{l+s+s2}{\PYGZdq{}}\PYG{l+s+s2}{CV Error depending on the number of selected features}\PYG{l+s+s2}{\PYGZdq{}}\PYG{p}{)}
\PYG{n}{ax}\PYG{o}{.}\PYG{n}{set\PYGZus{}xlabel}\PYG{p}{(}\PYG{l+s+s2}{\PYGZdq{}}\PYG{l+s+s2}{Selected features}\PYG{l+s+s2}{\PYGZdq{}}\PYG{p}{)}
\PYG{n}{ax}\PYG{o}{.}\PYG{n}{set\PYGZus{}ylabel}\PYG{p}{(}\PYG{l+s+s2}{\PYGZdq{}}\PYG{l+s+s2}{NMSE}\PYG{l+s+s2}{\PYGZdq{}}\PYG{p}{)}
\PYG{n}{ax}\PYG{o}{.}\PYG{n}{set\PYGZus{}xticks}\PYG{p}{(}\PYG{n}{ticks}\PYG{o}{=}\PYG{n}{x}\PYG{p}{,} \PYG{n}{labels}\PYG{o}{=}\PYG{n}{selected}\PYG{p}{,} \PYG{n}{rotation}\PYG{o}{=}\PYG{l+m+mi}{45}\PYG{p}{)}
\PYG{n}{plt}\PYG{o}{.}\PYG{n}{legend}\PYG{p}{(}\PYG{p}{)}
\PYG{n}{plt}\PYG{o}{.}\PYG{n}{grid}\PYG{p}{(}\PYG{l+s+s2}{\PYGZdq{}}\PYG{l+s+s2}{on}\PYG{l+s+s2}{\PYGZdq{}}\PYG{p}{)}
\PYG{n}{plt}\PYG{o}{.}\PYG{n}{show}\PYG{p}{(}\PYG{p}{)}
\end{sphinxVerbatim}

\end{sphinxuseclass}\end{sphinxVerbatimInput}
\begin{sphinxVerbatimOutput}

\begin{sphinxuseclass}{cell_output}
\noindent\sphinxincludegraphics{{13c7352b34a516f5e6fd05558c1d7b6097aec09642dc6dc9fc1bf0208461ab6f}.png}

\end{sphinxuseclass}\end{sphinxVerbatimOutput}

\end{sphinxuseclass}
\begin{sphinxuseclass}{cell}\begin{sphinxVerbatimInput}

\begin{sphinxuseclass}{cell_input}
\begin{sphinxVerbatim}[commandchars=\\\{\}]
\PYG{k+kn}{from}\PYG{+w}{ }\PYG{n+nn}{sklearn}\PYG{n+nn}{.}\PYG{n+nn}{neighbors}\PYG{+w}{ }\PYG{k+kn}{import} \PYG{n}{KNeighborsRegressor}

\PYG{n}{CV\PYGZus{}err\PYGZus{}lazy\PYGZus{}single\PYGZus{}model} \PYG{o}{=} \PYG{p}{[}\PYG{p}{]}
\PYG{k}{for} \PYG{n}{train\PYGZus{}index}\PYG{p}{,} \PYG{n}{test\PYGZus{}index} \PYG{o+ow}{in} \PYG{n}{kf}\PYG{o}{.}\PYG{n}{split}\PYG{p}{(}\PYG{n}{X}\PYG{p}{)}\PYG{p}{:}
    \PYG{n}{X\PYGZus{}tr}\PYG{p}{,} \PYG{n}{X\PYGZus{}ts} \PYG{o}{=} \PYG{n}{X}\PYG{o}{.}\PYG{n}{iloc}\PYG{p}{[}\PYG{n}{train\PYGZus{}index}\PYG{p}{]}\PYG{p}{,} \PYG{n}{X}\PYG{o}{.}\PYG{n}{iloc}\PYG{p}{[}\PYG{n}{test\PYGZus{}index}\PYG{p}{]}
    \PYG{n}{Y\PYGZus{}tr}\PYG{p}{,} \PYG{n}{Y\PYGZus{}ts} \PYG{o}{=} \PYG{n}{Y}\PYG{p}{[}\PYG{n}{train\PYGZus{}index}\PYG{p}{]}\PYG{p}{,} \PYG{n}{Y}\PYG{p}{[}\PYG{n}{test\PYGZus{}index}\PYG{p}{]}
    
    \PYG{n}{model} \PYG{o}{=} \PYG{n}{KNeighborsRegressor}\PYG{p}{(}\PYG{n}{n\PYGZus{}neighbors}\PYG{o}{=}\PYG{l+m+mi}{5}\PYG{p}{)}
    \PYG{n}{model}\PYG{o}{.}\PYG{n}{fit}\PYG{p}{(}\PYG{n}{X\PYGZus{}tr}\PYG{p}{,} \PYG{n}{Y\PYGZus{}tr}\PYG{p}{)}
    \PYG{n}{Y\PYGZus{}hat\PYGZus{}ts} \PYG{o}{=} \PYG{n}{model}\PYG{o}{.}\PYG{n}{predict}\PYG{p}{(}\PYG{n}{X\PYGZus{}ts}\PYG{p}{)}
    \PYG{n}{CV\PYGZus{}err\PYGZus{}lazy\PYGZus{}single\PYGZus{}model}\PYG{o}{.}\PYG{n}{append}\PYG{p}{(}\PYG{n}{nmse}\PYG{p}{(}\PYG{n}{Y\PYGZus{}ts}\PYG{p}{,} \PYG{n}{Y\PYGZus{}hat\PYGZus{}ts}\PYG{p}{)}\PYG{p}{)}

\PYG{n+nb}{print}\PYG{p}{(}\PYG{l+s+sa}{f}\PYG{l+s+s2}{\PYGZdq{}}\PYG{l+s+s2}{CV error: }\PYG{l+s+si}{\PYGZob{}}\PYG{n}{np}\PYG{o}{.}\PYG{n}{mean}\PYG{p}{(}\PYG{n}{CV\PYGZus{}err\PYGZus{}lazy\PYGZus{}single\PYGZus{}model}\PYG{p}{)}\PYG{l+s+si}{:}\PYG{l+s+s2}{.5f}\PYG{l+s+si}{\PYGZcb{}}\PYG{l+s+s2}{, std dev: }\PYG{l+s+si}{\PYGZob{}}\PYG{n}{np}\PYG{o}{.}\PYG{n}{std}\PYG{p}{(}\PYG{n}{CV\PYGZus{}err\PYGZus{}lazy\PYGZus{}single\PYGZus{}model}\PYG{p}{)}\PYG{l+s+si}{:}\PYG{l+s+s2}{.5f}\PYG{l+s+si}{\PYGZcb{}}\PYG{l+s+s2}{\PYGZdq{}}\PYG{p}{)}
\end{sphinxVerbatim}

\end{sphinxuseclass}\end{sphinxVerbatimInput}
\begin{sphinxVerbatimOutput}

\begin{sphinxuseclass}{cell_output}
\begin{sphinxVerbatim}[commandchars=\\\{\}]
CV error: 0.64745, std dev: 0.08826
\end{sphinxVerbatim}

\end{sphinxuseclass}\end{sphinxVerbatimOutput}

\end{sphinxuseclass}
\begin{sphinxuseclass}{cell}\begin{sphinxVerbatimInput}

\begin{sphinxuseclass}{cell_input}
\begin{sphinxVerbatim}[commandchars=\\\{\}]
\PYG{c+c1}{\PYGZsh{} Feature selection for KNN using mRMR}
\PYG{n}{CV\PYGZus{}err\PYGZus{}lazy\PYGZus{}single\PYGZus{}model\PYGZus{}fs} \PYG{o}{=} \PYG{n}{np}\PYG{o}{.}\PYG{n}{zeros}\PYG{p}{(}\PYG{p}{(}\PYG{n}{n\PYGZus{}variables}\PYG{p}{,}\PYG{l+m+mi}{10}\PYG{p}{)}\PYG{p}{)}
\PYG{n}{fold\PYGZus{}id}\PYG{o}{=}\PYG{l+m+mi}{0}
\PYG{k}{for} \PYG{n}{train\PYGZus{}index}\PYG{p}{,} \PYG{n}{test\PYGZus{}index} \PYG{o+ow}{in} \PYG{n}{kf}\PYG{o}{.}\PYG{n}{split}\PYG{p}{(}\PYG{n}{X}\PYG{p}{)}\PYG{p}{:}
    \PYG{n}{X\PYGZus{}tr}\PYG{p}{,} \PYG{n}{X\PYGZus{}ts} \PYG{o}{=} \PYG{n}{X}\PYG{o}{.}\PYG{n}{iloc}\PYG{p}{[}\PYG{n}{train\PYGZus{}index}\PYG{p}{]}\PYG{p}{,} \PYG{n}{X}\PYG{o}{.}\PYG{n}{iloc}\PYG{p}{[}\PYG{n}{test\PYGZus{}index}\PYG{p}{]}
    \PYG{n}{Y\PYGZus{}tr}\PYG{p}{,} \PYG{n}{Y\PYGZus{}ts} \PYG{o}{=} \PYG{n}{Y}\PYG{p}{[}\PYG{n}{train\PYGZus{}index}\PYG{p}{]}\PYG{p}{,} \PYG{n}{Y}\PYG{p}{[}\PYG{n}{test\PYGZus{}index}\PYG{p}{]}
    
    \PYG{n}{mutual\PYGZus{}info\PYGZus{}values} \PYG{o}{=} \PYG{n}{compute\PYGZus{}mi\PYGZus{}vector}\PYG{p}{(}\PYG{n}{X\PYGZus{}tr}\PYG{p}{,} \PYG{n}{Y\PYGZus{}tr}\PYG{p}{)}
    \PYG{n}{selected} \PYG{o}{=} \PYG{p}{[}\PYG{p}{]}
    \PYG{n}{candidates} \PYG{o}{=} \PYG{n+nb}{list}\PYG{p}{(}\PYG{n+nb}{range}\PYG{p}{(}\PYG{n}{n}\PYG{p}{)}\PYG{p}{)}
    
    \PYG{k}{for} \PYG{n}{j} \PYG{o+ow}{in} \PYG{n+nb}{range}\PYG{p}{(}\PYG{n}{n\PYGZus{}variables}\PYG{p}{)}\PYG{p}{:}
        \PYG{n}{redundancy\PYGZus{}score} \PYG{o}{=} \PYG{n}{np}\PYG{o}{.}\PYG{n}{zeros}\PYG{p}{(}\PYG{n+nb}{len}\PYG{p}{(}\PYG{n}{candidates}\PYG{p}{)}\PYG{p}{)}
        \PYG{k}{if} \PYG{n+nb}{len}\PYG{p}{(}\PYG{n}{selected}\PYG{p}{)}\PYG{o}{\PYGZgt{}}\PYG{l+m+mi}{0}\PYG{p}{:}
            \PYG{k}{for} \PYG{n}{ci}\PYG{p}{,} \PYG{n}{cidx} \PYG{o+ow}{in} \PYG{n+nb}{enumerate}\PYG{p}{(}\PYG{n}{candidates}\PYG{p}{)}\PYG{p}{:}
                \PYG{n}{col\PYGZus{}c} \PYG{o}{=} \PYG{n}{X\PYGZus{}tr}\PYG{o}{.}\PYG{n}{iloc}\PYG{p}{[}\PYG{p}{:}\PYG{p}{,} \PYG{n}{cidx}\PYG{p}{]}\PYG{o}{.}\PYG{n}{values}
                \PYG{n}{mis\PYGZus{}c} \PYG{o}{=} \PYG{p}{[}\PYG{p}{]}
                \PYG{k}{for} \PYG{n}{sidx} \PYG{o+ow}{in} \PYG{n}{selected}\PYG{p}{:}
                    \PYG{n}{col\PYGZus{}s} \PYG{o}{=} \PYG{n}{X\PYGZus{}tr}\PYG{o}{.}\PYG{n}{iloc}\PYG{p}{[}\PYG{p}{:}\PYG{p}{,} \PYG{n}{sidx}\PYG{p}{]}\PYG{o}{.}\PYG{n}{values}
                    \PYG{n}{cc} \PYG{o}{=} \PYG{n}{np}\PYG{o}{.}\PYG{n}{corrcoef}\PYG{p}{(}\PYG{n}{col\PYGZus{}s}\PYG{p}{,} \PYG{n}{col\PYGZus{}c}\PYG{p}{)}\PYG{p}{[}\PYG{l+m+mi}{0}\PYG{p}{,}\PYG{l+m+mi}{1}\PYG{p}{]}
                    \PYG{k}{if} \PYG{n+nb}{abs}\PYG{p}{(}\PYG{n}{cc}\PYG{p}{)}\PYG{o}{==}\PYG{l+m+mi}{1}\PYG{p}{:}
                        \PYG{n}{cc}\PYG{o}{=}\PYG{l+m+mf}{0.999999}
                    \PYG{n}{mis\PYGZus{}c}\PYG{o}{.}\PYG{n}{append}\PYG{p}{(}\PYG{o}{\PYGZhy{}}\PYG{l+m+mf}{0.5}\PYG{o}{*}\PYG{n}{np}\PYG{o}{.}\PYG{n}{log}\PYG{p}{(}\PYG{l+m+mi}{1}\PYG{o}{\PYGZhy{}}\PYG{n}{cc}\PYG{o}{*}\PYG{o}{*}\PYG{l+m+mi}{2}\PYG{p}{)}\PYG{p}{)}
                \PYG{n}{redundancy\PYGZus{}score}\PYG{p}{[}\PYG{n}{ci}\PYG{p}{]} \PYG{o}{=} \PYG{n}{np}\PYG{o}{.}\PYG{n}{mean}\PYG{p}{(}\PYG{n}{mis\PYGZus{}c}\PYG{p}{)}
        \PYG{n}{mRMR\PYGZus{}score} \PYG{o}{=} \PYG{n}{mutual\PYGZus{}info\PYGZus{}values}\PYG{p}{[}\PYG{n}{candidates}\PYG{p}{]} \PYG{o}{\PYGZhy{}} \PYG{n}{redundancy\PYGZus{}score}
        \PYG{n}{best\PYGZus{}idx} \PYG{o}{=} \PYG{n}{candidates}\PYG{p}{[}\PYG{n}{np}\PYG{o}{.}\PYG{n}{argmax}\PYG{p}{(}\PYG{n}{mRMR\PYGZus{}score}\PYG{p}{)}\PYG{p}{]}
        \PYG{n}{selected}\PYG{o}{.}\PYG{n}{append}\PYG{p}{(}\PYG{n}{best\PYGZus{}idx}\PYG{p}{)}
        \PYG{n}{candidates}\PYG{o}{.}\PYG{n}{remove}\PYG{p}{(}\PYG{n}{best\PYGZus{}idx}\PYG{p}{)}
        
    \PYG{k}{for} \PYG{n}{nb\PYGZus{}features} \PYG{o+ow}{in} \PYG{n+nb}{range}\PYG{p}{(}\PYG{l+m+mi}{1}\PYG{p}{,} \PYG{n}{n\PYGZus{}variables}\PYG{o}{+}\PYG{l+m+mi}{1}\PYG{p}{)}\PYG{p}{:}
        \PYG{n}{feats} \PYG{o}{=} \PYG{p}{[}\PYG{n}{X}\PYG{o}{.}\PYG{n}{columns}\PYG{p}{[}\PYG{n}{i}\PYG{p}{]} \PYG{k}{for} \PYG{n}{i} \PYG{o+ow}{in} \PYG{n}{selected}\PYG{p}{[}\PYG{p}{:}\PYG{n}{nb\PYGZus{}features}\PYG{p}{]}\PYG{p}{]}
        \PYG{n}{model} \PYG{o}{=} \PYG{n}{KNeighborsRegressor}\PYG{p}{(}\PYG{n}{n\PYGZus{}neighbors}\PYG{o}{=}\PYG{l+m+mi}{5}\PYG{p}{)}
        \PYG{n}{model}\PYG{o}{.}\PYG{n}{fit}\PYG{p}{(}\PYG{n}{X\PYGZus{}tr}\PYG{p}{[}\PYG{n}{feats}\PYG{p}{]}\PYG{p}{,} \PYG{n}{Y\PYGZus{}tr}\PYG{p}{)}
        \PYG{n}{Y\PYGZus{}hat\PYGZus{}ts} \PYG{o}{=} \PYG{n}{model}\PYG{o}{.}\PYG{n}{predict}\PYG{p}{(}\PYG{n}{X\PYGZus{}ts}\PYG{p}{[}\PYG{n}{feats}\PYG{p}{]}\PYG{p}{)}
        \PYG{n}{CV\PYGZus{}err\PYGZus{}lazy\PYGZus{}single\PYGZus{}model\PYGZus{}fs}\PYG{p}{[}\PYG{n}{nb\PYGZus{}features}\PYG{o}{\PYGZhy{}}\PYG{l+m+mi}{1}\PYG{p}{,} \PYG{n}{fold\PYGZus{}id}\PYG{p}{]} \PYG{o}{=} \PYG{n}{nmse}\PYG{p}{(}\PYG{n}{Y\PYGZus{}ts}\PYG{p}{,} \PYG{n}{Y\PYGZus{}hat\PYGZus{}ts}\PYG{p}{)}
    \PYG{n}{fold\PYGZus{}id}\PYG{o}{+}\PYG{o}{=}\PYG{l+m+mi}{1}

\PYG{n}{x} \PYG{o}{=} \PYG{n+nb}{list}\PYG{p}{(}\PYG{n+nb}{range}\PYG{p}{(}\PYG{l+m+mi}{1}\PYG{p}{,} \PYG{n}{n\PYGZus{}variables}\PYG{o}{+}\PYG{l+m+mi}{1}\PYG{p}{)}\PYG{p}{)}
\PYG{n}{errors} \PYG{o}{=} \PYG{n}{np}\PYG{o}{.}\PYG{n}{array}\PYG{p}{(}\PYG{p}{[}\PYG{n}{np}\PYG{o}{.}\PYG{n}{mean}\PYG{p}{(}\PYG{n}{CV\PYGZus{}err\PYGZus{}lazy\PYGZus{}single\PYGZus{}model\PYGZus{}fs}\PYG{p}{[}\PYG{n}{i}\PYG{p}{,} \PYG{p}{:}\PYG{p}{]}\PYG{p}{)} \PYG{k}{for} \PYG{n}{i} \PYG{o+ow}{in} \PYG{n+nb}{range}\PYG{p}{(}\PYG{n}{n\PYGZus{}variables}\PYG{p}{)}\PYG{p}{]}\PYG{p}{)}
\PYG{n}{stds} \PYG{o}{=} \PYG{n}{np}\PYG{o}{.}\PYG{n}{array}\PYG{p}{(}\PYG{p}{[}\PYG{n}{np}\PYG{o}{.}\PYG{n}{std}\PYG{p}{(}\PYG{n}{CV\PYGZus{}err\PYGZus{}lazy\PYGZus{}single\PYGZus{}model\PYGZus{}fs}\PYG{p}{[}\PYG{n}{i}\PYG{p}{,} \PYG{p}{:}\PYG{p}{]}\PYG{p}{)} \PYG{k}{for} \PYG{n}{i} \PYG{o+ow}{in} \PYG{n+nb}{range}\PYG{p}{(}\PYG{n}{n\PYGZus{}variables}\PYG{p}{)}\PYG{p}{]}\PYG{p}{)}

\PYG{n}{fig}\PYG{p}{,} \PYG{n}{ax} \PYG{o}{=} \PYG{n}{plt}\PYG{o}{.}\PYG{n}{subplots}\PYG{p}{(}\PYG{n}{figsize}\PYG{o}{=}\PYG{p}{(}\PYG{l+m+mi}{10}\PYG{p}{,} \PYG{l+m+mi}{6}\PYG{p}{)}\PYG{p}{)}
\PYG{n}{ax}\PYG{o}{.}\PYG{n}{plot}\PYG{p}{(}\PYG{n}{x}\PYG{p}{,} \PYG{n}{errors}\PYG{p}{,} \PYG{n}{label}\PYG{o}{=}\PYG{l+s+s2}{\PYGZdq{}}\PYG{l+s+s2}{Average NMSE}\PYG{l+s+s2}{\PYGZdq{}}\PYG{p}{)}
\PYG{n}{ax}\PYG{o}{.}\PYG{n}{fill\PYGZus{}between}\PYG{p}{(}\PYG{n}{x}\PYG{p}{,} \PYG{n}{errors} \PYG{o}{\PYGZhy{}} \PYG{n}{stds}\PYG{p}{,} \PYG{n}{errors} \PYG{o}{+} \PYG{n}{stds}\PYG{p}{,} \PYG{n}{alpha}\PYG{o}{=}\PYG{l+m+mf}{0.2}\PYG{p}{,} \PYG{n}{label}\PYG{o}{=}\PYG{l+s+s2}{\PYGZdq{}}\PYG{l+s+s2}{Standard deviation of NMSE}\PYG{l+s+s2}{\PYGZdq{}}\PYG{p}{)}
\PYG{n}{ax}\PYG{o}{.}\PYG{n}{set\PYGZus{}title}\PYG{p}{(}\PYG{l+s+s2}{\PYGZdq{}}\PYG{l+s+s2}{CV Error depending on the number of selected features}\PYG{l+s+s2}{\PYGZdq{}}\PYG{p}{)}
\PYG{n}{ax}\PYG{o}{.}\PYG{n}{set\PYGZus{}xlabel}\PYG{p}{(}\PYG{l+s+s2}{\PYGZdq{}}\PYG{l+s+s2}{Selected features}\PYG{l+s+s2}{\PYGZdq{}}\PYG{p}{)}
\PYG{n}{ax}\PYG{o}{.}\PYG{n}{set\PYGZus{}ylabel}\PYG{p}{(}\PYG{l+s+s2}{\PYGZdq{}}\PYG{l+s+s2}{NMSE}\PYG{l+s+s2}{\PYGZdq{}}\PYG{p}{)}
\PYG{n}{ax}\PYG{o}{.}\PYG{n}{set\PYGZus{}xticks}\PYG{p}{(}\PYG{n}{ticks}\PYG{o}{=}\PYG{n}{x}\PYG{p}{,} \PYG{n}{labels}\PYG{o}{=}\PYG{n}{selected}\PYG{p}{,} \PYG{n}{rotation}\PYG{o}{=}\PYG{l+m+mi}{45}\PYG{p}{)}
\PYG{n}{plt}\PYG{o}{.}\PYG{n}{legend}\PYG{p}{(}\PYG{p}{)}
\PYG{n}{plt}\PYG{o}{.}\PYG{n}{grid}\PYG{p}{(}\PYG{l+s+s2}{\PYGZdq{}}\PYG{l+s+s2}{on}\PYG{l+s+s2}{\PYGZdq{}}\PYG{p}{)}
\PYG{n}{plt}\PYG{o}{.}\PYG{n}{show}\PYG{p}{(}\PYG{p}{)}
\end{sphinxVerbatim}

\end{sphinxuseclass}\end{sphinxVerbatimInput}
\begin{sphinxVerbatimOutput}

\begin{sphinxuseclass}{cell_output}
\noindent\sphinxincludegraphics{{a909b253dc3ffbb594ba1979da16e6c99c8c1ded0752ab6f8fa0fc22f3b17ab2}.png}

\end{sphinxuseclass}\end{sphinxVerbatimOutput}

\end{sphinxuseclass}
\sphinxstepscope


\chapter{Classification}
\label{\detokenize{06:classification}}\label{\detokenize{06::doc}}
\sphinxAtStartPar
Let us define the basic imports and a \sphinxcode{\sphinxupquote{pprint}} function.

\begin{sphinxuseclass}{cell}\begin{sphinxVerbatimInput}

\begin{sphinxuseclass}{cell_input}
\begin{sphinxVerbatim}[commandchars=\\\{\}]
\PYG{k+kn}{import}\PYG{+w}{ }\PYG{n+nn}{numpy}\PYG{+w}{ }\PYG{k}{as}\PYG{+w}{ }\PYG{n+nn}{np}
\PYG{k+kn}{import}\PYG{+w}{ }\PYG{n+nn}{sklearn}
\PYG{k+kn}{import}\PYG{+w}{ }\PYG{n+nn}{warnings}
\PYG{n}{warnings}\PYG{o}{.}\PYG{n}{filterwarnings}\PYG{p}{(}\PYG{l+s+s1}{\PYGZsq{}}\PYG{l+s+s1}{ignore}\PYG{l+s+s1}{\PYGZsq{}}\PYG{p}{)}
\PYG{k+kn}{from}\PYG{+w}{ }\PYG{n+nn}{tqdm}\PYG{n+nn}{.}\PYG{n+nn}{auto}\PYG{+w}{ }\PYG{k+kn}{import} \PYG{n}{tqdm}
\PYG{k+kn}{from}\PYG{+w}{ }\PYG{n+nn}{IPython}\PYG{n+nn}{.}\PYG{n+nn}{display}\PYG{+w}{ }\PYG{k+kn}{import} \PYG{n}{display}\PYG{p}{,} \PYG{n}{Math}
\PYG{k+kn}{from}\PYG{+w}{ }\PYG{n+nn}{numpyarray\PYGZus{}to\PYGZus{}latex}\PYG{n+nn}{.}\PYG{n+nn}{jupyter}\PYG{+w}{ }\PYG{k+kn}{import} \PYG{n}{to\PYGZus{}ltx}
\PYG{k+kn}{import}\PYG{+w}{ }\PYG{n+nn}{matplotlib}\PYG{n+nn}{.}\PYG{n+nn}{pyplot}\PYG{+w}{ }\PYG{k}{as}\PYG{+w}{ }\PYG{n+nn}{plt}
\PYG{n}{plt}\PYG{o}{.}\PYG{n}{style}\PYG{o}{.}\PYG{n}{use}\PYG{p}{(}\PYG{p}{[}\PYG{l+s+s1}{\PYGZsq{}}\PYG{l+s+s1}{seaborn\PYGZhy{}v0\PYGZus{}8\PYGZhy{}muted}\PYG{l+s+s1}{\PYGZsq{}}\PYG{p}{,} \PYG{l+s+s1}{\PYGZsq{}}\PYG{l+s+s1}{practicals.mplstyle}\PYG{l+s+s1}{\PYGZsq{}}\PYG{p}{]}\PYG{p}{)}
\PYG{k}{def}\PYG{+w}{ }\PYG{n+nf}{pprint}\PYG{p}{(}\PYG{o}{*}\PYG{n}{args}\PYG{p}{)}\PYG{p}{:}
    \PYG{n}{res} \PYG{o}{=} \PYG{l+s+s2}{\PYGZdq{}}\PYG{l+s+s2}{\PYGZdq{}}
    \PYG{k}{for} \PYG{n}{i} \PYG{o+ow}{in} \PYG{n}{args}\PYG{p}{:}
        \PYG{k}{if} \PYG{n+nb}{type}\PYG{p}{(}\PYG{n}{i}\PYG{p}{)} \PYG{o}{==} \PYG{n}{np}\PYG{o}{.}\PYG{n}{ndarray}\PYG{p}{:}
            \PYG{n}{res} \PYG{o}{+}\PYG{o}{=} \PYG{n}{to\PYGZus{}ltx}\PYG{p}{(}\PYG{n}{i}\PYG{p}{,} \PYG{n}{brackets}\PYG{o}{=}\PYG{l+s+s1}{\PYGZsq{}}\PYG{l+s+s1}{[]}\PYG{l+s+s1}{\PYGZsq{}}\PYG{p}{)}
        \PYG{k}{elif} \PYG{n+nb}{type}\PYG{p}{(}\PYG{n}{i}\PYG{p}{)} \PYG{o}{==} \PYG{n+nb}{str}\PYG{p}{:}
            \PYG{n}{res} \PYG{o}{+}\PYG{o}{=} \PYG{n}{i}
    \PYG{n}{display}\PYG{p}{(}\PYG{n}{Math}\PYG{p}{(}\PYG{n}{res}\PYG{p}{)}\PYG{p}{)}
\end{sphinxVerbatim}

\end{sphinxuseclass}\end{sphinxVerbatimInput}

\end{sphinxuseclass}

\section{Introduction}
\label{\detokenize{06:introduction}}
\sphinxAtStartPar
In classification problem, the target is a category (a discrete variable) and its conditional distribution is
\begin{equation*}
\begin{split}P({y}|x)\end{split}
\end{equation*}
\sphinxAtStartPar
where:
\begin{itemize}
\item {} 
\sphinxAtStartPar
\(\mathbf{y} \in \{ c_1,\dots, c_K\}\) represents the output variable (also called target)

\item {} 
\sphinxAtStartPar
\(x \in \mathbb{R}^n\) represents the vector of inputs (also called features)

\end{itemize}

\sphinxAtStartPar
In classification, the goal of learning is to return an estimator
\begin{equation*}
\begin{split}h(x,\alpha)\end{split}
\end{equation*}
\sphinxAtStartPar
of the conditional probabiity \(P({\mathbf y}|x)\), where \(\alpha\) denotes the set of parameters of the model \(h\).

\sphinxAtStartPar
The learning of the estimator is done  on the basis of an available input/output
training set \(D_N\) made of \(N\) observation pairs \((x_i,y_i)\) where \(x_i \in \mathbb{R}^n\).


\section{Data generation}
\label{\detokenize{06:data-generation}}
\sphinxAtStartPar
As seen in the theoretical course, there are multiple ways to generate datasets for binary classification problems.


\subsection{Generative way:}
\label{\detokenize{06:generative-way}}
\sphinxAtStartPar
First, generate the class \(y_i \in \{0, 1\}\), \(i \in \{1,...,N\}\), using the \sphinxstyleemphasis{a priori} probability, then generate the \(x\) values, given the \sphinxstylestrong{class conditional density \(p(x | \mathbf{y}=y_i)\)}.

\begin{sphinxuseclass}{cell}\begin{sphinxVerbatimInput}

\begin{sphinxuseclass}{cell_input}
\begin{sphinxVerbatim}[commandchars=\\\{\}]
\PYG{n}{np}\PYG{o}{.}\PYG{n}{random}\PYG{o}{.}\PYG{n}{seed}\PYG{p}{(}\PYG{l+m+mi}{1}\PYG{p}{)}
\end{sphinxVerbatim}

\end{sphinxuseclass}\end{sphinxVerbatimInput}

\end{sphinxuseclass}
\sphinxAtStartPar
If we name the \sphinxstyleemphasis{a priori} probability \(p_1 = P(\mathbf{y} = 1)\), we can generate the \(Y\) dataset:

\begin{sphinxuseclass}{cell}\begin{sphinxVerbatimInput}

\begin{sphinxuseclass}{cell_input}
\begin{sphinxVerbatim}[commandchars=\\\{\}]
\PYG{n}{N} \PYG{o}{=} \PYG{l+m+mi}{2000}
\PYG{n}{n} \PYG{o}{=} \PYG{l+m+mi}{2}
\PYG{n}{p1} \PYG{o}{=} \PYG{l+m+mf}{0.4}

\PYG{n}{Y} \PYG{o}{=} \PYG{n}{np}\PYG{o}{.}\PYG{n}{random}\PYG{o}{.}\PYG{n}{choice}\PYG{p}{(}\PYG{p}{[}\PYG{l+m+mi}{0}\PYG{p}{,} \PYG{l+m+mi}{1}\PYG{p}{]}\PYG{p}{,} \PYG{n}{size}\PYG{o}{=}\PYG{n}{N}\PYG{p}{,} \PYG{n}{p}\PYG{o}{=}\PYG{p}{[}\PYG{l+m+mi}{1}\PYG{o}{\PYGZhy{}}\PYG{n}{p1}\PYG{p}{,} \PYG{n}{p1}\PYG{p}{]}\PYG{p}{)}
\end{sphinxVerbatim}

\end{sphinxuseclass}\end{sphinxVerbatimInput}

\end{sphinxuseclass}
\sphinxAtStartPar
Then, we can generate the \(X\) dataset, following two multivariate normal distribution (recall practical 1). We define two distributions, for each of the two classes. Then, depending on the \(y\) class of each item, we sample the corresponding multivariate normal distribution.

\sphinxAtStartPar
Here, to make sure that \sphinxcode{\sphinxupquote{cov0}} and \sphinxcode{\sphinxupquote{cov1}} are defined positive, we generate a random \(A\) matrix, then \sphinxcode{\sphinxupquote{cov}} is computed as \(A^TA\).

\begin{sphinxuseclass}{cell}\begin{sphinxVerbatimInput}

\begin{sphinxuseclass}{cell_input}
\begin{sphinxVerbatim}[commandchars=\\\{\}]
\PYG{n}{A0} \PYG{o}{=} \PYG{n}{np}\PYG{o}{.}\PYG{n}{random}\PYG{o}{.}\PYG{n}{rand}\PYG{p}{(}\PYG{n}{n}\PYG{p}{,} \PYG{n}{n}\PYG{p}{)}
\PYG{n}{cov0} \PYG{o}{=} \PYG{n}{A0}\PYG{o}{.}\PYG{n}{T} \PYG{o}{@} \PYG{n}{A0}
\PYG{n}{mu0} \PYG{o}{=} \PYG{p}{(}\PYG{l+m+mi}{0}\PYG{p}{,} \PYG{l+m+mi}{0}\PYG{p}{)}

\PYG{n}{A1} \PYG{o}{=} \PYG{n}{np}\PYG{o}{.}\PYG{n}{random}\PYG{o}{.}\PYG{n}{rand}\PYG{p}{(}\PYG{n}{n}\PYG{p}{,} \PYG{n}{n}\PYG{p}{)}
\PYG{n}{cov1} \PYG{o}{=} \PYG{n}{A1}\PYG{o}{.}\PYG{n}{T} \PYG{o}{@} \PYG{n}{A1}
\PYG{n}{mu1} \PYG{o}{=} \PYG{p}{(}\PYG{l+m+mf}{0.75}\PYG{p}{,} \PYG{l+m+mi}{0}\PYG{p}{)}
\end{sphinxVerbatim}

\end{sphinxuseclass}\end{sphinxVerbatimInput}

\end{sphinxuseclass}
\begin{sphinxuseclass}{cell}\begin{sphinxVerbatimInput}

\begin{sphinxuseclass}{cell_input}
\begin{sphinxVerbatim}[commandchars=\\\{\}]
\PYG{n}{X} \PYG{o}{=} \PYG{n}{np}\PYG{o}{.}\PYG{n}{zeros}\PYG{p}{(}\PYG{p}{(}\PYG{n}{N}\PYG{p}{,} \PYG{n}{n}\PYG{p}{)}\PYG{p}{)}
\PYG{k}{for} \PYG{n}{i} \PYG{o+ow}{in} \PYG{n}{np}\PYG{o}{.}\PYG{n}{arange}\PYG{p}{(}\PYG{n}{N}\PYG{p}{)}\PYG{p}{:}
    \PYG{k}{if} \PYG{n}{Y}\PYG{p}{[}\PYG{n}{i}\PYG{p}{]} \PYG{o}{==} \PYG{l+m+mi}{1}\PYG{p}{:}
        \PYG{n}{X}\PYG{p}{[}\PYG{n}{i}\PYG{p}{,} \PYG{p}{:}\PYG{p}{]} \PYG{o}{=} \PYG{n}{np}\PYG{o}{.}\PYG{n}{random}\PYG{o}{.}\PYG{n}{multivariate\PYGZus{}normal}\PYG{p}{(}\PYG{n}{mu1}\PYG{p}{,} \PYG{n}{cov1}\PYG{p}{,} \PYG{n}{size}\PYG{o}{=}\PYG{l+m+mi}{1}\PYG{p}{)}
    \PYG{k}{else}\PYG{p}{:}
        \PYG{n}{X}\PYG{p}{[}\PYG{n}{i}\PYG{p}{,} \PYG{p}{:}\PYG{p}{]} \PYG{o}{=} \PYG{n}{np}\PYG{o}{.}\PYG{n}{random}\PYG{o}{.}\PYG{n}{multivariate\PYGZus{}normal}\PYG{p}{(}\PYG{n}{mu0}\PYG{p}{,} \PYG{n}{cov0}\PYG{p}{,} \PYG{n}{size}\PYG{o}{=}\PYG{l+m+mi}{1}\PYG{p}{)}
\end{sphinxVerbatim}

\end{sphinxuseclass}\end{sphinxVerbatimInput}

\end{sphinxuseclass}
\sphinxAtStartPar
Then, we plot the generated data.

\begin{sphinxuseclass}{cell}\begin{sphinxVerbatimInput}

\begin{sphinxuseclass}{cell_input}
\begin{sphinxVerbatim}[commandchars=\\\{\}]
\PYG{n}{I0} \PYG{o}{=} \PYG{n}{np}\PYG{o}{.}\PYG{n}{where}\PYG{p}{(}\PYG{n}{Y}\PYG{o}{==}\PYG{l+m+mi}{0}\PYG{p}{)}
\PYG{n}{I1}\PYG{o}{=}\PYG{n}{np}\PYG{o}{.}\PYG{n}{where}\PYG{p}{(}\PYG{n}{Y}\PYG{o}{==}\PYG{l+m+mi}{1}\PYG{p}{)}

\PYG{n}{plt}\PYG{o}{.}\PYG{n}{figure}\PYG{p}{(}\PYG{n}{figsize}\PYG{o}{=}\PYG{p}{(}\PYG{l+m+mi}{10}\PYG{p}{,} \PYG{l+m+mi}{4}\PYG{p}{)}\PYG{p}{)}
\PYG{n}{plt}\PYG{o}{.}\PYG{n}{scatter}\PYG{p}{(}\PYG{n}{X}\PYG{p}{[}\PYG{n}{I0}\PYG{p}{,} \PYG{l+m+mi}{0}\PYG{p}{]}\PYG{p}{,} \PYG{n}{X}\PYG{p}{[}\PYG{n}{I0}\PYG{p}{,} \PYG{l+m+mi}{1}\PYG{p}{]}\PYG{p}{,} \PYG{n}{marker}\PYG{o}{=}\PYG{l+s+s1}{\PYGZsq{}}\PYG{l+s+s1}{.}\PYG{l+s+s1}{\PYGZsq{}}\PYG{p}{,} \PYG{n}{label}\PYG{o}{=}\PYG{l+s+s2}{\PYGZdq{}}\PYG{l+s+s2}{Class 0}\PYG{l+s+s2}{\PYGZdq{}}\PYG{p}{)}
\PYG{n}{plt}\PYG{o}{.}\PYG{n}{scatter}\PYG{p}{(}\PYG{n}{X}\PYG{p}{[}\PYG{n}{I1}\PYG{p}{,} \PYG{l+m+mi}{0}\PYG{p}{]}\PYG{p}{,} \PYG{n}{X}\PYG{p}{[}\PYG{n}{I1}\PYG{p}{,} \PYG{l+m+mi}{1}\PYG{p}{]}\PYG{p}{,} \PYG{n}{marker}\PYG{o}{=}\PYG{l+s+s1}{\PYGZsq{}}\PYG{l+s+s1}{.}\PYG{l+s+s1}{\PYGZsq{}}\PYG{p}{,} \PYG{n}{label}\PYG{o}{=}\PYG{l+s+s2}{\PYGZdq{}}\PYG{l+s+s2}{Class 1}\PYG{l+s+s2}{\PYGZdq{}}\PYG{p}{)}
\PYG{n}{plt}\PYG{o}{.}\PYG{n}{grid}\PYG{p}{(}\PYG{l+s+s2}{\PYGZdq{}}\PYG{l+s+s2}{on}\PYG{l+s+s2}{\PYGZdq{}}\PYG{p}{)}
\PYG{n}{plt}\PYG{o}{.}\PYG{n}{legend}\PYG{p}{(}\PYG{p}{)}
\PYG{n}{plt}\PYG{o}{.}\PYG{n}{title}\PYG{p}{(}\PYG{l+s+s2}{\PYGZdq{}}\PYG{l+s+s2}{Generative process}\PYG{l+s+s2}{\PYGZdq{}}\PYG{p}{)}
\PYG{n}{plt}\PYG{o}{.}\PYG{n}{show}\PYG{p}{(}\PYG{p}{)}
\end{sphinxVerbatim}

\end{sphinxuseclass}\end{sphinxVerbatimInput}
\begin{sphinxVerbatimOutput}

\begin{sphinxuseclass}{cell_output}
\noindent\sphinxincludegraphics{{c5bbed2461626dcc1598b77b3bc805eced1f5f8e33f76f6bf69a0cb9f8576da1}.png}

\end{sphinxuseclass}\end{sphinxVerbatimOutput}

\end{sphinxuseclass}

\subsection{Discriminative way}
\label{\detokenize{06:discriminative-way}}
\sphinxAtStartPar
Otherwise, if we have the conditional distributions \(P(y|x)\), we can first sample \(x\), then determine the class \(y\).

\begin{sphinxuseclass}{cell}\begin{sphinxVerbatimInput}

\begin{sphinxuseclass}{cell_input}
\begin{sphinxVerbatim}[commandchars=\\\{\}]
\PYG{k+kn}{from}\PYG{+w}{ }\PYG{n+nn}{scipy}\PYG{n+nn}{.}\PYG{n+nn}{stats}\PYG{+w}{ }\PYG{k+kn}{import} \PYG{n}{multivariate\PYGZus{}normal}
\end{sphinxVerbatim}

\end{sphinxuseclass}\end{sphinxVerbatimInput}

\end{sphinxuseclass}
\sphinxAtStartPar
First, we sample \(x\) according to the probability density function:
\begin{equation*}
\begin{split}p(x)= p_1*p(x|y=1) + (1-p_1)*p(x|y=0)\end{split}
\end{equation*}
\sphinxAtStartPar
Those are mixture of gaussians (recall: practical 1).

\begin{sphinxuseclass}{cell}\begin{sphinxVerbatimInput}

\begin{sphinxuseclass}{cell_input}
\begin{sphinxVerbatim}[commandchars=\\\{\}]
\PYG{n}{X} \PYG{o}{=} \PYG{n}{np}\PYG{o}{.}\PYG{n}{zeros}\PYG{p}{(}\PYG{p}{(}\PYG{n}{N}\PYG{p}{,}\PYG{n}{n}\PYG{p}{)}\PYG{p}{)}
\PYG{k}{for} \PYG{n}{i} \PYG{o+ow}{in} \PYG{n}{np}\PYG{o}{.}\PYG{n}{arange}\PYG{p}{(}\PYG{n}{N}\PYG{p}{)}\PYG{p}{:}
    \PYG{k}{if} \PYG{n}{np}\PYG{o}{.}\PYG{n}{random}\PYG{o}{.}\PYG{n}{uniform}\PYG{p}{(}\PYG{p}{)}\PYG{o}{\PYGZlt{}}\PYG{n}{p1}\PYG{p}{:}
        \PYG{n}{X}\PYG{p}{[}\PYG{n}{i}\PYG{p}{,}\PYG{p}{:}\PYG{p}{]} \PYG{o}{=} \PYG{n}{np}\PYG{o}{.}\PYG{n}{random}\PYG{o}{.}\PYG{n}{multivariate\PYGZus{}normal}\PYG{p}{(}\PYG{n}{mu1}\PYG{p}{,} \PYG{n}{cov1}\PYG{p}{,} \PYG{n}{size}\PYG{o}{=}\PYG{l+m+mi}{1}\PYG{p}{)}
    \PYG{k}{else}\PYG{p}{:}
       \PYG{n}{X}\PYG{p}{[}\PYG{n}{i}\PYG{p}{,}\PYG{p}{:}\PYG{p}{]} \PYG{o}{=} \PYG{n}{np}\PYG{o}{.}\PYG{n}{random}\PYG{o}{.}\PYG{n}{multivariate\PYGZus{}normal}\PYG{p}{(}\PYG{n}{mu0}\PYG{p}{,} \PYG{n}{cov0}\PYG{p}{,} \PYG{n}{size}\PYG{o}{=}\PYG{l+m+mi}{1}\PYG{p}{)}  
\end{sphinxVerbatim}

\end{sphinxuseclass}\end{sphinxVerbatimInput}

\end{sphinxuseclass}
\sphinxAtStartPar
Then, using the Bayes theorem, we compute the \(y\) classes.

\sphinxAtStartPar
Recall the Bayes theorem
\begin{equation*}
\begin{split}P(\mathbf{y}=y|x) = \frac{p_\mathbf{x}(x|\mathbf{y}=y)P(\mathbf{y}=y)}{p_\mathbf{x}(x)}\end{split}
\end{equation*}
\sphinxAtStartPar
In order to compute the term \(p_\mathbf{x}(x|\mathbf{y}=y)\), we need to make an hypothesis on the distribution of \(\mathbf{x}\) according to \(\mathbf{y}=y\). Here, since we know that the input features follow a gaussian law (if \(\mathbf{y} = 0\), then \(\mathbf{x} \sim \mathcal{N}(\mu_0, \Sigma_0)\) and if \(\mathbf{y} = 1\), then \(\mathbf{x} \sim \mathcal{N}(\mu_1, \Sigma_1)\)), we can us to use the \sphinxcode{\sphinxupquote{pdf}} function from the \sphinxcode{\sphinxupquote{multivariate\_normal}} library already presented.

\sphinxAtStartPar
In the binary case, we can rewrite the formula as
\begin{equation*}
\begin{split}P(\mathbf{y}=0|x) = \frac{p_\mathbf{x}(x|\mathbf{y}=0)P(\mathbf{y}=0)}{p_{\mathbf{x}}(x|\mathbf{y}=0)P(\mathbf{y}=0)+p_{\mathbf{x}}(x|\mathbf{y}=1
)P(\mathbf{y}=1)}\end{split}
\end{equation*}
\sphinxAtStartPar
The previous formula is directly implemented by:

\begin{sphinxuseclass}{cell}\begin{sphinxVerbatimInput}

\begin{sphinxuseclass}{cell_input}
\begin{sphinxVerbatim}[commandchars=\\\{\}]
\PYG{k}{def}\PYG{+w}{ }\PYG{n+nf}{condprob}\PYG{p}{(}\PYG{n}{x}\PYG{p}{,} \PYG{n}{cov1}\PYG{p}{,} \PYG{n}{mu1}\PYG{p}{,} \PYG{n}{cov0}\PYG{p}{,} \PYG{n}{mu0}\PYG{p}{,} \PYG{n}{p1}\PYG{p}{)}\PYG{p}{:}
    \PYG{k}{return} \PYG{n}{multivariate\PYGZus{}normal}\PYG{o}{.}\PYG{n}{pdf}\PYG{p}{(}\PYG{n}{x}\PYG{p}{,} \PYG{n}{mean}\PYG{o}{=}\PYG{n}{mu1}\PYG{p}{,} \PYG{n}{cov}\PYG{o}{=}\PYG{n}{cov1}\PYG{p}{)}\PYG{o}{*}\PYG{n}{p1}\PYG{o}{/}\PYG{p}{(}\PYG{n}{multivariate\PYGZus{}normal}\PYG{o}{.}\PYG{n}{pdf}\PYG{p}{(}\PYG{n}{x}\PYG{p}{,} \PYG{n}{mean}\PYG{o}{=}\PYG{n}{mu1}\PYG{p}{,} \PYG{n}{cov}\PYG{o}{=}\PYG{n}{cov1}\PYG{p}{)}\PYG{o}{*}\PYG{n}{p1}\PYG{o}{+}\PYG{n}{multivariate\PYGZus{}normal}\PYG{o}{.}\PYG{n}{pdf}\PYG{p}{(}\PYG{n}{x}\PYG{p}{,} \PYG{n}{mean}\PYG{o}{=}\PYG{n}{mu0}\PYG{p}{,} \PYG{n}{cov}\PYG{o}{=}\PYG{n}{cov0}\PYG{p}{)}\PYG{o}{*}\PYG{p}{(}\PYG{l+m+mi}{1}\PYG{o}{\PYGZhy{}}\PYG{n}{p1}\PYG{p}{)}\PYG{p}{)}
\end{sphinxVerbatim}

\end{sphinxuseclass}\end{sphinxVerbatimInput}

\end{sphinxuseclass}
\sphinxAtStartPar
where \(P(\mathbf{y}=y|x)\) is computed by \sphinxcode{\sphinxupquote{condprob}}.

\begin{sphinxuseclass}{cell}\begin{sphinxVerbatimInput}

\begin{sphinxuseclass}{cell_input}
\begin{sphinxVerbatim}[commandchars=\\\{\}]
\PYG{n}{Y} \PYG{o}{=} \PYG{n}{np}\PYG{o}{.}\PYG{n}{zeros}\PYG{p}{(}\PYG{p}{(}\PYG{n}{N}\PYG{p}{,} \PYG{l+m+mi}{1}\PYG{p}{)}\PYG{p}{)}
\PYG{k}{for} \PYG{n}{i} \PYG{o+ow}{in} \PYG{n}{np}\PYG{o}{.}\PYG{n}{arange}\PYG{p}{(}\PYG{n}{N}\PYG{p}{)}\PYG{p}{:}
    \PYG{n}{p1i} \PYG{o}{=} \PYG{n}{condprob}\PYG{p}{(}\PYG{n}{X}\PYG{p}{[}\PYG{n}{i}\PYG{p}{,} \PYG{p}{:}\PYG{p}{]}\PYG{p}{,} \PYG{n}{cov1}\PYG{p}{,} \PYG{n}{mu1}\PYG{p}{,} \PYG{n}{cov0}\PYG{p}{,} \PYG{n}{mu0}\PYG{p}{,} \PYG{n}{p1}\PYG{p}{)}
    \PYG{n}{Y}\PYG{p}{[}\PYG{n}{i}\PYG{p}{]} \PYG{o}{=} \PYG{n}{np}\PYG{o}{.}\PYG{n}{random}\PYG{o}{.}\PYG{n}{choice}\PYG{p}{(}\PYG{p}{[}\PYG{l+m+mi}{0}\PYG{p}{,} \PYG{l+m+mi}{1}\PYG{p}{]}\PYG{p}{,} \PYG{n}{size}\PYG{o}{=}\PYG{l+m+mi}{1}\PYG{p}{,} \PYG{n}{p}\PYG{o}{=}\PYG{p}{[}\PYG{l+m+mi}{1}\PYG{o}{\PYGZhy{}}\PYG{n}{p1i}\PYG{p}{,} \PYG{n}{p1i}\PYG{p}{]}\PYG{p}{)}
\end{sphinxVerbatim}

\end{sphinxuseclass}\end{sphinxVerbatimInput}

\end{sphinxuseclass}
\begin{sphinxuseclass}{cell}\begin{sphinxVerbatimInput}

\begin{sphinxuseclass}{cell_input}
\begin{sphinxVerbatim}[commandchars=\\\{\}]
\PYG{n}{I0} \PYG{o}{=} \PYG{n}{np}\PYG{o}{.}\PYG{n}{where}\PYG{p}{(}\PYG{n}{Y}\PYG{o}{==}\PYG{l+m+mi}{0}\PYG{p}{)}
\PYG{n}{I1} \PYG{o}{=} \PYG{n}{np}\PYG{o}{.}\PYG{n}{where}\PYG{p}{(}\PYG{n}{Y}\PYG{o}{==}\PYG{l+m+mi}{1}\PYG{p}{)}

\PYG{n}{plt}\PYG{o}{.}\PYG{n}{figure}\PYG{p}{(}\PYG{n}{figsize}\PYG{o}{=}\PYG{p}{(}\PYG{l+m+mi}{10}\PYG{p}{,} \PYG{l+m+mi}{4}\PYG{p}{)}\PYG{p}{)}
\PYG{n}{plt}\PYG{o}{.}\PYG{n}{scatter}\PYG{p}{(}\PYG{n}{X}\PYG{p}{[}\PYG{n}{I0}\PYG{p}{,} \PYG{l+m+mi}{0}\PYG{p}{]}\PYG{p}{,} \PYG{n}{X}\PYG{p}{[}\PYG{n}{I0}\PYG{p}{,} \PYG{l+m+mi}{1}\PYG{p}{]}\PYG{p}{,} \PYG{n}{marker}\PYG{o}{=}\PYG{l+s+s1}{\PYGZsq{}}\PYG{l+s+s1}{.}\PYG{l+s+s1}{\PYGZsq{}}\PYG{p}{,} \PYG{n}{label}\PYG{o}{=}\PYG{l+s+s2}{\PYGZdq{}}\PYG{l+s+s2}{Class 0}\PYG{l+s+s2}{\PYGZdq{}}\PYG{p}{)}
\PYG{n}{plt}\PYG{o}{.}\PYG{n}{scatter}\PYG{p}{(}\PYG{n}{X}\PYG{p}{[}\PYG{n}{I1}\PYG{p}{,} \PYG{l+m+mi}{0}\PYG{p}{]}\PYG{p}{,} \PYG{n}{X}\PYG{p}{[}\PYG{n}{I1}\PYG{p}{,} \PYG{l+m+mi}{1}\PYG{p}{]}\PYG{p}{,}\PYG{n}{marker}\PYG{o}{=}\PYG{l+s+s1}{\PYGZsq{}}\PYG{l+s+s1}{.}\PYG{l+s+s1}{\PYGZsq{}}\PYG{p}{,} \PYG{n}{label}\PYG{o}{=}\PYG{l+s+s2}{\PYGZdq{}}\PYG{l+s+s2}{Class 1}\PYG{l+s+s2}{\PYGZdq{}}\PYG{p}{)}
\PYG{n}{plt}\PYG{o}{.}\PYG{n}{grid}\PYG{p}{(}\PYG{l+s+s2}{\PYGZdq{}}\PYG{l+s+s2}{on}\PYG{l+s+s2}{\PYGZdq{}}\PYG{p}{)}
\PYG{n}{plt}\PYG{o}{.}\PYG{n}{legend}\PYG{p}{(}\PYG{p}{)}
\PYG{n}{plt}\PYG{o}{.}\PYG{n}{title}\PYG{p}{(}\PYG{l+s+s2}{\PYGZdq{}}\PYG{l+s+s2}{Discriminative process}\PYG{l+s+s2}{\PYGZdq{}}\PYG{p}{)}
\PYG{n}{plt}\PYG{o}{.}\PYG{n}{show}\PYG{p}{(}\PYG{p}{)}
\end{sphinxVerbatim}

\end{sphinxuseclass}\end{sphinxVerbatimInput}
\begin{sphinxVerbatimOutput}

\begin{sphinxuseclass}{cell_output}
\noindent\sphinxincludegraphics{{45860421de2b1a645864ad49e6e8807dac6bf25662ea66e2ee2407783f6d68e5}.png}

\end{sphinxuseclass}\end{sphinxVerbatimOutput}

\end{sphinxuseclass}
\sphinxAtStartPar
Of course, the computed datasets follow the same distribution, no matter the method used.


\section{Exercises}
\label{\detokenize{06:exercises}}\begin{itemize}
\item {} 
\sphinxAtStartPar
For the lastly defined distributions, given three points \(q=\{(0, 0), (3, 0), (0, 2)\}\), compute \(P(\mathbf{y}| \mathbf{x}=q)\) for each class \(y\). Compute \(p(q|\mathbf{y}=y)\) for each class \(y\) and \(P(\mathbf{y}=y)\) for each class \(y\).

\item {} 
\sphinxAtStartPar
Can you adapt the two procedures to get higher\sphinxhyphen{}dimensional data (i.e. \(x\) is not in \(\mathbb{R}^2\) but in \(\mathbb{R}^n\)) ?

\item {} 
\sphinxAtStartPar
Can you adapt the two procedures to get more\sphinxhyphen{}classes data (i.e. \(y\) is not in \(\{0, 1\}\) but in \(\{1,...,c_{k}\}\)) ?

\end{itemize}


\subsection{Question 1}
\label{\detokenize{06:question-1}}
\sphinxAtStartPar
Using Bayes’ theorem:
\begin{equation*}
\begin{split}P(\mathbf{y}= y | \mathbf{x} = q) = \frac{p(q | \mathbf{y} = y) P(\mathbf{y} = y)}{p(q)}\end{split}
\end{equation*}
\sphinxAtStartPar
where the marginal likelihood \(p(q)\) is given by:
\begin{equation*}
\begin{split}p(q) = \sum_{y} p(q | \mathbf{y} = y) P(\mathbf{y} = y).\end{split}
\end{equation*}
\sphinxAtStartPar
The likelihood \(p(q | \mathbf{y} = y)\) follows the multivariate normal density function (available in \sphinxcode{\sphinxupquote{scipy}}):
\begin{equation*}
\begin{split}p(q | \mathbf{y} = y) = \frac{1}{(2\pi)^{d/2} |\Sigma_y|^{1/2}} \exp\left(-\frac{1}{2} (q - \mu_y)^T \Sigma_y^{-1} (q - \mu_y)\right),\end{split}
\end{equation*}
\sphinxAtStartPar
where \(d = 2\) is the dimensionality.

\begin{sphinxuseclass}{cell}\begin{sphinxVerbatimInput}

\begin{sphinxuseclass}{cell_input}
\begin{sphinxVerbatim}[commandchars=\\\{\}]
\PYG{k}{for} \PYG{n}{q} \PYG{o+ow}{in} \PYG{p}{[}\PYG{p}{(}\PYG{l+m+mi}{0}\PYG{p}{,} \PYG{l+m+mi}{0}\PYG{p}{)}\PYG{p}{,} \PYG{p}{(}\PYG{l+m+mi}{3}\PYG{p}{,} \PYG{l+m+mi}{0}\PYG{p}{)}\PYG{p}{,} \PYG{p}{(}\PYG{l+m+mi}{0}\PYG{p}{,} \PYG{l+m+mi}{2}\PYG{p}{)}\PYG{p}{]}\PYG{p}{:}
    \PYG{n}{pq} \PYG{o}{=} \PYG{n}{multivariate\PYGZus{}normal}\PYG{o}{.}\PYG{n}{pdf}\PYG{p}{(}\PYG{n}{q}\PYG{p}{,} \PYG{n}{mean}\PYG{o}{=}\PYG{n}{mu0}\PYG{p}{,} \PYG{n}{cov}\PYG{o}{=}\PYG{n}{cov0}\PYG{p}{)}\PYG{o}{*}\PYG{p}{(}\PYG{l+m+mi}{1}\PYG{o}{\PYGZhy{}}\PYG{n}{p1}\PYG{p}{)} \PYG{o}{+} \PYG{n}{multivariate\PYGZus{}normal}\PYG{o}{.}\PYG{n}{pdf}\PYG{p}{(}\PYG{n}{q}\PYG{p}{,} \PYG{n}{mean}\PYG{o}{=}\PYG{n}{mu1}\PYG{p}{,} \PYG{n}{cov}\PYG{o}{=}\PYG{n}{cov1}\PYG{p}{)}\PYG{o}{*}\PYG{n}{p1}
    \PYG{n}{pprint}\PYG{p}{(}\PYG{l+s+sa}{f}\PYG{l+s+s2}{\PYGZdq{}}\PYG{l+s+s2}{P(y=0|x=}\PYG{l+s+si}{\PYGZob{}}\PYG{n}{q}\PYG{l+s+si}{\PYGZcb{}}\PYG{l+s+s2}{) = }\PYG{l+s+si}{\PYGZob{}}\PYG{n}{multivariate\PYGZus{}normal}\PYG{o}{.}\PYG{n}{pdf}\PYG{p}{(}\PYG{n}{q}\PYG{p}{,}\PYG{+w}{ }\PYG{n}{mean}\PYG{o}{=}\PYG{n}{mu0}\PYG{p}{,}\PYG{+w}{ }\PYG{n}{cov}\PYG{o}{=}\PYG{n}{cov0}\PYG{p}{)}\PYG{o}{*}\PYG{p}{(}\PYG{l+m+mi}{1}\PYG{o}{\PYGZhy{}}\PYG{n}{p1}\PYG{p}{)}\PYG{o}{/}\PYG{n}{pq}\PYG{l+s+si}{:}\PYG{l+s+s2}{.4e}\PYG{l+s+si}{\PYGZcb{}}\PYG{l+s+s2}{\PYGZdq{}}\PYG{p}{)}

\PYG{n+nb}{print}\PYG{p}{(}\PYG{p}{)}
\PYG{k}{for} \PYG{n}{q} \PYG{o+ow}{in} \PYG{p}{[}\PYG{p}{(}\PYG{l+m+mi}{0}\PYG{p}{,} \PYG{l+m+mi}{0}\PYG{p}{)}\PYG{p}{,} \PYG{p}{(}\PYG{l+m+mi}{3}\PYG{p}{,} \PYG{l+m+mi}{0}\PYG{p}{)}\PYG{p}{,} \PYG{p}{(}\PYG{l+m+mi}{0}\PYG{p}{,} \PYG{l+m+mi}{2}\PYG{p}{)}\PYG{p}{]}\PYG{p}{:}
    \PYG{n}{pq} \PYG{o}{=} \PYG{n}{multivariate\PYGZus{}normal}\PYG{o}{.}\PYG{n}{pdf}\PYG{p}{(}\PYG{n}{q}\PYG{p}{,} \PYG{n}{mean}\PYG{o}{=}\PYG{n}{mu0}\PYG{p}{,} \PYG{n}{cov}\PYG{o}{=}\PYG{n}{cov0}\PYG{p}{)}\PYG{o}{*}\PYG{p}{(}\PYG{l+m+mi}{1}\PYG{o}{\PYGZhy{}}\PYG{n}{p1}\PYG{p}{)} \PYG{o}{+} \PYG{n}{multivariate\PYGZus{}normal}\PYG{o}{.}\PYG{n}{pdf}\PYG{p}{(}\PYG{n}{q}\PYG{p}{,} \PYG{n}{mean}\PYG{o}{=}\PYG{n}{mu1}\PYG{p}{,} \PYG{n}{cov}\PYG{o}{=}\PYG{n}{cov1}\PYG{p}{)}\PYG{o}{*}\PYG{n}{p1}
    \PYG{n}{pprint}\PYG{p}{(}\PYG{l+s+sa}{f}\PYG{l+s+s2}{\PYGZdq{}}\PYG{l+s+s2}{P(y=1|x=}\PYG{l+s+si}{\PYGZob{}}\PYG{n}{q}\PYG{l+s+si}{\PYGZcb{}}\PYG{l+s+s2}{) = }\PYG{l+s+si}{\PYGZob{}}\PYG{n}{multivariate\PYGZus{}normal}\PYG{o}{.}\PYG{n}{pdf}\PYG{p}{(}\PYG{n}{q}\PYG{p}{,}\PYG{+w}{ }\PYG{n}{mean}\PYG{o}{=}\PYG{n}{mu1}\PYG{p}{,}\PYG{+w}{ }\PYG{n}{cov}\PYG{o}{=}\PYG{n}{cov1}\PYG{p}{)}\PYG{o}{*}\PYG{n}{p1}\PYG{o}{/}\PYG{n}{pq}\PYG{l+s+si}{:}\PYG{l+s+s2}{.4e}\PYG{l+s+si}{\PYGZcb{}}\PYG{l+s+s2}{\PYGZdq{}}\PYG{p}{)}

\PYG{n+nb}{print}\PYG{p}{(}\PYG{p}{)}
\PYG{k}{for} \PYG{n}{q} \PYG{o+ow}{in} \PYG{p}{[}\PYG{p}{(}\PYG{l+m+mi}{0}\PYG{p}{,} \PYG{l+m+mi}{0}\PYG{p}{)}\PYG{p}{,} \PYG{p}{(}\PYG{l+m+mi}{3}\PYG{p}{,} \PYG{l+m+mi}{0}\PYG{p}{)}\PYG{p}{,} \PYG{p}{(}\PYG{l+m+mi}{0}\PYG{p}{,} \PYG{l+m+mi}{2}\PYG{p}{)}\PYG{p}{]}\PYG{p}{:}
    \PYG{n}{pprint}\PYG{p}{(}\PYG{l+s+sa}{f}\PYG{l+s+s2}{\PYGZdq{}}\PYG{l+s+s2}{P(x=}\PYG{l+s+si}{\PYGZob{}}\PYG{n}{q}\PYG{l+s+si}{\PYGZcb{}}\PYG{l+s+s2}{|y=0) = }\PYG{l+s+si}{\PYGZob{}}\PYG{n}{multivariate\PYGZus{}normal}\PYG{o}{.}\PYG{n}{pdf}\PYG{p}{(}\PYG{n}{q}\PYG{p}{,}\PYG{+w}{ }\PYG{n}{mean}\PYG{o}{=}\PYG{n}{mu0}\PYG{p}{,}\PYG{+w}{ }\PYG{n}{cov}\PYG{o}{=}\PYG{n}{cov0}\PYG{p}{)}\PYG{l+s+si}{:}\PYG{l+s+s2}{.4e}\PYG{l+s+si}{\PYGZcb{}}\PYG{l+s+s2}{\PYGZdq{}}\PYG{p}{)}

\PYG{n+nb}{print}\PYG{p}{(}\PYG{p}{)}
\PYG{k}{for} \PYG{n}{q} \PYG{o+ow}{in} \PYG{p}{[}\PYG{p}{(}\PYG{l+m+mi}{0}\PYG{p}{,} \PYG{l+m+mi}{0}\PYG{p}{)}\PYG{p}{,} \PYG{p}{(}\PYG{l+m+mi}{3}\PYG{p}{,} \PYG{l+m+mi}{0}\PYG{p}{)}\PYG{p}{,} \PYG{p}{(}\PYG{l+m+mi}{0}\PYG{p}{,} \PYG{l+m+mi}{2}\PYG{p}{)}\PYG{p}{]}\PYG{p}{:}
    \PYG{n}{pprint}\PYG{p}{(}\PYG{l+s+sa}{f}\PYG{l+s+s2}{\PYGZdq{}}\PYG{l+s+s2}{P(x=}\PYG{l+s+si}{\PYGZob{}}\PYG{n}{q}\PYG{l+s+si}{\PYGZcb{}}\PYG{l+s+s2}{|y=1) = }\PYG{l+s+si}{\PYGZob{}}\PYG{n}{multivariate\PYGZus{}normal}\PYG{o}{.}\PYG{n}{pdf}\PYG{p}{(}\PYG{n}{q}\PYG{p}{,}\PYG{+w}{ }\PYG{n}{mean}\PYG{o}{=}\PYG{n}{mu1}\PYG{p}{,}\PYG{+w}{ }\PYG{n}{cov}\PYG{o}{=}\PYG{n}{cov1}\PYG{p}{)}\PYG{l+s+si}{:}\PYG{l+s+s2}{.4e}\PYG{l+s+si}{\PYGZcb{}}\PYG{l+s+s2}{\PYGZdq{}}\PYG{p}{)}

\PYG{n+nb}{print}\PYG{p}{(}\PYG{p}{)}
\PYG{n}{pprint}\PYG{p}{(}\PYG{l+s+sa}{f}\PYG{l+s+s2}{\PYGZdq{}}\PYG{l+s+s2}{P(y=0) = }\PYG{l+s+si}{\PYGZob{}}\PYG{l+m+mi}{1}\PYG{o}{\PYGZhy{}}\PYG{n}{p1}\PYG{l+s+si}{\PYGZcb{}}\PYG{l+s+s2}{\PYGZdq{}}\PYG{p}{)}
\PYG{n}{pprint}\PYG{p}{(}\PYG{l+s+sa}{f}\PYG{l+s+s2}{\PYGZdq{}}\PYG{l+s+s2}{P(y=1) = }\PYG{l+s+si}{\PYGZob{}}\PYG{n}{p1}\PYG{l+s+si}{\PYGZcb{}}\PYG{l+s+s2}{\PYGZdq{}}\PYG{p}{)}
\end{sphinxVerbatim}

\end{sphinxuseclass}\end{sphinxVerbatimInput}
\begin{sphinxVerbatimOutput}

\begin{sphinxuseclass}{cell_output}\begin{equation*}
\begin{split}\displaystyle P(y=0|x=(0, 0)) = 7.3086e-01\end{split}
\end{equation*}\begin{equation*}
\begin{split}\displaystyle P(y=0|x=(3, 0)) = 1.3140e-09\end{split}
\end{equation*}\begin{equation*}
\begin{split}\displaystyle P(y=0|x=(0, 2)) = 9.7560e-01\end{split}
\end{equation*}
\begin{sphinxVerbatim}[commandchars=\\\{\}]

\end{sphinxVerbatim}
\begin{equation*}
\begin{split}\displaystyle P(y=1|x=(0, 0)) = 2.6914e-01\end{split}
\end{equation*}\begin{equation*}
\begin{split}\displaystyle P(y=1|x=(3, 0)) = 1.0000e+00\end{split}
\end{equation*}\begin{equation*}
\begin{split}\displaystyle P(y=1|x=(0, 2)) = 2.4404e-02\end{split}
\end{equation*}
\begin{sphinxVerbatim}[commandchars=\\\{\}]

\end{sphinxVerbatim}
\begin{equation*}
\begin{split}\displaystyle P(x=(0, 0)|y=0) = 4.1069e-01\end{split}
\end{equation*}\begin{equation*}
\begin{split}\displaystyle P(x=(3, 0)|y=0) = 6.5642e-12\end{split}
\end{equation*}\begin{equation*}
\begin{split}\displaystyle P(x=(0, 2)|y=0) = 3.5164e-05\end{split}
\end{equation*}
\begin{sphinxVerbatim}[commandchars=\\\{\}]

\end{sphinxVerbatim}
\begin{equation*}
\begin{split}\displaystyle P(x=(0, 0)|y=1) = 2.2685e-01\end{split}
\end{equation*}\begin{equation*}
\begin{split}\displaystyle P(x=(3, 0)|y=1) = 7.4934e-03\end{split}
\end{equation*}\begin{equation*}
\begin{split}\displaystyle P(x=(0, 2)|y=1) = 1.3194e-06\end{split}
\end{equation*}
\begin{sphinxVerbatim}[commandchars=\\\{\}]

\end{sphinxVerbatim}
\begin{equation*}
\begin{split}\displaystyle P(y=0) = 0.6\end{split}
\end{equation*}\begin{equation*}
\begin{split}\displaystyle P(y=1) = 0.4\end{split}
\end{equation*}
\end{sphinxuseclass}\end{sphinxVerbatimOutput}

\end{sphinxuseclass}

\subsection{Question 2}
\label{\detokenize{06:question-2}}
\begin{sphinxuseclass}{cell}\begin{sphinxVerbatimInput}

\begin{sphinxuseclass}{cell_input}
\begin{sphinxVerbatim}[commandchars=\\\{\}]
\PYG{n}{n} \PYG{o}{=} \PYG{l+m+mi}{5}
\PYG{n}{Y} \PYG{o}{=} \PYG{n}{np}\PYG{o}{.}\PYG{n}{random}\PYG{o}{.}\PYG{n}{choice}\PYG{p}{(}\PYG{p}{[}\PYG{l+m+mi}{0}\PYG{p}{,} \PYG{l+m+mi}{1}\PYG{p}{]}\PYG{p}{,} \PYG{n}{size}\PYG{o}{=}\PYG{n}{N}\PYG{p}{,} \PYG{n}{p}\PYG{o}{=}\PYG{p}{[}\PYG{l+m+mi}{1}\PYG{o}{\PYGZhy{}}\PYG{n}{p1}\PYG{p}{,} \PYG{n}{p1}\PYG{p}{]}\PYG{p}{)}

\PYG{n}{mu0} \PYG{o}{=} \PYG{n}{np}\PYG{o}{.}\PYG{n}{zeros}\PYG{p}{(}\PYG{n}{n}\PYG{p}{)}
\PYG{n}{A0} \PYG{o}{=} \PYG{n}{np}\PYG{o}{.}\PYG{n}{random}\PYG{o}{.}\PYG{n}{rand}\PYG{p}{(}\PYG{n}{n}\PYG{p}{,} \PYG{n}{n}\PYG{p}{)}
\PYG{n}{cov0} \PYG{o}{=} \PYG{n}{A0}\PYG{o}{.}\PYG{n}{T} \PYG{o}{@} \PYG{n}{A0}

\PYG{n}{mu1} \PYG{o}{=} \PYG{n}{np}\PYG{o}{.}\PYG{n}{array}\PYG{p}{(}\PYG{p}{[}\PYG{l+m+mf}{0.75}\PYG{p}{]} \PYG{o}{+} \PYG{p}{[}\PYG{l+m+mi}{0}\PYG{p}{]} \PYG{o}{*} \PYG{p}{(}\PYG{n}{n} \PYG{o}{\PYGZhy{}} \PYG{l+m+mi}{1}\PYG{p}{)}\PYG{p}{)}
\PYG{n}{A1} \PYG{o}{=} \PYG{n}{np}\PYG{o}{.}\PYG{n}{random}\PYG{o}{.}\PYG{n}{rand}\PYG{p}{(}\PYG{n}{n}\PYG{p}{,} \PYG{n}{n}\PYG{p}{)}
\PYG{n}{cov1} \PYG{o}{=} \PYG{n}{A1}\PYG{o}{.}\PYG{n}{T} \PYG{o}{@} \PYG{n}{A1}

\PYG{n}{X} \PYG{o}{=} \PYG{n}{np}\PYG{o}{.}\PYG{n}{zeros}\PYG{p}{(}\PYG{p}{(}\PYG{n}{N}\PYG{p}{,} \PYG{n}{n}\PYG{p}{)}\PYG{p}{)}
\PYG{k}{for} \PYG{n}{i} \PYG{o+ow}{in} \PYG{n+nb}{range}\PYG{p}{(}\PYG{n}{N}\PYG{p}{)}\PYG{p}{:}
    \PYG{k}{if} \PYG{n}{Y}\PYG{p}{[}\PYG{n}{i}\PYG{p}{]} \PYG{o}{==} \PYG{l+m+mi}{1}\PYG{p}{:}
        \PYG{n}{X}\PYG{p}{[}\PYG{n}{i}\PYG{p}{,} \PYG{p}{:}\PYG{p}{]} \PYG{o}{=} \PYG{n}{np}\PYG{o}{.}\PYG{n}{random}\PYG{o}{.}\PYG{n}{multivariate\PYGZus{}normal}\PYG{p}{(}\PYG{n}{mu1}\PYG{p}{,} \PYG{n}{cov1}\PYG{p}{,} \PYG{n}{size}\PYG{o}{=}\PYG{l+m+mi}{1}\PYG{p}{)}
    \PYG{k}{else}\PYG{p}{:}
        \PYG{n}{X}\PYG{p}{[}\PYG{n}{i}\PYG{p}{,} \PYG{p}{:}\PYG{p}{]} \PYG{o}{=} \PYG{n}{np}\PYG{o}{.}\PYG{n}{random}\PYG{o}{.}\PYG{n}{multivariate\PYGZus{}normal}\PYG{p}{(}\PYG{n}{mu0}\PYG{p}{,} \PYG{n}{cov0}\PYG{p}{,} \PYG{n}{size}\PYG{o}{=}\PYG{l+m+mi}{1}\PYG{p}{)}
\end{sphinxVerbatim}

\end{sphinxuseclass}\end{sphinxVerbatimInput}

\end{sphinxuseclass}

\subsection{Question 3}
\label{\detokenize{06:question-3}}
\begin{sphinxuseclass}{cell}\begin{sphinxVerbatimInput}

\begin{sphinxuseclass}{cell_input}
\begin{sphinxVerbatim}[commandchars=\\\{\}]
\PYG{n}{np}\PYG{o}{.}\PYG{n}{random}\PYG{o}{.}\PYG{n}{seed}\PYG{p}{(}\PYG{l+m+mi}{1}\PYG{p}{)}
\PYG{n}{n} \PYG{o}{=} \PYG{l+m+mi}{2}
\PYG{n}{c\PYGZus{}k} \PYG{o}{=} \PYG{l+m+mi}{4}

\PYG{n}{class\PYGZus{}probs} \PYG{o}{=} \PYG{p}{[}\PYG{l+m+mf}{0.2}\PYG{p}{,} \PYG{l+m+mf}{0.3}\PYG{p}{,} \PYG{l+m+mf}{0.35}\PYG{p}{,} \PYG{l+m+mf}{0.15}\PYG{p}{]}
\PYG{n}{Y} \PYG{o}{=} \PYG{n}{np}\PYG{o}{.}\PYG{n}{random}\PYG{o}{.}\PYG{n}{choice}\PYG{p}{(}\PYG{n}{np}\PYG{o}{.}\PYG{n}{arange}\PYG{p}{(}\PYG{n}{c\PYGZus{}k}\PYG{p}{)}\PYG{p}{,} \PYG{n}{size}\PYG{o}{=}\PYG{n}{N}\PYG{p}{,} \PYG{n}{p}\PYG{o}{=}\PYG{n}{class\PYGZus{}probs}\PYG{p}{)}

\PYG{n}{mus} \PYG{o}{=} \PYG{p}{[}\PYG{n}{np}\PYG{o}{.}\PYG{n}{random}\PYG{o}{.}\PYG{n}{uniform}\PYG{p}{(}\PYG{o}{\PYGZhy{}}\PYG{l+m+mi}{1}\PYG{p}{,} \PYG{l+m+mi}{1}\PYG{p}{,} \PYG{n}{size}\PYG{o}{=}\PYG{n}{n}\PYG{p}{)} \PYG{k}{for} \PYG{n}{\PYGZus{}} \PYG{o+ow}{in} \PYG{n+nb}{range}\PYG{p}{(}\PYG{n}{c\PYGZus{}k}\PYG{p}{)}\PYG{p}{]}
\PYG{n}{covs} \PYG{o}{=} \PYG{p}{[}\PYG{n}{np}\PYG{o}{.}\PYG{n}{random}\PYG{o}{.}\PYG{n}{rand}\PYG{p}{(}\PYG{n}{n}\PYG{p}{,} \PYG{n}{n}\PYG{p}{)} \PYG{k}{for} \PYG{n}{\PYGZus{}} \PYG{o+ow}{in} \PYG{n+nb}{range}\PYG{p}{(}\PYG{n}{c\PYGZus{}k}\PYG{p}{)}\PYG{p}{]}
\PYG{n}{covs} \PYG{o}{=} \PYG{p}{[}\PYG{n}{C}\PYG{o}{.}\PYG{n}{T} \PYG{o}{@} \PYG{n}{C} \PYG{k}{for} \PYG{n}{C} \PYG{o+ow}{in} \PYG{n}{covs}\PYG{p}{]}

\PYG{n}{X} \PYG{o}{=} \PYG{n}{np}\PYG{o}{.}\PYG{n}{zeros}\PYG{p}{(}\PYG{p}{(}\PYG{n}{N}\PYG{p}{,} \PYG{n}{n}\PYG{p}{)}\PYG{p}{)}
\PYG{k}{for} \PYG{n}{i} \PYG{o+ow}{in} \PYG{n+nb}{range}\PYG{p}{(}\PYG{n}{N}\PYG{p}{)}\PYG{p}{:}
    \PYG{n+nb+bp}{cls} \PYG{o}{=} \PYG{n}{Y}\PYG{p}{[}\PYG{n}{i}\PYG{p}{]}
    \PYG{n}{X}\PYG{p}{[}\PYG{n}{i}\PYG{p}{,} \PYG{p}{:}\PYG{p}{]} \PYG{o}{=} \PYG{n}{np}\PYG{o}{.}\PYG{n}{random}\PYG{o}{.}\PYG{n}{multivariate\PYGZus{}normal}\PYG{p}{(}\PYG{n}{mus}\PYG{p}{[}\PYG{n+nb+bp}{cls}\PYG{p}{]}\PYG{p}{,} \PYG{n}{covs}\PYG{p}{[}\PYG{n+nb+bp}{cls}\PYG{p}{]}\PYG{p}{,} \PYG{n}{size}\PYG{o}{=}\PYG{l+m+mi}{1}\PYG{p}{)}
\end{sphinxVerbatim}

\end{sphinxuseclass}\end{sphinxVerbatimInput}

\end{sphinxuseclass}
\begin{sphinxuseclass}{cell}\begin{sphinxVerbatimInput}

\begin{sphinxuseclass}{cell_input}
\begin{sphinxVerbatim}[commandchars=\\\{\}]
\PYG{n}{Is} \PYG{o}{=} \PYG{p}{[}\PYG{n}{np}\PYG{o}{.}\PYG{n}{where}\PYG{p}{(}\PYG{n}{Y}\PYG{o}{==}\PYG{n}{i}\PYG{p}{)} \PYG{k}{for} \PYG{n}{i} \PYG{o+ow}{in} \PYG{n+nb}{range}\PYG{p}{(}\PYG{n}{c\PYGZus{}k}\PYG{p}{)}\PYG{p}{]}

\PYG{n}{plt}\PYG{o}{.}\PYG{n}{figure}\PYG{p}{(}\PYG{n}{figsize}\PYG{o}{=}\PYG{p}{(}\PYG{l+m+mi}{10}\PYG{p}{,} \PYG{l+m+mi}{4}\PYG{p}{)}\PYG{p}{)}
\PYG{k}{for} \PYG{n}{i} \PYG{o+ow}{in} \PYG{n+nb}{range}\PYG{p}{(}\PYG{n}{c\PYGZus{}k}\PYG{p}{)}\PYG{p}{:}
    \PYG{n}{plt}\PYG{o}{.}\PYG{n}{scatter}\PYG{p}{(}\PYG{n}{X}\PYG{p}{[}\PYG{n}{Is}\PYG{p}{[}\PYG{n}{i}\PYG{p}{]}\PYG{p}{,} \PYG{l+m+mi}{0}\PYG{p}{]}\PYG{p}{,} \PYG{n}{X}\PYG{p}{[}\PYG{n}{Is}\PYG{p}{[}\PYG{n}{i}\PYG{p}{]}\PYG{p}{,} \PYG{l+m+mi}{1}\PYG{p}{]}\PYG{p}{,} \PYG{n}{marker}\PYG{o}{=}\PYG{l+s+s1}{\PYGZsq{}}\PYG{l+s+s1}{.}\PYG{l+s+s1}{\PYGZsq{}}\PYG{p}{,} \PYG{n}{label}\PYG{o}{=}\PYG{l+s+sa}{f}\PYG{l+s+s2}{\PYGZdq{}}\PYG{l+s+s2}{Class }\PYG{l+s+si}{\PYGZob{}}\PYG{n}{i}\PYG{l+s+si}{\PYGZcb{}}\PYG{l+s+s2}{\PYGZdq{}}\PYG{p}{)}
\PYG{n}{plt}\PYG{o}{.}\PYG{n}{legend}\PYG{p}{(}\PYG{p}{)}
\PYG{n}{plt}\PYG{o}{.}\PYG{n}{grid}\PYG{p}{(}\PYG{l+s+s2}{\PYGZdq{}}\PYG{l+s+s2}{on}\PYG{l+s+s2}{\PYGZdq{}}\PYG{p}{)}
\PYG{n}{plt}\PYG{o}{.}\PYG{n}{title}\PYG{p}{(}\PYG{l+s+s2}{\PYGZdq{}}\PYG{l+s+s2}{More classes}\PYG{l+s+s2}{\PYGZdq{}}\PYG{p}{)}
\PYG{n}{plt}\PYG{o}{.}\PYG{n}{show}\PYG{p}{(}\PYG{p}{)}
\end{sphinxVerbatim}

\end{sphinxuseclass}\end{sphinxVerbatimInput}
\begin{sphinxVerbatimOutput}

\begin{sphinxuseclass}{cell_output}
\noindent\sphinxincludegraphics{{d43836de913a852e5180d8a560edd8d9a882d5a60848a61e18b5d8a29750e4e9}.png}

\end{sphinxuseclass}\end{sphinxVerbatimOutput}

\end{sphinxuseclass}
\sphinxAtStartPar
Finally we regenerate binary data for the rest of the practical:

\begin{sphinxuseclass}{cell}\begin{sphinxVerbatimInput}

\begin{sphinxuseclass}{cell_input}
\begin{sphinxVerbatim}[commandchars=\\\{\}]
\PYG{n}{np}\PYG{o}{.}\PYG{n}{random}\PYG{o}{.}\PYG{n}{seed}\PYG{p}{(}\PYG{l+m+mi}{1}\PYG{p}{)}
\PYG{n}{N} \PYG{o}{=} \PYG{l+m+mi}{4000}
\PYG{n}{N\PYGZus{}test} \PYG{o}{=} \PYG{l+m+mi}{1000}
\PYG{n}{n} \PYG{o}{=} \PYG{l+m+mi}{2}
\PYG{n}{p1} \PYG{o}{=} \PYG{l+m+mf}{0.4}

\PYG{n}{Y} \PYG{o}{=} \PYG{n}{np}\PYG{o}{.}\PYG{n}{random}\PYG{o}{.}\PYG{n}{choice}\PYG{p}{(}\PYG{p}{[}\PYG{l+m+mi}{0}\PYG{p}{,} \PYG{l+m+mi}{1}\PYG{p}{]}\PYG{p}{,} \PYG{n}{size}\PYG{o}{=}\PYG{n}{N}\PYG{p}{,} \PYG{n}{p}\PYG{o}{=}\PYG{p}{[}\PYG{l+m+mi}{1}\PYG{o}{\PYGZhy{}}\PYG{n}{p1}\PYG{p}{,} \PYG{n}{p1}\PYG{p}{]}\PYG{p}{)}
\PYG{n}{A0} \PYG{o}{=} \PYG{n}{np}\PYG{o}{.}\PYG{n}{array}\PYG{p}{(}\PYG{p}{[}\PYG{p}{[}\PYG{l+m+mf}{.45}\PYG{p}{,} \PYG{l+m+mf}{.5}\PYG{p}{]}\PYG{p}{,} \PYG{p}{[}\PYG{l+m+mf}{.05}\PYG{p}{,} \PYG{l+m+mf}{.35}\PYG{p}{]}\PYG{p}{]}\PYG{p}{)}
\PYG{n}{cov0} \PYG{o}{=} \PYG{n}{A0}\PYG{o}{.}\PYG{n}{T} \PYG{o}{@} \PYG{n}{A0}
\PYG{n}{mu0} \PYG{o}{=} \PYG{p}{(}\PYG{l+m+mi}{0}\PYG{p}{,} \PYG{l+m+mi}{0}\PYG{p}{)}

\PYG{n}{A1} \PYG{o}{=}  \PYG{n}{np}\PYG{o}{.}\PYG{n}{array}\PYG{p}{(}\PYG{p}{[}\PYG{p}{[}\PYG{l+m+mf}{.32}\PYG{p}{,} \PYG{o}{\PYGZhy{}}\PYG{l+m+mf}{.4}\PYG{p}{]}\PYG{p}{,} \PYG{p}{[}\PYG{o}{\PYGZhy{}}\PYG{l+m+mf}{.4}\PYG{p}{,} \PYG{l+m+mf}{.32}\PYG{p}{]}\PYG{p}{]}\PYG{p}{)}
\PYG{n}{cov1} \PYG{o}{=} \PYG{n}{A1}\PYG{o}{.}\PYG{n}{T} \PYG{o}{@} \PYG{n}{A1}
\PYG{n}{mu1} \PYG{o}{=} \PYG{p}{(}\PYG{l+m+mf}{.5}\PYG{p}{,} \PYG{l+m+mi}{0}\PYG{p}{)}
\PYG{n}{X} \PYG{o}{=} \PYG{n}{np}\PYG{o}{.}\PYG{n}{zeros}\PYG{p}{(}\PYG{p}{(}\PYG{n}{N}\PYG{p}{,} \PYG{n}{n}\PYG{p}{)}\PYG{p}{)}
\PYG{k}{for} \PYG{n}{i} \PYG{o+ow}{in} \PYG{n}{np}\PYG{o}{.}\PYG{n}{arange}\PYG{p}{(}\PYG{n}{N}\PYG{p}{)}\PYG{p}{:}
    \PYG{k}{if} \PYG{n}{Y}\PYG{p}{[}\PYG{n}{i}\PYG{p}{]} \PYG{o}{==} \PYG{l+m+mi}{1}\PYG{p}{:}
        \PYG{n}{X}\PYG{p}{[}\PYG{n}{i}\PYG{p}{,} \PYG{p}{:}\PYG{p}{]} \PYG{o}{=} \PYG{n}{np}\PYG{o}{.}\PYG{n}{random}\PYG{o}{.}\PYG{n}{multivariate\PYGZus{}normal}\PYG{p}{(}\PYG{n}{mu1}\PYG{p}{,} \PYG{n}{cov1}\PYG{p}{,} \PYG{n}{size}\PYG{o}{=}\PYG{l+m+mi}{1}\PYG{p}{)}
    \PYG{k}{else}\PYG{p}{:}
        \PYG{n}{X}\PYG{p}{[}\PYG{n}{i}\PYG{p}{,} \PYG{p}{:}\PYG{p}{]} \PYG{o}{=} \PYG{n}{np}\PYG{o}{.}\PYG{n}{random}\PYG{o}{.}\PYG{n}{multivariate\PYGZus{}normal}\PYG{p}{(}\PYG{n}{mu0}\PYG{p}{,} \PYG{n}{cov0}\PYG{p}{,} \PYG{n}{size}\PYG{o}{=}\PYG{l+m+mi}{1}\PYG{p}{)}

\PYG{n}{Y\PYGZus{}test} \PYG{o}{=} \PYG{n}{np}\PYG{o}{.}\PYG{n}{random}\PYG{o}{.}\PYG{n}{choice}\PYG{p}{(}\PYG{p}{[}\PYG{l+m+mi}{0}\PYG{p}{,} \PYG{l+m+mi}{1}\PYG{p}{]}\PYG{p}{,} \PYG{n}{size}\PYG{o}{=}\PYG{n}{N\PYGZus{}test}\PYG{p}{,} \PYG{n}{p}\PYG{o}{=}\PYG{p}{[}\PYG{l+m+mi}{1}\PYG{o}{\PYGZhy{}}\PYG{n}{p1}\PYG{p}{,} \PYG{n}{p1}\PYG{p}{]}\PYG{p}{)}
\PYG{n}{X\PYGZus{}test} \PYG{o}{=} \PYG{n}{np}\PYG{o}{.}\PYG{n}{zeros}\PYG{p}{(}\PYG{p}{(}\PYG{n}{N\PYGZus{}test}\PYG{p}{,} \PYG{n}{n}\PYG{p}{)}\PYG{p}{)}
\PYG{k}{for} \PYG{n}{i} \PYG{o+ow}{in} \PYG{n}{np}\PYG{o}{.}\PYG{n}{arange}\PYG{p}{(}\PYG{n}{N\PYGZus{}test}\PYG{p}{)}\PYG{p}{:}
    \PYG{k}{if} \PYG{n}{Y\PYGZus{}test}\PYG{p}{[}\PYG{n}{i}\PYG{p}{]} \PYG{o}{==} \PYG{l+m+mi}{1}\PYG{p}{:}
        \PYG{n}{X\PYGZus{}test}\PYG{p}{[}\PYG{n}{i}\PYG{p}{,} \PYG{p}{:}\PYG{p}{]} \PYG{o}{=} \PYG{n}{np}\PYG{o}{.}\PYG{n}{random}\PYG{o}{.}\PYG{n}{multivariate\PYGZus{}normal}\PYG{p}{(}\PYG{n}{mu1}\PYG{p}{,} \PYG{n}{cov1}\PYG{p}{,} \PYG{n}{size}\PYG{o}{=}\PYG{l+m+mi}{1}\PYG{p}{)}
    \PYG{k}{else}\PYG{p}{:}
        \PYG{n}{X\PYGZus{}test}\PYG{p}{[}\PYG{n}{i}\PYG{p}{,} \PYG{p}{:}\PYG{p}{]} \PYG{o}{=} \PYG{n}{np}\PYG{o}{.}\PYG{n}{random}\PYG{o}{.}\PYG{n}{multivariate\PYGZus{}normal}\PYG{p}{(}\PYG{n}{mu0}\PYG{p}{,} \PYG{n}{cov0}\PYG{p}{,} \PYG{n}{size}\PYG{o}{=}\PYG{l+m+mi}{1}\PYG{p}{)}
\end{sphinxVerbatim}

\end{sphinxuseclass}\end{sphinxVerbatimInput}

\end{sphinxuseclass}

\section{Metrics and model evaluations methods}
\label{\detokenize{06:metrics-and-model-evaluations-methods}}
\sphinxAtStartPar
In the previous practicals, when we were dealing with regression tasks, we defined different metrics for evaluating the quality of some model’s predictions. Remind: we defined the Mean Squared Error, then, we normalized it to define the Normalized Mean Squared Error. We spoke about other metrics, such as the MAE. All those quantities, while highly informative when applied to continuous quantities, are not as useful when dealing with classes.

\sphinxAtStartPar
Let us consider a practical example. You are assigned the task to determine if some rocket components are working (1) or not (0). In this task, it seems clear that the fact of assigning \sphinxstyleemphasis{working} to a \sphinxstyleemphasis{non\sphinxhyphen{}working} piece is worse than assigning \sphinxstyleemphasis{non\sphinxhyphen{}working} to a \sphinxstyleemphasis{working} piece (you could lose astronauts). However, this is not reflected by the MSE: in both cases, assigning the wrong label would lead to an error of \(1\).

\sphinxAtStartPar
In binary classification, however, we can determine entirely the cases possible after prediction:
\begin{itemize}
\item {} 
\sphinxAtStartPar
Expected: 0, Predicted: 0 \sphinxhyphen{}> True Negative (TN)

\item {} 
\sphinxAtStartPar
Expected: 1, Predicted: 1 \sphinxhyphen{}> True Positive (TP)

\item {} 
\sphinxAtStartPar
Expected: 0, Predicted: 1 \sphinxhyphen{}> False Positive (FP)

\item {} 
\sphinxAtStartPar
Expected: 1, Predicted: 0 \sphinxhyphen{}> False Negative (FN)

\end{itemize}

\sphinxAtStartPar
In this terminology, it is clear that the previous example suffers more from the False Positives than the False Negatives.

\sphinxAtStartPar
We can present those quantities in a matrix, called the \sphinxstyleemphasis{confusion matrix}:


\begin{savenotes}\sphinxattablestart
\sphinxthistablewithglobalstyle
\centering
\begin{tabulary}{\linewidth}[t]{TTT}
\sphinxtoprule

\sphinxAtStartPar

&\sphinxstyletheadfamily 
\sphinxAtStartPar
Real Class: Negative
&\sphinxstyletheadfamily 
\sphinxAtStartPar
Real Class: Positive
\\
\sphinxmidrule
\sphinxtableatstartofbodyhook
\sphinxAtStartPar
Classified as negative
&
\sphinxAtStartPar
TN
&
\sphinxAtStartPar
FN
\\
\sphinxhline
\sphinxAtStartPar
Classified as positive
&
\sphinxAtStartPar
FP
&
\sphinxAtStartPar
TP
\\
\sphinxbottomrule
\end{tabulary}
\sphinxtableafterendhook\par
\sphinxattableend\end{savenotes}

\sphinxAtStartPar
In \sphinxcode{\sphinxupquote{sklearn}}, the function \sphinxcode{\sphinxupquote{confusion\_matrix}} builds it from the set of expected and predicted labels. That is, if a model outputs a score (which is the case for most of them), for a given threshold on those scores the classes are determined, then the confusion matrix is built on top of those predictions. There exists a convenient \sphinxcode{\sphinxupquote{ConfusionMatrixDisplay}} class, but take care, the cases are inverted when compared to the course notation. A confusion matrix, however, is not so\sphinxhyphen{}to\sphinxhyphen{}say a scoring method. There are different values, derived from the confusion matrix quantities, that allows to score a given prediction set:
\begin{itemize}
\item {} 
\sphinxAtStartPar
The \sphinxstylestrong{precision}: \(\frac{\text{TP}}{\text{TP}+\text{FP}}\)

\item {} 
\sphinxAtStartPar
The \sphinxstylestrong{recall}: \(\frac{\text{TP}}{\text{TP}+\text{FN}}\)

\item {} 
\sphinxAtStartPar
The \sphinxstylestrong{F1 score}: \(\frac{2\text{TP}}{2\text{TP}+\text{FP}+\text{FN}}\) which is the harmonic mean between the precision and the recall scores

\item {} 
\sphinxAtStartPar
More generally, the \sphinxstylestrong{F Beta score}:  \(\frac{(1+\beta^2)\text{TP}}{(1+\beta^2)\text{TP}+\text{FP}+\beta^2\text{FN}}\)

\item {} 
\sphinxAtStartPar
The \sphinxstylestrong{accuracy}: \(\frac{\text{TP}+\text{TN}}{\text{TP}+\text{TN}+\text{FP}+\text{FN}}\)

\end{itemize}

\sphinxAtStartPar
However, those quantities often do not take into account the number of True Negatives. Note: most of those scoring methods are implemented in \sphinxcode{\sphinxupquote{sklearn}}.

\sphinxAtStartPar
Let us define furthermore the \sphinxstyleemphasis{True Positive Rate}, i.e., the proportion of True Positive on all the indeed positive samples, and, conversely, the  \sphinxstyleemphasis{False Positive Rate}, the proportion of False Positive on all the in fact positive samples.


\section{Bayes classifier}
\label{\detokenize{06:bayes-classifier}}
\sphinxAtStartPar
Remind: for discrete inputs, the Bayes theorem states that:
\begin{equation*}
\begin{split}P(y|x) = \frac{P(x|y)P(y)}{P(x)}\end{split}
\end{equation*}
\sphinxAtStartPar
and this can be used as a classifier. More specifically, for continuous inputs, we have
\begin{equation*}
\begin{split}P(\mathbf{y}=c_k|\mathbf{x}=x) = \frac{p_x(x|\mathbf{y}=c_k)P(\mathbf{y}=c_k)}{\sum_{k=1}^K p_x(x|\mathbf{y}=c_k)P(\mathbf{y}=c_k)}\end{split}
\end{equation*}
\sphinxAtStartPar
Note, it requires to know what is the density function \(p_x\). If we make the hypothesis that those follow a normal distribution (how could we test if this is valid ?). However, in practice, the \(\mu\) and \sphinxcode{\sphinxupquote{cov}} should be estimated from the data., we can use the \sphinxcode{\sphinxupquote{GaussianNB}} class from \sphinxcode{\sphinxupquote{sklearn}}.

\begin{sphinxuseclass}{cell}\begin{sphinxVerbatimInput}

\begin{sphinxuseclass}{cell_input}
\begin{sphinxVerbatim}[commandchars=\\\{\}]
\PYG{k+kn}{from}\PYG{+w}{ }\PYG{n+nn}{sklearn}\PYG{n+nn}{.}\PYG{n+nn}{naive\PYGZus{}bayes}\PYG{+w}{ }\PYG{k+kn}{import} \PYG{n}{GaussianNB}
\PYG{k+kn}{from}\PYG{+w}{ }\PYG{n+nn}{sklearn}\PYG{n+nn}{.}\PYG{n+nn}{metrics}\PYG{+w}{ }\PYG{k+kn}{import} \PYG{n}{ConfusionMatrixDisplay}
\PYG{k+kn}{from}\PYG{+w}{ }\PYG{n+nn}{sklearn}\PYG{n+nn}{.}\PYG{n+nn}{metrics}\PYG{+w}{ }\PYG{k+kn}{import} \PYG{n}{accuracy\PYGZus{}score}\PYG{p}{,} \PYG{n}{f1\PYGZus{}score}\PYG{p}{,} \PYG{n}{precision\PYGZus{}score}\PYG{p}{,} \PYG{n}{recall\PYGZus{}score}

\PYG{k}{def}\PYG{+w}{ }\PYG{n+nf}{plot\PYGZus{}field}\PYG{p}{(}\PYG{n}{X}\PYG{p}{,} \PYG{n}{Y}\PYG{p}{,} \PYG{n}{model}\PYG{p}{,} \PYG{n}{name}\PYG{o}{=}\PYG{l+s+s2}{\PYGZdq{}}\PYG{l+s+s2}{\PYGZdq{}}\PYG{p}{)}\PYG{p}{:}
    \PYG{n}{x0\PYGZus{}min}\PYG{p}{,} \PYG{n}{x0\PYGZus{}max} \PYG{o}{=} \PYG{n}{X}\PYG{p}{[}\PYG{p}{:}\PYG{p}{,} \PYG{l+m+mi}{0}\PYG{p}{]}\PYG{o}{.}\PYG{n}{min}\PYG{p}{(}\PYG{p}{)} \PYG{o}{\PYGZhy{}} \PYG{l+m+mi}{1}\PYG{p}{,} \PYG{n}{X}\PYG{p}{[}\PYG{p}{:}\PYG{p}{,} \PYG{l+m+mi}{0}\PYG{p}{]}\PYG{o}{.}\PYG{n}{max}\PYG{p}{(}\PYG{p}{)} \PYG{o}{+} \PYG{l+m+mi}{1}
    \PYG{n}{x1\PYGZus{}min}\PYG{p}{,} \PYG{n}{x1\PYGZus{}max} \PYG{o}{=} \PYG{n}{X}\PYG{p}{[}\PYG{p}{:}\PYG{p}{,} \PYG{l+m+mi}{1}\PYG{p}{]}\PYG{o}{.}\PYG{n}{min}\PYG{p}{(}\PYG{p}{)} \PYG{o}{\PYGZhy{}} \PYG{l+m+mi}{1}\PYG{p}{,} \PYG{n}{X}\PYG{p}{[}\PYG{p}{:}\PYG{p}{,} \PYG{l+m+mi}{1}\PYG{p}{]}\PYG{o}{.}\PYG{n}{max}\PYG{p}{(}\PYG{p}{)} \PYG{o}{+} \PYG{l+m+mi}{1}
    \PYG{n}{xx0}\PYG{p}{,} \PYG{n}{xx1} \PYG{o}{=} \PYG{n}{np}\PYG{o}{.}\PYG{n}{meshgrid}\PYG{p}{(}\PYG{n}{np}\PYG{o}{.}\PYG{n}{linspace}\PYG{p}{(}\PYG{n}{x0\PYGZus{}min}\PYG{p}{,} \PYG{n}{x0\PYGZus{}max}\PYG{p}{,} \PYG{l+m+mi}{1000}\PYG{p}{)}\PYG{p}{,}
                           \PYG{n}{np}\PYG{o}{.}\PYG{n}{linspace}\PYG{p}{(}\PYG{n}{x1\PYGZus{}min}\PYG{p}{,} \PYG{n}{x1\PYGZus{}max}\PYG{p}{,} \PYG{l+m+mi}{1000}\PYG{p}{)}\PYG{p}{)}
    
    \PYG{n}{Z} \PYG{o}{=} \PYG{n}{model}\PYG{o}{.}\PYG{n}{predict}\PYG{p}{(}\PYG{n}{np}\PYG{o}{.}\PYG{n}{c\PYGZus{}}\PYG{p}{[}\PYG{n}{xx0}\PYG{o}{.}\PYG{n}{ravel}\PYG{p}{(}\PYG{p}{)}\PYG{p}{,} \PYG{n}{xx1}\PYG{o}{.}\PYG{n}{ravel}\PYG{p}{(}\PYG{p}{)}\PYG{p}{]}\PYG{p}{)}\PYG{o}{.}\PYG{n}{reshape}\PYG{p}{(}\PYG{n}{xx0}\PYG{o}{.}\PYG{n}{shape}\PYG{p}{)}
    \PYG{n}{Y\PYGZus{}hat} \PYG{o}{=} \PYG{n}{model}\PYG{o}{.}\PYG{n}{predict}\PYG{p}{(}\PYG{n}{X}\PYG{p}{)} 
    \PYG{n}{fig}\PYG{p}{,} \PYG{n}{ax} \PYG{o}{=} \PYG{n}{plt}\PYG{o}{.}\PYG{n}{subplots}\PYG{p}{(}\PYG{l+m+mi}{1}\PYG{p}{,} \PYG{l+m+mi}{2}\PYG{p}{,} \PYG{n}{figsize}\PYG{o}{=}\PYG{p}{(}\PYG{l+m+mi}{20}\PYG{p}{,} \PYG{l+m+mi}{8}\PYG{p}{)}\PYG{p}{)}
    \PYG{n}{ax}\PYG{p}{[}\PYG{l+m+mi}{0}\PYG{p}{]}\PYG{o}{.}\PYG{n}{contourf}\PYG{p}{(}\PYG{n}{xx0}\PYG{p}{,} \PYG{n}{xx1}\PYG{p}{,} \PYG{n}{Z}\PYG{p}{,} \PYG{n}{alpha}\PYG{o}{=}\PYG{l+m+mf}{0.25}\PYG{p}{)}
    \PYG{n}{ax}\PYG{p}{[}\PYG{l+m+mi}{0}\PYG{p}{]}\PYG{o}{.}\PYG{n}{scatter}\PYG{p}{(}\PYG{n}{X}\PYG{p}{[}\PYG{p}{:}\PYG{p}{,} \PYG{l+m+mi}{0}\PYG{p}{]}\PYG{p}{,} \PYG{n}{X}\PYG{p}{[}\PYG{p}{:}\PYG{p}{,} \PYG{l+m+mi}{1}\PYG{p}{]}\PYG{p}{,} \PYG{n}{c}\PYG{o}{=}\PYG{n}{Y}\PYG{p}{,} \PYG{n}{marker}\PYG{o}{=}\PYG{l+s+s1}{\PYGZsq{}}\PYG{l+s+s1}{.}\PYG{l+s+s1}{\PYGZsq{}}\PYG{p}{)}
    \PYG{n}{ax}\PYG{p}{[}\PYG{l+m+mi}{0}\PYG{p}{]}\PYG{o}{.}\PYG{n}{grid}\PYG{p}{(}\PYG{l+s+s2}{\PYGZdq{}}\PYG{l+s+s2}{on}\PYG{l+s+s2}{\PYGZdq{}}\PYG{p}{)}
    
    \PYG{n}{ax}\PYG{p}{[}\PYG{l+m+mi}{0}\PYG{p}{]}\PYG{o}{.}\PYG{n}{set\PYGZus{}xlabel}\PYG{p}{(}\PYG{l+s+s2}{\PYGZdq{}}\PYG{l+s+s2}{\PYGZdl{}x\PYGZus{}0\PYGZdl{}}\PYG{l+s+s2}{\PYGZdq{}}\PYG{p}{)}
    \PYG{n}{ax}\PYG{p}{[}\PYG{l+m+mi}{0}\PYG{p}{]}\PYG{o}{.}\PYG{n}{set\PYGZus{}ylabel}\PYG{p}{(}\PYG{l+s+s2}{\PYGZdq{}}\PYG{l+s+s2}{\PYGZdl{}x\PYGZus{}1\PYGZdl{}}\PYG{l+s+s2}{\PYGZdq{}}\PYG{p}{)}
    \PYG{n}{ax}\PYG{p}{[}\PYG{l+m+mi}{0}\PYG{p}{]}\PYG{o}{.}\PYG{n}{set\PYGZus{}title}\PYG{p}{(}\PYG{l+s+sa}{f}\PYG{l+s+s2}{\PYGZdq{}}\PYG{l+s+s2}{Decision Boundary of }\PYG{l+s+si}{\PYGZob{}}\PYG{n}{name}\PYG{l+s+si}{\PYGZcb{}}\PYG{l+s+s2}{\PYGZdq{}}\PYG{p}{)}
    \PYG{n}{ConfusionMatrixDisplay}\PYG{o}{.}\PYG{n}{from\PYGZus{}predictions}\PYG{p}{(}\PYG{n}{Y}\PYG{p}{,} \PYG{n}{Y\PYGZus{}hat}\PYG{p}{,} \PYG{n}{ax}\PYG{o}{=}\PYG{n}{ax}\PYG{p}{[}\PYG{l+m+mi}{1}\PYG{p}{]}\PYG{p}{)}
    \PYG{n}{ax}\PYG{p}{[}\PYG{l+m+mi}{1}\PYG{p}{]}\PYG{o}{.}\PYG{n}{set\PYGZus{}title}\PYG{p}{(}\PYG{l+s+sa}{f}\PYG{l+s+s2}{\PYGZdq{}}\PYG{l+s+s2}{Confusion Matrix for }\PYG{l+s+si}{\PYGZob{}}\PYG{n}{name}\PYG{l+s+si}{\PYGZcb{}}\PYG{l+s+se}{\PYGZbs{}n}\PYG{l+s+s2}{Accuracy: }\PYG{l+s+si}{\PYGZob{}}\PYG{n}{accuracy\PYGZus{}score}\PYG{p}{(}\PYG{n}{Y}\PYG{p}{,}\PYG{+w}{ }\PYG{n}{Y\PYGZus{}hat}\PYG{p}{)}\PYG{l+s+si}{:}\PYG{l+s+s2}{.5f}\PYG{l+s+si}{\PYGZcb{}}\PYG{l+s+s2}{\PYGZdq{}}\PYG{p}{)}
    \PYG{n}{plt}\PYG{o}{.}\PYG{n}{show}\PYG{p}{(}\PYG{p}{)}
    
\PYG{n}{model} \PYG{o}{=} \PYG{n}{GaussianNB}\PYG{p}{(}\PYG{p}{)}\PYG{o}{.}\PYG{n}{fit}\PYG{p}{(}\PYG{n}{X}\PYG{p}{,} \PYG{n}{Y}\PYG{o}{.}\PYG{n}{ravel}\PYG{p}{(}\PYG{p}{)}\PYG{p}{)}

\PYG{n}{plot\PYGZus{}field}\PYG{p}{(}\PYG{n}{X\PYGZus{}test}\PYG{p}{,} \PYG{n}{Y\PYGZus{}test}\PYG{p}{,} \PYG{n}{model}\PYG{p}{,} \PYG{l+s+s2}{\PYGZdq{}}\PYG{l+s+s2}{Gaussian Naive Bayes}\PYG{l+s+s2}{\PYGZdq{}}\PYG{p}{)}
\end{sphinxVerbatim}

\end{sphinxuseclass}\end{sphinxVerbatimInput}
\begin{sphinxVerbatimOutput}

\begin{sphinxuseclass}{cell_output}
\noindent\sphinxincludegraphics{{b16bdff47a5296219fa9b4839d2d0742710ff8ecea2431f639f291385b81267d}.png}

\end{sphinxuseclass}\end{sphinxVerbatimOutput}

\end{sphinxuseclass}

\section{Linear model for classification}
\label{\detokenize{06:linear-model-for-classification}}
\sphinxAtStartPar
As it has been done for the regression tasks, we can use a linear model for classifying the previously generated points. As an example, in \sphinxcode{\sphinxupquote{sklearn}}, we can use the \sphinxcode{\sphinxupquote{LogisticRegression}} class:

\begin{sphinxuseclass}{cell}\begin{sphinxVerbatimInput}

\begin{sphinxuseclass}{cell_input}
\begin{sphinxVerbatim}[commandchars=\\\{\}]
\PYG{k+kn}{from}\PYG{+w}{ }\PYG{n+nn}{sklearn}\PYG{n+nn}{.}\PYG{n+nn}{linear\PYGZus{}model}\PYG{+w}{ }\PYG{k+kn}{import} \PYG{n}{LogisticRegression}

\PYG{n}{model} \PYG{o}{=} \PYG{n}{LogisticRegression}\PYG{p}{(}\PYG{p}{)}\PYG{o}{.}\PYG{n}{fit}\PYG{p}{(}\PYG{n}{X}\PYG{p}{,} \PYG{n}{Y}\PYG{o}{.}\PYG{n}{ravel}\PYG{p}{(}\PYG{p}{)}\PYG{p}{)}

\PYG{n}{plot\PYGZus{}field}\PYG{p}{(}\PYG{n}{X\PYGZus{}test}\PYG{p}{,} \PYG{n}{Y\PYGZus{}test}\PYG{p}{,} \PYG{n}{model}\PYG{p}{,} \PYG{l+s+s2}{\PYGZdq{}}\PYG{l+s+s2}{Logistic Regression}\PYG{l+s+s2}{\PYGZdq{}}\PYG{p}{)}
\end{sphinxVerbatim}

\end{sphinxuseclass}\end{sphinxVerbatimInput}
\begin{sphinxVerbatimOutput}

\begin{sphinxuseclass}{cell_output}
\noindent\sphinxincludegraphics{{ea495c0e6279aece620705135e54cb43c4f5d409fe048930c71ae23c8701f840}.png}

\end{sphinxuseclass}\end{sphinxVerbatimOutput}

\end{sphinxuseclass}

\subsection{Multi\sphinxhyphen{}layers perceptron}
\label{\detokenize{06:multi-layers-perceptron}}
\sphinxAtStartPar
Now, for fun, we can train another model to see how the space is then divided:

\begin{sphinxuseclass}{cell}\begin{sphinxVerbatimInput}

\begin{sphinxuseclass}{cell_input}
\begin{sphinxVerbatim}[commandchars=\\\{\}]
\PYG{k+kn}{from}\PYG{+w}{ }\PYG{n+nn}{sklearn}\PYG{n+nn}{.}\PYG{n+nn}{neural\PYGZus{}network}\PYG{+w}{ }\PYG{k+kn}{import} \PYG{n}{MLPClassifier}

\PYG{n}{model} \PYG{o}{=} \PYG{n}{MLPClassifier}\PYG{p}{(}\PYG{n}{max\PYGZus{}iter}\PYG{o}{=}\PYG{l+m+mi}{1000}\PYG{p}{)}\PYG{o}{.}\PYG{n}{fit}\PYG{p}{(}\PYG{n}{X}\PYG{p}{,} \PYG{n}{Y}\PYG{o}{.}\PYG{n}{ravel}\PYG{p}{(}\PYG{p}{)}\PYG{p}{)}

\PYG{n}{plot\PYGZus{}field}\PYG{p}{(}\PYG{n}{X\PYGZus{}test}\PYG{p}{,} \PYG{n}{Y\PYGZus{}test}\PYG{p}{,} \PYG{n}{model}\PYG{p}{,} \PYG{l+s+s2}{\PYGZdq{}}\PYG{l+s+s2}{MLP}\PYG{l+s+s2}{\PYGZdq{}}\PYG{p}{)}
\end{sphinxVerbatim}

\end{sphinxuseclass}\end{sphinxVerbatimInput}
\begin{sphinxVerbatimOutput}

\begin{sphinxuseclass}{cell_output}
\noindent\sphinxincludegraphics{{507eb4a84b1ecec2aa3211541ec86c62d403d8ac8ddf4bdc4a87d57b3f712e73}.png}

\end{sphinxuseclass}\end{sphinxVerbatimOutput}

\end{sphinxuseclass}

\section{Exercises}
\label{\detokenize{06:id1}}

\subsection{KNN}
\label{\detokenize{06:knn}}\begin{itemize}
\item {} 
\sphinxAtStartPar
Implement a KNN classifier (\sphinxcode{\sphinxupquote{sklearn.neighbors.KNeighborsClassifier}}).

\item {} 
\sphinxAtStartPar
Vary \(k\) (e.g., 1, 5, 20) and visualize decision boundaries.

\item {} 
\sphinxAtStartPar
Discuss underfitting/overfitting trade\sphinxhyphen{}offs.

\end{itemize}

\begin{sphinxuseclass}{cell}\begin{sphinxVerbatimInput}

\begin{sphinxuseclass}{cell_input}
\begin{sphinxVerbatim}[commandchars=\\\{\}]
\PYG{k+kn}{from}\PYG{+w}{ }\PYG{n+nn}{sklearn}\PYG{n+nn}{.}\PYG{n+nn}{neighbors}\PYG{+w}{ }\PYG{k+kn}{import} \PYG{n}{KNeighborsClassifier}

\PYG{k}{for} \PYG{n}{k} \PYG{o+ow}{in} \PYG{p}{(}\PYG{l+m+mi}{1}\PYG{p}{,} \PYG{l+m+mi}{5}\PYG{p}{,} \PYG{l+m+mi}{20}\PYG{p}{)}\PYG{p}{:}
    \PYG{n}{model} \PYG{o}{=} \PYG{n}{KNeighborsClassifier}\PYG{p}{(}\PYG{n}{k}\PYG{p}{)}\PYG{o}{.}\PYG{n}{fit}\PYG{p}{(}\PYG{n}{X}\PYG{p}{,} \PYG{n}{Y}\PYG{o}{.}\PYG{n}{ravel}\PYG{p}{(}\PYG{p}{)}\PYG{p}{)}
    \PYG{n}{plot\PYGZus{}field}\PYG{p}{(}\PYG{n}{X\PYGZus{}test}\PYG{p}{,} \PYG{n}{Y\PYGZus{}test}\PYG{p}{,} \PYG{n}{model}\PYG{p}{,} \PYG{l+s+sa}{f}\PYG{l+s+s2}{\PYGZdq{}}\PYG{l+s+si}{\PYGZob{}}\PYG{n}{k}\PYG{l+s+si}{\PYGZcb{}}\PYG{l+s+s2}{\PYGZhy{}NN}\PYG{l+s+s2}{\PYGZdq{}}\PYG{p}{)}
\end{sphinxVerbatim}

\end{sphinxuseclass}\end{sphinxVerbatimInput}
\begin{sphinxVerbatimOutput}

\begin{sphinxuseclass}{cell_output}
\noindent\sphinxincludegraphics{{e02f4188168cdfd22d35377f61d97c8d6b17278b5fe6384874f2b3db9de5223d}.png}

\noindent\sphinxincludegraphics{{212710ea9992795086741ad1ec8e94693a454c738386e207b28dbd9d0ea7c486}.png}

\noindent\sphinxincludegraphics{{784fbea7027fc434e8d91c77f7e8c7169fd39a98f37f9eb04abe20f396a42b8b}.png}

\end{sphinxuseclass}\end{sphinxVerbatimOutput}

\end{sphinxuseclass}

\subsection{Bayes}
\label{\detokenize{06:bayes}}\begin{itemize}
\item {} 
\sphinxAtStartPar
Compute the true Bayes decision boundary using the known densities \(p(x|y)\).

\item {} 
\sphinxAtStartPar
Compare it to the decision boundaries of logistic regression and KNN.

\item {} 
\sphinxAtStartPar
Discuss the optimality of Bayes classification.

\end{itemize}

\begin{sphinxuseclass}{cell}\begin{sphinxVerbatimInput}

\begin{sphinxuseclass}{cell_input}
\begin{sphinxVerbatim}[commandchars=\\\{\}]
\PYG{n}{p\PYGZus{}x\PYGZus{}given\PYGZus{}y0} \PYG{o}{=} \PYG{n}{multivariate\PYGZus{}normal}\PYG{p}{(}\PYG{n}{mean}\PYG{o}{=}\PYG{n}{mu0}\PYG{p}{,} \PYG{n}{cov}\PYG{o}{=}\PYG{n}{cov0}\PYG{p}{)}
\PYG{n}{p\PYGZus{}x\PYGZus{}given\PYGZus{}y1} \PYG{o}{=} \PYG{n}{multivariate\PYGZus{}normal}\PYG{p}{(}\PYG{n}{mean}\PYG{o}{=}\PYG{n}{mu1}\PYG{p}{,} \PYG{n}{cov}\PYG{o}{=}\PYG{n}{cov1}\PYG{p}{)}

\PYG{n}{x1}\PYG{p}{,} \PYG{n}{x2} \PYG{o}{=} \PYG{n}{np}\PYG{o}{.}\PYG{n}{meshgrid}\PYG{p}{(}\PYG{n}{np}\PYG{o}{.}\PYG{n}{linspace}\PYG{p}{(}\PYG{o}{\PYGZhy{}}\PYG{l+m+mi}{2}\PYG{p}{,} \PYG{l+m+mi}{3}\PYG{p}{,} \PYG{l+m+mi}{1000}\PYG{p}{)}\PYG{p}{,} \PYG{n}{np}\PYG{o}{.}\PYG{n}{linspace}\PYG{p}{(}\PYG{o}{\PYGZhy{}}\PYG{l+m+mi}{2}\PYG{p}{,} \PYG{l+m+mi}{2}\PYG{p}{,} \PYG{l+m+mi}{1000}\PYG{p}{)}\PYG{p}{)}
\PYG{n}{grid} \PYG{o}{=} \PYG{n}{np}\PYG{o}{.}\PYG{n}{c\PYGZus{}}\PYG{p}{[}\PYG{n}{x1}\PYG{o}{.}\PYG{n}{ravel}\PYG{p}{(}\PYG{p}{)}\PYG{p}{,} \PYG{n}{x2}\PYG{o}{.}\PYG{n}{ravel}\PYG{p}{(}\PYG{p}{)}\PYG{p}{]}

\PYG{n}{posterior\PYGZus{}1} \PYG{o}{=} \PYG{n}{p1} \PYG{o}{*} \PYG{n}{p\PYGZus{}x\PYGZus{}given\PYGZus{}y1}\PYG{o}{.}\PYG{n}{pdf}\PYG{p}{(}\PYG{n}{grid}\PYG{p}{)}
\PYG{n}{posterior\PYGZus{}0} \PYG{o}{=} \PYG{p}{(}\PYG{l+m+mi}{1} \PYG{o}{\PYGZhy{}} \PYG{n}{p1}\PYG{p}{)} \PYG{o}{*} \PYG{n}{p\PYGZus{}x\PYGZus{}given\PYGZus{}y0}\PYG{o}{.}\PYG{n}{pdf}\PYG{p}{(}\PYG{n}{grid}\PYG{p}{)}

\PYG{n}{bayes\PYGZus{}boundary} \PYG{o}{=} \PYG{p}{(}\PYG{n}{posterior\PYGZus{}1} \PYG{o}{\PYGZgt{}} \PYG{n}{posterior\PYGZus{}0}\PYG{p}{)}\PYG{o}{.}\PYG{n}{reshape}\PYG{p}{(}\PYG{n}{x1}\PYG{o}{.}\PYG{n}{shape}\PYG{p}{)}

\PYG{n}{plt}\PYG{o}{.}\PYG{n}{figure}\PYG{p}{(}\PYG{n}{figsize}\PYG{o}{=}\PYG{p}{(}\PYG{l+m+mi}{10}\PYG{p}{,} \PYG{l+m+mi}{10}\PYG{p}{)}\PYG{p}{)}
\PYG{n}{plt}\PYG{o}{.}\PYG{n}{contourf}\PYG{p}{(}\PYG{n}{x1}\PYG{p}{,} \PYG{n}{x2}\PYG{p}{,} \PYG{n}{bayes\PYGZus{}boundary}\PYG{o}{/}\PYG{l+m+mi}{2}\PYG{p}{,} \PYG{n}{alpha}\PYG{o}{=}\PYG{l+m+mf}{0.25}\PYG{p}{,} \PYG{n}{cmap}\PYG{o}{=}\PYG{l+s+s2}{\PYGZdq{}}\PYG{l+s+s2}{turbo}\PYG{l+s+s2}{\PYGZdq{}}\PYG{p}{)}
\PYG{n}{plt}\PYG{o}{.}\PYG{n}{scatter}\PYG{p}{(}\PYG{n}{X\PYGZus{}test}\PYG{p}{[}\PYG{p}{:}\PYG{p}{,} \PYG{l+m+mi}{0}\PYG{p}{]}\PYG{p}{,} \PYG{n}{X\PYGZus{}test}\PYG{p}{[}\PYG{p}{:}\PYG{p}{,} \PYG{l+m+mi}{1}\PYG{p}{]}\PYG{p}{,} \PYG{n}{c}\PYG{o}{=}\PYG{n}{Y\PYGZus{}test}\PYG{p}{,} \PYG{n}{marker}\PYG{o}{=}\PYG{l+s+s1}{\PYGZsq{}}\PYG{l+s+s1}{.}\PYG{l+s+s1}{\PYGZsq{}}\PYG{p}{,} \PYG{n}{cmap}\PYG{o}{=}\PYG{l+s+s2}{\PYGZdq{}}\PYG{l+s+s2}{turbo}\PYG{l+s+s2}{\PYGZdq{}}\PYG{p}{)}
\PYG{n}{plt}\PYG{o}{.}\PYG{n}{title}\PYG{p}{(}\PYG{l+s+s2}{\PYGZdq{}}\PYG{l+s+s2}{Bayes Decision Boundary}\PYG{l+s+s2}{\PYGZdq{}}\PYG{p}{)}
\PYG{n}{plt}\PYG{o}{.}\PYG{n}{show}\PYG{p}{(}\PYG{p}{)}
\end{sphinxVerbatim}

\end{sphinxuseclass}\end{sphinxVerbatimInput}
\begin{sphinxVerbatimOutput}

\begin{sphinxuseclass}{cell_output}
\noindent\sphinxincludegraphics{{5aa3f43d3e2ef9a313629d100891fca716c527fd98a55d95462dffbfa74a92be}.png}

\end{sphinxuseclass}\end{sphinxVerbatimOutput}

\end{sphinxuseclass}

\subsection{Other models}
\label{\detokenize{06:other-models}}\begin{itemize}
\item {} 
\sphinxAtStartPar
Use other models to classify tests’ \(X\) and \(Y\) and show their decision boundaries.

\item {} 
\sphinxAtStartPar
Compare their performances using the different metrics.

\end{itemize}

\begin{sphinxuseclass}{cell}\begin{sphinxVerbatimInput}

\begin{sphinxuseclass}{cell_input}
\begin{sphinxVerbatim}[commandchars=\\\{\}]
\PYG{k+kn}{from}\PYG{+w}{ }\PYG{n+nn}{sklearn}\PYG{n+nn}{.}\PYG{n+nn}{tree}\PYG{+w}{ }\PYG{k+kn}{import} \PYG{n}{DecisionTreeClassifier}

\PYG{n}{model} \PYG{o}{=} \PYG{n}{DecisionTreeClassifier}\PYG{p}{(}\PYG{p}{)}\PYG{o}{.}\PYG{n}{fit}\PYG{p}{(}\PYG{n}{X}\PYG{p}{,} \PYG{n}{Y}\PYG{o}{.}\PYG{n}{ravel}\PYG{p}{(}\PYG{p}{)}\PYG{p}{)}

\PYG{n}{plot\PYGZus{}field}\PYG{p}{(}\PYG{n}{X\PYGZus{}test}\PYG{p}{,} \PYG{n}{Y\PYGZus{}test}\PYG{p}{,} \PYG{n}{model}\PYG{p}{,} \PYG{l+s+s2}{\PYGZdq{}}\PYG{l+s+s2}{Random Forest}\PYG{l+s+s2}{\PYGZdq{}}\PYG{p}{)}
\end{sphinxVerbatim}

\end{sphinxuseclass}\end{sphinxVerbatimInput}
\begin{sphinxVerbatimOutput}

\begin{sphinxuseclass}{cell_output}
\noindent\sphinxincludegraphics{{8411f71039d80fd8d4446f4a9905f546c7a8c0715aae4764fcd2e126c8201aa4}.png}

\end{sphinxuseclass}\end{sphinxVerbatimOutput}

\end{sphinxuseclass}
\begin{sphinxuseclass}{cell}\begin{sphinxVerbatimInput}

\begin{sphinxuseclass}{cell_input}
\begin{sphinxVerbatim}[commandchars=\\\{\}]
\PYG{k+kn}{from}\PYG{+w}{ }\PYG{n+nn}{xgboost}\PYG{+w}{ }\PYG{k+kn}{import} \PYG{n}{XGBClassifier}

\PYG{n}{model} \PYG{o}{=} \PYG{n}{XGBClassifier}\PYG{p}{(}\PYG{p}{)}\PYG{o}{.}\PYG{n}{fit}\PYG{p}{(}\PYG{n}{X}\PYG{p}{,} \PYG{n}{Y}\PYG{o}{.}\PYG{n}{ravel}\PYG{p}{(}\PYG{p}{)}\PYG{p}{)}

\PYG{n}{plot\PYGZus{}field}\PYG{p}{(}\PYG{n}{X\PYGZus{}test}\PYG{p}{,} \PYG{n}{Y\PYGZus{}test}\PYG{p}{,} \PYG{n}{model}\PYG{p}{,} \PYG{l+s+s2}{\PYGZdq{}}\PYG{l+s+s2}{XGBoost}\PYG{l+s+s2}{\PYGZdq{}}\PYG{p}{)}
\end{sphinxVerbatim}

\end{sphinxuseclass}\end{sphinxVerbatimInput}
\begin{sphinxVerbatimOutput}

\begin{sphinxuseclass}{cell_output}
\noindent\sphinxincludegraphics{{7e511747db9de6f772846e3d313aadd777db47996ec4dfd12b775ce0389bc5ba}.png}

\end{sphinxuseclass}\end{sphinxVerbatimOutput}

\end{sphinxuseclass}
\begin{sphinxuseclass}{cell}\begin{sphinxVerbatimInput}

\begin{sphinxuseclass}{cell_input}
\begin{sphinxVerbatim}[commandchars=\\\{\}]
\PYG{k+kn}{from}\PYG{+w}{ }\PYG{n+nn}{sklearn}\PYG{n+nn}{.}\PYG{n+nn}{neural\PYGZus{}network}\PYG{+w}{ }\PYG{k+kn}{import} \PYG{n}{MLPClassifier}

\PYG{n}{model} \PYG{o}{=} \PYG{n}{MLPClassifier}\PYG{p}{(}\PYG{n}{hidden\PYGZus{}layer\PYGZus{}sizes}\PYG{o}{=}\PYG{p}{(}\PYG{l+m+mi}{8}\PYG{p}{,} \PYG{l+m+mi}{16}\PYG{p}{,} \PYG{l+m+mi}{8}\PYG{p}{)}\PYG{p}{,} \PYG{n}{max\PYGZus{}iter}\PYG{o}{=}\PYG{l+m+mi}{1000}\PYG{p}{,} \PYG{n}{random\PYGZus{}state}\PYG{o}{=}\PYG{l+m+mi}{1}\PYG{p}{)}\PYG{o}{.}\PYG{n}{fit}\PYG{p}{(}\PYG{n}{X}\PYG{p}{,} \PYG{n}{Y}\PYG{o}{.}\PYG{n}{ravel}\PYG{p}{(}\PYG{p}{)}\PYG{p}{)}

\PYG{n}{plot\PYGZus{}field}\PYG{p}{(}\PYG{n}{X\PYGZus{}test}\PYG{p}{,} \PYG{n}{Y\PYGZus{}test}\PYG{p}{,} \PYG{n}{model}\PYG{p}{,} \PYG{l+s+s2}{\PYGZdq{}}\PYG{l+s+s2}{MLP}\PYG{l+s+s2}{\PYGZdq{}}\PYG{p}{)}
\end{sphinxVerbatim}

\end{sphinxuseclass}\end{sphinxVerbatimInput}
\begin{sphinxVerbatimOutput}

\begin{sphinxuseclass}{cell_output}
\noindent\sphinxincludegraphics{{aac973f6abfbf98b2348092959fc9b1fa474a13a3056789bac5689fb39d3d1ad}.png}

\end{sphinxuseclass}\end{sphinxVerbatimOutput}

\end{sphinxuseclass}

\subsection{More classes}
\label{\detokenize{06:more-classes}}\begin{itemize}
\item {} 
\sphinxAtStartPar
Generate a dataset with two features \(x_1\), \(x_2\), with classes in \(\{0, 1, 2, 3\}\). The probability for each class is \([0.2, 0.3, 0.35, 0.15]\).

\item {} 
\sphinxAtStartPar
Use a Logisitc Regression model to classify the points.

\item {} 
\sphinxAtStartPar
Use a RandomForestClassifier and a MLPClassifier to classify the points.

\item {} 
\sphinxAtStartPar
For the three models, show the confusion matrix.

\item {} 
\sphinxAtStartPar
Using the same visualization function as before, show the decision boundaries of the three models: how are those boundaries related to the way those models work?

\end{itemize}

\begin{sphinxuseclass}{cell}\begin{sphinxVerbatimInput}

\begin{sphinxuseclass}{cell_input}
\begin{sphinxVerbatim}[commandchars=\\\{\}]
\PYG{n}{np}\PYG{o}{.}\PYG{n}{random}\PYG{o}{.}\PYG{n}{seed}\PYG{p}{(}\PYG{l+m+mi}{1}\PYG{p}{)}
\PYG{n}{n} \PYG{o}{=} \PYG{l+m+mi}{2}
\PYG{n}{c\PYGZus{}k} \PYG{o}{=} \PYG{l+m+mi}{4}

\PYG{n}{class\PYGZus{}probs} \PYG{o}{=} \PYG{p}{[}\PYG{l+m+mf}{0.2}\PYG{p}{,} \PYG{l+m+mf}{0.3}\PYG{p}{,} \PYG{l+m+mf}{0.35}\PYG{p}{,} \PYG{l+m+mf}{0.15}\PYG{p}{]}
\PYG{n}{Y4} \PYG{o}{=} \PYG{n}{np}\PYG{o}{.}\PYG{n}{random}\PYG{o}{.}\PYG{n}{choice}\PYG{p}{(}\PYG{n}{np}\PYG{o}{.}\PYG{n}{arange}\PYG{p}{(}\PYG{n}{c\PYGZus{}k}\PYG{p}{)}\PYG{p}{,} \PYG{n}{size}\PYG{o}{=}\PYG{n}{N}\PYG{p}{,} \PYG{n}{p}\PYG{o}{=}\PYG{n}{class\PYGZus{}probs}\PYG{p}{)}

\PYG{n}{mus} \PYG{o}{=} \PYG{p}{[}\PYG{n}{np}\PYG{o}{.}\PYG{n}{random}\PYG{o}{.}\PYG{n}{uniform}\PYG{p}{(}\PYG{o}{\PYGZhy{}}\PYG{l+m+mi}{1}\PYG{p}{,} \PYG{l+m+mi}{1}\PYG{p}{,} \PYG{n}{size}\PYG{o}{=}\PYG{n}{n}\PYG{p}{)} \PYG{k}{for} \PYG{n}{\PYGZus{}} \PYG{o+ow}{in} \PYG{n+nb}{range}\PYG{p}{(}\PYG{n}{c\PYGZus{}k}\PYG{p}{)}\PYG{p}{]}
\PYG{n}{covs} \PYG{o}{=} \PYG{p}{[}\PYG{n}{np}\PYG{o}{.}\PYG{n}{random}\PYG{o}{.}\PYG{n}{rand}\PYG{p}{(}\PYG{n}{n}\PYG{p}{,} \PYG{n}{n}\PYG{p}{)} \PYG{k}{for} \PYG{n}{\PYGZus{}} \PYG{o+ow}{in} \PYG{n+nb}{range}\PYG{p}{(}\PYG{n}{c\PYGZus{}k}\PYG{p}{)}\PYG{p}{]}
\PYG{n}{covs} \PYG{o}{=} \PYG{p}{[}\PYG{n}{C}\PYG{o}{.}\PYG{n}{T} \PYG{o}{@} \PYG{n}{C} \PYG{k}{for} \PYG{n}{C} \PYG{o+ow}{in} \PYG{n}{covs}\PYG{p}{]}

\PYG{n}{X4} \PYG{o}{=} \PYG{n}{np}\PYG{o}{.}\PYG{n}{zeros}\PYG{p}{(}\PYG{p}{(}\PYG{n}{N}\PYG{p}{,} \PYG{n}{n}\PYG{p}{)}\PYG{p}{)}
\PYG{k}{for} \PYG{n}{i} \PYG{o+ow}{in} \PYG{n+nb}{range}\PYG{p}{(}\PYG{n}{N}\PYG{p}{)}\PYG{p}{:}
    \PYG{n+nb+bp}{cls} \PYG{o}{=} \PYG{n}{Y4}\PYG{p}{[}\PYG{n}{i}\PYG{p}{]}
    \PYG{n}{X4}\PYG{p}{[}\PYG{n}{i}\PYG{p}{,} \PYG{p}{:}\PYG{p}{]} \PYG{o}{=} \PYG{n}{np}\PYG{o}{.}\PYG{n}{random}\PYG{o}{.}\PYG{n}{multivariate\PYGZus{}normal}\PYG{p}{(}\PYG{n}{mus}\PYG{p}{[}\PYG{n+nb+bp}{cls}\PYG{p}{]}\PYG{p}{,} \PYG{n}{covs}\PYG{p}{[}\PYG{n+nb+bp}{cls}\PYG{p}{]}\PYG{p}{,} \PYG{n}{size}\PYG{o}{=}\PYG{l+m+mi}{1}\PYG{p}{)}

\PYG{n}{Y4\PYGZus{}test} \PYG{o}{=} \PYG{n}{np}\PYG{o}{.}\PYG{n}{random}\PYG{o}{.}\PYG{n}{choice}\PYG{p}{(}\PYG{n}{np}\PYG{o}{.}\PYG{n}{arange}\PYG{p}{(}\PYG{n}{c\PYGZus{}k}\PYG{p}{)}\PYG{p}{,} \PYG{n}{size}\PYG{o}{=}\PYG{n}{N\PYGZus{}test}\PYG{p}{,} \PYG{n}{p}\PYG{o}{=}\PYG{n}{class\PYGZus{}probs}\PYG{p}{)}
\PYG{n}{X4\PYGZus{}test} \PYG{o}{=} \PYG{n}{np}\PYG{o}{.}\PYG{n}{zeros}\PYG{p}{(}\PYG{p}{(}\PYG{n}{N\PYGZus{}test}\PYG{p}{,} \PYG{n}{n}\PYG{p}{)}\PYG{p}{)}
\PYG{k}{for} \PYG{n}{i} \PYG{o+ow}{in} \PYG{n+nb}{range}\PYG{p}{(}\PYG{n}{N\PYGZus{}test}\PYG{p}{)}\PYG{p}{:}
    \PYG{n+nb+bp}{cls} \PYG{o}{=} \PYG{n}{Y4\PYGZus{}test}\PYG{p}{[}\PYG{n}{i}\PYG{p}{]}
    \PYG{n}{X4\PYGZus{}test}\PYG{p}{[}\PYG{n}{i}\PYG{p}{,} \PYG{p}{:}\PYG{p}{]} \PYG{o}{=} \PYG{n}{np}\PYG{o}{.}\PYG{n}{random}\PYG{o}{.}\PYG{n}{multivariate\PYGZus{}normal}\PYG{p}{(}\PYG{n}{mus}\PYG{p}{[}\PYG{n+nb+bp}{cls}\PYG{p}{]}\PYG{p}{,} \PYG{n}{covs}\PYG{p}{[}\PYG{n+nb+bp}{cls}\PYG{p}{]}\PYG{p}{,} \PYG{n}{size}\PYG{o}{=}\PYG{l+m+mi}{1}\PYG{p}{)}
\end{sphinxVerbatim}

\end{sphinxuseclass}\end{sphinxVerbatimInput}

\end{sphinxuseclass}
\begin{sphinxuseclass}{cell}\begin{sphinxVerbatimInput}

\begin{sphinxuseclass}{cell_input}
\begin{sphinxVerbatim}[commandchars=\\\{\}]
\PYG{n}{model} \PYG{o}{=} \PYG{n}{LogisticRegression}\PYG{p}{(}\PYG{p}{)}\PYG{o}{.}\PYG{n}{fit}\PYG{p}{(}\PYG{n}{X4}\PYG{p}{,} \PYG{n}{Y4}\PYG{o}{.}\PYG{n}{ravel}\PYG{p}{(}\PYG{p}{)}\PYG{p}{)}
\PYG{n}{plot\PYGZus{}field}\PYG{p}{(}\PYG{n}{X4\PYGZus{}test}\PYG{p}{,} \PYG{n}{Y4\PYGZus{}test}\PYG{p}{,} \PYG{n}{model}\PYG{p}{,} \PYG{l+s+s2}{\PYGZdq{}}\PYG{l+s+s2}{Logistic Regression}\PYG{l+s+s2}{\PYGZdq{}}\PYG{p}{)}
\end{sphinxVerbatim}

\end{sphinxuseclass}\end{sphinxVerbatimInput}
\begin{sphinxVerbatimOutput}

\begin{sphinxuseclass}{cell_output}
\noindent\sphinxincludegraphics{{8906d62e6059a2959aabdc64a4c7b9bcc6d73b93b1965612379ac961138acfcb}.png}

\end{sphinxuseclass}\end{sphinxVerbatimOutput}

\end{sphinxuseclass}
\begin{sphinxuseclass}{cell}\begin{sphinxVerbatimInput}

\begin{sphinxuseclass}{cell_input}
\begin{sphinxVerbatim}[commandchars=\\\{\}]
\PYG{k+kn}{from}\PYG{+w}{ }\PYG{n+nn}{sklearn}\PYG{n+nn}{.}\PYG{n+nn}{ensemble}\PYG{+w}{ }\PYG{k+kn}{import} \PYG{n}{RandomForestClassifier}
\PYG{n}{model} \PYG{o}{=} \PYG{n}{RandomForestClassifier}\PYG{p}{(}\PYG{p}{)}\PYG{o}{.}\PYG{n}{fit}\PYG{p}{(}\PYG{n}{X4}\PYG{p}{,} \PYG{n}{Y4}\PYG{o}{.}\PYG{n}{ravel}\PYG{p}{(}\PYG{p}{)}\PYG{p}{)}
\PYG{n}{plot\PYGZus{}field}\PYG{p}{(}\PYG{n}{X4\PYGZus{}test}\PYG{p}{,} \PYG{n}{Y4\PYGZus{}test}\PYG{p}{,} \PYG{n}{model}\PYG{p}{,} \PYG{l+s+s2}{\PYGZdq{}}\PYG{l+s+s2}{Random Forest}\PYG{l+s+s2}{\PYGZdq{}}\PYG{p}{)}

\PYG{n}{model} \PYG{o}{=} \PYG{n}{MLPClassifier}\PYG{p}{(}\PYG{n}{max\PYGZus{}iter}\PYG{o}{=}\PYG{l+m+mi}{1000}\PYG{p}{)}\PYG{o}{.}\PYG{n}{fit}\PYG{p}{(}\PYG{n}{X4}\PYG{p}{,} \PYG{n}{Y4}\PYG{o}{.}\PYG{n}{ravel}\PYG{p}{(}\PYG{p}{)}\PYG{p}{)}
\PYG{n}{plot\PYGZus{}field}\PYG{p}{(}\PYG{n}{X4\PYGZus{}test}\PYG{p}{,} \PYG{n}{Y4\PYGZus{}test}\PYG{p}{,} \PYG{n}{model}\PYG{p}{,} \PYG{l+s+s2}{\PYGZdq{}}\PYG{l+s+s2}{MLP}\PYG{l+s+s2}{\PYGZdq{}}\PYG{p}{)}
\end{sphinxVerbatim}

\end{sphinxuseclass}\end{sphinxVerbatimInput}
\begin{sphinxVerbatimOutput}

\begin{sphinxuseclass}{cell_output}
\noindent\sphinxincludegraphics{{f1e29db2576a39d7b38af4106a84734638799a754689d78405ec5861de3b8720}.png}

\noindent\sphinxincludegraphics{{990348dfda53dd87756fd5c1051405bead9b6690f5e07307cc7c68055fbf551a}.png}

\end{sphinxuseclass}\end{sphinxVerbatimOutput}

\end{sphinxuseclass}

\section{ROC curve}
\label{\detokenize{06:roc-curve}}
\sphinxAtStartPar
Let us define two important notions:
\begin{itemize}
\item {} 
\sphinxAtStartPar
The \sphinxstylestrong{sensibility}: \(\frac{\text{TP}}{\text{TP}+\text{FN}}\). On all the indeed positive cases, what is the portion of correctly assigned positive?

\item {} 
\sphinxAtStartPar
The \sphinxstylestrong{specificity}: \(\frac{\text{TN}}{\text{TN}+\text{FP}}\). On all the indeed negative cases, what is the portion of correctly assigned negative?

\end{itemize}

\sphinxAtStartPar
Now, let us use a threshold model. That is, for a given threshold, if the sample is below, classify it as 0, otherwise classify it as 1. For multiple values of this threshold, the \sphinxstylestrong{ROC} curve (for \sphinxstyleemphasis{receiver operating characteristic}) is a plot in which the False Positive Rate is put in abscissa, and the sensibility in ordinate. The \sphinxstylestrong{precision/recall} plot displays the precision vs. the recall. Finally, the \sphinxstylestrong{lift} curve shows the proportion of \sphinxstyleemphasis{assigned positives} among all the points (the \sphinxstyleemphasis{alert} level).

\sphinxAtStartPar
We can define the notion of \sphinxstyleemphasis{area under the ROC curve} (AUROC, often called AUC). In the best case, the area under the ROC curve is \(1\).


\subsection{Exercise}
\label{\detokenize{06:exercise}}\begin{itemize}
\item {} 
\sphinxAtStartPar
For the previous dataset, implement (by hand):
\begin{itemize}
\item {} 
\sphinxAtStartPar
A KNN Classifier

\item {} 
\sphinxAtStartPar
A Naïve Bayes Classifier

\end{itemize}

\item {} 
\sphinxAtStartPar
where both outputs a score

\item {} 
\sphinxAtStartPar
Then, for multiple threshold value, compute the True Positive rate and the False Positive rate of both models

\item {} 
\sphinxAtStartPar
Plot the ROC curve of both models.

\item {} 
\sphinxAtStartPar
Compute the AUROC for both models. How can you use this as a modelquality assessment ?

\end{itemize}

\begin{sphinxuseclass}{cell}\begin{sphinxVerbatimInput}

\begin{sphinxuseclass}{cell_input}
\begin{sphinxVerbatim}[commandchars=\\\{\}]
\PYG{k+kn}{from}\PYG{+w}{ }\PYG{n+nn}{scipy}\PYG{n+nn}{.}\PYG{n+nn}{stats}\PYG{+w}{ }\PYG{k+kn}{import} \PYG{n}{norm}

\PYG{k}{def}\PYG{+w}{ }\PYG{n+nf}{KNN}\PYG{p}{(}\PYG{n}{X}\PYG{p}{,} \PYG{n}{Y}\PYG{p}{,} \PYG{n}{q}\PYG{p}{,} \PYG{n}{k}\PYG{o}{=}\PYG{l+m+mi}{15}\PYG{p}{)}\PYG{p}{:}
    \PYG{n}{N\PYGZus{}local} \PYG{o}{=} \PYG{n}{X}\PYG{o}{.}\PYG{n}{shape}\PYG{p}{[}\PYG{l+m+mi}{0}\PYG{p}{]}
    \PYG{n}{d} \PYG{o}{=} \PYG{n}{np}\PYG{o}{.}\PYG{n}{sqrt}\PYG{p}{(}\PYG{n}{np}\PYG{o}{.}\PYG{n}{sum}\PYG{p}{(}\PYG{p}{(}\PYG{n}{X} \PYG{o}{\PYGZhy{}} \PYG{n}{q}\PYG{p}{)} \PYG{o}{*}\PYG{o}{*} \PYG{l+m+mi}{2}\PYG{p}{,} \PYG{n}{axis}\PYG{o}{=}\PYG{l+m+mi}{1}\PYG{p}{)}\PYG{p}{)}
    \PYG{n}{index} \PYG{o}{=} \PYG{n}{np}\PYG{o}{.}\PYG{n}{argsort}\PYG{p}{(}\PYG{n}{d}\PYG{p}{)}\PYG{p}{[}\PYG{p}{:}\PYG{n}{k}\PYG{p}{]}
    \PYG{k}{return} \PYG{n}{np}\PYG{o}{.}\PYG{n}{count\PYGZus{}nonzero}\PYG{p}{(}\PYG{n}{Y}\PYG{p}{[}\PYG{n}{index}\PYG{p}{]} \PYG{o}{==} \PYG{l+m+mi}{1}\PYG{p}{)} \PYG{o}{/} \PYG{n}{k}

\PYG{k}{def}\PYG{+w}{ }\PYG{n+nf}{NB}\PYG{p}{(}\PYG{n}{X}\PYG{p}{,} \PYG{n}{Y}\PYG{p}{,} \PYG{n}{q}\PYG{p}{)}\PYG{p}{:}
    \PYG{c+c1}{\PYGZsh{} Recall: we make the hypothesis that the the inputs are conditionnally independent}
    \PYG{n}{N} \PYG{o}{=} \PYG{n+nb}{len}\PYG{p}{(}\PYG{n}{Y}\PYG{p}{)}
    \PYG{n}{n} \PYG{o}{=} \PYG{n}{X}\PYG{o}{.}\PYG{n}{shape}\PYG{p}{[}\PYG{l+m+mi}{1}\PYG{p}{]}
    \PYG{n}{Yhat} \PYG{o}{=} \PYG{n}{np}\PYG{o}{.}\PYG{n}{zeros}\PYG{p}{(}\PYG{n}{N}\PYG{p}{)}
    \PYG{n}{I1} \PYG{o}{=} \PYG{n}{np}\PYG{o}{.}\PYG{n}{where}\PYG{p}{(}\PYG{n}{Y} \PYG{o}{==} \PYG{l+m+mi}{1}\PYG{p}{)}\PYG{p}{[}\PYG{l+m+mi}{0}\PYG{p}{]}
    \PYG{n}{I0} \PYG{o}{=} \PYG{n}{np}\PYG{o}{.}\PYG{n}{where}\PYG{p}{(}\PYG{n}{Y} \PYG{o}{==} \PYG{l+m+mi}{0}\PYG{p}{)}\PYG{p}{[}\PYG{l+m+mi}{0}\PYG{p}{]}
    \PYG{n}{p1} \PYG{o}{=} \PYG{n+nb}{len}\PYG{p}{(}\PYG{n}{I1}\PYG{p}{)} \PYG{o}{/} \PYG{n}{N}
    \PYG{n}{p0} \PYG{o}{=} \PYG{l+m+mi}{1} \PYG{o}{\PYGZhy{}} \PYG{n}{p1}
    \PYG{n}{p1x} \PYG{o}{=} \PYG{l+m+mi}{1}
    \PYG{n}{p0x} \PYG{o}{=} \PYG{l+m+mi}{1}
    \PYG{k}{for} \PYG{n}{j} \PYG{o+ow}{in} \PYG{n+nb}{range}\PYG{p}{(}\PYG{n}{n}\PYG{p}{)}\PYG{p}{:}
        \PYG{n}{mean\PYGZus{}I1} \PYG{o}{=} \PYG{n}{np}\PYG{o}{.}\PYG{n}{mean}\PYG{p}{(}\PYG{n}{X}\PYG{p}{[}\PYG{n}{I1}\PYG{p}{,} \PYG{n}{j}\PYG{p}{]}\PYG{p}{)}
        \PYG{n}{sd\PYGZus{}I1} \PYG{o}{=} \PYG{n}{np}\PYG{o}{.}\PYG{n}{std}\PYG{p}{(}\PYG{n}{X}\PYG{p}{[}\PYG{n}{I1}\PYG{p}{,} \PYG{n}{j}\PYG{p}{]}\PYG{p}{,} \PYG{n}{ddof}\PYG{o}{=}\PYG{l+m+mi}{1}\PYG{p}{)}
        \PYG{n}{mean\PYGZus{}I0} \PYG{o}{=} \PYG{n}{np}\PYG{o}{.}\PYG{n}{mean}\PYG{p}{(}\PYG{n}{X}\PYG{p}{[}\PYG{n}{I0}\PYG{p}{,} \PYG{n}{j}\PYG{p}{]}\PYG{p}{)}
        \PYG{n}{sd\PYGZus{}I0} \PYG{o}{=} \PYG{n}{np}\PYG{o}{.}\PYG{n}{std}\PYG{p}{(}\PYG{n}{X}\PYG{p}{[}\PYG{n}{I0}\PYG{p}{,} \PYG{n}{j}\PYG{p}{]}\PYG{p}{,} \PYG{n}{ddof}\PYG{o}{=}\PYG{l+m+mi}{1}\PYG{p}{)}
        \PYG{n}{p1x} \PYG{o}{=} \PYG{n}{p1x} \PYG{o}{*} \PYG{n}{norm}\PYG{o}{.}\PYG{n}{pdf}\PYG{p}{(}\PYG{n}{q}\PYG{p}{[}\PYG{n}{j}\PYG{p}{]}\PYG{p}{,} \PYG{n}{loc}\PYG{o}{=}\PYG{n}{mean\PYGZus{}I1}\PYG{p}{,} \PYG{n}{scale}\PYG{o}{=}\PYG{n}{sd\PYGZus{}I1}\PYG{p}{)}
        \PYG{n}{p0x} \PYG{o}{=} \PYG{n}{p0x} \PYG{o}{*} \PYG{n}{norm}\PYG{o}{.}\PYG{n}{pdf}\PYG{p}{(}\PYG{n}{q}\PYG{p}{[}\PYG{n}{j}\PYG{p}{]}\PYG{p}{,} \PYG{n}{loc}\PYG{o}{=}\PYG{n}{mean\PYGZus{}I0}\PYG{p}{,} \PYG{n}{scale}\PYG{o}{=}\PYG{n}{sd\PYGZus{}I0}\PYG{p}{)}
    \PYG{n}{Yhat} \PYG{o}{=} \PYG{n}{p1x} \PYG{o}{*} \PYG{n}{p1} \PYG{o}{/} \PYG{p}{(}\PYG{n}{p1x} \PYG{o}{*} \PYG{n}{p1} \PYG{o}{+} \PYG{n}{p0x} \PYG{o}{*} \PYG{n}{p0}\PYG{p}{)}
    \PYG{k}{return} \PYG{n}{Yhat}
\end{sphinxVerbatim}

\end{sphinxuseclass}\end{sphinxVerbatimInput}

\end{sphinxuseclass}
\begin{sphinxuseclass}{cell}\begin{sphinxVerbatimInput}

\begin{sphinxuseclass}{cell_input}
\begin{sphinxVerbatim}[commandchars=\\\{\}]
\PYG{n}{plt}\PYG{o}{.}\PYG{n}{figure}\PYG{p}{(}\PYG{n}{figsize}\PYG{o}{=}\PYG{p}{(}\PYG{l+m+mi}{10}\PYG{p}{,} \PYG{l+m+mi}{4}\PYG{p}{)}\PYG{p}{)}
\PYG{k}{for} \PYG{n}{model} \PYG{o+ow}{in} \PYG{p}{[}\PYG{n}{KNN}\PYG{p}{,} \PYG{n}{NB}\PYG{p}{]}\PYG{p}{:}
    \PYG{n}{Y\PYGZus{}hats\PYGZus{}prob} \PYG{o}{=} \PYG{p}{[}\PYG{p}{]}
    \PYG{c+c1}{\PYGZsh{} Leave\PYGZhy{}one\PYGZhy{}out procedure}
    \PYG{k}{for} \PYG{n}{i} \PYG{o+ow}{in} \PYG{n+nb}{range}\PYG{p}{(}\PYG{n}{X}\PYG{o}{.}\PYG{n}{shape}\PYG{p}{[}\PYG{l+m+mi}{0}\PYG{p}{]}\PYG{p}{)}\PYG{p}{:}
        \PYG{n}{X\PYGZus{}train} \PYG{o}{=} \PYG{n}{np}\PYG{o}{.}\PYG{n}{delete}\PYG{p}{(}\PYG{n}{X}\PYG{p}{,} \PYG{n}{i}\PYG{p}{,} \PYG{n}{axis}\PYG{o}{=}\PYG{l+m+mi}{0}\PYG{p}{)}
        \PYG{n}{Y\PYGZus{}train} \PYG{o}{=} \PYG{n}{np}\PYG{o}{.}\PYG{n}{delete}\PYG{p}{(}\PYG{n}{Y}\PYG{p}{,} \PYG{n}{i}\PYG{p}{,} \PYG{n}{axis}\PYG{o}{=}\PYG{l+m+mi}{0}\PYG{p}{)}
        \PYG{n}{Y\PYGZus{}hats\PYGZus{}prob}\PYG{o}{.}\PYG{n}{append}\PYG{p}{(}\PYG{n}{model}\PYG{p}{(}\PYG{n}{X\PYGZus{}train}\PYG{p}{,} \PYG{n}{Y\PYGZus{}train}\PYG{p}{,} \PYG{n}{X}\PYG{p}{[}\PYG{n}{i}\PYG{p}{]}\PYG{p}{)}\PYG{p}{)}
    
    \PYG{n}{Y\PYGZus{}hats\PYGZus{}prob} \PYG{o}{=} \PYG{n}{np}\PYG{o}{.}\PYG{n}{array}\PYG{p}{(}\PYG{n}{Y\PYGZus{}hats\PYGZus{}prob}\PYG{p}{)}

    \PYG{n}{thresholds} \PYG{o}{=} \PYG{n}{np}\PYG{o}{.}\PYG{n}{linspace}\PYG{p}{(}\PYG{l+m+mi}{0}\PYG{p}{,} \PYG{l+m+mi}{1}\PYG{p}{,} \PYG{l+m+mi}{100}\PYG{p}{)}
    \PYG{n}{TPR} \PYG{o}{=} \PYG{p}{[}\PYG{p}{]}
    \PYG{n}{FPR} \PYG{o}{=} \PYG{p}{[}\PYG{p}{]}
    \PYG{k}{for} \PYG{n}{t} \PYG{o+ow}{in} \PYG{n}{thresholds}\PYG{p}{:}
        \PYG{n}{Y\PYGZus{}hats} \PYG{o}{=} \PYG{n}{np}\PYG{o}{.}\PYG{n}{ones\PYGZus{}like}\PYG{p}{(}\PYG{n}{Y}\PYG{p}{)}
        \PYG{n}{Y\PYGZus{}hats}\PYG{p}{[}\PYG{n}{Y\PYGZus{}hats\PYGZus{}prob} \PYG{o}{\PYGZlt{}} \PYG{n}{t}\PYG{p}{]} \PYG{o}{=} \PYG{l+m+mi}{0}
        
        \PYG{n}{TP} \PYG{o}{=} \PYG{n}{np}\PYG{o}{.}\PYG{n}{sum}\PYG{p}{(}\PYG{p}{(}\PYG{n}{Y\PYGZus{}hats} \PYG{o}{==} \PYG{l+m+mi}{1}\PYG{p}{)} \PYG{o}{\PYGZam{}} \PYG{p}{(}\PYG{n}{Y} \PYG{o}{==} \PYG{l+m+mi}{1}\PYG{p}{)}\PYG{p}{)}
        \PYG{n}{TN} \PYG{o}{=} \PYG{n}{np}\PYG{o}{.}\PYG{n}{sum}\PYG{p}{(}\PYG{p}{(}\PYG{n}{Y\PYGZus{}hats} \PYG{o}{==} \PYG{l+m+mi}{0}\PYG{p}{)} \PYG{o}{\PYGZam{}} \PYG{p}{(}\PYG{n}{Y} \PYG{o}{==} \PYG{l+m+mi}{0}\PYG{p}{)}\PYG{p}{)}    
        \PYG{n}{FP} \PYG{o}{=} \PYG{n}{np}\PYG{o}{.}\PYG{n}{sum}\PYG{p}{(}\PYG{p}{(}\PYG{n}{Y\PYGZus{}hats} \PYG{o}{==} \PYG{l+m+mi}{1}\PYG{p}{)} \PYG{o}{\PYGZam{}} \PYG{p}{(}\PYG{n}{Y} \PYG{o}{==} \PYG{l+m+mi}{0}\PYG{p}{)}\PYG{p}{)}    
        \PYG{n}{FN} \PYG{o}{=} \PYG{n}{np}\PYG{o}{.}\PYG{n}{sum}\PYG{p}{(}\PYG{p}{(}\PYG{n}{Y\PYGZus{}hats} \PYG{o}{==} \PYG{l+m+mi}{0}\PYG{p}{)} \PYG{o}{\PYGZam{}} \PYG{p}{(}\PYG{n}{Y} \PYG{o}{==} \PYG{l+m+mi}{1}\PYG{p}{)}\PYG{p}{)}
    
        \PYG{n}{TPR}\PYG{o}{.}\PYG{n}{append}\PYG{p}{(}\PYG{n}{TP} \PYG{o}{/} \PYG{p}{(}\PYG{n}{TP} \PYG{o}{+} \PYG{n}{FN}\PYG{p}{)} \PYG{k}{if} \PYG{p}{(}\PYG{n}{TP} \PYG{o}{+} \PYG{n}{FN}\PYG{p}{)} \PYG{o}{\PYGZgt{}} \PYG{l+m+mi}{0} \PYG{k}{else} \PYG{l+m+mi}{0}\PYG{p}{)}
        \PYG{n}{FPR}\PYG{o}{.}\PYG{n}{append}\PYG{p}{(}\PYG{n}{FP} \PYG{o}{/} \PYG{p}{(}\PYG{n}{FP} \PYG{o}{+} \PYG{n}{TN}\PYG{p}{)} \PYG{k}{if} \PYG{p}{(}\PYG{n}{FP} \PYG{o}{+} \PYG{n}{TN}\PYG{p}{)} \PYG{o}{\PYGZgt{}} \PYG{l+m+mi}{0} \PYG{k}{else} \PYG{l+m+mi}{0}\PYG{p}{)}
    \PYG{n}{AUROC} \PYG{o}{=} \PYG{l+m+mi}{0}
    \PYG{k}{for} \PYG{n}{i} \PYG{o+ow}{in} \PYG{n+nb}{range}\PYG{p}{(}\PYG{n+nb}{len}\PYG{p}{(}\PYG{n}{TPR}\PYG{p}{)}\PYG{o}{\PYGZhy{}}\PYG{l+m+mi}{1}\PYG{p}{)}\PYG{p}{:}
       \PYG{n}{AUROC} \PYG{o}{+}\PYG{o}{=} \PYG{n}{TPR}\PYG{p}{[}\PYG{n}{i}\PYG{o}{+}\PYG{l+m+mi}{1}\PYG{p}{]}\PYG{o}{*}\PYG{p}{(}\PYG{o}{\PYGZhy{}}\PYG{n}{FPR}\PYG{p}{[}\PYG{n}{i}\PYG{o}{+}\PYG{l+m+mi}{1}\PYG{p}{]}\PYG{o}{+}\PYG{n}{FPR}\PYG{p}{[}\PYG{n}{i}\PYG{p}{]}\PYG{p}{)}
    \PYG{n+nb}{print}\PYG{p}{(}\PYG{n}{AUROC}\PYG{p}{)}
    \PYG{n}{plt}\PYG{o}{.}\PYG{n}{scatter}\PYG{p}{(}\PYG{n}{FPR}\PYG{p}{,} \PYG{n}{TPR}\PYG{p}{)}
    
\PYG{n}{plt}\PYG{o}{.}\PYG{n}{plot}\PYG{p}{(}\PYG{p}{[}\PYG{l+m+mi}{0}\PYG{p}{,} \PYG{l+m+mi}{1}\PYG{p}{]}\PYG{p}{,} \PYG{p}{[}\PYG{l+m+mi}{0}\PYG{p}{,} \PYG{l+m+mi}{1}\PYG{p}{]}\PYG{p}{,} \PYG{n}{linestyle}\PYG{o}{=}\PYG{l+s+s2}{\PYGZdq{}}\PYG{l+s+s2}{dashed}\PYG{l+s+s2}{\PYGZdq{}}\PYG{p}{,} \PYG{n}{color}\PYG{o}{=}\PYG{l+s+s2}{\PYGZdq{}}\PYG{l+s+s2}{gray}\PYG{l+s+s2}{\PYGZdq{}}\PYG{p}{)}
\PYG{n}{plt}\PYG{o}{.}\PYG{n}{legend}\PYG{p}{(}\PYG{p}{[}\PYG{l+s+s2}{\PYGZdq{}}\PYG{l+s+s2}{KNN}\PYG{l+s+s2}{\PYGZdq{}}\PYG{p}{,} \PYG{l+s+s2}{\PYGZdq{}}\PYG{l+s+s2}{NB}\PYG{l+s+s2}{\PYGZdq{}}\PYG{p}{]}\PYG{p}{)}
\PYG{n}{plt}\PYG{o}{.}\PYG{n}{xlabel}\PYG{p}{(}\PYG{l+s+s2}{\PYGZdq{}}\PYG{l+s+s2}{FPR}\PYG{l+s+s2}{\PYGZdq{}}\PYG{p}{)}
\PYG{n}{plt}\PYG{o}{.}\PYG{n}{ylabel}\PYG{p}{(}\PYG{l+s+s2}{\PYGZdq{}}\PYG{l+s+s2}{TPR}\PYG{l+s+s2}{\PYGZdq{}}\PYG{p}{)}
\PYG{n}{plt}\PYG{o}{.}\PYG{n}{title}\PYG{p}{(}\PYG{l+s+s2}{\PYGZdq{}}\PYG{l+s+s2}{ROC curve for KNN and NB}\PYG{l+s+s2}{\PYGZdq{}}\PYG{p}{)}
\PYG{n}{plt}\PYG{o}{.}\PYG{n}{grid}\PYG{p}{(}\PYG{l+s+s2}{\PYGZdq{}}\PYG{l+s+s2}{on}\PYG{l+s+s2}{\PYGZdq{}}\PYG{p}{)}
\PYG{n}{plt}\PYG{o}{.}\PYG{n}{show}\PYG{p}{(}\PYG{p}{)}
\end{sphinxVerbatim}

\end{sphinxuseclass}\end{sphinxVerbatimInput}
\begin{sphinxVerbatimOutput}

\begin{sphinxuseclass}{cell_output}
\begin{sphinxVerbatim}[commandchars=\\\{\}]
0.972294562756058
\end{sphinxVerbatim}

\begin{sphinxVerbatim}[commandchars=\\\{\}]
0.7735446279338773
\end{sphinxVerbatim}

\noindent\sphinxincludegraphics{{1b1179e048990209509a1fa302a3d6e4ace994433a6dec6c8ee5866652f79b4d}.png}

\end{sphinxuseclass}\end{sphinxVerbatimOutput}

\end{sphinxuseclass}
\sphinxAtStartPar
For recall, those are a few models for classification that output a score. Try to display their ROC curve as well. See slides 28\sphinxhyphen{}33 from the theoretical course for more information.


\subsection{More models}
\label{\detokenize{06:more-models}}
\sphinxAtStartPar
First, this model compute the mean of the two classes clusters, and then computes the score as a function of the distance between the query point and those two means.

\begin{sphinxuseclass}{cell}\begin{sphinxVerbatimInput}

\begin{sphinxuseclass}{cell_input}
\begin{sphinxVerbatim}[commandchars=\\\{\}]
\PYG{k}{def}\PYG{+w}{ }\PYG{n+nf}{MD}\PYG{p}{(}\PYG{n}{X}\PYG{p}{,} \PYG{n}{Y}\PYG{p}{,} \PYG{n}{q}\PYG{p}{)}\PYG{p}{:}
    \PYG{n}{N\PYGZus{}local} \PYG{o}{=} \PYG{n+nb}{len}\PYG{p}{(}\PYG{n}{Y}\PYG{p}{)}
    \PYG{n}{n} \PYG{o}{=} \PYG{n}{X}\PYG{o}{.}\PYG{n}{shape}\PYG{p}{[}\PYG{l+m+mi}{1}\PYG{p}{]}
    \PYG{n}{I1\PYGZus{}local} \PYG{o}{=} \PYG{n}{np}\PYG{o}{.}\PYG{n}{where}\PYG{p}{(}\PYG{n}{Y} \PYG{o}{==} \PYG{l+m+mi}{1}\PYG{p}{)}\PYG{p}{[}\PYG{l+m+mi}{0}\PYG{p}{]}
    \PYG{n}{I0\PYGZus{}local} \PYG{o}{=} \PYG{n}{np}\PYG{o}{.}\PYG{n}{where}\PYG{p}{(}\PYG{n}{Y} \PYG{o}{==} \PYG{l+m+mi}{0}\PYG{p}{)}\PYG{p}{[}\PYG{l+m+mi}{0}\PYG{p}{]}
    \PYG{n}{mu1} \PYG{o}{=} \PYG{n}{np}\PYG{o}{.}\PYG{n}{mean}\PYG{p}{(}\PYG{n}{X}\PYG{p}{[}\PYG{n}{I1\PYGZus{}local}\PYG{p}{,} \PYG{p}{:}\PYG{p}{]}\PYG{p}{,} \PYG{n}{axis}\PYG{o}{=}\PYG{l+m+mi}{0}\PYG{p}{)}
    \PYG{n}{mu0} \PYG{o}{=} \PYG{n}{np}\PYG{o}{.}\PYG{n}{mean}\PYG{p}{(}\PYG{n}{X}\PYG{p}{[}\PYG{n}{I0\PYGZus{}local}\PYG{p}{,} \PYG{p}{:}\PYG{p}{]}\PYG{p}{,} \PYG{n}{axis}\PYG{o}{=}\PYG{l+m+mi}{0}\PYG{p}{)}
    \PYG{n}{d1} \PYG{o}{=} \PYG{n}{np}\PYG{o}{.}\PYG{n}{sum}\PYG{p}{(}\PYG{p}{(}\PYG{n}{q} \PYG{o}{\PYGZhy{}} \PYG{n}{mu1}\PYG{p}{)} \PYG{o}{*}\PYG{o}{*} \PYG{l+m+mi}{2}\PYG{p}{)}
    \PYG{n}{d0} \PYG{o}{=} \PYG{n}{np}\PYG{o}{.}\PYG{n}{sum}\PYG{p}{(}\PYG{p}{(}\PYG{n}{q} \PYG{o}{\PYGZhy{}} \PYG{n}{mu0}\PYG{p}{)} \PYG{o}{*}\PYG{o}{*} \PYG{l+m+mi}{2}\PYG{p}{)}
    \PYG{k}{return} \PYG{n}{np}\PYG{o}{.}\PYG{n}{exp}\PYG{p}{(}\PYG{o}{\PYGZhy{}}\PYG{n}{d1}\PYG{p}{)} \PYG{o}{/} \PYG{p}{(}\PYG{n}{np}\PYG{o}{.}\PYG{n}{exp}\PYG{p}{(}\PYG{o}{\PYGZhy{}}\PYG{n}{d1}\PYG{p}{)} \PYG{o}{+} \PYG{n}{np}\PYG{o}{.}\PYG{n}{exp}\PYG{p}{(}\PYG{o}{\PYGZhy{}}\PYG{n}{d0}\PYG{p}{)}\PYG{p}{)}
\end{sphinxVerbatim}

\end{sphinxuseclass}\end{sphinxVerbatimInput}

\end{sphinxuseclass}
\begin{sphinxuseclass}{cell}\begin{sphinxVerbatimInput}

\begin{sphinxuseclass}{cell_input}
\begin{sphinxVerbatim}[commandchars=\\\{\}]
\PYG{k}{def}\PYG{+w}{ }\PYG{n+nf}{QDA}\PYG{p}{(}\PYG{n}{Xtr}\PYG{p}{,} \PYG{n}{Ytr}\PYG{p}{,} \PYG{n}{q}\PYG{p}{)}\PYG{p}{:}
    \PYG{n}{n} \PYG{o}{=} \PYG{n}{Xtr}\PYG{o}{.}\PYG{n}{shape}\PYG{p}{[}\PYG{l+m+mi}{1}\PYG{p}{]}
    \PYG{n}{Ntr} \PYG{o}{=} \PYG{n}{Xtr}\PYG{o}{.}\PYG{n}{shape}\PYG{p}{[}\PYG{l+m+mi}{0}\PYG{p}{]}
    \PYG{n}{I0\PYGZus{}local} \PYG{o}{=} \PYG{n}{np}\PYG{o}{.}\PYG{n}{where}\PYG{p}{(}\PYG{n}{Ytr} \PYG{o}{==} \PYG{l+m+mi}{0}\PYG{p}{)}\PYG{p}{[}\PYG{l+m+mi}{0}\PYG{p}{]}
    \PYG{n}{I1\PYGZus{}local} \PYG{o}{=} \PYG{n}{np}\PYG{o}{.}\PYG{n}{where}\PYG{p}{(}\PYG{n}{Ytr} \PYG{o}{==} \PYG{l+m+mi}{1}\PYG{p}{)}\PYG{p}{[}\PYG{l+m+mi}{0}\PYG{p}{]}
    \PYG{n}{P0} \PYG{o}{=} \PYG{n+nb}{len}\PYG{p}{(}\PYG{n}{I0\PYGZus{}local}\PYG{p}{)} \PYG{o}{/} \PYG{n}{Ntr}
    \PYG{n}{P1} \PYG{o}{=} \PYG{l+m+mi}{1} \PYG{o}{\PYGZhy{}} \PYG{n}{P0}
    \PYG{n}{mu0} \PYG{o}{=} \PYG{n}{np}\PYG{o}{.}\PYG{n}{mean}\PYG{p}{(}\PYG{n}{Xtr}\PYG{p}{[}\PYG{n}{I0\PYGZus{}local}\PYG{p}{,} \PYG{p}{:}\PYG{p}{]}\PYG{p}{,} \PYG{n}{axis}\PYG{o}{=}\PYG{l+m+mi}{0}\PYG{p}{)}
    \PYG{n}{Sigma0} \PYG{o}{=} \PYG{n}{np}\PYG{o}{.}\PYG{n}{cov}\PYG{p}{(}\PYG{n}{Xtr}\PYG{p}{[}\PYG{n}{I0\PYGZus{}local}\PYG{p}{,} \PYG{p}{:}\PYG{p}{]}\PYG{p}{,} \PYG{n}{rowvar}\PYG{o}{=}\PYG{k+kc}{False}\PYG{p}{)}
    \PYG{n}{mu1} \PYG{o}{=} \PYG{n}{np}\PYG{o}{.}\PYG{n}{mean}\PYG{p}{(}\PYG{n}{Xtr}\PYG{p}{[}\PYG{n}{I1\PYGZus{}local}\PYG{p}{,} \PYG{p}{:}\PYG{p}{]}\PYG{p}{,} \PYG{n}{axis}\PYG{o}{=}\PYG{l+m+mi}{0}\PYG{p}{)}
    \PYG{n}{Sigma1} \PYG{o}{=} \PYG{n}{np}\PYG{o}{.}\PYG{n}{cov}\PYG{p}{(}\PYG{n}{Xtr}\PYG{p}{[}\PYG{n}{I1\PYGZus{}local}\PYG{p}{,} \PYG{p}{:}\PYG{p}{]}\PYG{p}{,} \PYG{n}{rowvar}\PYG{o}{=}\PYG{k+kc}{False}\PYG{p}{)}
    \PYG{n}{xts} \PYG{o}{=} \PYG{n}{np}\PYG{o}{.}\PYG{n}{array}\PYG{p}{(}\PYG{p}{[}\PYG{n}{q}\PYG{p}{]}\PYG{p}{)}
    \PYG{n}{diff0} \PYG{o}{=} \PYG{n}{xts} \PYG{o}{\PYGZhy{}} \PYG{n}{mu0}
    \PYG{n}{inv\PYGZus{}Sigma0} \PYG{o}{=} \PYG{n}{np}\PYG{o}{.}\PYG{n}{linalg}\PYG{o}{.}\PYG{n}{inv}\PYG{p}{(}\PYG{n}{Sigma0}\PYG{p}{)}
    \PYG{n}{quad0} \PYG{o}{=} \PYG{n}{diff0}\PYG{o}{.}\PYG{n}{dot}\PYG{p}{(}\PYG{n}{inv\PYGZus{}Sigma0}\PYG{p}{)}\PYG{o}{.}\PYG{n}{dot}\PYG{p}{(}\PYG{n}{diff0}\PYG{o}{.}\PYG{n}{T}\PYG{p}{)}\PYG{p}{[}\PYG{l+m+mi}{0}\PYG{p}{,} \PYG{l+m+mi}{0}\PYG{p}{]}
    \PYG{n}{g0} \PYG{o}{=} \PYG{o}{\PYGZhy{}}\PYG{l+m+mf}{0.5} \PYG{o}{*} \PYG{n}{quad0} \PYG{o}{\PYGZhy{}} \PYG{n}{n}\PYG{o}{/}\PYG{l+m+mi}{2} \PYG{o}{*} \PYG{n}{np}\PYG{o}{.}\PYG{n}{log}\PYG{p}{(}\PYG{l+m+mi}{2} \PYG{o}{*} \PYG{n}{np}\PYG{o}{.}\PYG{n}{pi}\PYG{p}{)} \PYG{o}{\PYGZhy{}} \PYG{l+m+mf}{0.5} \PYG{o}{*} \PYG{n}{np}\PYG{o}{.}\PYG{n}{log}\PYG{p}{(}\PYG{n}{np}\PYG{o}{.}\PYG{n}{linalg}\PYG{o}{.}\PYG{n}{det}\PYG{p}{(}\PYG{n}{Sigma0}\PYG{p}{)}\PYG{p}{)} \PYG{o}{+} \PYG{n}{np}\PYG{o}{.}\PYG{n}{log}\PYG{p}{(}\PYG{n}{P0}\PYG{p}{)}
    \PYG{n}{diff1} \PYG{o}{=} \PYG{n}{xts} \PYG{o}{\PYGZhy{}} \PYG{n}{mu1}
    \PYG{n}{inv\PYGZus{}Sigma1} \PYG{o}{=} \PYG{n}{np}\PYG{o}{.}\PYG{n}{linalg}\PYG{o}{.}\PYG{n}{inv}\PYG{p}{(}\PYG{n}{Sigma1}\PYG{p}{)}
    \PYG{n}{quad1} \PYG{o}{=} \PYG{n}{diff1}\PYG{o}{.}\PYG{n}{dot}\PYG{p}{(}\PYG{n}{inv\PYGZus{}Sigma1}\PYG{p}{)}\PYG{o}{.}\PYG{n}{dot}\PYG{p}{(}\PYG{n}{diff1}\PYG{o}{.}\PYG{n}{T}\PYG{p}{)}\PYG{p}{[}\PYG{l+m+mi}{0}\PYG{p}{,} \PYG{l+m+mi}{0}\PYG{p}{]}
    \PYG{n}{g1} \PYG{o}{=} \PYG{o}{\PYGZhy{}}\PYG{l+m+mf}{0.5} \PYG{o}{*} \PYG{n}{quad1} \PYG{o}{\PYGZhy{}} \PYG{n}{n}\PYG{o}{/}\PYG{l+m+mi}{2} \PYG{o}{*} \PYG{n}{np}\PYG{o}{.}\PYG{n}{log}\PYG{p}{(}\PYG{l+m+mi}{2} \PYG{o}{*} \PYG{n}{np}\PYG{o}{.}\PYG{n}{pi}\PYG{p}{)} \PYG{o}{\PYGZhy{}} \PYG{l+m+mf}{0.5} \PYG{o}{*} \PYG{n}{np}\PYG{o}{.}\PYG{n}{log}\PYG{p}{(}\PYG{n}{np}\PYG{o}{.}\PYG{n}{linalg}\PYG{o}{.}\PYG{n}{det}\PYG{p}{(}\PYG{n}{Sigma1}\PYG{p}{)}\PYG{p}{)} \PYG{o}{+} \PYG{n}{np}\PYG{o}{.}\PYG{n}{log}\PYG{p}{(}\PYG{n}{P1}\PYG{p}{)}
    \PYG{k}{return} \PYG{n}{np}\PYG{o}{.}\PYG{n}{exp}\PYG{p}{(}\PYG{n}{g1}\PYG{p}{)} \PYG{o}{/} \PYG{p}{(}\PYG{n}{np}\PYG{o}{.}\PYG{n}{exp}\PYG{p}{(}\PYG{n}{g0}\PYG{p}{)} \PYG{o}{+} \PYG{n}{np}\PYG{o}{.}\PYG{n}{exp}\PYG{p}{(}\PYG{n}{g1}\PYG{p}{)}\PYG{p}{)}
\end{sphinxVerbatim}

\end{sphinxuseclass}\end{sphinxVerbatimInput}

\end{sphinxuseclass}
\begin{sphinxuseclass}{cell}\begin{sphinxVerbatimInput}

\begin{sphinxuseclass}{cell_input}
\begin{sphinxVerbatim}[commandchars=\\\{\}]
\PYG{k}{def}\PYG{+w}{ }\PYG{n+nf}{LDA}\PYG{p}{(}\PYG{n}{Xtr}\PYG{p}{,} \PYG{n}{Ytr}\PYG{p}{,} \PYG{n}{q}\PYG{p}{)}\PYG{p}{:}
    \PYG{n}{n} \PYG{o}{=} \PYG{n}{Xtr}\PYG{o}{.}\PYG{n}{shape}\PYG{p}{[}\PYG{l+m+mi}{1}\PYG{p}{]}
    \PYG{n}{Ntr} \PYG{o}{=} \PYG{n}{Xtr}\PYG{o}{.}\PYG{n}{shape}\PYG{p}{[}\PYG{l+m+mi}{0}\PYG{p}{]}
    \PYG{n}{I0\PYGZus{}local} \PYG{o}{=} \PYG{n}{np}\PYG{o}{.}\PYG{n}{where}\PYG{p}{(}\PYG{n}{Ytr} \PYG{o}{==} \PYG{l+m+mi}{0}\PYG{p}{)}\PYG{p}{[}\PYG{l+m+mi}{0}\PYG{p}{]}
    \PYG{n}{I1\PYGZus{}local} \PYG{o}{=} \PYG{n}{np}\PYG{o}{.}\PYG{n}{where}\PYG{p}{(}\PYG{n}{Ytr} \PYG{o}{==} \PYG{l+m+mi}{1}\PYG{p}{)}\PYG{p}{[}\PYG{l+m+mi}{0}\PYG{p}{]}
    \PYG{n}{P0} \PYG{o}{=} \PYG{n+nb}{len}\PYG{p}{(}\PYG{n}{I0\PYGZus{}local}\PYG{p}{)} \PYG{o}{/} \PYG{n}{Ntr}
    \PYG{n}{P1} \PYG{o}{=} \PYG{l+m+mi}{1} \PYG{o}{\PYGZhy{}} \PYG{n}{P0}
    \PYG{n}{sigma2} \PYG{o}{=} \PYG{n}{np}\PYG{o}{.}\PYG{n}{mean}\PYG{p}{(}\PYG{n}{np}\PYG{o}{.}\PYG{n}{var}\PYG{p}{(}\PYG{n}{Xtr}\PYG{p}{,} \PYG{n}{axis}\PYG{o}{=}\PYG{l+m+mi}{0}\PYG{p}{,} \PYG{n}{ddof}\PYG{o}{=}\PYG{l+m+mi}{1}\PYG{p}{)}\PYG{p}{)}
    \PYG{n}{mu0} \PYG{o}{=} \PYG{n}{np}\PYG{o}{.}\PYG{n}{mean}\PYG{p}{(}\PYG{n}{Xtr}\PYG{p}{[}\PYG{n}{I0\PYGZus{}local}\PYG{p}{,} \PYG{p}{:}\PYG{p}{]}\PYG{p}{,} \PYG{n}{axis}\PYG{o}{=}\PYG{l+m+mi}{0}\PYG{p}{)}
    \PYG{n}{mu1} \PYG{o}{=} \PYG{n}{np}\PYG{o}{.}\PYG{n}{mean}\PYG{p}{(}\PYG{n}{Xtr}\PYG{p}{[}\PYG{n}{I1\PYGZus{}local}\PYG{p}{,} \PYG{p}{:}\PYG{p}{]}\PYG{p}{,} \PYG{n}{axis}\PYG{o}{=}\PYG{l+m+mi}{0}\PYG{p}{)}
    \PYG{n}{g0} \PYG{o}{=} \PYG{o}{\PYGZhy{}}\PYG{l+m+mi}{1}\PYG{o}{/}\PYG{p}{(}\PYG{l+m+mi}{2}\PYG{o}{*}\PYG{n}{sigma2}\PYG{p}{)} \PYG{o}{*} \PYG{n}{np}\PYG{o}{.}\PYG{n}{linalg}\PYG{o}{.}\PYG{n}{norm}\PYG{p}{(}\PYG{n}{q}\PYG{o}{\PYGZhy{}}\PYG{n}{mu0}\PYG{p}{)} \PYG{o}{+} \PYG{n}{np}\PYG{o}{.}\PYG{n}{log}\PYG{p}{(}\PYG{n}{P0}\PYG{p}{)}
    \PYG{n}{g1} \PYG{o}{=} \PYG{o}{\PYGZhy{}}\PYG{l+m+mi}{1}\PYG{o}{/}\PYG{p}{(}\PYG{l+m+mi}{2}\PYG{o}{*}\PYG{n}{sigma2}\PYG{p}{)} \PYG{o}{*} \PYG{n}{np}\PYG{o}{.}\PYG{n}{linalg}\PYG{o}{.}\PYG{n}{norm}\PYG{p}{(}\PYG{n}{q}\PYG{o}{\PYGZhy{}}\PYG{n}{mu1}\PYG{p}{)} \PYG{o}{+} \PYG{n}{np}\PYG{o}{.}\PYG{n}{log}\PYG{p}{(}\PYG{n}{P1}\PYG{p}{)}
    \PYG{k}{return} \PYG{n}{np}\PYG{o}{.}\PYG{n}{exp}\PYG{p}{(}\PYG{n}{g1}\PYG{p}{)} \PYG{o}{/} \PYG{p}{(}\PYG{n}{np}\PYG{o}{.}\PYG{n}{exp}\PYG{p}{(}\PYG{n}{g0}\PYG{p}{)} \PYG{o}{+} \PYG{n}{np}\PYG{o}{.}\PYG{n}{exp}\PYG{p}{(}\PYG{n}{g1}\PYG{p}{)}\PYG{p}{)}
\end{sphinxVerbatim}

\end{sphinxuseclass}\end{sphinxVerbatimInput}

\end{sphinxuseclass}

\section{Unequal costs}
\label{\detokenize{06:unequal-costs}}
\sphinxAtStartPar
When the cost associated with False Negatives and False Positive are not equal, the computation cannot be used as is. Let us define two costs: \(L_{FP}\) and \(L_{FN}\). Keep in mind, here, that we consider the two other costs \(L_{TP}, L_{TN} = 0\). The expected cost for each possible decision is \(L(1|x) = P(\mathbf{y}=0|x)L_{FP}\) and \(L(0|x) = P(\mathbf{y}=1|x)L_{FN}\). The two \sphinxstyleemphasis{a posteriori} probabilities \(P(y|x)\) can be estimated from the data. If we have \(P(\mathbf{y}=1|x)\), for a sample \(x\), we want to assign the class that minimizes the expected cost.

\sphinxAtStartPar
We predict
\begin{equation*}
\begin{split}
\hat{y} =
\begin{cases} 
1, & \text{if } L(1|x) < L(0|x) \Leftrightarrow (1 - P(y=1|x)) L_{FP} < P(y=1|x) L_{FN}, \\
0, & \text{otherwise}.
\end{cases}
\end{split}
\end{equation*}
\sphinxAtStartPar
We can rewrite this as
\begin{equation*}
\begin{split}P(\mathbf{y}=1|x) > \frac{L_{FN}}{L_{FP}+L_{FN}}\end{split}
\end{equation*}\begin{itemize}
\item {} 
\sphinxAtStartPar
If \(L_{FP} = L_{FN}\), then \(\frac{L_{FN}}{L_{FP}+L_{FN}} = \frac12\). This is the classic case, where we output the maximum a posteriori class.

\item {} 
\sphinxAtStartPar
If \(L_{FP} > L_{FN}\), then it is better to output \(y=1\), this reduces the number of FP.

\item {} 
\sphinxAtStartPar
Otherwise, it is better to output \(y=0\).

\end{itemize}


\subsection{Exercise}
\label{\detokenize{06:id2}}\begin{itemize}
\item {} 
\sphinxAtStartPar
In the generated data, suppose that having a FN has a cost of \(1\) while having a FP has a costof \(0.2\). What is the threshold and the optimal class ?

\end{itemize}

\sphinxAtStartPar
The threshold is \(\frac{1}{1.2} = 0.833\). The optimal class is therefore the class \(0\).

\sphinxstepscope


\chapter{Conclusion}
\label{\detokenize{conclusion:conclusion}}\label{\detokenize{conclusion::doc}}
\sphinxAtStartPar
Throughout this course, you’ve taken a deep dive into the \sphinxstylestrong{statistical and algorithmic foundations} of machine learning. Starting with probabilistic methods and Monte Carlo simulations, you’ve built up your understanding step by step:
\begin{itemize}
\item {} 
\sphinxAtStartPar
You learned how to simulate data and use probabilistic tools to model uncertainty.

\item {} 
\sphinxAtStartPar
You developed \sphinxstylestrong{linear models} and extended them using regularization and basis transformations.

\item {} 
\sphinxAtStartPar
You explored \sphinxstylestrong{tree\sphinxhyphen{}based models}, pipelines, and saw the importance of preprocessing and feature engineering.

\item {} 
\sphinxAtStartPar
You implemented \sphinxstylestrong{neural networks} for regression tasks and trained them using PyTorch.

\item {} 
\sphinxAtStartPar
You investigated \sphinxstylestrong{ensemble methods}, \sphinxstylestrong{feature selection}, and how to combine multiple models.

\item {} 
\sphinxAtStartPar
Finally, you tackled \sphinxstylestrong{classification problems}, including model evaluation using precision, recall, and ROC curves.

\end{itemize}

\sphinxAtStartPar
What you’ve acquired isn’t just a toolkit of models, it’s a \sphinxstylestrong{methodology}. The real power lies in knowing \sphinxstylestrong{how to ask the right questions} and \sphinxstylestrong{choose the right tools} for the task at hand.


\bigskip\hrule\bigskip



\section{A Practical Checklist: Building a Machine Learning Pipeline}
\label{\detokenize{conclusion:a-practical-checklist-building-a-machine-learning-pipeline}}
\sphinxAtStartPar
Whenever you’re faced with a new problem in the real world, keep these questions in mind. They form a structured approach for building and evaluating any ML model:


\subsection{Understanding the Task}
\label{\detokenize{conclusion:understanding-the-task}}\begin{itemize}
\item {} 
\sphinxAtStartPar
\sphinxstylestrong{Is this a classification or regression problem?}
\begin{itemize}
\item {} 
\sphinxAtStartPar
If it’s classification:
\begin{itemize}
\item {} 
\sphinxAtStartPar
How many classes are there?

\item {} 
\sphinxAtStartPar
Is it a binary or multiclass problem? If multiclass, should we use one\sphinxhyphen{}vs\sphinxhyphen{}all or another strategy?

\end{itemize}

\item {} 
\sphinxAtStartPar
If it’s regression:
\begin{itemize}
\item {} 
\sphinxAtStartPar
Should the output be normalized?

\end{itemize}

\end{itemize}

\end{itemize}


\subsection{Understanding the Inputs}
\label{\detokenize{conclusion:understanding-the-inputs}}\begin{itemize}
\item {} 
\sphinxAtStartPar
\sphinxstylestrong{How many features are there?}

\item {} 
\sphinxAtStartPar
\sphinxstylestrong{Can/should we normalize them?}

\item {} 
\sphinxAtStartPar
\sphinxstylestrong{Are some features categorical?}
\begin{itemize}
\item {} 
\sphinxAtStartPar
If yes, can we encode them (e.g., one\sphinxhyphen{}hot encoding, ordinal encoding)?

\end{itemize}

\item {} 
\sphinxAtStartPar
\sphinxstylestrong{Can we reduce dimensionality?}
\begin{itemize}
\item {} 
\sphinxAtStartPar
How can we detect and discard useless or redundant features?

\end{itemize}

\end{itemize}


\subsection{Exploring the Data}
\label{\detokenize{conclusion:exploring-the-data}}\begin{itemize}
\item {} 
\sphinxAtStartPar
\sphinxstylestrong{What is the size of the dataset?} (number of rows and columns)

\item {} 
\sphinxAtStartPar
\sphinxstylestrong{What are the main characteristics of the data?}
\begin{itemize}
\item {} 
\sphinxAtStartPar
Is there correlation or redundancy between features?

\end{itemize}

\item {} 
\sphinxAtStartPar
\sphinxstylestrong{What does the data look like when plotted?}
\begin{itemize}
\item {} 
\sphinxAtStartPar
What types of visualizations best reveal its structure?

\end{itemize}

\end{itemize}


\subsection{Evaluating Models}
\label{\detokenize{conclusion:evaluating-models}}\begin{itemize}
\item {} 
\sphinxAtStartPar
\sphinxstylestrong{What performance metrics are appropriate?}
\begin{itemize}
\item {} 
\sphinxAtStartPar
Accuracy, MSE, precision, recall, ROC AUC, etc.

\end{itemize}

\item {} 
\sphinxAtStartPar
\sphinxstylestrong{How do we assess model quality?}
\begin{itemize}
\item {} 
\sphinxAtStartPar
Visual inspection of residuals?

\item {} 
\sphinxAtStartPar
Cross\sphinxhyphen{}validation scores?

\end{itemize}

\item {} 
\sphinxAtStartPar
\sphinxstylestrong{What does the bias/variance tradeoff look like for our model?}

\item {} 
\sphinxAtStartPar
\sphinxstylestrong{Which model is most justifiable and interpretable for the task?}

\end{itemize}


\subsection{Thinking Critically}
\label{\detokenize{conclusion:thinking-critically}}\begin{itemize}
\item {} 
\sphinxAtStartPar
\sphinxstylestrong{How do we validate the model’s generalization capacity?}
\begin{itemize}
\item {} 
\sphinxAtStartPar
Do we have a test set or cross\sphinxhyphen{}validation strategy?

\end{itemize}

\item {} 
\sphinxAtStartPar
\sphinxstylestrong{What assumptions are we making by choosing this model?}
\begin{itemize}
\item {} 
\sphinxAtStartPar
Linearity? Independence? Normality?

\end{itemize}

\item {} 
\sphinxAtStartPar
\sphinxstylestrong{Is there a simpler solution that performs just as well?}
\begin{itemize}
\item {} 
\sphinxAtStartPar
Occam’s razor applies here too.

\end{itemize}

\item {} 
\sphinxAtStartPar
\sphinxstylestrong{What is the cost of being wrong?}
\begin{itemize}
\item {} 
\sphinxAtStartPar
Especially important in medical, legal, or financial applications.

\end{itemize}

\end{itemize}


\bigskip\hrule\bigskip



\section{Final Words}
\label{\detokenize{conclusion:final-words}}
\sphinxAtStartPar
Machine learning is not just about applying models: it’s about \sphinxstylestrong{thinking critically}, \sphinxstylestrong{analyzing data}, and \sphinxstylestrong{selecting tools} wisely. This course has given you the statistical backbone to understand what’s happening \sphinxstyleemphasis{under the hood}, and the programming practice to start applying it to real problems.

\sphinxAtStartPar
\sphinxstylestrong{Your next steps}: dive deeper into advanced topics like unsupervised learning, probabilistic graphical models, or deep learning. But always return to the foundations when things get unclear: they’re your compass.

\sphinxAtStartPar
And remember: \sphinxstylestrong{models may change, but methodology stays}.

\sphinxstepscope


\chapter{Bibliography}
\label{\detokenize{bibliography:bibliography}}\label{\detokenize{bibliography::doc}}
\begin{sphinxthebibliography}{FH51}
\bibitem[FH51]{bibliography:id2}
\sphinxAtStartPar
E. Fix and J.L. Hodges. \sphinxstyleemphasis{Discriminatory Analysis: Nonparametric Discrimination: Consistency Properties}. USAF School of Aviation Medicine, 1951. URL: \sphinxurl{https://books.google.be/books?id=4XwytAEACAAJ}.
\bibitem[Ho98]{bibliography:id3}
\sphinxAtStartPar
Tin Kam Ho. The random subspace method for constructing decision forests. \sphinxstyleemphasis{IEEE Transactions on Pattern Analysis and Machine Intelligence}, 20(8):832–844, 1998. \sphinxhref{https://doi.org/10.1109/34.709601}{doi:10.1109/34.709601}.
\end{sphinxthebibliography}







\renewcommand{\indexname}{Index}
\printindex
\end{document}